% Options for packages loaded elsewhere
\PassOptionsToPackage{unicode}{hyperref}
\PassOptionsToPackage{hyphens}{url}
\documentclass[
]{book}
\usepackage{xcolor}
\usepackage[left=2cm,right=2cm,top=2cm,bottom=2cm]{geometry}
\usepackage{amsmath,amssymb}
\setcounter{secnumdepth}{-\maxdimen} % remove section numbering
\usepackage{iftex}
\ifPDFTeX
  \usepackage[T1]{fontenc}
  \usepackage[utf8]{inputenc}
  \usepackage{textcomp} % provide euro and other symbols
\else % if luatex or xetex
  \usepackage{unicode-math} % this also loads fontspec
  \defaultfontfeatures{Scale=MatchLowercase}
  \defaultfontfeatures[\rmfamily]{Ligatures=TeX,Scale=1}
\fi
\usepackage{lmodern}
\ifPDFTeX\else
  % xetex/luatex font selection
\fi
% Use upquote if available, for straight quotes in verbatim environments
\IfFileExists{upquote.sty}{\usepackage{upquote}}{}
\IfFileExists{microtype.sty}{% use microtype if available
  \usepackage[]{microtype}
  \UseMicrotypeSet[protrusion]{basicmath} % disable protrusion for tt fonts
}{}
\makeatletter
\@ifundefined{KOMAClassName}{% if non-KOMA class
  \IfFileExists{parskip.sty}{%
    \usepackage{parskip}
  }{% else
    \setlength{\parindent}{0pt}
    \setlength{\parskip}{6pt plus 2pt minus 1pt}}
}{% if KOMA class
  \KOMAoptions{parskip=half}}
\makeatother
\usepackage{color}
\usepackage{fancyvrb}
\newcommand{\VerbBar}{|}
\newcommand{\VERB}{\Verb[commandchars=\\\{\}]}
\DefineVerbatimEnvironment{Highlighting}{Verbatim}{commandchars=\\\{\}}
% Add ',fontsize=\small' for more characters per line
\usepackage{framed}
\definecolor{shadecolor}{RGB}{248,248,248}
\newenvironment{Shaded}{\begin{snugshade}}{\end{snugshade}}
\newcommand{\AlertTok}[1]{\textcolor[rgb]{0.94,0.16,0.16}{#1}}
\newcommand{\AnnotationTok}[1]{\textcolor[rgb]{0.56,0.35,0.01}{\textbf{\textit{#1}}}}
\newcommand{\AttributeTok}[1]{\textcolor[rgb]{0.13,0.29,0.53}{#1}}
\newcommand{\BaseNTok}[1]{\textcolor[rgb]{0.00,0.00,0.81}{#1}}
\newcommand{\BuiltInTok}[1]{#1}
\newcommand{\CharTok}[1]{\textcolor[rgb]{0.31,0.60,0.02}{#1}}
\newcommand{\CommentTok}[1]{\textcolor[rgb]{0.56,0.35,0.01}{\textit{#1}}}
\newcommand{\CommentVarTok}[1]{\textcolor[rgb]{0.56,0.35,0.01}{\textbf{\textit{#1}}}}
\newcommand{\ConstantTok}[1]{\textcolor[rgb]{0.56,0.35,0.01}{#1}}
\newcommand{\ControlFlowTok}[1]{\textcolor[rgb]{0.13,0.29,0.53}{\textbf{#1}}}
\newcommand{\DataTypeTok}[1]{\textcolor[rgb]{0.13,0.29,0.53}{#1}}
\newcommand{\DecValTok}[1]{\textcolor[rgb]{0.00,0.00,0.81}{#1}}
\newcommand{\DocumentationTok}[1]{\textcolor[rgb]{0.56,0.35,0.01}{\textbf{\textit{#1}}}}
\newcommand{\ErrorTok}[1]{\textcolor[rgb]{0.64,0.00,0.00}{\textbf{#1}}}
\newcommand{\ExtensionTok}[1]{#1}
\newcommand{\FloatTok}[1]{\textcolor[rgb]{0.00,0.00,0.81}{#1}}
\newcommand{\FunctionTok}[1]{\textcolor[rgb]{0.13,0.29,0.53}{\textbf{#1}}}
\newcommand{\ImportTok}[1]{#1}
\newcommand{\InformationTok}[1]{\textcolor[rgb]{0.56,0.35,0.01}{\textbf{\textit{#1}}}}
\newcommand{\KeywordTok}[1]{\textcolor[rgb]{0.13,0.29,0.53}{\textbf{#1}}}
\newcommand{\NormalTok}[1]{#1}
\newcommand{\OperatorTok}[1]{\textcolor[rgb]{0.81,0.36,0.00}{\textbf{#1}}}
\newcommand{\OtherTok}[1]{\textcolor[rgb]{0.56,0.35,0.01}{#1}}
\newcommand{\PreprocessorTok}[1]{\textcolor[rgb]{0.56,0.35,0.01}{\textit{#1}}}
\newcommand{\RegionMarkerTok}[1]{#1}
\newcommand{\SpecialCharTok}[1]{\textcolor[rgb]{0.81,0.36,0.00}{\textbf{#1}}}
\newcommand{\SpecialStringTok}[1]{\textcolor[rgb]{0.31,0.60,0.02}{#1}}
\newcommand{\StringTok}[1]{\textcolor[rgb]{0.31,0.60,0.02}{#1}}
\newcommand{\VariableTok}[1]{\textcolor[rgb]{0.00,0.00,0.00}{#1}}
\newcommand{\VerbatimStringTok}[1]{\textcolor[rgb]{0.31,0.60,0.02}{#1}}
\newcommand{\WarningTok}[1]{\textcolor[rgb]{0.56,0.35,0.01}{\textbf{\textit{#1}}}}
\usepackage{longtable,booktabs,array}
\usepackage{calc} % for calculating minipage widths
% Correct order of tables after \paragraph or \subparagraph
\usepackage{etoolbox}
\makeatletter
\patchcmd\longtable{\par}{\if@noskipsec\mbox{}\fi\par}{}{}
\makeatother
% Allow footnotes in longtable head/foot
\IfFileExists{footnotehyper.sty}{\usepackage{footnotehyper}}{\usepackage{footnote}}
\makesavenoteenv{longtable}
\usepackage{graphicx}
\makeatletter
\newsavebox\pandoc@box
\newcommand*\pandocbounded[1]{% scales image to fit in text height/width
  \sbox\pandoc@box{#1}%
  \Gscale@div\@tempa{\textheight}{\dimexpr\ht\pandoc@box+\dp\pandoc@box\relax}%
  \Gscale@div\@tempb{\linewidth}{\wd\pandoc@box}%
  \ifdim\@tempb\p@<\@tempa\p@\let\@tempa\@tempb\fi% select the smaller of both
  \ifdim\@tempa\p@<\p@\scalebox{\@tempa}{\usebox\pandoc@box}%
  \else\usebox{\pandoc@box}%
  \fi%
}
% Set default figure placement to htbp
\def\fps@figure{htbp}
\makeatother
% definitions for citeproc citations
\NewDocumentCommand\citeproctext{}{}
\NewDocumentCommand\citeproc{mm}{%
  \begingroup\def\citeproctext{#2}\cite{#1}\endgroup}
\makeatletter
 % allow citations to break across lines
 \let\@cite@ofmt\@firstofone
 % avoid brackets around text for \cite:
 \def\@biblabel#1{}
 \def\@cite#1#2{{#1\if@tempswa , #2\fi}}
\makeatother
\newlength{\cslhangindent}
\setlength{\cslhangindent}{1.5em}
\newlength{\csllabelwidth}
\setlength{\csllabelwidth}{3em}
\newenvironment{CSLReferences}[2] % #1 hanging-indent, #2 entry-spacing
 {\begin{list}{}{%
  \setlength{\itemindent}{0pt}
  \setlength{\leftmargin}{0pt}
  \setlength{\parsep}{0pt}
  % turn on hanging indent if param 1 is 1
  \ifodd #1
   \setlength{\leftmargin}{\cslhangindent}
   \setlength{\itemindent}{-1\cslhangindent}
  \fi
  % set entry spacing
  \setlength{\itemsep}{#2\baselineskip}}}
 {\end{list}}
\usepackage{calc}
\newcommand{\CSLBlock}[1]{\hfill\break\parbox[t]{\linewidth}{\strut\ignorespaces#1\strut}}
\newcommand{\CSLLeftMargin}[1]{\parbox[t]{\csllabelwidth}{\strut#1\strut}}
\newcommand{\CSLRightInline}[1]{\parbox[t]{\linewidth - \csllabelwidth}{\strut#1\strut}}
\newcommand{\CSLIndent}[1]{\hspace{\cslhangindent}#1}
\setlength{\emergencystretch}{3em} % prevent overfull lines
\providecommand{\tightlist}{%
  \setlength{\itemsep}{0pt}\setlength{\parskip}{0pt}}
\usepackage{booktabs}
\usepackage{longtable}
\usepackage{array}
\usepackage{multirow}
\usepackage{wrapfig}
\usepackage{float}
\usepackage{colortbl}
\usepackage{pdflscape}
\usepackage{tabu}
\usepackage{threeparttable}
\usepackage{threeparttablex}
\usepackage[normalem]{ulem}
\usepackage{makecell}
\usepackage{xcolor}
\usepackage{bookmark}
\IfFileExists{xurl.sty}{\usepackage{xurl}}{} % add URL line breaks if available
\urlstyle{same}
\hypersetup{
  pdftitle={Design and Analysis of Experiments},
  pdfauthor={Larry Winner},
  hidelinks,
  pdfcreator={LaTeX via pandoc}}

\title{Design and Analysis of Experiments}
\author{Larry Winner}
\date{2026-06-18}

\begin{document}
\frontmatter
\maketitle

\mainmatter
\chapter*{Preface}\label{preface}
\addcontentsline{toc}{chapter}{Preface}

These notes are for STA 4211 at the University of Florida. All examples
are based on published articles, though most data has been generated to
match summary statistics.

Some commonly used notation used throughout the course notes with
respect to random variables and estimators is as follows.

\[E\left\{\bullet\right\} - \mbox{ Expected Value or Expectation} \]

\[V\left\{\bullet\right\} - \mbox{ Variance} \]\\
\[SE\left\{\bullet\right\} - \mbox{ Standard Error (square root of Variance)} \]
\[\mbox{COV}\left\{\bullet,\bullet\right\} - \mbox{ Covariance}\]
\[\hat{V}\left\{\bullet\right\} - \mbox{ Estimated Variance} \]
\[\hat{SE}\left\{\bullet\right\} - \mbox{ Estimated Standard Error} \]
\[\hat{\mbox{COV}}\left\{\bullet, \bullet\right\} - \mbox{ Estimated Covariance}\]

For \textbf{quantiles} of distributions used for constructing Confidence
Intervals or setting up rejection regions for tests, the probability
will be the \textbf{lower tail} probability. Some common examples are
given here.

\[
z_{.025}=-1.96 \qquad  z_{.975}=1.96 \qquad \chi_{.95;1}^2=3.84\qquad \chi_{.95;2}^2=5.99
\]

\begin{verbatim}
## Warning: package 'ggplot2' was built under R version 4.5.3
\end{verbatim}

\begin{verbatim}
## -- Attaching core tidyverse packages ------------------------ tidyverse 2.0.0 --
## v dplyr     1.1.4     v readr     2.1.5
## v forcats   1.0.0     v stringr   1.5.1
## v ggplot2   4.0.3     v tibble    3.3.0
## v lubridate 1.9.4     v tidyr     1.3.1
## v purrr     1.0.4     
## -- Conflicts ------------------------------------------ tidyverse_conflicts() --
## x dplyr::filter() masks stats::filter()
## x dplyr::lag()    masks stats::lag()
## i Use the conflicted package (<http://conflicted.r-lib.org/>) to force all conflicts to become errors
## 
## Attaching package: 'kableExtra'
## 
## 
## The following object is masked from 'package:dplyr':
## 
##     group_rows
\end{verbatim}

\begin{verbatim}
## Warning: package 'effectsize' was built under R version 4.5.3
\end{verbatim}

\chapter{Comparing Two Treatments or
Populations}\label{comparing-two-treatments-or-populations}

In this Chapter, we introduce two designs for comparing two treatments
or populations. These could be measurement devices, drug formulations,
brands of food, or any pair of conditions. The two designs are based on
\textbf{independent} or \textbf{paired} samples. We will develop the
\(t\)-tests for each design in this chapter to compare means. We will
also include methods of comparing population variances for independent
or paired samples.

We begin with linear functions of random variables, then introduce
Student's \(t\)-distribution for normally distributed data. Next, we
cover \(t\)-tests for both paired and independent samples. Then we
consider tests for equal variances for both independent and paired
samples. Finally, we will apply these methods directly and using
\textbf{R} functions.

Note that we will use capital letters for estimators that are random
variables based on random samples, such as \(\overline{Y}\),
\(\overline{D}\), and \(S^2\). We will use lower case letters to
represent estimators from a specific sample, such as \(\overline{y}\),
\(\overline{d}\), and \(s^2\), resulting from an experiment or
observational study, being used to construct a Confidence Interval or
test a hypothesis.

\section{Linear Functions of Random
Variables}\label{linear-functions-of-random-variables}

Let \(Y_1,\ldots,Y_n\) be a sample of random variables with the
following properties.

\[
E\left\{Y_i\right\} = \mu_i  \qquad V\left\{Y_i\right\} = \sigma_i^2  \qquad
\mbox{COV}\left\{Y_i,Y_{i'}\right\}=\sigma_{ii'} \qquad  i=1,\ldots,n \quad i\neq i'
\]

Now, let \(a_1,\ldots,a_n\) be a set of constants (not random
variables), and define \(U=\sum_{i=1}^na_iY_i\). Then, we have the
following results regarding the distribution of \(U\).

\[
E\{U\} = E\left\{\sum_{i=1}^na_iY_i\right\}=\sum_{i=1}^na_i\mu_i  \qquad
V\{U\}=V\left\{\sum_{i=1}^na_iY_i\right\}=\sum_{i=1}^na_i^2\sigma_i^2+
2\sum_{i=1}^{n-1}\sum_{i'=i+1}^na_ia_{i'}\sigma_{ii'}
\]

Note that when \(Y_1,\ldots,Y_n\) are independent, then
\(V\{U\}=V\left\{\sum_{i=1}^na_iY_i\right\}=\sum_{i=1}^na_i^2\sigma_i^2\).
Also, if \(Y_1,\ldots,Y_n\) are normally distributed, then \(U\) is as
well.

\textbf{Sample Mean for a Random Sample from a Normal Distribution}

Let \(Y_1,\ldots,Y_n\) be a random sample from a normal distribution
with mean \(\mu\) and variance \(\sigma^2\). Then we are interested in
the distribution of the sample mean \(\overline{Y}\), which is a special
case of \(U\) described above, with \(a_i=1/n\).

\[
E\left\{Y_i\right\}=\mu \quad V\left\{Y_i\right\}=\sigma^2 \quad
\mbox{COV}\left\{Y_i,Y_{i'}\right\}=0 \qquad
\overline{Y}=\frac{\sum_{i=1}^nY_i}{n}=\sum_{i=1}^n\left(\frac{1}{n}\right)Y_i 
\] \[
 E\left\{\overline{Y}\right\}=
E\left\{\frac{\sum_{i=1}^nY_i}{n}\right\}=
\sum_{i=1}^n\left(\frac{1}{n}\right)E\left\{Y_i\right\}=
\sum_{i=1}^n\left(\frac{1}{n}\right)\mu=
n\left(\frac{1}{n}\right)\mu=\mu
\] \[
V\left\{\overline{Y}\right\}=
V\left\{\sum_{i=1}^n\left(\frac{1}{n}\right)Y_i\right\}=
\] \[
\sum_{i=1}^n\left(\frac{1}{n}\right)^2V\left\{Y_i\right\}+
2\sum_{i=1}^{n-1}\sum_{i'=i+1}^n\left(\frac{1}{n}\right)\left(\frac{1}{n}\right)
\mbox{COV}\left\{Y_i,Y_{i'}\right\}=
\] \[
n\left(\frac{1}{n}\right)^2\sigma^2+2\frac{n(n-1)}{2}\left(\frac{1}{n}\right)^2(0)=
\frac{\sigma^2}{n}  \quad \Rightarrow \quad 
\overline{Y} \sim N\left(\mu,\frac{\sigma^2}{n}\right)
\]

Thus for the sample mean in this setting, it is normally distributed
with mean \(\mu\) and variance \(\sigma^2/n\). Further we say that it's
\textbf{Standard Error} is
\(SE\left\{\overline{Y}\right\}=\sqrt{\sigma^2/n}\).

\textbf{Distribution of Paired Differences and Their Mean Difference}

Suppose we are observing measurements on a sample of experimental units
under two conditions. For example, we may be observing vertical leap on
a sample of basketball players, with each player being measured wearing
two different brands of shoes (two measurements per player). Suppose the
brands have means \(\mu_1\) and \(\mu_2\), their variances are equal
\(\sigma_1^2=\sigma_2^2=\sigma^2\), and the covariance of vertical leaps
within players is \(\sigma_{12}\). The covariance would expected to be
positive, as players will tend to be above or below the mean for both
brands, depending on their natural jumping ability. Suppose there are a
sample of \(n\) players in the experiment. We observe the pairs
\(\left\{\left(Y_{11},Y_{21}\right),\ldots,\left(Y_{1n},Y_{2n}\right)\right\}\)
and we are interested in the difference in the leaps for the players
\(D_i=Y_{1i}-Y_{2i}\). Note that by definition, the within player
covariance is \(\sigma_{12}=\rho\sigma_1\sigma_2\), where \(\rho\) is
the correlation between the within player measurements. Now we derive
the distribution of the within player differences, and the mean
difference, where we assume differences across players are independent.

\[
E\left\{Y_{1i}\right\}=\mu_1 \quad E\left\{Y_{2i}\right\}=\mu_2 \qquad
V\left\{Y_{1i}\right\}=V\left\{Y_{2i}\right\}=\sigma^2 \qquad
\mbox{COV}\left\{Y_{1i},Y_{2i}\right\}=\sigma_{12}=\rho\sigma^2
\] \[
D_i=Y_{1i}-Y_{2i}=(1)Y_{1i} + (-1)Y_{2i} \quad \Rightarrow \quad a_1=1,a_2=-1
\] \[
E\left\{D_i\right\}=E\left\{Y_{1i}-Y_{2i}\right\}=
(1)
E\left\{Y_{1i}\right\}+(-1)E\left\{Y_{2i}\right\}=\mu_1-\mu_2=\mu_D
\] \[
V\left\{D_i\right\}=V\left\{Y_{1i}-Y_{2i}\right\}=
(1)^2V\left\{Y_{1i}\right\}+(-1)^2V\left\{Y_{2i}\right\}+
2(1)(-1)\mbox{COV}\left\{Y_{1i},Y_{2i}\right\}=
\] \[
\sigma_1^2+\sigma_2^2-2\rho\sigma_1\sigma_2=2\sigma^2-2\rho\sigma^2=
2(1-\rho)\sigma^2=\sigma_D^2
\]

Now, we are interested in the mean of the \(n\) differences
\(\overline{D}\). Making use of the results from the distribution of the
sample mean above, we have the following results. Further, we will
assume that measurements are (jointly) normally distributed among
players.

\[
E\left\{\overline{D}\right\}=\mu_D=\mu_1-\mu_2 \qquad 
V\left\{\overline{D}\right\}=\frac{\sigma_D^2}{n}=\frac{2(1-\rho)\sigma^2}{n} \qquad
\overline{D}\sim N\left(\mu_D,\frac{\sigma_D^2}{n}\right)
\]

\section{\texorpdfstring{Student's \(t\)-Distribution for Making
Inference on Population
Means}{Student's t-Distribution for Making Inference on Population Means}}\label{students-t-distribution-for-making-inference-on-population-means}

Now consider the sample variance \(S^2\), assuming independent, normally
distributed data with population mean \(\mu\) and variance \(\sigma^2\).

\[
S^2 = \frac{\sum_{i=1}^n\left(Y_i-\overline{Y}\right)^2}{n-1}  \qquad \qquad
W=\frac{(n-1)S^2}{\sigma^2}\sim \chi_{n-1}^2
\] That is, the scaled sample variance follows a Chi-square distribution
with \(n-1\) degrees of freedom. Further, under the normal model, the
sample mean \(\overline{Y}\) and sample variance \(S^2\) are
independent. We define two random variables that are independent. \(Z\)
follows a standard normal distribution and \(W\) follows a Chi-square
distribution with \(\nu\) degrees of freedom, and the resulting random
variable \(T\) follows Student's \(t\)-distribution with \(\nu\) degrees
of freedom.

\[
Z \sim N(0,1) \qquad W\sim \chi_{\nu}^2 \qquad Z\perp W \quad \Rightarrow \quad
T=\frac{Z}{\sqrt{W/\nu}}\sim t_{\nu}
\]

\textbf{Inference Concerning a Population Mean Based on the Sample Mean
and Variance}

Returning to the random sample from a normal distribution with
population mean \(\mu\) and variance \(\sigma^2\), we have the following
results for the sample mean \(\overline{Y}\) and variance \(S^2\).

\[
\overline{Y} \sim N\left(\mu,\frac{\sigma^2}{n}\right) \quad \Rightarrow \quad
Z=\frac{\overline{Y}-\mu}{\sqrt{\sigma^2/n}} \sim N(0,1)  \qquad
W=\frac{(n-1)S^2}{\sigma^2} \sim \chi_{n-1}^2 \qquad Z\perp W
\] \[
T=\frac{Z}{\sqrt{W/(n-1)}} = 
\frac{\frac{\overline{Y}-\mu}{\sqrt{\sigma^2/n}}}{\sqrt{\frac{(n-1)S^2}{\sigma^2}/(n-1)}}
= \frac{\overline{Y}-\mu}{\sqrt{S^2/n}}=
\frac{\overline{Y}-\mu}{\hat{SE}\left\{\overline{Y}\right\}} \sim t_{n-1}
\]

That is when we divide \(\overline{Y}-\mu\) by its estimated standard
error, we obtain a random variable that follows Student's
\(t\)-distribution with \(n-1\) degrees of freedom. This allows
inference for the population mean, without having to know the population
variance. We get the following result when we take repeated samples of
size \(n\), observing \(\overline{Y}\) and \(S^2\).

\[
1-\alpha=P\left(t_{\alpha/2;n-1} \leq 
\frac{\overline{Y}-\mu}{\hat{SE}\left\{\overline{Y}\right\}} \leq
t_{1-\alpha/2;n-1}
\right)=
\] \[
P\left(\overline{Y}-t_{1-\alpha/2;n-1}\hat{SE}\left\{\overline{Y}\right\}\leq
\mu\leq \overline{Y}-t_{\alpha/2;n-1}\hat{SE}\left\{\overline{Y}\right\}\right)
\]

We can construct a \((1-\alpha)100\%\) Confidence Interval for \(\mu\)
or test whether \(\mu=\mu_0\) from this result based on a particular
sample, with \(\overline{y}\) and \(s^2\), the observed sample mean and
variance.

For the hypothesis test, we denote the Test Statistic as \(TS\) and the
Rejection Region as \(RR\).

\[
(1-\alpha)100\% \mbox{ CI for } \mu:
\left[\overline{y}-t_{1-\alpha/2;n-1}\sqrt{\frac{s^2}{n}},
\overline{y}-t_{\alpha/2;n-1}\sqrt{\frac{s^2}{n}}\right]
\qquad \hat{SE}\left\{\overline{Y}\right\}=\sqrt{\frac{s^2}{n}}
\]

\[
H_0:\mu=\mu_0 \qquad H_A:\mu \neq \mu_0 \qquad 
TS:t^*=\frac{\overline{y}-\mu_0}{\sqrt{\frac{s^2}{n}}}
\] \[
RR:|t^*|\geq t_{1-\alpha/2;n-1} 
\qquad P=2P\left(t_{n-1}\geq|t^*|\right)
\]

\section{Inference for the Population Mean
Difference}\label{inference-for-the-population-mean-difference}

In this section, we use the results from Sections 1.1 and 1.2 to make
inferences for population mean differences. First, the paired difference
case is described, as it is a direct application of the previous case
with differences replacing individual observations. Next we consider
independent samples, first assuming equal variances, then when variances
are unequal in the two conditions.

\subsection{Inference for the Population Mean Difference for Paired
Samples}\label{inference-for-the-population-mean-difference-for-paired-samples}

Note that for the paired data case described above, we have the
following results.

\[
D_i=Y_{1i}-Y_{2i} \quad i=1,\ldots,n \qquad \overline{D}=\frac{1}{n}\sum_{i=1}^nD_i   \qquad  
S_D^2 = \frac{1}{n-1}\sum_{i=1}^n\left(D_i-\overline{D}\right)^2
\] \[
\overline{D} \sim N\left(\mu_D,\frac{\sigma_D^2}{n}\right)  \qquad
\frac{(n-1)S_D^2}{\sigma_D^2}\sim \chi_{n-1}^2  \qquad \overline{D} \perp S_D^2
\quad \Rightarrow \quad \frac{\overline{D}-\mu_D}{\sqrt{S_D^2/n}}\sim t_{n-1}
\]

From this, we can construct Confidence Intervals and test hypotheses for
\(\mu_D\). As we did for the sample mean \(\overline{Y}\), we can write
the estimated standard error for \(\overline{D}\) as
\(\hat{SE}\left\{\overline{D}\right\}=\sqrt{S_D^2/n}\).

\[
(1-\alpha)100\% \mbox{ CI for } \mu_D:
\left[\overline{d}-t_{1-\alpha/2;n-1}\sqrt{\frac{s_d^2}{n}},
\overline{d}-t_{\alpha/2;n-1}\sqrt{\frac{s_d^2}{n}}\right]
\qquad \hat{SE}\left\{\overline{D}\right\}=\sqrt{\frac{s_d^2}{n}}
\]

\[
H_0:\mu_D=\Delta_0 \qquad H_A:\mu_D \neq \Delta_0 \qquad 
TS:t^*=\frac{\overline{d}-\Delta_0}{\sqrt{\frac{s_d^2}{n}}}
\]

\[
RR:|t^*|\geq t_{1-\alpha/2;n-1} 
\qquad P=2P\left(t_{n-1}\geq|t^*|\right)
\]

In practice, when testing for \(\mu_D\), typically \(\Delta_0=0\) as we
are testing for a difference in means.

\subsection{Inference for the Population Mean Difference for Independent
Samples}\label{inference-for-the-population-mean-difference-for-independent-samples}

In this section, we consider inference for \(\mu_1-\mu_2\) based on
independent samples. We consider two cases, equal variances
\(\left(\sigma_1^2=\sigma_2^2=\sigma^2\right)\) and unequal variances
\(\left(\sigma_1^2\neq\sigma_2^2\right)\). In each case we assume the
following data structure.

\[Y_{1i} \sim NID\left(\mu_1,\sigma_1^2\right) \quad i=1,\ldots,n_1 \qquad 
Y_{2i} \sim NID\left(\mu_2,\sigma_2^2\right) \quad i=1,\ldots,n_2
\] \[
\overline{Y}_1 \sim N\left(\mu_1,\frac{\sigma_1^2}{n_1}\right)  \qquad
\overline{Y}_2 \sim N\left(\mu_2,\frac{\sigma_2^2}{n_2}\right) \qquad 
\overline{Y}_1 \perp \overline{Y}_2
\] \[
\overline{Y}_1 - \overline{Y}_2 \sim 
N\left(\mu_1-\mu_2,\frac{\sigma_1^2}{n_1}+\frac{\sigma_2^2}{n_2}\right)
\]

\subsubsection{Equal Variances}\label{equal-variances}

When the population variances are equal
\(\left(\sigma_1^2=\sigma_2^2=\sigma^2\right)\), we obtain the following
sampling distribution for the difference in means.

\[
\overline{Y}_1 - \overline{Y}_2 \sim 
N\left(\mu_1-\mu_2,\sigma^2\left[\frac{1}{n_1}+\frac{1}{n_2}\right]\right)
\] We can use the pooled variance \(S_p^2\) as the estimator for the
common \(\sigma^2\). Note that for this model
\(\overline{Y}_1 - \overline{Y}_2\) and \(S_p^2\) are independent.

\[
S_p^2=\frac{\left(n_1-1\right)S_1^2+\left(n_2-1\right)S_2^2}{n_1+n_2-2} \qquad \qquad
\frac{\left(n_1+n_2-2\right)S_p^2}{\sigma^2} \sim \chi_{n_1+n_2-2}^2
\] Then we can obtain the estimated standard error of
\(\overline{Y}_1 - \overline{Y}_2\), by replacing \(\sigma^2\) with
\(S_p^2\).

\[
\frac{\overline{Y}_1 - \overline{Y}_2}{
\sqrt{S_p^2\left[\frac{1}{n_1}+\frac{1}{n_2}\right]}}=
\frac{\overline{Y}_1 - \overline{Y}_2}{
\hat{SE}\left\{\overline{Y}_1 - \overline{Y}_2\right\}} \sim t_{n_1+n_2-2}
\] From this, we can construct Confidence Intervals and conduct
hypothesis tests for \(\mu_1-\mu_2\).

\[
(1-\alpha)100\% \mbox{ CI for } \mu_1-\mu_2:
\left(\overline{y}_1-\overline{y}_2\right) \pm 
t_{1-\alpha/2;n_1+n_2-2}\sqrt{s_p^2\left[\frac{1}{n_1}+\frac{1}{n_2}\right]}
\] \[
\hat{SE}\left\{\overline{Y}_1-\overline{Y}_2\right\}=
\sqrt{s_p^2\left[\frac{1}{n_1}+\frac{1}{n_2}\right]}
\]

\[
H_0:\mu_1-\mu_2=\Delta_0 \qquad H_A:\mu_1-\mu_2 \neq \Delta_0 \qquad 
TS:t^*=\frac{\left(\overline{y}_1-\overline{y}_2\right)-\Delta_0}{
\sqrt{s_p^2\left[\frac{1}{n_1}+\frac{1}{n_2}\right]}}
\]

\[
RR:|t^*|\geq t_{1-\alpha/2;n_1+n_2-2} 
\qquad P=2P\left(t_{n_1+n_2-2}\geq|t^*|\right)
\]

As with the paired case, in practice, typically \(\Delta_0=0\) as we are
comparing two means.

\subsubsection{Unequal Variances}\label{unequal-variances}

When the population variances are unequal
\(\left(\sigma_1^2\neq\sigma_2^2\right)\), we obtain the following
sampling distribution for the difference in means.

\[
\overline{Y}_1 - \overline{Y}_2 \sim 
N\left(\mu_1-\mu_2,\frac{\sigma^2_1}{n_1}+\frac{\sigma_2^2}{n_2}\right)
\]

We can estimate \(V\left\{\overline{Y}_1 - \overline{Y}_2\right\}\) by
replacing \(\sigma_1^2\) and \(\sigma_2^2\) with \(S_1^2\) and
\(S_2^2\). The problem is that we don't obtain a Chi-Square distribution
for a function of \(S_1^2\) and \(S_2^2\) as we did in the equal
variance case.

\textbf{Welch's \(t\)-test} replaces the population variances with the
sample variances in the estimated standard error of
\(\overline{Y}_1 - \overline{Y}_2\) and then makes an adjustment to the
degrees of freedom for the approximate Chi-Square distribution, matching
the mean and variance of the quantity in the denominator of the
\(T=Z/\sqrt{W/{\nu}}\). We briefly describe the method here.

\[
Z=\frac{\left(\overline{Y}_1 - \overline{Y}_2\right)-\left(\mu_1 - \mu_2\right)}{
\sqrt{\frac{\sigma^2_1}{n_1}+\frac{\sigma_2^2}{n_2}}} \sim N(0,1)
\] We next want to replace the population variances in the denominator
with the sample variances. For \(W\sim \chi_{\nu}^2\) the mean and
variance for \(W\) are \(\nu\) and \(2\nu\), respectively. Now, consider
the ratio we will need in the denominator for our approximate \(T\)
ratio. First, we consider the distributions of \(S_1^2\) and \(S_2^2\).

\[
j=1,2: \quad W_j=\frac{\left(n_j-1\right)S_j^2}{\sigma_j^2} \sim \chi_{n_j-1}^2
\quad\Rightarrow\quad E\left\{W_j\right\}=n_j-1 \quad V\left\{W_j\right\}=
2\left(n_j-1\right)
\] \[
\quad\Rightarrow\quad E\left\{S_j^2\right\}=
\frac{\left(n_j-1\right)\sigma_j^2}{\left(n_j-1\right)}=\sigma_j^2 \qquad
V\left\{S_j^2\right\}=\frac{2\left(n_j-1\right)\sigma_j^4}{\left(n_j-1\right)^2}
=\frac{2\sigma_j^4}{\left(n_j-1\right)}
\]

\[
\frac{W^*}{\nu^*}=
\frac{\frac{S_1^2}{n_1}+\frac{S_2^2}{n_2}}{\frac{\sigma_1^2}{n_1}+\frac{\sigma_2^2}{n_2}} \quad 
\Rightarrow \quad E\left\{\frac{W^*}{\nu^*}\right\}=
\frac{\frac{\sigma^2_1}{n_1}+\frac{\sigma_2^2}{n_2}}{
\frac{\sigma^2_1}{n_1}+\frac{\sigma_2^2}{n_2}} = 1
\]

\[
V\left\{\frac{W^*}{\nu^*}\right\} = 
\left(\frac{1}{\nu^*}\right)^2V\left\{W^*\right\}=
\frac{2}{\nu^*}=
\frac{1}{\left(\frac{\sigma^2_1}{n_1}+\frac{\sigma_2^2}{n_2}\right)^2} \left[
\frac{1}{n_1^2}\frac{2\sigma_1^4}{\left(n_1-1\right)} + 
\frac{1}{n_2^2}\frac{2\sigma_2^4}{\left(n_2-1\right)}
\right]
\] \[
\quad\Rightarrow\quad \nu^*=
\frac{\left(\frac{\sigma^2_1}{n_1}+\frac{\sigma_2^2}{n_2}\right)^2}{
\frac{1}{n_1^2}\frac{\sigma_1^4}{\left(n_1-1\right)} + 
\frac{1}{n_2^2}\frac{\sigma_2^4}{\left(n_2-1\right)}}  =
\frac{\left(\frac{\sigma^2_1}{n_1}+\frac{\sigma_2^2}{n_2}\right)^2}{
\frac{\left(\sigma_1^2/n_1\right)^2}{\left(n_1-1\right)} + 
\frac{\left(\sigma_2^2/n_2\right)^2}{\left(n_2-1\right)}}
\] Replacing the unknown \(\sigma_j^2\) with \(S_j^2\), we get
\textbf{Satterthwaite's} approximation for the degrees of freedom for
Welch's \(t\)-test.

\[
\hat{\nu}^*=
\frac{\left(\frac{S_1^2}{n_1}+\frac{S_2^2}{n_2}\right)^2}{
\frac{\left(S_1^2/n_1\right)^2}{\left(n_1-1\right)} + 
\frac{\left(S_2^2/n_2\right)^2}{\left(n_2-1\right)}} = 
\frac{\left(\sum_{j=1}^2g_jMS_j\right)^2}{\sum_{j=1}^2
\frac{\left(g_jMS_j\right)^2}{\nu_j}}  \quad \mbox{ where } \quad
g_j=\frac{1}{n_j}  \quad MS_j=S_j^2 \quad \nu_j=n_j-1
\]

The approximate \((1-\alpha)100\%\) Confidence Interval and test of
hypothesis for the population mean difference for the unequal variance
case are as follow.

\[
(1-\alpha)100\% \mbox{ CI for } \mu_1-\mu_2:
\left(\overline{y}_1-\overline{y}_2\right) \pm 
t_{1-\alpha/2;\hat{\nu}^*}\sqrt{\frac{s_1^2}{n_1}+\frac{s_2^2}{n_2}}
\] \[
\hat{SE}\left\{\overline{Y}_1-\overline{Y}_2\right\}=
\sqrt{\frac{s_1^2}{n_1}+\frac{s_2^2}{n_2}}
\] \[
H_0:\mu_1-\mu_2=\Delta_0 \qquad H_A:\mu_1-\mu_2 \neq \Delta_0 \qquad 
TS:t^*=\frac{\left(\overline{y}_1-\overline{y}_2\right)-\Delta_0}{
\sqrt{\frac{s_1^2}{n_1}+\frac{s_2^2}{n_2}}}
\] \[
RR:|t^*|\geq t_{1-\alpha/2;\hat{\nu}^*} 
\qquad P=2P\left(t_{\hat{\nu}^*}\geq|t^*|\right)
\]

\section{Inference for the Population Variance
Ratios}\label{inference-for-the-population-variance-ratios}

In this section, we describe tests for testing equal variances for two
conditions for independent and paired samples. We will consider the
independent sample case first, as that is much more widely used and
covered in standard statistical textbooks. Then we will briefly cover a
method that can be used for paired samples.

\subsection{Inferences for Variance Ratios for Independent
Samples}\label{inferences-for-variance-ratios-for-independent-samples}

Returning to the estimated variances for independent samples, we have
the following results for independent random samples from normal
distributions.

\[
j=1,2: \quad S_j^2=\frac{1}{n_j-1}\sum_{i=1}^{n_j}\left(Y_{ji}-\overline{Y}_j\right)^2
\qquad \frac{\left(n_j-1\right)S_j^2}{\sigma_j^2}\sim \chi_{n_j-1}^2 \qquad
S_1^2 \perp S_2^2
\]

When we have two independent Chi-Square random variables with \(\nu_1\),
and \(\nu_2\) degrees of freedom, respectively, we obtain the following
distribution result, which allows us to make inference concerning the
variance ratio.

\[
W_1 \sim \chi_{\nu_1}^2 \qquad W_2 \sim \chi_{\nu_2}^2 \qquad W_1\perp W_2
\quad \Rightarrow \quad F=\frac{W_1/{\nu_1}}{W_2/{\nu_2}}\sim F_{\nu_1,\nu_2}
\]

Where \(F_{\nu_1,\nu_2}\) is the \(F\)-distribution with \(\nu_1\) and
\(\nu_2\) numerator and denominator degrees of freedom, respectively. In
terms of this case of comparing variances for two conditions with
independent samples, we have the following results.

\[
F=\frac{W_1/{\nu_1}}{W_2/{\nu_2}}=\frac{\frac{\left(n_1-1\right)S_1^2}{\sigma_1^2}/
\left(n_1-1\right)}{\frac{\left(n_2-1\right)S_2^2}{\sigma_2^2}/
\left(n_2-1\right)}=\frac{S_1^2/\sigma_1^2}{S_2^2/\sigma_2^2}=
\frac{S_1^2/S_2^2}{\sigma_1^2/\sigma_2^2} \sim F_{n_1-1,n_2-1}
\] This leads to the following results that allow for constructing a
Confidence Interval for \(\sigma_1^2/\sigma_2^2\) and testing whether
that ratio is 1 (equal population variances).

\[
1-\alpha = P\left(F_{\alpha/2;n_1-1,n_2-1} \leq 
\frac{S_1^2/S_2^2}{\sigma_1^2/\sigma_2^2} \leq 
F_{1-\alpha/2;n_1-1,n_2-1}
\right) = 
\] \[
P\left(
\frac{S_1^2/S_2^2}{F_{1-\alpha/2;n_1-1,n_2-1}}\leq \frac{\sigma_1^2}{\sigma_2^2} \leq
\frac{S_1^2/S_2^2}{F_{\alpha/2;n_1-1,n_2-1}}
\right)
\] Thus, a \(100(1-\alpha)\%\) Confidence Interval for
\(\sigma_1^2/\sigma_2^2\) will be obtained as follows.

\[
(1-\alpha)100\% \mbox{ CI for } \frac{\sigma_1^2}{\sigma_2^2}: \qquad
\left[\frac{s_1^2/s_2^2}{F_{1-\alpha/2;n_1-1,n_2-1}},
\frac{s_1^2/s_2^2}{F_{\alpha/2;n_1-1,n_2-1}}\right]
\]

If we want to test for equal variances, we have the following test
procedure.

\[
H_0:\sigma_1^2=\sigma_2^2 \qquad H_A:\sigma_1^2\neq\sigma_2^2 \qquad 
TS: F^*=\frac{s_1^2}{s_2^2}
\] \[
RR: F^*\leq F_{\alpha/2;n_1-1,n_2-1} \cup 
F^*\geq F_{1-\alpha/2;n_1-1,n_2-1} 
\] \[
P=2\mbox{min}\left[P\left(F_{n-1,n-2}\leq F^*\right),P\left(F_{n-1,n-2}\geq F^*\right)\right]
\]

\subsection{Test for Equal Variances with Paired
Samples}\label{test-for-equal-variances-with-paired-samples}

A test for equal variances is given in
(\citeproc{ref-snedecor_1967}{Snedecor and Cochran 1967}), which we
will\\
describe here. The test is based on bivariate normal variables \(Y_1\)
and \(Y_2\) with standard deviations \(\sigma_1\), and \(\sigma_2\), and
correlation \(\rho\). Define the difference \(D\) and the sum \(S\) as
follow, along with their variances and covariance.

\[D=Y_1-Y_2 \qquad  S=Y_1+Y_2\]
\[\sigma_D^2 = V\{Y_1-Y_2\}=V\{Y_1\}+V\{Y_2\}-2\mbox{COV}\{Y_1,Y_2\}=
\sigma_1^2+\sigma_2^2-2\rho_{12}\sigma_1\sigma_2\]

\[\sigma_S^2 = V\{Y_1+Y_2\}=V\{Y_1\}+V\{Y_2\}+2\mbox{COV}\{Y_1,Y_2\}=
\sigma_1^2+\sigma_2^2+2\rho_{12}\sigma_1\sigma_2\]

\[\mbox{COV}\{D,S\}=\mbox{COV}\{Y_1-Y_2,Y_1+Y_2\}=\]
\[\mbox{COV}\{Y_1,Y_1\}+\mbox{COV}\{Y_1,Y_2\}-\mbox{COV}\{Y_2,Y_1\}-
\mbox{COV}\{Y_2,Y_2\}=\sigma_1^2-\sigma_2^2\]

The correlation of \(D\) and \(S\) is the covariance divided by the
square root of the product of their variances.

\[\rho_{DS}=
\frac{\mbox{COV}\{D,S\}}{
\sqrt{\left(\sigma_1^2+\sigma_2^2-2\rho_{12}\sigma_1\sigma_2\right)
\left(\sigma_1^2+\sigma_2^2+2\rho_{12}\sigma_1\sigma_2\right)}}=
\]

\[
\frac{\sigma_1^2-\sigma_2^2}{
\sqrt{\left(\sigma_1^2+\sigma_2^2\right)^2-
4\rho_{12}^2\sigma_1^2\sigma_2^2}}
\]

Dividing the numerator and denominator by \(\sigma_2^2\) and defining
\(\phi=\sigma_1^2/\sigma_2^2\), we obtain the following result for
\(\rho_{DS}\).

\[\rho_{DS}=\frac{\frac{\sigma_1^2}{\sigma_2^2}-1}{
\sqrt{\frac{\left(\sigma_1^2+\sigma_2^2\right)^2}{\left(\sigma_2^2\right)^2}-4\rho_{12}^2\left(\frac{\sigma_1^2}{\sigma_2^2}\right)
\left(\frac{\sigma_2^2}{\sigma_2^2}\right)}}=
\frac{\phi-1}{
\sqrt{(\phi+1)^2-4\rho_{12}^2\phi}
}
\]

Now, letting \(F=s_1^2/s_2^2\), the ratio of the sample variances, we
compute the sample version of \(\rho_{DS}\), labelled \(r_{DS}\).

\[r_{DS}=\frac{F-1}{\sqrt{(F+1)^2-4r_{12}^2F}}\] where \(r_{12}\) is the
sample correlation between \(Y_1\) and \(Y_2\). Then the test of whether
\(\rho_{DS}=0\), which implies \(\sigma_1^2-\sigma_2^2=0\) is conducted
with the following \(t\)-test,which is of the same form as testing that
a Pearson correlation \(\rho=0\).

\[H_0:\sigma_1^2=\sigma_2^2 \qquad H_A:\sigma_1^2 \neq \sigma_2^2 \qquad
\mbox{TS: } t_{obs}=\frac{r_{DS}}{\sqrt{\frac{1-r_{DS}^2}{n-2}}}\]

\[ \mbox{RR: } |t_{obs}|\geq t_{1-\alpha/2,n-2}  \qquad
P=2P\left(t_{n-2}\geq |t_{obs}|\right)  \]

\section{Examples}\label{examples}

In this section we will consider examples of probability distributions
and inferences for differences among population means and variances.

In the first example we will consider an example of a population of two
random variables being observed on the same units, and examine paired
differences.

\textbf{Example 1.1: WNBA Team Actual and Predicted Points - Paired
Data}

This example concerns all team-games for the WNBA 2010-2024 regular
seasons. There were 3059 games played, each with two teams, so there are
6118 units. Closing betting lines provide a point spread and Over/Under
line for each team game. We use the convention used in the sports
betting literature as assigning favored teams a positive spread
(opposite of how it's posted on betting sites).

Suppose the home team is favored by \(HS=+5\) points (and away team
underdog by 5 points), and the Over/Under is \(OU=150\) points. The
betting line Predicted scores for the home team \(HP\) and away team
\(AP\) can be obtained as follows.

\[ HP = \frac{OU+HS}{2}=\frac{150+5}{2}=77.5   \qquad  
AP = \frac{OU-HS}{2}=\frac{150-5}{2}=72.5   \qquad
\] This preserves the sum (Over/Under) and home team spread.

Then we observe the Actual score for each team-game. We are interested
in the difference \(D=A-P\) for each team-game, where \(A\) is the
Actual score (\(Y_1\) in the formulas) and \(P\) is the Predicted score
(\(Y_2\) in the formulas). Note that if \(D\) is positive, the team
scored higher than predicted by the betting lines.

For this example, we will do the following steps.

\begin{itemize}
\tightlist
\item
  Treating this as a population, obtain parameters for \(A\), \(P\), and
  \(D=A-P\) along with Shapiro-Wilk tests for normality and a histogram
  with superimposed normal distribution for \(D=A-P\).
\item
  Treating this as a single sample from a \textbf{metaphorical
  population} where these 6118 games could be replayed many times,
  obtain a 95\% Confidence Interval for \(\mu_D\) and test
  \(H_0:\mu_D=0\), implying that on average, the predictions are
  correct.
\item
  Treating this as a population, take 100000 (pseudo) random samples of
  size \(n=20\) and observe the quantity
  \(t=\left(\overline{d}-\mu_d\right)/\sqrt{s_d^2/n}\) and plot the
  superimposed \(t_{19}\) distribution. Further, obtain the coverage
  rate for the 95\% Confidence Intervals for \(\mu_D\).
\item
  For the 100000 samples described above, conduct the test of equal
  variances and consider the distribution of observed correlations
  \(r_{DS}\)
\end{itemize}

\textbf{Population Parameters, Normality Tests, and Histogram}

Population parameters for \(A\), \(P\), and \(D=A-P\) are given in Table
@ref(tab:chapter1a1). and a histogram with a superimposed normal
distribution for the differences is given in Figure
@ref(fig:chapter1a1).

These differences are approximately normally distributed with mean
\(\mu=0.0814\) and standard deviation \(\sigma=10.0813\) points.

As the \textbf{shapiro.test} function does not report \(P\)-values for
sample sizes above 5000 (as they are considered unreliable for
relatively large samples), we take 500 random samples, each of
\(n=1000\) games, and obtain the average \(P\)-values for those 500
samples. For the Actual points, the average \(P\)-value is .1369, for
Predicted, it is .0584, and for the Difference, it is .4550. Thus, there
is not strong evidence against normalty, particularly for the
Differences.

\begin{verbatim}
##  [1] "team"         "date"         "year"         "tmGmYr"       "opGmYr"      
##  [6] "home"         "teamPts"      "oppPts"       "OT"           "tmSprd"      
## [11] "OU"           "tmSpDev"      "oppTeam"      "tmWin"        "tmATSWin"    
## [16] "tmOver"       "tmWinSeas"    "tmLosSeas"    "tmWinStrk"    "tmLosStrk"   
## [21] "tmATSWinSeas" "tmATSLosSeas" "tmATSWinStrk" "tmATSLosStrk" "tmOVRSeas"   
## [26] "tmUDRSeas"    "tmOVRStrk"    "tmUDRStrk"    "opWin"        "opATSWin"    
## [31] "opOver"       "opWinSeas"    "opLosSeas"    "opWinStrk"    "opLosStrk"   
## [36] "opATSWinSeas" "opATSLosSeas" "opATSWinStrk" "opATSLosStrk" "opOVRSeas"   
## [41] "opUDRSeas"    "opOVRStrk"    "opUDRStrk"    "day"          "month"       
## [46] "year.1"       "dmy"          "code"
\end{verbatim}

\begin{longtable}[]{@{}lrrr@{}}
\caption{Predicted and Actual Points and the Actual-Predicted difference
- WNBA 2010/2024}\tabularnewline
\toprule\noalign{}
& Predicted & Actual & Diff (A-P) \\
\midrule\noalign{}
\endfirsthead
\toprule\noalign{}
& Predicted & Actual & Diff (A-P) \\
\midrule\noalign{}
\endhead
\bottomrule\noalign{}
\endlastfoot
N & 6118.0000 & 6118.0000 & 6118.0000 \\
Mean & 79.7896 & 79.8710 & 0.0814 \\
Variance & 33.0773 & 131.8456 & 101.6167 \\
Standard Deviation & 5.7513 & 11.4824 & 10.0805 \\
Correlation & 0.4793 & 0.4793 & NA \\
Standard Error & 0.0735 & 0.1468 & 0.1289 \\
Average Shapiro-Wilk P-value & 0.0584 & 0.1369 & 0.4550 \\
\end{longtable}

\begin{figure}

{\centering \includegraphics[width=0.8\linewidth]{bookdownproj_STA4211Notes1_files/figure-latex/chapter1a1-1} 

}

\caption{Distribution of team point prediction errors - WNBA 2010/2024}\label{fig:chapter1a1}
\end{figure}

Note that the Actual scores have mean \(\mu_A=79.87\), variance
\(\sigma_A^2=131.85\), and standard deviation \(\sigma_A=11.48\), and
the Predicted scores have mean \(\mu_P=79.79\), variance
\(\sigma_P^2=33.08\), and standard deviation \(\sigma_P=5.75\), and the
correlation between actual and predicted scores is \(\rho_{AP}=.479\).
This leads to the following direct computation of the mean, variance,
and standard deviations among the team point prediction errors \(D=A-P\)
based on rules for linear functions of random variables.

\[ E\{D\}=\mu_D=E\{A-P\}=\mu_A-\mu_P=79.87-79.79=0.08 \]
\[ V\{D\}=\sigma_D^2=V\{A-P\}=\sigma_A^2+\sigma_P^2-2\rho_{AP}\sigma_A\sigma_P = \]
\[=131.85+33.08-2(.479)(11.48)(5.75) = 101.62 
 \quad \Rightarrow \quad \sigma_D = 10.08 \]

\textbf{Treating these Team-Games as a Sample from a Conceptual or
``Metaphorical'' Population }

We obtain a 95\% Confidence Interval for \(\mu_D\) and a test of
\(H_0:\mu_D=0\) for the population mean Difference as follow with
\(\overline{d}=0.0814\), \(s_d^2=101.6167\), \(n=6118\),
\(\sqrt{s_d^2/n}=0.1289\),and \(t_{.975;6118-1}=1.9604\).

\[ \overline{d} \pm t_{.025;n-1}\sqrt{\frac{s_d^2}{n}}  \equiv  0.0814 \pm 1.9604(0.1289)  \equiv  0.0814 \pm 0.2527
 \equiv  (-0.1713, 0.3341) \]

The interval contains 0, showing no evidence of a systematic ``bias'' in
team point predictions.

\[
H_0:\mu_D=0 \qquad H_A:\mu_D \neq 0 \qquad TS: t^*=\frac{0.0814}{0.1289}=0.6315 
\] \[
RR: |t^*| \geq 1.9604  \qquad  P=2P\left(t_{6117} \geq 0.3221\right)=.5277
\]

\begin{Shaded}
\begin{Highlighting}[]
\FunctionTok{t.test}\NormalTok{(wnba}\SpecialCharTok{$}\NormalTok{teamPts, wnba}\SpecialCharTok{$}\NormalTok{tmPred, }\AttributeTok{paired=}\ConstantTok{TRUE}\NormalTok{)}
\end{Highlighting}
\end{Shaded}

\begin{verbatim}
## 
##  Paired t-test
## 
## data:  wnba$teamPts and wnba$tmPred
## t = 0.63155, df = 6117, p-value = 0.5277
## alternative hypothesis: true mean difference is not equal to 0
## 95 percent confidence interval:
##  -0.1712672  0.3340655
## sample estimates:
## mean difference 
##      0.08139915
\end{verbatim}

Results by year are given in Table @ref(tab:chapter1a2) and displayed in
Figure @ref(fig:chapter1a2). Only for 2011, when the average team point
prediction error was close to \(-1\), does the 95\% Confidence Interval
fail to contain 0.

\begin{longtable}[]{@{}lrrrrr@{}}
\caption{95\% Confidence Intervals for mean prediction error by year -
WNBA 2010/2024}\tabularnewline
\toprule\noalign{}
& n & mean & SD & LB & UB \\
\midrule\noalign{}
\endfirsthead
\toprule\noalign{}
& n & mean & SD & LB & UB \\
\midrule\noalign{}
\endhead
\bottomrule\noalign{}
\endlastfoot
2010 & 408 & 0.2684 & 10.2038 & -0.7247 & 1.2614 \\
2011 & 408 & -0.9975 & 9.8182 & -1.9531 & -0.0420 \\
2012 & 408 & 0.1275 & 9.6650 & -0.8132 & 1.0681 \\
2013 & 408 & -0.8480 & 10.1696 & -1.8378 & 0.1417 \\
2014 & 408 & 0.2990 & 9.7827 & -0.6531 & 1.2511 \\
2015 & 408 & -0.2892 & 10.4884 & -1.3100 & 0.7315 \\
2016 & 408 & 0.9350 & 9.9229 & -0.0307 & 1.9008 \\
2017 & 408 & 0.2880 & 10.3358 & -0.7179 & 1.2939 \\
2018 & 406 & 0.6613 & 10.4139 & -0.3547 & 1.6773 \\
2019 & 408 & 0.0833 & 10.1730 & -0.9067 & 1.0734 \\
2020 & 264 & 0.7273 & 10.2293 & -0.5124 & 1.9669 \\
2021 & 384 & -0.5260 & 9.7460 & -1.5039 & 0.4518 \\
2022 & 432 & 0.5822 & 10.2943 & -0.3913 & 1.5557 \\
2023 & 480 & 0.6625 & 10.4264 & -0.2726 & 1.5976 \\
2024 & 480 & -0.5740 & 9.4645 & -1.4228 & 0.2749 \\
\end{longtable}

\begin{figure}

{\centering \includegraphics[width=0.8\linewidth]{bookdownproj_STA4211Notes1_files/figure-latex/chapter1a2-1} 

}

\caption{95\% Confidence Intervals for mean prediction error by year - WNBA 2010/2024}\label{fig:chapter1a2}
\end{figure}

\textbf{Taking 100000 Random Samples of \(n=20\) from this Population of
Team-Games }

Now we treat this as a population of games, and consider taking (pseudo)
random samples and observing the properties of the estimated mean
differences and 95\% Confidence Intervals in repeated sampling. We take
100000 random samples of size \(n=20\), implying an estimated standard
error of \(\hat{SE}_{\overline{Y}}=\sqrt{s_d^2/n}\), where we are using
the sample variance in the standard error.

First, we compute the value
\(t=\left(\overline{d}-\mu_D\right)/\sqrt{s_d^2/n}\) for each sample,
where \(\mu_D=0.0814\) and \(n=20\). Figure @ref(fig:chapter1a3) gives a
histogram of the \(t\)-values, with a superimposed \(t_{19}\)
distribution. The model fits exceedingly well.

When we construct 95\% Confidence Intervals for \(\mu_D\) of the form
\(\overline{d} \pm t_{.975;20-1}\sqrt{s_d^2/20}\), we find that 95.06\%
include the true \(\mu_D=0.0814\). Note that different samples (based on
different seeds) will have slightly different results.

For the first random sample of the 100000 generated, we observe
\(\overline{d}=-0.300\) and \(s_d^2=121.320\).

\[
\overline{d}  \pm t_{.975;,n-1}\sqrt{\frac{s_d^2}{n}}  \equiv 
-0.300  \pm 2.093\sqrt{\frac{121.320}{20}} \equiv -0.300\pm 5.155 \equiv 
\left[-5.545,4.855\right]
\] For the test of equal variances for the paired data
\(\left(H_0:\sigma_A^2=\sigma_P^2\right)\), we should suspect that the
test will reject often, given the the true ratio is
\(\sigma_A^2/\sigma_P^2=131.85/33.08=4\). We will conduct the test
anyway to observe the \textbf{power} of the test, the proportion of
samples that reject the null hypothesis. Further, we will save the
sample correlations \(r_{AP}\) for the samples and obtain a histogram
with a line at the population correlation \(\rho_{AP}=0.4793\).

The first sample of \(n=20\) produced \(s_A^2=117.474\),
\(s_P^2=20.057\), and \(r_{AP}=0.167\). We now conduct the test
described in Section 1.4.2. First, we compute the \(F\)-statistic:
\(F=s_A^2/s_P^2=117.474/20.057=5.857\).

\[
r_{DS}=\frac{F-1}{\sqrt{(F+1)^2-4r_{AP}^2F}} = \frac{5.857-1}{\sqrt{(5.857+1)^2-4(0.167)^2(5.857)}} = 
\frac{4.857}{6.809}=0.713
\] Now, we conduct the \(t\)-test for equal variances.

\[
H_0:\sigma_A^2=\sigma_P^2 \quad H_A:\sigma_A^2\neq\sigma_P^2 \quad 
TS:t^*=\frac{0.713}{\sqrt{\frac{1-0.713^2}{20-2}}}=\frac{0.713}{0.165}=4.321
\] \[
RR:|t^*| \geq t_{.975,18}=2.093  \qquad  2P=P\left(t_{18}\geq4.321\right)=.0004
\] We conclude the variances are not equal. Of the 100000 samples, the
null hypothesis of equal variances was rejected in 90.67\% of the
samples.

Figure @ref(fig:chapter1a4) gives a histogram of the sample correlations
\(r_{AP}\) of the Actual and Predicted scores with a line at the
population correlation \(\rho_{AP}=0.479\). As is the case when the
population correlation is positive, the sampling distribution is skewed
left.

\begin{verbatim}
## [1] -0.3
\end{verbatim}

\begin{verbatim}
## [1] 121.3197
\end{verbatim}

\begin{verbatim}
## [1] 0.95059
\end{verbatim}

\begin{figure}

{\centering \includegraphics[width=0.8\linewidth]{bookdownproj_STA4211Notes1_files/figure-latex/chapter1a3-1} 

}

\caption{Distribution of t values for team point prediction errors - WNBA 2010/2024}\label{fig:chapter1a3}
\end{figure}

\begin{verbatim}
## [1] 117.4737
\end{verbatim}

\begin{verbatim}
## [1] 20.05658
\end{verbatim}

\begin{verbatim}
## [1] 0.1669815
\end{verbatim}

\begin{verbatim}
## [1] 0.90672
\end{verbatim}

\begin{figure}

{\centering \includegraphics[width=0.8\linewidth]{bookdownproj_STA4211Notes1_files/figure-latex/chapter1a4-1} 

}

\caption{Distribution of of correlations between Actual and Predicted Points - WNBA 2010/2024}\label{fig:chapter1a4}
\end{figure}

\[ \nabla  \]

Next we consider Independent samples for NBA and WNBA regular season
games for seasons ending in 2010-2024. The analyses will be done in a
similar manner to Example 1.1.

\textbf{Example 1.2: NBA and WNBA Combined Total Points and Over/Under
Line - Independent Data}

This example concerns all games for the NBA 2009/10-2023/24 and WNBA
2010-2024 regular seasons. There were 19116 NBA and 3059 WNBA games
played. For each game, we will compare the total combined points for the
two teams with the Over/Under betting line and observe the difference
(Actual - OU). We will consider the differences for the two
``associations.'' Positive differences mean the teams combine for more
points than the Over/Under line, and negative differences mean the teams
combine for fewer points than the Over/Under line. Note that \(Y\) is
the difference score for the games, but we are comparing the differences
for the NBA and WNBA, which will be independent samples.

For this example, we will do the following steps.

\begin{itemize}
\tightlist
\item
  Treating these as two populations, obtain parameters for the
  differences by league and histograms with superimposed normal
  distributions for the differences.
\item
  Treating these as single samples from two \textbf{metaphorical
  populations} where these 19116 NBA and 3059 WNBA games could be
  replayed many times, obtain 95\% Confidence Intervals for
  \(\mu_1-\mu_2\) and test \(H_0:\mu_1-\mu_2=0\), implying that on
  average, the prediction errors are equivalent for the two leagues,
  based on equal and unequal variance assumptions.
\item
  Treating these as populations, take 100000 (pseudo) random samples of
  size \(n_1=40\) and \(n_2=20\) and observe the quantity
  \(t=\left[\left(\overline{y}_1-\overline{y}_2\right)-\left(\mu_1-\mu_2\right)\right]/\sqrt{s_p^2\left(1/n_1+1/n_2\right)}\)
  and plot the superimposed \(t_{58}\) distribution. Further, obtain the
  coverage rate for the 95\% Confidence Intervals for \(\mu_1-\mu_2\),
  based on the equal variance and unequal variance assumptions.
\item
  For the 100000 samples described above, conduct the test of equal
  variances and obtain the empirical power of the \(F\)-test.
\end{itemize}

\textbf{Population Parameters and Histograms}

Population parameters for the league total prediction errors are given
in Table @ref(tab:chapter1b1) and histograms with superimposed normal
distributions for the differences are given in Figure
@ref(fig:chapter1b1).

For the NBA, we have \(N_1=19116\), \(\mu_1=0.443\), and
\(\sigma_1=17.72\). For the WNBA, we have \(N_2=3059\), \(\mu_2=0.163\),
and \(\sigma_2=16.27\). On average, both leagues Over/Under predictions
are very accurate (mean differences less than half a point), but the
variability is quite large.

\begin{verbatim}
##   team season1 tmGmYr opGmYr day month year   mdy opTeam home neutral tmPts
## 1  ATL    2006      1      1   1    11 2006 39022    PHI    0       0    75
##   opPts ot tmWin tmWinSea tmLosSea tmWinStr tmLosStr tmSprd tmWinATS
## 1    88  0     0        0        0        0        0    5.5       -1
##   tmWinATSSea tmLosATSSea tmWinATSStr tmLosATSStr tmOU tmWinOvr tmWinOvrSea
## 1           0           0           0           0  199       -1           0
##   tmLosOvrSea tmWinOvrStr tmLosOvrStr opWin opWinSea opLosSea opWinStr opLosStr
## 1           0           0           0     1        0        0        0        0
##   opSprd opWinATS opWinATSSea opLosATSSea opWinATSStr opLosATSStr opOU opWinOvr
## 1   -5.5        1           0           0           0           0  199       -1
##   opWinOvrSea opLosOvrSea opWinOvrStr opLosOvrStr                    code X
## 1           0           0           0           0 1PHIATL2006\\\\11\\\\01
\end{verbatim}

\begin{verbatim}
##   team      date year tmGmYr opGmYr home teamPts oppPts OT tmSprd    OU tmSpDev
## 1  ATL 5/15/2010 2010      1      1    0      75     70  0    5.5 168.5    10.5
##   oppTeam tmWin tmATSWin tmOver tmWinSeas tmLosSeas tmWinStrk tmLosStrk
## 1      LV     1        1     -1         0         0         0         0
##   tmATSWinSeas tmATSLosSeas tmATSWinStrk tmATSLosStrk tmOVRSeas tmUDRSeas
## 1            0            0            0            0         0         0
##   tmOVRStrk tmUDRStrk opWin opATSWin opOver opWinSeas opLosSeas opWinStrk
## 1         0         0     0       -1     -1         0         0         0
##   opLosStrk opATSWinSeas opATSLosSeas opATSWinStrk opATSLosStrk opOVRSeas
## 1         0            0            0            0            0         0
##   opUDRSeas opOVRStrk opUDRStrk day month year.1       dmy            code
## 1         0         0         0  15     5   2010 15/5/2010 1ATLLV15/5/2010
\end{verbatim}

\begin{verbatim}
## 
##  Two Sample t-test
## 
## data:  nba_tot_diff and wnba_tot_diff
## t = 0.8212, df = 22173, p-value = 0.4115
## alternative hypothesis: true difference in means is not equal to 0
## 95 percent confidence interval:
##  -0.3886726  0.9491923
## sample estimates:
## mean of x mean of y 
## 0.4430582 0.1627983
\end{verbatim}

\begin{verbatim}
## 
##  Welch Two Sample t-test
## 
## data:  nba_tot_diff and wnba_tot_diff
## t = 0.87327, df = 4303.3, p-value = 0.3826
## alternative hypothesis: true difference in means is not equal to 0
## 95 percent confidence interval:
##  -0.3489294  0.9094492
## sample estimates:
## mean of x mean of y 
## 0.4430582 0.1627983
\end{verbatim}

\begin{longtable}[]{@{}lrr@{}}
\caption{NBA and WNBA Combined Points - Over/Under
2010-2024}\tabularnewline
\toprule\noalign{}
& NBA & WNBA \\
\midrule\noalign{}
\endfirsthead
\toprule\noalign{}
& NBA & WNBA \\
\midrule\noalign{}
\endhead
\bottomrule\noalign{}
\endlastfoot
N & 19116.0000 & 3059.0000 \\
Mean & 0.4431 & 0.1628 \\
Variance & 313.8896 & 264.7473 \\
Standard Deviation & 17.7169 & 16.2711 \\
Standard Error & 0.1281 & 0.2942 \\
\end{longtable}

\begin{figure}

{\centering \includegraphics[width=0.8\linewidth]{bookdownproj_STA4211Notes1_files/figure-latex/chapter1b1-1} 

}

\caption{Distribution of combined point prediction errors - NBa and WNBA 2010/2024}\label{fig:chapter1b1}
\end{figure}

\textbf{Treating these Games as Samples from Conceptual or
``Metaphorical'' Populations }

First, we consider the equal variance case, noting that the standard
deviations are fairly similar (17.72 and 16.27, respectively), however
with these large samples, we will easily reject
\(H_0:\sigma_1^2=\sigma_2^2\).

\[
TS: F^*=\frac{313.89}{264.75}=1.186 \qquad RR: F^*\leq F_{.025,19915,3058}= 0.948 \cup
F^*\geq F_{.025,19915,3058}= 1.056
\]

Then we will consider the unequal variance case.

The summary statistics and calculations are given below.

\[
\mbox{NBA: } n_1=19116 \quad \overline{y}_1=0.443 \quad s_1^2 = 313.89 \qquad 
\mbox{WNBA: } n_2=3059 \quad \overline{y}_2=0.163 \quad s_2^2 = 264.75
\] \[
s_P^2 = \frac{(n_1-1)s_1^2+(n_2-1)s_2^2}{n_1+n_2-2}=
\frac{(19116-1)(313.89)+(3059-1)(264.75)}{19119+3059-2} = 307.11
\] \[
\hat{SE}\left\{\overline{Y}_1-\overline{Y}_2\right\}=
\sqrt{s_P^2\left[\frac{1}{n_1}+\frac{1}{n_2}\right]}
=\sqrt{307.11\left[\frac{1}{19116}+\frac{1}{3059}\right]}=0.341
\] We obtain a 95\% Confidence Interval for \(\mu_1-\mu_2\) and a test
of \(H_0:\mu_1-\mu_2=0\) for the population mean Difference as follow
with \(\overline{y}_1-\overline{y}_2=0.280\),
\(\hat{SE}\left\{\overline{Y}_1-\overline{Y}_2\right\}=0.341\), and
\(t_{.975;19116+3059-2}=1.9601\).

\[ \overline{y}_1-\overline{y}_2 \pm t_{.975;n_1+n_2-2}\hat{SE}\left\{\overline{Y}_1-\overline{Y}_2\right\}  \equiv  0.280 \pm 1.9601(0.341)  \equiv  0.280 \pm 0.668
 \equiv  (-0.388,0.948) 
 \] \[
H_0:\mu_1-\mu_2=0 \quad H_A:\mu_1-\mu_2\neq 0 \qquad
TS:t^*=\frac{\overline{y}_1-\overline{y}_2}{\hat{SE}\left\{\overline{Y}_1-\overline{Y}_2\right\}}=\frac{0.280}{0.341}=0.821
\] \[
RR: |t^*| \geq t_{.975;22173} = 1.9601 
P=2P\left(t_{22173}\geq|t^*|\right)=.4117
\]

Thus, even with these large ``samples'', we fail to reject
\(H_0: \mu_1-\mu_2=0\).

Now, we consider the unequal variance case (Welch's Test, with
Satterthwaite's Approximation for the degrees of freedom). The sample
mean difference doesn't change. We obtain the estimated standard error
of the mean difference and approximate degrees of freedom here.

\[
\hat{SE}\left\{\overline{Y}_1-\overline{Y}_2\right\}=
\sqrt{\frac{s_1^2}{n_1}+\frac{s_2^2}{n_2}}=\sqrt{\frac{313.89}{19116}+\frac{264.75}{3059}}
=0.321
\] \[
\hat{\nu}^* = \frac{\left(\frac{s_1^2}{n_1}+\frac{s_2^2}{n_2}\right)^2}{
\frac{\left(\frac{s_1^2}{n_1}\right)^2}{n_1-1} + 
\frac{\left(\frac{s_2^2}{n_2}\right)^2}{n_2-1}
} = 
\frac{\left(\frac{313.89}{19116}+\frac{264.75}{3059}\right)^2}{
\frac{\left(\frac{313.89}{19116}\right)^2}{19116-1} + 
\frac{\left(\frac{264.75}{3059}\right)^2}{3059-1}
}=4303.65
\]

Now, \(t_{.975;4303.65}=1.9605\), and we obtain a 95\% Confidence
Interval for \(\mu_1-\mu_2\) and test \(H_0:\mu_1-\mu_2=0\).

\[
0.280 \pm 1.9605(0.321)\equiv 0.280\pm 0.629 \equiv (-0.349,0.909)
\] \[
H_0: \mu_1-\mu_2=0  \qquad  TS: t^*=\frac{0.280}{0.321}=0.872
\] \[
RR: |t^*| \geq 1.9605 \qquad P=2P\left(t_{4303.65}\geq |t^*|\right)=.3833
\]

These results are in agreement with \textbf{t.test} functions in the
previous code block (up to rounding error).

\textbf{Taking 100000 Random Samples of \(n_1=40\) and \(n_2=20\) from
these Populations of Games}

Now we treat these as two populations of games, and consider taking
(pseudo) random samples and observing the properties of the estimated
mean differences and 95\% Confidence Intervals in repeated sampling,
based on both the equal and unequal variance cases. We take 100000
random samples of sizes \(n_1=40\) NBA and \(n_2=20\) WNBA games.

First, we compute the value
\(t=\left[\left(\overline{y}_1-\overline{y}_2\right)-\left(\mu_1-\mu_2\right)\right]/\hat{SE}\left\{\overline{Y}_1-\overline{Y}_2\right\}\)
for each sample, where \(\mu_1-\mu_2=0.280\) with \(n_1=40\) and
\(n_2=20\). Figure @ref(fig:chapter1b2) gives a histogram of the
\(t\)-values, with a superimposed \(t_{58}\) distribution. The model
gives a reasonable fit, but does not fit as well as in Example 1.1, due
to unequal variances and sample sizes.

As we have computed Confidence Intervals and tests for the ``full
sample'', we will not repeat the calculations for our first random
samples here.

When we construct 95\% Confidence Intervals for \(\mu_1-\mu_2\), we find
that 92.88\% of the Equal variance case intervals include the true
\(\mu_1-\mu_2=0.280\). Based on the unequal variance case, 94.55\% cover
0.280. Note that different samples (based on different seeds) will have
slightly different results.

For testing \(H_0:\sigma_1^2=\sigma_2^2\), \(F_{.025,39,19}=0.477\) and
\(F_{.975,39,19}=2.338\), so we will reject \(H_0\) if \(s_1^2/s_2^2\)
is less than 0.477 or greater than 2.338. Of these 100000 samples,
\(H_0\) is rejected in 36.15\% of samples.

\begin{verbatim}
## [1] 0.92876
\end{verbatim}

\begin{verbatim}
## [1] 0.94554
\end{verbatim}

\begin{verbatim}
##      Min.   1st Qu.    Median      Mean   3rd Qu.      Max. 
## -5.537120 -0.695593  0.001199  0.017788  0.715335  6.380007
\end{verbatim}

\begin{figure}

{\centering \includegraphics[width=0.8\linewidth]{bookdownproj_STA4211Notes1_files/figure-latex/chapter1b2-1} 

}

\caption{Distribution of t values for combined team point prediction errors - NBA and WNBA 2010/2024}\label{fig:chapter1b2}
\end{figure}

\begin{verbatim}
## [1] 0.36154
\end{verbatim}

\[
\nabla
\]

\chapter{Introduction}\label{introduction}

In this chapter, we will introduce terminology and briefly describe
experimental designs used in a wide variety of research fields.

\section{Terminology}\label{terminology}

Studies involve making observations on individual units under various
conditions. When units have been randomly assigned to a treatment
condition, they may be referred to as \textbf{Experimental Units}. When
units have been sampled and observed from an existing population they
may be referred to as \textbf{Observational Units}. In some studies,
larger blocks of units may be randomized to a treatment, with subunits
being observed (as when classrooms are randomized to various
conditions), in this setting the measured units (e.g.~students) may be
referred to as \textbf{Measurement Units}.

\begin{itemize}
\tightlist
\item
  \textbf{Experimental Studies} - Observational units are assigned at
  random to treatments/conditions

  \begin{itemize}
  \tightlist
  \item
    \textbf{Experimental Factors} - Conditions with two or more levels
    that are assigned (at random) to units. Many experiments include
    multiple factors and treatments are the combinations of factor
    levels.
  \end{itemize}
\item
  \textbf{Observational Studies} - Observational units are sampled from
  various populations/subpopulations

  \begin{itemize}
  \tightlist
  \item
    \textbf{Observational Factors} - Set of populations/subpopulations
    observed in a study
  \end{itemize}
\item
  \textbf{Mixed Studies} - Studies that have both experimental and
  observational factors (not to be confused with Mixed Effects Designs)
\end{itemize}

\textbf{Example 1.1 - Effect of Container Size on Food Intake}

A study was conducted to compare three conditions on food intake in
students (\citeproc{ref-marchiori_2012}{Marchiori, Cornielle, and Klein
2012}). A sample of 88 subjects was obtained, and randomly assigned to
one of 3 conditions involving bowl size and portion of M\&Ms while
watching a television program: 1) Medium portion size/Small container
\((n_1 = 30)\), 2) Medium portion/Large container \((n_2 = 29)\), and 3)
Large portion/Large container \((n_3 = 29)\). Researchers measured the
food intake among the students. Note that this is an
\textbf{Experimental Study}, as Subjects were assigned at random to
treatments.

\[\nabla\]

\textbf{Example 1.2 - Waste in the Mediterranean Sea}

A study measured the amounts of Natural and Artificial floating debris
at samples of transects at 14 locations in the Mediterranean Sea
(\citeproc{ref-suaria_2014}{Suaria and Aliani 2014}). The researchers
wished to compare the amounts of debris of each type among the
locations. Note this is an \textbf{Observational Study}, as transects
were sampled within the selected locations.

\[\nabla\]

\textbf{Example 1.3 - Quilting Layers in Body Armour}

A study was conducted to determine whether the number of quilting layers
improved the fragment protective performance of body armour
(\citeproc{ref-carr_2012}{Carr et al. 2012}). The researchers sampled 36
specimens of each number of layers (1,2,3, and 5), assigning 12 at
random to each of 3 bullet impacts (slow, fast, and edge). The energy
absorbed by each specimen was measured. Note this is a \textbf{Mixed
Study}, as Layers is Observational, and Impact is Experimental.

\[\nabla\]

\section{Basics of Controlled
Experiments}\label{basics-of-controlled-experiments}

In this section we describe some aspects and terminology of controlled
experiments and give brief examples of them.

\begin{itemize}
\tightlist
\item
  \textbf{Explanatory Factors} -- Conditions (with 2 or more levels)
  that are assigned to units.
\item
  \textbf{Crossed Factors} -- Factors with levels that are the same
  within levels of the other factor(s)
\item
  \textbf{Nested Factors} -- Factors with levels that are different
  within levels of the other factor(s)
\item
  \textbf{Treatments} -- Combinations of factor levels given to units
\item
  \textbf{Experimental Units} -- Units used in the study, which are
  subject to randomization to treatments.
\item
  \textbf{Randomization Process} - Use of random number generator to
  assign units to treatments
\item
  \textbf{Response(s)} - Outcome measurement(s) obtained from treated
  units
\end{itemize}

\textbf{Example 1.4 - Reading Times on 3 Electronic Readers at 4
Illumination Levels}

An experiment was conducted to measure reading times on 3 e-reader
devices at 4 illumination levels (\citeproc{ref-chang_2013}{Chang, Chou,
and Shieh 2013}). A sample of 60 subjects were randomly assigned so that
5 received each of 12 treatments (combinations of 3 e-reader models and
4 illumination levels). The experiment was \textbf{Crossed} in the sense
that each e-reader model was set at the same 4 illumination levels (200,
500, 1000, 1500Lx). The time to complete a reading task was the measured
response.

\[\nabla\]

\textbf{Example 1.5 - Combability of Hair for Two Shampoo Formulations}

An experiment was conducted to compare two shampoo formulations with
respect to combability of hair (\citeproc{ref-garcia_1976}{Garcia and
Diaz 1976}). A sample of 16 hair swatches were created and randomly
assigned, such that 8 received shampoo A, and 8 received shampoo B. Each
swatch was washed 5 times, and the combability was measured. The
experiment was \textbf{Nested}, as the swatches receiving shampoo A were
different from the swatches receiving shampoo B. Note that the swatches
are the experimental units, as they are randomly assigned to treatments
(shampoos). The replicates measured are the measurement units.

\[\nabla\]

\section{Completely Randomized Design
(CRD)}\label{completely-randomized-design-crd}

In the \textbf{Completely Randomized Design}, experimental units are
randomly assigned to treatments, and responses are recorded on the units
after treatments are applied. We will refer to the number of treatments
as \(r\), with the number of replicates for the \(i^{th}\) treatment
being \(n_i\). When all treatments have the same number of replicates
\((n_1=\ldots =n_r=n)\), the design is said to be \textbf{balanced}. The
total sample size across all treatments will be labelled
\(n_T=n_1+\cdots+n_r\). When the experiment is balanced, \(n_T=rn\).

The statistical model for the One-Way Analysis of Variance based on the
Completely Randomized Design is as follows, where the subscript \(i\)
represents the treatment and \(j\) represents the replicate within the
treatment.

\[ Y_{ij}=\mu_i + \epsilon_{ij} = \mu + \tau_i + \epsilon_{ij}
\qquad i=1,\ldots,r; \quad j=1,\ldots, n_i \qquad \epsilon_{ij} \sim NID\left(0,\sigma^2\right) \]

\textbf{Example 1.6 - Anthocyanin Extractability in Cabernet Franc
Grapes}

In a study conducted by researchers in France and Italy, Cabernet Franc
grapes were harvested at \(r=6\) different classes of sugar content
(176.5, 192.6, 209.3, 225.0, 242.1, and 258.5 grams/litre). While these
are numeric levels, the authors treated sugar content as a factor
variable. There were \(n=15\) berries within each treatment for a total
of \(n_T=6(15)=90\) berries included in the study. Various physical,
textural, and anthocyanin extractability measurements were made. We will
focus on extraction yield of anthocyanin, which was labelled in the
paper as EA\%, (\citeproc{ref-zouid_2013}{Zouid et al. 2013}). A plot of
the data is given in Figure @ref(fig:ch0101).

\begin{verbatim}
## [1] "sugar"   "anthExt"
\end{verbatim}

\begin{figure}

{\centering \includegraphics[width=0.8\linewidth]{bookdownproj_STA4211Notes1_files/figure-latex/ch0101-1} 

}

\caption{Berry texture scores by sugar content}\label{fig:ch0101}
\end{figure}

\[\nabla\]

\section{Randomized Complete Block Design (RCBD or
RBD)}\label{randomized-complete-block-design-rcbd-or-rbd}

In the \textbf{Randomized Complete Block Design} (which is often simply
referred to as the \textbf{Randomized Block Design}), experimental units
are blocked into ``groups'' of homogeneous units. Within each block,
units are randomly assigned to the treatments, with each treatment being
applied to one unit within each block. We will continue using \(r\) as
the number of treatments and will use \(b\) for the number of blocks.
The goal is to remove the heterogeneity across blocks to obtain more
precise comparisons among treatments, when possible. In many instances,
blocks will be the same individual that will receive each treatment when
this is feasible. Blocks are typically treated as a random factor, in
the sense that results are to be generalized across a population of such
blocks or individuals. When the blocks are individuals who receive each
treatment, this design is often referred to as a \textbf{Repeated
Measures Design} or a \textbf{Crossover Design}.

The statistical model can be written as follows where the subscript
\(i\) represents the treatment and \(j\) represents the block.

\[ Y_{ij}=\mu_i + \beta_j + \epsilon_{ij} = \mu + \tau_i + \beta_j + \epsilon_{ij}
\qquad i=1,\ldots,r; \quad j=1,\ldots, b \qquad \epsilon_{ij} \sim NID\left(0,\sigma^2\right) \]

\textbf{Example 1.7 - Comparison of 4 Treadmill Models for User
Satisfaction}

Researchers in Italy and New Zealand conducted an experiment to compare
\(r=4\) treadmill models (Life Fitness, Precor, Matrix, Technogym) among
\(b=57\) trained runners (\citeproc{ref-carraro_2019}{Carraro, Elliott,
and Gobbi 2019}). Each runner rated each treadmill model in terms of
seven characteristics: Running Surface, Controls, Stability and Safety,
Physical Interaction, Console Readability, Aesthetic Appeal, and
Enjoyment of Use. The responses were measured on visual Analogue Scales
(VAS) from very unpleasant to very pleasant. We will consider Enjoyment
of Use (enjoyUse) in Figure @ref(fig:ch0102) and Figure
@ref(fig:ch0102a).

\begin{verbatim}
##   trdMill subject runSurf controls stblSfty physIntrct consRead aesthApp
## 1       1       1  8.5473   6.2358   5.0478     6.8996   6.2044   4.5798
## 2       1       2  7.4018   7.2364   6.0081     2.5457   7.1048   6.4683
## 3       1       3  5.1259   8.6687   8.8688     6.4669   4.6827   9.3432
## 4       1       4  4.8110   4.8361   9.5405     5.1269   6.7633   6.4748
## 5       1       5 10.1249   4.2111   9.5724     6.4366   6.6144   7.0896
## 6       1       6  5.1529   6.5803   8.0583     5.8038   3.7302   6.3881
##   enjoyUse
## 1   8.9383
## 2   4.8944
## 3   7.9970
## 4   7.4555
## 5   3.7710
## 6   5.2169
\end{verbatim}

\begin{figure}

{\centering \includegraphics[width=0.8\linewidth]{bookdownproj_STA4211Notes1_files/figure-latex/ch0102-1} 

}

\caption{Enjoyment of use of 4 treadmills for 57 runners}\label{fig:ch0102}
\end{figure}

\begin{figure}

{\centering \includegraphics[width=0.8\linewidth]{bookdownproj_STA4211Notes1_files/figure-latex/ch0102a-1} 

}

\caption{Boxplots of Enjoyment of use of 4 treadmills for 57 runners}\label{fig:ch0102a}
\end{figure}

\[\nabla\]

\section{Overview of Some Standard Experimental
Designs}\label{overview-of-some-standard-experimental-designs}

In this section we list some commonly used experimental designs as well
as some examples of them.

\begin{itemize}
\tightlist
\item
  \textbf{Completely Randomized Design (CRD)} -- Units randomized to
  treatments with no restrictions on randomization process
\item
  \textbf{Factorial Experiments} -- CRD with two or more crossed
  factors. Treatment effects are made up of main factor effects and
  interaction effects
\item
  \textbf{Randomized Complete Block Design (RCBD)} -- Units are grouped
  into blocks. Treatments are randomly assigned to units within blocks
\item
  \textbf{Nested Designs} -- Levels of Factor B differ across levels of
  Factor A
\item
  \textbf{Crossed/Nested Designs} -- Designs with both crossed and
  nested factors
\item
  \textbf{Repeated Measures Designs} -- Each unit is measured multiple
  times

  \begin{itemize}
  \tightlist
  \item
    \textbf{RBD} - Each subject receives each treatment once
  \item
    \textbf{CRD} Each subject receives only one treatment, but is
    measured at multiple time points
  \end{itemize}
\item
  \textbf{Split-Plot Designs} -- Two (or more) sizes of experimental
  units due to randomization restrictions for factors
\item
  \textbf{Incomplete Block Designs} -- Block Designs with block sizes
  smaller than the number of treatments
\item
  \textbf{2-Level Factorial Experiments} -- Several (possibly many)
  factors, each at 2 levels (low/high). With k factors, there will be
  \(2^k\) treatments
\item
  \textbf{2-Level Fractional Factorial Designs} -- Experiments with only
  a subset of all \(2^k\) treatments to reduce cost, but still obtain
  estimates of main effects and lower-order interactions
\item
  \textbf{Response Surface Designs} -- Designs used to fit polynomial
  regression models to optimize responses for numeric factors
\item
  \textbf{Mixture Designs} -- Designs used to fit models to optimize
  responses among mixtures (components sum to 1) of numeric factors.
\end{itemize}

\textbf{Example 1.8 - Advertising Messaging Strategy and Attitude to the
Firm}

An experiment was conducted to compare 4 advertisement conditions
(\citeproc{ref-hyllegard_2009}{Hyllegard, Ogle, and Yan 2009}). A sample
of 425 students were selected and randomly assigned to one of 4
conditions. The ads were:

\begin{itemize}
\tightlist
\item
  Ad1: Firm as ``pioneer of industry standards in social
  responsibility'' and US location
\item
  Ad2: Young woman partially clothed in shower, winner of wet t-shirt
  contest
\item
  Ad3: Female co-founder of porn mag for women, in jogging shorts/hoodie
\item
  Ad4: Female and male partially clothed couple in bed, faces cropped
  out of image. Female on top of male.
\end{itemize}

The response was an overall attitude toward the firm based on a series
of rating items. Note that each student was exposed to only one
condition. Condition 1: Ad1 Only, Condition 2: Ad1\&Ad2, Condition 3:
Ad1\&Ad3, Condition 4: Ad1\&Ad4.

\[\nabla\]

\textbf{Example 1.9 - Energy Efficiency of 4 Dryer Types and 3 Clothing
Categories}

A study compared combinations of 4 dryer types and 3 clothing categories
on energy efficiency (\citeproc{ref-to_2007}{To et al. 2007}). The dryer
types were (1=Electric Dryer, 2=Bi-directional Electric dryer, 3=Town
Gas-Fired Dryer, 4=LPG-Fired dryer) and the clothing categories were
(1=Towels, 2=Jeans, 3=Thermal Clothing). The response was Energy
Efficiency (kWh/kg), and there were 3 replications per treatment.

\[\nabla\]

\textbf{Example 1.10 - Comparison of 6 Chopstick Lengths on Feeding
Efficiency}

A study compared 6 chopstick lengths (180, 210, 240, 270, 300, 330mm) in
terms of the numbers of peanuts picked up and placed in a cup
(\citeproc{ref-hsu_1991}{Hsu and Wu 1991}). There were 31 subjects, and
each subject used each chopstick. The subjects act as blocks. This can
also be treated as a Repeated Measures Design, where each subject
receives each treatment.

\[\nabla\]

\textbf{Example 1.11 - Caffeine Content of Coke and Pepsi Products at
Various Restaurants}

A study compared Coca-Cola and Pepsi-Cola at various restaurants with
respect to caffeine content (\citeproc{ref-grand_1997}{Grand and Bell
1997}). There were 5 restaurants that sold Coca-Cola brand (1=Red
Lobster, 2=Applebees, 3=McDs, 4=BK, 5=Hardees) and 7 that sold Pepsi
products (6=Arbys, 7=Subway2, 8=Subway1, 9=KFC, 10=PizzaHut,
11=TacoBell, 12=Wendys). Each restaurant sold both sugar and diet
formulations. There were 10 measurements per restaurant per formulation.
Note that restaurant is nested within brand, but crossed with
formulation.

\[\nabla\]

\textbf{Example 1.12 - Zylkene v Placebo for Cats with Anxiety Over
Time}

A study compared the effects of Zylkene v Placebo in cats with anxiety
(\citeproc{ref-beata_2007}{Beata, Cordel, and Marlois 2007}). A sample
of 34 cats with anxiety was obtained, and randomized to receive either
Zylkene or Placebo (17 cats per treatment). Each cat was observed on a
global anxiety scale at each of 5 time points. The goal is to compare
the treatments and determine whether time effects occur, and whether the
treatment effects differ over time.

\[\nabla\]

\textbf{Example 1.13 - Effects of Seeding Rates of Hardinggrass and
Ryegrass on Growth}

An experiment was conducted to measure the effects of 4 rates of seeding
of the perennial hardinggrass and 6 rates for ryegrass
(\citeproc{ref-schultz_1952}{Schultz and Biswell 1952}). Hardinggrass
was measured on whole (larger) plots, with levels of 1,2,3,4 pounds per
acre) while ryegrass which was measured on subplots within the whole
plots with levels of 0,3,6,9,12,15 pounds per acre). The experiment was
conducted in 3 blocks (replicates). The response measured was the
density of hardinggrass. Note that the experimental units for levels of
hardinggrass are larger than those for levels of ryegrass.

\[\nabla\]

\textbf{Example 1.14 - Consumer Liking of 16 Dark Chocolate
Formulations}

An experiment was conducted to compare consumer liking among 16 dark
chocolate formulations (\citeproc{ref-hinneh_2020}{Hinneh et al. 2020}).
As consumers have quite different tastes and preferences, they are
treated as blocks, and there were 16 raters. However, due to fatigue,
the researchers had each consumer rate only 6 dark chocolate
formulations. In a balanced experiment, each formulation would be rated
by the same numbers of consumers, and each pair of formulations would be
tasted by equal numbers of consumers. In this case, each chocolate was
rated by 6 raters, and each pair of chocolates were rated by 2 raters.
This represents a Balanced Incomplete Block Design.

\[\nabla\]

\textbf{Example 1.15 - Using Seaweed to Extract Phenol from Aqueous
Solution}

An experiment was conducted to study the effects of 3 factors (pH (3,
9), adsorbent dosage (1, 10 g/L), and temperature (30, 60C)) on phenol
extraction efficiency (\%) from an aqueous solution
(\citeproc{ref-ranthinam_2011}{Ranthinam, Rao, and Nair 2011}). Dried
seaweed was treated with zinc chloride, then applied to one of the 8
combinations of the 3 factors. There were two replicates at each factor
level. This is an example of a \(2^3\) full-factorial design.

\[\nabla\]

\textbf{Example 1.16 - Factors Affecting Damage to Motorcycle Wheels}

An experiment was conducted to determine the effect of 5 factors on the
Crush Radius on the front wheel of a motorcycle
(\citeproc{ref-tan_2009}{Tan et al. 2009}). The factors were: Impact
speed (3, 6), Impact Mass (51.18, 101.33), Tire Pressure (148, 252),
Striker Contact Geometry (0.03, 0.10), and Impact Offset Distance (0,
0.108). Although there are \(2^5 = 32\) combinations of the factors, the
experimenters ran it in \(2^{5-1}=16\) combinations to reduce costs.
There were four replicates at each combination of factor settings. This
is an example of a \(2^{5-1}\) fractional factorial design.

\[\nabla\]

\textbf{Example 1.17 - Optimizing Qualities of Potato Chips}

An experiment varied 3 factors to optimize responses regarding potato
chips (\citeproc{ref-song_2007}{Song, Zhang, and Mujumdar 2007}). The
factors (and levels) were: Vacuum microwave pre-drying time (0.95, 3, 6,
9, 11.05 minutes), Vacuum temperature (83.18, 90, 100, 110, 116.82C),
and Frying Time (11.59, 15, 20, 25, 28.41 minutes). Three responses were
measured (analyzed one at a time): Moisture Content, Fat Content, and
Breaking Force. The goal was to choose factor levels that optimize the
response. This is an example of a response surface design.

\[\nabla\]

\textbf{Example 1.18 - Optimizing Antibacterial Effect of Mixtures of 3
Essential Oils}

An experiment varied mixtures of 3 types of essential oils in terms of
optimizing a response (\citeproc{ref-ouedrhiria_2016}{Ouedrhiria et al.
2016}). The three oils were: (O. compactum, O. majorana, and T.
serypyllum). Three responses were measured (analyzed one-at-a-time):
minimum inhibitory concentration (MIC \%) of 3 types of bacterium: B.
subtilis. S. aureus, and E. coli. The goal was to choose the mixture of
the 3 types of essential oils that minimizes the MIC of the bacterium.

\[\nabla\]

\section{Overview of Observational Study
Designs}\label{overview-of-observational-study-designs}

\begin{itemize}
\tightlist
\item
  \textbf{Cross-Sectional Studies} -- Observations made from
  populations/subpopulations at a single time point or interval.
\item
  \textbf{Prospective Studies} -- Groups are formed by levels of a
  potential causal factor, then observed over time for some measurable
  outcome.
\item
  \textbf{Retrospective Studies} -- Studies where subjects are
  identified based on the outcome of interest and potential risk factors
  are identified that previously occurred.
\item
  \textbf{Matching} -- Subjects from different populations are matched,
  based on external factors, similar to blocking in experimental
  studies.
\end{itemize}

\textbf{Example 1.19 - Medical Profession Students Attitudes Toward
Interdisciplinary Studies}

A survey was conducted to measure students' in medical professions
Readiness for Inter-Professional Learning
(\citeproc{ref-keshtkaran_2014}{Keshtkaran, Sharif, and Rambod 2014}).
Students were sampled from 3 groups: Nursing, Science in Surgical
Technology, and Medicine, and given the scale measuring their readiness.
This was cross-sectional in the sense that it was taken at one point in
time.

\[\nabla\]

\textbf{Example 1.20 - Blood Transfusions and Caesarean Deliveries in 3
Pakistan Hospitals}

A prospective study was conducted to compare birth deliveries in 3
Pakistan hospitals over the period of January-June 2010
(\citeproc{ref-ismail_2014}{Ismail et al. 2014}). In particular, the
authors were interested in whether the mother had a Caesarean Section
and whether the patient had a subsequent blood transfusion.

\[\nabla\]

\textbf{Example 1.21 - Fertilization Time-Lapse Variables and Embryo
Sex}

A study at a University-affiliated private fertility center considered
the gender of an embryo, as well as various cleavage timing variables,
from the time of the fertilization, retrospectively
(\citeproc{ref-bronet_2015}{Bronet et al. 2015}). The researchers were
interested in determining whether any of the timing variables could
predict the eventual sex of the embryo.

\[\nabla\]

\textbf{Example 1.22 - Recidivism Rates for Juvenile Offenders}

A study compared recidivism rates among juvenile offenders in 2
conditions: transferred to adult court and not transferred to adult
court, that is, tried in juvenile court
(\citeproc{ref-bishop_1996}{Bishop et al. 1996}). A database of past
criminal record including number and severity of prior and current
charges, gender, and age was created, and matches were created where
within each pair, one had been transferred, the other had not.
Subsequent recidivism was observed within each pair.

\[\nabla\]

\chapter{Single Factor Studies}\label{single-factor-studies}

In this chapter, we consider the case where there is a single factor.
The factor can be \textbf{qualitative} or \textbf{quantitative}.
Qualitative factors can be \textbf{nominal} or \textbf{ordinal}. Nominal
factors have no inherent ordering, while ordinal factors have an
underlying ordering that is non-numeric. Quantitative factors can be
used in linear or nonlinear regression models, however in some cases
researchers don't wish to place a structure to the model and treat
levels as categories.

The factor levels, unless specified otherwise, will be treated as all
levels of interest to researchers. These are referred to as
\textbf{fixed factors}. Later in the course, we will consider factors
with levels that are a sample from a population of levels, which will be
referred to as \textbf{random factors}.

We will let the number of levels of the factor be represented as \(r\)
where \(r \geq 2\). When \(r=2\), the analysis can be analyzed as an
independent sample \(t\)-test, which we will see gives identical
conclusions as the 1-Way Analysis of Variance described below.

The analysis will be the same whether the data are obtained from a
Controlled Experiment or Observational Study, however interpretations of
cause and effect are stronger for Controlled Experiments.

In this chapter, we will focus on models with independent, normally
distributed errors with constant variance. We will consider other tests
in subsequent chapters.

\section{\texorpdfstring{Models with \(r=2\) Factor
Levels}{Models with r=2 Factor Levels}}\label{models-with-r2-factor-levels}

In this section, we consider a Controlled Experiment and an
Observational Study, each with two groups. We use the independent sample
\(t\)-test in each case to compare the two means.

\textbf{Example 2.1 - Tai Chi and Strength Training in Hip Replacement
Patients}

A controlled experiment was conducted to compare post-operative effects
of an in-home training treatment for patients who have received hip
replacement surgery (\citeproc{ref-zeng_2015}{Zeng et al. 2015}). The
final analysis was based on \(n_T=59\) subjects who completed the
protocol, 32 in the treatment group (Tai Chi and strength training) and
27 in the control group. Patients were randomly assigned to the \(r=2\)
conditions. One response that was reported was Timed Up and Go Test
(TUG) to measure mobility (lower times are better). The means, standard
deviations, and sample sizes are given below for the treatment and
control groups (the units are seconds).

\[ \mbox{Treatment: } \overline{y}_T=14.61 \quad s_T= 2.60 \quad n_T=32 \]
\[ \mbox{Control: } \overline{y}_C=19.06 \quad s_C=3.37 \quad n_C=27 \]

For the 2-sample \(t\)-test (aka independent sample \(t\)-test), we will
first compute the pooled variance then the \(t\)-statistic, critical
value \((\alpha=0.05)\) and \(P\)-value for a 2-sided test.

\[ \mbox{Pooled Variance:} \quad s_p^2 = \frac{(32-1)(2.60)^2+(27-1)(3.37)^2}{32+27-2}
= 8.86  \qquad s_p=\sqrt{8.86}=2.98\]

\[
\mbox{Test Statistic:} \quad t^* = \frac{14.61-19.06}{\sqrt{8.86\left(\frac{1}{32}+\frac{1}{27}\right)}}
=\frac{-4.45}{0.78}=-5.71 
\]

\[
\mbox{Rejection Region:} \quad |t^*| \geq t_{.975,57}=2.002  \qquad
\mbox{P-value:} \quad 2P\left(t_{57}\geq |-5.71|\right)<.0001 
\] \[
95\% \mbox{ Confidence Interval:} \quad -5.71 \pm 2.002(0.78) \quad \equiv \quad
-5.71 \pm 1.56 \quad \equiv \quad (-7.27,-4.15)
\]

The \textbf{Effect Size} is the absolute difference in the sample means
in units of the pooled standard deviation. In this example,
\(e=4.45/2.98=1.49\) which is considered a large effect size.

\[\nabla\]

\textbf{Example 2.2 - Firefighter Air Consumption}

An observational study was conducted to compare firefighters that were
classified into two categories of air consumption, slow and fast
(\citeproc{ref-wohlgemuth_2024}{Wohlgemuth et al. 2024}). The
firefighters were classified based on scores for 10 tasks pertinent to
firefighting. Once the \(n_T=160\) firefighters were classified into the
two groups (94 fast and 66 slow), the groups were compared with respect
to several physical characteristics including Age, Body Mass Index, and
Peak Heart Rate (none of these were used in classifying as Fast/Slow Air
Consumption). Summary statistics for the response Body Fat Percentage
and the corresponding \(t\)-test are given below.

\[
\mbox{Fast: } \overline{y}_F=20.71 \quad s_F=6.65 \quad n_F=94 \qquad
\mbox{Slow: } \overline{y}_S=25.11 \quad s_S=8.93 \quad n_S=66 \]

\[
s_p^2 = \frac{(94-1)(6.65)^2+(66-1)(8.93)^2}{94+66-2}=58.84  \qquad  s_p=7.67\]

\[
t^*=\frac{20.71-25.11}{\sqrt{58.85\left(\frac{1}{94}+\frac{1}{66}\right)}} =
\frac{-4.40}{1.23}=-3.57  \qquad t_{.975,158}=1.975 \qquad P=.0005 
\] \[
95\% \mbox{ Confidence Interval:} \quad -4.40 \pm 1.975(1.23) \quad \equiv \quad
-4.40 \pm 2.43 \quad \equiv \quad (-6.83,-1.97)
\]

For this example, the effect size is \(e=4.40/7.67=0.57\).

Idealized normal distributions for the Tai Chi and Firefighter Air
Consumption studies are given in Figure @ref(fig:ch0201) and Figure
@ref(fig:ch0201a).

\begin{figure}

{\centering \includegraphics[width=0.8\linewidth]{bookdownproj_STA4211Notes1_files/figure-latex/ch0201-1} 

}

\caption{Idealized normal distibutions for the Tai Chi study}\label{fig:ch0201}
\end{figure}

\begin{figure}

{\centering \includegraphics[width=0.8\linewidth]{bookdownproj_STA4211Notes1_files/figure-latex/ch0201a-1} 

}

\caption{Idealized normal distibutions for the Firefighter air consumption study}\label{fig:ch0201a}
\end{figure}

\[\nabla\]

There are various ways to parameterize single factor models. We will
consider the following three formulations in this chapter.

\begin{itemize}
\tightlist
\item
  \textbf{Cell Means Model} - Each treatment mean is its own parameter
  (no intercept term)
\item
  \textbf{Treatment Effect Model} - Intercept is the mean of the group
  means, each treatment has an effect relative to the mean, with effects
  summing to 0
\item
  \textbf{Reference Group Model} - Intercept corresponds to the mean for
  Group 1, remaining parameters are differences between remaining group
  means and group 1 mean
\end{itemize}

\section{Three Model Formulations}\label{three-model-formulations}

\subsection{Cell Means Model}\label{cell-means-model}

The model and assumptions are given below for the \textbf{Cell Means
Model}. Note that the total sample size is
\(n_T=n_1+\cdots+n_r=\sum_{i=1}^rn_i\). When all sample sizes are equal
\((n_i=n)\), the design is said to be \textbf{balanced}.

\[ Y_{ij}=\mu_i + \epsilon_{ij} 
\qquad i=1,\ldots,r; \quad j=1,\ldots, n_i \qquad \epsilon_{ij} \sim NID\left(0,\sigma^2\right) \]

\[ E\{Y_{ij}\} = E\{\mu_i + \epsilon_{ij}\} = \mu_i+0=\mu_i \qquad
V\{Y_{ij}\} = V\{\mu_i + \epsilon_{ij}\} = 0 +\sigma^2=\sigma^2 \]

For this model, the treatment/population mean for the \(i^{th}\)
condition is \(\mu_i\) and the variance of the individual measurements
within each condition is assumed to be \(\sigma^2\) (standard deviation
= \(\sigma\)) . Further, the individual measurements within each
condition are assumed to be independent and normally distributed.

The point estimators for the population means, which are derived below,
are the sample means for the various groups. The estimator for the
population variance is the pooled variance among the individual groups,
which is the same for all three model forms.

\subsection{Treatment Effects Model}\label{treatment-effects-model}

Let \(\mu_i\) be the mean for treatment \(i\), and let \(\mu_{\bullet}\)
be the unweighted mean of the \(\mu_i^s\).

\[ 
\mu_{\bullet}=\frac{\sum_{i=1}^r\mu_i}{r}  \qquad
\mu_i = \mu_{\bullet} + \left(\mu_i-\mu_{\bullet}\right) = \mu_{\bullet} + \tau_i
\quad \Rightarrow \quad \tau_i = \mu_i-\mu_{\bullet}
\] \[
\mu_1=\cdots=\mu_r=\mu_{\bullet} \quad \Rightarrow \quad
\tau_1=\cdots=\tau_r = 0
\]

\[
Y_{ij}=\mu_{\bullet} +\tau_i + \epsilon_{ij} 
\qquad i=1,\ldots,r; \quad j=1,\ldots, n_i \qquad \epsilon_{ij} \sim NID\left(0,\sigma^2\right) \]
\[ E\{Y_{ij}\} = E\{\mu_{\bullet} +\tau_i + \epsilon_{ij}\} = \mu_{\bullet} +\tau_i+0=\mu_{\bullet} +\tau_i \qquad
V\{Y_{ij}\} = V\{\mu_{\bullet} +\tau_i + \epsilon_{ij}\} = 0 +\sigma^2=\sigma^2 
\]

\[
\sum_{i=1}^{r}\tau_i=0 \quad \Rightarrow \quad \tau_r = -\sum_{i=1}^{r-1}\tau_i
\]

The point estimator for the overall population mean is the unweighted
mean of group sample means. The estimator for the effect of group \(i\)
is the difference between the sample mean for group \(i\) and the
overall unweighted mean.

\subsection{Reference Group Model}\label{reference-group-model}

This model, which is the default model for the \textbf{lm} and
\textbf{aov} functions in R, uses group 1 (based on alpha-numeric
ordering of the group levels) as the intercept. The remaining parameters
are the differences between the remaining \(r-1\) group means and the
mean for group 1. The point estimators are based on the corresponding
sample means.

An example where \(r=3\) conditions have means of \(\mu_1=40\),
\(\mu_2=50\), and \(\mu_3=60\), respectively and standard deviation
\(\sigma=10\) is given in Figure @ref(fig:ch0202).

\begin{figure}

{\centering \includegraphics[width=0.8\linewidth]{bookdownproj_STA4211Notes1_files/figure-latex/ch0202-1} 

}

\caption{Normal distibutions with means=40,50,60 and standard deviation=10}\label{fig:ch0202}
\end{figure}

\textbf{Example 2.3 - Virtual Training for a Lifeboat Launching Task} A
study in South Korea compared \(r=4\) methods of training to conduct a
lifeboat launching task (\citeproc{ref-jung_2018}{Jung and Ahn 2018}).
The treatments and their labels from the paper are given below and the
response \(Y\) was a procedural knowledge score for subjects post
training.

\begin{itemize}
\tightlist
\item
  \textbf{Lecture/Materials} - Traditional Lecture with no computer
  component (LEC/MAT)
\item
  \textbf{Monitor/Keyboard} - Trained virtually with a monitor,
  keyboard, and mouse (MON/KEY)
\item
  \textbf{Head-Mounted Display/Joypad} - Trained virtually with HMD and
  joypad (HMD/JOY)
\item
  \textbf{Head-Mounted Display/Wearable Sensors} - (HMD/WEA)
\end{itemize}

There were a total of \(n_T=64\) subjects and they were randomized so
that 16 subjects received each treatment
\(\left(n_1=n_2=n_3=n_4=16\right)\). Data that have been generated to
match the authors' reported means and standard deviations are given in
Table @ref(tab:ch0203). Note that when we later run this analysis in R,
there will be a variable for treatment and a variable for procedural
score. Figure @ref(fig:ch0203) gives the generated observed values
(dots), treatment means (green lines), and the overall mean (red line).

\begin{longtable}[]{@{}lrrrr@{}}
\caption{Lifeboat Virtual Training Data}\tabularnewline
\toprule\noalign{}
& LEC/MAT (i=1) & MON/KEY (i=2) & HMD/JOY (i=3) & HMD/WEA (i=4) \\
\midrule\noalign{}
\endfirsthead
\toprule\noalign{}
& LEC/MAT (i=1) & MON/KEY (i=2) & HMD/JOY (i=3) & HMD/WEA (i=4) \\
\midrule\noalign{}
\endhead
\bottomrule\noalign{}
\endlastfoot
j=1 & 4.5614 & 5.4532 & 1.8773 & 9.8257 \\
j=2 & 6.6593 & 8.9208 & 6.1884 & 8.9988 \\
j=3 & 5.6427 & 8.3489 & 6.8029 & 5.4112 \\
j=4 & 6.4394 & 8.3175 & 9.4535 & 4.9396 \\
j=5 & 4.8635 & 9.7848 & 8.7312 & 8.6486 \\
j=6 & 0.3268 & 6.2697 & 6.4198 & 5.1459 \\
j=7 & 5.6990 & 7.1327 & 8.4955 & 7.0656 \\
j=8 & 6.3545 & 7.0886 & 5.9008 & 7.9552 \\
j=9 & 6.7509 & 5.0733 & 3.3336 & 5.7284 \\
j=10 & 6.6019 & 8.3865 & 1.6530 & 6.6314 \\
j=11 & 3.2365 & 9.4227 & 9.4710 & 7.5014 \\
j=12 & 6.1655 & 8.9142 & 7.2818 & 8.2456 \\
j=13 & 2.4060 & 8.5150 & 9.7649 & 4.2465 \\
j=14 & 1.9851 & 6.8198 & 3.4931 & 3.2286 \\
j=15 & 5.2198 & 8.6390 & 8.9848 & 6.5520 \\
j=16 & 5.9765 & 6.2467 & 9.9261 & 9.8754 \\
n & 16.0000 & 16.0000 & 16.0000 & 16.0000 \\
Mean & 4.9306 & 7.7083 & 6.7361 & 6.8750 \\
SD & 1.9400 & 1.4300 & 2.8200 & 1.9900 \\
\end{longtable}

\begin{verbatim}
##   grp.trt procKnow
## 1       1   4.5614
## 2       1   6.6593
\end{verbatim}

\begin{figure}

{\centering \includegraphics[width=0.8\linewidth]{bookdownproj_STA4211Notes1_files/figure-latex/ch0203-1} 

}

\caption{Procedural knowledge scores for lifeboat training study}\label{fig:ch0203}
\end{figure}

Using the summary data from the table, we will obtain point estimates
and standard errors directly for the three model formulations directly,
then using R. First, we compute the pooled standard deviation (aka
residual standard error).

\[
s=\sqrt{\frac{\sum_{i=1}^r\left(n_i-1\right)s_i^2}{n_T-r}} =
\sqrt{\frac{(16-1)\left(1.94^2+1.43^2+2.82^2+1.99^2\right)}{64-4}} = 2.105
\]

\textbf{Cell Means Model}

\[
\hat{\mu}_1=\overline{y}_1=4.931 \quad 
\hat{\mu}_2=\overline{y}_2=7.708 \quad 
\hat{\mu}_3=\overline{y}_3=6.736 \quad 
\hat{\mu}_4=\overline{y}_4=6.875 \]

\[
\hat{SE}\left\{\hat{\mu}_i\right\}=\frac{s}{\sqrt{n_i}}=\frac{2.105}{\sqrt{16}}=0.526
\qquad i=1,2,3,4
\]

\textbf{Treatment Effects Model}

\[
\hat{\mu}=\frac{4.931+7.708+6.736+6.875}{4}=6.563 \]

\[
\hat{SE}\left\{\hat{\mu}\right\}=s\sqrt{\frac{1}{r^2}\sum_{i=1}^r\frac{1}{n_i}} =
2.105\sqrt{\frac{1}{4^2}4\left(\frac{1}{16}\right)}=0.263
\] \[
\hat{\tau}_1=4.931-6.563=-1.632
\quad \hat{\tau}_2=1.145
\quad \hat{\tau}_3=0.173
\quad \hat{\tau}_4=0.312 \]

To obtain the standard error for \(\hat{\tau}_k\), we can re-write it as
follows and make use of independence among group means.

\[
\hat{\tau}_k = \overline{y}_k-\hat{\mu} =
\overline{y}_k - \frac{1}{r}\sum_{i=1}^r\overline{y}_i
=\frac{r-1}{r} \overline{y}_k - \frac{1}{r}\sum_{i=1 \atop i \not=k}^r\overline{y}_i
\]

\[
\hat{SE}\left\{\hat{\tau}_k\right\} =
s\sqrt{\frac{1}{r^2}\left(\frac{(r-1)^2}{n_k}+
\sum_{i=1 \atop i \not=k}^r\frac{1}{n_i}\right)} =
2.105\sqrt{\frac{1}{4^2}\left(\frac{(4-1)^2}{16}+3\frac{1}{16}\right)}
=0.456 
\]

\textbf{Reference Group Model}

\[
\hat{\mu}_1 = \overline{y}_1 = 4.931 \qquad \hat{SE}\left\{\hat{\mu}_1\right\}=
\frac{s}{\sqrt{n_1}}=\frac{2.105}{\sqrt{16}}=0.526
\]

\[
i=2,\ldots, r \quad \hat{\tau}_i = \overline{y}_i-\overline{y}_1
\qquad \hat{SE}\left\{\hat{\tau}_i\right\}=s \sqrt{\frac{1}{n_i}+\frac{1}{n_1}}
\] \[
\hat{\tau}_2=7.708-4.931=2.777 \quad
\hat{\tau}_3=6.736-4.931=1.805 \quad
\hat{\tau}_4=6.875-4.931=1.944
\] \[
\hat{SE}\left\{\hat{\tau}_i\right\}=2.105 \sqrt{\frac{1}{16}+\frac{1}{16}} =
0.744
\]

\begin{Shaded}
\begin{Highlighting}[]
\NormalTok{vt }\OtherTok{\textless{}{-}} \FunctionTok{read.csv}\NormalTok{(}\StringTok{"data/virtual\_training.csv"}\NormalTok{)}
\FunctionTok{head}\NormalTok{(vt,}\DecValTok{2}\NormalTok{)}
\end{Highlighting}
\end{Shaded}

\begin{verbatim}
##   grp.trt procKnow
## 1       1   4.5614
## 2       1   6.6593
\end{verbatim}

\begin{Shaded}
\begin{Highlighting}[]
\DocumentationTok{\#\# Cell Means Model (Separate means for each group, no intercept)}
\NormalTok{vt.mod1 }\OtherTok{\textless{}{-}} \FunctionTok{lm}\NormalTok{(procKnow }\SpecialCharTok{\textasciitilde{}} \FunctionTok{factor}\NormalTok{(grp.trt) }\SpecialCharTok{{-}} \DecValTok{1}\NormalTok{, }\AttributeTok{data=}\NormalTok{vt)}
\FunctionTok{summary}\NormalTok{(vt.mod1)}
\end{Highlighting}
\end{Shaded}

\begin{verbatim}
## 
## Call:
## lm(formula = procKnow ~ factor(grp.trt) - 1, data = vt)
## 
## Residuals:
##     Min      1Q  Median      3Q     Max 
## -5.0831 -1.4444  0.5774  1.5495  3.1900 
## 
## Coefficients:
##                  Estimate Std. Error t value Pr(>|t|)    
## factor(grp.trt)1   4.9306     0.5262    9.37 2.37e-13 ***
## factor(grp.trt)2   7.7083     0.5262   14.65  < 2e-16 ***
## factor(grp.trt)3   6.7361     0.5262   12.80  < 2e-16 ***
## factor(grp.trt)4   6.8750     0.5262   13.06  < 2e-16 ***
## ---
## Signif. codes:  0 '***' 0.001 '**' 0.01 '*' 0.05 '.' 0.1 ' ' 1
## 
## Residual standard error: 2.105 on 60 degrees of freedom
## Multiple R-squared:  0.9139, Adjusted R-squared:  0.9082 
## F-statistic: 159.2 on 4 and 60 DF,  p-value: < 2.2e-16
\end{verbatim}

\begin{Shaded}
\begin{Highlighting}[]
\DocumentationTok{\#\# Treatment Effects Model}
\FunctionTok{options}\NormalTok{(}\AttributeTok{contrasts=}\FunctionTok{c}\NormalTok{(}\StringTok{"contr.sum"}\NormalTok{,}\StringTok{"contr.poly"}\NormalTok{))}
\NormalTok{vt.mod2 }\OtherTok{\textless{}{-}} \FunctionTok{lm}\NormalTok{(procKnow }\SpecialCharTok{\textasciitilde{}} \FunctionTok{factor}\NormalTok{(grp.trt), }\AttributeTok{data=}\NormalTok{vt)}
\FunctionTok{summary}\NormalTok{(vt.mod2)}
\end{Highlighting}
\end{Shaded}

\begin{verbatim}
## 
## Call:
## lm(formula = procKnow ~ factor(grp.trt), data = vt)
## 
## Residuals:
##     Min      1Q  Median      3Q     Max 
## -5.0831 -1.4444  0.5774  1.5495  3.1900 
## 
## Coefficients:
##                  Estimate Std. Error t value Pr(>|t|)    
## (Intercept)        6.5625     0.2631  24.943  < 2e-16 ***
## factor(grp.trt)1  -1.6319     0.4557  -3.581 0.000686 ***
## factor(grp.trt)2   1.1458     0.4557   2.514 0.014622 *  
## factor(grp.trt)3   0.1736     0.4557   0.381 0.704573    
## ---
## Signif. codes:  0 '***' 0.001 '**' 0.01 '*' 0.05 '.' 0.1 ' ' 1
## 
## Residual standard error: 2.105 on 60 degrees of freedom
## Multiple R-squared:  0.1981, Adjusted R-squared:  0.158 
## F-statistic: 4.941 on 3 and 60 DF,  p-value: 0.003931
\end{verbatim}

\begin{Shaded}
\begin{Highlighting}[]
\DocumentationTok{\#\# Reference Group Model (Default in R)}
\FunctionTok{options}\NormalTok{(}\AttributeTok{contrasts=}\FunctionTok{c}\NormalTok{(}\StringTok{"contr.treatment"}\NormalTok{,}\StringTok{"contr.poly"}\NormalTok{))}
\NormalTok{vt.mod3 }\OtherTok{\textless{}{-}} \FunctionTok{lm}\NormalTok{(procKnow }\SpecialCharTok{\textasciitilde{}} \FunctionTok{factor}\NormalTok{(grp.trt), }\AttributeTok{data=}\NormalTok{vt)}
\FunctionTok{summary}\NormalTok{(vt.mod3)}
\end{Highlighting}
\end{Shaded}

\begin{verbatim}
## 
## Call:
## lm(formula = procKnow ~ factor(grp.trt), data = vt)
## 
## Residuals:
##     Min      1Q  Median      3Q     Max 
## -5.0831 -1.4444  0.5774  1.5495  3.1900 
## 
## Coefficients:
##                  Estimate Std. Error t value Pr(>|t|)    
## (Intercept)        4.9306     0.5262   9.370 2.37e-13 ***
## factor(grp.trt)2   2.7778     0.7442   3.733 0.000423 ***
## factor(grp.trt)3   1.8056     0.7442   2.426 0.018277 *  
## factor(grp.trt)4   1.9444     0.7442   2.613 0.011329 *  
## ---
## Signif. codes:  0 '***' 0.001 '**' 0.01 '*' 0.05 '.' 0.1 ' ' 1
## 
## Residual standard error: 2.105 on 60 degrees of freedom
## Multiple R-squared:  0.1981, Adjusted R-squared:  0.158 
## F-statistic: 4.941 on 3 and 60 DF,  p-value: 0.003931
\end{verbatim}

\[\nabla\]

\section{\texorpdfstring{The Analysis of Variance and
\(F\)-test}{The Analysis of Variance and F-test}}\label{the-analysis-of-variance-and-f-test}

In this section, we obtain a decomposition of the total variation of the
individual measurements around the overall mean, the \textbf{Total Sum
of Squares}, which we will denote \(SSTO\). The \textbf{Degrees of
Freedom} associated with the total sum of squares is \(df_{TO}=n_T-1\).

\[
\overline{Y}_{\bullet\bullet}=
\frac{\sum_{i=1}^r\sum_{j=1}^{n_i}Y_{ij}}{n_T}
\qquad \qquad
SSTO = \sum_{i=1}^r\sum_{j=1}^{n_i}
  \left(Y_{ij}-\overline{Y}_{\bullet\bullet}\right)^2
\qquad  df_{TO}=n_T-1
\]

The total sum of squares can be partitioned into the \textbf{Error Sum
of Squares} and the \textbf{Treatment Sum of Squares}. The error sum of
squares can be thought of as how variable are measurements within the
same treatment or group. The treatment sum of squares represents how
variable the treatment or group means are around the overall mean
(weighted by their sample sizes). The sums of squares are given below
(note that in R, for linear models, the error sum of squares is labelled
as \(RSS\), for the residual sum of squares). As with linear regression
models, the error sum of squares is the sum of the squared distances
from the observed values to their predicted values.

\[
SSE = \sum_{i=1}^r\sum_{j=1}^{n_i}  \left(Y_{ij}-\hat{Y}_{ij}\right)^2
\] Noting that for this model, \(\hat{Y}_{ij}=\overline{Y}_{i\bullet}\),
we obtain the following formula for \(SSE\). The degrees of freedom for
the error sum of squares is \(df_E=n_T-r\).

\[
SSE=\sum_{i=1}^r\sum_{j=1}^{n_i} \left(Y_{ij}-\overline{Y}_{i\bullet}\right)^2
\qquad df_E=n_T-r
\]

Another useful way of computing \(SSE\) based on published summary
statistics is given below. Suppose the authors have published the sample
sizes \(\left\{n_i\right\}\), the sample means
\(\left\{\overline{y}_{i\bullet}\right\}\) using lower case \(y\) as
these are observed means from a realization of this experiment, and the
sample standard deviations \(\left\{s_i\right\}\) for the \(r\)
treatments. Then \(SSE\) can be computed as below.

\[ 
s_i = \sqrt{\frac{\sum_{j=1}^{n_i}\left(y_{ij}-\overline{y}_{i\bullet}\right)^2}{n_i-1}}
\quad \Rightarrow \quad
SSE = \sum_{i=1}^r\left(n_i-1\right)s_i^2 
\]

The \textbf{Treatment Sum of Squares}, denoted as \(SSTR\), acts like
the regression sum of squares for linear regression models. That is,
\(SSTR\) represents the sum of squared differences from the predicted
values to the overall mean. The degrees of freedom for the treatment sum
of squares is \(df_{TR}=r-1\). Keeping in mind that the fitted value for
\(Y_{ij}\) is \(\hat{Y_{ij}} = \overline{Y}_{i\bullet}\), we obtain SSTR
as follows.

\[
SSTR = \sum_{i=1}^r\sum_{j=1}^{n_i}
\left(\overline{Y}_{i\bullet}-
  \overline{Y}_{\bullet\bullet}\right)^2
 = \sum_{i=1}^r n_i \left(\overline{Y}_{i\bullet}-
  \overline{Y}_{\bullet\bullet}\right)^2 
\qquad df_{TR}=r-1
\]

For each sum of squares, the \textbf{Mean Square} is the sum of squares
divided by its corresponding degrees of freedom, and are used for
inferences among the means.

In single factor models, the goal is typically to test for differences
among the treatment or group means. The null hypothesis is that the
means are all equal, while the alternative hypothesis is that there are
differences among the means. In the next chapter, we consider specific
comparisons among treatments.

\[
H_0: \mu_1=\cdots = \mu_r \quad \Rightarrow \quad \tau_1=\cdots = \tau_r = 0
\] The test statistic is the \(F\)-statistic from the Analysis of
Variance and will be used in many forms throughout this course. Under
the null hypothesis, the \(F\)-statistic follows the \(F\)-distribution,
based on its numerator and denominator degrees of freedom. Large values
of the \(F\)-statistic are evidence against the null hypothesis.
Technical details of the test are given in Section 2.4.

\[
\mbox{Test Statistic:} \quad F^* = 
\frac{\left[\frac{SSTR}{df_{TR}}\right]}
{\left[\frac{SSE}{df_{E}}\right]}
\quad = \quad \frac{MSTR}{MSE} 
\]

\[
\mbox{Rejection Region:} \quad F^* \geq F_{1-\alpha;r-1,n_T-r}  \qquad
P=P\left(F_{r-1,n_T-r} \geq F^*\right)
\]

The \textbf{effect size} for the One-Way ANOVA model is
\(\eta^2=SSTR/SSTO\), that is, the fraction of the total variation that
is ``attributable'' to differences in the group means.

\textbf{Example 2.4 - Virtual Training for a Lifeboat Launching Task}

\[
n_1=n_2=n_3=n_4=16 \qquad \overline{y}_{1\bullet}=4.931 \quad
\overline{y}_{2\bullet}=7.708 \quad
\overline{y}_{3\bullet}=6.736 \quad
\overline{y}_{4\bullet}=6.875 \quad
\overline{y}_{\bullet\bullet}=6.563
\]

\[
SSTR=16\left[(4.931-6.563)^2+(7.708-6.563)^2+(6.736-6.563)^2+(6.875-6.563)^2\right]
=65.628
\]

\[
s_1=1.94 \quad s_2=1.43 \quad s_3=2.82 \quad s_4=1.99
\] \[
\quad \Rightarrow \quad SSE=(16-1)\left(1.94^2+1.43^2+2.82^2+1.99^2\right)=265.815
\]

Now, we test for differences among the effects of the \(r=4\) teaching
methods. The degrees of freedom for treatments is \(df_{TR}=r-1=3\) and
for error is \(df_E=n_T-r=60\).

\[
H_0: \mu_1=\cdots = \mu_4 \quad \Rightarrow \quad \tau_1=\cdots = \tau_4 = 0
\] \[
\mbox{Test Statistic:} \quad F^* = 
\frac{\left[\frac{SSTR}{df_{TR}}\right]}
{\left[\frac{SSE}{df_{E}}\right]}
=\frac{\left[\frac{65.789}{3}\right]}
{\left[\frac{265.815}{60}\right]}
= \frac{21.930}{4.430} = 4.950
\] \[
\mbox{Rejection Region:} \quad F^* \geq F_{.95;3,60}=2.578  \qquad
P=P\left(F_{3,60} \geq 4.950\right) = .0039
\] \[
\eta^2=\frac{SSTR}{SSTO}=\frac{SSTR}{SSTR+SSE}=\frac{65.789}{65.789+265.815}=
\frac{65.789}{331.604}=0.1984
\] Approximately 20\% of the total variation is due to differences in
the training methods.

We obtain the Sums of Squares for Treatments and Error below using the
\textbf{anova} of the \textbf{lm} model fit. A warning involves the fact
that when using the cell means model (with no intercept) the Treatment
Sum of Squares and Degrees of Freedom are incorrect here. The Treatment
Sum of Squares in this form also includes the correction for the mean,
\(n_T\overline{y}_{\bullet\bullet}^2\). We are testing
\(H_0:\mu_1=\mu_2=\mu_3=\mu_4=0\), which is a much stronger hypothesis
than testing \(H_0:\mu_1=\mu_2=\mu_3=\mu_4\). We will use the Treatment
Effects model or the Reference Group model function to obtain the
correct degrees of freedom and \(F\)-statistic below.

\begin{Shaded}
\begin{Highlighting}[]
\NormalTok{vt }\OtherTok{\textless{}{-}} \FunctionTok{read.csv}\NormalTok{(}\StringTok{"data/virtual\_training.csv"}\NormalTok{)}
\FunctionTok{head}\NormalTok{(vt,}\DecValTok{2}\NormalTok{)}
\end{Highlighting}
\end{Shaded}

\begin{verbatim}
##   grp.trt procKnow
## 1       1   4.5614
## 2       1   6.6593
\end{verbatim}

\begin{Shaded}
\begin{Highlighting}[]
\NormalTok{vt.mod1 }\OtherTok{\textless{}{-}} \FunctionTok{lm}\NormalTok{(procKnow }\SpecialCharTok{\textasciitilde{}} \FunctionTok{factor}\NormalTok{(grp.trt) }\SpecialCharTok{{-}} \DecValTok{1}\NormalTok{, }\AttributeTok{data=}\NormalTok{vt)}
\FunctionTok{summary}\NormalTok{(vt.mod1)}
\end{Highlighting}
\end{Shaded}

\begin{verbatim}
## 
## Call:
## lm(formula = procKnow ~ factor(grp.trt) - 1, data = vt)
## 
## Residuals:
##     Min      1Q  Median      3Q     Max 
## -5.0831 -1.4444  0.5774  1.5495  3.1900 
## 
## Coefficients:
##                  Estimate Std. Error t value Pr(>|t|)    
## factor(grp.trt)1   4.9306     0.5262    9.37 2.37e-13 ***
## factor(grp.trt)2   7.7083     0.5262   14.65  < 2e-16 ***
## factor(grp.trt)3   6.7361     0.5262   12.80  < 2e-16 ***
## factor(grp.trt)4   6.8750     0.5262   13.06  < 2e-16 ***
## ---
## Signif. codes:  0 '***' 0.001 '**' 0.01 '*' 0.05 '.' 0.1 ' ' 1
## 
## Residual standard error: 2.105 on 60 degrees of freedom
## Multiple R-squared:  0.9139, Adjusted R-squared:  0.9082 
## F-statistic: 159.2 on 4 and 60 DF,  p-value: < 2.2e-16
\end{verbatim}

\begin{Shaded}
\begin{Highlighting}[]
\FunctionTok{anova}\NormalTok{(vt.mod1)    }\DocumentationTok{\#\# Incorrect df, SS, and MS for Treatments and F{-}stat}
\end{Highlighting}
\end{Shaded}

\begin{verbatim}
## Analysis of Variance Table
## 
## Response: procKnow
##                 Df  Sum Sq Mean Sq F value    Pr(>F)    
## factor(grp.trt)  4 2821.91  705.48  159.24 < 2.2e-16 ***
## Residuals       60  265.81    4.43                      
## ---
## Signif. codes:  0 '***' 0.001 '**' 0.01 '*' 0.05 '.' 0.1 ' ' 1
\end{verbatim}

\begin{Shaded}
\begin{Highlighting}[]
\NormalTok{vt.mod2 }\OtherTok{\textless{}{-}} \FunctionTok{aov}\NormalTok{(procKnow }\SpecialCharTok{\textasciitilde{}} \FunctionTok{factor}\NormalTok{(grp.trt), }\AttributeTok{data=}\NormalTok{vt)}
\FunctionTok{anova}\NormalTok{(vt.mod2)}
\end{Highlighting}
\end{Shaded}

\begin{verbatim}
## Analysis of Variance Table
## 
## Response: procKnow
##                 Df  Sum Sq Mean Sq F value   Pr(>F)   
## factor(grp.trt)  3  65.664 21.8880  4.9406 0.003931 **
## Residuals       60 265.815  4.4302                    
## ---
## Signif. codes:  0 '***' 0.001 '**' 0.01 '*' 0.05 '.' 0.1 ' ' 1
\end{verbatim}

\begin{Shaded}
\begin{Highlighting}[]
\FunctionTok{options}\NormalTok{(}\AttributeTok{es.use\_symbols=}\ConstantTok{TRUE}\NormalTok{)}

\FunctionTok{eta\_squared}\NormalTok{(vt.mod1)}
\end{Highlighting}
\end{Shaded}

\begin{verbatim}
## For one-way between subjects designs, partial eta squared is equivalent
##   to eta squared. Returning eta squared.
\end{verbatim}

\begin{verbatim}
## # Effect Size for ANOVA
## 
## Parameter       |   η² |       95% CI
## -------------------------------------
## factor(grp.trt) | 0.91 | [0.88, 1.00]
## 
## - One-sided CIs: upper bound fixed at [1.00].
\end{verbatim}

\begin{Shaded}
\begin{Highlighting}[]
\FunctionTok{eta\_squared}\NormalTok{(vt.mod2)}
\end{Highlighting}
\end{Shaded}

\begin{verbatim}
## For one-way between subjects designs, partial eta squared is equivalent
##   to eta squared. Returning eta squared.
\end{verbatim}

\begin{verbatim}
## # Effect Size for ANOVA
## 
## Parameter       |   η² |       95% CI
## -------------------------------------
## factor(grp.trt) | 0.20 | [0.05, 1.00]
## 
## - One-sided CIs: upper bound fixed at [1.00].
\end{verbatim}

Note that for the cell means model (vt.mod1) the sum of squares for
treatments includes the ``correction for the mean,'' thus it is larger
than \(SSTR\) by \(n_T\overline{y}_{\bullet\bullet}^2\). The effect size
for that model is also too large
\(\eta^2=2821.91/(2821.91+265.81)=0.91\).

Here we will conduct the analysis directly from the raw data, computing
\(\left\{n_i\right\}\), \(\left\{\overline{y}_{i\bullet}\right\}\), and
\(\left\{s_i\right\}\) as well as sums of squares and degrees of
freedom.

\begin{Shaded}
\begin{Highlighting}[]
\NormalTok{vt }\OtherTok{\textless{}{-}} \FunctionTok{read.csv}\NormalTok{(}\StringTok{"data/virtual\_training.csv"}\NormalTok{)}
\FunctionTok{head}\NormalTok{(vt,}\DecValTok{2}\NormalTok{)}
\end{Highlighting}
\end{Shaded}

\begin{verbatim}
##   grp.trt procKnow
## 1       1   4.5614
## 2       1   6.6593
\end{verbatim}

\begin{Shaded}
\begin{Highlighting}[]
\FunctionTok{attach}\NormalTok{(vt)}

\DocumentationTok{\#\#\#\# "Brute{-}Force" Calculations  }
\NormalTok{n.trt }\OtherTok{\textless{}{-}} \FunctionTok{as.vector}\NormalTok{(}\FunctionTok{tapply}\NormalTok{(procKnow, grp.trt, length))   }\DocumentationTok{\#\# Treatment sample sizes}
\NormalTok{mean.trt }\OtherTok{\textless{}{-}} \FunctionTok{as.vector}\NormalTok{(}\FunctionTok{tapply}\NormalTok{(procKnow, grp.trt, mean))  }\DocumentationTok{\#\# Treatment means}
\NormalTok{sd.trt }\OtherTok{\textless{}{-}} \FunctionTok{as.vector}\NormalTok{(}\FunctionTok{tapply}\NormalTok{(procKnow, grp.trt, sd))      }\DocumentationTok{\#\# Treatment SDs}
\NormalTok{mean.all }\OtherTok{\textless{}{-}} \FunctionTok{mean}\NormalTok{(procKnow)                              }\DocumentationTok{\#\# Overall Mean}
\NormalTok{r }\OtherTok{\textless{}{-}} \FunctionTok{length}\NormalTok{(n.trt)                                      }\DocumentationTok{\#\# Number of Treatments}
\NormalTok{n\_T }\OtherTok{\textless{}{-}} \FunctionTok{sum}\NormalTok{(n.trt)                                       }\DocumentationTok{\#\# Overall Mean}

\DocumentationTok{\#\#\#\# Sums of Squares, degrees of freedom, Mean Squares, F{-}test}
\NormalTok{SSTR }\OtherTok{\textless{}{-}} \FunctionTok{sum}\NormalTok{(n.trt }\SpecialCharTok{*}\NormalTok{ (mean.trt }\SpecialCharTok{{-}}\NormalTok{ mean.all)}\SpecialCharTok{\^{}}\DecValTok{2}\NormalTok{);    dfTR }\OtherTok{\textless{}{-}}\NormalTok{ r}\DecValTok{{-}1}      
\NormalTok{SSE }\OtherTok{\textless{}{-}} \FunctionTok{sum}\NormalTok{((n.trt}\DecValTok{{-}1}\NormalTok{) }\SpecialCharTok{*}\NormalTok{ sd.trt}\SpecialCharTok{\^{}}\DecValTok{2}\NormalTok{);                dfE }\OtherTok{\textless{}{-}}\NormalTok{ n\_T}\SpecialCharTok{{-}}\NormalTok{r}
\NormalTok{SSTO }\OtherTok{\textless{}{-}}\NormalTok{ SSTR }\SpecialCharTok{+}\NormalTok{ SSE;                              dfTO }\OtherTok{\textless{}{-}}\NormalTok{ n\_T}\DecValTok{{-}1}
\NormalTok{MSTR }\OtherTok{\textless{}{-}}\NormalTok{ SSTR }\SpecialCharTok{/}\NormalTok{ dfTR}
\NormalTok{MSE }\OtherTok{\textless{}{-}}\NormalTok{ SSE }\SpecialCharTok{/}\NormalTok{ dfE}
\NormalTok{F\_star }\OtherTok{\textless{}{-}}\NormalTok{ MSTR }\SpecialCharTok{/}\NormalTok{ MSE}
\NormalTok{F}\FloatTok{.95} \OtherTok{\textless{}{-}} \FunctionTok{qf}\NormalTok{(.}\DecValTok{95}\NormalTok{,dfTR,dfE)}
\NormalTok{pF }\OtherTok{\textless{}{-}} \DecValTok{1} \SpecialCharTok{{-}} \FunctionTok{pf}\NormalTok{(F\_star, dfTR, dfE)}
\NormalTok{eta2 }\OtherTok{\textless{}{-}}\NormalTok{ SSTR}\SpecialCharTok{/}\NormalTok{SSTO}

\DocumentationTok{\#\#\#\# Set{-}up output}
\NormalTok{trt.out }\OtherTok{\textless{}{-}} \FunctionTok{cbind}\NormalTok{(dfTR, SSTR, MSTR, F\_star, F}\FloatTok{.95}\NormalTok{, pF, eta2)}
\NormalTok{err.out }\OtherTok{\textless{}{-}} \FunctionTok{cbind}\NormalTok{(dfE,  SSE,  MSE,  }\ConstantTok{NA}\NormalTok{,     }\ConstantTok{NA}\NormalTok{,   }\ConstantTok{NA}\NormalTok{, }\ConstantTok{NA}\NormalTok{)}
\NormalTok{tot.out }\OtherTok{\textless{}{-}} \FunctionTok{cbind}\NormalTok{(dfTO, SSTO, }\ConstantTok{NA}\NormalTok{,   }\ConstantTok{NA}\NormalTok{,     }\ConstantTok{NA}\NormalTok{,   }\ConstantTok{NA}\NormalTok{, }\ConstantTok{NA}\NormalTok{)}
\NormalTok{aov.out }\OtherTok{\textless{}{-}} \FunctionTok{rbind}\NormalTok{(trt.out, err.out, tot.out)}
\FunctionTok{colnames}\NormalTok{(aov.out) }\OtherTok{\textless{}{-}} \FunctionTok{c}\NormalTok{(}\StringTok{"df"}\NormalTok{, }\StringTok{"SS"}\NormalTok{, }\StringTok{"MS"}\NormalTok{, }\StringTok{"F*"}\NormalTok{, }\StringTok{"F(.95)"}\NormalTok{, }\StringTok{"P(\textgreater{}F*)"}\NormalTok{, }\StringTok{"eta2"}\NormalTok{)}
\FunctionTok{rownames}\NormalTok{(aov.out) }\OtherTok{\textless{}{-}} \FunctionTok{c}\NormalTok{(}\StringTok{"Treatment"}\NormalTok{, }\StringTok{"Error"}\NormalTok{, }\StringTok{"Total"}\NormalTok{)}
\FunctionTok{round}\NormalTok{(aov.out,}\DecValTok{4}\NormalTok{)}
\end{Highlighting}
\end{Shaded}

\begin{verbatim}
##           df       SS      MS     F* F(.95) P(>F*)   eta2
## Treatment  3  65.6639 21.8880 4.9406 2.7581 0.0039 0.1981
## Error     60 265.8149  4.4302     NA     NA     NA     NA
## Total     63 331.4788      NA     NA     NA     NA     NA
\end{verbatim}

\begin{Shaded}
\begin{Highlighting}[]
\FunctionTok{detach}\NormalTok{(vt)}
\end{Highlighting}
\end{Shaded}

\[\nabla\]

\subsection{General Linear Test
Approach}\label{general-linear-test-approach}

We begin with \(r\) treatments with no restrictions on their population
means \(\left\{\mu_i\right\}\). We then place a restriction on the
population means. Both models are fit, and the Error Sums of Squares and
Degrees of Freedom are obtained for each model.

\textbf{Example 2.5 - Virtual Training for a Lifeboat Launching Task}

Consider the following restrictions in light of the virtual training
experiment.

\begin{itemize}
\tightlist
\item
  No treatment differences among the 4 treatments:
  \(\mu_1=\mu_2=\mu_3=\mu_4\)
\item
  No treatment differences among the 3 virtual reality treatments:
  \(\mu_2=\mu_3=\mu_4\)
\item
  No treatment differences among the 2 Head-Mounted Display treatments:
  \(\mu_3=\mu_4\)
\end{itemize}

All of these can be tested using a general linear test. Note that the
first test is the same as the \(F\)-test conducted from the Analysis of
Variance.

The form of the test is as follows. Let the Complete model be the model
with all \(r\) means allowed to be different values. In the virtual
reality case we have the following numbers of restrictions.

\begin{itemize}
\tightlist
\item
  \(H_0:\mu_1=\mu_2=\mu_3=\mu_4\) - There are 3 restrictions
\item
  \(H_0:\mu_1, \quad \mu_2=\mu_3=\mu_4\) - There are 2 restrictions
\item
  \(H_0:\mu_1, \mu_2, \quad \mu_3=\mu_4\) - There is 1 restriction
\end{itemize}

Once we fit the Complete (No restriction) and Reduced (Restricted)
models, we can obtain the Test Statistic as follows. Start with the
first restriction: \(\mu_1=\mu_2=\mu_3=\mu_4=\mu_c\) where \(\mu_c\) is
the common mean.

\[
\mbox{Reduced Model: (1 Mean)} \qquad \hat{\mu}_c = \overline{Y}_{\bullet\bullet} = \hat{Y}_{ij}(R)
\]

\[
  \quad \Rightarrow \quad SSE(R)=\sum_{i=1}^r\sum_{j=1}^{n_i}\left(Y_{ij}-\hat{Y}_{ij}(R)\right)^2
  = \sum_{i=1}^r\sum_{j=1}^{n_i}\left(Y_{ij}-\overline{Y}_{\bullet\bullet}\right)^2 = SSTO
  \qquad df_R = n_T-1
  \]

\[
\mbox{Complete Model: (4 Means)} \qquad \hat{\mu}_i = \overline{Y}_{i\bullet}=\hat{Y}_{ij}(C)
  \] \[
  \quad \Rightarrow \quad SSE(C)=\sum_{i=1}^r\sum_{j=1}^{n_i}\left(Y_{ij}-\hat{Y}_{ij}(C)\right)^2
  = \sum_{i=1}^r\sum_{j=1}^{n_i}\left(Y_{ij}-\overline{Y}_{i\bullet}\right)^2 = SSE
  \qquad df_C = n_T-r
  \]

The general linear \(F\)-test is conducted as follows once the previous
models have been fit.

\[
  F^*=\frac{\left[\frac{SSE(R)-SSE(C)}{df_R-df_C}\right]}{\left[\frac{SSE(C)}{df_C}\right]}
  \qquad \mbox{Reject } H_0 \mbox{ if: } \quad F^* \geq F_{1-\alpha,df_R-df_C, df_C}
  \]

In the case of the test of \(H_0:\mu_1=\mu_2=\cdots = \mu_r\), we obtain
the following result where we reproduce the ANOVA \(F\)-test.

\[
SSE(R)=SSTO \quad SSE(C) = SSE \quad \Rightarrow \quad SSE(R)-SSE(C)=SSTR
\]

\[ 
df_R = n_T-1 \quad df_C = n_T-r \quad \Rightarrow \quad df_R-df_C =
\left(n_T-1\right)-\left(n_T-r\right) = r-1
\]

\[
 F^*=\frac{\left[\frac{SSE(R)-SSE(C)}{df_R-df_C}\right]}{\left[\frac{SSE(C)}{df_C}\right]}=
 \frac{\left[\frac{SSTR}{r-1}\right]}{\left[\frac{SSE}{n_T-r}\right]}=\frac{MSTR}{MSE}
\]

For the second model described above (equality of virtual reality
means), the reduced model would have a single mean for the LEC/MAT
control group and common means for the three virtual reality treatments.

\[
 \hat{Y}_{1j}=\overline{Y}_{1\bullet}  \qquad \qquad 
 \hat{Y}_{2j}=\hat{Y}_{3j}=\hat{Y}_{4j}=\frac{\sum_{i=2}^4\sum_{j=1}^{n_i}Y_{ij}}{n_2+n_3+n_4} = \frac{\sum_{i=2}^4n_i\overline{Y}_{i\bullet}}{\sum_{i=2}^4 n_i}
\]

The third model involves leaving treatments 1 and 2 as individual groups
and combines treatments 3 and 4 into a single group.

\[
 \hat{Y}_{1j}=\overline{Y}_{1\bullet}  \qquad
 \hat{Y}_{2j}=\overline{Y}_{2\bullet}  \qquad
 \hat{Y}_{3j}=\hat{Y}_{4j}=  \frac{\sum_{i=3}^4n_i\overline{Y}_{i\bullet}}{\sum_{i=3}^4 n_i}
\] We will fit these models in R and use the \textbf{anova} command to
conduct the \(F\)-tests from the various model fits. The trick is to
create new treatment factors for the various restrictions, and fit an
\textbf{lm} or \textbf{aov} object.

\begin{Shaded}
\begin{Highlighting}[]
\NormalTok{vt }\OtherTok{\textless{}{-}} \FunctionTok{read.csv}\NormalTok{(}\StringTok{"data/virtual\_training.csv"}\NormalTok{)}
\FunctionTok{head}\NormalTok{(vt,}\DecValTok{2}\NormalTok{)}
\end{Highlighting}
\end{Shaded}

\begin{verbatim}
##   grp.trt procKnow
## 1       1   4.5614
## 2       1   6.6593
\end{verbatim}

\begin{Shaded}
\begin{Highlighting}[]
\DocumentationTok{\#\# Create grouping variable that separates control (1) vs virtual (234)}
\NormalTok{vt}\SpecialCharTok{$}\NormalTok{grp.trt\_1\_234 }\OtherTok{\textless{}{-}} \FunctionTok{ifelse}\NormalTok{(vt}\SpecialCharTok{$}\NormalTok{grp.trt}\SpecialCharTok{==}\DecValTok{1}\NormalTok{,}\DecValTok{1}\NormalTok{,}\DecValTok{234}\NormalTok{)}
\DocumentationTok{\#\# Create grouping variable that separates control (1) vs MON/KEY (2) vs HMD (34)}
\NormalTok{vt}\SpecialCharTok{$}\NormalTok{grp.trt\_1\_2\_34 }\OtherTok{\textless{}{-}} \FunctionTok{ifelse}\NormalTok{(vt}\SpecialCharTok{$}\NormalTok{grp.trt}\SpecialCharTok{==}\DecValTok{1}\NormalTok{,}\DecValTok{1}\NormalTok{,}\FunctionTok{ifelse}\NormalTok{(vt}\SpecialCharTok{$}\NormalTok{grp.trt}\SpecialCharTok{==}\DecValTok{2}\NormalTok{,}\DecValTok{2}\NormalTok{,}\DecValTok{34}\NormalTok{))}

\FunctionTok{table}\NormalTok{(vt}\SpecialCharTok{$}\NormalTok{grp.trt\_1\_234)}
\end{Highlighting}
\end{Shaded}

\begin{verbatim}
## 
##   1 234 
##  16  48
\end{verbatim}

\begin{Shaded}
\begin{Highlighting}[]
\FunctionTok{table}\NormalTok{(vt}\SpecialCharTok{$}\NormalTok{grp.trt\_1\_2\_34)}
\end{Highlighting}
\end{Shaded}

\begin{verbatim}
## 
##  1  2 34 
## 16 16 32
\end{verbatim}

\begin{Shaded}
\begin{Highlighting}[]
\DocumentationTok{\#\#\#\# Testing H0\_1:\textbackslash{}mu\_1=\textbackslash{}mu\_2=\textbackslash{}mu\_3=\textbackslash{}mu\_4}
\NormalTok{vt.mod0 }\OtherTok{\textless{}{-}} \FunctionTok{lm}\NormalTok{(procKnow }\SpecialCharTok{\textasciitilde{}} \DecValTok{1}\NormalTok{, }\AttributeTok{data=}\NormalTok{vt)                 }\DocumentationTok{\#\# 1 mean  (H0)}
\NormalTok{vt.mod1 }\OtherTok{\textless{}{-}} \FunctionTok{lm}\NormalTok{(procKnow }\SpecialCharTok{\textasciitilde{}} \FunctionTok{factor}\NormalTok{(grp.trt), }\AttributeTok{data=}\NormalTok{vt)   }\DocumentationTok{\#\# 4 means}
\FunctionTok{anova}\NormalTok{(vt.mod0)}
\end{Highlighting}
\end{Shaded}

\begin{verbatim}
## Analysis of Variance Table
## 
## Response: procKnow
##           Df Sum Sq Mean Sq F value Pr(>F)
## Residuals 63 331.48  5.2616
\end{verbatim}

\begin{Shaded}
\begin{Highlighting}[]
\FunctionTok{anova}\NormalTok{(vt.mod1)}
\end{Highlighting}
\end{Shaded}

\begin{verbatim}
## Analysis of Variance Table
## 
## Response: procKnow
##                 Df  Sum Sq Mean Sq F value   Pr(>F)   
## factor(grp.trt)  3  65.664 21.8880  4.9406 0.003931 **
## Residuals       60 265.815  4.4302                    
## ---
## Signif. codes:  0 '***' 0.001 '**' 0.01 '*' 0.05 '.' 0.1 ' ' 1
\end{verbatim}

\begin{Shaded}
\begin{Highlighting}[]
\FunctionTok{anova}\NormalTok{(vt.mod0, vt.mod1)}
\end{Highlighting}
\end{Shaded}

\begin{verbatim}
## Analysis of Variance Table
## 
## Model 1: procKnow ~ 1
## Model 2: procKnow ~ factor(grp.trt)
##   Res.Df    RSS Df Sum of Sq      F   Pr(>F)   
## 1     63 331.48                                
## 2     60 265.81  3    65.664 4.9406 0.003931 **
## ---
## Signif. codes:  0 '***' 0.001 '**' 0.01 '*' 0.05 '.' 0.1 ' ' 1
\end{verbatim}

\begin{Shaded}
\begin{Highlighting}[]
\DocumentationTok{\#\#\#\# Testing H02: \textbackslash{}mu\_2=\textbackslash{}mu\_3=\textbackslash{}mu\_4}
\NormalTok{vt.mod2 }\OtherTok{\textless{}{-}} \FunctionTok{lm}\NormalTok{(procKnow }\SpecialCharTok{\textasciitilde{}} \FunctionTok{factor}\NormalTok{(grp.trt\_1\_234), }\AttributeTok{data=}\NormalTok{vt)   }\DocumentationTok{\#\# 2 means}
\FunctionTok{anova}\NormalTok{(vt.mod2)}
\end{Highlighting}
\end{Shaded}

\begin{verbatim}
## Analysis of Variance Table
## 
## Response: procKnow
##                       Df  Sum Sq Mean Sq F value   Pr(>F)    
## factor(grp.trt_1_234)  1  56.816  56.816  12.825 0.000672 ***
## Residuals             62 274.663   4.430                     
## ---
## Signif. codes:  0 '***' 0.001 '**' 0.01 '*' 0.05 '.' 0.1 ' ' 1
\end{verbatim}

\begin{Shaded}
\begin{Highlighting}[]
\FunctionTok{anova}\NormalTok{(vt.mod2, vt.mod1)}
\end{Highlighting}
\end{Shaded}

\begin{verbatim}
## Analysis of Variance Table
## 
## Model 1: procKnow ~ factor(grp.trt_1_234)
## Model 2: procKnow ~ factor(grp.trt)
##   Res.Df    RSS Df Sum of Sq      F Pr(>F)
## 1     62 274.66                           
## 2     60 265.81  2    8.8479 0.9986 0.3744
\end{verbatim}

\begin{Shaded}
\begin{Highlighting}[]
\DocumentationTok{\#\#\#\# Testing H03: \textbackslash{}mu\_3=\textbackslash{}mu\_4}
\NormalTok{vt.mod3 }\OtherTok{\textless{}{-}} \FunctionTok{lm}\NormalTok{(procKnow }\SpecialCharTok{\textasciitilde{}} \FunctionTok{factor}\NormalTok{(grp.trt\_1\_2\_34), }\AttributeTok{data=}\NormalTok{vt)   }\DocumentationTok{\#\# 3 means}
\FunctionTok{anova}\NormalTok{(vt.mod3)}
\end{Highlighting}
\end{Shaded}

\begin{verbatim}
## Analysis of Variance Table
## 
## Response: procKnow
##                        Df Sum Sq Mean Sq F value   Pr(>F)   
## factor(grp.trt_1_2_34)  2  65.51  32.755  7.5123 0.001212 **
## Residuals              61 265.97   4.360                    
## ---
## Signif. codes:  0 '***' 0.001 '**' 0.01 '*' 0.05 '.' 0.1 ' ' 1
\end{verbatim}

\begin{Shaded}
\begin{Highlighting}[]
\FunctionTok{anova}\NormalTok{(vt.mod3, vt.mod1)}
\end{Highlighting}
\end{Shaded}

\begin{verbatim}
## Analysis of Variance Table
## 
## Model 1: procKnow ~ factor(grp.trt_1_2_34)
## Model 2: procKnow ~ factor(grp.trt)
##   Res.Df    RSS Df Sum of Sq      F Pr(>F)
## 1     61 265.97                           
## 2     60 265.81  1   0.15432 0.0348 0.8526
\end{verbatim}

These are the interpretations of the results of the three tests from the
models fit above.

\begin{itemize}
\tightlist
\item
  \(H_0:\mu_1=\mu_2=\mu_3=\mu_4\) - Conclude the 4 means are NOT all
  equal
\item
  \(H_0:\mu_1,\quad \mu_2=\mu_3=\mu_4\) - Fail to reject that the 3
  virtual means are equal
\item
  \(H_0:\mu_1, \mu_2,\quad  \mu_3=\mu_4\) - Fail to reject that the 2
  HMD means are equal
\end{itemize}

\[\nabla\]

Later we will see formally how to make contrasts among treatments to
answer specific research questions.

\subsection{Power Calculations}\label{power-calculations}

When the null hypothesis \(H_0:\mu_1=\cdots =\mu_r\) is false, the
\(F\)-statistic is distributed non-central \(F\) with degrees of freedom
\(\nu_1=r-1\) and \(\nu_2=n_T-r\), and non-centrality parameter
\(\lambda\) defined below. We will use the notation
\(F^* \sim F_{\nu_1,\nu_2}\left(\lambda\right)\).

\[
\lambda = \frac{\sum_{i=1}^rn_i\left(\mu_i-\mu_{\bullet}\right)^2}{\sigma^2}
\qquad \qquad \mu_{\bullet}=\frac{\sum_{i=1}^rn_i\mu_i}{n_T}
\] For a given configuration \(\left\{n_i\right\}\),
\(\left\{\mu_i\right\}\), and \(\sigma^2\), the \textbf{power} of the
test, the probability that it will reject \(H_0\) can be computed as
follows.

\begin{itemize}
\tightlist
\item
  Compute the non-centrality parameter \(\lambda\)
\item
  Obtain the critical value for the \(F\)-statistic:
  \(F_{1-\alpha;r-1,n_T-r}\)
\item
  Obtain the area in the non-central \(F\)-distribution that exceeds the
  critical value in the previous step:
  \(P\left(F_{\nu_1,\nu_2}\left(\lambda\right) \ge F_{1-\alpha;r-1,n_T-r} \right)\)
\end{itemize}

In many practical applications, researchers are interested in detecting
a certain effect size among treatments and must iteratively solve for
sample sizes that assure a certain power, such as 0.80.

\textbf{Example 2.6 - Obtaining the Power of the \(F\)-test}

Consider the example described previously, and suppose we have samples
of size \(n_i=5\) from each treatment.

\[
n_1=n_2=n_3=5 \qquad \mu_1=40 \quad \mu_2=50 \quad \mu_3=60
\quad \mu_{\bullet}=50\qquad \sigma^2=10^2=100 
\]

\[ 
\lambda = \frac{5(40-50)^2+5(50-50)^2+5(60-50)^2}{100}=\frac{1000}{100}=10
\]

If we were conducting a test of equal means, we would reject the null
hypothesis if the \(F\)-statistic \(MSTR/MSE\) exceeds the \(1-\alpha\)
quantile of the \(F_{r-1,n_T-r}\) distribution. The following plot gives
the central and non-central F-distributions for this scenario with with
critical value for \(\alpha=0.05\) significance level
\(F_{.95;2,12}=3.885\). Note that the area above 3.885 in the
non-central F-distribution represents the \textbf{power}, the
probability that we reject \(H_0\) under this scenario for \(r\),
\(\left\{\mu_i\right\}\), \(\sigma\), and \(\left\{n_i\right\}\). Using
the \textbf{pf} function in R, we find the power to be .7015. In many
experimental settings, the goal is to reach a power of .80 to detect a
particular effect. In that case, given the settings for
\(\left\{\mu_i\right\}\), and \(\sigma\), we would need to increase the
sample sizes within treatments, \(\left\{n_i\right\}\).

\begin{Shaded}
\begin{Highlighting}[]
\NormalTok{n }\OtherTok{\textless{}{-}} \FunctionTok{rep}\NormalTok{(}\DecValTok{5}\NormalTok{,}\DecValTok{3}\NormalTok{)}
\NormalTok{mu }\OtherTok{\textless{}{-}} \FunctionTok{c}\NormalTok{(}\DecValTok{40}\NormalTok{, }\DecValTok{50}\NormalTok{, }\DecValTok{60}\NormalTok{)}
\NormalTok{sigma }\OtherTok{\textless{}{-}} \DecValTok{10}
\NormalTok{r }\OtherTok{\textless{}{-}} \FunctionTok{length}\NormalTok{(mu)}
\NormalTok{n\_T }\OtherTok{\textless{}{-}} \FunctionTok{sum}\NormalTok{(n)}
\NormalTok{nu\_1 }\OtherTok{\textless{}{-}}\NormalTok{ r}\DecValTok{{-}1}
\NormalTok{nu\_2 }\OtherTok{\textless{}{-}}\NormalTok{ n\_T}\SpecialCharTok{{-}}\NormalTok{r}
\NormalTok{(F}\FloatTok{.95} \OtherTok{\textless{}{-}} \FunctionTok{qf}\NormalTok{(}\DecValTok{1}\FloatTok{{-}.05}\NormalTok{,nu\_1, nu\_2))}
\end{Highlighting}
\end{Shaded}

\begin{verbatim}
## [1] 3.885294
\end{verbatim}

\begin{Shaded}
\begin{Highlighting}[]
\NormalTok{(mu\_dot }\OtherTok{\textless{}{-}} \FunctionTok{sum}\NormalTok{(n}\SpecialCharTok{*}\NormalTok{mu) }\SpecialCharTok{/} \FunctionTok{sum}\NormalTok{(n))}
\end{Highlighting}
\end{Shaded}

\begin{verbatim}
## [1] 50
\end{verbatim}

\begin{Shaded}
\begin{Highlighting}[]
\NormalTok{(lambda }\OtherTok{\textless{}{-}} \FunctionTok{sum}\NormalTok{(n}\SpecialCharTok{*}\NormalTok{(mu}\SpecialCharTok{{-}}\NormalTok{mu\_dot)}\SpecialCharTok{\^{}}\DecValTok{2}\NormalTok{) }\SpecialCharTok{/}\NormalTok{ sigma}\SpecialCharTok{\^{}}\DecValTok{2}\NormalTok{) }
\end{Highlighting}
\end{Shaded}

\begin{verbatim}
## [1] 10
\end{verbatim}

\begin{Shaded}
\begin{Highlighting}[]
\NormalTok{(power\_HA }\OtherTok{\textless{}{-}} \DecValTok{1}\SpecialCharTok{{-}}\FunctionTok{pf}\NormalTok{(F}\FloatTok{.95}\NormalTok{, nu\_1, nu\_2, lambda))}
\end{Highlighting}
\end{Shaded}

\begin{verbatim}
## [1] 0.7015083
\end{verbatim}

\begin{Shaded}
\begin{Highlighting}[]
\NormalTok{f.grid }\OtherTok{\textless{}{-}} \FunctionTok{seq}\NormalTok{(}\DecValTok{0}\NormalTok{,}\DecValTok{8}\NormalTok{,.}\DecValTok{01}\NormalTok{)}

\FunctionTok{plot}\NormalTok{(f.grid, }\FunctionTok{df}\NormalTok{(f.grid, nu\_1, nu\_2), }\AttributeTok{type=}\StringTok{"l"}\NormalTok{, }\AttributeTok{col=}\StringTok{"red3"}\NormalTok{,}
     \AttributeTok{xlab=}\StringTok{"f"}\NormalTok{, }\AttributeTok{ylab=}\StringTok{"density"}\NormalTok{)}
\FunctionTok{lines}\NormalTok{(f.grid, }\FunctionTok{df}\NormalTok{(f.grid, nu\_1, nu\_2, lambda), }\AttributeTok{col=}\StringTok{"blue4"}\NormalTok{)}
\FunctionTok{abline}\NormalTok{(}\AttributeTok{v=}\NormalTok{F}\FloatTok{.95}\NormalTok{)}
\FunctionTok{legend}\NormalTok{(}\StringTok{"topright"}\NormalTok{, }\FunctionTok{c}\NormalTok{(}\StringTok{"Central F"}\NormalTok{, }\StringTok{"Non{-}Central F"}\NormalTok{), }\AttributeTok{lty=}\DecValTok{1}\NormalTok{,}
       \AttributeTok{col=}\FunctionTok{c}\NormalTok{(}\StringTok{"red4"}\NormalTok{, }\StringTok{"blue4"}\NormalTok{))}
\end{Highlighting}
\end{Shaded}

\begin{figure}

{\centering \includegraphics[width=0.8\linewidth]{bookdownproj_STA4211Notes1_files/figure-latex/ch0208-1} 

}

\caption{Central and non-central F-distributions used for testing for treatment effects and power calculations}\label{fig:ch0208}
\end{figure}

\[\nabla\]

\section{Technical Details}\label{technical-details}

\subsection{Least Squares Estimation for the Cell Means Model
(Regression
Approach)}\label{least-squares-estimation-for-the-cell-means-model-regression-approach}

The cell means model can be written as a regression model in matrix form
as follows. Note that there will not be an intercept in this
formulation. We assume the data are stacked by treatment, then replicate
within treatment in the \(n_T\times 1\) vector \(Y\). The vector
\(\beta\) representing the regression coefficients is the \(r\times 1\)
vector containing \(\mu_1, \ldots,\mu_r\) and the \(X\) matrix is the
\(n_T\times r\) matrix containing a 1 if the observation in the current
row is in the treatment (column) \(i\) and 0 otherwise. We will
formalize this for the simple case of \(r=3\) treatments and
\(n_1=n_2=n_3=2\) replicates per treatment, then we will generalize it.
As with the linear regression model, we will write the model in matrix
form as follows.

\[ Y = X\beta + \epsilon \] where:

\[ 
Y = \begin{bmatrix}
Y_{11} \\ Y_{12} \\ Y_{21} \\ Y_{22} \\
Y_{31} \\ Y_{32} \\ \end{bmatrix}
\qquad 
X = \begin{bmatrix}
1 & 0 & 0 \\ 1 & 0 & 0 \\
0 & 1 & 0 \\ 0 & 1 & 0 \\
0 & 0 & 1 \\ 0 & 0 & 1 \\ \end{bmatrix}
\qquad
\beta = \begin{bmatrix}
\mu_1 \\ \mu_2 \\ \mu_3 \\ \end{bmatrix}
\qquad 
\epsilon = \begin{bmatrix}
\epsilon_{11} \\ \epsilon_{12} \\ \epsilon_{21} \\ \epsilon_{22} \\
\epsilon_{31} \\ \epsilon_{32} \\ \end{bmatrix}
\] Recall that in linear regression, we obtain the ordinary least
squares estimator of \(\beta\), which is labelled \(b\) as follows.
Further the variance-covariance matrix of \(b\) is given as well.

\[ b = \left(X'X\right)^{-1}X'Y \qquad V\{b\} = 
\sigma^2\left(X'X\right)^{-1} \]

For this example, we obtain the following results.

\[ 
X'X = 
\begin{bmatrix}
1 & 1 & 0 & 0 & 0 & 0 \\
0 & 0 & 1 & 1 & 0 & 0 \\
0 & 0 & 0 & 0 & 1 & 1 \\ \end{bmatrix}
\begin{bmatrix}
1 & 0 & 0 \\ 1 & 0 & 0 \\
0 & 1 & 0 \\ 0 & 1 & 0 \\
0 & 0 & 1 \\ 0 & 0 & 1 \\ \end{bmatrix}
= 
\begin{bmatrix}
2 & 0 & 0 \\ 0 & 2 & 0 \\ 0 & 0 & 2 \\ \end{bmatrix}
\]

\[ 
\quad \Rightarrow \quad \left(X'X\right)^{-1}=
\begin{bmatrix}
\frac{1}{2} & 0 & 0 \\ 0 & \frac{1}{2} & 0 \\ 0 & 0 & \frac{1}{2} \\ \end{bmatrix}
\]

Continuing on the computations, we obtain the following results.

\[ 
X'Y = 
\begin{bmatrix}
1 & 1 & 0 & 0 & 0 & 0 \\
0 & 0 & 1 & 1 & 0 & 0 \\
0 & 0 & 0 & 0 & 1 & 1 \\ \end{bmatrix}
\begin{bmatrix}
Y_{11} \\ Y_{12} \\ Y_{21} \\ Y_{22} \\
Y_{31} \\ Y_{32} \\ \end{bmatrix}
=
\begin{bmatrix}
Y_{11}+Y_{12} \\ Y_{21}+Y_{22} \\ Y_{31}+Y_{32}\\ \end{bmatrix}
\]

\[
\quad \Rightarrow \quad
b = \left(X'X\right)^{-1}X'Y = 
\begin{bmatrix}
\frac{1}{2} & 0 & 0 \\ 0 & \frac{1}{2} & 0 \\ 0 & 0 & \frac{1}{2} \\ \end{bmatrix}
\begin{bmatrix}
Y_{11}+Y_{12} \\ Y_{21}+Y_{22} \\ Y_{31}+Y_{32}\\ \end{bmatrix}
=
\begin{bmatrix}
\frac{Y_{11}+Y_{12}}{2} \\
\frac{Y_{21}+Y_{22}}{2} \\
\frac{Y_{31}+Y_{32}}{2} \\
\end{bmatrix}
\]

\[
V\{b\} = 
\sigma^2\left(X'X\right)^{-1}
=\sigma^2 \begin{bmatrix}
\frac{1}{2} & 0 & 0 \\ 0 & \frac{1}{2} & 0 \\ 0 & 0 & \frac{1}{2} \\ \end{bmatrix}
\]

So, this leads to the (hopefully) not surprising general results.

\[ 
\overline{Y}_{i\bullet} = \frac{\sum_{j=1}^{n_i}Y_{ij}}{n_i}
\qquad \qquad
b = \begin{bmatrix}
\overline{Y}_{1\bullet} \\ \overline{Y}_{2\bullet} \\ \vdots \\
\overline{Y}_{r-1,\bullet} \\
\overline{Y}_{r\bullet} \\ \end{bmatrix}
\]

\[
V\{b\} = \sigma^2
\begin{bmatrix}
\frac{1}{n_1} & 0 & \cdots & 0 & 0 \\
0 & \frac{1}{n_2} & \cdots & 0 & 0\\ 
\vdots & \vdots & \ddots & \vdots & \vdots \\
0 & 0 & \cdots & \frac{1}{n_{r-1}} & 0 \\
0 & 0 & \cdots & 0 & \frac{1}{n_r}  \\ \end{bmatrix}
\]

\textbf{Example 2.7 - Virtual Training for a Lifeboat Launching Task}

For this example, we will set up the \(X\) matrix and \(Y\) vector as
follow. Let \(1_{16}\) represent a \(16\times 1\) vector of \(1^s\) and
\(0_{16}\) represent a \(16\times 1\) vector of \(0^s\) and \(Y_i\) be
the \(16\times 1\) vector of \(Y_{i1},\ldots,Y_{i,16}\) for
\(i=1,2,3,4\). Further, let \(Y_{i\bullet}=\sum_{j=1}^{n_i}Y_{ij}\) and
\(\overline{Y}_{i\bullet} = Y_{i\bullet}/n_i\).

\[
X = 
\begin{bmatrix}
1_{16} & 0_{16} & 0_{16} & 0_{16}  \\
0_{16} & 1_{16} & 0_{16} & 0_{16}  \\
0_{16} & 0_{16} & 1_{16} & 0_{16}  \\
0_{16} & 0_{16} & 0_{16} & 1_{16}  \\
\end{bmatrix}
\qquad 
X'X = 
\begin{bmatrix}
16 & 0 & 0 & 0 \\
0 & 16 & 0 & 0\\
0 & 0 & 16 & 0 \\
0 & 0 & 0 & 16\\
\end{bmatrix}
\]

\[ 
Y = 
\begin{bmatrix}
Y_1 \\ Y_2 \\ Y_3 \\ Y_4\\
\end{bmatrix}
\qquad
X'Y =
\begin{bmatrix}
Y_{1\bullet} \\
Y_{2\bullet} \\
Y_{3\bullet} \\
Y_{4\bullet} \\
\end{bmatrix}
=
\begin{bmatrix}
78.8888 \\ 123.3334 \\ 107.7777 \\109.9999 \\
\end{bmatrix}
\]

\[
b=(X'X)^{-1}X'Y
=
\begin{bmatrix}
\frac{1}{16} & 0 & 0 & 0 \\
0 & \frac{1}{16} & 0 & 0\\
0 & 0 & \frac{1}{16} & 0 \\
0 & 0 & 0 & \frac{1}{16}\\
\end{bmatrix}
\begin{bmatrix}
78.8888 \\ 123.3334 \\ 107.7777 \\109.9999 \\
\end{bmatrix}
=
\begin{bmatrix}
4.9306 \\ 7.7083 \\ 6.7361 \\ 6.8750 \\
\end{bmatrix}
\] The variance of each element of \(b\) is \(\sigma^2/16\), and the
standard error is \(\sigma/\sqrt{16}=\sigma/4\). Here, we make use of
the \textbf{lm} function in R to obtain point estimates and estimated
standard errors for the cell means model. Note that the intercept is
removed \((-1)\) in the command and we have to define the treatment as a
factor variable (so as not to fit a straight line).

\begin{Shaded}
\begin{Highlighting}[]
\NormalTok{vt }\OtherTok{\textless{}{-}} \FunctionTok{read.csv}\NormalTok{(}\StringTok{"data/virtual\_training.csv"}\NormalTok{)}
\FunctionTok{head}\NormalTok{(vt, }\DecValTok{2}\NormalTok{)}
\end{Highlighting}
\end{Shaded}

\begin{verbatim}
##   grp.trt procKnow
## 1       1   4.5614
## 2       1   6.6593
\end{verbatim}

\begin{Shaded}
\begin{Highlighting}[]
\NormalTok{vt.mod1 }\OtherTok{\textless{}{-}} \FunctionTok{lm}\NormalTok{(procKnow }\SpecialCharTok{\textasciitilde{}} \FunctionTok{factor}\NormalTok{(grp.trt) }\SpecialCharTok{{-}} \DecValTok{1}\NormalTok{, }\AttributeTok{data=}\NormalTok{vt)}
\FunctionTok{summary}\NormalTok{(vt.mod1)}
\end{Highlighting}
\end{Shaded}

\begin{verbatim}
## 
## Call:
## lm(formula = procKnow ~ factor(grp.trt) - 1, data = vt)
## 
## Residuals:
##     Min      1Q  Median      3Q     Max 
## -5.0831 -1.4444  0.5774  1.5495  3.1900 
## 
## Coefficients:
##                  Estimate Std. Error t value Pr(>|t|)    
## factor(grp.trt)1   4.9306     0.5262    9.37 2.37e-13 ***
## factor(grp.trt)2   7.7083     0.5262   14.65  < 2e-16 ***
## factor(grp.trt)3   6.7361     0.5262   12.80  < 2e-16 ***
## factor(grp.trt)4   6.8750     0.5262   13.06  < 2e-16 ***
## ---
## Signif. codes:  0 '***' 0.001 '**' 0.01 '*' 0.05 '.' 0.1 ' ' 1
## 
## Residual standard error: 2.105 on 60 degrees of freedom
## Multiple R-squared:  0.9139, Adjusted R-squared:  0.9082 
## F-statistic: 159.2 on 4 and 60 DF,  p-value: < 2.2e-16
\end{verbatim}

The summary output includes the residual standard error, which is our
estimate of \(\sigma\), the standard deviation of the model errors.
Derivation of the unbiased estimator of \(\sigma^2\) is described below.

\[\nabla\]

\subsection{Scalar Estimation for the Cell Means
Model}\label{scalar-estimation-for-the-cell-means-model}

We re-write the model here in scalar form.

\[ Y_{ij}=\mu_i + \epsilon_{ij} 
\qquad i=1,\ldots,r; \quad j=1,\ldots, n_i \qquad \epsilon_{ij} \sim NID\left(0,\sigma^2\right) \]

Let \(Q=\sum_{i=1}^r\sum_{j=1}^{n_i}\epsilon_{ij}^2\). We choose
\(\hat{\mu}_k\) that minimizes \(Q\) for \(k=1,\ldots ,r\). This is done
by taking the derivative of \(Q\) with respect to \(\mu_k\), setting it
equal to zero and solving for the least squares estimator.

\[
Q=\sum_{i=1}^r\sum_{j=1}^{n_i}\epsilon_{ij}^2 
= \sum_{i=1}^r\sum_{j=1}^{n_i}\left(Y_{ij}-\mu_i\right)^2 
\]

\[
\quad \Rightarrow \quad 
\frac{\partial Q}{\partial \mu_k} = 
\frac{\partial}{\partial \mu_k}\sum_{i=1}^r\sum_{j=1}^{n_i}\left(Y_{ij}-\mu_i\right)^2
=2\sum_{j=1}^{n_k}\left(Y_{kj}-\mu_k\right)(-1)
\]

Note that this only involves the \(n_k\) observations for treatment
\(k\). Now we set the partial derivative to zero for the least squares
(and maximum likelihood) estimator of \(\mu_k\).

\[
-2\sum_{j=1}^{n_k}\left(Y_{kj}-\hat{\mu_k}\right) = 0
\quad \Rightarrow \quad \sum_{j=1}^{n_k}Y_{kj} = n_k\hat{\mu_k}
\quad \Rightarrow \quad \hat{\mu_k}= \frac{\sum_{j=1}^{n_k}Y_{kj}}{n_k}
=\overline{Y}_{k\bullet}
\]

Thus the fitted values for \(Y_{ij}\) are
\(\hat{Y}_{ij}=\overline{Y}_{i\bullet}\).

\subsection{Estimation for the Treatment Effects
Model}\label{estimation-for-the-treatment-effects-model}

Let \(\mu_i\) be the population mean for treatment \(i\), and let
\(\mu_{\bullet}\) be the unweighted mean of the \(\mu_i^s\). Note that
this will differ from \(\mu_{\bullet}\) used in the cell means model
when the sample sizes are not all equal.

\[ 
\mu_{\bullet}=\frac{\sum_{i=1}^r\mu_i}{r}  \qquad
\mu_i = \mu_{\bullet} + \left(\mu_i-\mu_{\bullet}\right) = \mu_{\bullet} + \tau_i
\quad \Rightarrow \quad \tau_i = \mu_i-\mu_{\bullet}
\]

We will refer to \(\tau_i\) as the effect of treatment \(i\) relative to
the mean among all treatments. The sum of \(\tau_i\) is zero, when we
use the unweighted mean \(\mu_{\bullet}\).

\[
\mu_1=\cdots=\mu_r=\mu_{\bullet} \quad \Rightarrow \quad
\tau_1=\cdots=\tau_r = 0
\]

\[
Y_{ij}=\mu_{\bullet} +\tau_i + \epsilon_{ij} 
\qquad i=1,\ldots,r; \quad j=1,\ldots, n_i \qquad \epsilon_{ij} \sim NID\left(0,\sigma^2\right) \]

\[ E\{Y_{ij}\} = E\{\mu_{\bullet} +\tau_i + \epsilon_{ij}\} = \mu_{\bullet} +\tau_i+0=\mu_{\bullet} +\tau_i \qquad
V\{Y_{ij}\} = V\{\mu_{\bullet} +\tau_i + \epsilon_{ij}\} = 0 +\sigma^2=\sigma^2 
\]

\[
\sum_{i=1}^{r}\tau_i=0 \quad \Rightarrow \quad \tau_r = -\sum_{i=1}^{r-1}\tau_i
\]

\subsection{Least Squares Estimation (Regression
Approach)}\label{least-squares-estimation-regression-approach}

Let \(Y\) be defined as before. We will parameterize \(\beta\) in terms
of \(\mu_{\bullet}\) and \(\tau_1,\ldots,\tau_{r-1}\), keeping in mind
the result above regarding \(\tau_r\). First consider the case as we did
previously with \(r=3\) and \(n_1=n_2=n_3=2\).

\[ 
Y = \begin{bmatrix}
Y_{11} \\ Y_{12} \\ Y_{21} \\ Y_{22} \\
Y_{31} \\ Y_{32} \\ \end{bmatrix}
\qquad 
X = \begin{bmatrix}
1 & \phantom{-}1 & \phantom{-}0 \\ 1 & \phantom{-}1 & \phantom{-}0 \\
1 & \phantom{-}0 & \phantom{-}1 \\ 1 & \phantom{-}0 & \phantom{-}1 \\
1 & -1 & -1 \\ 1 & -1 & -1 \\ \end{bmatrix}
\qquad
\beta = \begin{bmatrix}
\mu_{\bullet} \\ \tau_1 \\ \tau_2 \\ \end{bmatrix}\phantom{-}
\qquad 
\epsilon = \begin{bmatrix}
\epsilon_{11} \\ \epsilon_{12} \\ \epsilon_{21} \\ \epsilon_{22} \\
\epsilon_{31} \\ \epsilon_{32} \\ \end{bmatrix}
\]

\[ b = \left(X'X\right)^{-1}X'Y \qquad V\{b\} = 
\sigma^2\left(X'X\right)^{-1} \]

For this example, we obtain the following results.

\[ 
X'X = 
\begin{bmatrix}
1 & 1 & 1 & 1 & \phantom{-}1 & \phantom{-}1 \\
1 & 1 & 0 & 0 & -1 & -1 \\
0 & 0 & 1 & 1 & -1 & -1 \\ \end{bmatrix}
\begin{bmatrix}
1 & \phantom{-}1 & \phantom{-}0 \\ 1 & \phantom{-}1 & \phantom{-}0 \\
1 & \phantom{-}0 & \phantom{-}1 \\ 1 & \phantom{-}0 & \phantom{-}1 \\
1 & -1 & -1 \\ 1 & -1 & -1 \\ \end{bmatrix}
= 
\begin{bmatrix}
6 & 0 & 0 \\ 0 & 4 & 2 \\ 0 & 2 & 4 \\ \end{bmatrix}
\]

Making use of cofactors, we obtain \(\left(X'X\right)^{-1}\) below.

\[ 
\quad \Rightarrow \quad \left(X'X\right)^{-1}=
\begin{bmatrix}
\frac{1}{6} & \phantom{-}0 & \phantom{-}0 \\ 
0 & \phantom{-}\frac{1}{3} & -\frac{1}{6} \\ 
0 & -\frac{1}{6} & \phantom{-}\frac{1}{3} \\ \end{bmatrix} =
\frac{1}{6}
\begin{bmatrix}
1 & \phantom{-}0 & \phantom{-}0 \\ 0 & \phantom{-}2 & -1 \\ 0 & -1 & \phantom{-}2 \\
\end{bmatrix}
\]

Continuing on the computations, we obtain the following results.

\[ 
X'Y = 
\begin{bmatrix}
1 & 1 & 1 & 1 & \phantom{-}1 & \phantom{-}1 \\
1 & 1 & 0 & 0 & -1 & -1 \\
0 & 0 & 1 & 1 & -1 & -1 \\ \end{bmatrix}
\begin{bmatrix}
Y_{11} \\ Y_{12} \\ Y_{21} \\ Y_{22} \\
Y_{31} \\ Y_{32} \\ \end{bmatrix}
=
\begin{bmatrix}
Y_{11}+Y_{12}+Y_{21}+Y_{22}+Y_{31}+Y_{32} \\ 
Y_{11}+Y_{12}-Y_{31}-Y_{32} \\ 
Y_{21}+Y_{22}-Y_{31}-Y_{32}\\ \end{bmatrix}
\]

When the dust settles, the estimator \(b\) for \(\beta\) is obtained
below.

\[
b=\left(X'X\right)^{-1}X'Y =
\frac{1}{6}
\begin{bmatrix}
1 & \phantom{-}0 & \phantom{-}0 \\ 0 & \phantom{-}2 & -1 \\ 0 & -1 & \phantom{-}2 \\
\end{bmatrix}
\begin{bmatrix}
Y_{11}+Y_{12}+Y_{21}+Y_{22}+Y_{31}+Y_{32} \\ 
Y_{11}+Y_{12}-Y_{31}-Y_{32} \\ 
Y_{21}+Y_{22}-Y_{31}-Y_{32}\\ \end{bmatrix}
= 
\begin{bmatrix}
\overline{Y}_{\bullet\bullet} \\
\overline{Y}_{1\bullet} - \overline{Y}_{\bullet\bullet} \\
\overline{Y}_{2\bullet} - \overline{Y}_{\bullet\bullet} \\
\end{bmatrix}
\] The general result is as follows, allowing for unequal sample sizes.

\[
b =
\begin{bmatrix}
\frac{1}{r}\sum_{i=1}^r \overline{Y}_{i\bullet} \\
\frac{r-1}{r}\overline{Y}_{1\bullet}-\frac{1}{r}\sum_{i=2}^r\overline{Y}_{i\bullet} \\
\vdots \\
\frac{r-1}{r}\overline{Y}_{r-1,\bullet}-\frac{1}{r}\sum_{i=1 \atop i \ne r-1}^r\overline{Y}_{i\bullet}
\end{bmatrix}
\qquad
\hat{\tau}_r=-\sum_{i=1}^{r-1}\hat{\tau}_i=
\frac{r-1}{r}\overline{Y}_{r\bullet}-\frac{1}{r}\sum_{i=1}^{r-1}\overline{Y}_{i\bullet}
\] From these expressions, we obtain the following variances for
\(\hat{\mu}\) and \(\hat{\tau}_k\).

\[
V\left\{\hat{\mu}\right\} = V\left\{\frac{1}{r}\sum_{i=1}^r \overline{Y}_{i\bullet}\right\}=
\frac{1}{r^2}\sum_{i=1}^rV\left\{\overline{Y}_{i\bullet}\right\}=
\frac{\sigma^2}{r^2}\sum_{i=1}^r\frac{1}{n_i}
\] \[V\left\{\hat{\tau}_k\right\} =
V\left\{\frac{r-1}{r}\overline{Y}_{k\bullet}-\frac{1}{r}\sum_{i=1 \atop i \ne k}^r\overline{Y}_{i\bullet}\right\}=
\frac{\sigma^2}{r^2}\left[
\frac{(r-1)^2}{n_k} + \sum_{i=1 \atop i \ne k}^r\frac{1}{n_i}
\right]
\]

\subsection{The Analysis of Variance}\label{the-analysis-of-variance}

In this section, we obtain a decomposition of the total variation of the
individual measurements around the overall mean, the \textbf{Total Sum
of Squares}, which we will denote \(SSTO\). Further, we will derive the
expectations of the error and treatment sums of squares (and mean
squares).

\[
\overline{Y}_{\bullet\bullet}=
\frac{\sum_{i=1}^r\sum_{j=1}^{n_i}Y_{ij}}{n_T}
\qquad \qquad
SSTO = \sum_{i=1}^r\sum_{j=1}^{n_i}
  \left(Y_{ij}-\overline{Y}_{\bullet\bullet}\right)^2
\] A useful expansion of this is helpful for derivations below.

\[
SSTO = \sum_{i=1}^r\sum_{j=1}^{n_i}
  \left(Y_{ij}^2 + \overline{Y}_{\bullet\bullet}^2
  -2Y_{ij}\overline{Y}_{\bullet\bullet}\right)
\] \[
\quad \Rightarrow \quad
SSTO = \sum_{i=1}^r\sum_{j=1}^{n_i} Y_{ij}^2 +
  n_T\overline{Y}_{\bullet\bullet}^2 -
  2\overline{Y}_{\bullet\bullet}\sum_{i=1}^r\sum_{j=1}^{n_i} Y_{ij}=
  \sum_{i=1}^r\sum_{j=1}^{n_i} Y_{ij}^2 +
  n_T\overline{Y}_{\bullet\bullet}^2 -
  2\overline{Y}_{\bullet\bullet}n_T\overline{Y}_{\bullet\bullet}
  \] \[
  \quad \Rightarrow \quad
SSTO 
  = \sum_{i=1}^r\sum_{j=1}^{n_i} Y_{ij}^2 -
  n_T\overline{Y}_{\bullet\bullet}^2
  \]

As was the case for the linear regression model, we obtain fitted values
for the observations in the form \(\hat{Y}=Xb\). This leads to the error
sum of squares, labelled as
\(SSE=\left(Y-\hat{Y}\right)'\left(Y-\hat{Y}\right)\) here (note that in
R, for linear models, it is labelled as \(RSS\), for the residual sum of
squares).

\[
SSE = \sum_{i=1}^r\sum_{j=1}^{n_i}  \left(Y_{ij}-\hat{Y}_{ij}\right)^2
\] Noting that for this model, \(\hat{Y}_{ij}=\overline{Y}_{i\bullet}\),
we obtain the following formula for \(SSE\).

\[
SSE=\sum_{i=1}^r\sum_{j=1}^{n_i} \left(Y_{ij}-\overline{Y}_{i\bullet}\right)^2
\] We can expand this form of \(SSE\) which will be used to derive the
unbiased estimator for \(\sigma^2\), the error variance.

\[
SSE=\sum_{i=1}^r\sum_{j=1}^{n_i}\left(
Y_{ij}^2 + \overline{Y}_{i\bullet}^2 - 
2Y_{ij}\overline{Y}_{i\bullet}\right)
\] \[
\quad \Rightarrow \quad
SSE = \sum_{i=1}^r\sum_{j=1}^{n_i} Y_{ij}^2
+ \sum_{i=1}^r n_i\overline{Y}_{i\bullet}^2
-2\sum_{i=1}^r\overline{Y}_{i\bullet}\sum_{j=1}^{n_i} Y_{ij}=
\sum_{i=1}^r\sum_{j=1}^{n_i} Y_{ij}^2
+ \sum_{i=1}^r n_i\overline{Y}_{i\bullet}^2
-2\sum_{i=1}^r  n_i\overline{Y}_{i\bullet}^2
\] \[
\quad \Rightarrow \quad 
SSE =\sum_{i=1}^r\sum_{j=1}^{n_i} Y_{ij}^2 -
\sum_{i=1}^r n_i\overline{Y}_{i\bullet}^2
\]

To obtain an unbiased estimator for \(\sigma^2\), consider the following
results from the Cell Means Model.

\[
E\left\{Y_{ij}\right\}=\mu_i 
\qquad
V\left\{Y_{ij}\right\}= \sigma^2 
\]

\[
\quad \Rightarrow \quad
E\left\{Y_{ij}^2\right\} = 
V\left\{Y_{ij}\right\} + 
\left[E\left\{Y_{ij}\right\}\right]^2
=
\sigma^2 + \mu_i^2
\]

Similarly, we obtain the expected values of the squares of the treatment
means below.

\[
E\left\{\overline{Y}_{i\bullet}\right\}
=\frac{1}{n_i}\sum_{j=1}^{n_i}E\left\{Y_{ij}\right\}
=\mu_i
\qquad 
V\left\{\overline{Y}_{i\bullet}\right\}
=\frac{1}{n_i^2}\sum_{j=1}^{n_i}V\left\{Y_{ij}\right\}
=\frac{\sigma^2}{n_i}
\] \[
\quad \Rightarrow \quad
E\left\{\overline{Y}_{i\bullet}^2\right\}=
\frac{\sigma^2}{n_i} + \mu_i^2
\] Next, we obtain \(E\left\{SSE\right\}\) making use of these results.

\[
E\left\{SSE\right\} =
\sum_{i=1}^r\sum_{j=1}^{n_i} E\left\{Y_{ij}^2\right\} -
\sum_{i=1}^r n_iE\left\{\overline{Y}_{i\bullet}^2\right\}
\]

\[
\quad \Rightarrow \quad
E\left\{SSE\right\} =
\sum_{i=1}^r n_i\left(\sigma^2 + \mu_i^2\right)
- \sum_{i=1}^r n_i\left(\frac{\sigma^2}{n_i} + \mu_i^2\right)
\]

\[
\quad \Rightarrow \quad
E\left\{SSE\right\} 
=n_T\sigma^2 - r\sigma^2 = \left(n_T-r\right)\sigma^2
\]

Thus, if we divide \(SSE\) by \(n_T-r\), we obtain an unbiased estimator
of \(\sigma^2\). This is the \textbf{Mean Square Error}, denoted as
\(MSE\), and often labelled as \(S^2\) (as an estimator) and \(s^2\)
when computed on an observed set of data. The denominator \(n_T-r\) is
the \textbf{Error Degrees of Freedom}.

\[
MSE = \frac{SSE}{n_T-r}  \qquad E\left\{MSE\right\} = \sigma^2
\]

The \textbf{Treatment Sum of Squares}, denoted as \(SSTR\), acts like
the regression sum of squares for linear regression models. That is,
\(SSTR = \left(\hat{Y}-\overline{Y}1\right)'\left(\hat{Y}-\overline{Y}1\right)\).
Keeping in mind that the fitted value for \(Y_{ij}\) is
\(\hat{Y_{ij}} = \overline{Y}_{i\bullet}\), we obtain SSTR as follows.

\[
SSTR = \sum_{i=1}^r\sum_{j=1}^{n_i}
\left(\overline{Y}_{i\bullet}-
  \overline{Y}_{\bullet\bullet}\right)^2
 = \sum_{i=1}^r n_i \left(\overline{Y}_{i\bullet}-
  \overline{Y}_{\bullet\bullet}\right)^2
\]

Again, it is helpful to expand this as before.

\[
SSTR = \sum_{i=1}^r n_i \overline{Y}_{i\bullet}^2
- n_T\overline{Y}_{\bullet\bullet}^2 
\]

We know the expectation of the first term from the derivation of
\(E\left\{SSE\right\}\), so we just need the second term.

\[
E\left\{\overline{Y}_{\bullet\bullet}\right\}
=\frac{\sum_{i=1}^r\sum_{j=1}^{n_i}E\left\{Y_{ij}\right\}}{n_T}
=\frac{\sum_{i=1}^rn_i\mu_i}{n_T} = \mu_{\bullet}
\]

\[
V\left\{\overline{Y}_{\bullet\bullet}\right\}
=\frac{1}{n_T^2}\sum_{i=1}^r\sum_{j=1}^{n_i}
V\left\{Y_{ij}\right\} = \frac{\sigma^2}{n_T}
\]

\[
\quad \Rightarrow \quad
E\left\{SSTR\right\} =
\sum_{i=1}^r n_i\left(\frac{\sigma^2}{n_i} + \mu_i^2\right) -
\left(\sigma^2 + n_T \mu_{\bullet}^2\right)
\]

\[
\quad \Rightarrow \quad
E\left\{SSTR\right\} =
\sigma^2(r-1) + \sum_{i=1}^rn_i\mu_i^2 - n_T\mu_{\bullet}^2
\]

\[
\quad \Rightarrow \quad
E\left\{SSTR\right\}
= \sigma^2(r-1) + \sum_{i=1}^rn_i\left(\mu_i-\mu_{\bullet}\right)^2
\] The last step involving the \(\mu^s\) is easiest to see as follows.

\[
\sum_{i=1}^rn_i\left(\mu_i-\mu_{\bullet}\right)^2=
\sum_{i=1}^rn_i\left(\mu_i^2+\mu_{\bullet}^2-2\mu_i\mu_{\bullet}\right)=
\sum_{i=1}^rn_i \mu_i^2+n_T\mu_{\bullet}^2-2\mu_{\bullet}\sum_{i=1}^rn_i \mu_i
\] \[
\mbox{ where } \sum_{i=1}^rn_i \mu_i = n_T\mu_{\bullet} \quad \Rightarrow \quad
\sum_{i=1}^rn_i\left(\mu_i-\mu_{\bullet}\right)^2=
\sum_{i=1}^rn_i \mu_i^2-n_T\mu_{\bullet}^2
\]

The treatment degrees of freedom are \(r-1\), the number of treatments
minus one. We compute the \textbf{Treatment Mean Square} by dividing the
Treatment Sum of Squares by its degrees of freedom and obtain its
expectation below.

\[ 
MSTR = \frac{SSTR}{r-1} \qquad 
E\left\{MSTR\right\} = \sigma^2 + \frac{\sum_{i=1}^rn_i\left(\mu_i-\mu_{\bullet}\right)^2}{r-1}
\]

When all treatment means are equal, that is,
\(\mu_1=\cdots=\mu_r=\mu_{\bullet}\), then the expected mean squares for
treatment and error are equal
\(\left(E\left\{MSTR\right\} = E\left\{MSE\right\}=\sigma^2\right)\),
otherwise it is higher for treatments than for error. Further, under the
assumptions of the model thus far (normal, independent errors with
constant variance) we obtain the following results.

\[
\frac{SSE}{\sigma^2}=\frac{\left(n_T-r\right)MSE}{\sigma^2} \sim \chi_{n_T-r}^2
\]

\[
\frac{SSTR}{\sigma^2}=\frac{(r-1)MSTR}{\sigma^2} \sim \chi_{r-1}^2(\lambda)
\]

where \(\lambda\) is the non-centrality parameter of the chi-square
distribution. In the One-Way Analysis of Variance, \(\lambda\) is
computed as follow.

\[
\lambda = \frac{\sum_{i=1}^rn_i\left(\mu_i-\mu_{\bullet}\right)^2}{\sigma^2}
\]

Further \(MSTR\) and \(MSE\) are independent. When all \(\mu_i\) are
equal, then the chi-square distribution for \(SSTR/\sigma^2\) is
``central'' chi-square, and we can use the (central) \(F\)-distribution
to test the hypothesis \(H_0:\mu_1=\cdots = \mu_r\).

We will make use of the following results when testing for differences
among means. Suppose \(W_1\) is chi-square with \(\nu_1\) degrees of
freedom, \(W_2\) is chi-square with \(\nu_2\) degrees of freedom and
\(W_1\) and \(W_2\) are independent, then we get the following
distributional result.

\[
W_1 \sim \chi_{\nu_1}^2 \quad
W_2 \sim \chi_{\nu_2}^2 \quad W_1 \bot W_2 \quad \Rightarrow \quad
F = \frac{W_1/\nu_1}{W_2/\nu_2} \sim F_{\nu_1,\nu_2}
\]

When \(W_1\) is chi-square with \(\nu_1\) degrees of freedom and
non-centrality parameter \(\lambda\) and \(W_2\) is chi-square with
\(\nu_2\) degrees of freedom and \(W_1\) and \(W_2\) are independent,
then we get the following distributional result.

\[
W_1 \sim \chi_{\nu_1}^2(\lambda) \quad
W_2 \sim \chi_{\nu_2}^2 \quad W_1 \bot W_2 \quad \Rightarrow \quad
F = \frac{W_1/\nu_1}{W_2/\nu_2} \sim F_{\nu_1,\nu_2}(\lambda)
\]

This is the non-central \(F\)-distribution. We can use R to obtain
probabilities, densities, quantiles, and random samples for non-central
\(F\)-distributions.

\begin{verbatim}
## Warning: package 'agricolae' was built under R version 4.5.3
\end{verbatim}

\begin{verbatim}
## Warning: package 'multcompView' was built under R version 4.5.3
\end{verbatim}

\chapter{Analysis of Treatment Means}\label{analysis-of-treatment-means}

In this chapter, we will consider estimation and inference concerning
treatment/group means. We will first consider estimating individual
treatment means. Then, we will consider \textbf{Contrasts} among
treatment means. Finally, we consider all pairwise comparisons.

\section{Individual Treatment Means}\label{individual-treatment-means}

For the cell means model, we may be interested in making inference
regarding the individual treatment.group means. The model is of the
following form. Note that this is a shortened version of material
covered in Chapter 1.

\[ 
Y_{ij}=\mu_i + \epsilon_{ij} 
\qquad i=1,\ldots,r; \quad j=1,\ldots, n_i \qquad \epsilon_{ij} \sim NID\left(0,\sigma^2\right) 
\]

The least squares estimator for \(\mu_i\) is
\(\hat{\mu}_i=\overline{Y}_{i\bullet}\).

\[
\hat{\mu}_i = \overline{Y}_{i\bullet} = \frac{1}{n_i}\sum_{j=1}^{n_i} Y_{ij} \qquad
E\left\{\overline{Y}_{i\bullet}\right\}=\frac{1}{n_i}n_i\mu_i = \mu_i \qquad
V\left\{\overline{Y}_{i\bullet}\right\}=\frac{1}{n_i^2}n_i\sigma^2=
\frac{\sigma^2}{n_i}
\] For this model, the treatment means are independent and we have the
following distributional properties.

\[
\overline{Y}_{i\bullet} \sim N\left(\mu_i,\frac{\sigma^2}{n_i}\right)
\qquad \qquad
\frac{SSE}{\sigma^2}=\frac{\left(n_T-r\right)MSE}{\sigma^2} \sim \chi_{n_T-r}^2
\]

Further, \(\overline{Y}_{i\bullet}\) and \(MSE\) are independent.

Consider two independent random variables, \(Z \sim N(0,1)\) and
\(W\sim \chi_{\nu}^2\). The ratio \(T=Z/\sqrt{W/\nu}\) is distributed as
\textbf{Student's t} with \(\nu\) degrees of freedom. In this
application, we have the following results. \[
Z=\frac{\overline{Y}_{i\bullet}-\mu_i}{\sqrt{\sigma^2/n_i}} \sim N(0,1) \qquad \qquad 
W=\frac{\left(n_T-r\right)MSE}{\sigma^2} \sim \chi_{n_T-r}^2 
\] \[
T=\frac{Z}{\sqrt{W/\nu}}=
\frac{\frac{\overline{Y}_{i\bullet}-\mu_i}{\sqrt{\sigma^2/n_i}}}{
\sqrt{\frac{\left(n_T-r\right)MSE}{\sigma^2}/\left(n_T-r\right)}
}=\frac{\overline{Y}_{i\bullet}-\mu_i}{\sqrt{MSE/n_i}} \sim t_{n_T-r}
\] Since \(\sigma^2\) is unknown in practice and must be estimated with
\(s^2=MSE\), we obtain this result, where \(t_{n-r}\) represents
Student's \(t\)-distribution with \(n_T-r\) degrees of freedom.

\[
\hat{V}\left\{\overline{Y}_{i\bullet}\right\}=\frac{MSE}{n_i} \qquad  
\hat{SE}\left\{\overline{Y}_{i\bullet}\right\}=\sqrt{\frac{MSE}{n_i}}
\qquad
\frac{\overline{Y}_{i\bullet}-\mu_i}{\hat{SE}\left\{\overline{Y}_{i\bullet}\right\}}
\sim t_{n_T-r}
\]

This leads to a \((1-\alpha)100\%\) Confidence Interval for \(\mu_i\) of
this form.

\[\overline{Y}_{i\bullet} \pm t_{1-\alpha/2;n_T-r}\hat{SE}\left\{\overline{Y}_{i\bullet}\right\}
\quad \equiv \quad
\overline{Y}_{i\bullet} \pm t_{1-\alpha/2;n_T-r}\sqrt{\frac{MSE}{n_i}}
\]

In the rare situation that we wish to test \(H_0:\mu_i=c\) for some
pre-specified constant \(c\), we can compute the \(t\)-statistic as
follows.

\[
\mbox{Test Statistic:} \quad t_i^* = \frac{\overline{Y}_{i\bullet}-c}{\hat{SE}\left\{\overline{Y}_{i\bullet}\right\}} \qquad
\mbox{Rejection Region:} \quad \left|t_i^*\right| \ge t_{1-\alpha/2;n_T-r}
\]

\section{Comparing Two Treatment
Means}\label{comparing-two-treatment-means}

In some cases, researchers may wish to compare two specific treatment
means. In this case, suppose the goal is to compare treatments \(i\) and
\(i'\). The parameter of interest is \(D=\mu_i-\mu_{i'}\) and the
estimator is
\(\hat{D}=\overline{Y}_{i\bullet}-\overline{Y}_{i'\bullet}\). The mean
and variance of \(\hat{D}\) are given below.

\[
E\left\{\hat{D}\right\}=E\left\{\overline{Y}_{i\bullet}-\overline{Y}_{i'\bullet}\right\}=\mu_i-\mu_{i'}=D 
\] \[
V\left\{\hat{D}\right\}=
V\left\{\overline{Y}_{i\bullet}\right\} +
V\left\{\overline{Y}_{i'\bullet}\right\}
=\sigma^2\left(\frac{1}{n_i}+\frac{1}{n_{i'}}\right)
\] As \(\sigma^2\) is not known in practice, the estimated variance of
\(\hat{D}\) replaces \(\sigma^2\) with \(MSE\).

\[
\hat{V}\left\{\hat{D}\right\}=
\hat{V}\left\{\overline{Y}_{i\bullet}\right\} +
\hat{V}\left\{\overline{Y}_{i'\bullet}\right\}
=MSE\left(\frac{1}{n_i}+\frac{1}{n_{i'}}\right)
\]

A \((1-\alpha)100\%\) Confidence Interval for \(D=\mu_i-\mu_{i'}\) is
obtained as follows, as well as a test of \(H_0:D=0\) versus
\(H_A:D\ne 0\).

\[
\mbox{Confidence Interval:} \quad
\hat{D} \pm t_{1-\alpha/2;n_T-r}\hat{SE}\left\{\hat{D}\right\}  \equiv
\hat{D} \pm t_{1-\alpha/2;n_T-r}
\sqrt{MSE\left(\frac{1}{n_i}+\frac{1}{n_{i'}}\right)}
\]

\[
\mbox{Test Statistic:} \quad t^*=\frac{\hat{D}}{\hat{SE}\left\{\hat{D}\right\}}
\qquad \mbox{Rejection Region:} \quad
\left|t^*\right| \ge t_{1-\alpha/2;n_T-r}
\]

The only difference between this and the 2-sample \(t\)-test is that all
\(r\) treatments/groups are used to estimate \(\sigma^2\), and
increasing the degrees of freedom to \(n_T-r\). We will consider
adjustments for comparing all pairs of treatments later in the chapter.

\section{Contrasts Among Treatment
Means}\label{contrasts-among-treatment-means}

A \textbf{Contrast} is any linear function of treatment means such that
the coefficients sum to 0. Thus, for instance \(D=\mu_i-\mu_{i'}\) is a
contrast among treatment means that only involves the difference between
treatments \(i\) and \(i'\). In general, we define contrasts as follows.

\[
L = \sum_{i=1}^rc_i\mu_i \qquad \mbox{such that: } \sum_{i=1}^rc_i=0
\]

The estimator for \(L\) is \(\hat{L}\) that replaces the
\(\left\{\mu_i\right\}\) with the treatment means
\(\left\{\overline{Y}_{i\bullet}\right\}\).

\[
\hat{L}=\sum_{i=1}^rc_i\overline{Y}_{i\bullet} \qquad
E\left\{\hat{L}\right\}=\sum_{i=1}^rc_i\mu_i=L
\]

\[
V\left\{\hat{L}\right\}=\sum_{i=1}^r c_i^2\frac{\sigma^2}{n_i}
\qquad \qquad
\hat{V}\left\{\hat{L}\right\}=MSE\sum_{i=1}^r\frac{c_i^2}{n_i}
\] We can make use of the \(t\)-distribution to obtain confidence
intervals and tests regarding contrasts. Also, we can use an \(F\)-test
as well.

\[
\frac{\hat{L}-L}{\hat{SE}\left\{\hat{L}\right\}} \sim t_{n_T-r}
\quad \Rightarrow \quad (1-\alpha)100\% \mbox{ CI for L:} \quad
\hat{L} \pm t_{1-\alpha/2;n_T-r}\hat{SE}\left\{\hat{L}\right\}
\]

For testing \(H_0:L=0\) versus \(H_A:L\ne 0\), we can use either a
\(t\)-test or \(F\)-test, which provide the same conclusions and
\(P\)-values. The \(t\)-test is given here.

\[
\mbox{Test Statistic:} \quad t^*=\frac{\hat{L}}{\hat{SE}\left\{\hat{L}\right\}}
\qquad \mbox{Rejection Region:}\quad \left|t^*\right| \ge t_{1-\alpha/2;n_T-r}
\]

For the \(F\)-test, we first define the sum of squares for the contrast,
then construct the \(F\)-statistic.

\[
\mbox{Contrast Sum of Squares:} \quad
SSL = \frac{\left(\hat{L}\right)^2}{\sum_{i=1}^r\frac{c_i^2}{n_i}}
\] \[n_1=\cdots =n_r=n \quad \Rightarrow \quad
SSL=n\frac{\left(\hat{L}\right)^2}{\sum_{i=1}^rc_i^2}
\]

\[
\mbox{Test Statistic:} \quad F^*=\frac{SSL}{MSE} \qquad
\mbox{Rejection Region:} \quad F^* \ge F_{1-\alpha;1,n_T-r}
\]

Two contrasts \(L_1=\sum_{i=1}^ra_i\mu_i\) and
\(L_2=\sum_{i=1}^rb_i\mu_i\) where \(\sum_{i=1}^ra_i=\sum_{i=1}^rb_i=0\)
are said to be \textbf{Orthogonal Contrasts} if the product of their
coefficients divided by their sample sizes is 0.

\[
L_1,L_2\mbox{ are orthogonal if:} \quad \sum_{i=1}^r\frac{a_ib_i}{n_i}=0
\]

When the experiment is balanced, this simplifies to
\(\sum_{i=1}^ra_ib_i=0\). Among \(r\) treatments/groups, we can
construct \(r-1\) pairwise orthogonal contrasts, whose sums of squares
and degrees of freedom sum to \(SSTR\) and \(df_{TR}\), respectively.

\textbf{Example 3.1 - Virtual Training for a Lifeboat Launching Task}

Here, we re-describe the virtual training study to set up some
interesting contrasts among the training methods. A study in South Korea
compared \(r=4\) methods of training to conduct a lifeboat launching
task (\citeproc{ref-jung_2018}{Jung and Ahn 2018}). The treatments and
their labels from the paper are given below and the response \(Y\) was a
procedural knowledge score for subjects post training.

\begin{itemize}
\tightlist
\item
  \textbf{Lecture/Materials} - Traditional Lecture with no computer
  component (LEC/MAT)
\item
  \textbf{Monitor/Keyboard} - Trained virtually with a monitor,
  keyboard, and mouse (MON/KEY)
\item
  \textbf{Head-Mounted Display/Joypad} - Trained virtually with HMD and
  joypad (HMD/JOY)
\item
  \textbf{Head-Mounted Display/Wearable Sensors} - (HMD/WEA)
\end{itemize}

There were a total of \(n_T=64\) subjects and they were randomized so
that 16 subjects received each treatment
\(\left(n_1=n_2=n_3=n_4=16\right)\).

Consider the following 3 contrasts.

\begin{itemize}
\tightlist
\item
  \textbf{Lecture Materials versus Virtual Training}
  \(a_1=3,a_2=a_3=a_4=-1\)
\item
  \textbf{Monitor/Keyboard versus Head Mounted Display}
  \(b_1=0,b_2=2,b_3=b_4=-1\)
\item
  \textbf{Head Mounted Displays: Joypad versus Wearables}
  \(d_1=d_2=0,d_3=1,d_4=-1\)
\end{itemize}

All of these are contrasts, as their coefficients sum to zero. Further,
they are pairwise orthogonal.

\[
\sum_{i=1}^ra_ib_i=3(0)+(-1)(2)+(-1)(-1)+(-1)(-1)=0-2+1+1=0
\]

\[
\sum_{i=1}^ra_id_i=3(0)+2(0)+(-1)(1)+(-1)(-1)=0+0-1+1=0
\]

\[
\sum_{i=1}^rb_id_i=0(0)+2(0)+(-1)(1)+(-1)(-1)=0+0-1+1=0
\] Table @ref(tab:ch0301) gives calculations necessary to obtain the
Confidence Intervals, \(t\)-tests, sums of squares and \(F\)-tests for
the three contrasts. Recall from the last chapter, the following
quantities.

\[
SSTR=65.628 \quad df_{TR}=3 \qquad SSE=265.815 \quad df_E=60 \quad MSE=4.430
\] The critical \(t\)-value for \(\alpha=.05\) and \(df_E=60\) is
\(t_{.975,60}=2.000\).

For contrast \(L_1\), comparing the virtual training methods to the
``control'' treatment has the following values.

\[
\hat{L}_1=3\overline{Y}_{1\bullet}-\overline{Y}_{2\bullet}
-\overline{Y}_{3\bullet}-\overline{Y}_{4\bullet}=-6.526
\] \[
\hat{SE}\left\{\hat{L}_1\right\}=\sqrt{MSE\sum_{i=1}^4\frac{a_i^2}{n_i}}=
\sqrt{4.430(0.75)}=1.823
\] \[
95\% \mbox{ CI for } L_1: \quad -6.526 \pm 2.000(1.823)  \equiv
-6.526 \pm 3.646  \equiv  (-10.172,-2.880)
\] \[
\mbox{Test Statistic:} \quad t_1^* = \frac{-6.526}{1.823}=-3.580
\] \[
\mbox{Rejection Region:}\quad \left|t_1^*\right|\ge 2.000
\qquad
P=2P\left(t_{60}\ge |-3.580|\right)=.0007
\]

\[
SSL_1 = \frac{\left(\hat{L}_1\right)^2}{\sum_{i=1}^r\frac{a_i^2}{n_i}}=
\frac{(-6.526)^2}{0.75}=56.785
\]

\[
\mbox{Test Statistic:} \quad F_1^* = \frac{SSL_1}{MSE}=\frac{56.785}{4.430}=12.818 
\]

\[
\mbox{Rejection Region:}\quad F_1^*\ge F_{.95,1,60}=4.001
\qquad
P=P\left(F_{1,60}\ge 12.818\right)=.0007
\]

Without going through the calculations, we obtain the following values
for \(L_2\) and \(L_3\).

\[
L_2: \quad \hat{L}_2=-1.805 \quad \hat{SE}\left\{\hat{L}_2\right\}=1.289 \quad
SSL_2=8.688
\] \[
L_3: \quad \hat{L}_3=-0.139 \quad \hat{SE}\left\{\hat{L}_3\right\}=0.744 \quad
SSL_3=0.155
\] The sums of squares for the three pairwise orthogonal contrasts sum
up to the treatment sum of squares.

\[
SSL_1+SSL_2+SSL_3=56.785+8.688+0.155=65.628=SSTR
\] The main take away is that the three virtual training methods (as a
group) have significantly higher scores than the Lecture/Material
control group. There appears to be no significant differences among the
virtual training methods.

\begin{longtable}[]{@{}
  >{\raggedright\arraybackslash}p{(\linewidth - 22\tabcolsep) * \real{0.1818}}
  >{\raggedright\arraybackslash}p{(\linewidth - 22\tabcolsep) * \real{0.0390}}
  >{\raggedright\arraybackslash}p{(\linewidth - 22\tabcolsep) * \real{0.0779}}
  >{\raggedright\arraybackslash}p{(\linewidth - 22\tabcolsep) * \real{0.0390}}
  >{\raggedright\arraybackslash}p{(\linewidth - 22\tabcolsep) * \real{0.0390}}
  >{\raggedright\arraybackslash}p{(\linewidth - 22\tabcolsep) * \real{0.0390}}
  >{\raggedright\arraybackslash}p{(\linewidth - 22\tabcolsep) * \real{0.0909}}
  >{\raggedright\arraybackslash}p{(\linewidth - 22\tabcolsep) * \real{0.0909}}
  >{\raggedright\arraybackslash}p{(\linewidth - 22\tabcolsep) * \real{0.0909}}
  >{\raggedright\arraybackslash}p{(\linewidth - 22\tabcolsep) * \real{0.1039}}
  >{\raggedright\arraybackslash}p{(\linewidth - 22\tabcolsep) * \real{0.1039}}
  >{\raggedright\arraybackslash}p{(\linewidth - 22\tabcolsep) * \real{0.1039}}@{}}
\caption{Lifeboat Virtual Training Contrast Calculations}\tabularnewline
\toprule\noalign{}
\begin{minipage}[b]{\linewidth}\raggedright
\end{minipage} & \begin{minipage}[b]{\linewidth}\raggedright
n
\end{minipage} & \begin{minipage}[b]{\linewidth}\raggedright
Mean
\end{minipage} & \begin{minipage}[b]{\linewidth}\raggedright
a
\end{minipage} & \begin{minipage}[b]{\linewidth}\raggedright
b
\end{minipage} & \begin{minipage}[b]{\linewidth}\raggedright
d
\end{minipage} & \begin{minipage}[b]{\linewidth}\raggedright
L1
\end{minipage} & \begin{minipage}[b]{\linewidth}\raggedright
L2
\end{minipage} & \begin{minipage}[b]{\linewidth}\raggedright
L3
\end{minipage} & \begin{minipage}[b]{\linewidth}\raggedright
\(a^2/n\)
\end{minipage} & \begin{minipage}[b]{\linewidth}\raggedright
\(b^2/n\)
\end{minipage} & \begin{minipage}[b]{\linewidth}\raggedright
\(d^2/n\)
\end{minipage} \\
\midrule\noalign{}
\endfirsthead
\toprule\noalign{}
\begin{minipage}[b]{\linewidth}\raggedright
\end{minipage} & \begin{minipage}[b]{\linewidth}\raggedright
n
\end{minipage} & \begin{minipage}[b]{\linewidth}\raggedright
Mean
\end{minipage} & \begin{minipage}[b]{\linewidth}\raggedright
a
\end{minipage} & \begin{minipage}[b]{\linewidth}\raggedright
b
\end{minipage} & \begin{minipage}[b]{\linewidth}\raggedright
d
\end{minipage} & \begin{minipage}[b]{\linewidth}\raggedright
L1
\end{minipage} & \begin{minipage}[b]{\linewidth}\raggedright
L2
\end{minipage} & \begin{minipage}[b]{\linewidth}\raggedright
L3
\end{minipage} & \begin{minipage}[b]{\linewidth}\raggedright
\(a^2/n\)
\end{minipage} & \begin{minipage}[b]{\linewidth}\raggedright
\(b^2/n\)
\end{minipage} & \begin{minipage}[b]{\linewidth}\raggedright
\(d^2/n\)
\end{minipage} \\
\midrule\noalign{}
\endhead
\bottomrule\noalign{}
\endlastfoot
LEC/MAT (i=1) & 16 & 4.931 & 3 & 0 & 0 & 14.793 & 0 & 0 & 0.5625 & 0 &
0 \\
MON/KEY (i=2) & 16 & 7.708 & -1 & 2 & 0 & -7.708 & 15.416 & 0 & 0.0625 &
0.25 & 0 \\
HMD/JOY (i=3) & 16 & 6.736 & -1 & -1 & 1 & -6.736 & -6.736 & 6.736 &
0.0625 & 0.0625 & 0.0625 \\
HMD/WEA (i=4) & 16 & 6.875 & -1 & -1 & -1 & -6.875 & -6.875 & 6.875 &
0.0625 & 0.0625 & 0.0625 \\
Sum & 64 & & 0 & 0 & 0 & -6.526 & 1.805 & -0.139 & 0.75 & 0.375 &
0.125 \\
\end{longtable}

While it is possible to construct contrasts directly in R, it is very
easy to compute them directly, as we show below.

\begin{Shaded}
\begin{Highlighting}[]
\NormalTok{vt }\OtherTok{\textless{}{-}} \FunctionTok{read.csv}\NormalTok{(}\StringTok{"data/virtual\_training.csv"}\NormalTok{)}
\FunctionTok{head}\NormalTok{(vt,}\DecValTok{2}\NormalTok{)}
\end{Highlighting}
\end{Shaded}

\begin{verbatim}
##   grp.trt procKnow
## 1       1   4.5614
## 2       1   6.6593
\end{verbatim}

\begin{Shaded}
\begin{Highlighting}[]
\FunctionTok{attach}\NormalTok{(vt)}

\NormalTok{vt.n }\OtherTok{\textless{}{-}} \FunctionTok{as.vector}\NormalTok{(}\FunctionTok{tapply}\NormalTok{(procKnow, grp.trt, length))     }\DocumentationTok{\#\# Obtain n\_i}
\NormalTok{vt.mean }\OtherTok{\textless{}{-}} \FunctionTok{as.vector}\NormalTok{(}\FunctionTok{tapply}\NormalTok{(procKnow, grp.trt, mean))    }\DocumentationTok{\#\# Obtain ybar\_i}
\NormalTok{vt.var }\OtherTok{\textless{}{-}} \FunctionTok{as.vector}\NormalTok{(}\FunctionTok{tapply}\NormalTok{(procKnow, grp.trt, var))      }\DocumentationTok{\#\# Obtain s\_i\^{}2}
\FunctionTok{cbind}\NormalTok{(vt.n, vt.mean, vt.var)}
\end{Highlighting}
\end{Shaded}

\begin{verbatim}
##      vt.n  vt.mean   vt.var
## [1,]   16 4.930550 3.763633
## [2,]   16 7.708337 2.044874
## [3,]   16 6.736106 7.952409
## [4,]   16 6.874994 3.960078
\end{verbatim}

\begin{Shaded}
\begin{Highlighting}[]
\NormalTok{n\_T }\OtherTok{\textless{}{-}} \FunctionTok{sum}\NormalTok{(vt.n)}
\NormalTok{r }\OtherTok{\textless{}{-}} \FunctionTok{length}\NormalTok{(vt.n)}
\NormalTok{ybar. }\OtherTok{\textless{}{-}} \FunctionTok{sum}\NormalTok{(vt.n}\SpecialCharTok{*}\NormalTok{vt.mean)}\SpecialCharTok{/}\NormalTok{n\_T}
\NormalTok{SSTR }\OtherTok{\textless{}{-}} \FunctionTok{sum}\NormalTok{(vt.n}\SpecialCharTok{*}\NormalTok{(vt.mean}\SpecialCharTok{{-}}\NormalTok{ybar.)}\SpecialCharTok{\^{}}\DecValTok{2}\NormalTok{)}
\NormalTok{SSE }\OtherTok{\textless{}{-}} \FunctionTok{sum}\NormalTok{((vt.n}\DecValTok{{-}1}\NormalTok{)}\SpecialCharTok{*}\NormalTok{vt.var)}
\NormalTok{MSE }\OtherTok{\textless{}{-}}\NormalTok{ SSE}\SpecialCharTok{/}\NormalTok{(n\_T}\SpecialCharTok{{-}}\NormalTok{r)}

\NormalTok{L1.a }\OtherTok{\textless{}{-}} \FunctionTok{c}\NormalTok{(}\DecValTok{3}\NormalTok{,}\SpecialCharTok{{-}}\DecValTok{1}\NormalTok{,}\SpecialCharTok{{-}}\DecValTok{1}\NormalTok{,}\SpecialCharTok{{-}}\DecValTok{1}\NormalTok{)}
\NormalTok{L2.b }\OtherTok{\textless{}{-}} \FunctionTok{c}\NormalTok{(}\DecValTok{0}\NormalTok{,}\DecValTok{2}\NormalTok{,}\SpecialCharTok{{-}}\DecValTok{1}\NormalTok{,}\SpecialCharTok{{-}}\DecValTok{1}\NormalTok{)}
\NormalTok{L3.c }\OtherTok{\textless{}{-}} \FunctionTok{c}\NormalTok{(}\DecValTok{0}\NormalTok{,}\DecValTok{0}\NormalTok{,}\DecValTok{1}\NormalTok{,}\SpecialCharTok{{-}}\DecValTok{1}\NormalTok{)}

\NormalTok{L1.hat }\OtherTok{\textless{}{-}} \FunctionTok{sum}\NormalTok{(L1.a}\SpecialCharTok{*}\NormalTok{vt.mean)}
\NormalTok{L2.hat }\OtherTok{\textless{}{-}} \FunctionTok{sum}\NormalTok{(L2.b}\SpecialCharTok{*}\NormalTok{vt.mean)}
\NormalTok{L3.hat }\OtherTok{\textless{}{-}} \FunctionTok{sum}\NormalTok{(L3.c}\SpecialCharTok{*}\NormalTok{vt.mean)}

\NormalTok{SE.L1.hat }\OtherTok{\textless{}{-}} \FunctionTok{sqrt}\NormalTok{(MSE}\SpecialCharTok{*}\FunctionTok{sum}\NormalTok{(L1.a}\SpecialCharTok{\^{}}\DecValTok{2}\SpecialCharTok{/}\NormalTok{vt.n))}
\NormalTok{SE.L2.hat }\OtherTok{\textless{}{-}} \FunctionTok{sqrt}\NormalTok{(MSE}\SpecialCharTok{*}\FunctionTok{sum}\NormalTok{(L2.b}\SpecialCharTok{\^{}}\DecValTok{2}\SpecialCharTok{/}\NormalTok{vt.n))}
\NormalTok{SE.L3.hat }\OtherTok{\textless{}{-}} \FunctionTok{sqrt}\NormalTok{(MSE}\SpecialCharTok{*}\FunctionTok{sum}\NormalTok{(L3.c}\SpecialCharTok{\^{}}\DecValTok{2}\SpecialCharTok{/}\NormalTok{vt.n))}

\NormalTok{SSL}\FloatTok{.1} \OtherTok{\textless{}{-}}\NormalTok{ L1.hat}\SpecialCharTok{\^{}}\DecValTok{2}\SpecialCharTok{/}\FunctionTok{sum}\NormalTok{(L1.a}\SpecialCharTok{\^{}}\DecValTok{2}\SpecialCharTok{/}\NormalTok{vt.n)}
\NormalTok{SSL}\FloatTok{.2} \OtherTok{\textless{}{-}}\NormalTok{ L2.hat}\SpecialCharTok{\^{}}\DecValTok{2}\SpecialCharTok{/}\FunctionTok{sum}\NormalTok{(L2.b}\SpecialCharTok{\^{}}\DecValTok{2}\SpecialCharTok{/}\NormalTok{vt.n)}
\NormalTok{SSL}\FloatTok{.3} \OtherTok{\textless{}{-}}\NormalTok{ L3.hat}\SpecialCharTok{\^{}}\DecValTok{2}\SpecialCharTok{/}\FunctionTok{sum}\NormalTok{(L3.c}\SpecialCharTok{\^{}}\DecValTok{2}\SpecialCharTok{/}\NormalTok{vt.n)}

\NormalTok{L1.out }\OtherTok{\textless{}{-}} \FunctionTok{cbind}\NormalTok{(L1.hat, SE.L1.hat, L1.hat}\SpecialCharTok{/}\NormalTok{SE.L1.hat,}
\NormalTok{                L1.hat}\SpecialCharTok{{-}}\FunctionTok{qt}\NormalTok{(.}\DecValTok{975}\NormalTok{,n\_T}\SpecialCharTok{{-}}\NormalTok{r)}\SpecialCharTok{*}\NormalTok{SE.L1.hat, }
\NormalTok{                L1.hat}\SpecialCharTok{+}\FunctionTok{qt}\NormalTok{(.}\DecValTok{975}\NormalTok{,n\_T}\SpecialCharTok{{-}}\NormalTok{r)}\SpecialCharTok{*}\NormalTok{SE.L1.hat,}
\NormalTok{                SSL}\FloatTok{.1}\NormalTok{, SSL}\FloatTok{.1}\SpecialCharTok{/}\NormalTok{MSE, }\DecValTok{1}\SpecialCharTok{{-}}\FunctionTok{pf}\NormalTok{(SSL}\FloatTok{.1}\SpecialCharTok{/}\NormalTok{MSE,}\DecValTok{1}\NormalTok{,n\_T}\SpecialCharTok{{-}}\NormalTok{r))}
\NormalTok{L2.out }\OtherTok{\textless{}{-}} \FunctionTok{cbind}\NormalTok{(L2.hat, SE.L2.hat, L2.hat}\SpecialCharTok{/}\NormalTok{SE.L2.hat,}
\NormalTok{                L2.hat}\SpecialCharTok{{-}}\FunctionTok{qt}\NormalTok{(.}\DecValTok{975}\NormalTok{,n\_T}\SpecialCharTok{{-}}\NormalTok{r)}\SpecialCharTok{*}\NormalTok{SE.L2.hat, }
\NormalTok{                L2.hat}\SpecialCharTok{+}\FunctionTok{qt}\NormalTok{(.}\DecValTok{975}\NormalTok{,n\_T}\SpecialCharTok{{-}}\NormalTok{r)}\SpecialCharTok{*}\NormalTok{SE.L2.hat,}
\NormalTok{                SSL}\FloatTok{.2}\NormalTok{, SSL}\FloatTok{.2}\SpecialCharTok{/}\NormalTok{MSE, }\DecValTok{1}\SpecialCharTok{{-}}\FunctionTok{pf}\NormalTok{(SSL}\FloatTok{.2}\SpecialCharTok{/}\NormalTok{MSE,}\DecValTok{1}\NormalTok{,n\_T}\SpecialCharTok{{-}}\NormalTok{r))}
\NormalTok{L3.out }\OtherTok{\textless{}{-}} \FunctionTok{cbind}\NormalTok{(L3.hat, SE.L3.hat, L3.hat}\SpecialCharTok{/}\NormalTok{SE.L3.hat,}
\NormalTok{                L3.hat}\SpecialCharTok{{-}}\FunctionTok{qt}\NormalTok{(.}\DecValTok{975}\NormalTok{,n\_T}\SpecialCharTok{{-}}\NormalTok{r)}\SpecialCharTok{*}\NormalTok{SE.L3.hat, }
\NormalTok{                L3.hat}\SpecialCharTok{+}\FunctionTok{qt}\NormalTok{(.}\DecValTok{975}\NormalTok{,n\_T}\SpecialCharTok{{-}}\NormalTok{r)}\SpecialCharTok{*}\NormalTok{SE.L3.hat,}
\NormalTok{                SSL}\FloatTok{.3}\NormalTok{, SSL}\FloatTok{.3}\SpecialCharTok{/}\NormalTok{MSE, }\DecValTok{1}\SpecialCharTok{{-}}\FunctionTok{pf}\NormalTok{(SSL}\FloatTok{.3}\SpecialCharTok{/}\NormalTok{MSE,}\DecValTok{1}\NormalTok{,n\_T}\SpecialCharTok{{-}}\NormalTok{r))}
\NormalTok{L.out }\OtherTok{\textless{}{-}} \FunctionTok{rbind}\NormalTok{(L1.out, L2.out, L3.out)}
\FunctionTok{colnames}\NormalTok{(L.out) }\OtherTok{\textless{}{-}} \FunctionTok{c}\NormalTok{(}\StringTok{"Lhat"}\NormalTok{, }\StringTok{"Std Err"}\NormalTok{, }\StringTok{"t*"}\NormalTok{, }\StringTok{"LB"}\NormalTok{, }\StringTok{"UB"}\NormalTok{, }\StringTok{"SS"}\NormalTok{, }\StringTok{"F*"}\NormalTok{, }\StringTok{"P(\textgreater{}F*)"}\NormalTok{)}
\FunctionTok{rownames}\NormalTok{(L.out) }\OtherTok{\textless{}{-}} \FunctionTok{c}\NormalTok{(}\StringTok{"L1:(3,{-}1,{-}1,{-}1)"}\NormalTok{,}\StringTok{"L2:(0,2,{-}1,{-}1)"}\NormalTok{,}\StringTok{"L3:(0,0,1,{-}1)"}\NormalTok{)}
\FunctionTok{round}\NormalTok{(L.out,}\DecValTok{3}\NormalTok{)}
\end{Highlighting}
\end{Shaded}

\begin{verbatim}
##                   Lhat Std Err     t*      LB     UB     SS     F* P(>F*)
## L1:(3,-1,-1,-1) -6.528   1.823 -3.581 -10.174 -2.882 56.816 12.825  0.001
## L2:(0,2,-1,-1)   1.806   1.289  1.401  -0.773  4.384  8.694  1.962  0.166
## L3:(0,0,1,-1)   -0.139   0.744 -0.187  -1.627  1.350  0.154  0.035  0.853
\end{verbatim}

\begin{Shaded}
\begin{Highlighting}[]
\FunctionTok{detach}\NormalTok{(vt)}
\end{Highlighting}
\end{Shaded}

\[\nabla\]

\section{Simultaneous Comparisons}\label{simultaneous-comparisons}

Often when fitting models, we wish to make multiple comparisons
simultaneously. If we have \(k\) independent comparisons and we make
each one with significance level of \(\alpha\), the chances that all
conclusions are correct is \((1-\alpha)^k\) which is less than
\(1-\alpha\).

Other issues arise when researchers wait until observing the data to
choose which comparisons to make (like comparing the maximum and minimum
observed means). This practice has the effect of inflating the Type I
error rate, \(\alpha\), and is referred to as \textbf{data snooping}.

There are various methods for making simultaneous corrections. One
common method is to do each comparison with \(\alpha^*=\alpha/k\), and
then \(\left(1-\alpha^*\right)^k \ge (1-\alpha)\).

For instance, if we wish to compute \(k=5\) confidence intervals, and we
want them all to contain their true parameter with \(95\%\) confidence,
\(\alpha=0.05\) and \(\alpha^*=\alpha/k=0.01\), then we should construct
99\% confidence intervals for the individual parameters. Note that
\((1-.01)^5=.9510>.95\).

In this section, several procedures are described for making
simultaneous comparisons.

\subsection{Tukey's Honest Significant Difference
(HSD)}\label{tukeys-honest-significant-difference-hsd}

This method is widely used for ANOVA models. It makes use of the
\textbf{Studentized Range Distribution} with critical values given in
many textbooks, online, and on the course slides. In R, the
\textbf{qtukey} function can be used to obtain critical values and the
\textbf{ptukey} function can be used to obtain adjusted \(P\)-values.

The basis for the distribution is as follows, where \(Y_1,\ldots,Y_r\)
are independent and normally distributed with mean \(\mu\) and variance
\(\sigma^2\). Further, let \(s^2\) be an unbiased estimator of
\(\sigma^2\) that is independent of \(Y_1,\ldots,Y_r\), with degrees of
freedom \(\nu\). Note that \(s^2\) cannot be computed from
\(Y_1,\ldots,Y_r\) in this situation, but it does work when we are
comparing treatment/group means below. Let \(w\) be the range of
\(Y_1,\ldots,Y_r\) and \(w/s\) be the studentized range, indexed by the
sample size, \(r\) and the degrees of freedom for \(s\), \(\nu\). We
will use \(q_{1-\alpha;r,\nu}\) represent the \(1-\alpha\) quantile of
the Studentized Range distribution.

\[
w=\max\left(Y_1,\ldots,Y_r\right)-\min\left(Y_1,\ldots,Y_r\right)
\qquad \frac{w}{s}=q(r,\nu)
\]

\[
P\left(\frac{w}{s}=q(r,\nu)\le q_{1-\alpha;r,\nu}\right)=1-\alpha
\] \[
\quad \Rightarrow \quad P\left(\frac{\left|Y_i-Y_{i'}\right|}{s}
\le q_{1-\alpha;r,\nu}\right)=1-\alpha \quad \mbox{for all } i,i'
\]

Now, consider making all pairwise comparisons among means under the
assumption of equal means \(\left(\mu_1=\cdots = \mu_r = \mu\right)\)
and equal sample sizes \(\left(n_1=\cdots = n_r=n\right)\).

Then, we have the following results for the sampling distribution of
\(\overline{Y}_{i\bullet}\) and its estimated variance, which is
unbiased with \(n_T-r\) degrees of freedom. Furthermore the sample means
\(\left\{\overline{Y}_{i\bullet}\right\}\) and \(MSE\) are independent.

\[
\overline{Y}_{i\bullet} \sim N\left(\mu,\frac{\sigma^2}{n}\right)
\qquad \qquad
\hat{V}\left\{\overline{Y}_{i\bullet}\right\}=\frac{MSE}{n}
\] \[
\quad \Rightarrow \quad 
P\left(\frac{\left|\overline{Y}_{i\bullet}-
\overline{Y}_{i'\bullet}\right|}{\sqrt{MSE/n}}
\le q_{1-\alpha;r,\nu}\right)=1-\alpha \quad \mbox{for all } i,i'
\]

Then we conclude \(\mu_i \ne \mu_{i'}\) if the following criteria is
met.

\[
\frac{\left|\overline{Y}_{i\bullet}-
\overline{Y}_{i'\bullet}\right|}{\sqrt{MSE/n}}
\ge q_{1-\alpha;r,\nu}
\quad \Rightarrow \quad
\left|\overline{Y}_{i\bullet}-
\overline{Y}_{i'\bullet}\right| 
\ge q_{1-\alpha;r,\nu}\sqrt{\frac{MSE}{n}}
\] Subsequently this method was generalized to allow for unequal sample
sizes with the following decision rule. Conclude \(\mu_i \ne \mu_{i'}\)
if the following result holds.

\[
\left|\overline{Y}_{i\bullet}-
\overline{Y}_{i'\bullet}\right| 
\ge \frac{q_{1-\alpha;r,\nu}}{\sqrt{2}}\sqrt{MSE
\left(\frac{1}{n_i}+\frac{1}{n_{i'}}\right)}
\] An equivalent way of writing this test, in terms of
\(D=\mu_i-\mu_{i'}\) and
\(\hat{D}=\overline{Y}_{i\bullet}-\overline{Y}_{i'\bullet}\) is to
reject \(H_0: D=0\) based on the following test.

\[
\mbox{Test Statistic:} \quad 
q^*=\frac{\sqrt{2}\hat{D}}{\hat{SE}\left\{\hat{D}\right\}} \qquad
\mbox{Rejection Region:} \quad \left|q^*\right| \ge
q_{1-\alpha;r,\nu}
\] The adjusted \(P\)-value is the area above \(\left|q^*\right|\) in
the Studentized Range distribution.

To obtain simultaneous \((1-\alpha)100\%\) Confidence Intervals for all
\(r(r-1)/2\) mean differences we can simply invert the test as follows.

\[
(1-\alpha)100\% \mbox{ for } \mu_i-\mu_{i'}: \quad
\left(\overline{Y}_{i\bullet}-
\overline{Y}_{i'\bullet}\right) \pm
\frac{1}{\sqrt{2}}q_{1-\alpha;r,n_T-r}
\sqrt{MSE
\left(\frac{1}{n_i}+\frac{1}{n_{i'}}\right)}
\]

\textbf{Example 3.2 - Virtual Training for a Lifeboat Launching Task}

First, we obtain \(r=4\), \(\nu=64-4=60\), and \(q_{.95;4,60}=3.737\)
for this study. Using calculations from before, we have the following
results.

\[
n_1=n_2=n_3=n_4=16 \qquad \overline{y}_{1\bullet}=4.931 \quad
\overline{y}_{2\bullet}=7.708 \quad
\overline{y}_{3\bullet}=6.736 \quad
\overline{y}_{4\bullet}=6.875 
\]

\[
MSE = 4.430 \qquad \hat{SE}\left\{\hat{D}\right\}=
\sqrt{4.430\left(\frac{1}{16}+\frac{1}{16}\right)}=0.744
\]

Next we compute the Tukey HSD for all pairs of means (as this is a
balanced design).

\[
HSD_{ii'}=\frac{3.737}{\sqrt{2}}(0.744)=2.642(0.744)=1.966
\]

That is, any pairs of means that differ by more than 1.966 in absolute
value will be declared significantly different. We also obtain
simultaneous 95\% Confidence Intervals by adding and subtracting 1.966
to each difference \(\hat{D}\). Here we compute the values for the
comparison of treatments 1 (LEC/MAT) and 2 (MON/KEY). Then we will use
the R function \textbf{TukeyHSD} after fitting the model with an
\textbf{aov} object.

\[
\hat{D}_{21}=\overline{Y}_{2\bullet}-\overline{Y}_{1\bullet}=
7.708-4.931=2.777
\]

\[
95\% \mbox{ Simultaneous CI for } \mu_2-\mu_1: \quad
2.777 \pm 1.966 \quad \equiv \quad (0.811,4.743) 
\]

Thus, we can conclude that the Monitor/Keyboard condition provides
higher scores on average than the Lecture/Material condition. That is
\(\mu_2>\mu_1\). We next obtain the adjusted \(P\)-value, making use of
the \textbf{ptukey} function in R. First, though, we have to scale it
back to a studentized range, \(q^*\).

\[
q^* = \frac{\sqrt{2}\hat{D}}{\hat{SE}\left\{\hat{D}\right\}}=
\frac{\sqrt{2}(2.777)}{0.744}=5.279  \qquad 
P= P\left(q_{4,60}\ge |5.279|\right)=.0023
\]

Now, we use R to run all possible comparisons. Note that we have to fit
an \textbf{aov} object and not a \textbf{lm} object to use the
\textbf{TukeyHSD} function.

\begin{Shaded}
\begin{Highlighting}[]
\NormalTok{vt }\OtherTok{\textless{}{-}} \FunctionTok{read.csv}\NormalTok{(}\StringTok{"data/virtual\_training.csv"}\NormalTok{)}
\FunctionTok{head}\NormalTok{(vt,}\DecValTok{2}\NormalTok{)}
\end{Highlighting}
\end{Shaded}

\begin{verbatim}
##   grp.trt procKnow
## 1       1   4.5614
## 2       1   6.6593
\end{verbatim}

\begin{Shaded}
\begin{Highlighting}[]
\NormalTok{vt}\SpecialCharTok{$}\NormalTok{grp.trt.f }\OtherTok{\textless{}{-}} \FunctionTok{factor}\NormalTok{(vt}\SpecialCharTok{$}\NormalTok{grp.trt, }\AttributeTok{levels=}\DecValTok{1}\SpecialCharTok{:}\DecValTok{4}\NormalTok{,}
                       \AttributeTok{labels=}\FunctionTok{c}\NormalTok{(}\StringTok{"LEC/MAT"}\NormalTok{,}\StringTok{"MON/KEY"}\NormalTok{,}\StringTok{"HMD/JOY"}\NormalTok{, }\StringTok{"HMD/WEA"}\NormalTok{))}

\NormalTok{vt.mod1 }\OtherTok{\textless{}{-}} \FunctionTok{aov}\NormalTok{(procKnow }\SpecialCharTok{\textasciitilde{}}\NormalTok{ grp.trt.f, }\AttributeTok{data=}\NormalTok{vt)}
\FunctionTok{summary.lm}\NormalTok{(vt.mod1)}
\end{Highlighting}
\end{Shaded}

\begin{verbatim}
## 
## Call:
## aov(formula = procKnow ~ grp.trt.f, data = vt)
## 
## Residuals:
##     Min      1Q  Median      3Q     Max 
## -5.0831 -1.4444  0.5774  1.5495  3.1900 
## 
## Coefficients:
##                  Estimate Std. Error t value Pr(>|t|)    
## (Intercept)        4.9306     0.5262   9.370 2.37e-13 ***
## grp.trt.fMON/KEY   2.7778     0.7442   3.733 0.000423 ***
## grp.trt.fHMD/JOY   1.8056     0.7442   2.426 0.018277 *  
## grp.trt.fHMD/WEA   1.9444     0.7442   2.613 0.011329 *  
## ---
## Signif. codes:  0 '***' 0.001 '**' 0.01 '*' 0.05 '.' 0.1 ' ' 1
## 
## Residual standard error: 2.105 on 60 degrees of freedom
## Multiple R-squared:  0.1981, Adjusted R-squared:  0.158 
## F-statistic: 4.941 on 3 and 60 DF,  p-value: 0.003931
\end{verbatim}

\begin{Shaded}
\begin{Highlighting}[]
\FunctionTok{anova}\NormalTok{(vt.mod1)}
\end{Highlighting}
\end{Shaded}

\begin{verbatim}
## Analysis of Variance Table
## 
## Response: procKnow
##           Df  Sum Sq Mean Sq F value   Pr(>F)   
## grp.trt.f  3  65.664 21.8880  4.9406 0.003931 **
## Residuals 60 265.815  4.4302                    
## ---
## Signif. codes:  0 '***' 0.001 '**' 0.01 '*' 0.05 '.' 0.1 ' ' 1
\end{verbatim}

\begin{Shaded}
\begin{Highlighting}[]
\NormalTok{HSD.vt }\OtherTok{\textless{}{-}} \FunctionTok{TukeyHSD}\NormalTok{(vt.mod1)}
\NormalTok{HSD.vt}
\end{Highlighting}
\end{Shaded}

\begin{verbatim}
##   Tukey multiple comparisons of means
##     95% family-wise confidence level
## 
## Fit: aov(formula = procKnow ~ grp.trt.f, data = vt)
## 
## $grp.trt.f
##                       diff        lwr       upr     p adj
## MON/KEY-LEC/MAT  2.7777875  0.8113167 4.7442583 0.0023332
## HMD/JOY-LEC/MAT  1.8055562 -0.1609146 3.7720271 0.0829543
## HMD/WEA-LEC/MAT  1.9444437 -0.0220271 3.9109146 0.0537097
## HMD/JOY-MON/KEY -0.9722312 -2.9387021 0.9942396 0.5624876
## HMD/WEA-MON/KEY -0.8333437 -2.7998146 1.1331271 0.6788308
## HMD/WEA-HMD/JOY  0.1388875 -1.8275833 2.1053583 0.9976681
\end{verbatim}

\begin{Shaded}
\begin{Highlighting}[]
\FunctionTok{plot}\NormalTok{(HSD.vt, }\AttributeTok{las=}\DecValTok{1}\NormalTok{, }\AttributeTok{cex.axis=}\FloatTok{0.4}\NormalTok{, }\AttributeTok{col=}\StringTok{"purple"}\NormalTok{)}
\end{Highlighting}
\end{Shaded}

\pandocbounded{\includegraphics[keepaspectratio]{bookdownproj_STA4211Notes1_files/figure-latex/ch0303-1.pdf}}

Note that only treatments 2 and 1 are significantly different. The
Confidence Interval does not contain 0 and the adjusted \(P\)-value is
below 0.05. Treatments 1 and 4 are close to being significantly
different, but the Confidence Interval contains 0 and the \(P\)-value is
above 0.05.

\[\nabla\]

\subsection{Scheffe's Method for Multiple
Comparisons}\label{scheffes-method-for-multiple-comparisons}

This method is very conservative, but due to its procedure, it can be
applied to all possible contrasts among treatment means. Using previous
results on contrasts, we have the following set up for the procedure.

\[
L=\sum_{i=1}^rc_i \mu_i \qquad \mbox{ such that } \sum_{i=1}^rc_i=0
\]

\[
\hat{L}=\sum_{i=1}^rc_i\overline{Y}_{i\bullet}  \qquad
\hat{SE}\left\{\hat{L}\right\}=\sqrt{MSE\sum_{i=1}^r\frac{c_i^2}{n_i}}
\]

Then we construct simultaneous intervals of the following form.

\[
(1-\alpha)100\% \mbox{ Simultaneous CI for }L: \quad
\hat{L} \pm \left(\sqrt{(r-1)F_{1-\alpha;r-1,n_T-r}}\right)
\hat{SE}\left\{\hat{L}\right\}
\] The critical \(F\)-value is the same critical value for testing
\(H_0:\mu_1=\cdots = \mu_r\) for the Analysis of Variance. Note that
when this hypothesis is true, all contrasts are 0.

We can also simultaneously test whether a set of contrasts are 0.

\(H_0:L=\sum_{i=1}^r c_i\mu_i=0\) vs
\(H_A:L=\sum_{i=1}^r c_i\mu_i \ne 0\)

\[
\mbox{Test Statistic:} \quad 
F^*=\frac{\left(\hat{L}\right)^2}{(r-1)\hat{V}\left\{\hat{L}\right\}}
\qquad
\mbox{Rejection Region:} \quad F^*\geq F_{1-\alpha;r-1,n_T-r}
\]

\textbf{Example 3.3 - Virtual Training for a Lifeboat Launching Task}

Consider the three pairwise orthogonal contrasts that we previously
estimated.

\begin{itemize}
\tightlist
\item
  \(L_1\): Lecture/Materials vs Virtual Training
\item
  \(L_2\): Monitor/Keyboard versus Head Mounted Displays
\item
  \(L_3\): HMD/Joypad versus HMD/Wearables
\end{itemize}

The critical \(F\)-value we need is \(F_{.95;3,60}=2.758\), leading to
the following critical value for the Scheffe simultaneous confidence
intervals.

\[
\sqrt{(r-1)F_{1-\alpha;r-1,n_T-r}}=
\sqrt{(4-1)2.758}=2.876
\] Next, we give the estimated contrasts, their standard errors, and
simultaneous 95\% Confidence Intervals for the population based
contrasts.

\[
\hat{L}_1=-6.526 \qquad \hat{SE}\left\{\hat{L}_1\right\}=1.823 
\] \[
\hat{L}_2=1.805 \qquad \hat{SE}\left\{\hat{L}_2\right\}=1.289 
\]

\[
\hat{L}_3=-0.139 \qquad \hat{SE}\left\{\hat{L}_3\right\}=0.744 
\]

\[
95\% \mbox{ CI for: }L_1=3\mu_1-\mu_2-\mu_3-\mu_4:
\quad -6.526 \pm 2.876(1.823)
\]

\[
\equiv -6.526\pm 5.243 \equiv
(-11.769,-1.283)
\]

\[
95\% \mbox{ CI for: }L_2=2\mu_2-\mu_3-\mu_4:
\quad 1.805 \pm 2.876(1.289)
\]

\[
\equiv 1.805\pm 3.707 \equiv
(-1.902,5.512)
\]

\[
95\% \mbox{ CI for: }L_3=\mu_3-\mu_4:
\quad -0.139 \pm 2.876(0.744)
\]

\[
\equiv -0.139\pm 2.140 \equiv
(-2.779,2.001)
\]

As a comparison between no adjustment for multiple Confidence Intervals
and the Scheffe adjustment, consider the case of \(L_1\). When there was
no adjustment, the interval was \((-10.172,-2.880)\). When the Scheffe
adjustment was used, it was \((-11.769,-1.283)\). Again, Scheffe's
method can be used for all possible contrasts simultaneously.

\[\nabla\]

\subsection{Bonferroni Method for Multiple
Comparisons}\label{bonferroni-method-for-multiple-comparisons}

The Bonferroni method is widely used in many situations where multiple
tests are conducted (not just among treatment means). If we have \(g\)
pre-planned comparisons or contrasts to be made, we make each comparison
at the \(\alpha^*=\alpha/g\) significance level. Note that a special
case of this is making \(g=r(r-1)/2\) comparisons among all pairs of
treatment means. In that case, Tukey's HSD will provide more powerful
tests than the Bonferroni method (unless \(r\)=2, and there is only one
comparison).

\[
L=\sum_{i=1}^rc_i \mu_i \qquad \mbox{ such that } \sum_{i=1}^rc_i=0
\]

\[
\hat{L}=\sum_{i=1}^rc_i\overline{Y}_{i\bullet}  \qquad
\hat{SE}\left\{\hat{L}\right\}=\sqrt{MSE\sum_{i=1}^r\frac{c_i^2}{n_i}}
\]

Then we construct simultaneous confidence intervals of the following
form.

\[
(1-\alpha)100\% \mbox{ Simultaneous CI for }L: \quad
\hat{L} \pm t_{1-\alpha/(2g);n_T-r}
\hat{SE}\left\{\hat{L}\right\}
\]

Tables that give critical \(t\)-values for the Bonferroni method are
available in textbooks, online, and on the course slides.

Simultaneous tests for \(g\) pre-planned contrasts can be conducted to
test whether the contrast is equal to zero.

\[
H_0:L=\sum_{i=1}^rc_i\mu_i=0 \qquad H_A:L=\sum_{i=1}^rc_i\mu_i \ne 0
\] \[
\mbox{Test Statistic:} \quad t^*=\frac{\hat{L}}{\hat{SE}\left\{\hat{L}\right\}}
\qquad \mbox{Rejection Region:} \quad \left|t^*\right| \ge t_{1-\alpha/(2g),n_T-r}
\] \[
\mbox{Adjusted P-value:} \quad P=\min\left(1,g\times 2P\left(t_{n_T-r}\ge \left|t^*\right|\right)\right)
\]

\textbf{Example 3.4 - Virtual Training for a Lifeboat Launching Task}

Consider the \(g=3\) pairwise orthogonal contrasts that we previously
estimated.

\begin{itemize}
\tightlist
\item
  \(L_1\) Lecture/Materials vs Virtual Training
\item
  \(L_2\) Monitor/Keyboard versus Head Mounted Displays
\item
  \(L_3\) HMD/Joypad versus HMD/Wearables
\end{itemize}

The critical \(t\)-value we need is \(t_{1-.05/(2(3));60}=2.463\). We
will conduct the three simultaneous \(t\)-tests and obtain their
adjusted \(P\)-values.

Next, we give the estimated contrasts, their standard errors, and
simultaneous \(t\)-tests with an overall experimentwise error rate of
\(\alpha=0.05\) for the population based contrasts. Note that the
individual tests are being conducted at \(\alpha^*=.05/3=.0167\) error
rate.

\[
\hat{L}_1=-6.526 \qquad \hat{SE}\left\{\hat{L}_1\right\}=1.823 \qquad t_1^*=\frac{-6.526}{1.823}=-3.580
\] \[
\hat{L}_2=1.805 \qquad \hat{SE}\left\{\hat{L}_2\right\}=1.289 \qquad t_2^*=\frac{1.805}{1.289}=1.400
\]

\[
\hat{L}_3=-0.139 \qquad \hat{SE}\left\{\hat{L}_3\right\}=0.744 \qquad t_3^*=\frac{-0.139}{0.744}=-0.187
\] Thus, again, only the first contrast is significant. We compute the 3
adjusted \(P\)-values by first obtaining the 2-sided tail areas for the
\(g=3\) \(t\)-statistics, then applying the Bonferroni adjustments.

\[
2P\left(t_{60}\ge|3.580|\right) = 2(.00034)=.0007
\qquad P_1 = \min\left(1,3(.0007)\right) = .0021
\] \[
2P\left(t_{60}\ge|1.400|\right) = 2(.0833)=.1666
\qquad
P_2 = \min\left(1,3(.1666)\right) = .4998
\]

\[
2P\left(t_{60}\ge|-0.187|\right) = 2(.4261)=.8522
\qquad
P_3 = \min\left(1,3(.8522)\right) = 1
\]

Note that if we used the Bonferroni method for comparing all
\(g=4(4-1)/2=6\) pairs of treatment means, the critical \(t\)-value
would be \(t_{1-.05/(2(6));60}=2.729\). Since the standard error for the
difference in two mean for this example is
\(\hat{SE}\left\{\hat{D}\right\}=0.744\), the simultaneous confidence
intervals would be of the following form.

\[
\hat{D}_{ii'} \pm 2.729(0.744) \quad \equiv \quad \hat{D}_{ii'} \pm 2.030 
\]

These intervals are slightly wider than the Tukey based intervals
\(\left(\hat{D}_{ii'} \pm 1.966\right)\). The same conclusions among
treatment means are made.

\[\nabla\]

\subsection{Multiple Range Methods}\label{multiple-range-methods}

In this section, we briefly describe two methods that use multiple
ranges for comparisons: \textbf{Student-Newman-Keuls'} (often shortened
to SNK) and \textbf{Duncan's} methods. They both are based on ordering
the observed means. Neither can be used to construct simultaneous
confidence intervals.

\[
\overline{Y}_{(1)\bullet} \le \cdots \le \overline{Y}_{(r)\bullet} \qquad n_{(i)} \mbox{ is the sample size for the ordered groups}
\]

Each method uses different critical values for testing pairs of
treatment means, depending on how many means are between them. They both
begin by comparing the extreme means, which encompass \(r\) means.

For the SNK method, we reject \(H_0:\mu_{(i)}-\mu_{(i')}=0\) where
\(i>i'\) if the following result holds, where \(q_{1-\alpha;k,n_T-r}\)
is obtained from the studentized range distribution.

\[
\overline{Y}_{(i)\bullet}- \overline{Y}_{(i')\bullet} \ge
\frac{q_{1-\alpha;k,n_T-r}}{\sqrt{2}}
\sqrt{MSE\left(\frac{1}{n_{(i)}}+\frac{1}{n_{(i')}}\right)} \qquad k=i-i'+1
\]

If we fail to reject this hypothesis, no means between these are
compared.

For Duncan's method, the error rates are increased as the means are
farther apart in terms of the number of means between them. This makes
it ``easier'' to detect differences at the cost of increasing overall
error rates. In this case we reject \(H_0:\mu_{(i)}-\mu_{(i')}=0\) if
the following result holds. As with the SNK method, if two means are not
significantly different, no means between them are compared.

\[
\overline{Y}_{(i)\bullet}- \overline{Y}_{(i')\bullet} \ge
\frac{q_{(1-\alpha)^{k-1};k,n_T-r}}{\sqrt{2}}
\sqrt{MSE\left(\frac{1}{n_{(i)}}+\frac{1}{n_{(i')}}\right)} \qquad k=i-i'+1
\]

\textbf{Example 3.5 - Virtual Training for a Lifeboat Launching Task}

For this example, we have the following values.

\[
r=4 \qquad n_1=\cdots=n_r=16 \qquad MSE=4.430 \qquad \hat{SE}\left\{\hat{D}\right\}=
\sqrt{4.430\left(\frac{1}{16}+\frac{1}{16}\right)}=0.744
\]

\[
\overline{Y}_{(1)\bullet}=4.931 \quad
\overline{Y}_{(2)\bullet}=6.736 \quad
\overline{Y}_{(3)\bullet}=6.875 \quad
\overline{Y}_{(4)\bullet}=7.708 
\] We first test \(H_0:\mu_{(4)}-\mu_{(1)}=0\) with range \(k=4-1+1=4\).
Using the \textbf{qtukey} function in R, we obtain the following
critical values and rejection regions (RR) based on the studentized
range distribution with overall experimentwise error rate
\(\alpha=.05\).

\[
\mbox{SNK:} \quad q(.95;4,60)=3.737   \qquad \mbox{RR:} \quad
\overline{Y}_{(4)\bullet}- \overline{Y}_{(1)\bullet} \ge
\frac{3.737}{\sqrt{2}}(0.744)=1.966
\]

\[
\mbox{Duncan:} \quad q((.95)^{4-1};4,60)=3.073   \qquad \mbox{RR:} \quad
\overline{Y}_{(4)\bullet}- \overline{Y}_{(1)\bullet} \ge
\frac{3.073}{\sqrt{2}}(0.744)=1.617
\] The observed difference is
\(\overline{y}_{(4)\bullet}- \overline{y}_{(1)\bullet}=7.708-4.931=2.777\).
Both methods conclude these means are significantly different.

Next, we obtain the critical values and rejection regions for testing
\(H_0:\mu_{(4)}-\mu_{(2)}=0\) and \(H_0:\mu_{(3)}-\mu_{(1)}=0\), each
with a range of \(k=3\).

\[
\mbox{SNK:} \quad q(.95;3,60)=3.399   \qquad \mbox{RR:} \quad
\overline{Y}_{(i)\bullet}- \overline{Y}_{(i')\bullet} \ge
\frac{3.399}{\sqrt{2}}(0.744)=1.788
\]

\[
\mbox{Duncan:} \quad q((.95)^{3-1};3,60)=2.976   \qquad \mbox{RR:} \quad
\overline{Y}_{(i)\bullet}- \overline{Y}_{(i')\bullet} \ge
\frac{2.976}{\sqrt{2}}(0.744)=1.566
\]

The observed differences are as follow.

\[
\overline{y}_{(4)\bullet}- \overline{y}_{(2)\bullet}=7.708-6.736=0.972 \qquad
\overline{y}_{(3)\bullet}- \overline{y}_{(1)\bullet}=6.875-4.931=1.944
\]

Neither test rejects \(H_0:\mu_{(4)}-\mu_{(2)}=0\) and both reject
\(H_0:\mu_{(3)}-\mu_{(1)}=0\). Thus, we do not compare \(\mu_{(4)}\) and
\(\mu_{(3)}\), or compare \(\mu_{(3)}\) and \(\mu_{(2)}\), as both of
these are between \(\mu_{(4)}\) and \(\mu_{(2)}\) . We will conduct the
following test though, with \(k=2\), \(H_0:\mu_{(2)}-\mu_{(1)}=0\)

\[
\mbox{SNK:} \quad q(.95,2,60)=2.829   \qquad \mbox{RR:} \quad
\overline{Y}_{(i)\bullet}- \overline{Y}_{(i')\bullet} \ge
\frac{2.829}{\sqrt{2}}(0.744)=1.488
\]

\[
\mbox{Duncan:} \quad q((.95)^{2-1},2,60)=2.829   \qquad \mbox{RR:} \quad
\overline{Y}_{(2)\bullet}- \overline{Y}_{(1)\bullet} \ge
\frac{2.829}{\sqrt{2}}(0.744)=1.488
\]

The observed differences are as follow.

\[
\overline{y}_{(2)\bullet}- \overline{y}_{(1)\bullet}=6.736-4.931=1.805
\]

Thus, we reject \(H_0:\mu_{(2)}-\mu_{(1)}=0\).

Based on both the SNK and Duncan methods, we conclude that all three
virtual reality methods perform better than the control
(Lecture/Material) method. These methods are available in the
\textbf{agricolae} package in R. Note that the treatment factor must be
defined as a factor outside of \textbf{aov} or \textbf{lm} object to be
used in the \textbf{SNK.test} and \textbf{duncan.test} functions.

\begin{Shaded}
\begin{Highlighting}[]
\NormalTok{vt }\OtherTok{\textless{}{-}} \FunctionTok{read.csv}\NormalTok{(}\StringTok{"data/virtual\_training.csv"}\NormalTok{)}
\FunctionTok{head}\NormalTok{(vt,}\DecValTok{2}\NormalTok{)}
\end{Highlighting}
\end{Shaded}

\begin{verbatim}
##   grp.trt procKnow
## 1       1   4.5614
## 2       1   6.6593
\end{verbatim}

\begin{Shaded}
\begin{Highlighting}[]
\NormalTok{vt}\SpecialCharTok{$}\NormalTok{grp.trt.f }\OtherTok{\textless{}{-}} \FunctionTok{factor}\NormalTok{(vt}\SpecialCharTok{$}\NormalTok{grp.trt)}
\NormalTok{vt.mod1 }\OtherTok{\textless{}{-}} \FunctionTok{lm}\NormalTok{(procKnow }\SpecialCharTok{\textasciitilde{}}\NormalTok{ grp.trt.f, }\AttributeTok{data=}\NormalTok{vt)}
\FunctionTok{summary}\NormalTok{(vt.mod1)}
\end{Highlighting}
\end{Shaded}

\begin{verbatim}
## 
## Call:
## lm(formula = procKnow ~ grp.trt.f, data = vt)
## 
## Residuals:
##     Min      1Q  Median      3Q     Max 
## -5.0831 -1.4444  0.5774  1.5495  3.1900 
## 
## Coefficients:
##             Estimate Std. Error t value Pr(>|t|)    
## (Intercept)   4.9306     0.5262   9.370 2.37e-13 ***
## grp.trt.f2    2.7778     0.7442   3.733 0.000423 ***
## grp.trt.f3    1.8056     0.7442   2.426 0.018277 *  
## grp.trt.f4    1.9444     0.7442   2.613 0.011329 *  
## ---
## Signif. codes:  0 '***' 0.001 '**' 0.01 '*' 0.05 '.' 0.1 ' ' 1
## 
## Residual standard error: 2.105 on 60 degrees of freedom
## Multiple R-squared:  0.1981, Adjusted R-squared:  0.158 
## F-statistic: 4.941 on 3 and 60 DF,  p-value: 0.003931
\end{verbatim}

\begin{Shaded}
\begin{Highlighting}[]
\FunctionTok{anova}\NormalTok{(vt.mod1)}
\end{Highlighting}
\end{Shaded}

\begin{verbatim}
## Analysis of Variance Table
## 
## Response: procKnow
##           Df  Sum Sq Mean Sq F value   Pr(>F)   
## grp.trt.f  3  65.664 21.8880  4.9406 0.003931 **
## Residuals 60 265.815  4.4302                    
## ---
## Signif. codes:  0 '***' 0.001 '**' 0.01 '*' 0.05 '.' 0.1 ' ' 1
\end{verbatim}

\begin{Shaded}
\begin{Highlighting}[]
\FunctionTok{library}\NormalTok{(agricolae)}
\FunctionTok{SNK.test}\NormalTok{(vt.mod1, }\StringTok{"grp.trt.f"}\NormalTok{, }\AttributeTok{console=}\ConstantTok{TRUE}\NormalTok{)}
\end{Highlighting}
\end{Shaded}

\begin{verbatim}
## 
## Study: vt.mod1 ~ "grp.trt.f"
## 
## Student Newman Keuls Test
## for procKnow 
## 
## Mean Square Error:  4.430248 
## 
## grp.trt.f,  means
## 
##   procKnow      std  r        se    Min    Max      Q25     Q50      Q75
## 1 4.930550 1.940008 16 0.5262039 0.3268 6.7509 4.230175 5.67085 6.375725
## 2 7.708337 1.429991 16 0.5262039 5.0733 9.7848 6.682275 8.33320 8.707800
## 3 6.736106 2.820002 16 0.5262039 1.6530 9.9261 5.298875 7.04235 9.101975
## 4 6.874994 1.989995 16 0.5262039 3.2286 9.8754 5.344875 6.84850 8.346350
## 
## Alpha: 0.05 ; DF Error: 60 
## 
## Critical Range
##        2        3        4 
## 1.488551 1.788389 1.966471 
## 
## Means with the same letter are not significantly different.
## 
##   procKnow groups
## 2 7.708337      a
## 4 6.874994      a
## 3 6.736106      a
## 1 4.930550      b
\end{verbatim}

\begin{Shaded}
\begin{Highlighting}[]
\FunctionTok{duncan.test}\NormalTok{(vt.mod1, }\StringTok{"grp.trt.f"}\NormalTok{, }\AttributeTok{console=}\ConstantTok{TRUE}\NormalTok{)}
\end{Highlighting}
\end{Shaded}

\begin{verbatim}
## 
## Study: vt.mod1 ~ "grp.trt.f"
## 
## Duncan's new multiple range test
## for procKnow 
## 
## Mean Square Error:  4.430248 
## 
## grp.trt.f,  means
## 
##   procKnow      std  r        se    Min    Max      Q25     Q50      Q75
## 1 4.930550 1.940008 16 0.5262039 0.3268 6.7509 4.230175 5.67085 6.375725
## 2 7.708337 1.429991 16 0.5262039 5.0733 9.7848 6.682275 8.33320 8.707800
## 3 6.736106 2.820002 16 0.5262039 1.6530 9.9261 5.298875 7.04235 9.101975
## 4 6.874994 1.989995 16 0.5262039 3.2286 9.8754 5.344875 6.84850 8.346350
## 
## Alpha: 0.05 ; DF Error: 60 
## 
## Critical Range
##        2        3        4 
## 1.488551 1.565916 1.616955 
## 
## Means with the same letter are not significantly different.
## 
##   procKnow groups
## 2 7.708337      a
## 4 6.874994      a
## 3 6.736106      a
## 1 4.930550      b
\end{verbatim}

\[\nabla\]

\begin{verbatim}
## Loading required package: carData
\end{verbatim}

\begin{verbatim}
## 
## Attaching package: 'car'
\end{verbatim}

\begin{verbatim}
## The following object is masked from 'package:dplyr':
## 
##     recode
\end{verbatim}

\begin{verbatim}
## The following object is masked from 'package:purrr':
## 
##     some
\end{verbatim}

\begin{verbatim}
## Warning: package 'PMCMRplus' was built under R version 4.5.3
\end{verbatim}

\chapter{Alternative Tests for Treatment
Effects}\label{alternative-tests-for-treatment-effects}

In this chapter, we consider tests for the equality of variances among
treatments/groups, as well as alternatives to the \(F\)-test that are
used in practice.

The tests for equal variances are \textbf{Levene's Test},
\textbf{Hartley's Test}, and \textbf{Bartlett's Test}.

Alternatives to the \(F\)-test are \textbf{Randomization/Permutation
Tests}, the \textbf{Kruskal-Wallis Test}, and \textbf{Welch's Test}.
Also, \textbf{Weighted Least Squares} can be implemented with unequal
variances.

Three departures of the model assumptions are briefly described below.

\begin{itemize}
\tightlist
\item
  \textbf{Non-normal Errors} - This is generally not too problematic for
  the \(F\)-test as long as data are not too far from normal with
  reasonable sample sizes
\item
  \textbf{Unequal Error Variances} - As long as the sample sizes are
  approximately equal this is not a large problem for the \(F\)-test,
  but can be an issue with very unbalanced designs.
\item
  \textbf{Non-independence of Error Terms} This can cause issues for the
  \(F\)-test. If each unit receives each treatment, we can get a
  stronger test based on a Repeated Measures Design.
\end{itemize}

\section{Randomization/Permutation
Tests}\label{randomizationpermutation-tests}

These tests are based on resampling, and make no assumptions on the
distribution of the error terms. The test treats the units observed in
the study as a finite population of units with fixed error terms
\(\epsilon_{ij}\). The model and hypothesis are of the following forms.

\[
Y_{ij}=\mu_{\bullet} + \tau_i + \epsilon_{ij} \qquad \qquad  H_0:\tau_1=\cdots = \tau_r=0 
\] The idea is that if this null hypothesis true, the observed values
are due only to the units themselves and not to treatment/group
differences. A statistic, such as \(F^*\) or \(SSTR\) is computed and
saved. Then, many permutations of the \(n_T\) responses are randomly
assigned to the treatment labels of sizes \(n_1,\ldots,n_r\). The
statistic is computed for each permutation and saved. The \(P\)-value is
the number of more extreme values than the observed statistic plus 1,
divided by the number of permutations plus 1. That is, the observed
value counts in both the numerator and denominator. Often the number of
permutations is 9999 in practice.

This test is applied to the lifeboat training data here.

\textbf{Example 4.1 - Virtual Training for a Lifeboat Launching Task}

Here we will apply the Randomization test to the Lifeboat training data.
For the original sample, \(SSTR=65.664\). Note that for each
permutation, the total sum of squares will be the same,
\(SSTO=331.479\), so we can simply save \(SSTR\) from each permuation,
as opposed to computing \(F^*\). The R code and output are given below.

\begin{Shaded}
\begin{Highlighting}[]
\NormalTok{vt }\OtherTok{\textless{}{-}} \FunctionTok{read.csv}\NormalTok{(}\StringTok{"data/virtual\_training.csv"}\NormalTok{)}
\FunctionTok{head}\NormalTok{(vt,}\DecValTok{2}\NormalTok{)}
\end{Highlighting}
\end{Shaded}

\begin{verbatim}
##   grp.trt procKnow
## 1       1   4.5614
## 2       1   6.6593
\end{verbatim}

\begin{Shaded}
\begin{Highlighting}[]
\FunctionTok{attach}\NormalTok{(vt)}

\DocumentationTok{\#\#\# Randomization Test}
\NormalTok{n.trt }\OtherTok{\textless{}{-}} \FunctionTok{as.vector}\NormalTok{(}\FunctionTok{tapply}\NormalTok{(procKnow, grp.trt, length))  }\DocumentationTok{\#\#\# Obtain n\_1,...,n\_4}
\NormalTok{mean.trt }\OtherTok{\textless{}{-}} \FunctionTok{as.vector}\NormalTok{(}\FunctionTok{tapply}\NormalTok{(procKnow, grp.trt, mean)) }\DocumentationTok{\#\#\# Obtain ybar\_1,...,ybar\_4}
\NormalTok{mean.all }\OtherTok{\textless{}{-}} \FunctionTok{mean}\NormalTok{(procKnow)                             }\DocumentationTok{\#\#\# Obtain ybar}
\NormalTok{r }\OtherTok{\textless{}{-}} \FunctionTok{length}\NormalTok{(n.trt)                                     }\DocumentationTok{\#\#\# r = \#trts}
\NormalTok{n\_T }\OtherTok{\textless{}{-}} \FunctionTok{sum}\NormalTok{(n.trt)                                      }\DocumentationTok{\#\#\# n\_T=Total sample size}
\NormalTok{(SSTR.obs }\OtherTok{\textless{}{-}} \FunctionTok{sum}\NormalTok{(n.trt }\SpecialCharTok{*}\NormalTok{ (mean.trt }\SpecialCharTok{{-}}\NormalTok{ mean.all)}\SpecialCharTok{\^{}}\DecValTok{2}\NormalTok{))     }\DocumentationTok{\#\#\# Compute the observed SSTR}
\end{Highlighting}
\end{Shaded}

\begin{verbatim}
## [1] 65.66393
\end{verbatim}

\begin{Shaded}
\begin{Highlighting}[]
\NormalTok{N.perm }\OtherTok{\textless{}{-}} \DecValTok{9999}                                         \DocumentationTok{\#\#\# \# of permutation samples}
\FunctionTok{set.seed}\NormalTok{(}\DecValTok{1234}\NormalTok{)                                         }\DocumentationTok{\#\#\# Set seed for reproducibility}
\NormalTok{SSTR.perm }\OtherTok{\textless{}{-}} \FunctionTok{rep}\NormalTok{(}\DecValTok{0}\NormalTok{, N.perm)                            }\DocumentationTok{\#\#\# Placeholder for SSTR\textquotesingle{}s}
\ControlFlowTok{for}\NormalTok{ (i }\ControlFlowTok{in} \DecValTok{1}\SpecialCharTok{:}\NormalTok{N.perm) \{}
\NormalTok{  perm }\OtherTok{\textless{}{-}} \FunctionTok{sample}\NormalTok{(}\DecValTok{1}\SpecialCharTok{:}\NormalTok{n\_T, n\_T, }\AttributeTok{replace=}\ConstantTok{FALSE}\NormalTok{)            }\DocumentationTok{\#\# Random permutation of 1:n\_T}
\NormalTok{  Y.perm }\OtherTok{\textless{}{-}}\NormalTok{ procKnow[perm]                             }\DocumentationTok{\#\# "Sorts" Y values by perm}
\NormalTok{  mean.trt.perm }\OtherTok{\textless{}{-}} \FunctionTok{as.vector}\NormalTok{(}\FunctionTok{tapply}\NormalTok{(Y.perm, grp.trt, mean))  }\DocumentationTok{\#\# Obtain Perm trt means }
\NormalTok{  SSTR.perm[i] }\OtherTok{\textless{}{-}} \FunctionTok{sum}\NormalTok{(n.trt }\SpecialCharTok{*}\NormalTok{ (mean.trt.perm }\SpecialCharTok{{-}}\NormalTok{ mean.all)}\SpecialCharTok{\^{}}\DecValTok{2}\NormalTok{)  }\DocumentationTok{\#\# SSTR for perm sample}
\NormalTok{\}}

\NormalTok{(perm.pvalue }\OtherTok{\textless{}{-}}\NormalTok{ (}\FunctionTok{sum}\NormalTok{(SSTR.perm }\SpecialCharTok{\textgreater{}=}\NormalTok{ SSTR.obs) }\SpecialCharTok{+} \DecValTok{1}\NormalTok{) }\SpecialCharTok{/}\NormalTok{ (N.perm }\SpecialCharTok{+} \DecValTok{1}\NormalTok{))  }\DocumentationTok{\#\# Permutation P{-}value}
\end{Highlighting}
\end{Shaded}

\begin{verbatim}
## [1] 0.005
\end{verbatim}

\begin{Shaded}
\begin{Highlighting}[]
\DocumentationTok{\#\# Histogram with vertical line at observed SSTR}
\FunctionTok{hist}\NormalTok{(SSTR.perm, }\AttributeTok{breaks=}\StringTok{"Scott"}\NormalTok{, }\AttributeTok{col=}\StringTok{"orange"}\NormalTok{, }\AttributeTok{border=}\StringTok{"blue4"}\NormalTok{)}
\FunctionTok{abline}\NormalTok{(}\AttributeTok{v =}\NormalTok{ SSTR.obs, }\AttributeTok{col=}\StringTok{"green4"}\NormalTok{, }\AttributeTok{lwd=}\DecValTok{2}\NormalTok{)}
\end{Highlighting}
\end{Shaded}

\pandocbounded{\includegraphics[keepaspectratio]{bookdownproj_STA4211Notes1_files/figure-latex/ch0401-1.pdf}}

\begin{Shaded}
\begin{Highlighting}[]
\FunctionTok{detach}\NormalTok{(vt)}
\end{Highlighting}
\end{Shaded}

The \(P\)-value is 0.005, which is well below 0.05. Only 49 permutations
out of 9999 had larger values of \(SSTR\) than the observed data. There
is strong evidence of treatment effects. \[\nabla \]

\section{Tests for Equal Error
Variances}\label{tests-for-equal-error-variances}

In this section, we consider three widely used tests of equal variances
among the treatments/groups. They are based on various assumptions and
do not always provide the same conclusion. The three tests covered here
are \textbf{Hartley's Test}, \textbf{Levene's Test}, and
\textbf{Bartlett's Test}. All tests are testing the following
hypotheses.

\[
H_0:\sigma_1^2=\cdots = \sigma_r^2 \qquad
H_A: \mbox{ Not all } \sigma_i^2 \mbox{ are equal}
\]

\subsection{Hartley's Test}\label{hartleys-test}

This test is dependent on the errors being normally distributed.
Further, its derivation is based on equal sample sizes. It uses a
special table of critical values available in many textbooks and online.
There is a R package that purports to generate the critical values, but
I have had issues with certain sample sizes/numbers of treatments in
obtaining results.

The test statistic \((H)\) is simply the ratio of the largest to
smallest variance (not standard deviation), and the critical values are
indexed by the number of treatments/groups \((r)\) and the number of
degrees of freedom for each sample variance \((n-1)\), where \(n\) is
the number of units per treatment.

\textbf{Example 4.2 - Virtual Training for a Lifeboat Launching Task}

For the lifeboat training experiment, there were \(r=4\) treatments and
each sample size was \(n=16\). The sample standard deviations are given
below. The critical value for this test with \(\alpha=0.05\), \(r=4\),
and \(n-1=15\) is \(H_{.95;4,15}=4.01\).

\[
s_1=1.94 \quad s_2=1.43 \quad s_3=2.82 \quad s_4=1.99
\qquad H^*=\frac{\max\left(s_i^2\right)}{\min\left(s_i^2\right)}=
\frac{(2.82)^2}{(1.43)^2}=3.89
\]

The test statistic is just below the critical value, so we fail to
reject the hypothesis of equal variances.

\[ \nabla \]

\subsection{Levene's Test}\label{levenes-test}

This test, also known as the Brown-Forsythe test is robust to outliers
and does not require equal sample sizes. There are various versions,
based on whether deviations from group means or medians are used. For
each observation, we compute \(d_{ij}\) as described here.

\[
d_{ij}= \left|Y_{ij}-\tilde{Y}_i\right| \qquad  
\tilde{Y}_i=\mbox{median}\left(Y_{i1},\ldots ,Y_{in_i}\right)
\]

Then an Analysis of Variance is applied to the \(d_{ij}\), with
\(\overline{d}_{i\bullet}\) and \(\overline{d}_{\bullet\bullet}\) being
the treatment and overall means.

\[
SSTR_L = \sum_{i=1}^rn_i\left(\overline{d}_{i\bullet}-
\overline{d}_{\bullet\bullet}\right)^2  \qquad df_{TR}=r-1 
\]

\[
SSE_L = \sum_{i=1}^r\sum_{j=1}^{n_i}\left(d_{ij}-\overline{d}_{i\bullet}\right)^2
\qquad df_E = n_T-r
\]

Then, just as when we tested for equality of means, we compute the
\(F\)-statistic and compare it with \(F_{1-\alpha;r-1,n_T-r}\).

\[
\mbox{Test Statistic:} \quad F_L^* = \frac{\left[\frac{SSTR_L}{r-1}\right]}
           {\left[\frac{SSE_L}{n_T-r}\right]} \qquad
\mbox{Rejection Region:} \quad F_L^* \ge F_{1-\alpha;r-1,n_T-r}
\]

\textbf{Example 4.3 - Virtual Training for a Lifeboat Launching Task }

We will first create the \(d_{ij}\) ``brute-force'', then use the
\textbf{leveneTest} function in the \textbf{car} package.

\begin{Shaded}
\begin{Highlighting}[]
\NormalTok{vt }\OtherTok{\textless{}{-}} \FunctionTok{read.csv}\NormalTok{(}\StringTok{"data/virtual\_training.csv"}\NormalTok{)}
\NormalTok{n.trt }\OtherTok{\textless{}{-}} \FunctionTok{as.vector}\NormalTok{(}\FunctionTok{tapply}\NormalTok{(vt}\SpecialCharTok{$}\NormalTok{procKnow, vt}\SpecialCharTok{$}\NormalTok{grp.trt, length))      }\DocumentationTok{\#\# n\_1,...,n\_4}
\NormalTok{med.trt }\OtherTok{\textless{}{-}} \FunctionTok{as.vector}\NormalTok{(}\FunctionTok{tapply}\NormalTok{(vt}\SpecialCharTok{$}\NormalTok{procKnow, vt}\SpecialCharTok{$}\NormalTok{grp.trt, median))    }\DocumentationTok{\#\# med\_1,...,med\_4}
\NormalTok{vt}\SpecialCharTok{$}\NormalTok{med.trt }\OtherTok{\textless{}{-}} \FunctionTok{rep}\NormalTok{(med.trt, n.trt)                    }\DocumentationTok{\#\# Repeat med\_1 n\_1 times, ...}
\FunctionTok{head}\NormalTok{(vt, }\DecValTok{3}\NormalTok{); }\FunctionTok{tail}\NormalTok{(vt, }\DecValTok{3}\NormalTok{)                  }
\end{Highlighting}
\end{Shaded}

\begin{verbatim}
##   grp.trt procKnow med.trt
## 1       1   4.5614 5.67085
## 2       1   6.6593 5.67085
## 3       1   5.6427 5.67085
\end{verbatim}

\begin{verbatim}
##    grp.trt procKnow med.trt
## 62       4   3.2286  6.8485
## 63       4   6.5520  6.8485
## 64       4   9.8754  6.8485
\end{verbatim}

\begin{Shaded}
\begin{Highlighting}[]
\NormalTok{vt}\SpecialCharTok{$}\NormalTok{d }\OtherTok{\textless{}{-}} \FunctionTok{abs}\NormalTok{(vt}\SpecialCharTok{$}\NormalTok{procKnow}\SpecialCharTok{{-}}\NormalTok{vt}\SpecialCharTok{$}\NormalTok{med.trt)                  }\DocumentationTok{\#\# Compute d\_ij}

\NormalTok{lev.mod }\OtherTok{\textless{}{-}} \FunctionTok{lm}\NormalTok{(d }\SpecialCharTok{\textasciitilde{}} \FunctionTok{factor}\NormalTok{(grp.trt), }\AttributeTok{data=}\NormalTok{vt)          }\DocumentationTok{\#\# Run aov function on d\_ij\textquotesingle{}s}
\FunctionTok{anova}\NormalTok{(lev.mod)                                       }\DocumentationTok{\#\# Levene\textquotesingle{}s F{-}test output}
\end{Highlighting}
\end{Shaded}

\begin{verbatim}
## Analysis of Variance Table
## 
## Response: d
##                 Df  Sum Sq Mean Sq F value Pr(>F)
## factor(grp.trt)  3  11.108  3.7026  2.1463 0.1038
## Residuals       60 103.506  1.7251
\end{verbatim}

\begin{Shaded}
\begin{Highlighting}[]
\DocumentationTok{\#\# leveneTest (using Y\_ij\textquotesingle{}s) from car package}
\FunctionTok{leveneTest}\NormalTok{(procKnow }\SpecialCharTok{\textasciitilde{}} \FunctionTok{factor}\NormalTok{(grp.trt), }\StringTok{"median"}\NormalTok{, }\AttributeTok{data=}\NormalTok{vt)  }
\end{Highlighting}
\end{Shaded}

\begin{verbatim}
## Levene's Test for Homogeneity of Variance (center = "median")
##       Df F value Pr(>F)
## group  3  2.1463 0.1038
##       60
\end{verbatim}

We fail to reject the null hypothesis of equal variances. Note that we
ran the \textbf{leveneTest} function on the \(Y_{ij}\) values, not on
the with the \(d_{ij}\) values (it creates them internally).
\[ \nabla \]

\subsection{Bartlett's Test}\label{bartletts-test}

Bartlett's test is based on the residuals being normally distributed and
can be used in applications beyond ANOVA models, including regression
models with categorical predictors. It makes use of the individual
variances \(\left\{s_i^2\right\}\) and the pooled \(MSE\) to obtain a
chi-square statistic that is defined here.

\[
\mbox{Test Statistic:} \quad X_B^2 =
\frac{1}{C}\left[\left(n_T-r\right)\ln(MSE)-\sum_{i=1}^r\left(n_i-1\right)\ln\left(s_i^2\right)\right]
\] \[
\mbox{where:}\quad 
C=1+\frac{1}{3(r-1)}\left[\left(\sum_{i=1}^r\frac{1}{n_i-1}\right)-\left(\frac{1}{n_T-r}\right)\right]
\]

\[
\mbox{Rejection Region:} \quad 
X_B^2 \ge \chi_{1-\alpha;r-1}^2 \qquad
P = P\left(\chi_{r-1}^2 \ge X_B^2\right)
\]

\textbf{Example 4.4 - Virtual Training for a Lifeboat Launching Task}

For the lifeboat training experiment, there were \(r=4\) treatments and
each sample size was \(n=16\). The sample variances are given below. The
critical value for this test with \(\alpha=0.05\) and \(r=4\) is
\(\chi_{.95;4-1}=7.815\).

\[
s_1^2=3.764 \quad s_2^2=2.045 \quad s_3^2=7.952 \quad s_4^2=3.960
\qquad MSE=4.430
\]

\[
\left(n_T-r\right)\ln(MSE) = 89.304 \qquad
\sum_{i=1}^r\left(n_i-1\right)\ln\left(s_i^2\right)=82.358
\]

\[
C=1+\frac{1}{3(r-1)}\left[\left(\sum_{i=1}^r\frac{1}{n_i-1}\right)-\left(\frac{1}{n_T-r}\right)\right]=1+\frac{1}{9}\left[\frac{4}{16-1}-\frac{1}{64-4}\right]=1.028
\]

\[
\mbox{Test Statistic:} \quad X_B^2 = \frac{1}{1.028}(89.304-82.358)=6.757 \qquad
P=P\left(\chi_3^2 \ge 6.757\right)=.0801
\]

Finally, we use the \textbf{bartlett.test} function in R to conduct the
test.

\begin{Shaded}
\begin{Highlighting}[]
\NormalTok{vt }\OtherTok{\textless{}{-}} \FunctionTok{read.csv}\NormalTok{(}\StringTok{"data/virtual\_training.csv"}\NormalTok{)}

\FunctionTok{bartlett.test}\NormalTok{(procKnow }\SpecialCharTok{\textasciitilde{}} \FunctionTok{factor}\NormalTok{(grp.trt), }\AttributeTok{data=}\NormalTok{vt)}
\end{Highlighting}
\end{Shaded}

\begin{verbatim}
## 
##  Bartlett test of homogeneity of variances
## 
## data:  procKnow by factor(grp.trt)
## Bartlett's K-squared = 6.7626, df = 3, p-value = 0.07986
\end{verbatim}

\[
\nabla
\]

\section{Remedial Measures}\label{remedial-measures}

In this section, several procedures are considered for issues of
normality and/or unequal variances among the error terms. Methods will
include \textbf{Weighted Least Squares} and \textbf{Welch's Test} when
the errors are normal, but the error variances are unequal
(heteroscedastic).

When variances are non-normal and heteroscedastic, transformations can
be applied to \(Y\) that may solve both problems. These include the
\textbf{Box-Cox Transformation}, and were covered in the Linear
Regression course.

When the distributions are highly skewed, non-parametric tests can be
conducted based on ranks. We will describe the \textbf{Kruskal-Wallis
Test}.

\subsection{Weighted Least Squares and Welch's
Test}\label{weighted-least-squares-and-welchs-test}

Consider the following scenario and simulation.

\begin{itemize}
\tightlist
\item
  There are \(r=3\) Treatments with equal means
  \(\left(\mu_1=\mu_2=\mu_3\right)\)
\item
  The Treatments have unequal standard deviations (and thus variances)
\item
  The Treatments have equal or unequal sample sizes
\end{itemize}

Note that the null hypothesis \(H_0:\mu_1=\mu_2=\mu_3\) holds, so
optimally, the rejection rate based on the \(F\)-test with significance
level \(\alpha\) should be approximately \(\alpha\). We will simulate
datasets for the following scenarios, all with \(\mu=100\) for all
treatments, and normal error terms.

\begin{itemize}
\tightlist
\item
  \textbf{Scenario 1} -
  \(n_1=n_2=n_3=20\quad \sigma_1=\sigma_2=\sigma_3=20\) (standard
  assumptions)
\item
  \textbf{Scenario 2} -
  \(n_1=n_2=n_3=20 \quad \sigma_1=10,\sigma_2=20,\sigma_3=30\)
\item
  \textbf{Scenario 3} -
  \(n_1=10,n_2=20,n_3=30 \quad \sigma_1=10,\sigma_2=20,\sigma_3=30\)
\item
  \textbf{Scenario 4} -
  \(n_1=30,n_2=20,n_3=10 \quad \sigma_1=10,\sigma_2=20,\sigma_3=30\)
\end{itemize}

Each scenario will be run, and the rejection rate will be reported in
each case. The program below is the same for each scenario with only the
inputs for \textbf{n} and \textbf{sigma} changing.

\begin{Shaded}
\begin{Highlighting}[]
\DocumentationTok{\#\# Scenario 3}
\NormalTok{mu }\OtherTok{\textless{}{-}} \FunctionTok{rep}\NormalTok{(}\DecValTok{100}\NormalTok{,}\DecValTok{3}\NormalTok{)}
\NormalTok{sigma }\OtherTok{\textless{}{-}} \FunctionTok{c}\NormalTok{(}\DecValTok{10}\NormalTok{,}\DecValTok{20}\NormalTok{,}\DecValTok{30}\NormalTok{)   }\DocumentationTok{\#\# Change these for various scenarios}
\NormalTok{n }\OtherTok{\textless{}{-}} \FunctionTok{c}\NormalTok{(}\DecValTok{10}\NormalTok{,}\DecValTok{20}\NormalTok{,}\DecValTok{30}\NormalTok{)       }\DocumentationTok{\#\# Change these for various scenarios}
\NormalTok{y.grid }\OtherTok{\textless{}{-}} \FunctionTok{seq}\NormalTok{(}\DecValTok{0}\NormalTok{,}\DecValTok{200}\NormalTok{,.}\DecValTok{025}\NormalTok{)}

\FunctionTok{plot}\NormalTok{(y.grid, }\FunctionTok{dnorm}\NormalTok{(y.grid,mu[}\DecValTok{1}\NormalTok{],sigma[}\DecValTok{1}\NormalTok{]), }\AttributeTok{type=}\StringTok{"l"}\NormalTok{, }\AttributeTok{col=}\StringTok{"red4"}\NormalTok{,}\AttributeTok{lwd=}\FloatTok{1.5}\NormalTok{,}
     \AttributeTok{xlab=}\StringTok{"y"}\NormalTok{, }\AttributeTok{ylab=}\StringTok{"density"}\NormalTok{)}
\FunctionTok{lines}\NormalTok{(y.grid, }\FunctionTok{dnorm}\NormalTok{(y.grid,mu[}\DecValTok{2}\NormalTok{],sigma[}\DecValTok{2}\NormalTok{]), }\AttributeTok{col=}\StringTok{"blue4"}\NormalTok{,}\AttributeTok{lwd=}\FloatTok{1.5}\NormalTok{)}
\FunctionTok{lines}\NormalTok{(y.grid, }\FunctionTok{dnorm}\NormalTok{(y.grid,mu[}\DecValTok{3}\NormalTok{],sigma[}\DecValTok{3}\NormalTok{]), }\AttributeTok{col=}\StringTok{"green4"}\NormalTok{,}\AttributeTok{lwd=}\FloatTok{1.5}\NormalTok{)}
\FunctionTok{legend}\NormalTok{(}\StringTok{"topright"}\NormalTok{, }\FunctionTok{c}\NormalTok{(}\StringTok{"n=10,sigma=10"}\NormalTok{,}\StringTok{"n=20,sigma=20"}\NormalTok{,}\StringTok{"n=30,sigma=30"}\NormalTok{),}
     \AttributeTok{col=}\FunctionTok{c}\NormalTok{(}\StringTok{"red4"}\NormalTok{, }\StringTok{"blue4"}\NormalTok{, }\StringTok{"green4"}\NormalTok{), }\AttributeTok{lty=}\DecValTok{1}\NormalTok{, }\AttributeTok{lwd=}\FloatTok{1.5}\NormalTok{)}
\end{Highlighting}
\end{Shaded}

\pandocbounded{\includegraphics[keepaspectratio]{bookdownproj_STA4211Notes1_files/figure-latex/ch0404-1.pdf}}

\begin{Shaded}
\begin{Highlighting}[]
\FunctionTok{set.seed}\NormalTok{(}\DecValTok{32611}\NormalTok{)}
\NormalTok{N.samp }\OtherTok{\textless{}{-}} \DecValTok{10000}
\NormalTok{n\_T }\OtherTok{\textless{}{-}} \FunctionTok{sum}\NormalTok{(n)}
\NormalTok{r }\OtherTok{\textless{}{-}} \FunctionTok{length}\NormalTok{(n)}
\NormalTok{trt.grp }\OtherTok{\textless{}{-}} \FunctionTok{rep}\NormalTok{(}\DecValTok{1}\SpecialCharTok{:}\DecValTok{3}\NormalTok{,}\FunctionTok{c}\NormalTok{(n[}\DecValTok{1}\NormalTok{],n[}\DecValTok{2}\NormalTok{],n[}\DecValTok{3}\NormalTok{]))}
\NormalTok{F.stat }\OtherTok{\textless{}{-}} \FunctionTok{rep}\NormalTok{(}\DecValTok{0}\NormalTok{,N.samp)}

\ControlFlowTok{for}\NormalTok{ (i1 }\ControlFlowTok{in} \DecValTok{1}\SpecialCharTok{:}\NormalTok{N.samp) \{  }
\NormalTok{  y1 }\OtherTok{\textless{}{-}} \FunctionTok{rnorm}\NormalTok{(n[}\DecValTok{1}\NormalTok{],mu[}\DecValTok{1}\NormalTok{],sigma[}\DecValTok{1}\NormalTok{])}
\NormalTok{  y2 }\OtherTok{\textless{}{-}} \FunctionTok{rnorm}\NormalTok{(n[}\DecValTok{2}\NormalTok{],mu[}\DecValTok{2}\NormalTok{],sigma[}\DecValTok{2}\NormalTok{])}
\NormalTok{  y3 }\OtherTok{\textless{}{-}} \FunctionTok{rnorm}\NormalTok{(n[}\DecValTok{3}\NormalTok{],mu[}\DecValTok{3}\NormalTok{],sigma[}\DecValTok{3}\NormalTok{])}
\NormalTok{  y123 }\OtherTok{\textless{}{-}} \FunctionTok{c}\NormalTok{(y1,y2,y3)}
\NormalTok{  ybar1 }\OtherTok{\textless{}{-}} \FunctionTok{mean}\NormalTok{(y1); var1 }\OtherTok{\textless{}{-}} \FunctionTok{var}\NormalTok{(y1)}
\NormalTok{  ybar2 }\OtherTok{\textless{}{-}} \FunctionTok{mean}\NormalTok{(y2); var2 }\OtherTok{\textless{}{-}} \FunctionTok{var}\NormalTok{(y2)}
\NormalTok{  ybar3 }\OtherTok{\textless{}{-}} \FunctionTok{mean}\NormalTok{(y3); var3 }\OtherTok{\textless{}{-}} \FunctionTok{var}\NormalTok{(y3)}
\NormalTok{  ybar123 }\OtherTok{\textless{}{-}} \FunctionTok{mean}\NormalTok{(y123)}
\NormalTok{  SSTR }\OtherTok{\textless{}{-}}\NormalTok{ n[}\DecValTok{1}\NormalTok{]}\SpecialCharTok{*}\NormalTok{(ybar1}\SpecialCharTok{{-}}\NormalTok{ybar123)}\SpecialCharTok{\^{}}\DecValTok{2} \SpecialCharTok{+}\NormalTok{ n[}\DecValTok{2}\NormalTok{]}\SpecialCharTok{*}\NormalTok{(ybar2}\SpecialCharTok{{-}}\NormalTok{ybar123)}\SpecialCharTok{\^{}}\DecValTok{2} \SpecialCharTok{+} 
\NormalTok{          n[}\DecValTok{3}\NormalTok{]}\SpecialCharTok{*}\NormalTok{(ybar3}\SpecialCharTok{{-}}\NormalTok{ybar123)}\SpecialCharTok{\^{}}\DecValTok{2}
\NormalTok{  SSE }\OtherTok{\textless{}{-}}\NormalTok{ (n[}\DecValTok{1}\NormalTok{]}\SpecialCharTok{{-}}\DecValTok{1}\NormalTok{)}\SpecialCharTok{*}\NormalTok{var1 }\SpecialCharTok{+}\NormalTok{ (n[}\DecValTok{2}\NormalTok{]}\SpecialCharTok{{-}}\DecValTok{1}\NormalTok{)}\SpecialCharTok{*}\NormalTok{var2 }\SpecialCharTok{+}\NormalTok{ (n[}\DecValTok{3}\NormalTok{]}\SpecialCharTok{{-}}\DecValTok{1}\NormalTok{)}\SpecialCharTok{*}\NormalTok{var3}
\NormalTok{  F.stat[i1] }\OtherTok{\textless{}{-}}\NormalTok{ ( (SSTR}\SpecialCharTok{/}\NormalTok{(r}\DecValTok{{-}1}\NormalTok{)) }\SpecialCharTok{/}\NormalTok{ (SSE}\SpecialCharTok{/}\NormalTok{(n\_T}\SpecialCharTok{{-}}\NormalTok{r)) )}
\NormalTok{\}}

\FunctionTok{summary}\NormalTok{(F.stat)}
\end{Highlighting}
\end{Shaded}

\begin{verbatim}
##      Min.   1st Qu.    Median      Mean   3rd Qu.      Max. 
## 4.730e-05 1.601e-01 4.145e-01 6.918e-01 9.044e-01 1.322e+01
\end{verbatim}

\begin{Shaded}
\begin{Highlighting}[]
\FunctionTok{mean}\NormalTok{(F.stat }\SpecialCharTok{\textgreater{}=} \FunctionTok{qf}\NormalTok{(.}\DecValTok{95}\NormalTok{,r}\DecValTok{{-}1}\NormalTok{,n\_T}\SpecialCharTok{{-}}\NormalTok{r))}
\end{Highlighting}
\end{Shaded}

\begin{verbatim}
## [1] 0.022
\end{verbatim}

Each scenario was run on 10000 simulated data sets with \(\alpha=0.05\).
The following results were obtained based on the seed 32611.

\begin{itemize}
\tightlist
\item
  \textbf{Scenario 1} - Equal variances, Equal sample sizes. Rejection
  Rate = .0513
\item
  \textbf{Scenario 2} - Unequal variances, Equal sample sizes. Rejection
  Rate = .0573
\item
  \textbf{Scenario 3} - Increasing variances, Increasing sample sizes.
  Rejection Rate = .0220.
\item
  \textbf{Scenario 4} - Increasing variances, Decreasing sample sizes.
  Rejection Rate = .1637
\end{itemize}

The primary result is that as long as sample sizes are equal (or very
similar), the \(F\)-test performs well whether the variances are equal
or not (both Scenarios 1 and 2 have rejection rates close to
\(\alpha\)).

When the sample sizes are positively correlated with the variance
(Scenario 3), the rejection rate was well below \(\alpha\)
\((.0220<.05)\). Thus, when the larger (smaller) samples correspond to
the larger (smaller) variances the \(F\)-test rejects less often than it
should.

When the sample sizes are negatively correlated with the variance
(Scenario 4), the rejection rate was well above \(\alpha\)
\((.1637>.05)\). Thus, when the smaller (larger) samples correspond to
the larger (smaller) variances the \(F\)-test rejects more often than it
should.

\textbf{Estimated Weighted Least Squares} can be used when the error
variances are unequal. The estimated weights for the individual
measurements is the reciprocal of the variance of the measurements
within that treatment/group, \(w_{ij}=1/s_i^2\). It is much easier doing
direct computations in matrix form than in scalar form. A sketch of the
method is given below for the cell means model with balanced data:
\(n_1=\cdots=n_r=n\), although it can be easily generalized for unequal
sample sizes. Let \(I_n\) be the \(n\times n\) identity matrix and
\(0_n\) be the \(n\times n\) matrix of \(0^s\). The \(n_T\times r\)
matrix \(X\) and the \(n_T\times 1\) vector \(Y\) are described in
Chapter 2.

\[
\hat{V}=
\begin{bmatrix}
s_1^2I_n & 0_n & \cdots & 0_n \\
0_n & s_2^2I_n & \cdots & 0_n \\
\vdots & \vdots & \ddots & \vdots \\
0_n & 0_n & \cdots & s_r^2I_n \\
\end{bmatrix}
\qquad
X^*=\hat{V}^{-1/2}X
\qquad
Y^*=\hat{V}^{-1/2}Y
\]

\[
H^*=
X^*\left(X^{*'}X^*\right)
^{-1}X^{*'} 
\qquad
SSE_W(F)=Y^{*'}\left(I_{n_T}-H^*\right)Y^* \quad df_E(F)=n_T-r
\] Under the null hypothesis \(H_0:\mu_1=\cdots\mu_r=\mu\), the \(X_0\)
matrix is the \(n_T\times 1\) vector of \(1^s\). In this (null) case we
obtain the reduced \(SSE\) below.

\[
X_0^*=\hat{V}^{-1/2}X_0 \quad 
H_0^*=
X_0^*\left(X_0^{*'}X_0^*\right)
^{-1}X_0^{*'} 
\qquad
SSE_W(R)=Y^{*'}\left(I_{n_T}-H_0^*\right)Y^* \quad df_E(R)=n_T-1
\] The test statistic and rejection region for the general linear test
are given below.

\[
\mbox{Test Statistic:} \quad F_{EWLS}^*=
\frac{ \left[\frac{\left(SSE_W(R)-SSE_W(F)\right)}
      {\left(n_T-1\right)-\left(n_T-r\right)}\right]}
      {\left[\frac{SSE_W(F)}{n_T-r}\right]}
\qquad
\mbox{Rejection Region:} F_{EWLS}^* \ge F_{1-\alpha,r-1,n_T-r}
\]

\textbf{Example 4.5 - Virtual Training for a Lifeboat Launching Task}

For the lifeboat training experiment, there were \(r=4\) treatments and
each sample size was \(n=16\). The following R code runs the matrix form
of the general linear test and obtains the sample standard deviation by
treatment and creates \(\hat{V}^{-1/2}\) matrix and runs the generalized
linear test directly. Then the model is fit using the \textbf{weight}
command in the \textbf{aov} function.

\begin{Shaded}
\begin{Highlighting}[]
\NormalTok{vt }\OtherTok{\textless{}{-}} \FunctionTok{read.csv}\NormalTok{(}\StringTok{"data/virtual\_training.csv"}\NormalTok{)}
\DocumentationTok{\#\#\# Matrix form}
\NormalTok{Vhat\_12 }\OtherTok{\textless{}{-}} \FunctionTok{matrix}\NormalTok{(}\FunctionTok{rep}\NormalTok{(}\DecValTok{0}\NormalTok{,}\DecValTok{64}\SpecialCharTok{\^{}}\DecValTok{2}\NormalTok{),}\AttributeTok{ncol=}\DecValTok{64}\NormalTok{)        }\DocumentationTok{\#\# Construct estimated V\^{}\{{-}1/2\} matrix}
\NormalTok{trt.sd }\OtherTok{\textless{}{-}} \FunctionTok{as.vector}\NormalTok{(}\FunctionTok{tapply}\NormalTok{(vt}\SpecialCharTok{$}\NormalTok{procKnow, vt}\SpecialCharTok{$}\NormalTok{grp.trt, sd))    }\DocumentationTok{\#\# Obtain s\_1,...,s\_4}
\DocumentationTok{\#\# V\^{}\{{-}1/2\} matrix = diag\{1/s\_i\}}
\NormalTok{Vhat\_12[}\DecValTok{1}\SpecialCharTok{:}\DecValTok{16}\NormalTok{,}\DecValTok{1}\SpecialCharTok{:}\DecValTok{16}\NormalTok{] }\OtherTok{\textless{}{-}}\NormalTok{ (}\DecValTok{1}\SpecialCharTok{/}\NormalTok{trt.sd[}\DecValTok{1}\NormalTok{])}\SpecialCharTok{*}\FunctionTok{diag}\NormalTok{(}\DecValTok{16}\NormalTok{)             }
\NormalTok{Vhat\_12[}\DecValTok{17}\SpecialCharTok{:}\DecValTok{32}\NormalTok{,}\DecValTok{17}\SpecialCharTok{:}\DecValTok{32}\NormalTok{] }\OtherTok{\textless{}{-}}\NormalTok{ (}\DecValTok{1}\SpecialCharTok{/}\NormalTok{trt.sd[}\DecValTok{2}\NormalTok{])}\SpecialCharTok{*}\FunctionTok{diag}\NormalTok{(}\DecValTok{16}\NormalTok{)}
\NormalTok{Vhat\_12[}\DecValTok{33}\SpecialCharTok{:}\DecValTok{48}\NormalTok{,}\DecValTok{33}\SpecialCharTok{:}\DecValTok{48}\NormalTok{] }\OtherTok{\textless{}{-}}\NormalTok{ (}\DecValTok{1}\SpecialCharTok{/}\NormalTok{trt.sd[}\DecValTok{3}\NormalTok{])}\SpecialCharTok{*}\FunctionTok{diag}\NormalTok{(}\DecValTok{16}\NormalTok{)}
\NormalTok{Vhat\_12[}\DecValTok{49}\SpecialCharTok{:}\DecValTok{64}\NormalTok{,}\DecValTok{49}\SpecialCharTok{:}\DecValTok{64}\NormalTok{] }\OtherTok{\textless{}{-}}\NormalTok{ (}\DecValTok{1}\SpecialCharTok{/}\NormalTok{trt.sd[}\DecValTok{4}\NormalTok{])}\SpecialCharTok{*}\FunctionTok{diag}\NormalTok{(}\DecValTok{16}\NormalTok{)}

\NormalTok{one }\OtherTok{\textless{}{-}} \FunctionTok{matrix}\NormalTok{(}\FunctionTok{rep}\NormalTok{(}\DecValTok{1}\NormalTok{,}\DecValTok{16}\NormalTok{),}\AttributeTok{ncol=}\DecValTok{1}\NormalTok{)         }\DocumentationTok{\#\# 16x1 vector of 1\textquotesingle{}s}
\NormalTok{zero }\OtherTok{\textless{}{-}} \FunctionTok{matrix}\NormalTok{(}\FunctionTok{rep}\NormalTok{(}\DecValTok{0}\NormalTok{,}\DecValTok{16}\NormalTok{),}\AttributeTok{ncol=}\DecValTok{1}\NormalTok{)        }\DocumentationTok{\#\# 16x1 vector of 0\textquotesingle{}s}
\DocumentationTok{\#\# Construct X matrix for cell means model}
\NormalTok{X }\OtherTok{\textless{}{-}} \FunctionTok{cbind}\NormalTok{(}\FunctionTok{rbind}\NormalTok{(one,zero,zero,zero),}
           \FunctionTok{rbind}\NormalTok{(zero,one,zero,zero),}
           \FunctionTok{rbind}\NormalTok{(zero,zero,one,zero),}
           \FunctionTok{rbind}\NormalTok{(zero,zero,zero,one))}
\NormalTok{Y }\OtherTok{\textless{}{-}}\NormalTok{ vt}\SpecialCharTok{$}\NormalTok{procKnow}
\NormalTok{Xstar }\OtherTok{\textless{}{-}}\NormalTok{ Vhat\_12 }\SpecialCharTok{\%*\%}\NormalTok{ X                                     }\DocumentationTok{\#\# X*}
\NormalTok{Ystar }\OtherTok{\textless{}{-}}\NormalTok{ Vhat\_12 }\SpecialCharTok{\%*\%}\NormalTok{ Y                                     }\DocumentationTok{\#\# Y*}
\NormalTok{Hstar }\OtherTok{\textless{}{-}}\NormalTok{ Xstar }\SpecialCharTok{\%*\%} \FunctionTok{solve}\NormalTok{(}\FunctionTok{t}\NormalTok{(Xstar)}\SpecialCharTok{\%*\%}\NormalTok{Xstar) }\SpecialCharTok{\%*\%} \FunctionTok{t}\NormalTok{(Xstar)    }\DocumentationTok{\#\# H*}
\NormalTok{I\_64 }\OtherTok{\textless{}{-}} \FunctionTok{diag}\NormalTok{(}\DecValTok{64}\NormalTok{)                                           }\DocumentationTok{\#\# I\_64}
\NormalTok{SSE\_W.F }\OtherTok{\textless{}{-}} \FunctionTok{t}\NormalTok{(Ystar) }\SpecialCharTok{\%*\%}\NormalTok{ (I\_64}\SpecialCharTok{{-}}\NormalTok{Hstar) }\SpecialCharTok{\%*\%}\NormalTok{ Ystar             }\DocumentationTok{\#\# SSE\_W(F)}
\NormalTok{dfE.F }\OtherTok{\textless{}{-}} \DecValTok{64{-}4}                                              \DocumentationTok{\#\# df\_E(F)}

\NormalTok{X\_0 }\OtherTok{\textless{}{-}} \FunctionTok{rbind}\NormalTok{(one,one,one,one)                                      }\DocumentationTok{\#\# X\_0}
\NormalTok{Xstar\_0 }\OtherTok{\textless{}{-}}\NormalTok{ Vhat\_12 }\SpecialCharTok{\%*\%}\NormalTok{ X\_0                                         }\DocumentationTok{\#\# X\_0*}
\NormalTok{Hstar\_0 }\OtherTok{\textless{}{-}}\NormalTok{ Xstar\_0 }\SpecialCharTok{\%*\%} \FunctionTok{solve}\NormalTok{(}\FunctionTok{t}\NormalTok{(Xstar\_0)}\SpecialCharTok{\%*\%}\NormalTok{Xstar\_0) }\SpecialCharTok{\%*\%} \FunctionTok{t}\NormalTok{(Xstar\_0)  }\DocumentationTok{\#\# H\_0*}
\NormalTok{SSE\_W.R }\OtherTok{\textless{}{-}} \FunctionTok{t}\NormalTok{(Ystar) }\SpecialCharTok{\%*\%}\NormalTok{ (I\_64}\SpecialCharTok{{-}}\NormalTok{Hstar\_0) }\SpecialCharTok{\%*\%}\NormalTok{ Ystar                   }\DocumentationTok{\#\# SSE\_W(R)}
\NormalTok{dfE.R }\OtherTok{\textless{}{-}} \DecValTok{64{-}1}                                                      \DocumentationTok{\#\# df\_E(R)}

\DocumentationTok{\#\# EWLS F{-}test and Output}
\NormalTok{F\_EWLS }\OtherTok{\textless{}{-}}\NormalTok{ ((SSE\_W.R}\SpecialCharTok{{-}}\NormalTok{SSE\_W.F) }\SpecialCharTok{/}\NormalTok{ (dfE.R}\SpecialCharTok{{-}}\NormalTok{dfE.F)) }\SpecialCharTok{/}\NormalTok{ (SSE\_W.F}\SpecialCharTok{/}\NormalTok{dfE.F)}
\NormalTok{WLS.out }\OtherTok{\textless{}{-}} \FunctionTok{cbind}\NormalTok{(dfE.F,dfE.R,SSE\_W.F,SSE\_W.R,F\_EWLS,}\FunctionTok{qf}\NormalTok{(.}\DecValTok{95}\NormalTok{,dfE.R}\SpecialCharTok{{-}}\NormalTok{dfE.F,dfE.F),}
                 \DecValTok{1}\SpecialCharTok{{-}}\FunctionTok{pf}\NormalTok{(F\_EWLS,dfE.R}\SpecialCharTok{{-}}\NormalTok{dfE.F,dfE.F))}
\FunctionTok{colnames}\NormalTok{(WLS.out) }\OtherTok{\textless{}{-}} \FunctionTok{c}\NormalTok{(}\StringTok{"df(F)"}\NormalTok{, }\StringTok{"df(R)"}\NormalTok{, }\StringTok{"SSE(F)"}\NormalTok{, }\StringTok{"SSE(R)"}\NormalTok{, }\StringTok{"F*"}\NormalTok{,}\StringTok{"F(.95)"}\NormalTok{, }\StringTok{"P(\textgreater{}F*)"}\NormalTok{)}
\FunctionTok{rownames}\NormalTok{(WLS.out) }\OtherTok{\textless{}{-}} \FunctionTok{c}\NormalTok{(}\StringTok{"EWLS"}\NormalTok{)}
\FunctionTok{round}\NormalTok{(WLS.out,}\DecValTok{4}\NormalTok{)}
\end{Highlighting}
\end{Shaded}

\begin{verbatim}
##      df(F) df(R) SSE(F)  SSE(R)     F* F(.95) P(>F*)
## EWLS    60    63     60 81.3196 7.1065 2.7581  4e-04
\end{verbatim}

\begin{Shaded}
\begin{Highlighting}[]
\DocumentationTok{\#\#\#\#\#\#\#\#\#\#\#\#\#\#\#\#\#\#\#\#\#\#\#\#\#\#\#\#\#\#\#\#\#\#\#\#\#\#\#\#\#\#\#\#\#\#\#\#\#\#\#\#\#\#\#\#\#\#\#\#\#\#\#\#\#\#}
\DocumentationTok{\#\# Using aov function with treatment variances as weights}
\DocumentationTok{\#\# Obtain treatment sample sizes and variances  }
\NormalTok{n.trt }\OtherTok{\textless{}{-}} \FunctionTok{as.vector}\NormalTok{(}\FunctionTok{tapply}\NormalTok{(vt}\SpecialCharTok{$}\NormalTok{procKnow, vt}\SpecialCharTok{$}\NormalTok{grp.trt, length))}
\NormalTok{var.trt }\OtherTok{\textless{}{-}} \FunctionTok{as.vector}\NormalTok{(}\FunctionTok{tapply}\NormalTok{(vt}\SpecialCharTok{$}\NormalTok{procKnow, vt}\SpecialCharTok{$}\NormalTok{grp.trt, var))}

\DocumentationTok{\#\# String out weights = 1/var so that weight i is replicated n\_i times}
\NormalTok{vt}\SpecialCharTok{$}\NormalTok{virt.wt }\OtherTok{\textless{}{-}} \FunctionTok{rep}\NormalTok{(}\DecValTok{1}\SpecialCharTok{/}\NormalTok{var.trt, n.trt)}

\DocumentationTok{\#\# Weighted Least Squares ANOVA}
\NormalTok{aov.wt }\OtherTok{\textless{}{-}} \FunctionTok{aov}\NormalTok{(procKnow }\SpecialCharTok{\textasciitilde{}} \FunctionTok{factor}\NormalTok{(grp.trt), }\AttributeTok{weight=}\NormalTok{virt.wt, }\AttributeTok{data=}\NormalTok{vt)}
\FunctionTok{anova}\NormalTok{(aov.wt)}
\end{Highlighting}
\end{Shaded}

\begin{verbatim}
## Analysis of Variance Table
## 
## Response: procKnow
##                 Df Sum Sq Mean Sq F value    Pr(>F)    
## factor(grp.trt)  3  21.32  7.1065  7.1065 0.0003655 ***
## Residuals       60  60.00  1.0000                      
## ---
## Signif. codes:  0 '***' 0.001 '**' 0.01 '*' 0.05 '.' 0.1 ' ' 1
\end{verbatim}

\begin{Shaded}
\begin{Highlighting}[]
\DocumentationTok{\#\# Ordinary Least Squares ANOVA}
\NormalTok{aov.nwt }\OtherTok{\textless{}{-}} \FunctionTok{aov}\NormalTok{(procKnow }\SpecialCharTok{\textasciitilde{}} \FunctionTok{factor}\NormalTok{(grp.trt), }\AttributeTok{data=}\NormalTok{vt)}
\FunctionTok{anova}\NormalTok{(aov.nwt)}
\end{Highlighting}
\end{Shaded}

\begin{verbatim}
## Analysis of Variance Table
## 
## Response: procKnow
##                 Df  Sum Sq Mean Sq F value   Pr(>F)   
## factor(grp.trt)  3  65.664 21.8880  4.9406 0.003931 **
## Residuals       60 265.815  4.4302                    
## ---
## Signif. codes:  0 '***' 0.001 '**' 0.01 '*' 0.05 '.' 0.1 ' ' 1
\end{verbatim}

\[ \nabla \]

\textbf{Welch's Test} was developed to handle the issues that arise in
Scenarios 3 and 4 above. The method involves computing a weighted
\(F\)-statistic for testing for treatment effects. The weights for the
treatment means are the reciprocals of their variances,
\(w_i=n_i/s_i^2\). The procedure (which is easy to run in a spreadsheet
or statistical package) computes the weighted \(F\)-statistic and
generates multiples of that statistic and the error degrees of freedom
to achieve an approximate \(F\)-distribution for the test statistic.

A sketch of the calculations is given below. Note that all calculations
are based on the treatment/group sample sizes, means, and variances.

\[
w_i = \frac{n_i}{s_i^2}  \quad w_{\bullet}=\sum_{i=1}^rw_i \quad
C_W=\sum_{i=1}^r\left[\frac{1}{n_i-1}\left(1-\frac{w_i}{w_{\bullet}}\right)
\right] \quad m_W=\left[1+\frac{2(r-2)}{r^2-1}C_W\right]^{-1}
\]

\[
F^*=\frac{1}{r-1}\left[\sum_{i=1}^rw_i\left(\overline{y}_{i\bullet}\right)^2
- \frac{\left(\sum_{i=1}^rw_i\overline{y}_{i\bullet}\right)^2}{w_{\bullet}}\right]
\qquad\qquad \nu_W=\frac{r^2-1}{3C_W}
\]

\[
F_W=m_WF^* \stackrel{\cdot}{\sim}F_{r-1,\nu_W}
\]

\textbf{Example 4.6 - Virtual Training for a Lifeboat Launching Task}

For the lifeboat training experiment, there were \(r=4\) treatments and
each sample size was \(n=16\). The following R code produces Welch's
\(F_W\)-test directly from the treatment means, variances, and sample
sizes as well as using the \textbf{oneway.test} function in R.

\begin{Shaded}
\begin{Highlighting}[]
\NormalTok{vt }\OtherTok{\textless{}{-}} \FunctionTok{read.csv}\NormalTok{(}\StringTok{"data/virtual\_training.csv"}\NormalTok{)}

\DocumentationTok{\#\#\# Brute{-}Force Computation of Welch\textquotesingle{}s F\_W{-}test}
\NormalTok{vt.n }\OtherTok{\textless{}{-}} \FunctionTok{as.vector}\NormalTok{(}\FunctionTok{tapply}\NormalTok{(vt}\SpecialCharTok{$}\NormalTok{procKnow,vt}\SpecialCharTok{$}\NormalTok{grp.trt,length))}
\NormalTok{vt.mean }\OtherTok{\textless{}{-}} \FunctionTok{as.vector}\NormalTok{(}\FunctionTok{tapply}\NormalTok{(vt}\SpecialCharTok{$}\NormalTok{procKnow,vt}\SpecialCharTok{$}\NormalTok{grp.trt,mean))}
\NormalTok{vt.var }\OtherTok{\textless{}{-}} \FunctionTok{as.vector}\NormalTok{(}\FunctionTok{tapply}\NormalTok{(vt}\SpecialCharTok{$}\NormalTok{procKnow,vt}\SpecialCharTok{$}\NormalTok{grp.trt,var))}
\NormalTok{vt.r }\OtherTok{\textless{}{-}} \FunctionTok{length}\NormalTok{(vt.n)}
\NormalTok{vt.n\_T }\OtherTok{\textless{}{-}} \FunctionTok{sum}\NormalTok{(vt.n)}
\NormalTok{vt.wt }\OtherTok{\textless{}{-}}\NormalTok{ vt.n }\SpecialCharTok{/}\NormalTok{ vt.var}
\NormalTok{vt.wt.sum }\OtherTok{\textless{}{-}} \FunctionTok{sum}\NormalTok{(vt.wt)}
\NormalTok{vt.C\_W }\OtherTok{\textless{}{-}} \FunctionTok{sum}\NormalTok{((}\DecValTok{1}\SpecialCharTok{/}\NormalTok{(vt.n}\DecValTok{{-}1}\NormalTok{)) }\SpecialCharTok{*}\NormalTok{ (}\DecValTok{1}\SpecialCharTok{{-}}\NormalTok{vt.wt}\SpecialCharTok{/}\NormalTok{vt.wt.sum)}\SpecialCharTok{\^{}}\DecValTok{2}\NormalTok{)}
\NormalTok{vt.m\_W }\OtherTok{\textless{}{-}} \DecValTok{1} \SpecialCharTok{/}\NormalTok{ (}\DecValTok{1} \SpecialCharTok{+}\NormalTok{ vt.C\_W}\SpecialCharTok{*}\DecValTok{2}\SpecialCharTok{*}\NormalTok{(vt.r}\DecValTok{{-}2}\NormalTok{)}\SpecialCharTok{/}\NormalTok{(vt.r}\SpecialCharTok{\^{}}\DecValTok{2{-}1}\NormalTok{))}
\NormalTok{vt.Fstar }\OtherTok{\textless{}{-}}\NormalTok{ (}\DecValTok{1}\SpecialCharTok{/}\NormalTok{(vt.r}\DecValTok{{-}1}\NormalTok{)) }\SpecialCharTok{*}\NormalTok{ (}\FunctionTok{sum}\NormalTok{(vt.wt}\SpecialCharTok{*}\NormalTok{vt.mean}\SpecialCharTok{\^{}}\DecValTok{2}\NormalTok{) }\SpecialCharTok{{-}}\NormalTok{ (}\FunctionTok{sum}\NormalTok{(vt.wt}\SpecialCharTok{*}\NormalTok{vt.mean)}\SpecialCharTok{\^{}}\DecValTok{2}\NormalTok{)}\SpecialCharTok{/}\NormalTok{vt.wt.sum)}
\NormalTok{vt.nu\_W }\OtherTok{\textless{}{-}}\NormalTok{ (vt.r}\SpecialCharTok{\^{}}\DecValTok{2{-}1}\NormalTok{) }\SpecialCharTok{/}\NormalTok{ (}\DecValTok{3}\SpecialCharTok{*}\NormalTok{vt.C\_W)}
\NormalTok{vt.F\_W }\OtherTok{\textless{}{-}}\NormalTok{ vt.m\_W }\SpecialCharTok{*}\NormalTok{ vt.Fstar}

\NormalTok{welch.out }\OtherTok{\textless{}{-}} \FunctionTok{cbind}\NormalTok{(vt.r}\DecValTok{{-}1}\NormalTok{,vt.nu\_W,vt.F\_W,}\FunctionTok{qf}\NormalTok{(.}\DecValTok{95}\NormalTok{,vt.r}\DecValTok{{-}1}\NormalTok{,vt.nu\_W),}\DecValTok{1}\SpecialCharTok{{-}}\FunctionTok{pf}\NormalTok{(vt.F\_W,vt.r}\DecValTok{{-}1}\NormalTok{,vt.nu\_W))}
\FunctionTok{colnames}\NormalTok{(welch.out) }\OtherTok{\textless{}{-}} \FunctionTok{c}\NormalTok{(}\StringTok{"df1"}\NormalTok{, }\StringTok{"df2"}\NormalTok{, }\StringTok{"F\_W*"}\NormalTok{, }\StringTok{"F(.95)"}\NormalTok{, }\StringTok{"P(\textgreater{}F\_W*)"}\NormalTok{)}
\FunctionTok{rownames}\NormalTok{(welch.out) }\OtherTok{\textless{}{-}} \FunctionTok{c}\NormalTok{(}\StringTok{"Welch\textquotesingle{}s F{-}test"}\NormalTok{)}
\FunctionTok{round}\NormalTok{(welch.out,}\DecValTok{4}\NormalTok{)}
\end{Highlighting}
\end{Shaded}

\begin{verbatim}
##                df1     df2  F_W* F(.95) P(>F_W*)
## Welch's F-test   3 32.5622 6.827 2.8957   0.0011
\end{verbatim}

\begin{Shaded}
\begin{Highlighting}[]
\DocumentationTok{\#\#\# Using oneway.test function}
\FunctionTok{oneway.test}\NormalTok{(procKnow }\SpecialCharTok{\textasciitilde{}} \FunctionTok{factor}\NormalTok{(grp.trt), }\AttributeTok{var.equal=}\ConstantTok{FALSE}\NormalTok{, }\AttributeTok{data=}\NormalTok{vt)}
\end{Highlighting}
\end{Shaded}

\begin{verbatim}
## 
##  One-way analysis of means (not assuming equal variances)
## 
## data:  procKnow and factor(grp.trt)
## F = 6.827, num df = 3.000, denom df = 32.562, p-value = 0.001072
\end{verbatim}

\subsection{Post-hoc Comparisons with Unequal
Variances}\label{post-hoc-comparisons-with-unequal-variances}

The \textbf{Games-Howell} method can be used to make pairwise
comparisons among all pairs of treatments when variances are unequal.
Like Tukey's method, it makes use of the studentized range distribution.
It makes use of \textbf{Satterthwaite's approximation} for the degrees
of freedom for each pair of means.

\[
\nu_{ii'}=\frac{\left(\frac{s_i^2}{n_i}+\frac{s_{i'}^2}{n_{i'}}\right)^2}
{\frac{\left(\frac{s_i^2}{n_i}\right)^2}{n_i-1} + \frac{\left(\frac{s_{i'}^2}{n_{i'}}\right)^2}{n_{i'}-1}}
\qquad \qquad
GH_{ii'}=\frac{q_{1-\alpha;r,\nu_{ii'}}}{\sqrt{2}}\sqrt{\frac{s_i^2}{n_i}+\frac{s_{i'}^2}{n_{i'}}}
\]

Conclude \(\mu_i\ne\mu_{i'}\) if
\(\left|\overline{y}_{i\bullet}-\overline{y}_{i'\bullet}\right| \ge GH_{ii'}\).
Simultaneous \((1-\alpha)100\%\) Confidence Intervals for
\(\mu_i-\mu_{i'}\) are of the form
\(\left(\overline{y}_{i\bullet}-\overline{y}_{i'\bullet}\right)\pm GH_{ii'}\).

\textbf{Example 4.7 - Virtual Training for a Lifeboat Launching Task}

For the lifeboat training experiment, there were \(r=4\) treatments and
each sample size was \(n=16\). The following R code computes
Simultaneous 95\% Confidence Intervals and adjusted \(P\)-values based
on the Games-Howell method.

\begin{Shaded}
\begin{Highlighting}[]
\NormalTok{vt }\OtherTok{\textless{}{-}} \FunctionTok{read.csv}\NormalTok{(}\StringTok{"data/virtual\_training.csv"}\NormalTok{)}
\FunctionTok{head}\NormalTok{(vt,}\DecValTok{2}\NormalTok{)}
\end{Highlighting}
\end{Shaded}

\begin{verbatim}
##   grp.trt procKnow
## 1       1   4.5614
## 2       1   6.6593
\end{verbatim}

\begin{Shaded}
\begin{Highlighting}[]
\DocumentationTok{\#\#\# Brute{-}Force Computation of Games{-}Howell Comparisons}
\NormalTok{games\_howell }\OtherTok{\textless{}{-}} \ControlFlowTok{function}\NormalTok{(y,trt,alpha) \{}

\NormalTok{ybar }\OtherTok{\textless{}{-}} \FunctionTok{mean}\NormalTok{(y)}
\NormalTok{ybar.trt }\OtherTok{\textless{}{-}} \FunctionTok{as.vector}\NormalTok{(}\FunctionTok{tapply}\NormalTok{(y,trt,mean))}
\NormalTok{var.trt }\OtherTok{\textless{}{-}} \FunctionTok{as.vector}\NormalTok{(}\FunctionTok{tapply}\NormalTok{(y,trt,var))}
\NormalTok{rep.trt }\OtherTok{\textless{}{-}} \FunctionTok{as.vector}\NormalTok{(}\FunctionTok{tapply}\NormalTok{(y,trt,length))}
\NormalTok{num.trt }\OtherTok{\textless{}{-}} \FunctionTok{length}\NormalTok{(ybar.trt)}
\NormalTok{num.pairs }\OtherTok{\textless{}{-}}\NormalTok{ num.trt}\SpecialCharTok{*}\NormalTok{(num.trt}\DecValTok{{-}1}\NormalTok{)}\SpecialCharTok{/}\DecValTok{2}
\NormalTok{sse }\OtherTok{\textless{}{-}} \FunctionTok{sum}\NormalTok{((rep.trt}\DecValTok{{-}1}\NormalTok{)}\SpecialCharTok{*}\NormalTok{var.trt)}
\NormalTok{mse }\OtherTok{\textless{}{-}}\NormalTok{ sse}\SpecialCharTok{/}\NormalTok{(}\FunctionTok{sum}\NormalTok{(rep.trt)}\SpecialCharTok{{-}}\NormalTok{num.trt)}

\NormalTok{i1.i2 }\OtherTok{\textless{}{-}} \FunctionTok{matrix}\NormalTok{(}\FunctionTok{rep}\NormalTok{(}\DecValTok{0}\NormalTok{,}\DecValTok{2}\SpecialCharTok{*}\NormalTok{num.pairs),}\AttributeTok{ncol=}\DecValTok{2}\NormalTok{)}
\NormalTok{ybar.diff }\OtherTok{\textless{}{-}} \FunctionTok{rep}\NormalTok{(}\DecValTok{0}\NormalTok{,num.pairs)}
\NormalTok{gh.df }\OtherTok{\textless{}{-}} \FunctionTok{rep}\NormalTok{(}\DecValTok{0}\NormalTok{,num.pairs)}
\NormalTok{gh.msd }\OtherTok{\textless{}{-}} \FunctionTok{rep}\NormalTok{(}\DecValTok{0}\NormalTok{,num.pairs)}
\NormalTok{gh.LL }\OtherTok{\textless{}{-}} \FunctionTok{rep}\NormalTok{(}\DecValTok{0}\NormalTok{,num.pairs)}
\NormalTok{gh.UL }\OtherTok{\textless{}{-}} \FunctionTok{rep}\NormalTok{(}\DecValTok{0}\NormalTok{,num.pairs)}
\NormalTok{gh.adjpv }\OtherTok{\textless{}{-}} \FunctionTok{rep}\NormalTok{(}\DecValTok{0}\NormalTok{,num.pairs)}
\NormalTok{pair }\OtherTok{\textless{}{-}} \DecValTok{0}

\ControlFlowTok{for}\NormalTok{(i1 }\ControlFlowTok{in} \DecValTok{2}\SpecialCharTok{:}\NormalTok{num.trt) \{}
  \ControlFlowTok{for}\NormalTok{ (i2 }\ControlFlowTok{in} \DecValTok{1}\SpecialCharTok{:}\NormalTok{(i1}\DecValTok{{-}1}\NormalTok{)) \{}
\NormalTok{    pair }\OtherTok{\textless{}{-}}\NormalTok{ pair }\SpecialCharTok{+} \DecValTok{1}
\NormalTok{    i1.i2[pair,] }\OtherTok{\textless{}{-}} \FunctionTok{cbind}\NormalTok{(i1,i2)}
\NormalTok{    ybar.diff[pair] }\OtherTok{\textless{}{-}}\NormalTok{ ybar.trt[i1] }\SpecialCharTok{{-}}\NormalTok{ ybar.trt[i2]}
\NormalTok{    gh.df[pair] }\OtherTok{\textless{}{-}} 
\NormalTok{      (var.trt[i1]}\SpecialCharTok{/}\NormalTok{rep.trt[i1] }\SpecialCharTok{+}\NormalTok{ var.trt[i2]}\SpecialCharTok{/}\NormalTok{rep.trt[i2])}\SpecialCharTok{\^{}}\DecValTok{2} \SpecialCharTok{/} 
\NormalTok{      ((var.trt[i1]}\SpecialCharTok{/}\NormalTok{rep.trt[i1])}\SpecialCharTok{\^{}}\DecValTok{2}\SpecialCharTok{/}\NormalTok{(rep.trt[i1]}\SpecialCharTok{{-}}\DecValTok{1}\NormalTok{) }\SpecialCharTok{+} 
\NormalTok{       (var.trt[i2]}\SpecialCharTok{/}\NormalTok{rep.trt[i2])}\SpecialCharTok{\^{}}\DecValTok{2}\SpecialCharTok{/}\NormalTok{(rep.trt[i2]}\SpecialCharTok{{-}}\DecValTok{1}\NormalTok{))}
\NormalTok{    gh.msd[pair] }\OtherTok{\textless{}{-}} 
\NormalTok{      (}\FunctionTok{qtukey}\NormalTok{(}\DecValTok{1}\SpecialCharTok{{-}}\NormalTok{alpha,num.trt,gh.df[pair])}\SpecialCharTok{/}\FunctionTok{sqrt}\NormalTok{(}\DecValTok{2}\NormalTok{)) }\SpecialCharTok{*}
      \FunctionTok{sqrt}\NormalTok{(var.trt[i1]}\SpecialCharTok{/}\NormalTok{rep.trt[i1]}\SpecialCharTok{+}\NormalTok{var.trt[i2]}\SpecialCharTok{/}\NormalTok{rep.trt[i2])}
\NormalTok{    gh.LL[pair] }\OtherTok{\textless{}{-}}\NormalTok{ ybar.diff[pair]}\SpecialCharTok{{-}}\NormalTok{gh.msd[pair]}
\NormalTok{    gh.UL[pair] }\OtherTok{\textless{}{-}}\NormalTok{ ybar.diff[pair]}\SpecialCharTok{+}\NormalTok{gh.msd[pair]}
\NormalTok{    gh.q }\OtherTok{\textless{}{-}} \FunctionTok{sqrt}\NormalTok{(}\DecValTok{2}\NormalTok{)}\SpecialCharTok{*}\NormalTok{ybar.diff[pair]}\SpecialCharTok{/}\FunctionTok{sqrt}\NormalTok{(var.trt[i1]}\SpecialCharTok{/}\NormalTok{rep.trt[i1]}\SpecialCharTok{+}\NormalTok{var.trt[i2]}\SpecialCharTok{/}\NormalTok{rep.trt[i2])}
\NormalTok{    gh.adjpv[pair] }\OtherTok{\textless{}{-}} \DecValTok{1}\SpecialCharTok{{-}}\FunctionTok{ptukey}\NormalTok{(}\FunctionTok{abs}\NormalTok{(gh.q),num.trt,gh.df[pair])}
\NormalTok{    \}}
\NormalTok{  \}}
\NormalTok{gh.out }\OtherTok{\textless{}{-}} \FunctionTok{cbind}\NormalTok{(i1.i2,ybar.diff,gh.df,gh.msd,gh.LL,gh.UL,gh.adjpv)}
\FunctionTok{colnames}\NormalTok{(gh.out) }\OtherTok{\textless{}{-}} \FunctionTok{c}\NormalTok{(}\StringTok{"Trt i"}\NormalTok{, }\StringTok{"Trt i\textquotesingle{}"}\NormalTok{, }\StringTok{"Ybar diff"}\NormalTok{,}\StringTok{"df"}\NormalTok{, }\StringTok{"MSD"}\NormalTok{, }\StringTok{"LL"}\NormalTok{, }\StringTok{"UL"}\NormalTok{, }\StringTok{"Adj P"}\NormalTok{)}
\FunctionTok{round}\NormalTok{(gh.out,}\DecValTok{4}\NormalTok{)}
\NormalTok{\}}

\FunctionTok{games\_howell}\NormalTok{(vt}\SpecialCharTok{$}\NormalTok{procKnow, vt}\SpecialCharTok{$}\NormalTok{grp.trt, .}\DecValTok{05}\NormalTok{)}
\end{Highlighting}
\end{Shaded}

\begin{verbatim}
##      Trt i Trt i' Ybar diff      df    MSD      LL     UL  Adj P
## [1,]     2      1    2.7778 27.5847 1.6466  1.1312 4.4244 0.0005
## [2,]     3      1    1.8056 26.5999 2.3440 -0.5384 4.1495 0.1759
## [3,]     3      2   -0.9722 22.2357 2.1931 -3.1654 1.2209 0.6150
## [4,]     4      1    1.9444 29.9806 1.8893  0.0552 3.8337 0.0418
## [5,]     4      2   -0.8333 27.2301 1.6756 -2.5089 0.8422 0.5340
## [6,]     4      3    0.1389 26.9707 2.3614 -2.2226 2.5003 0.9985
\end{verbatim}

\begin{Shaded}
\begin{Highlighting}[]
\DocumentationTok{\#\#\# Using the gamesHowellTest function in PMCMRplus package (gives adjusted P{-}values)}
\FunctionTok{gamesHowellTest}\NormalTok{(procKnow }\SpecialCharTok{\textasciitilde{}}\NormalTok{ grp.trt, }\AttributeTok{data=}\NormalTok{vt)}
\end{Highlighting}
\end{Shaded}

\begin{verbatim}
## 
##  Pairwise comparisons using Games-Howell test
\end{verbatim}

\begin{verbatim}
## data: procKnow by grp.trt
\end{verbatim}

\begin{verbatim}
##   1       2       3      
## 2 0.00046 -       -      
## 3 0.17589 0.61503 -      
## 4 0.04183 0.53405 0.99848
\end{verbatim}

\begin{verbatim}
## 
## P value adjustment method: none
\end{verbatim}

\begin{verbatim}
## alternative hypothesis: two.sided
\end{verbatim}

Treatment 2 (Monitor/Keyboard) is significantly higher than Treatment 1
(Lecture/Materials) and Treatment 4 (Head Mounted Display/Wearable
sensors) is significantly higher than Treatment 1.

\[
\nabla
\]

Finally, we simulate Welch's Test for Scenario 4 described previously
where \(\mu_1=\mu_2=\mu_3=100\),
\(\sigma_1=10,\sigma_2=20,\sigma_3=30\), and \(n_1=30,n_2=20,n_3=10\).

\begin{verbatim}
##      Min.   1st Qu.    Median      Mean   3rd Qu.      Max. 
##  0.000788  0.419304  1.053916  1.754281  2.301317 37.070308
\end{verbatim}

\begin{verbatim}
## [1] 0.1637
\end{verbatim}

\begin{verbatim}
## [1] 0.0562
\end{verbatim}

The ANOVA \(F\)-test rejected the null hypothesis in 16.37\% of the
samples, while Welch's \(F_W\)-test rejected it in 5.62\% of the
samples, much closer to the target of 5\%.

\section{Nonparametric Test for Non-normal
Data}\label{nonparametric-test-for-non-normal-data}

The \textbf{Kruskal-Wallis Test} is a rank based test used in cases
where the errors are not normally distributed. In particular, outlying
observations, which have large impacts on the means and standard
deviations have smaller impact when measurements are replaced by ranks.

The \(n_T\) observations are ranked from 1 (smallest) to \(n_T\)
(largest), with ties being given the average of the ranks they would
have received if not exactly equal. For instance, if two observations
tie for the smallest value, they would each receive the rank of 1.5, and
the next smallest observation would get the rank of 3. The sum of the
ranks is \(1+\cdots+n_T=n_T\left(n_T+1\right)/2\).

Once the individual measurements are ranked, rank sums are obtained for
each treatment/group, as well as the Kruskal-Wallis \(H\)-statistic,
which under the null hypothesis of equal medians \(H_0:M_1=\cdots=M_r\)
follows a chi-square distribution with \(r-1\) degrees of freedom
approximately.

\[
R_{ij} = \mbox{rank}\left(Y_{ij}\right) \mbox{ among } Y_{11},\ldots Y_{rn_r}
\qquad R_{i\bullet} = \sum_{j=1}^{n_i}R_{ij} 
\] \[
\mbox{Test Statistic:} \quad 
H^*=\left[\frac{12}{n_T\left(n_T+1\right)}\sum_{i=1}^r\frac{R_{i\bullet}^2}{n_i}\right] -
3\left(n_T+1\right) 
\] \[
\mbox{Rejection Region:} \quad H^* \ge \chi_{r-1}^2   \qquad
P=P\left(\chi_{r-1}^2 \geq H^*\right)
\] When there are ties, an adjustment can be made to \(H^*\), that
rarely makes a large effect. Let the number of groups with ties be \(g\)
and let the \(k^{th}\) group have \(t_k\) tied observations. Then, the
adjusted test statistic is given below.

\[
\mbox{Test Statistic:} \quad 
H^{*'} = 
\frac{H^*}
{\left[1-
\frac{\sum_{k=1}^g\left(t_k-1\right)t_k\left(t_k+1\right)}
{\left(n_T-1\right)n_T\left(n_T+1\right)}
\right]}
\]

When the null hypothesis of equal medians is rejected, approximate
Confidence Intervals for the differences in mean ranks among the
\(g=r(r-1)/2\) pairs of treatments are computed as follows.

\[
\left(\overline{R}_{i\bullet}-\overline{R}_{i'\bullet}\right) \pm
z_{1-\alpha/(2g)} \sqrt{\frac{n_T\left(n_T+1\right)}{2}
\left(\frac{1}{n_i}+\frac{1}{n_{i'}}\right)}
\]

\textbf{Example 4.8 - Virtual Training for a Lifeboat Launching Task}

For the lifeboat training experiment, there were \(r=4\) treatments and
each sample size was \(n=16\). We will directly compute the ranks, test,
and multiple comparisons. Then the \textbf{kruskal.test} function will
be used. Note that the \textbf{rank} function in R correctly adjusts
ranks for ties. The null hypothesis is that the population medians are
equal for the \(r=4\) treatments.

\begin{Shaded}
\begin{Highlighting}[]
\NormalTok{vt }\OtherTok{\textless{}{-}} \FunctionTok{read.csv}\NormalTok{(}\StringTok{"data/virtual\_training.csv"}\NormalTok{)}
\FunctionTok{head}\NormalTok{(vt,}\DecValTok{2}\NormalTok{)}
\end{Highlighting}
\end{Shaded}

\begin{verbatim}
##   grp.trt procKnow
## 1       1   4.5614
## 2       1   6.6593
\end{verbatim}

\begin{Shaded}
\begin{Highlighting}[]
\DocumentationTok{\#\#\# Direct calculations}
    
\NormalTok{vt}\SpecialCharTok{$}\NormalTok{rankY }\OtherTok{\textless{}{-}} \FunctionTok{rank}\NormalTok{(vt}\SpecialCharTok{$}\NormalTok{procKnow)                               }\DocumentationTok{\#\# Ranks of Y values}
\NormalTok{sum.rank }\OtherTok{\textless{}{-}} \FunctionTok{as.vector}\NormalTok{(}\FunctionTok{tapply}\NormalTok{(vt}\SpecialCharTok{$}\NormalTok{rankY, vt}\SpecialCharTok{$}\NormalTok{grp.trt, sum))   }\DocumentationTok{\#\# Rank sums R\_1,...,R\_4}
\NormalTok{n.rank }\OtherTok{\textless{}{-}} \FunctionTok{as.vector}\NormalTok{(}\FunctionTok{tapply}\NormalTok{(vt}\SpecialCharTok{$}\NormalTok{rankY, vt}\SpecialCharTok{$}\NormalTok{grp.trt, length))   }\DocumentationTok{\#\# n\_1,...,n\_4}
\NormalTok{r }\OtherTok{\textless{}{-}} \FunctionTok{length}\NormalTok{(n.rank)                                         }
\NormalTok{n\_T }\OtherTok{\textless{}{-}} \FunctionTok{sum}\NormalTok{(n.rank)}

\NormalTok{KW.stat }\OtherTok{\textless{}{-}}\NormalTok{ (}\DecValTok{12}\SpecialCharTok{/}\NormalTok{(n\_T}\SpecialCharTok{*}\NormalTok{(n\_T}\SpecialCharTok{+}\DecValTok{1}\NormalTok{))) }\SpecialCharTok{*} \FunctionTok{sum}\NormalTok{(sum.rank}\SpecialCharTok{\^{}}\DecValTok{2}\SpecialCharTok{/}\NormalTok{n.rank) }\SpecialCharTok{{-}} \DecValTok{3}\SpecialCharTok{*}\NormalTok{(n\_T}\SpecialCharTok{+}\DecValTok{1}\NormalTok{)  }\DocumentationTok{\#\# K{-}W H* statistic}
\DocumentationTok{\#\# Set{-}up and print output}
\NormalTok{KW.out }\OtherTok{\textless{}{-}} \FunctionTok{cbind}\NormalTok{(KW.stat, r}\DecValTok{{-}1}\NormalTok{, }\FunctionTok{qchisq}\NormalTok{(.}\DecValTok{95}\NormalTok{,r}\DecValTok{{-}1}\NormalTok{), }\DecValTok{1}\SpecialCharTok{{-}}\FunctionTok{pchisq}\NormalTok{(KW.stat,r}\DecValTok{{-}1}\NormalTok{))}
\FunctionTok{colnames}\NormalTok{(KW.out) }\OtherTok{\textless{}{-}} \FunctionTok{c}\NormalTok{(}\StringTok{"KW stat"}\NormalTok{, }\StringTok{"df"}\NormalTok{, }\StringTok{"X2(.95)"}\NormalTok{, }\StringTok{"P(\textgreater{}KW)"}\NormalTok{)}
\FunctionTok{rownames}\NormalTok{(KW.out) }\OtherTok{\textless{}{-}} \FunctionTok{c}\NormalTok{(}\StringTok{"Kruskal{-}Wallis Test"}\NormalTok{)}
\FunctionTok{round}\NormalTok{(KW.out,}\DecValTok{4}\NormalTok{)}
\end{Highlighting}
\end{Shaded}

\begin{verbatim}
##                     KW stat df X2(.95) P(>KW)
## Kruskal-Wallis Test 12.9151  3  7.8147 0.0048
\end{verbatim}

\begin{Shaded}
\begin{Highlighting}[]
\DocumentationTok{\#\# Function for pairwise comparisons of medians}
\NormalTok{KW.CI }\OtherTok{\textless{}{-}} \ControlFlowTok{function}\NormalTok{(y, trt, alpha) \{}
\NormalTok{rank.y }\OtherTok{\textless{}{-}} \FunctionTok{rank}\NormalTok{(y)}
\NormalTok{rank.trt }\OtherTok{\textless{}{-}} \FunctionTok{as.vector}\NormalTok{(}\FunctionTok{tapply}\NormalTok{(rank.y,trt,mean))}
\NormalTok{rep.trt }\OtherTok{\textless{}{-}} \FunctionTok{as.vector}\NormalTok{(}\FunctionTok{tapply}\NormalTok{(rank.y,trt,length))}
\NormalTok{num.trt }\OtherTok{\textless{}{-}} \FunctionTok{length}\NormalTok{(rank.trt)}
\NormalTok{num.pairs }\OtherTok{\textless{}{-}}\NormalTok{ num.trt}\SpecialCharTok{*}\NormalTok{(num.trt}\DecValTok{{-}1}\NormalTok{)}\SpecialCharTok{/}\DecValTok{2}

\NormalTok{i1.i2 }\OtherTok{\textless{}{-}} \FunctionTok{matrix}\NormalTok{(}\FunctionTok{rep}\NormalTok{(}\DecValTok{0}\NormalTok{,}\DecValTok{2}\SpecialCharTok{*}\NormalTok{num.pairs),}\AttributeTok{ncol=}\DecValTok{2}\NormalTok{)}
\NormalTok{rbar.diff }\OtherTok{\textless{}{-}} \FunctionTok{rep}\NormalTok{(}\DecValTok{0}\NormalTok{,num.pairs)}
\NormalTok{r.msd }\OtherTok{\textless{}{-}} \FunctionTok{rep}\NormalTok{(}\DecValTok{0}\NormalTok{,num.pairs)}
\NormalTok{r.LL }\OtherTok{\textless{}{-}} \FunctionTok{rep}\NormalTok{(}\DecValTok{0}\NormalTok{,num.pairs)}
\NormalTok{r.UL }\OtherTok{\textless{}{-}} \FunctionTok{rep}\NormalTok{(}\DecValTok{0}\NormalTok{,num.pairs)}
\NormalTok{r.adjpv }\OtherTok{\textless{}{-}} \FunctionTok{rep}\NormalTok{(}\DecValTok{0}\NormalTok{,num.pairs)}
\NormalTok{pair }\OtherTok{\textless{}{-}} \DecValTok{0}

\ControlFlowTok{for}\NormalTok{(i1 }\ControlFlowTok{in} \DecValTok{2}\SpecialCharTok{:}\NormalTok{num.trt) \{}
  \ControlFlowTok{for}\NormalTok{ (i2 }\ControlFlowTok{in} \DecValTok{1}\SpecialCharTok{:}\NormalTok{(i1}\DecValTok{{-}1}\NormalTok{)) \{}
\NormalTok{    pair }\OtherTok{\textless{}{-}}\NormalTok{ pair }\SpecialCharTok{+} \DecValTok{1}
\NormalTok{    i1.i2[pair,] }\OtherTok{\textless{}{-}} \FunctionTok{cbind}\NormalTok{(i1,i2)}
\NormalTok{    rbar.diff[pair] }\OtherTok{\textless{}{-}}\NormalTok{ rank.trt[i1] }\SpecialCharTok{{-}}\NormalTok{ rank.trt[i2]}
\NormalTok{    r.msd[pair] }\OtherTok{\textless{}{-}} 
      \FunctionTok{qnorm}\NormalTok{(}\DecValTok{1}\SpecialCharTok{{-}}\NormalTok{alpha}\SpecialCharTok{/}\NormalTok{(}\DecValTok{2}\SpecialCharTok{*}\NormalTok{num.pairs)) }\SpecialCharTok{*}
      \FunctionTok{sqrt}\NormalTok{((n\_T}\SpecialCharTok{*}\NormalTok{(n\_T}\SpecialCharTok{+}\DecValTok{1}\NormalTok{)}\SpecialCharTok{/}\DecValTok{12}\NormalTok{)}\SpecialCharTok{*}\NormalTok{(}\DecValTok{1}\SpecialCharTok{/}\NormalTok{rep.trt[i1]}\SpecialCharTok{+}\DecValTok{1}\SpecialCharTok{/}\NormalTok{rep.trt[i2]))}
\NormalTok{    r.LL[pair] }\OtherTok{\textless{}{-}}\NormalTok{ rbar.diff[pair]}\SpecialCharTok{{-}}\NormalTok{r.msd[pair]}
\NormalTok{    r.UL[pair] }\OtherTok{\textless{}{-}}\NormalTok{ rbar.diff[pair]}\SpecialCharTok{+}\NormalTok{r.msd[pair]}
\NormalTok{     r.pv}\OtherTok{\textless{}{-}} 
      \DecValTok{2}\SpecialCharTok{*}\NormalTok{(}\DecValTok{1}\SpecialCharTok{{-}}\FunctionTok{pnorm}\NormalTok{(}\FunctionTok{abs}\NormalTok{(rbar.diff[pair]}\SpecialCharTok{/}\FunctionTok{sqrt}\NormalTok{((n\_T}\SpecialCharTok{*}\NormalTok{(n\_T}\SpecialCharTok{+}\DecValTok{1}\NormalTok{)}\SpecialCharTok{/}\DecValTok{12}\NormalTok{)}\SpecialCharTok{*}\NormalTok{(}\DecValTok{1}\SpecialCharTok{/}\NormalTok{rep.trt[i1]}\SpecialCharTok{+}\DecValTok{1}\SpecialCharTok{/}\NormalTok{rep.trt[i2])))))}
\NormalTok{    r.adjpv[pair] }\OtherTok{\textless{}{-}} \FunctionTok{min}\NormalTok{(}\DecValTok{1}\NormalTok{,num.pairs}\SpecialCharTok{*}\NormalTok{r.pv)}
\NormalTok{    \}}
\NormalTok{  \}}
\NormalTok{r.out }\OtherTok{\textless{}{-}} \FunctionTok{cbind}\NormalTok{(i1.i2,rbar.diff,r.msd,r.LL,r.UL,r.adjpv)}
\FunctionTok{colnames}\NormalTok{(r.out) }\OtherTok{\textless{}{-}} \FunctionTok{c}\NormalTok{(}\StringTok{"Trt i"}\NormalTok{, }\StringTok{"Trt i\textquotesingle{}"}\NormalTok{, }\StringTok{"Rank diff"}\NormalTok{, }\StringTok{"MSD"}\NormalTok{, }\StringTok{"LL"}\NormalTok{, }\StringTok{"UL"}\NormalTok{, }\StringTok{"Adj P"}\NormalTok{)}
\FunctionTok{round}\NormalTok{(r.out,}\DecValTok{4}\NormalTok{)}
\NormalTok{\}}

\FunctionTok{KW.CI}\NormalTok{(vt}\SpecialCharTok{$}\NormalTok{procKnow, vt}\SpecialCharTok{$}\NormalTok{grp.trt, }\FloatTok{0.05}\NormalTok{)}
\end{Highlighting}
\end{Shaded}

\begin{verbatim}
##      Trt i Trt i' Rank diff     MSD       LL      UL  Adj P
## [1,]     2      1    22.625 17.3671   5.2579 39.9921 0.0035
## [2,]     3      1    16.875 17.3671  -0.4921 34.2421 0.0622
## [3,]     3      2    -5.750 17.3671 -23.1171 11.6171 1.0000
## [4,]     4      1    15.250 17.3671  -2.1171 32.6171 0.1231
## [5,]     4      2    -7.375 17.3671 -24.7421  9.9921 1.0000
## [6,]     4      3    -1.625 17.3671 -18.9921 15.7421 1.0000
\end{verbatim}

\begin{Shaded}
\begin{Highlighting}[]
\DocumentationTok{\#\#\# kruskal.test function}

\FunctionTok{kruskal.test}\NormalTok{(procKnow }\SpecialCharTok{\textasciitilde{}}\NormalTok{ grp.trt, }\AttributeTok{data=}\NormalTok{vt)}
\end{Highlighting}
\end{Shaded}

\begin{verbatim}
## 
##  Kruskal-Wallis rank sum test
## 
## data:  procKnow by grp.trt
## Kruskal-Wallis chi-squared = 12.915, df = 3, p-value = 0.004824
\end{verbatim}

The test concludes that the medians are not all equal. Based on the
simultaneous comparisons, conclude that the median for treatment 2
(Monitor/Keyboard) is significantly higher than that for treatment 1
(Lecture/Materials). No other medians are significantly different.

\[\nabla \]

\chapter{Balanced Two Factor Designs}\label{balanced-two-factor-designs}

In this chapter, we introduce two factor designs with equal numbers of
replicates \((n)\) per treatment (factor combination). The factors will
be generically labelled as \textbf{A} with \(a\) levels and \textbf{B}
with \(b\) levels. The total number of observations is \(n_T=abn\).

\section{Introduction}\label{introduction-1}

In a controlled experiment, experimental units are obtained and randomly
assigned to the \(ab\) treatments, with \(n\) units per treatment. In
observational studies, \(n\) units are sampled from each of the \(ab\)
populations/subpopulations.

The \textbf{1-factor-at-a-time} method is used in some fields of study.
This involves selecting a particular level of factor A and obtain the
best level of factor B. Then, at the ``best'' level of factor B, obtain
the best level of factor A. This method has problems, in particular it
does not permit studying interactions between the factors. It is best
(when feasible) to conduct the experiment at all \(ab\) treatment
levels.

We first consider the population mean structures for \textbf{additive}
and \textbf{interaction} models based on hypothetical values for an
experiment of the ``halo effect''.

\textbf{Example 5.1 - Halo Effect}

In this study, researchers were interested in observing the ``halo
effect'' where individuals who are known to be good/bad in a particular
dimension and assumed to also be good/bad in another dimension. This
phenomenon has been studied involving people, consumer products, and in
many other subject areas. This hypothetical data is taken from a
psychology experiment where students were given a picture of a student
who had written an essay (Good, Bad, No Picture), and an essay that was
supposedly written by the student (Good, Poor). The response measured
was the score that the student assigned to the essay. Each student rater
was assigned to one Picture/Essay Quality condition
(\citeproc{ref-landy_1974}{Landy and Sigall 1974}).

For this example, factor A is the Picture condition with \(a=3\) levels
and factor B is Essay quality with \(b=2\) levels. Let \(\mu_{ij}\) be
the mean when factor A is at level \(i\) and factor \(B\) is at level
\(j\). Further, let \(\mu_{i\bullet}\) be the mean when factor A is at
level \(i\) (across levels of factor B) and \(\mu_{\bullet j}\) be the
mean when factor B is at level \(j\). Finally \(\mu_{\bullet\bullet}\)
is the overall mean across all levels of factors A and B.

Consider the following hypothetical means.

\begin{Shaded}
\begin{Highlighting}[]
\NormalTok{means }\OtherTok{\textless{}{-}} \FunctionTok{matrix}\NormalTok{(}\FunctionTok{c}\NormalTok{(}\DecValTok{25}\NormalTok{,}\DecValTok{18}\NormalTok{,}\DecValTok{20}\NormalTok{,}\DecValTok{21}\NormalTok{, }\DecValTok{17}\NormalTok{,}\DecValTok{10}\NormalTok{,}\DecValTok{12}\NormalTok{,}\DecValTok{13}\NormalTok{, }\DecValTok{21}\NormalTok{,}\DecValTok{14}\NormalTok{,}\DecValTok{16}\NormalTok{,}\DecValTok{17}\NormalTok{), }\AttributeTok{ncol=}\DecValTok{3}\NormalTok{)}
\FunctionTok{colnames}\NormalTok{(means) }\OtherTok{\textless{}{-}} \FunctionTok{c}\NormalTok{(}\StringTok{"j=1"}\NormalTok{, }\StringTok{"j=2"}\NormalTok{, }\StringTok{"Row Mean"}\NormalTok{)}
\FunctionTok{rownames}\NormalTok{(means) }\OtherTok{\textless{}{-}} \FunctionTok{c}\NormalTok{(}\StringTok{"i=1"}\NormalTok{, }\StringTok{"i=2"}\NormalTok{, }\StringTok{"i=3"}\NormalTok{, }\StringTok{"Col Mean"}\NormalTok{)}
\NormalTok{means }
\end{Highlighting}
\end{Shaded}

\begin{verbatim}
##          j=1 j=2 Row Mean
## i=1       25  17       21
## i=2       18  10       14
## i=3       20  12       16
## Col Mean  21  13       17
\end{verbatim}

For the \textbf{additive effects} model, we define the following
parameters based on the population means.

\[
\mu_{ij} = \mu_{\bullet\bullet} + \alpha_i + \beta_j \qquad \mbox{ s.t. }
\sum_{i=1}^a \alpha_i=\sum_{j=1}^b\beta_j=0 
\qquad \mu_{\bullet\bullet}=\frac{1}{ab}\sum_{i=1}^a \sum_{j=1}^b \mu_{ij}
\] \[
\mu_{i\bullet} = \frac{1}{b}\sum_{j=1}^b \mu_{ij}
= \frac{1}{b}\sum_{j=1}^b \left(\mu_{\bullet\bullet} + \alpha_i + \beta_j\right)
=\frac{1}{b}\left[b\mu_{\bullet\bullet} + b\alpha_i + \sum_{j=1}^b\beta_j\right]=
\mu_{\bullet\bullet} + \alpha_i  \qquad
\mu_{\bullet j}=\mu_{\bullet\bullet}+\beta_j
\] \[
\alpha_i = \mu_{i\bullet} - \mu_{\bullet\bullet} \qquad
\beta_j = \mu_{\bullet j} - \mu_{\bullet\bullet} \qquad \qquad
\mu_{\bullet\bullet} = \frac{1}{a}\sum_{i=1}^a\mu_{i\bullet} =
\frac{1}{b}\sum_{ij1}^b\mu_{\bullet j}
\]

For the halo effect values, we obtain the following parameters.

\[
\alpha_1=\mu_{1\bullet} - \mu_{\bullet\bullet} = 21-17=4 \quad
\alpha_2=14-17=-3 \quad \alpha_3=16-17=-1 \qquad \alpha_1+\alpha_2+\alpha_3=0
\] \[
\beta_1 = \mu_{\bullet 1} - \mu_{\bullet\bullet} = 21-17=4 \qquad
\beta_2=13-17 = -4 \qquad \beta_1+\beta_2=0
\]

Plots of the means for each cell are given in Figure @ref(fig:ch0502).
The first plot gives means versus Picture with separate lines for each
Essay Quality. The second plot gives means versus Essay Quality with
separate lines for each Picture. In both cases lines are parallel,
consistent with an additive model. This is an example of an
\textbf{interaction plot} (though for these means, there is no
interaction).

\begin{figure}

{\centering \includegraphics[width=0.8\linewidth]{bookdownproj_STA4211Notes1_files/figure-latex/ch0502-1} 

}

\caption{Halo effect means - Interaction plots with additive effects}\label{fig:ch0502}
\end{figure}

The effect of Good versus Poor Essay Quality is the same for each
Picture condition and the effects of the Picture conditions are the same
for each Essay Quality.

Now, consider the following mean structure, consistent with an
\textbf{interaction} model.

\begin{verbatim}
##          j=1 j=2 Row Mean
## i=1       23  19       21
## i=2       20   8       14
## i=3       20  12       16
## Col Mean  21  13       17
\end{verbatim}

For the \textbf{interaction effects} model, we define the following
parameters based on the population means.

\[
\mu_{ij} = \mu_{\bullet\bullet} + \alpha_i + \beta_j + \left(\alpha\beta\right)_{ij}\qquad \mbox{ s.t. }
\sum_{i=1}^a \alpha_i=\sum_{j=1}^b\beta_j=\sum_{i=1}^a\left(\alpha\beta\right)_{ij} =
\sum_{j=1}^b\left(\alpha\beta\right)_{ij} =0 
\] \[
\left(\alpha\beta\right)_{ij} =
\mu_{ij} - \mu_{\bullet\bullet} - \alpha_i - \beta_j
=\mu_{ij} - \mu_{\bullet\bullet} - \left(\mu_{i\bullet} - \mu_{\bullet\bullet}\right)
-\left(\mu_{\bullet j} - \mu_{\bullet\bullet}\right) =
\mu_{ij} - \mu_{i\bullet} - \mu_{\bullet j} +  \mu_{\bullet\bullet}
\] For the halo effect example, we obtain the following values for the
interaction effects.

\[ 
\left(\alpha\beta\right)_{11} = 23-21-21+17=-2 \qquad
\left(\alpha\beta\right)_{12} = 19-21-13+17=2
\] \[ 
\left(\alpha\beta\right)_{21} = 20-14-21+17=2 \qquad
\left(\alpha\beta\right)_{22} = 8-14-13+17=-2
\]

\[ 
\left(\alpha\beta\right)_{31} = 20-16-21+17=0 \qquad
\left(\alpha\beta\right)_{32} = 12-16-13+17=0
\]

Interaction plots are given in Figure @ref(fig:ch0504). In this case,
the lines are not parallel.

\begin{figure}

{\centering \includegraphics[width=0.8\linewidth]{bookdownproj_STA4211Notes1_files/figure-latex/ch0504-1} 

}

\caption{Halo effect means - Interaction plots with non-additive effects}\label{fig:ch0504}
\end{figure}

\[
\nabla
\] In practice, we don't know the structure of the population means, and
must estimate them and make inferences regarding them based on sample
data. Before describing the data model, we make a few comments regarding
interactions.

\begin{itemize}
\tightlist
\item
  In many situations, the interaction among effects is small relative to
  main effects and can be ignored.
\item
  Transformations (as in the Box-Cox transformation) can be applied
  which may remove an interaction.
\item
  In many research settings, interactions may be hypothesized and have
  interesting theoretical interpretations.
\item
  When factors are ordinal or quantitative, interactions that are
  synergistic or antagonistic can be observed.
\end{itemize}

\section{Two Factor Analysis of
Variance}\label{two-factor-analysis-of-variance}

In this section, we describe the 2-factor Analysis of Variance. We will
first consider the \textbf{Cell Means} model, where interest is only on
the \(ab\) means, typically trying to determine the ``best'' treatment.
Then we consider the \textbf{Factor Effects} model, which measures main
effects and interactions among the treatments. In each model,
\(Y_{ijk}\) represents a random outcome when factor A is at level \(i\),
factor B is at level \(j\), and replicate is \(k\).

\subsection{Cell Means and Factor Effects
Models}\label{cell-means-and-factor-effects-models}

In this model, the focus is on the \(ab\) cell means among the
combination of levels of factors A and B.

\[
Y_{ijk}=\mu_{ij}+\epsilon_{ijk} \quad i=1,\ldots,a; \quad j=1,\ldots,b;
\quad k=1,\ldots,n \quad n_T=abn
\]

\[
\mu_{ij} \equiv \mbox{ mean when factor A is at level } i \mbox{ and factor B is at level } j   \qquad
\epsilon_{ijk}\sim NID\left(0,\sigma^2\right)
\]

The model can be written in matrix form with \(Y\) being
\(n_T\times 1\), \(X\) being \(n_T \times ab\), \(\beta\) being
\(ab \times 1\), and \(\epsilon\) being \(n_T\times 1\). For the case
where \(a=b=n=2\), we have the following structure.

\[
Y=
\begin{bmatrix}
Y_{111} \\Y_{112}\\ Y_{121}\\Y_{122}\\Y_{211}\\Y_{212}\\Y_{221}\\Y_{222}\\
\end{bmatrix}
\quad
X=
\begin{bmatrix}
1 & 0 & 0 & 0 \\ 1 & 0 & 0 & 0\\
0 & 1 & 0 & 0 \\ 0 & 1 & 0 & 0\\
0 & 0 & 1 & 0 \\ 0 & 0 & 1 & 0\\
0 & 0 & 0 & 1 \\ 0 & 0 & 0 & 1\\
\end{bmatrix}
\quad
\beta=
\begin{bmatrix}
\mu_{11} \\ \mu_{12}  \\ \mu_{21} \\ \mu_{22}\\
\end{bmatrix}
\quad
\epsilon =
\begin{bmatrix}
\epsilon_{111} \\\epsilon_{112}\\ \epsilon_{121}\\\epsilon_{122}\\
\epsilon_{211}\\\epsilon_{212}\\\epsilon_{221}\\\epsilon_{222}\\
\end{bmatrix}
\]

Under this model, \(V\{Y\}=V\{\epsilon\}=\sigma^2I_{n_T}\).

The \textbf{Factor Effects} model decomposes \(\mu_{ij}\) into main
effects among the levels of factors A and B, as well as the interactions
among the combinations of levels of factors A and B.

\[
Y_{ijk}=\mu_{\bullet\bullet}+\alpha_i+\beta_j+\left(\alpha\beta\right)_{ij}+\epsilon_{ijk} \quad i=1,\ldots,a;\quad j= 1,\ldots,b; \quad k=1,\ldots,n
\]

\[
 \mu_{\bullet\bullet}=\frac{1}{ab}\sum_{i=1}^a\sum_{j=1}^b\mu_{ij} 
 \qquad
 \mu_{i\bullet}=\frac{1}{b}\sum_{j=1}^b\mu_{ij}
 \qquad
\mu_{\bullet j}=\frac{1}{a}\sum_{i=1}^a\mu_{ij}
\] \[
\mbox{Main effect of level } i \mbox{ factor A:} \quad
\alpha_i = \mu_{i\bullet} - \mu_{\bullet\bullet} \qquad
\sum_{i=1}^a \alpha_i=0
\] \[
\mbox{Main effect of level } j \mbox{ factor B:} \quad
\beta_j = \mu_{\bullet j} - \mu_{\bullet\bullet} \qquad
\sum_{j=1}^b \beta_j=0
\] \[
\mbox{Interaction effect of A at level } i \mbox{ and B at level } j:
\quad
\left(\alpha\beta\right)_{ij} =
\mu_{ij}-\mu_{i\bullet}-\mu_{\bullet j}+\mu_{\bullet\bullet}
\] \[
\sum_{i=1}^a\left(\alpha\beta\right)_{ij} =
\sum_{j=1}^b\left(\alpha\beta\right)_{ij}=0
\]

\[
Y_{ijk} \sim N\left(\mu_{\bullet\bullet}+\alpha_i+\beta_j+\left(\alpha\beta\right)_{ij},\sigma^2\right) \quad \mbox{independent}
\] For the case where \(a=b=n=2\), we have the following structure.

\[
Y=
\begin{bmatrix}
Y_{111} \\Y_{112}\\ Y_{121}\\Y_{122}\\Y_{211}\\Y_{212}\\Y_{221}\\Y_{222}\\
\end{bmatrix}
\quad
X=
\begin{bmatrix}
1 & \phantom{-}1 & \phantom{-}1 & \phantom{-}1 \\ 
1 & \phantom{-}1 & \phantom{-}1 & \phantom{-}1 \\
1 & \phantom{-}1 & -1 & -1 \\ 1 & \phantom{-}1 & -1 & -1\\
1 & -1 & \phantom{-}1 & -1 \\ 1 & -1 & \phantom{-}1 & -1\\
1 & -1 & -1 & \phantom{-}1 \\ 1 & -1 & -1 & \phantom{-}1\\
\end{bmatrix}
\quad
\beta=
\begin{bmatrix}
\mu_{\bullet\bullet} \\ \alpha_1  \\ \beta_1 \\ (\alpha\beta)_{11}\\
\end{bmatrix}
\quad
\epsilon =
\begin{bmatrix}
\epsilon_{111} \\\epsilon_{112}\\ \epsilon_{121}\\\epsilon_{122}\\
\epsilon_{211}\\\epsilon_{212}\\\epsilon_{221}\\\epsilon_{222}\\
\end{bmatrix}
\]

\subsection{Least Squares/Maximum Likelihood
Estimators}\label{least-squaresmaximum-likelihood-estimators}

First, we define the following means.

\[
\overline{Y}_{ij\bullet} = \frac{1}{n}\sum_{k=1}^nY_{ijk} \quad
\overline{Y}_{i\bullet\bullet} = 
\frac{1}{bn}\sum_{j=1}^b\sum_{k=1}^nY_{ijk} \quad
\overline{Y}_{\bullet j \bullet} = 
\frac{1}{an}\sum_{i=1}^a\sum_{k=1}^nY_{ijk} 
\] \[
\overline{Y}_{\bullet \bullet \bullet} = 
\frac{1}{abn}\sum_{i=1}^a\sum_{j=1}^b\sum_{k=1}^nY_{ijk}
\]

For the cell means model, the least squares estimator for \(\mu_{ij}\)
is simply the sample mean for that cell, \(\overline{Y}_{ij\bullet}\).

\[
Q=\sum_{i=1}^a\sum_{j=1}^b\sum_{k=1}^n\epsilon_{ijk}^2 =
\sum_{i=1}^a\sum_{j=1}^b\sum_{k=1}^n\left(Y_{ijk}-\mu_{ij}\right)^2
\quad
\frac{\partial Q}{\partial \mu_{ij}}=
-2\sum_{k=1}^n\left(Y_{ijk}-\mu_{ij}\right)\]

\[\mbox{Set: }
\frac{\partial Q}{\partial \mu_{ij}}=0
\quad \Rightarrow \quad
\sum_{k=1}^nY_{ijk}=n\hat{\mu}_{ij}
\quad \Rightarrow \quad
\hat{\mu}_{ij}=\overline{Y}_{ij\bullet}
\] The fitted values and residuals for the cell means model are given
below.

\[
\hat{Y}_{ijk}=\hat{\mu}_{ij}=\overline{Y}_{ij\bullet}
\qquad
e_{ijk}=Y_{ijk}-\hat{Y}_{ijk} = Y_{ijk}-\overline{Y}_{ij\bullet}
\]

For the factor effects model, we obtain the following least squares
estimators.

\[
Q=\sum_{i=1}^a\sum_{j=1}^b\sum_{k=1}^n
\left(Y_{ijk}-\mu_{\bullet\bullet}-\alpha_i-\beta_j-(\alpha\beta)_{ij}\right)^2
\] \[
\frac{\partial Q}{\partial \mu_{\bullet\bullet}}=-2
\sum_{i=1}^a\sum_{j=1}^b\sum_{k=1}^n
\left(Y_{ijk}-\mu_{\bullet\bullet}-\alpha_i-\beta_j-(\alpha\beta)_{ij}\right)
\quad \mbox{Set } \frac{\partial Q}{\partial \mu_{\bullet\bullet}}=0:
\] \[
\quad \Rightarrow \quad \sum_{i=1}^a\sum_{j=1}^b\sum_{k=1}^n Y_{ijk} = 
abn\hat{\mu}_{\bullet\bullet}+bn\sum_{i=1}^a\alpha_i+an\sum_{j=1}^b\beta_j+
n\sum_{i=1}^a\sum_{j=1}^b(\alpha\beta)_{ij}=abn\hat{\mu}_{\bullet\bullet}+0+0+0
\] \[
\quad \Rightarrow \quad \hat{\mu}_{\bullet\bullet} = \frac{1}{abn}
\sum_{i=1}^a\sum_{j=1}^b\sum_{k=1}^n Y_{ijk}=
\overline{Y}_{\bullet\bullet\bullet}
\] \[
\frac{\partial Q}{\partial \alpha_i}=-2\sum_{j=1}^b\sum_{k=1}^n
\left(Y_{ijk}-\mu_{\bullet\bullet}-\alpha_i-\beta_j-(\alpha\beta)_{ij}\right)
\quad \mbox{Set } \frac{\partial Q}{\partial \alpha_i}=0
\] \[
\quad \Rightarrow \quad \sum_{j=1}^b\sum_{k=1}^n Y_{ijk} =
bn\hat{\mu}_{\bullet\bullet} + bn\hat{\alpha}_i + n\sum_{j=1}^b\beta_j+
n\sum_{j=1}^b(\alpha\beta)_{ij}=
bn\overline{Y}_{\bullet\bullet\bullet}+bn\hat{\alpha}_i+0+0
\] \[
\quad \Rightarrow \quad \hat{\alpha}_i = 
\overline{Y}_{i\bullet\bullet}-\overline{Y}_{\bullet\bullet\bullet}
\quad \mbox { by similar argument } \quad 
\hat{\beta}_j = 
\overline{Y}_{\bullet j\bullet}-\overline{Y}_{\bullet\bullet\bullet}
\] \[
\frac{\partial Q}{\partial (\alpha\beta)_{ij}}=-2\sum_{k=1}^n
\left(Y_{ijk}-\mu_{\bullet\bullet}-\alpha_i-\beta_j-(\alpha\beta)_{ij}\right)
\quad \mbox{Set } \frac{\partial Q}{\partial (\alpha\beta)_{ij}}=0
\] \[
\quad \Rightarrow \quad \sum_{k=1}^n Y_{ijk} =
n\hat{\mu}_{\bullet\bullet} + n\hat{\alpha}_i + n\hat{\beta}_j+
n\hat{(\alpha\beta)}_{ij}
\] \[
\quad\Rightarrow\quad \hat{\left(\alpha\beta\right)}_{ij}=
\overline{Y}_{ij\bullet} - \overline{Y}_{i\bullet\bullet}
-\overline{Y}_{\bullet j \bullet} + \overline{Y}_{\bullet\bullet\bullet}
\]

The fitted values and the residuals for the factor effects model are the
same as for the cell means model.

\[
\hat{Y}_{ijk}=\hat{\mu}_{\bullet\bullet} + \hat{\alpha}_i
+\hat{\beta}_j+\hat{\left(\alpha\beta\right)}_{ij}=
\overline{Y}_{ij\bullet}
\qquad 
e_{ijk}=Y_{ijk}-\hat{Y}_{ijk}=Y_{ijk}-\overline{Y}_{ij\bullet}
\]

\subsection{Sums of Squares and the Analysis of
Variance}\label{sums-of-squares-and-the-analysis-of-variance}

For the cell means model, the 2-Way ANOVA simplifies to a 1-Way ANOVA
with \(r=ab\) treatments.

\[
Y_{ijk}-\overline{Y}_{\bullet\bullet\bullet} =
\left(Y_{ijk}-\overline{Y}_{ij\bullet}\right)+
\left(\overline{Y}_{ij\bullet}-
      \overline{Y}_{\bullet\bullet\bullet}\right)
\] This identity leads to the sums of squares, as the cross-products sum
to zero.

\[
\sum_{i=1}^a\sum_{j=1}^b\sum_{k=1}^n
\left(Y_{ijk}-\overline{Y}_{\bullet\bullet\bullet}\right)^2
=
\sum_{i=1}^a\sum_{j=1}^b\sum_{k=1}^n
\left(Y_{ijk}-\overline{Y}_{ij\bullet}\right)^2
+
\sum_{i=1}^a\sum_{j=1}^b\sum_{k=1}^n
\left(\overline{Y}_{ij\bullet}-
      \overline{Y}_{\bullet\bullet\bullet}\right)^2
\]

\[
SSTO=SSE+SSTR \qquad df_{TO}=n_T-1=\left(n_T-ab\right) +
\left(ab-1\right)=df_E + df_{TR}
\]

For the factor effects model, we decompose the treatment deviations and
sums of squares into main effects and interactions among levels of
factors A and B.

\[
Y_{ijk}-\overline{Y}_{\bullet\bullet\bullet} =
\] \[
\left(Y_{ijk}-\overline{Y}_{ij\bullet}\right) +
\left(\overline{Y}_{i\bullet\bullet}-
\overline{Y}_{\bullet\bullet\bullet}\right) +
\left(\overline{Y}_{\bullet j \bullet}-
\overline{Y}_{\bullet\bullet\bullet}\right) +
\left(\overline{Y}_{ij\bullet}-
      \overline{Y}_{i\bullet\bullet} -
      \overline{Y}_{\bullet j \bullet} +
      \overline{Y}_{\bullet\bullet\bullet}\right)
\]

Again, sums of cross-product terms are zero and we partition the sums of
squares as follows.

\[
\sum_{i=1}^a\sum_{j=1}^b\sum_{k=1}^n
\left(Y_{ijk}-\overline{Y}_{\bullet\bullet\bullet}\right)^2
=
\]

\[
\sum_{i=1}^a\sum_{j=1}^b\sum_{k=1}^n
\left(Y_{ijk}-\overline{Y}_{ij\bullet}\right)^2 +
\sum_{i=1}^a\sum_{j=1}^b\sum_{k=1}^n
\left(\overline{Y}_{i\bullet\bullet}-
\overline{Y}_{\bullet\bullet\bullet}\right)^2 +
\sum_{i=1}^a\sum_{j=1}^b\sum_{k=1}^n
\left(\overline{Y}_{\bullet j \bullet}-
\overline{Y}_{\bullet\bullet\bullet}\right)^2 +
\]

\[
+ \sum_{i=1}^a\sum_{j=1}^b\sum_{k=1}^n
\left(\overline{Y}_{ij\bullet}-
      \overline{Y}_{i\bullet\bullet} -
      \overline{Y}_{\bullet j \bullet} +
      \overline{Y}_{\bullet\bullet\bullet}\right)^2
\]

\[
SSTO = SSE + SSA + SSB + SSAB
\]

\[
df_{TO}=n_T-1 = \left(n_T-ab\right) + (a-1) + (b-1) + (a-1)(b-1)=
\] \[
=df_E + df_A + df_B + df_{AB}
\]

For the factor effects model, we summarize the following sums of
squares, degrees of freedom, and expected mean squares.

\[
\mbox{Factor A:} \quad
SSA=bn\sum_{i=1}^a
\left(\overline{Y}_{i\bullet\bullet}-
\overline{Y}_{\bullet\bullet\bullet}\right)^2
\qquad
df_A=a-1
\qquad
MSA=\frac{SSA}{a-1}
\] \[
 E\{MSA\}=\sigma^2+\frac{bn\sum_{i=1}^a\alpha_i^2}{a-1}
\]

\[
\mbox{Factor B:} \quad
SSB=an\sum_{j=1}^b
\left(\overline{Y}_{\bullet j \bullet}-
\overline{Y}_{\bullet\bullet\bullet}\right)^2
\qquad
df_B=b-1
\qquad
MSB=\frac{SSB}{b-1}
\] \[
 E\{MSB\}=\sigma^2+\frac{an\sum_{j=1}^b\beta_j^2}{b-1}
\]

\[
\mbox{AB Interaction:} \quad
SSAB=n\sum_{i=1}^a\sum_{j=1}^b
\left(\overline{Y}_{ij\bullet}-
      \overline{Y}_{i\bullet\bullet} -
      \overline{Y}_{\bullet j \bullet} +
      \overline{Y}_{\bullet\bullet\bullet}\right)^2
\] \[
df_{AB}=(a-1)(b-1)
\qquad
MSAB=\frac{SSAB}{(a-1)(b-1)}
\]

\[
E\{MSAB\}=\sigma^2 + 
\frac{n\sum_{i=1}^a\sum_{j=1}^b(\alpha\beta)_{ij}^2}{(a-1)(b-1)}
\]

\[
\mbox{Error:} \quad
SSE=
\sum_{i=1}^a\sum_{j=1}^b\sum_{k=1}^n
\left(Y_{ijk}-\overline{Y}_{ij\bullet}\right)^2
\qquad
df_E=ab(n-1)=n_T-ab
\] \[
MSE=\frac{SSE}{ab(n-1)}
\qquad
E\{MSE\}=\sigma^2
\]

When working with published results, the error sum of squares can be
obtained from the treatment standard deviations as in the 1-Way ANOVA.

\[
SSE=\sum_{i=1}^a\sum_{j=1}^b\sum_{k=1}^n
\left(Y_{ijk}-\overline{Y}_{ij\bullet}\right)^2=
(n-1)\sum_{i=1}^a\sum_{j=1}^bs_{ij}^2
\]

When testing for the effects of the factors, begin with a test for
interaction, then test for main effects. If the interaction is
important, the tests for the main effects may be deceiving, as the
effects of factors A and B depend of the level of the other factor.

The test for interaction is given here.

\[
H_0^{AB}:\left(\alpha\beta\right)_{11}=\cdots =\left(\alpha\beta\right)_{ab}=0
\qquad
H_A^{AB}: \mbox{ Not all } \left(\alpha\beta\right)_{ij}=0
\]

\[
\mbox{Test Statistic:} \quad F_{AB}^*=\frac{MSAB}{MSE}
\qquad
\mbox{Rejection Region:} \quad F_{AB}^* \ge F_{1-\alpha;(a-1)(b-1),ab(n-1)}
\]

The tests for the main effects of factors A and B are conducted as
follow.

\[
H_0^A: \alpha_1 = \cdots = \alpha_a=0 \qquad
H_A^A: \mbox { Not all } \alpha_i=0
\]

\[
\mbox{Test Statistic:} \quad F_{A}^*=\frac{MSA}{MSE}
\qquad
\mbox{Rejection Region:} \quad F_{A}^* \ge F_{1-\alpha;a-1,ab(n-1)}
\]

\[
H_0^B: \beta_1 = \cdots = \beta_b=0 \qquad
H_A^B: \mbox { Not all } \beta_j=0
\]

\[
\mbox{Test Statistic:} \quad F_{B}^*=\frac{MSB}{MSE}
\qquad
\mbox{Rejection Region:} \quad F_{B}^* \ge F_{1-\alpha;b-1,ab(n-1)}
\]

Effect sizes for the main effects and interaction can be obtained from
the sums of squares in the Analysis of Variance. The effect size
\(\eta^2\) is obtained by dividing each sum of squares by the total sum
of squares. Partial \(\eta^2\) is obtained by dividing each sum of
squares by the sum of it and the error sum of squares.

\[
\eta_A^2 = \frac{SSA}{SSTO}  \qquad \qquad \eta_B^2 = \frac{SSB}{SSTO} \qquad \qquad \eta_{AB}^2=\frac{SSAB}{SSTO}
\]

\[
\mbox{Partial-} \eta_A^2 = \frac{SSA}{SSA+SSE}   \qquad
\mbox{Partial-} \eta_B^2 = \frac{SSB}{SSA+SSE}   \qquad
\mbox{Partial-} \eta_{AB}^2 = \frac{SSAB}{SSAB+SSE}   \qquad
\]

\textbf{Example 5.2 - Drumstick Weights of Broiler Chickens}

A study was conducted to compare weights of chicken parts under 4 diets
(\citeproc{ref-aksu_2007}{Aksu, Imik, and Karoglu 2007}). The 4 diets
were combinations of two factors, each with two levels \((a=b=2)\).
Factor A was base diet (Sorghum \((i=1)\) and Corn \((i=2)\)). Factor B
was methionine supplement (Absent \((j=1)\) and Present \((j=2)\)).
There were \(n=60\) chickens receiving each diet for a total of
\((n_T=abn=2(2)(60)=240)\).

The observed means and standard deviations for the diets are given
below. We will use five decimal places to reduce rounding differences
with the R output

\[
\overline{y}_{11\bullet}=106.07983 \quad 
\overline{y}_{12\bullet}=93.67017 \quad 
\overline{y}_{21\bullet}=101.16983 \quad 
\overline{y}_{22\bullet}=108.82967
\]

\[
s_{11}=15.04009 \quad s_{12}=12.28969 \quad s_{21}=16.74029 \quad s_{22}=20.92978
\] Here we compute the marginal means for each level of factors A and B,
as well as the overall mean, and then obtain least squares estimates and
the Analysis of Variance.

\[
\overline{y}_{1\bullet\bullet}=\frac{106.07983+93.67017}{2}=99.87500 \quad
\overline{y}_{2\bullet\bullet}=\frac{101.16983+108.82967}{2}=104.99975 
\] \[
\overline{y}_{\bullet 1 \bullet}=\frac{106.07983+101.16983}{2}=103.62483 \quad
\overline{y}_{\bullet 2 \bullet}=\frac{93.67017+108.82967}{2}=101.24992 
\] \[
\overline{y}_{\bullet\bullet\bullet}=\frac{106.07983+93.67017+101.16983+108.82967}{4}=
102.437375
\] The main effect and interacion effects are computed from the means
above.

\[
\hat{\alpha}_1=99.87500-102.437375=-2.562375 \quad
\hat{\alpha}_2=104.99975-102.437375=2.562375 
\]

\[
\hat{\beta}_1=103.62483-102.437375=1.187455 \quad
\hat{\beta}_2=101.24992-102.4375=-1.187455
\]

\[
\hat{(\alpha\beta)}_{11}=106.07983-99.87500-103.62483+102.437375=5.017375
\]

\[
\hat{(\alpha\beta)}_{12}=93.67017-99.87500-101.24992+102.437375=-5.017375
\]

\[
\hat{(\alpha\beta)}_{21}=101.16983-104.99975-103.62483+102.437375=-5.017375
\]

\[
\hat{(\alpha\beta)}_{22}=108.82967-104.99975-101.24992+102.437375=5.017375
\]

Next, we compute the sums of squares, making use of the estimated
effects for \(SSA\), \(SSB\), and \(SSAB\).

\[
SSA=bn\sum_{i=1}^a\hat{\alpha}_i^2 = 
2(60)\left[(-2.562375)^2+(2.562375)^2\right]=1575.78375
\] \[
df_A=2-1 \qquad
MSA=1575.78375
\]

\[
SSB=an\sum_{j=1}^b\hat{\beta}_j^2 = 
2(60)\left[(1.187455)^2+(-1.187455)^2\right]=338.41185
\] \[
df_B=2-1 \qquad
MSB=338.41185
\] \[
SSAB=n\sum_{i=1}^a\sum_{j=1}^b\hat{(\alpha\beta)}_{ij}^2 =
\] \[
60\left[(5.017375)^2+(-5.017375)^2+(-5.017375)^2+(5.017375)^2\right]
=6041.77245
\] \[
df_{AB}=(2-1)(2-1)=1 \qquad MSAB=6041.77245
\]

\[
SSE=(n-1)\sum_{i=1}^a\sum_{j=1}^bs_{ij}^2=
59\left[(15.04009)^2+(12.28969)^2+(16.74029)^2+(20.92978)^2\right]=64636.49347
\]

\[
df_E=2(2)(60-1)=236 \qquad MSE=273.88345
\]

First, test for an interaction between base diet and methionine
supplement.

\[
H_0^{AB}: (\alpha\beta)_{11} = \cdots = (\alpha\beta)_{22} = 0 
\]

\[
\mbox{Test Statistic:} \quad F_{AB}^*=\frac{MSAB}{MSE}=\frac{6041.77245}{273.88345} = 22.060
\] \[
\mbox{Rejection Region:} \quad F_{AB}^* \ge F_{.95;1,236}=3.881
\qquad P=P\left(F_{1,236}\ge 22.060\right) < .0001
\]

As there is strong evidence of an interaction (adding methionine has a
large negative effect when sorghum is used, and has a large positive
effect when corn is used), we will conduct the main effects test for
completeness.

\[
H_0^A: \alpha_1 =  \alpha_2=0 \qquad
H_A^A: \mbox { Not all } \alpha_i=0
\]

\[
\mbox{Test Statistic:} \quad F_{A}^*=
\frac{MSA}{MSE}=\frac{1575.78375}{273.88345}=5.753
\]

\[
\mbox{Rejection Region:} \quad F_{A}^* \ge F_{.95;1,236}=3.881
\qquad P=P\left(F_{1,236}\ge 5.753\right)=.0172
\]

\[
H_0^B: \beta_1 = \beta_2=0 \qquad
H_A^B: \mbox { Not all } \beta_j=0
\]

\[
\mbox{Test Statistic:} \quad F_{B}^*=\frac{MSB}{MSE}=\frac{338.41185}{273.88345}
=1.236
\]

\[
\mbox{Rejection Region:} \quad F_{A}^* \ge F_{.95;1,236}=3.881
\qquad P=P\left(F_{1,236}\ge 1.236\right)=.2674
\]

The test for factor A is significant (Corn does better marginally than
Sorghum). The test for factor B is not significant (marginally
Methionine Absent and Present give similar results). This last result is
mis-leading as Methionine has a large positve effect for Corn (7.66 gram
increase) but a large negative effect for Sorghum (12.41 gram decrease),
so marginally they ``cancel out.''

The effect sizes and partial effect sizes for the base diet, methionine,
and their interaction are obtained below, along with the total sum of
squares.

\[
SSTO=SSA+SSB+SSAB+SSE=1575.94+338.44+6445.58+64636.75=72996.71
\]

\[
\eta_A^2=\frac{1575.94}{72996.71}=0.0216 \qquad
\eta_B^2=\frac{338.44}{72996.71}=0.0046 \qquad
\eta_{AB}^2=\frac{6445.58}{72996.71}=0.0883 
\] \[
\mbox{Partial-} \eta_A^2=\frac{1575.94}{1575.94+64636.75}=0.0238 \qquad
\mbox{Partial-} \eta_B^2=\frac{338.44}{338.44+64636.75}=0.0052 
\] \[
\mbox{Partial-} \eta_{AB}^2=\frac{6445.58}{6445.58+64636.75}=0.0907
\]

The interaction has a much larger effect size than either of the main
effects. Based on \(\eta_{AB}^2\), approximately 9\% of the total
variance is attributable to the interaction.

The following R code uses the \textbf{aov} function. As before, we set
the options to have the effects sum to zero.

\begin{Shaded}
\begin{Highlighting}[]
\NormalTok{broiler }\OtherTok{\textless{}{-}} \FunctionTok{read.csv}\NormalTok{(}\StringTok{"data/drumstick.csv"}\NormalTok{)}
\FunctionTok{head}\NormalTok{(broiler,}\DecValTok{2}\NormalTok{)}
\end{Highlighting}
\end{Shaded}

\begin{verbatim}
##   diet base methionine drumStkWt
## 1    1    1          1    116.33
## 2    1    1          1     99.43
\end{verbatim}

\begin{Shaded}
\begin{Highlighting}[]
\NormalTok{broiler}\SpecialCharTok{$}\NormalTok{base.f }\OtherTok{\textless{}{-}} \FunctionTok{factor}\NormalTok{(broiler}\SpecialCharTok{$}\NormalTok{base, }\AttributeTok{levels=}\DecValTok{1}\SpecialCharTok{:}\DecValTok{2}\NormalTok{, }
                         \AttributeTok{labels=}\FunctionTok{c}\NormalTok{(}\StringTok{"Sorghum"}\NormalTok{, }\StringTok{"Corn"}\NormalTok{))}
\NormalTok{broiler}\SpecialCharTok{$}\NormalTok{meth.f }\OtherTok{\textless{}{-}} \FunctionTok{factor}\NormalTok{(broiler}\SpecialCharTok{$}\NormalTok{meth, }\AttributeTok{levels=}\DecValTok{1}\SpecialCharTok{:}\DecValTok{2}\NormalTok{, }
                         \AttributeTok{labels=}\FunctionTok{c}\NormalTok{(}\StringTok{"Absent"}\NormalTok{, }\StringTok{"Present"}\NormalTok{))}

\FunctionTok{interaction.plot}\NormalTok{(broiler}\SpecialCharTok{$}\NormalTok{base.f, broiler}\SpecialCharTok{$}\NormalTok{meth.f, broiler}\SpecialCharTok{$}\NormalTok{drumStkWt,}
                 \AttributeTok{xlab=}\StringTok{"Base Diet"}\NormalTok{, }\AttributeTok{ylab=}\StringTok{"Weight"}\NormalTok{,}
                 \AttributeTok{col=}\FunctionTok{c}\NormalTok{(}\StringTok{"blue"}\NormalTok{, }\StringTok{"red"}\NormalTok{), }\AttributeTok{lty=}\FunctionTok{c}\NormalTok{(}\DecValTok{1}\NormalTok{,}\DecValTok{1}\NormalTok{))}
\end{Highlighting}
\end{Shaded}

\pandocbounded{\includegraphics[keepaspectratio]{bookdownproj_STA4211Notes1_files/figure-latex/ch0505-1.pdf}}

\begin{Shaded}
\begin{Highlighting}[]
\FunctionTok{options}\NormalTok{(}\AttributeTok{contrasts=}\FunctionTok{c}\NormalTok{(}\StringTok{"contr.sum"}\NormalTok{,}\StringTok{"contr.poly"}\NormalTok{))}
\NormalTok{ds.mod1 }\OtherTok{\textless{}{-}} \FunctionTok{aov}\NormalTok{(drumStkWt }\SpecialCharTok{\textasciitilde{}}\NormalTok{ base.f}\SpecialCharTok{*}\NormalTok{meth.f, }\AttributeTok{data=}\NormalTok{broiler)}
\FunctionTok{summary.lm}\NormalTok{(ds.mod1)}
\end{Highlighting}
\end{Shaded}

\begin{verbatim}
## 
## Call:
## aov(formula = drumStkWt ~ base.f * meth.f, data = broiler)
## 
## Residuals:
##    Min     1Q Median     3Q    Max 
## -59.16 -10.94   0.78   9.58  57.21 
## 
## Coefficients:
##                 Estimate Std. Error t value Pr(>|t|)    
## (Intercept)      102.437      1.068  95.892  < 2e-16 ***
## base.f1           -2.562      1.068  -2.399   0.0172 *  
## meth.f1            1.187      1.068   1.112   0.2674    
## base.f1:meth.f1    5.017      1.068   4.697 4.49e-06 ***
## ---
## Signif. codes:  0 '***' 0.001 '**' 0.01 '*' 0.05 '.' 0.1 ' ' 1
## 
## Residual standard error: 16.55 on 236 degrees of freedom
## Multiple R-squared:  0.1096, Adjusted R-squared:  0.09828 
## F-statistic: 9.683 on 3 and 236 DF,  p-value: 4.731e-06
\end{verbatim}

\begin{Shaded}
\begin{Highlighting}[]
\FunctionTok{anova}\NormalTok{(ds.mod1)}
\end{Highlighting}
\end{Shaded}

\begin{verbatim}
## Analysis of Variance Table
## 
## Response: drumStkWt
##                Df Sum Sq Mean Sq F value    Pr(>F)    
## base.f          1   1576  1575.8  5.7535   0.01723 *  
## meth.f          1    338   338.4  1.2356   0.26745    
## base.f:meth.f   1   6042  6041.8 22.0596 4.488e-06 ***
## Residuals     236  64636   273.9                      
## ---
## Signif. codes:  0 '***' 0.001 '**' 0.01 '*' 0.05 '.' 0.1 ' ' 1
\end{verbatim}

\begin{Shaded}
\begin{Highlighting}[]
\FunctionTok{options}\NormalTok{(}\AttributeTok{es.use\_symbols=}\ConstantTok{TRUE}\NormalTok{)}
\FunctionTok{eta\_squared}\NormalTok{(ds.mod1)}
\end{Highlighting}
\end{Shaded}

\begin{verbatim}
## # Effect Size for ANOVA (Type I)
## 
## Parameter     | η² (partial) |       95% CI
## -------------------------------------------
## base.f        |         0.02 | [0.00, 1.00]
## meth.f        |     5.21e-03 | [0.00, 1.00]
## base.f:meth.f |         0.09 | [0.04, 1.00]
## 
## - One-sided CIs: upper bound fixed at [1.00].
\end{verbatim}

\begin{Shaded}
\begin{Highlighting}[]
\FunctionTok{eta\_squared}\NormalTok{(ds.mod1, }\AttributeTok{partial=}\ConstantTok{FALSE}\NormalTok{)}
\end{Highlighting}
\end{Shaded}

\begin{verbatim}
## # Effect Size for ANOVA (Type I)
## 
## Parameter     |       η² |       95% CI
## ---------------------------------------
## base.f        |     0.02 | [0.00, 1.00]
## meth.f        | 4.66e-03 | [0.00, 1.00]
## base.f:meth.f |     0.08 | [0.04, 1.00]
## 
## - One-sided CIs: upper bound fixed at [1.00].
\end{verbatim}

\[ \nabla \]

\textbf{Example 5.3 - Reading Times on Electronic Readers}

Next, consider a study that compared \(a=3\) electronic reader models,
each at \(b=4\) illuminance levels (\citeproc{ref-chang_2013}{Chang,
Chou, and Shieh 2013}). The three reader models are described as
follows.

\begin{itemize}
\tightlist
\item
  Sony PRS 700 - 6'' diagonal screen size \((i=1)\)
\item
  Amazon Kindle DX - 9.7'' diagonal screen size \((i=2)\)
\item
  iRex 1000S - 10.2'' diagonal screen size \((i=3)\)
\end{itemize}

The illuminance levels studied were 200 lx \((j=1)\), 500 lx \((j=2)\),
1000 lx \((j=3)\), and 1500 lx \((j=4)\). A total study group of
\(n_T=60\) subjects was obtained with \(n=5\) being measured on the
\(ab=3(4)=12\) combinations of device and illumination level. The
response variable \(Y\) is the reading time in seconds.

R code and output are given below.

\begin{Shaded}
\begin{Highlighting}[]
\NormalTok{ereader }\OtherTok{\textless{}{-}} \FunctionTok{read.csv}\NormalTok{(}\StringTok{"data/ereader1.csv"}\NormalTok{)}
\FunctionTok{head}\NormalTok{(ereader,}\DecValTok{2}\NormalTok{)}
\end{Highlighting}
\end{Shaded}

\begin{verbatim}
##   device illum readTime
## 1      1     1  1656.26
## 2      1     1  1405.92
\end{verbatim}

\begin{Shaded}
\begin{Highlighting}[]
\NormalTok{ereader}\SpecialCharTok{$}\NormalTok{device.f }\OtherTok{\textless{}{-}} \FunctionTok{factor}\NormalTok{(ereader}\SpecialCharTok{$}\NormalTok{device, }\AttributeTok{levels=}\DecValTok{1}\SpecialCharTok{:}\DecValTok{3}\NormalTok{,}
                           \AttributeTok{labels=}\FunctionTok{c}\NormalTok{(}\StringTok{"Sony"}\NormalTok{, }\StringTok{"Amaz"}\NormalTok{, }\StringTok{"iRex"}\NormalTok{))}
\NormalTok{ereader}\SpecialCharTok{$}\NormalTok{illum.f }\OtherTok{\textless{}{-}} \FunctionTok{factor}\NormalTok{(ereader}\SpecialCharTok{$}\NormalTok{illum, }\AttributeTok{levels=}\DecValTok{1}\SpecialCharTok{:}\DecValTok{4}\NormalTok{,}
                  \AttributeTok{labels=}\FunctionTok{c}\NormalTok{(}\StringTok{"200"}\NormalTok{, }\StringTok{"500"}\NormalTok{, }\StringTok{"1000"}\NormalTok{, }\StringTok{"1500"}\NormalTok{))}

\FunctionTok{interaction.plot}\NormalTok{(ereader}\SpecialCharTok{$}\NormalTok{illum.f, ereader}\SpecialCharTok{$}\NormalTok{device.f, ereader}\SpecialCharTok{$}\NormalTok{readTime,}
                 \AttributeTok{xlab=}\StringTok{"Illuminination"}\NormalTok{, }\AttributeTok{ylab=}\StringTok{"Reading Time"}\NormalTok{,}
                 \AttributeTok{col=}\FunctionTok{c}\NormalTok{(}\StringTok{"green3"}\NormalTok{, }\StringTok{"purple"}\NormalTok{, }\StringTok{"red"}\NormalTok{), }\AttributeTok{lty=}\FunctionTok{c}\NormalTok{(}\DecValTok{1}\NormalTok{,}\DecValTok{1}\NormalTok{,}\DecValTok{1}\NormalTok{))}
\end{Highlighting}
\end{Shaded}

\pandocbounded{\includegraphics[keepaspectratio]{bookdownproj_STA4211Notes1_files/figure-latex/ch0506-1.pdf}}

\begin{Shaded}
\begin{Highlighting}[]
\FunctionTok{tapply}\NormalTok{(ereader}\SpecialCharTok{$}\NormalTok{readTime, }\FunctionTok{list}\NormalTok{(ereader}\SpecialCharTok{$}\NormalTok{device.f, ereader}\SpecialCharTok{$}\NormalTok{illum.f), mean)}
\end{Highlighting}
\end{Shaded}

\begin{verbatim}
##           200      500     1000    1500
## Sony 1462.158 1385.960 1093.694 1069.49
## Amaz 1250.188 1096.594  914.006  868.01
## iRex 1208.350 1102.146  904.400  841.10
\end{verbatim}

\begin{Shaded}
\begin{Highlighting}[]
\FunctionTok{options}\NormalTok{(}\AttributeTok{contrasts=}\FunctionTok{c}\NormalTok{(}\StringTok{"contr.sum"}\NormalTok{, }\StringTok{"contr.poly"}\NormalTok{))}

\NormalTok{ereader.mod1 }\OtherTok{\textless{}{-}} \FunctionTok{aov}\NormalTok{(readTime }\SpecialCharTok{\textasciitilde{}}\NormalTok{ device.f}\SpecialCharTok{*}\NormalTok{illum.f, }\AttributeTok{data=}\NormalTok{ereader)}
\FunctionTok{summary.lm}\NormalTok{(ereader.mod1)}
\end{Highlighting}
\end{Shaded}

\begin{verbatim}
## 
## Call:
## aov(formula = readTime ~ device.f * illum.f, data = ereader)
## 
## Residuals:
##     Min      1Q  Median      3Q     Max 
## -497.41 -187.20  -20.66  206.88  463.53 
## 
## Coefficients:
##                    Estimate Std. Error t value Pr(>|t|)    
## (Intercept)        1099.675     35.601  30.889  < 2e-16 ***
## device.f1           153.151     50.347   3.042  0.00380 ** 
## device.f2           -67.475     50.347  -1.340  0.18649    
## illum.f1            207.224     61.663   3.361  0.00153 ** 
## illum.f2             95.225     61.663   1.544  0.12909    
## illum.f3           -128.975     61.663  -2.092  0.04179 *  
## device.f1:illum.f1    2.108     87.204   0.024  0.98081    
## device.f2:illum.f1   10.765     87.204   0.123  0.90227    
## device.f1:illum.f2   37.909     87.204   0.435  0.66572    
## device.f2:illum.f2  -30.831     87.204  -0.354  0.72523    
## device.f1:illum.f3  -30.157     87.204  -0.346  0.73099    
## device.f2:illum.f3   10.781     87.204   0.124  0.90212    
## ---
## Signif. codes:  0 '***' 0.001 '**' 0.01 '*' 0.05 '.' 0.1 ' ' 1
## 
## Residual standard error: 275.8 on 48 degrees of freedom
## Multiple R-squared:  0.3771, Adjusted R-squared:  0.2343 
## F-statistic: 2.641 on 11 and 48 DF,  p-value: 0.01001
\end{verbatim}

\begin{Shaded}
\begin{Highlighting}[]
\FunctionTok{anova}\NormalTok{(ereader.mod1)}
\end{Highlighting}
\end{Shaded}

\begin{verbatim}
## Analysis of Variance Table
## 
## Response: readTime
##                  Df  Sum Sq Mean Sq F value    Pr(>F)    
## device.f          2  706968  353484  4.6483 0.0142790 *  
## illum.f           3 1481064  493688  6.4920 0.0008906 ***
## device.f:illum.f  6   21543    3591  0.0472 0.9995253    
## Residuals        48 3650203   76046                      
## ---
## Signif. codes:  0 '***' 0.001 '**' 0.01 '*' 0.05 '.' 0.1 ' ' 1
\end{verbatim}

\begin{Shaded}
\begin{Highlighting}[]
\NormalTok{HSD.dev }\OtherTok{\textless{}{-}} \FunctionTok{TukeyHSD}\NormalTok{(ereader.mod1, }\StringTok{"device.f"}\NormalTok{)}
\NormalTok{HSD.ill }\OtherTok{\textless{}{-}} \FunctionTok{TukeyHSD}\NormalTok{(ereader.mod1, }\StringTok{"illum.f"}\NormalTok{)}
\NormalTok{HSD.dev}
\end{Highlighting}
\end{Shaded}

\begin{verbatim}
##   Tukey multiple comparisons of means
##     95% family-wise confidence level
## 
## Fit: aov(formula = readTime ~ device.f * illum.f, data = ereader)
## 
## $device.f
##                diff       lwr        upr     p adj
## Amaz-Sony -220.6260 -431.5285  -9.723488 0.0384849
## iRex-Sony -238.8265 -449.7290 -27.923988 0.0230557
## iRex-Amaz  -18.2005 -229.1030 192.702012 0.9762840
\end{verbatim}

\begin{Shaded}
\begin{Highlighting}[]
\NormalTok{HSD.ill}
\end{Highlighting}
\end{Shaded}

\begin{verbatim}
##   Tukey multiple comparisons of means
##     95% family-wise confidence level
## 
## Fit: aov(formula = readTime ~ device.f * illum.f, data = ereader)
## 
## $illum.f
##                diff       lwr          upr     p adj
## 500-200   -111.9987 -379.9852  155.9878183 0.6838249
## 1000-200  -336.1987 -604.1852  -68.2121817 0.0085849
## 1500-200  -380.6987 -648.6852 -112.7121817 0.0023697
## 1000-500  -224.2000 -492.1865   43.7864849 0.1306514
## 1500-500  -268.7000 -536.6865   -0.7135151 0.0491639
## 1500-1000  -44.5000 -312.4865  223.4864849 0.9708523
\end{verbatim}

\begin{Shaded}
\begin{Highlighting}[]
\FunctionTok{plot}\NormalTok{(HSD.dev, }\AttributeTok{las=}\DecValTok{0}\NormalTok{, }\AttributeTok{cex.axis=}\FloatTok{0.75}\NormalTok{, }\AttributeTok{col=}\StringTok{"red"}\NormalTok{)}
\end{Highlighting}
\end{Shaded}

\pandocbounded{\includegraphics[keepaspectratio]{bookdownproj_STA4211Notes1_files/figure-latex/ch0506-2.pdf}}

\begin{Shaded}
\begin{Highlighting}[]
\FunctionTok{plot}\NormalTok{(HSD.ill, }\AttributeTok{las=}\DecValTok{1}\NormalTok{, }\AttributeTok{cex.axis=}\FloatTok{0.60}\NormalTok{, }\AttributeTok{col=}\StringTok{"blue"}\NormalTok{)}
\end{Highlighting}
\end{Shaded}

\pandocbounded{\includegraphics[keepaspectratio]{bookdownproj_STA4211Notes1_files/figure-latex/ch0506-3.pdf}}

\begin{Shaded}
\begin{Highlighting}[]
\FunctionTok{options}\NormalTok{(}\AttributeTok{es.use\_symbols=}\ConstantTok{TRUE}\NormalTok{)}

\FunctionTok{eta\_squared}\NormalTok{(ereader.mod1)}
\end{Highlighting}
\end{Shaded}

\begin{verbatim}
## # Effect Size for ANOVA (Type I)
## 
## Parameter        | η² (partial) |       95% CI
## ----------------------------------------------
## device.f         |         0.16 | [0.02, 1.00]
## illum.f          |         0.29 | [0.10, 1.00]
## device.f:illum.f |     5.87e-03 | [0.00, 1.00]
## 
## - One-sided CIs: upper bound fixed at [1.00].
\end{verbatim}

\begin{Shaded}
\begin{Highlighting}[]
\FunctionTok{eta\_squared}\NormalTok{(ereader.mod1, }\AttributeTok{partial=}\ConstantTok{FALSE}\NormalTok{)}
\end{Highlighting}
\end{Shaded}

\begin{verbatim}
## # Effect Size for ANOVA (Type I)
## 
## Parameter        |       η² |       95% CI
## ------------------------------------------
## device.f         |     0.12 | [0.00, 1.00]
## illum.f          |     0.25 | [0.07, 1.00]
## device.f:illum.f | 3.68e-03 | [0.00, 1.00]
## 
## - One-sided CIs: upper bound fixed at [1.00].
\end{verbatim}

Unlike the chicken broilers example, there is no evidence (whatsoever)
of a significant interaction between device and illumination level
\(F_{AB}^*=0.0472\), \(P_{AB}=.9995\). There is evidence of device
differences \(F_A^*=4.648\), \(P_A=.0143\) and illumination differences
\(F_B^*=6.492\), \(P_B=.0009\).

The residual standard error, which estimates \(\sigma\), is 275.8
seconds, a little over 4.5 minutes. The variation in people's reading
times (within treatments) is quite large. Based on \(\eta^2\), the
illumination levels explains 25\% of the total variation in reading
times, while device explains 12\%.

Based on Tukey comparisons among devices and among illumination levels
we obtain the following significant differences. Formulas will be given
in the next section.

\begin{itemize}
\tightlist
\item
  The larger Amazon and iRex readers have significantly lower reading
  times than the smaller Sony model.
\item
  1000 lx and 1500 lx have significantly lower reading times than 200 lx
\item
  1500 lx has significantly lower reading times than 500 lx
\end{itemize}

\[ \nabla \]

Finally, some guidelines are given for testing and modelling strategy.

If the interaction is significant, and the goal was to demonstrate
interactions (as is often the case in behavioral studies), describe the
interaction effects in terms of cell means, as in the Chicken Diet
example. If the interaction is significant, and the goal is to simplify
the model, try some power (Box-Cox) transformations on \(Y\) and see if
that simplifies the model.

If the interaction isn't significant or important, test for main effects
for factors A and B. Also, make comparisons among the marginal means of
factors A and B, as in the E-reader example.

\section{Factor Effect Contrasts}\label{factor-effect-contrasts}

In this section, we consider contrasts among means when interaction is
absent and when it is present.

\subsection{Contrasts when there is No
Interaction}\label{contrasts-when-there-is-no-interaction}

This is the simpler case, where contrasts are made among levels of
factors A and B, as was done previously for single factor models.

For contrasts among levels for factor A, we summarize the results from
the 1-Way ANOVA. Keep in mind that the marginal means for the levels of
factor A are based on \(an\) levels.

\[
\mbox{Contrasts among Factor A Levels:} \quad
L=\sum_{i=1}^a c_i\mu_{i\bullet} \qquad
\sum_{i=1}^a c_i=0
\]

\[
\mbox{Estimator:} \quad
\hat{L}=\sum_{i=1}^a c_i\overline{Y}_{i\bullet\bullet} \qquad
\mbox{Standard Error:} \quad
\hat{SE}\left\{\hat{L}\right\}=\sqrt{\frac{MSE}{bn}\sum_{i=1}^a c_i^2}
\]

\[
(1-\alpha)100\% \mbox{ Confidence Interval for }L: \quad
\hat{L} \pm t_{1-\alpha/2; ab(n-1)} \hat{SE}\left\{\hat{L}\right\}
\]

\[
\mbox{Scheffe (all contrasts):}\quad
\hat{L} \pm \left(\sqrt{(a-1)F_{1-\alpha;a-1,ab(n-1)}}\right)
\hat{SE}\left\{\hat{L}\right\}
\]

\[
\mbox{Bonferroni (} g \mbox{ pre-planned) contrasts:} \quad
\hat{L} \pm t_{1-\alpha/(2g); ab(n-1)} \hat{SE}\left\{\hat{L}\right\}
\] \[
\mbox{Tukey (All pairs):} \quad
\hat{L} \pm \frac{q_{1-\alpha; a, ab(n-1)}}{\sqrt{2}} \hat{SE}\left\{\hat{L}\right\}
\] Similar results for factor B occur, with the following adjustments,
as well as critical values for Scheffe's and Tukey's methods.

\[
L=\sum_{j=1}^b c_j\mu_{\bullet j} \qquad
\sum_{j=1}^b c_j=0  \qquad
\hat{L}=\sum_{j=1}^b c_j\overline{Y}_{\bullet j \bullet} \qquad
\hat{SE}\left\{\hat{L}\right\}=\sqrt{\frac{MSE}{an}\sum_{j=1}^b c_j^2}
\]

\subsection{Contrasts when Interaction is
Present}\label{contrasts-when-interaction-is-present}

In this case, contrasts are made among cell means as opposed to marginal
means.

\[
\mbox{Contrasts:} \quad L=\sum_{i=1}^a\sum_{j=1}^bc_{ij}\mu_{ij} \qquad
\sum_{i=1}^a\sum_{j=1}^bc_{ij}=0
\]

\[
\mbox{Estimator:} \quad
\hat{L}=\sum_{i=1}^a\sum_{j=1}^b c_{ij}\overline{Y}_{ij\bullet} \qquad
\mbox{Std Error:} \quad
\hat{SE}\left\{\hat{L}\right\}=
\sqrt{\frac{MSE}{n}\sum_{i=1}^a\sum_{j=1}^b c_{ij}^2}
\]

\[
(1-\alpha)100\% \mbox{ Confidence Interval for }L: \quad
\hat{L} \pm t_{1-\alpha/2; ab(n-1)} \hat{SE}\left\{\hat{L}\right\}
\]

\[
\mbox{Scheffe (all contrasts):}\quad
\hat{L} \pm \left(\sqrt{(ab-1)F_{1-\alpha;ab-1,ab(n-1)}}\right)
\hat{SE}\left\{\hat{L}\right\}
\]

\[
\mbox{Bonferroni (g pre-planned) contrasts:} \quad
\hat{L} \pm t_{1-\alpha/(2g); ab(n-1)} \hat{SE}\left\{\hat{L}\right\}
\] \[
\mbox{Tukey (All pairs):} \quad
\hat{L} \pm \frac{q_{1-\alpha; ab, ab(n-1)}}{\sqrt{2}} \hat{SE}\left\{\hat{L}\right\}
\]

\begin{verbatim}
## Warning: package 'additivityTests' was built under R version 4.5.3
\end{verbatim}

\chapter{Two-Factor Designs with 1 Observation per
Treatment}\label{two-factor-designs-with-1-observation-per-treatment}

In some cases, there is only a single observation within each
combination of factor levels. This causes a problem when considering
models with interaction effects. The error degrees of freedom are
\(df_E=ab(n-1)=ab(1-1)=0\) and the error sum of squares is \(SSE=0\).
Several methods have been developed for testing for interaction under
particular restrictions. In this chapter, we describe \textbf{Tukey's
One Degree of Freedom for Non-Additivity Test} (ODOFNA). For this test,
the form of the interaction and the resulting test are given below.

\[
\left(\alpha\beta\right)_{ij} = D\alpha_i\beta_j \qquad \qquad H_0:D=0 \quad H_A:D\ne 0
\]

Intuitively, the test involves estimating \(\mu\), as well as
\(\alpha_i\) and \(\beta_j\) from the marginal means for factors A and
B. Then \(D\) is estimated by fitting a regression through the origin
relating \(Y^*\) (\(Y\) minus the sum of the mean and the main effects)
to \(X^*\) where \(X^*\) is the product of the estimates of
\(\alpha_i\), and \(\beta_j\).

\[
\hat{\mu}=\overline{Y}_{\bullet\bullet}\qquad
\hat{\alpha}_i = \overline{Y}_{i\bullet} - \overline{Y}_{\bullet\bullet}\qquad
\hat{\beta}_j = \overline{Y}_{\bullet j} - \overline{Y}_{\bullet\bullet} \qquad
Y_{ij}^*=Y_{ij}-\hat{\mu}-\hat{\alpha}_i-\hat{\beta}_j \qquad
 X_{ij}^*=\hat{\alpha}_i\hat{\beta}_j
\]

For regression through the origin, we obtain the estimate and test
whether \(D=0\) as follows, where the model with \(D\) is the complete
model.

\[
\hat{D}=\frac{\sum_{i=1}^a\sum_{j=1}^bX_{ij}^*Y_{ij}^*}{\sum_{i=1}^a\sum_{j=1}^bX_{ij}^{*2}}
\qquad  \hat{Y}_{ij}(C)=\hat{\mu}_{\bullet\bullet}+\hat{\alpha}_i + \hat{\beta}_j +
\hat{D}\hat{\alpha}_i\hat{\beta}_j \qquad i=1,\ldots,a; \quad j=1,\ldots,b
\] \[
SSE(C)=\sum_{i=1}^a\sum_{j=1}^b\left(Y_{ij}-\hat{Y}_{ij}(C)\right)^2  \qquad
df_E(C)=ab-1-(a-1)-(b-1)-1=ab-a-b
\] Under the reduced model, \(D=0\), and we obtain the following fitted
values and error sum of squares and degrees of freedom.

\[
\hat{Y}_{ij}(R)=\hat{\mu} + \hat{\alpha}_i + \hat{\beta}_j
\qquad SSE(R)=\sum_{i=1}^a\sum_{j=1}^b\left(Y_{ij}-\hat{Y}_{ij}(R)\right)^2
\] \[ df_E(R)=ab-1-(a-1)-(b-1)=ab-a-b+1=(a-1)(b-1)
\]

Then Tukey's ODOFNA test is a special case of the general linear test
with \(H_0:D=0\) vs \(H_A:D \ne 0\).

\[
\mbox{Test Statistic:} \quad
F^*=\frac{\left[\frac{SSE(R)-SSE(C)}{df_E(R)-df_E(C)}\right]}
         {\left[\frac{SSE(C)}{df_E(C)}\right]}
\qquad
\mbox{Rejection Region:} \quad F^*\ge F_{1-\alpha;1,ab-a-b}
\]

\textbf{Example 6.1 - Economic Indices for 8 Sources Over 18 Years}

A study considered U.S. economic indices for \(a=8\) business news
sources over a \(b=18\) year period (\citeproc{ref-smith_1969}{K. V.
Smith 1969}). The sources were: DJIA, Poors, NYSE, GNP, CPI, Forbes,
Business Week, and Money Magazine. The years were 1948-1965. There was
one index per source per year. The following R code runs calculations
directly, then structures the data in matrix form and uses the
\textbf{tukey.test} function in the \textbf{additivityTests} package.
Note that we order the data frame first by source, then by year within
source, as \textbf{tapply} output is sorted by the alphanumeric levels
of the grouping variable(s).

\begin{Shaded}
\begin{Highlighting}[]
\NormalTok{biz1 }\OtherTok{\textless{}{-}} \FunctionTok{read.csv}\NormalTok{(}\StringTok{"data/jb42.csv"}\NormalTok{)}
\FunctionTok{head}\NormalTok{(biz1,}\DecValTok{2}\NormalTok{)}
\end{Highlighting}
\end{Shaded}

\begin{verbatim}
##   source year index
## 1   DJIA 1965 1.103
## 2   DJIA 1964 1.145
\end{verbatim}

\begin{Shaded}
\begin{Highlighting}[]
\DocumentationTok{\#\# Create a new data frame that orders first by source then year within source}
\NormalTok{biz }\OtherTok{\textless{}{-}}\NormalTok{ biz1[}\FunctionTok{order}\NormalTok{(biz1}\SpecialCharTok{$}\NormalTok{source, biz1}\SpecialCharTok{$}\NormalTok{year),]}
\FunctionTok{head}\NormalTok{(biz,}\DecValTok{2}\NormalTok{)}
\end{Highlighting}
\end{Shaded}

\begin{verbatim}
##     source year index
## 126  BWEEK 1948 1.019
## 125  BWEEK 1949 0.995
\end{verbatim}

\begin{Shaded}
\begin{Highlighting}[]
\FunctionTok{tail}\NormalTok{(biz,}\DecValTok{2}\NormalTok{)}
\end{Highlighting}
\end{Shaded}

\begin{verbatim}
##    source year index
## 20   POOR 1964 1.131
## 19   POOR 1965 1.099
\end{verbatim}

\begin{Shaded}
\begin{Highlighting}[]
\NormalTok{all.mean }\OtherTok{\textless{}{-}} \FunctionTok{mean}\NormalTok{(biz}\SpecialCharTok{$}\NormalTok{index)}
\NormalTok{source.mean }\OtherTok{\textless{}{-}} \FunctionTok{as.vector}\NormalTok{(}\FunctionTok{tapply}\NormalTok{(biz}\SpecialCharTok{$}\NormalTok{index, biz}\SpecialCharTok{$}\NormalTok{source, mean))}
\NormalTok{year.mean }\OtherTok{\textless{}{-}} \FunctionTok{as.vector}\NormalTok{(}\FunctionTok{tapply}\NormalTok{(biz}\SpecialCharTok{$}\NormalTok{index, biz}\SpecialCharTok{$}\NormalTok{year, mean))}
\NormalTok{a }\OtherTok{\textless{}{-}} \FunctionTok{length}\NormalTok{(source.mean)}
\NormalTok{b }\OtherTok{\textless{}{-}} \FunctionTok{length}\NormalTok{(year.mean)}

\NormalTok{biz}\SpecialCharTok{$}\NormalTok{mu }\OtherTok{\textless{}{-}} \FunctionTok{rep}\NormalTok{(all.mean, a}\SpecialCharTok{*}\NormalTok{b)}
\NormalTok{biz}\SpecialCharTok{$}\NormalTok{alpha }\OtherTok{\textless{}{-}} \FunctionTok{rep}\NormalTok{(source.mean}\SpecialCharTok{{-}}\NormalTok{all.mean, }\AttributeTok{each=}\NormalTok{b)}
\NormalTok{biz}\SpecialCharTok{$}\NormalTok{beta }\OtherTok{\textless{}{-}} \FunctionTok{rep}\NormalTok{(year.mean}\SpecialCharTok{{-}}\NormalTok{all.mean, }\AttributeTok{times=}\NormalTok{a)}

\NormalTok{biz}\SpecialCharTok{$}\NormalTok{Ystar }\OtherTok{\textless{}{-}}\NormalTok{ biz}\SpecialCharTok{$}\NormalTok{index }\SpecialCharTok{{-}}\NormalTok{ biz}\SpecialCharTok{$}\NormalTok{mu }\SpecialCharTok{{-}}\NormalTok{ biz}\SpecialCharTok{$}\NormalTok{alpha }\SpecialCharTok{{-}}\NormalTok{ biz}\SpecialCharTok{$}\NormalTok{beta}
\NormalTok{biz}\SpecialCharTok{$}\NormalTok{Xstar }\OtherTok{\textless{}{-}}\NormalTok{ biz}\SpecialCharTok{$}\NormalTok{alpha }\SpecialCharTok{*}\NormalTok{ biz}\SpecialCharTok{$}\NormalTok{beta}

\NormalTok{(Dhat }\OtherTok{\textless{}{-}} \FunctionTok{sum}\NormalTok{(biz}\SpecialCharTok{$}\NormalTok{Xstar}\SpecialCharTok{*}\NormalTok{biz}\SpecialCharTok{$}\NormalTok{Ystar) }\SpecialCharTok{/} \FunctionTok{sum}\NormalTok{(biz}\SpecialCharTok{$}\NormalTok{Xstar}\SpecialCharTok{\^{}}\DecValTok{2}\NormalTok{))}
\end{Highlighting}
\end{Shaded}

\begin{verbatim}
## [1] 23.6517
\end{verbatim}

\begin{Shaded}
\begin{Highlighting}[]
\NormalTok{biz}\SpecialCharTok{$}\NormalTok{Yhat.C }\OtherTok{\textless{}{-}}\NormalTok{ biz}\SpecialCharTok{$}\NormalTok{mu }\SpecialCharTok{+}\NormalTok{ biz}\SpecialCharTok{$}\NormalTok{alpha }\SpecialCharTok{+}\NormalTok{ biz}\SpecialCharTok{$}\NormalTok{beta }\SpecialCharTok{+}\NormalTok{ Dhat}\SpecialCharTok{*}\NormalTok{biz}\SpecialCharTok{$}\NormalTok{Xstar}
\NormalTok{biz}\SpecialCharTok{$}\NormalTok{Yhat.R }\OtherTok{\textless{}{-}}\NormalTok{ biz}\SpecialCharTok{$}\NormalTok{mu }\SpecialCharTok{+}\NormalTok{ biz}\SpecialCharTok{$}\NormalTok{alpha }\SpecialCharTok{+}\NormalTok{ biz}\SpecialCharTok{$}\NormalTok{beta}

\NormalTok{SSE.C }\OtherTok{\textless{}{-}} \FunctionTok{sum}\NormalTok{((biz}\SpecialCharTok{$}\NormalTok{index}\SpecialCharTok{{-}}\NormalTok{biz}\SpecialCharTok{$}\NormalTok{Yhat.C)}\SpecialCharTok{\^{}}\DecValTok{2}\NormalTok{)}
\NormalTok{SSE.R }\OtherTok{\textless{}{-}} \FunctionTok{sum}\NormalTok{((biz}\SpecialCharTok{$}\NormalTok{index}\SpecialCharTok{{-}}\NormalTok{biz}\SpecialCharTok{$}\NormalTok{Yhat.R)}\SpecialCharTok{\^{}}\DecValTok{2}\NormalTok{)}
\NormalTok{df\_E.C }\OtherTok{\textless{}{-}}\NormalTok{ a}\SpecialCharTok{*}\NormalTok{b }\SpecialCharTok{{-}}\NormalTok{ a }\SpecialCharTok{{-}}\NormalTok{ b}
\NormalTok{df\_E.R }\OtherTok{\textless{}{-}}\NormalTok{ a}\SpecialCharTok{*}\NormalTok{b }\SpecialCharTok{{-}}\NormalTok{ a }\SpecialCharTok{{-}}\NormalTok{ b }\SpecialCharTok{+} \DecValTok{1}

\NormalTok{Fstar }\OtherTok{\textless{}{-}}\NormalTok{ ((SSE.R}\SpecialCharTok{{-}}\NormalTok{SSE.C)}\SpecialCharTok{/}\NormalTok{(df\_E.R}\SpecialCharTok{{-}}\NormalTok{df\_E.C)) }\SpecialCharTok{/}\NormalTok{ (SSE.C}\SpecialCharTok{/}\NormalTok{df\_E.C)}

\NormalTok{odofna.out }\OtherTok{\textless{}{-}} \FunctionTok{cbind}\NormalTok{(df\_E.R, df\_E.C, SSE.R, SSE.C, Fstar, }\FunctionTok{qf}\NormalTok{(.}\DecValTok{95}\NormalTok{,}\DecValTok{1}\NormalTok{,df\_E.C),}
                    \DecValTok{1}\SpecialCharTok{{-}}\FunctionTok{pf}\NormalTok{(Fstar,}\DecValTok{1}\NormalTok{,df\_E.C))}
\FunctionTok{colnames}\NormalTok{(odofna.out) }\OtherTok{\textless{}{-}} \FunctionTok{c}\NormalTok{(}\StringTok{"df(R)"}\NormalTok{, }\StringTok{"df(C)"}\NormalTok{, }\StringTok{"SSE(R)"}\NormalTok{, }\StringTok{"SSE(C)"}\NormalTok{, }\StringTok{"F*"}\NormalTok{, }\StringTok{"F(.95)"}\NormalTok{, }\StringTok{"P(\textgreater{}F*)"}\NormalTok{)}
\FunctionTok{rownames}\NormalTok{(odofna.out) }\OtherTok{\textless{}{-}} \FunctionTok{c}\NormalTok{(}\StringTok{"ODOFNA"}\NormalTok{)}
\FunctionTok{round}\NormalTok{(odofna.out, }\DecValTok{4}\NormalTok{)}
\end{Highlighting}
\end{Shaded}

\begin{verbatim}
##        df(R) df(C) SSE(R) SSE(C)       F* F(.95) P(>F*)
## ODOFNA   119   118 0.9305 0.4075 151.4522 3.9215      0
\end{verbatim}

\begin{Shaded}
\begin{Highlighting}[]
\FunctionTok{plot}\NormalTok{(biz}\SpecialCharTok{$}\NormalTok{Ystar }\SpecialCharTok{\textasciitilde{}}\NormalTok{ biz}\SpecialCharTok{$}\NormalTok{Xstar, }\AttributeTok{pch=}\DecValTok{16}\NormalTok{, }
     \AttributeTok{xlab=}\FunctionTok{expression}\NormalTok{(}\FunctionTok{paste}\NormalTok{(}\StringTok{"X\textquotesingle{} = "}\NormalTok{, }\FunctionTok{hat}\NormalTok{(alpha)}\SpecialCharTok{*}\FunctionTok{hat}\NormalTok{(beta))), }
     \AttributeTok{ylab=}\FunctionTok{expression}\NormalTok{(}\FunctionTok{paste}\NormalTok{(}\StringTok{"Y\textquotesingle{} = "}\NormalTok{, Y}\SpecialCharTok{{-}}\FunctionTok{hat}\NormalTok{(mu)}\SpecialCharTok{{-}}\FunctionTok{hat}\NormalTok{(alpha)}\SpecialCharTok{{-}}\FunctionTok{hat}\NormalTok{(beta))))}
\FunctionTok{abline}\NormalTok{(}\FunctionTok{lm}\NormalTok{(Ystar }\SpecialCharTok{\textasciitilde{}}\NormalTok{ Xstar }\SpecialCharTok{{-}} \DecValTok{1}\NormalTok{,}\AttributeTok{data=}\NormalTok{biz), }\AttributeTok{col=}\StringTok{"blue3"}\NormalTok{, }\AttributeTok{lwd=}\FloatTok{1.5}\NormalTok{)}
\end{Highlighting}
\end{Shaded}

\pandocbounded{\includegraphics[keepaspectratio]{bookdownproj_STA4211Notes1_files/figure-latex/ch.06.01-1.pdf}}

\begin{Shaded}
\begin{Highlighting}[]
\DocumentationTok{\#\#\# biz is sorted by index, then year (stacked COLUMNS from spreadsheet)}
\NormalTok{(Y.mat }\OtherTok{\textless{}{-}} \FunctionTok{matrix}\NormalTok{(biz}\SpecialCharTok{$}\NormalTok{index,}\AttributeTok{byrow=}\NormalTok{F,}\AttributeTok{ncol=}\DecValTok{8}\NormalTok{))}
\end{Highlighting}
\end{Shaded}

\begin{verbatim}
##        [,1]  [,2]  [,3]  [,4]  [,5]  [,6]  [,7]  [,8]
##  [1,] 1.019 1.077 0.972 1.041 1.110 0.985 0.972 0.996
##  [2,] 0.995 0.990 1.114 0.945 0.996 0.994 1.102 1.091
##  [3,] 1.213 1.009 1.179 1.157 1.102 1.057 1.211 1.247
##  [4,] 1.012 1.080 1.174 1.085 1.156 1.055 1.132 1.178
##  [5,] 1.141 1.022 1.074 1.037 1.055 1.038 1.065 1.109
##  [6,] 0.952 1.008 0.965 1.083 1.053 1.010 0.938 0.925
##  [7,] 1.074 1.004 1.393 0.940 0.994 1.029 1.426 1.497
##  [8,] 1.125 0.997 1.231 1.126 1.095 1.020 1.222 1.301
##  [9,] 1.011 1.015 1.000 1.034 1.055 1.012 1.026 1.034
## [10,] 0.906 1.035 0.833 1.008 1.056 0.993 0.866 0.856
## [11,] 1.099 1.028 1.425 0.930 1.044 1.038 1.366 1.376
## [12,] 1.073 1.008 1.184 1.127 1.086 1.006 1.097 1.094
## [13,] 0.917 1.016 0.896 1.029 1.041 0.924 0.976 0.953
## [14,] 1.154 1.011 1.207 1.099 1.032 1.031 1.240 1.231
## [15,] 1.018 1.012 0.890 1.078 1.072 1.013 0.880 0.872
## [16,] 1.060 1.012 1.169 1.051 1.050 1.038 1.180 1.201
## [17,] 1.073 1.013 1.145 1.064 1.066 1.043 1.143 1.131
## [18,] 1.093 1.017 1.103 1.083 1.086 1.048 1.095 1.099
\end{verbatim}

\begin{Shaded}
\begin{Highlighting}[]
\FunctionTok{tukey.test}\NormalTok{(Y.mat)}
\end{Highlighting}
\end{Shaded}

\begin{verbatim}
## 
## Tukey test on 5% alpha-level:
## 
## Test statistic: 151.5 
## Critival value: 3.921 
## The additivity hypothesis was rejected.
\end{verbatim}

There is strong evidence of an interaction between source and year, of
this form.

\[
(\alpha\beta)_{ij} = D\alpha_i\beta_j \quad \mbox{ with } \hat{D}=23.65
\]

\[ \nabla  \]

\begin{verbatim}
## Warning: package 'emmeans' was built under R version 4.5.3
\end{verbatim}

\begin{verbatim}
## Welcome to emmeans.
## Caution: You lose important information if you filter this package's results.
## See '? untidy'
\end{verbatim}

\chapter{Unbalanced Two-Factor Analysis of
Variance}\label{unbalanced-two-factor-analysis-of-variance}

When sample sizes are not all equal, the sums of squares cannot be
obtained simply among the cell, marginal, and overall means. Even in
well planned controlled experiments, observations may not be used due to
malfunction or subjects dropping out. In observational studies, it may
not be feasible to get equal number of units from the various
sub-populations.

The model can be fit based on regression models in scalar and matrix
form. In this chapter, we describe the process based on the treatment
effects model. The balanced case can be formed in this manner as well.

\section{Statistical Model and the Analysis of
Variance}\label{statistical-model-and-the-analysis-of-variance}

\[
Y_{ijk}=\mu_{\bullet\bullet}+\alpha_i+\beta_j+(\alpha\beta)_{ij}+\epsilon_{ijk}\qquad
i=1,\ldots,a; \quad j=1,\ldots,b; \quad k=1,\ldots,n_{ij}\]
\[\epsilon_{ijk} \sim NID\left(0,\sigma^2\right) 
\]

\[
n_{i\bullet} = \sum_{j=1}^b n_{ij}  \quad
n_{\bullet j} = \sum_{i=1}^a n_{ij}  \quad
n_T=\sum_{i=1}^a\sum_{j=1}^b n_{ij}  \qquad
Y_{ij\bullet}=\sum_{k=1}^{n_{ij}}Y_{ijk} \quad
\overline{Y}_{ij\bullet} = \frac{Y_{ij\bullet}}{n_{ij}}
\]

\[
\sum_{i=1}^a\alpha_i = \sum_{j=1}^b\beta_j = \sum_{i=1}^a(\alpha\beta)_{ij}
= \sum_{j=1}^b(\alpha\beta)_{ij} = 0
\] \[
\alpha_a=-\sum_{i=1}^{a-1}\alpha_i \qquad\qquad
\beta_b = -\sum_{j=1}^{b-1}\beta_j 
\] \[
(\alpha\beta)_{aj}=-\sum_{i=1}^{a-1}(\alpha\beta)_{ij} \quad j=1,\ldots,b
\qquad 
(\alpha\beta)_{ib}=-\sum_{j=1}^{b-1}(\alpha\beta)_{ij} \quad i=1,\ldots,a
\]

To fit this model, construct \(a-1\) \(X\) variables for the factor A
effects, \(b-1\) \(X\) variables for the factor B effects, and then
obtain the \((a-1)(b-1)\) cross-products of these variables for the
interaction effects.

\[
X_{ijk1}^A=\left\{ \begin{array}{r@{\quad:\quad}l} 
    1 & i=1 \\ 0 & i=2,\ldots,a-1\\ -1 & i=a 
    \end{array}  \right.
    \qquad \cdots \qquad
X_{ijk,a-1}^A=\left\{ \begin{array}{r@{\quad:\quad}l} 
    1 & i=a-1 \\ 0 & i=1,\ldots,a-2\\ -1 & i=a 
    \end{array}  \right.    
\]

\[
X_{ijk1}^B=\left\{ \begin{array}{r@{\quad:\quad}l} 
    1 & j=1 \\ 0 & j=2,\ldots,b-1\\ -1 & j=b 
    \end{array}  \right.
    \qquad \cdots \qquad
X_{ijk,b-1}^B=\left\{ \begin{array}{r@{\quad:\quad}l} 
    1 & j=b-1 \\ 0 & j=1,\ldots,b-2\\ -1 & j=b 
    \end{array}  \right.    
\] \[
Y_{ijk}=\mu_{\bullet\bullet} + \alpha_1X_{ijk1}^A + \cdots +\alpha_{a-1}X_{ijk,a-1}^A +
\beta_1X_{ijk1}^B + \cdots + \beta_{b-1}X_{ijk,b-1}^B + 
\] \[
+ (\alpha\beta)_{11}X_{ijk1}^AX_{ijk1}^B + \cdots +
(\alpha\beta)_{a-1,b-1}X_{ijk,a-1}^AX_{ijk,b-1}^B + \epsilon_{ijk}
\]

The testing strategy is done as it was in the balanced case. Fit various
models and compare them using the general linear test procedure.

\begin{itemize}
\tightlist
\item
  \textbf{Model 1} - Include main effects for factors A and B as well as
  interaction effects.
\item
  \textbf{Model 2} - Include main effects for factors A and B, but no
  interaction effects
\item
  \textbf{Model 3} - Include main effects for factor B and interaction
  effects, but no factor A effects
\item
  \textbf{Model 4} - Include main effects for factor A and interaction
  effects, but no factor B effects
\end{itemize}

To test for interaction effects, controlling for main effects for
factors A and B, we compare Model 1 (Complete) and Model 2 (Reduced).

To test for factor A effects, controlling for factor B effects and
interaction effects, we compare Model 1 (Complete) and Model 3
(Reduced).

To test for factor B effects, controlling for factor A effects and
interaction effects, we compare Model 1 (Complete) and Model 4
(Reduced).

When the test for interaction is ``very non-significant'' (high
\(P\)-value), often tests for factor A and B effects use Model 2 as the
Complete Model.

\textbf{Example 7.1 - Age at Peak for Famous Writers}

A study considered the age at career peak for famous writers
(\citeproc{ref-simonson_2007}{Simonson 2007}). The factors were
\textbf{Style} (Factor A: Conceptualists \((i=1)\), Experimentalists
\((i=2)\)) and Writer \textbf{Type} (Factor B:Poets \((j=1)\), Novelists
\((j=2)\)). Thus, there are \(a=b=2\) levels for the factors, there will
only be \(a-1=1\) \(X\) variable for factor A and \(b-1=1\) \(X\)
variable for factor B, and \((a-1)(b-1)=1\) cross-product term for the
AB interaction. The numbers of writers for the \(ab=4\) conditions are
\(n_{11}=n_{12}=5,n_{21}=6,n_{22}=7\). The data and the generated X
variables are given in Table @ref(tab:ch0701). The regression model is
fit directly, then with the \textbf{aov} function. Note that in the
original dataset the levels were 0 and 1 for style and type, not 1 and
2.

\begin{Shaded}
\begin{Highlighting}[]
\NormalTok{ap1 }\OtherTok{\textless{}{-}} \FunctionTok{read.csv}\NormalTok{(}\StringTok{"data/creativelifecycles1.csv"}\NormalTok{)}
\NormalTok{knitr}\SpecialCharTok{::}\FunctionTok{kable}\NormalTok{(ap1, }\AttributeTok{caption =} \StringTok{\textquotesingle{}Author Age at Peak Data\textquotesingle{}}\NormalTok{)}
\end{Highlighting}
\end{Shaded}

\begin{longtable}[]{@{}lrrrrr@{}}
\caption{Author Age at Peak Data}\tabularnewline
\toprule\noalign{}
author & Type & Style & ageDth & cageDth & agePeak \\
\midrule\noalign{}
\endfirsthead
\toprule\noalign{}
author & Type & Style & ageDth & cageDth & agePeak \\
\midrule\noalign{}
\endhead
\bottomrule\noalign{}
\endlastfoot
Eliot & 0 & 0 & 77 & 8.57 & 23 \\
Cummings & 0 & 0 & 68 & -0.43 & 26 \\
Plath & 0 & 0 & 31 & -37.43 & 30 \\
Pound & 0 & 0 & 87 & 18.57 & 30 \\
Wilber & 0 & 0 & 86 & 17.57 & 34 \\
Williams & 0 & 1 & 80 & 11.57 & 40 \\
Bishop & 0 & 1 & 68 & -0.43 & 29 \\
Moore & 0 & 1 & 85 & 16.57 & 32 \\
Lowell & 0 & 1 & 60 & -8.43 & 41 \\
Stevens & 0 & 1 & 76 & 7.57 & 42 \\
Frost & 0 & 1 & 89 & 20.57 & 48 \\
Fitzgerald & 1 & 0 & 44 & -24.43 & 29 \\
Hemingway & 1 & 0 & 62 & -6.43 & 30 \\
Melville & 1 & 0 & 72 & 3.57 & 32 \\
Lawrence & 1 & 0 & 45 & -23.43 & 35 \\
Joyce & 1 & 0 & 59 & -9.43 & 40 \\
James & 1 & 1 & 73 & 4.57 & 38 \\
Faulkner & 1 & 1 & 65 & -3.43 & 39 \\
Dickens & 1 & 1 & 58 & -10.43 & 41 \\
Woolf & 1 & 1 & 59 & -9.43 & 45 \\
Conrad & 1 & 1 & 67 & -1.43 & 47 \\
Twain & 1 & 1 & 75 & 6.57 & 50 \\
Hardy & 1 & 1 & 88 & 19.57 & 51 \\
\end{longtable}

\begin{Shaded}
\begin{Highlighting}[]
\NormalTok{ap1}\SpecialCharTok{$}\NormalTok{w.style }\OtherTok{\textless{}{-}} \FunctionTok{factor}\NormalTok{(ap1}\SpecialCharTok{$}\NormalTok{Style, }\AttributeTok{levels=}\DecValTok{0}\SpecialCharTok{:}\DecValTok{1}\NormalTok{, }\AttributeTok{labels=}\FunctionTok{c}\NormalTok{(}\StringTok{"C"}\NormalTok{, }\StringTok{"E"}\NormalTok{))}
\NormalTok{ap1}\SpecialCharTok{$}\NormalTok{w.type }\OtherTok{\textless{}{-}} \FunctionTok{factor}\NormalTok{(ap1}\SpecialCharTok{$}\NormalTok{Type, }\AttributeTok{levels=}\DecValTok{0}\SpecialCharTok{:}\DecValTok{1}\NormalTok{, }\AttributeTok{labels=}\FunctionTok{c}\NormalTok{(}\StringTok{"P"}\NormalTok{, }\StringTok{"N"}\NormalTok{))}

\NormalTok{ap1}\SpecialCharTok{$}\NormalTok{X1A }\OtherTok{\textless{}{-}} \FunctionTok{ifelse}\NormalTok{(ap1}\SpecialCharTok{$}\NormalTok{Style}\SpecialCharTok{==}\DecValTok{0}\NormalTok{,}\DecValTok{1}\NormalTok{,}\SpecialCharTok{{-}}\DecValTok{1}\NormalTok{)}
\NormalTok{ap1}\SpecialCharTok{$}\NormalTok{X1B }\OtherTok{\textless{}{-}} \FunctionTok{ifelse}\NormalTok{(ap1}\SpecialCharTok{$}\NormalTok{Type}\SpecialCharTok{==}\DecValTok{0}\NormalTok{,}\DecValTok{1}\NormalTok{,}\SpecialCharTok{{-}}\DecValTok{1}\NormalTok{)}
\NormalTok{ap1}\SpecialCharTok{$}\NormalTok{X1AX1B }\OtherTok{\textless{}{-}}\NormalTok{ ap1}\SpecialCharTok{$}\NormalTok{X1A}\SpecialCharTok{*}\NormalTok{ap1}\SpecialCharTok{$}\NormalTok{X1B}
\DocumentationTok{\#\# Regression Model }

\NormalTok{ap.reg1 }\OtherTok{\textless{}{-}} \FunctionTok{lm}\NormalTok{(agePeak }\SpecialCharTok{\textasciitilde{}}\NormalTok{ X1A }\SpecialCharTok{+}\NormalTok{ X1B }\SpecialCharTok{+}\NormalTok{ X1AX1B, }\AttributeTok{data=}\NormalTok{ap1)}
\FunctionTok{summary}\NormalTok{(ap.reg1)}
\end{Highlighting}
\end{Shaded}

\begin{verbatim}
## 
## Call:
## lm(formula = agePeak ~ X1A + X1B + X1AX1B, data = ap1)
## 
## Residuals:
##    Min     1Q Median     3Q    Max 
## -9.667 -3.814  1.333  2.952  9.333 
## 
## Coefficients:
##             Estimate Std. Error t value Pr(>|t|)    
## (Intercept)  36.2238     1.1402  31.769  < 2e-16 ***
## X1A          -5.3238     1.1402  -4.669 0.000167 ***
## X1B          -2.5905     1.1402  -2.272 0.034902 *  
## X1AX1B        0.2905     1.1402   0.255 0.801652    
## ---
## Signif. codes:  0 '***' 0.001 '**' 0.01 '*' 0.05 '.' 0.1 ' ' 1
## 
## Residual standard error: 5.415 on 19 degrees of freedom
## Multiple R-squared:  0.5978, Adjusted R-squared:  0.5343 
## F-statistic: 9.413 on 3 and 19 DF,  p-value: 0.0005033
\end{verbatim}

\begin{Shaded}
\begin{Highlighting}[]
\FunctionTok{anova}\NormalTok{(ap.reg1)}
\end{Highlighting}
\end{Shaded}

\begin{verbatim}
## Analysis of Variance Table
## 
## Response: agePeak
##           Df Sum Sq Mean Sq F value    Pr(>F)    
## X1A        1 667.75  667.75 22.7758 0.0001324 ***
## X1B        1 158.26  158.26  5.3979 0.0314109 *  
## X1AX1B     1   1.90    1.90  0.0649 0.8016516    
## Residuals 19 557.05   29.32                      
## ---
## Signif. codes:  0 '***' 0.001 '**' 0.01 '*' 0.05 '.' 0.1 ' ' 1
\end{verbatim}

\begin{Shaded}
\begin{Highlighting}[]
\FunctionTok{drop1}\NormalTok{(ap.reg1, }\AttributeTok{test=}\StringTok{"F"}\NormalTok{)}
\end{Highlighting}
\end{Shaded}

\begin{verbatim}
## Single term deletions
## 
## Model:
## agePeak ~ X1A + X1B + X1AX1B
##        Df Sum of Sq     RSS    AIC F value    Pr(>F)    
## <none>               557.05 81.305                      
## X1A     1    639.14 1196.19 96.882 21.8001 0.0001672 ***
## X1B     1    151.33  708.37 84.832  5.1615 0.0349023 *  
## X1AX1B  1      1.90  558.95 79.383  0.0649 0.8016516    
## ---
## Signif. codes:  0 '***' 0.001 '**' 0.01 '*' 0.05 '.' 0.1 ' ' 1
\end{verbatim}

\begin{Shaded}
\begin{Highlighting}[]
\DocumentationTok{\#\# ANOVA model (Treatment effects)}

\FunctionTok{options}\NormalTok{(}\AttributeTok{contrasts=}\FunctionTok{c}\NormalTok{(}\StringTok{"contr.sum"}\NormalTok{, }\StringTok{"contr.poly"}\NormalTok{))}
\NormalTok{ap.aov1 }\OtherTok{\textless{}{-}} \FunctionTok{aov}\NormalTok{(agePeak }\SpecialCharTok{\textasciitilde{}}\NormalTok{ w.style}\SpecialCharTok{*}\NormalTok{w.type, }\AttributeTok{data=}\NormalTok{ap1)}
\FunctionTok{summary.lm}\NormalTok{(ap.aov1)}
\end{Highlighting}
\end{Shaded}

\begin{verbatim}
## 
## Call:
## aov(formula = agePeak ~ w.style * w.type, data = ap1)
## 
## Residuals:
##    Min     1Q Median     3Q    Max 
## -9.667 -3.814  1.333  2.952  9.333 
## 
## Coefficients:
##                  Estimate Std. Error t value Pr(>|t|)    
## (Intercept)       36.2238     1.1402  31.769  < 2e-16 ***
## w.style1          -5.3238     1.1402  -4.669 0.000167 ***
## w.type1           -2.5905     1.1402  -2.272 0.034902 *  
## w.style1:w.type1   0.2905     1.1402   0.255 0.801652    
## ---
## Signif. codes:  0 '***' 0.001 '**' 0.01 '*' 0.05 '.' 0.1 ' ' 1
## 
## Residual standard error: 5.415 on 19 degrees of freedom
## Multiple R-squared:  0.5978, Adjusted R-squared:  0.5343 
## F-statistic: 9.413 on 3 and 19 DF,  p-value: 0.0005033
\end{verbatim}

\begin{Shaded}
\begin{Highlighting}[]
\FunctionTok{anova}\NormalTok{(ap.aov1)}
\end{Highlighting}
\end{Shaded}

\begin{verbatim}
## Analysis of Variance Table
## 
## Response: agePeak
##                Df Sum Sq Mean Sq F value    Pr(>F)    
## w.style         1 667.75  667.75 22.7758 0.0001324 ***
## w.type          1 158.26  158.26  5.3979 0.0314109 *  
## w.style:w.type  1   1.90    1.90  0.0649 0.8016516    
## Residuals      19 557.05   29.32                      
## ---
## Signif. codes:  0 '***' 0.001 '**' 0.01 '*' 0.05 '.' 0.1 ' ' 1
\end{verbatim}

\begin{Shaded}
\begin{Highlighting}[]
\FunctionTok{drop1}\NormalTok{(ap.aov1, .}\SpecialCharTok{\textasciitilde{}}\NormalTok{., }\AttributeTok{test=}\StringTok{"F"}\NormalTok{)}
\end{Highlighting}
\end{Shaded}

\begin{verbatim}
## Single term deletions
## 
## Model:
## agePeak ~ w.style * w.type
##                Df Sum of Sq     RSS    AIC F value    Pr(>F)    
## <none>                       557.05 81.305                      
## w.style         1    639.14 1196.19 96.882 21.8001 0.0001672 ***
## w.type          1    151.33  708.37 84.832  5.1615 0.0349023 *  
## w.style:w.type  1      1.90  558.95 79.383  0.0649 0.8016516    
## ---
## Signif. codes:  0 '***' 0.001 '**' 0.01 '*' 0.05 '.' 0.1 ' ' 1
\end{verbatim}

We see from the regression coefficients (that control for all other
factors) that there is no evidence of an interaction, but there is
evidence of main effects for Style and Writer Type. Formal tests are
given below. Given that each test has 1 degree of freedom, we use
\(t\)-tests.

\[
H_0^{AB}: (\alpha\beta)_{11}=(\alpha\beta)_{12}=(\alpha\beta)_{21}=(\alpha\beta)_{22}=0
\qquad
t_{AB}^*=\frac{0.2905}{1.1402}=0.255
\qquad P_{AB}=.8017
\] \[
H_0^A:\alpha_1=\alpha_2=0 \qquad
t_A^*=\frac{-5.3238}{1.1402}=-4.669 \qquad P_A=.0002
\] \[
H_0^B: \beta_1=\beta_2=0 \qquad
t_B^*=\frac{-2.5905}{1.1402}=-2.272   \qquad P_B=.0349
\] As there is strong evidence of main effects and no suggestion of an
interaction, we will fit Model 2 for making comparisons, and not fit
Models 3 or 4.

\begin{verbatim}
## 
## Call:
## lm(formula = agePeak ~ X1A + X1B, data = ap1)
## 
## Residuals:
##    Min     1Q Median     3Q    Max 
## -9.940 -4.027  1.060  2.933  9.060 
## 
## Coefficients:
##             Estimate Std. Error t value Pr(>|t|)    
## (Intercept)   36.234      1.113  32.566  < 2e-16 ***
## X1A           -5.334      1.113  -4.794 0.000111 ***
## X1B           -2.628      1.104  -2.380 0.027397 *  
## ---
## Signif. codes:  0 '***' 0.001 '**' 0.01 '*' 0.05 '.' 0.1 ' ' 1
## 
## Residual standard error: 5.287 on 20 degrees of freedom
## Multiple R-squared:  0.5964, Adjusted R-squared:  0.5561 
## F-statistic: 14.78 on 2 and 20 DF,  p-value: 0.0001146
\end{verbatim}

\begin{verbatim}
## 
## Call:
## aov(formula = agePeak ~ factor(Style) + factor(Type), data = ap1)
## 
## Residuals:
##    Min     1Q Median     3Q    Max 
## -9.940 -4.027  1.060  2.933  9.060 
## 
## Coefficients:
##                Estimate Std. Error t value Pr(>|t|)    
## (Intercept)      36.234      1.113  32.566  < 2e-16 ***
## factor(Style)1   -5.334      1.113  -4.794 0.000111 ***
## factor(Type)1    -2.628      1.104  -2.380 0.027397 *  
## ---
## Signif. codes:  0 '***' 0.001 '**' 0.01 '*' 0.05 '.' 0.1 ' ' 1
## 
## Residual standard error: 5.287 on 20 degrees of freedom
## Multiple R-squared:  0.5964, Adjusted R-squared:  0.5561 
## F-statistic: 14.78 on 2 and 20 DF,  p-value: 0.0001146
\end{verbatim}

\begin{verbatim}
## Analysis of Variance Table
## 
## Response: agePeak
##               Df Sum Sq Mean Sq F value    Pr(>F)    
## factor(Style)  1 667.75  667.75 23.8930 8.894e-05 ***
## factor(Type)   1 158.26  158.26  5.6627    0.0274 *  
## Residuals     20 558.95   27.95                      
## ---
## Signif. codes:  0 '***' 0.001 '**' 0.01 '*' 0.05 '.' 0.1 ' ' 1
\end{verbatim}

Based on the additive model (Model 2), we obtain the following estimates
and predicted ages at peak.

\[
\hat{\mu}_{\bullet\bullet}=36.234 \qquad \hat{\alpha}_1=-5.334 \quad
\hat{\alpha}_2=5.334 \qquad \hat{\beta}_1=-2.628 \quad
\hat{\beta}_2=2.628
\] \[
\mbox{Conceptualists/Poets:} \quad 
\hat{Y}_{11}=\hat{\mu}_{\bullet\bullet}+\hat{\alpha}_1+\hat{\beta}_1=
36.234-5.334-2.628=28.272
\]

\[
\mbox{Conceptualists/Novelists:} \quad 
\hat{Y}_{12}=\hat{\mu}_{\bullet\bullet}+\hat{\alpha}_1+\hat{\beta}_2=
36.234-5.334+2.628=33.528
\]

\[
\mbox{Experimantalists/Poets:} \quad 
\hat{Y}_{21}=\hat{\mu}_{\bullet\bullet}+\hat{\alpha}_2+\hat{\beta}_1=
36.234+5.334-2.628=38.940
\]

\[
\mbox{Experimantalists/Novelists:} \quad 
\hat{Y}_{21}=\hat{\mu}_{\bullet\bullet}+\hat{\alpha}_2+\hat{\beta}_2=
36.234+5.334+2.628=44.196
\]

Experimentalists take longer on average to reach their peak than
conceptualists, and novelists take longer than poets.

\[\nabla \]

A second example with more levels of factor B is included here.

\textbf{Example 7.2 - Makiwara Punching Boards}

An experiment was conducted to compare four wood types of and two board
types, of karate boards (\citeproc{ref-smith_2010}{P. K. Smith, Niiler,
and McCullough 2010}). The wood types were: cherry \((i=1)\), ash
\((i=2)\), fir \((i=3)\), and oak \((i=4)\). The board types were
stacked \((j=1)\) and tapered \((j=2)\). Apparently due to breakage,
there were unequal numbers of replicates among the \(ab=2(4)=8\)
treatments. We use the \textbf{aov} function in R with the treatment
effects form to fit Models 1 (Interaction) and 2 (Additive) to test
whether there is an interaction between wood and board types. The
response \(Y\) is the deflection (millimeters) for the Makiwara board,
with a total of \(n_T=336\) measurements.

\begin{Shaded}
\begin{Highlighting}[]
\NormalTok{mb1 }\OtherTok{\textless{}{-}} \FunctionTok{read.csv}\NormalTok{(}\StringTok{"data/karate\_board.csv"}\NormalTok{)}
\FunctionTok{head}\NormalTok{(mb1,}\DecValTok{2}\NormalTok{)}
\end{Highlighting}
\end{Shaded}

\begin{verbatim}
##   Trt Wood Bdtype repnum Deflect X1A X2A X3A X1B X1AX1B X2AX1B X3AX1B
## 1   1    1      1      1  144.27   1   0   0   1      1      0      0
## 2   1    1      1      2  125.89   1   0   0   1      1      0      0
\end{verbatim}

\begin{Shaded}
\begin{Highlighting}[]
\FunctionTok{tail}\NormalTok{(mb1,}\DecValTok{2}\NormalTok{)}
\end{Highlighting}
\end{Shaded}

\begin{verbatim}
##     Trt Wood Bdtype repnum Deflect X1A X2A X3A X1B X1AX1B X2AX1B X3AX1B
## 335   8    4      2     44   73.32  -1  -1  -1  -1      1      1      1
## 336   8    4      2     45   44.86  -1  -1  -1  -1      1      1      1
\end{verbatim}

\begin{Shaded}
\begin{Highlighting}[]
\NormalTok{mb1}\SpecialCharTok{$}\NormalTok{Wood.f }\OtherTok{\textless{}{-}} \FunctionTok{factor}\NormalTok{(mb1}\SpecialCharTok{$}\NormalTok{Wood, }\AttributeTok{levels=}\DecValTok{1}\SpecialCharTok{:}\DecValTok{4}\NormalTok{, }\AttributeTok{labels=}\FunctionTok{c}\NormalTok{(}\StringTok{"cherry"}\NormalTok{, }\StringTok{"ash"}\NormalTok{, }\StringTok{"fir"}\NormalTok{, }\StringTok{"oak"}\NormalTok{))}
\NormalTok{mb1}\SpecialCharTok{$}\NormalTok{Bdtype.f }\OtherTok{\textless{}{-}} \FunctionTok{factor}\NormalTok{(mb1}\SpecialCharTok{$}\NormalTok{Bdtype, }\AttributeTok{levels=}\DecValTok{1}\SpecialCharTok{:}\DecValTok{2}\NormalTok{, }\AttributeTok{labels=}\FunctionTok{c}\NormalTok{(}\StringTok{"stacked"}\NormalTok{, }\StringTok{"tapered"}\NormalTok{))}

\FunctionTok{tapply}\NormalTok{(mb1}\SpecialCharTok{$}\NormalTok{Deflect, }\FunctionTok{list}\NormalTok{(mb1}\SpecialCharTok{$}\NormalTok{Wood.f, mb1}\SpecialCharTok{$}\NormalTok{Bdtype.f), mean)}
\end{Highlighting}
\end{Shaded}

\begin{verbatim}
##          stacked  tapered
## cherry 117.99976 77.99978
## ash    106.00056 75.99933
## fir     94.00059 55.00000
## oak     89.99956 50.00067
\end{verbatim}

\begin{Shaded}
\begin{Highlighting}[]
\FunctionTok{tapply}\NormalTok{(mb1}\SpecialCharTok{$}\NormalTok{Deflect, }\FunctionTok{list}\NormalTok{(mb1}\SpecialCharTok{$}\NormalTok{Wood.f, mb1}\SpecialCharTok{$}\NormalTok{Bdtype.f), length)}
\end{Highlighting}
\end{Shaded}

\begin{verbatim}
##        stacked tapered
## cherry      41      45
## ash         36      45
## fir         34      45
## oak         45      45
\end{verbatim}

\begin{Shaded}
\begin{Highlighting}[]
\FunctionTok{options}\NormalTok{(}\AttributeTok{contrasts=}\FunctionTok{c}\NormalTok{(}\StringTok{"contr.sum"}\NormalTok{, }\StringTok{"contr.poly"}\NormalTok{))}

\NormalTok{mb.mod1 }\OtherTok{\textless{}{-}} \FunctionTok{aov}\NormalTok{(Deflect }\SpecialCharTok{\textasciitilde{}}\NormalTok{ Wood.f}\SpecialCharTok{*}\NormalTok{Bdtype.f, }\AttributeTok{data=}\NormalTok{mb1)}
\FunctionTok{summary.lm}\NormalTok{(mb.mod1)}
\end{Highlighting}
\end{Shaded}

\begin{verbatim}
## 
## Call:
## aov(formula = Deflect ~ Wood.f * Bdtype.f, data = mb1)
## 
## Residuals:
##    Min     1Q Median     3Q    Max 
## -97.59 -35.81  -7.23  25.69 281.13 
## 
## Coefficients:
##                   Estimate Std. Error t value Pr(>|t|)    
## (Intercept)        83.3750     3.0838  27.036  < 2e-16 ***
## Wood.f1            14.6247     5.2834   2.768  0.00596 ** 
## Wood.f2             7.6249     5.4085   1.410  0.15954    
## Wood.f3            -8.8747     5.4678  -1.623  0.10553    
## Bdtype.f1          18.6251     3.0838   6.040 4.19e-09 ***
## Wood.f1:Bdtype.f1   1.3749     5.2834   0.260  0.79485    
## Wood.f2:Bdtype.f1  -3.6245     5.4085  -0.670  0.50324    
## Wood.f3:Bdtype.f1   0.8752     5.4678   0.160  0.87293    
## ---
## Signif. codes:  0 '***' 0.001 '**' 0.01 '*' 0.05 '.' 0.1 ' ' 1
## 
## Residual standard error: 56.2 on 328 degrees of freedom
## Multiple R-squared:  0.1359, Adjusted R-squared:  0.1174 
## F-statistic: 7.366 on 7 and 328 DF,  p-value: 3.18e-08
\end{verbatim}

\begin{Shaded}
\begin{Highlighting}[]
\FunctionTok{anova}\NormalTok{(mb.mod1)}
\end{Highlighting}
\end{Shaded}

\begin{verbatim}
## Analysis of Variance Table
## 
## Response: Deflect
##                  Df  Sum Sq Mean Sq F value    Pr(>F)    
## Wood.f            3   45089   15030  4.7582  0.002913 ** 
## Bdtype.f          1  116342  116342 36.8329 3.552e-09 ***
## Wood.f:Bdtype.f   3    1441     480  0.1521  0.928296    
## Residuals       328 1036035    3159                      
## ---
## Signif. codes:  0 '***' 0.001 '**' 0.01 '*' 0.05 '.' 0.1 ' ' 1
\end{verbatim}

\begin{Shaded}
\begin{Highlighting}[]
\FunctionTok{drop1}\NormalTok{(mb.mod1, .}\SpecialCharTok{\textasciitilde{}}\NormalTok{., }\AttributeTok{test=}\StringTok{"F"}\NormalTok{)}
\end{Highlighting}
\end{Shaded}

\begin{verbatim}
## Single term deletions
## 
## Model:
## Deflect ~ Wood.f * Bdtype.f
##                 Df Sum of Sq     RSS    AIC F value    Pr(>F)    
## <none>                       1036035 2715.4                      
## Wood.f           3     45205 1081239 2723.7  4.7705  0.002866 ** 
## Bdtype.f         1    115217 1151251 2748.8 36.4766 4.187e-09 ***
## Wood.f:Bdtype.f  3      1441 1037476 2709.8  0.1521  0.928296    
## ---
## Signif. codes:  0 '***' 0.001 '**' 0.01 '*' 0.05 '.' 0.1 ' ' 1
\end{verbatim}

\begin{Shaded}
\begin{Highlighting}[]
\NormalTok{mb.mod2 }\OtherTok{\textless{}{-}} \FunctionTok{aov}\NormalTok{(Deflect }\SpecialCharTok{\textasciitilde{}}\NormalTok{ Wood.f}\SpecialCharTok{+}\NormalTok{Bdtype.f, }\AttributeTok{data=}\NormalTok{mb1)}
\FunctionTok{summary.lm}\NormalTok{(mb.mod2)}
\end{Highlighting}
\end{Shaded}

\begin{verbatim}
## 
## Call:
## aov(formula = Deflect ~ Wood.f + Bdtype.f, data = mb1)
## 
## Residuals:
##     Min      1Q  Median      3Q     Max 
## -96.212 -35.819  -7.721  26.743 282.446 
## 
## Coefficients:
##             Estimate Std. Error t value Pr(>|t|)    
## (Intercept)   83.434      3.067  27.206  < 2e-16 ***
## Wood.f1       14.505      5.252   2.762  0.00607 ** 
## Wood.f2        7.976      5.358   1.488  0.13760    
## Wood.f3       -9.047      5.407  -1.673  0.09520 .  
## Bdtype.f1     18.684      3.067   6.092 3.09e-09 ***
## ---
## Signif. codes:  0 '***' 0.001 '**' 0.01 '*' 0.05 '.' 0.1 ' ' 1
## 
## Residual standard error: 55.99 on 331 degrees of freedom
## Multiple R-squared:  0.1346, Adjusted R-squared:  0.1242 
## F-statistic: 12.88 on 4 and 331 DF,  p-value: 9.386e-10
\end{verbatim}

\begin{Shaded}
\begin{Highlighting}[]
\FunctionTok{anova}\NormalTok{(mb.mod2)}
\end{Highlighting}
\end{Shaded}

\begin{verbatim}
## Analysis of Variance Table
## 
## Response: Deflect
##            Df  Sum Sq Mean Sq F value    Pr(>F)    
## Wood.f      3   45089   15030  4.7951  0.002769 ** 
## Bdtype.f    1  116342  116342 37.1181 3.088e-09 ***
## Residuals 331 1037476    3134                      
## ---
## Signif. codes:  0 '***' 0.001 '**' 0.01 '*' 0.05 '.' 0.1 ' ' 1
\end{verbatim}

\begin{Shaded}
\begin{Highlighting}[]
\FunctionTok{drop1}\NormalTok{(mb.mod2, .}\SpecialCharTok{\textasciitilde{}}\NormalTok{., }\AttributeTok{test=}\StringTok{"F"}\NormalTok{)}
\end{Highlighting}
\end{Shaded}

\begin{verbatim}
## Single term deletions
## 
## Model:
## Deflect ~ Wood.f + Bdtype.f
##          Df Sum of Sq     RSS    AIC F value    Pr(>F)    
## <none>                1037476 2709.8                      
## Wood.f    3     45948 1083424 2718.4  4.8864  0.002448 ** 
## Bdtype.f  1    116342 1153818 2743.5 37.1181 3.088e-09 ***
## ---
## Signif. codes:  0 '***' 0.001 '**' 0.01 '*' 0.05 '.' 0.1 ' ' 1
\end{verbatim}

When comparing the mean deflection differences (stacked-tapered) within
the 4 wood types, based on the \textbf{tapply} function, we observe: 40
(cherry), 30 (ash), 39 (fir), and 41 (oak). While these differences are
not exactly the same, they are very consistent. To test for interaction,
we directly use the ANOVA table from Model 1.

\[
H_0: (\alpha\beta)_{11}= \cdots = (\alpha\beta)_{42}=0 \qquad
H_A:\mbox{ Not all } (\alpha\beta)_{ij}= 0
\]

\[
\mbox{Test Statistic:} \quad F_{AB}^*=\frac{MSAB}{MSE}=\frac{481}{3159}=0.1524
\qquad
\mbox{Rejection Region:}\quad F_{AB}^*\ge F_{.95;3,328}=2.632
\]

The \(P\)-value is large, \(P_{AB}=.9281\), giving no evidence of any
interaction effects.

Based on this, we test for main effects for Board and Wood Effects,
based on Model 2, the additive model. The (rounded) partial Mean Squares
for Model 2 are \(MSA=45948/3=15316\) for wood type and \(MSB=116342/1\)
for board type, obtained from the \textbf{drop1} function.

For Wood effects, \(H_0^A:\alpha_1=\cdots = \alpha_4=0\), we obtain the
following test.

\[
\mbox{Test Statistic:} \quad F_{A}^*=\frac{MSA}{MSE}=\frac{15316}{3134}=4.89
\qquad
\mbox{Rejection Region:}\quad F_{A}^*\ge F_{.95;3,331}=2.632
\] For Board effects, \(H_0^B:\beta_1=\beta_2=0\), we obtain the
following test.

\[
\mbox{Test Statistic:} \quad F_{B}^*=\frac{MSB}{MSE}=\frac{116342}{3134}=37.12
\qquad
\mbox{Rejection Region:}\quad F_{B}^*\ge F_{.95;1,331}=3.870
\]

Both results are significant, with \(P_A=.0024\) and \(P_B<.0001\).

\[\nabla\]

\section{Least Squares Estimators and Contrasts/Linear Functions Among
Means}\label{least-squares-estimators-and-contrastslinear-functions-among-means}

In this section, the formulas for least squares estimators and estimated
standard errors for means and linear functions are given. For treatment
(cell) means, we have the following results.

\[ 
\mbox{Parameter:} \quad \mu_{ij} \qquad 
\mbox{Estimator:} \quad \hat{\mu}_{ij} = \frac{\sum_{k=1}^{n_{ij}}Y_{ijk}}{n_{ij}}=
\overline{Y}_{ij\bullet}
\] \[
\mbox{Estimated Standard Error:} \quad \hat{SE}\left\{\hat{\mu}_{ij}\right\}=\sqrt{\frac{MSE}{n_{ij}}}
\] For factor A level means, we have the following results. \[ 
\mbox{Parameter:} \quad \mu_{i\bullet} = \frac{\sum_{j=1}^b\mu_{ij}}{b} \qquad 
\mbox{Estimator:} \quad \hat{\mu}_{i\bullet} = 
\frac{\sum_{j=1}^b\overline{Y}_{ij\bullet}}{b} 
\] \[
\mbox{Estimated Standard Error:} \quad \hat{SE}\left\{\hat{\mu}_{i\bullet}\right\}=
\sqrt{\frac{MSE}{b^2}\sum_{j=1}^b\frac{1}{n_{ij}}}
\] For factor B level means, we have the following results. \[ 
\mbox{Parameter:} \quad \mu_{\bullet j} = \frac{\sum_{i=1}^a\mu_{ij}}{a} \qquad 
\mbox{Estimator:} \quad \hat{\mu}_{\bullet j} = 
\frac{\sum_{i=1}^a\overline{Y}_{ij\bullet}}{a} 
\] \[
\mbox{Estimated Standard Error:} \quad \hat{SE}\left\{\hat{\mu}_{\bullet j}\right\}=
\sqrt{\frac{MSE}{a^2}\sum_{i=1}^a\frac{1}{n_{ij}}}
\]

For contrasts or linear functions among factor A and B levels, we have
the following parameters, estimators, and estimated standard errors.

\[
L_A = \sum_{i=1}^a c_i\mu_{i\bullet} \qquad
\hat{L}_A = \sum_{i=1}^a c_i\hat{\mu}_{i\bullet} \qquad
\hat{SE}\left\{\hat{L}_A\right\} = \sqrt{\frac{MSE}{b^2}\sum_{i=1}^ac_i^2\sum_{j=1}^b\frac{1}{n_{ij}}}
\]

\[
L_B = \sum_{j=1}^b c_j\mu_{\bullet j} \qquad
\hat{L}_B = \sum_{j=1}^b c_j\hat{\mu}_{\bullet j} \qquad
\hat{SE}\left\{\hat{L}_B\right\} = \sqrt{\frac{MSE}{a^2}\sum_{j=1}^bc_j^2\sum_{i=1}^a\frac{1}{n_{ij}}}
\] Finally, the results are given for contrasts and linear functions
among treatment (cell) means.

\[
L_{AB} = \sum_{i=1}^a \sum_{j=1}^b c_{ij}\mu_{ij} \qquad
\hat{L}_{AB} = \sum_{i=1}^a\sum_{j=1}^b c_{ij}\hat{\mu}_{ij} \qquad
\hat{SE}\left\{\hat{L}_{AB}\right\} = \sqrt{MSE\sum_{i=1}^a\sum_{j=1}^b\frac{c_{ij}^2}{n_{ij}}}
\]

The standard error multipliers for contrasts/linear functions based on
the Scheffe (all comparisons), Bonferroni (\(g\) pre-planned
comparisons), and Tukey (all pairwise comparisons) methods are given
here.

For single comparisons, use \(t_{1-\alpha/2;n_T-ab}\).

For comparisons among treatment (cell) means, use the following
multipliers.

\[
\mbox{Scheffe:} \quad S=\sqrt{(ab-1)F_{1-\alpha;ab-1,n_T-ab}}
\] \[
\mbox{Bonferroni:} \quad B=t_{1-\alpha/(2g); n_T-ab}
\] \[
\mbox{Tukey:} \quad T=\frac{1}{\sqrt{2}}q_{1-\alpha;ab,n_T-ab}
\]

Standard error multipliers for contrasts and linear functions among
levels of factors A and B are given here.

\[
\mbox{Scheffe:} \qquad \mbox{A:} \quad S_A=\sqrt{(a-1)F_{1-\alpha;a-1,n_T-ab}} \qquad
\quad \mbox{B:} \quad S_B=\sqrt{(b-1)F_{1-\alpha;b-1,n_T-ab}}
\] \[
\mbox{Bonferroni (Factors A and B):} \quad B=t_{1-\alpha/(2g); n_T-ab}
\] \[
\mbox{Tukey:} \qquad \mbox{A:} \quad T_A=\frac{1}{\sqrt{2}}q_{1-\alpha;a,n_T-ab} \qquad
\mbox{B:} \quad T_B=\frac{1}{\sqrt{2}}q_{1-\alpha;b,n_T-ab}
\] Be warned that the \textbf{TukeyHSD} function in R does not use the
least squares means when making pairwise comparisons among means. For
the first factor given in the \textbf{aov} object, it computes
differences among the \textbf{marginal} means. For each level of factor
A, the effect \(a_i\) is obtained. Then, the \(a_i\) are subtracted from
the original data, for observations ``adjusted'' for factor A. Then the
marginal means for the adjusted data are obtained for each level of
factor B.

\[
\overline{Y}_{i\bullet\bullet}=\frac{\sum_{j=1}^b\sum_{k=1}^nY_{ijk}}{n_{i\bullet}}
\qquad a_i=\overline{Y}_{i\bullet\bullet}-\overline{Y}_{\bullet\bullet\bullet}
\qquad Y_{ijk}^*=Y_{ijk}-a_i  \qquad 
\overline{Y}_{\bullet j\bullet}^*=\frac{\sum_{i=1}^a\sum_{k=1}^nY_{ijk}^*}{n_{\bullet} j}
\]

The \textbf{emmeans} package makes use of the least squares (estimated
marginal) means, and the formulas given in this section. For the
interaction comparisons among cell means, the results are the same.

\textbf{Example 7.3 - Age at Peak for Famous Writers}

Although the interaction was not significant for the Age at Peak of the
writers, we will use the Tukey method to compare all pairs of Style/Type
means. The Mean Square Error for the interaction model was \(MSE=29.32\)
with degrees of freedom \(n_T-ab=23-2(2)=19\).

\[
\overline{y}_{11\bullet} = 28.60 \qquad n_{11}=5 \qquad 
\hat{SE}\left\{\overline{Y}_{11\bullet}\right\}=\sqrt{\frac{MSE}{n_{11}}}=\sqrt{\frac{29.32}{5}}=2.42
\] \[
\overline{y}_{12\bullet} = 33.20 \qquad n_{12}=5 \qquad 
\hat{SE}\left\{\overline{Y}_{12\bullet}\right\}=\sqrt{\frac{MSE}{n_{12}}}=\sqrt{\frac{29.32}{5}}=2.42
\] \[
\overline{y}_{21\bullet} = 38.67 \qquad n_{21}=6 \qquad 
\hat{SE}\left\{\overline{Y}_{21\bullet}\right\}=\sqrt{\frac{MSE}{n_{21}}}=\sqrt{\frac{29.32}{6}}=2.21
\] \[
\overline{y}_{22\bullet} = 44.43 \qquad n_{22}=7 \qquad 
\hat{SE}\left\{\overline{Y}_{22\bullet}\right\}=\sqrt{\frac{MSE}{n_{22}}}=\sqrt{\frac{29.32}{7}}=2.05
\] \[
T = \frac{1}{\sqrt{2}}q_{.95;4,19}=\frac{1}{\sqrt{2}}(3.977)=2.812 
\qquad
\hat{SE}\left\{\overline{Y}_{ij\bullet}-\overline{Y}_{i'j'\bullet}\right\}=
\sqrt{MSE\left(\frac{1}{n_{ij}}+\frac{1}{n_{i'j'}}\right)}
\] \[
\overline{y}_{11\bullet}-\overline{y}_{12\bullet} = 28.60-33.20=-4.60 \qquad
\hat{SE}\left\{\overline{Y}_{11\bullet}-\overline{Y}_{12\bullet}\right\} =
\sqrt{29.32\left(\frac{1}{5}+\frac{1}{5}\right)}=3.43 
\] \[
\overline{y}_{11\bullet}-\overline{y}_{21\bullet} = 28.60-38.67=-10.07 \qquad
\hat{SE}\left\{\overline{Y}_{11\bullet}-\overline{Y}_{21\bullet}\right\} =
\sqrt{29.32\left(\frac{1}{5}+\frac{1}{6}\right)}=3.28 
\] \[
\overline{y}_{11\bullet}-\overline{y}_{22\bullet} = 28.60-44.43=-15.83 \qquad
\hat{SE}\left\{\overline{Y}_{11\bullet}-\overline{Y}_{22\bullet}\right\} =
\sqrt{29.32\left(\frac{1}{5}+\frac{1}{7}\right)}=3.17 
\] \[
\overline{y}_{12\bullet}-\overline{y}_{21\bullet} = 33.20-38.67=-5.47 \qquad
\hat{SE}\left\{\overline{Y}_{12\bullet}-\overline{Y}_{21\bullet}\right\} =
\sqrt{29.32\left(\frac{1}{5}+\frac{1}{6}\right)}=3.28 
\] \[
\overline{y}_{12\bullet}-\overline{y}_{22\bullet} = 33.20-44.43=-11.23 \qquad
\hat{SE}\left\{\overline{Y}_{12\bullet}-\overline{Y}_{22\bullet}\right\} =
\sqrt{29.32\left(\frac{1}{5}+\frac{1}{7}\right)}=3.17 
\] \[
\overline{y}_{21\bullet}-\overline{y}_{22\bullet} = 38.67-44.43=-5.76 \qquad
\hat{SE}\left\{\overline{Y}_{21\bullet}-\overline{Y}_{22\bullet}\right\} =
\sqrt{29.32\left(\frac{1}{6}+\frac{1}{7}\right)}=3.01 
\]

\[
\mbox{Conc/Poet-Conc/Novel:} \quad  HSD=2.812(3.43)=9.65 \qquad 
\overline{y}_{11\bullet}-\overline{y}_{12\bullet} =-4.60
\] \[
\mbox{Conc/Poet-Exp/Poet:} \quad  HSD=2.812(3.28)=9.22 \qquad 
\overline{y}_{11\bullet}-\overline{y}_{21\bullet} = 28.60-38.67=-10.07
\] \[
\mbox{Conc/Poet-Exp/Novel:} \quad  HSD=2.812(3.17)=8.92 \qquad 
\overline{y}_{11\bullet}-\overline{y}_{22\bullet} =-15.83
\] \[
\mbox{Conc/Nov-Exp/Poet:} \quad  HSD=2.812(3.28)=9.22 \qquad 
\overline{y}_{12\bullet}-\overline{y}_{21\bullet} = -5.47
\] \[
\mbox{Conc/Nov-Exp/Nov:} \quad  HSD=2.812(3.17)=8.92 \qquad 
\overline{y}_{12\bullet}-\overline{y}_{22\bullet} = -11.23
\] \[
\mbox{Exp/Poet-Exp/Nov:} \quad  HSD=2.812(3.01)=8.47 \qquad 
\overline{y}_{21\bullet}-\overline{y}_{22\bullet} = -5.76
\]

The groups that are significantly different, based on Tukey's HSD are as
follow.

\begin{itemize}
\tightlist
\item
  Conceptualist/Poet is significantly lower than Experimentalist/Poet
\item
  Conceptualist/Poet is significantly lower than
  Experimentalist/Novelist
\item
  Conceptualist/Novelist is significantly lower than
  Experimentalist/Novelist
\end{itemize}

We now obtain the comparisons of Styles (Conceptualist-Experimentalist)
and Types (Poet-Novelist) based on Model 1 (even though the interaction
is not significant). Very similar results would be obtained from Model 2
(with adjustments to MSE and error degrees of freedom).

\[
\hat{\mu}_{1\bullet}=\frac{\overline{Y}_{11\bullet}+\overline{Y}_{12\bullet}}{2}=
\frac{28.6000+33.2000}{2}=30.9000  \]
\[\hat{\mu}_{2\bullet}=\frac{\overline{Y}_{21\bullet}+\overline{Y}_{22\bullet}}{2}=
\frac{38.6667+44.4286}{2}=41.5477
\] \[
\hat{SE}\left\{\hat{\mu}_{1\bullet}-\hat{\mu}_{2\bullet}\right\}=
\sqrt{\frac{MSE}{2^2}\left(
\frac{1}{n_{11}}+\frac{1}{n_{12}}+\frac{1}{n_{21}}+\frac{1}{n_{22}}\right)}=
\sqrt{\frac{29.32}{4}\left(
\frac{1}{5}+\frac{1}{5}+\frac{1}{6}+\frac{1}{7}\right)}=2.281
\] \[
95\% \mbox{CI for } \mu_{1\bullet}-\mu_{2\bullet}: \quad
(30.900-41.548) \pm 2.093(2.281)\equiv -10.648\pm 4.774\equiv (-15.422,-5.874)
\] Note that in this rare case where \(a=b=2\), the standard error of
\(\hat{\mu}_{\bullet 1}-\hat{\mu}_{\bullet 2}\) is the same as that of
\(\hat{\mu}_{1\bullet}-\hat{\mu}_{2\bullet}\), we have the following
results.

\[
\hat{\mu}_{\bullet 1}=\frac{\overline{Y}_{11\bullet}+\overline{Y}_{21\bullet}}{2}=
\frac{28.6000+38.6667}{2}=33.6334  \]
\[\hat{\mu}_{\bullet 2}=\frac{\overline{Y}_{12\bullet}+\overline{Y}_{22\bullet}}{2}=
\frac{33.2000+44.4286}{2}=38.8143
\]

\[
95\% \mbox{CI for } \mu_{\bullet}-\mu_{\bullet}: \quad (33.633-38.814)\pm
2.093(2.281)\equiv -5.181\pm 4.774\equiv (-9.955,-0.407)
\]

The following R program runs Tukey's HSD on the marginal means for
factors A and B, and on the cell means as well. Note, however, that the
\textbf{TukeyHSD} function in R is not based on the least squares (aka
estimated marginal) means. We also obtain the least squares means,
standard errors and pairwise comparisons described in this section with
the \textbf{emmeans} package.

\begin{Shaded}
\begin{Highlighting}[]
\DocumentationTok{\#\# Continuation of ch0701, ch0702}

\FunctionTok{options}\NormalTok{(}\AttributeTok{contrasts=}\FunctionTok{c}\NormalTok{(}\StringTok{"contr.sum"}\NormalTok{, }\StringTok{"contr.poly"}\NormalTok{))}
\NormalTok{ap.aov1 }\OtherTok{\textless{}{-}} \FunctionTok{aov}\NormalTok{(agePeak }\SpecialCharTok{\textasciitilde{}}\NormalTok{ w.style}\SpecialCharTok{*}\NormalTok{w.type, }\AttributeTok{data=}\NormalTok{ap1)}
\FunctionTok{TukeyHSD}\NormalTok{(ap.aov1)}
\end{Highlighting}
\end{Shaded}

\begin{verbatim}
##   Tukey multiple comparisons of means
##     95% family-wise confidence level
## 
## Fit: aov(formula = agePeak ~ w.style * w.type, data = ap1)
## 
## $w.style
##         diff      lwr      upr     p adj
## E-C 10.86923 6.102333 15.63613 0.0001324
## 
## $w.type
##         diff       lwr      upr     p adj
## N-P 5.247378 0.5167309 9.978024 0.0315205
## 
## $`w.style:w.type`
##              diff         lwr       upr     p adj
## E:P-C:P 10.066667   0.8473951 19.285938 0.0293965
## C:N-C:P  4.600000  -5.0292152 14.229215 0.5481270
## E:N-C:P 15.828571   6.9136505 24.743492 0.0004327
## C:N-E:P -5.466667 -14.6859383  3.752605 0.3672669
## E:N-E:P  5.761905  -2.7085734 14.232383 0.2559999
## E:N-C:N 11.228571   2.3136505 20.143492 0.0107243
\end{verbatim}

\begin{Shaded}
\begin{Highlighting}[]
\NormalTok{ap.emm1 }\OtherTok{\textless{}{-}} \FunctionTok{emmeans}\NormalTok{(ap.aov1, }\SpecialCharTok{\textasciitilde{}}\NormalTok{ w.style)}
\end{Highlighting}
\end{Shaded}

\begin{verbatim}
## NOTE: Results may be misleading due to involvement in interactions
\end{verbatim}

\begin{Shaded}
\begin{Highlighting}[]
\FunctionTok{emmeans}\NormalTok{(ap.emm1, }\SpecialCharTok{\textasciitilde{}}\NormalTok{ w.style)}
\end{Highlighting}
\end{Shaded}

\begin{verbatim}
## NOTE: Results may be misleading due to involvement in interactions
\end{verbatim}

\begin{verbatim}
##  w.style emmean   SE df lower.CL upper.CL
##  C         30.9 1.71 19     27.3     34.5
##  E         41.5 1.51 19     38.4     44.7
## 
## Results are averaged over the levels of: w.type 
## Confidence level used: 0.95
\end{verbatim}

\begin{Shaded}
\begin{Highlighting}[]
\FunctionTok{contrast}\NormalTok{(ap.emm1, }\AttributeTok{method=}\StringTok{"pairwise"}\NormalTok{, }\AttributeTok{adjust=}\StringTok{"tukey"}\NormalTok{)}
\end{Highlighting}
\end{Shaded}

\begin{verbatim}
##  contrast estimate   SE df t.ratio p.value
##  C - E       -10.6 2.28 19  -4.669  0.0002
## 
## Results are averaged over the levels of: w.type
\end{verbatim}

\begin{Shaded}
\begin{Highlighting}[]
\NormalTok{ap.emm2 }\OtherTok{\textless{}{-}} \FunctionTok{emmeans}\NormalTok{(ap.aov1,  }\SpecialCharTok{\textasciitilde{}}\NormalTok{ w.type)}
\end{Highlighting}
\end{Shaded}

\begin{verbatim}
## NOTE: Results may be misleading due to involvement in interactions
\end{verbatim}

\begin{Shaded}
\begin{Highlighting}[]
\FunctionTok{emmeans}\NormalTok{(ap.emm2, }\SpecialCharTok{\textasciitilde{}}\NormalTok{ w.type)}
\end{Highlighting}
\end{Shaded}

\begin{verbatim}
## NOTE: Results may be misleading due to involvement in interactions
\end{verbatim}

\begin{verbatim}
##  w.type emmean   SE df lower.CL upper.CL
##  P        33.6 1.64 19     30.2     37.1
##  N        38.8 1.59 19     35.5     42.1
## 
## Results are averaged over the levels of: w.style 
## Confidence level used: 0.95
\end{verbatim}

\begin{Shaded}
\begin{Highlighting}[]
\FunctionTok{contrast}\NormalTok{(ap.emm2, }\AttributeTok{method=}\StringTok{"pairwise"}\NormalTok{, }\AttributeTok{adjust=}\StringTok{"tukey"}\NormalTok{)}
\end{Highlighting}
\end{Shaded}

\begin{verbatim}
##  contrast estimate   SE df t.ratio p.value
##  P - N       -5.18 2.28 19  -2.272  0.0349
## 
## Results are averaged over the levels of: w.style
\end{verbatim}

\begin{Shaded}
\begin{Highlighting}[]
\NormalTok{ap.emm3 }\OtherTok{\textless{}{-}} \FunctionTok{emmeans}\NormalTok{(ap.aov1, }\SpecialCharTok{\textasciitilde{}}\NormalTok{ w.style}\SpecialCharTok{:}\NormalTok{w.type)}
\FunctionTok{emmeans}\NormalTok{(ap.emm3, }\SpecialCharTok{\textasciitilde{}}\NormalTok{ w.style}\SpecialCharTok{:}\NormalTok{w.type)}
\end{Highlighting}
\end{Shaded}

\begin{verbatim}
##  w.style w.type emmean   SE df lower.CL upper.CL
##  C       P        28.6 2.42 19     23.5     33.7
##  E       P        38.7 2.21 19     34.0     43.3
##  C       N        33.2 2.42 19     28.1     38.3
##  E       N        44.4 2.05 19     40.1     48.7
## 
## Confidence level used: 0.95
\end{verbatim}

\begin{Shaded}
\begin{Highlighting}[]
\FunctionTok{contrast}\NormalTok{(ap.emm3, }\AttributeTok{method=}\StringTok{"pairwise"}\NormalTok{, }\AttributeTok{adjust=}\StringTok{"tukey"}\NormalTok{)}
\end{Highlighting}
\end{Shaded}

\begin{verbatim}
##  contrast  estimate   SE df t.ratio p.value
##  C P - E P   -10.07 3.28 19  -3.070  0.0294
##  C P - C N    -4.60 3.42 19  -1.343  0.5481
##  C P - E N   -15.83 3.17 19  -4.992  0.0004
##  E P - C N     5.47 3.28 19   1.667  0.3673
##  E P - E N    -5.76 3.01 19  -1.913  0.2560
##  C N - E N   -11.23 3.17 19  -3.542  0.0107
## 
## P value adjustment: tukey method for comparing a family of 4 estimates
\end{verbatim}

\begin{Shaded}
\begin{Highlighting}[]
\FunctionTok{emmip}\NormalTok{(ap.emm3, w.style }\SpecialCharTok{\textasciitilde{}}\NormalTok{ w.type)}
\end{Highlighting}
\end{Shaded}

\pandocbounded{\includegraphics[keepaspectratio]{bookdownproj_STA4211Notes1_files/figure-latex/ch0704-1.pdf}}

\[\nabla\]

\textbf{Example 7.4 - Makiwara Punching Boards}

For the Karate board experiment, we will conduct tests among levels of
factors A and B, respectively. The Mean Square Error was \(MSE=3159\)
with degrees of freedom \(n_T-ab=336-2(4)=328\) for Model 1 (with
interaction).

For comparing the Board Types (Stacked-Tapered), we obtain the following
cell means from the \textbf{tapply} function previously.

\[
\overline{y}_{11\bullet}=118 \quad \overline{y}_{12\bullet}=106 \quad
\overline{y}_{13\bullet}=94 \quad \overline{y}_{14\bullet}=90 \quad \Rightarrow \quad
\hat{\mu}_{1\bullet}=\frac{118+106+94+90}{4}=102
\] \[
\overline{y}_{21\bullet}=78 \quad \overline{y}_{22\bullet}=76 \quad
\overline{y}_{23\bullet}=55 \quad \overline{y}_{24\bullet}=50 \quad \Rightarrow \quad
\hat{\mu}_{2\bullet}=\frac{78+76+55+50}{4}=64.75
\]

The standard errors are computed here, where the sample sizes for each
cell are given first.

\[
n_{11}=41 \quad n_{12}=36 \quad n_{13}=34 \quad n_{14}=n_{21}=n_{22}=n_{23}=n_{24}=45
\] \[
\hat{SE}\left\{\hat{\mu}_{1\bullet}\right\}=
\sqrt{\frac{3159}{4^2}\left(\frac{1}{41}+\frac{1}{36}+\frac{1}{34}+\frac{1}{45}\right)}
=4.53
\] \[
\hat{SE}\left\{\hat{\mu}_{2\bullet}\right\}=
\sqrt{\frac{3159}{4^2}\left(\frac{1}{45}+\frac{1}{45}+\frac{1}{45}+\frac{1}{45}\right)}
=4.19
\] The standard error of the difference
\(\hat{\mu}_{1\bullet}-\hat{\mu}_{2\bullet}\) is the square root of the
sum of their variances, as these are independent samples. Also, we
obtain the critical \(t\)-value, and the 95\% Confidence Interval for
\(\mu_{1\bullet}-\mu_{2\bullet}\).

\[
\hat{SE}\left\{\hat{\mu}_{1\bullet}-\hat{\mu}_{2\bullet}\right\} = \sqrt{(4.53)^2+(4.19)^2}=6.17
\qquad
t_{.975;328}=1.967
\] \[
95\% \mbox{ Confidence Interval:} \quad (102-64.75) \pm 1.967(6.17) \quad \equiv \quad
37.25 \pm 12.14 \quad \equiv \quad (25.11,49.39)
\] Thus, the mean deflection is higher for the stacked board than the
tapered on average (likely by somewhere between 25-50 millimeters). We
could also compare all pairs of Wood Types. The computations are
similar, but there are 4(3)/2 = 6 possible pairs. We will use the
\textbf{TukeyHSD} function along with the \textbf{aov} object in R
below.

\begin{Shaded}
\begin{Highlighting}[]
\NormalTok{mb1 }\OtherTok{\textless{}{-}} \FunctionTok{read.csv}\NormalTok{(}\StringTok{"data/karate\_board.csv"}\NormalTok{)}
\FunctionTok{head}\NormalTok{(mb1,}\DecValTok{2}\NormalTok{)}
\end{Highlighting}
\end{Shaded}

\begin{verbatim}
##   Trt Wood Bdtype repnum Deflect X1A X2A X3A X1B X1AX1B X2AX1B X3AX1B
## 1   1    1      1      1  144.27   1   0   0   1      1      0      0
## 2   1    1      1      2  125.89   1   0   0   1      1      0      0
\end{verbatim}

\begin{Shaded}
\begin{Highlighting}[]
\FunctionTok{tail}\NormalTok{(mb1,}\DecValTok{2}\NormalTok{)}
\end{Highlighting}
\end{Shaded}

\begin{verbatim}
##     Trt Wood Bdtype repnum Deflect X1A X2A X3A X1B X1AX1B X2AX1B X3AX1B
## 335   8    4      2     44   73.32  -1  -1  -1  -1      1      1      1
## 336   8    4      2     45   44.86  -1  -1  -1  -1      1      1      1
\end{verbatim}

\begin{Shaded}
\begin{Highlighting}[]
\NormalTok{mb1}\SpecialCharTok{$}\NormalTok{Wood.f }\OtherTok{\textless{}{-}} \FunctionTok{factor}\NormalTok{(mb1}\SpecialCharTok{$}\NormalTok{Wood, }\AttributeTok{levels=}\DecValTok{1}\SpecialCharTok{:}\DecValTok{4}\NormalTok{, }\AttributeTok{labels=}\FunctionTok{c}\NormalTok{(}\StringTok{"cherry"}\NormalTok{, }\StringTok{"ash"}\NormalTok{, }\StringTok{"fir"}\NormalTok{, }\StringTok{"oak"}\NormalTok{))}
\NormalTok{mb1}\SpecialCharTok{$}\NormalTok{Bdtype.f }\OtherTok{\textless{}{-}} \FunctionTok{factor}\NormalTok{(mb1}\SpecialCharTok{$}\NormalTok{Bdtype, }\AttributeTok{levels=}\DecValTok{1}\SpecialCharTok{:}\DecValTok{2}\NormalTok{, }\AttributeTok{labels=}\FunctionTok{c}\NormalTok{(}\StringTok{"stacked"}\NormalTok{, }\StringTok{"tapered"}\NormalTok{))}

\FunctionTok{tapply}\NormalTok{(mb1}\SpecialCharTok{$}\NormalTok{Deflect, }\FunctionTok{list}\NormalTok{(mb1}\SpecialCharTok{$}\NormalTok{Wood.f, mb1}\SpecialCharTok{$}\NormalTok{Bdtype.f), mean)}
\end{Highlighting}
\end{Shaded}

\begin{verbatim}
##          stacked  tapered
## cherry 117.99976 77.99978
## ash    106.00056 75.99933
## fir     94.00059 55.00000
## oak     89.99956 50.00067
\end{verbatim}

\begin{Shaded}
\begin{Highlighting}[]
\NormalTok{mb.mod1 }\OtherTok{\textless{}{-}} \FunctionTok{aov}\NormalTok{(Deflect }\SpecialCharTok{\textasciitilde{}}\NormalTok{ Wood.f }\SpecialCharTok{*}\NormalTok{ Bdtype.f , }\AttributeTok{data=}\NormalTok{mb1)}

\NormalTok{mb.emm1 }\OtherTok{\textless{}{-}} \FunctionTok{emmeans}\NormalTok{(mb.mod1, }\SpecialCharTok{\textasciitilde{}}\NormalTok{ Wood.f)}
\end{Highlighting}
\end{Shaded}

\begin{verbatim}
## NOTE: Results may be misleading due to involvement in interactions
\end{verbatim}

\begin{Shaded}
\begin{Highlighting}[]
\FunctionTok{emmeans}\NormalTok{(mb.emm1, }\SpecialCharTok{\textasciitilde{}}\NormalTok{ Wood.f)}
\end{Highlighting}
\end{Shaded}

\begin{verbatim}
## NOTE: Results may be misleading due to involvement in interactions
\end{verbatim}

\begin{verbatim}
##  Wood.f emmean   SE  df lower.CL upper.CL
##  cherry   98.0 6.07 328     86.1    109.9
##  ash      91.0 6.28 328     78.6    103.4
##  fir      74.5 6.39 328     61.9     87.1
##  oak      70.0 5.92 328     58.3     81.7
## 
## Results are averaged over the levels of: Bdtype.f 
## Confidence level used: 0.95
\end{verbatim}

\begin{Shaded}
\begin{Highlighting}[]
\FunctionTok{contrast}\NormalTok{(mb.emm1, }\AttributeTok{method=}\StringTok{"pairwise"}\NormalTok{, }\AttributeTok{adjust=}\StringTok{"tukey"}\NormalTok{)}
\end{Highlighting}
\end{Shaded}

\begin{verbatim}
##  contrast     estimate   SE  df t.ratio p.value
##  cherry - ash      7.0 8.73 328   0.801  0.8537
##  cherry - fir     23.5 8.81 328   2.668  0.0398
##  cherry - oak     28.0 8.48 328   3.302  0.0058
##  ash - fir        16.5 8.96 328   1.842  0.2557
##  ash - oak        21.0 8.64 328   2.432  0.0732
##  fir - oak         4.5 8.71 328   0.517  0.9551
## 
## Results are averaged over the levels of: Bdtype.f 
## P value adjustment: tukey method for comparing a family of 4 estimates
\end{verbatim}

\begin{Shaded}
\begin{Highlighting}[]
\NormalTok{mb.emm2 }\OtherTok{\textless{}{-}} \FunctionTok{emmeans}\NormalTok{(mb.mod1, }\SpecialCharTok{\textasciitilde{}}\NormalTok{ Bdtype.f)}
\end{Highlighting}
\end{Shaded}

\begin{verbatim}
## NOTE: Results may be misleading due to involvement in interactions
\end{verbatim}

\begin{Shaded}
\begin{Highlighting}[]
\FunctionTok{emmeans}\NormalTok{(mb.emm2, }\SpecialCharTok{\textasciitilde{}}\NormalTok{ Bdtype.f)}
\end{Highlighting}
\end{Shaded}

\begin{verbatim}
## NOTE: Results may be misleading due to involvement in interactions
\end{verbatim}

\begin{verbatim}
##  Bdtype.f emmean   SE  df lower.CL upper.CL
##  stacked   102.0 4.53 328     93.1      111
##  tapered    64.7 4.19 328     56.5       73
## 
## Results are averaged over the levels of: Wood.f 
## Confidence level used: 0.95
\end{verbatim}

\begin{Shaded}
\begin{Highlighting}[]
\FunctionTok{contrast}\NormalTok{(mb.emm2, }\AttributeTok{method=}\StringTok{"pairwise"}\NormalTok{, }\AttributeTok{adjust=}\StringTok{"tukey"}\NormalTok{)}
\end{Highlighting}
\end{Shaded}

\begin{verbatim}
##  contrast          estimate   SE  df t.ratio p.value
##  stacked - tapered     37.3 6.17 328   6.040 <0.0001
## 
## Results are averaged over the levels of: Wood.f
\end{verbatim}

\begin{Shaded}
\begin{Highlighting}[]
\NormalTok{mb.emm3 }\OtherTok{\textless{}{-}} \FunctionTok{emmeans}\NormalTok{(mb.mod1, }\SpecialCharTok{\textasciitilde{}}\NormalTok{ Wood.f}\SpecialCharTok{:}\NormalTok{Bdtype.f)}
\FunctionTok{emmeans}\NormalTok{(mb.emm3, }\SpecialCharTok{\textasciitilde{}}\NormalTok{ Bdtype.f)}
\end{Highlighting}
\end{Shaded}

\begin{verbatim}
## NOTE: Results may be misleading due to involvement in interactions
\end{verbatim}

\begin{verbatim}
##  Bdtype.f emmean   SE  df lower.CL upper.CL
##  stacked   102.0 4.53 328     93.1      111
##  tapered    64.7 4.19 328     56.5       73
## 
## Results are averaged over the levels of: Wood.f 
## Confidence level used: 0.95
\end{verbatim}

\begin{Shaded}
\begin{Highlighting}[]
\FunctionTok{contrast}\NormalTok{(mb.emm3, }\AttributeTok{method=}\StringTok{"pairwise"}\NormalTok{, }\AttributeTok{adjust=}\StringTok{"tukey"}\NormalTok{)}
\end{Highlighting}
\end{Shaded}

\begin{verbatim}
##  contrast                        estimate   SE  df t.ratio p.value
##  cherry stacked - ash stacked          12 12.8 328   0.935  0.9825
##  cherry stacked - fir stacked          24 13.0 328   1.841  0.5925
##  cherry stacked - oak stacked          28 12.1 328   2.308  0.2926
##  cherry stacked - cherry tapered       40 12.1 328   3.297  0.0239
##  cherry stacked - ash tapered          42 12.1 328   3.461  0.0140
##  cherry stacked - fir tapered          63 12.1 328   5.192 <0.0001
##  cherry stacked - oak tapered          68 12.1 328   5.604 <0.0001
##  ash stacked - fir stacked             12 13.4 328   0.893  0.9866
##  ash stacked - oak stacked             16 12.6 328   1.273  0.9082
##  ash stacked - cherry tapered          28 12.6 328   2.228  0.3377
##  ash stacked - ash tapered             30 12.6 328   2.387  0.2511
##  ash stacked - fir tapered             51 12.6 328   4.058  0.0016
##  ash stacked - oak tapered             56 12.6 328   4.456  0.0003
##  fir stacked - oak stacked              4 12.8 328   0.313  1.0000
##  fir stacked - cherry tapered          16 12.8 328   1.253  0.9151
##  fir stacked - ash tapered             18 12.8 328   1.410  0.8525
##  fir stacked - fir tapered             39 12.8 328   3.054  0.0495
##  fir stacked - oak tapered             44 12.8 328   3.445  0.0147
##  oak stacked - cherry tapered          12 11.8 328   1.013  0.9723
##  oak stacked - ash tapered             14 11.8 328   1.182  0.9367
##  oak stacked - fir tapered             35 11.8 328   2.954  0.0656
##  oak stacked - oak tapered             40 11.8 328   3.376  0.0185
##  cherry tapered - ash tapered           2 11.8 328   0.169  1.0000
##  cherry tapered - fir tapered          23 11.8 328   1.941  0.5236
##  cherry tapered - oak tapered          28 11.8 328   2.363  0.2632
##  ash tapered - fir tapered             21 11.8 328   1.772  0.6393
##  ash tapered - oak tapered             26 11.8 328   2.194  0.3579
##  fir tapered - oak tapered              5 11.8 328   0.422  0.9999
## 
## P value adjustment: tukey method for comparing a family of 8 estimates
\end{verbatim}

\begin{Shaded}
\begin{Highlighting}[]
\FunctionTok{emmip}\NormalTok{(mb.emm3, Wood.f }\SpecialCharTok{\textasciitilde{}}\NormalTok{ Bdtype.f)}
\end{Highlighting}
\end{Shaded}

\pandocbounded{\includegraphics[keepaspectratio]{bookdownproj_STA4211Notes1_files/figure-latex/ch0705-1.pdf}}
In terms of the Wood Types, fir has significantly lower deflection than
cherry, and oak has significantly lower deflection than cherry.

\[\nabla\]

\chapter{Higher Order Models}\label{higher-order-models}

Many experiments include three or more factors. Main effects and two-way
interactions have the same interpretations as in two factor studies. A
three-way interaction is when the patterns of the two-way interactions
between two factors differ at different levels of the other factor(s).

In screening experiments, there can be any number of factors, typically
with two or three levels each. The goal in these experiments is to
determine which factors have the larger effects, then conduct new
experiments focusing on these factors.

\section{Three Factor Models - Mean
Structure}\label{three-factor-models---mean-structure}

In this section, we consider balanced three factor models, with labels
A, B, and C, and with numbers of levels \(a\), \(b\), and \(c\),
respectively. The number of replicates per treatment is \(n\), with an
overall sample size of \(n_T=abcn\). The same adjustments can be made
for unbalanced data as was described in the previous chapter.

The population mean when factor A is at level \(i\), B is at level
\(j\), and C is at level \(k\) is \(\mu_{ijk}\). The marginal and
overall means are given below.

\[
\mu_{ij\bullet}=\frac{\sum_{k=1}^c\mu_{ijk}}{c}  \qquad
\mu_{i\bullet k}=\frac{\sum_{j=1}^b\mu_{ijk}}{b}  \qquad
\mu_{\bullet jk}=\frac{\sum_{i=1}^a\mu_{ijk}}{a}
\] \[
\mu_{i\bullet\bullet}=\frac{\sum_{j=1}^b\sum_{k=1}^c\mu_{ijk}}{bc} \qquad
\mu_{\bullet j\bullet}=\frac{\sum_{i=1}^a\sum_{k=1}^c\mu_{ijk}}{ac} \qquad
\mu_{\bullet\bullet k}=\frac{\sum_{i=1}^a\sum_{j=1}^b\mu_{ijk}}{ab} 
\] \[
\mu_{\bullet\bullet\bullet}=\frac{\sum_{i=1}^a\sum_{j=1}^b\sum_{k=1}^c\mu_{ijk}}{abc}
\]

Main effects and interactions are defined below.

\[
\alpha_i = \mu_{i\bullet\bullet}-\mu_{\bullet\bullet\bullet} \qquad
\beta_j = \mu_{\bullet j\bullet}-\mu_{\bullet\bullet\bullet} \qquad
\gamma_k = \mu_{\bullet\bullet k}-\mu_{\bullet\bullet\bullet} 
\] \[
(\alpha\beta)_{ij}=\mu_{ij\bullet}-\mu_{i\bullet\bullet}-
\mu_{\bullet j\bullet} + \mu_{\bullet\bullet\bullet} \qquad
(\alpha\gamma)_{ik}=\mu_{i\bullet k}-\mu_{i\bullet\bullet}-
\mu_{\bullet\bullet k} + \mu_{\bullet\bullet\bullet}
\] \[
(\beta\gamma)_{jk}=\mu_{\bullet jk}-\mu_{\bullet j \bullet}-
\mu_{\bullet\bullet k} + \mu_{\bullet\bullet\bullet}
\] \[
(\alpha\beta\gamma)_{ijk}=\mu_{ijk}-\mu_{ij\bullet}-\mu_{i\bullet k}-
\mu_{\bullet jk} + \mu_{i\bullet\bullet} + \mu_{\bullet j\bullet} +
\mu_{\bullet\bullet k} - \mu_{\bullet\bullet\bullet}
\]

The algorithm for obtaining the main effects and interaction effects
goes as follows.

\begin{itemize}
\tightlist
\item
  Begin with the mean corresponding to the effect.
\item
  Cover each subscript, one-at-a-time and multiply by \((-1)^1=-1\).
\item
  Cover subscripts two-at-a-time and multiply by \((-1)^2=1\).
\item
  Continue until all subscripts have been covered.
\item
  Coefficients \(\pm 1\) will sum to 0.
\end{itemize}

The following constraints are obtained for the effects and their
estimators.

\[
\sum_{i=1}^a\alpha_i=\sum_{j=1}^b\beta_j=\sum_{k=1}^c\gamma_k=0
\] \[
\sum_{i=1}^a(\alpha\beta)_{ij}=\sum_{j=1}^b(\alpha\beta)_{ij}=
\sum_{i=1}^a(\alpha\gamma)_{ik}=\sum_{k=1}^c(\alpha\gamma)_{ij}=
\sum_{j=1}^b(\beta\gamma)_{jk}=\sum_{k=1}^c(\beta\gamma)_{jk}=0
\] \[
\sum_{i=1}^a(\alpha\beta\gamma)_{ijk}=
\sum_{j=1}^b(\alpha\beta\gamma)_{ijk}=
\sum_{k=1}^c(\alpha\beta\gamma)_{ijk}=0
\]

\textbf{Example 8.1 - Flavonoids in Wine}

The following ``population means'' are based on observed experimental
values to apply the formulas above and have been adjusted to have
integer values for all effects (\citeproc{ref-genova_2012}{Genova et al.
2012}). There were three factors, each at two levels, described below
and the response was total flavonoids.

\begin{itemize}
\tightlist
\item
  Factor A - \textbf{Grape Type} Sangiovese (Red, \(i=1\)) and Muscat
  (White, \(i=2\))
\item
  Factor B - \textbf{Storage Temperature} 4C \((j=1)\) and -20C
  \((j=2)\)
\item
  Factor C - \textbf{Storage Time} 24 hours (1 day, \(k=1\)) and 2 weeks
  (14 days, \(k=2\))
\end{itemize}

The means are given below.

\[
\mbox{Sangiovese:} \quad
\mu_{111}=78 \quad \mu_{112}=68 \quad \mu_{121}=84 \quad \mu_{122}=90
\] \[
\mbox{Muscat:} \quad
\mu_{211}=22 \quad \mu_{212}=12 \quad \mu_{221}=24 \quad \mu_{222}=22
\] \[
\mu_{11\bullet}=\frac{78+68}{2}=73 \qquad 
\mu_{12\bullet}=\frac{84+90}{2}=87
\] \[
\mu_{1\bullet 1}=\frac{78+84}{2}=81 \qquad
\mu_{1\bullet 2}=\frac{68+90}{2}=79
\] \[
\mu_{21\bullet}=\frac{22+12}{2}=17 \qquad 
\mu_{22\bullet}=\frac{24+22}{2}=23
\] \[
\mu_{2\bullet 1}=\frac{22+24}{2}=23 \qquad
\mu_{2\bullet 2}=\frac{12+22}{2}=17
\] \[
\mu_{\bullet 11}=\frac{78+22}{2}=50 \qquad
\mu_{\bullet 12}=\frac{68+12}{2}=40
\] \[
\mu_{\bullet 21}=\frac{84+24}{2}=54 \qquad
\mu_{\bullet 22}=\frac{90+22}{2}=56
\] \[
\mu_{1\bullet\bullet}=\frac{78+68+84+90}{4}=80 \quad
\mu_{2\bullet\bullet}=\frac{22+12+24+22}{4}=20
\] \[
\mu_{\bullet 1\bullet}=\frac{78+68+22+12}{4}=45 \quad
\mu_{\bullet 2\bullet}=\frac{84+90+24+22}{4}=55
\] \[
\mu_{\bullet\bullet 1}=\frac{78+84+22+24}{4}=52 \quad
\mu_{\bullet\bullet 2}=\frac{68+90+12+22}{4}=48
\] \[
\mu_{\bullet\bullet\bullet} =
\frac{78+68+84+90+22+12+24+22}{8}=50
\] The main effects and interaction effects are obtained below.

\[
\mbox{Grape:} \quad \alpha_1=\mu_{1\bullet\bullet}-
\mu_{\bullet\bullet\bullet} = 80-50=30
\qquad
\alpha_2=20-50=-30
\] \[
\mbox{Temperature:} \quad
\beta_1=\mu_{\bullet 1\bullet}-
\mu_{\bullet\bullet\bullet} = 45-50=-5 \qquad
\beta_2=55-50=5
\] \[
\mbox{Time:} \quad
\gamma_1=\mu_{\bullet\bullet1}-
\mu_{\bullet\bullet\bullet} = 52-50=2 \qquad
\gamma_2=48-50=-2
\] \[
\mbox{Grape/Temp:} \quad
(\alpha\beta)_{11}=
\mu_{11\bullet}-\mu_{1\bullet\bullet}-\mu_{\bullet 1 \bullet} +
\mu_{\bullet\bullet\bullet} = 73-80-45+50=-2
\] \[
(\alpha\beta)_{12}=2 \qquad
(\alpha\beta)_{21}=2 \qquad
(\alpha\beta)_{22}=-2
\] \[
\mbox{Grape/Time:} \quad
(\alpha\gamma)_{11}=
\mu_{1\bullet 1}-\mu_{1\bullet\bullet}-\mu_{\bullet\bullet 1} +
\mu_{\bullet\bullet\bullet} = 81-80-52+50=-1
\] \[
(\alpha\gamma)_{12}=1 \qquad
(\alpha\gamma)_{21}=1 \qquad
(\alpha\gamma)_{22}=-1
\] \[
\mbox{Temp/Time:} \quad
(\beta\gamma)_{11}=
\mu_{\bullet  11}-\mu_{\bullet 1\bullet}-\mu_{\bullet\bullet 1} +
\mu_{\bullet\bullet\bullet} = 50-45-52+50=3
\] \[
(\beta\gamma)_{12}=-3 \qquad
(\beta\gamma)_{21}=-3 \qquad
(\beta\gamma)_{22}=3
\] \[
(\alpha\beta_\gamma)_{111}=78-73-81-50+80+45+52-50=1
\] \[
(\alpha\beta_\gamma)_{112}=(\alpha\beta_\gamma)_{121}=
(\alpha\beta_\gamma)_{211}=(\alpha\beta_\gamma)_{222}=-1
\quad
(\alpha\beta_\gamma)_{122}=(\alpha\beta_\gamma)_{212}=
(\alpha\beta_\gamma)_{221}=1
\] Interaction plots are given in Figure @ref(fig:ch0801).

\begin{figure}

{\centering \includegraphics[width=0.8\linewidth]{bookdownproj_STA4211Notes1_files/figure-latex/ch0801-1} 

}

\caption{Interaction plots for flavonoid mean structure}\label{fig:ch0801}
\end{figure}

The interactions are very small. By far the largest effects are red
versus white grapes.

\[\nabla\]

\section{Statistical Model and the Analysis of
Variance}\label{statistical-model-and-the-analysis-of-variance-1}

In this section, the model is defined, least squares estimators are
given, as well as the sums of squares and \(F\)-tests.

The model for \(Y_{ijkl}\) is given below in terms of the parameters
described in the previous section.

\[
Y_{ijkl} = \mu_{\bullet\bullet\bullet} + \alpha_i + \beta_j + \gamma_k +
(\alpha\beta)_{ij} + (\alpha\gamma)_{ik} + (\beta\gamma)_{jk} +
(\alpha\beta\gamma)_{ijk} + \epsilon_{ijkl}
\] \[
\epsilon_{ijkl} \sim N\left(0,\sigma^2\right) \qquad
i=1,\ldots,a; \quad j=1,\ldots,b; \quad k=1,\ldots,c; \quad
l=1,\ldots,n
\]

The parameters and their least squares estimators are given here.

\[
\mbox{Overall Mean:} \quad  \mu_{\bullet\bullet\bullet}
\qquad  \qquad  \hat{\mu}_{\bullet\bullet\bullet}=
\overline{Y}_{\bullet\bullet\bullet\bullet}
\]

\[
\mbox{Factor A Effects:} \quad \alpha_i \qquad
\hat{\alpha}_i = \overline{Y}_{i\bullet\bullet\bullet} -
\overline{Y}_{\bullet\bullet\bullet\bullet} \qquad i=1,\ldots,a
\]

\[
\mbox{Factor B Effects:} \quad \beta_j \qquad
\hat{\beta}_j = \overline{Y}_{\bullet j \bullet\bullet} -
\overline{Y}_{\bullet\bullet\bullet\bullet} \qquad j=1,\ldots,b
\]

\[
\mbox{Factor C Effects:} \quad \gamma_k \qquad
\hat{\gamma}_k = \overline{Y}_{\bullet\bullet k \bullet} -
\overline{Y}_{\bullet\bullet\bullet\bullet} \qquad k=1,\ldots,c
\]

\[
\mbox{AB Interactions:} \quad (\alpha\beta)_{ij} \quad
\hat{(\alpha\beta)}_{ij}=
\overline{Y}_{ij\bullet\bullet} -
\overline{Y}_{i\bullet\bullet\bullet} -
\overline{Y}_{\bullet j \bullet\bullet} +
\overline{Y}_{\bullet\bullet\bullet\bullet}
\]

\[
\mbox{AC Interactions:} \quad (\alpha\gamma)_{ik} \quad
\hat{(\alpha\gamma)}_{ik}=
\overline{Y}_{i\bullet k\bullet} -
\overline{Y}_{i\bullet\bullet\bullet} -
\overline{Y}_{\bullet\bullet k\bullet} +
\overline{Y}_{\bullet\bullet\bullet\bullet}
\]

\[
\mbox{BC Interactions:} \quad (\beta\gamma)_{jk} \quad
\hat{(\beta\gamma)}_{ij}=
\overline{Y}_{\bullet jk\bullet} -
\overline{Y}_{\bullet j\bullet\bullet} -
\overline{Y}_{\bullet\bullet k\bullet} +
\overline{Y}_{\bullet\bullet\bullet\bullet}
\]

\[
\mbox{ABC Interactions:} \quad (\alpha\beta\gamma)_{ijk}
\] \[
\hat{(\alpha\beta\gamma)}_{ijk}=
\overline{Y}_{ijk\bullet} -
\overline{Y}_{ij\bullet \bullet} -
\overline{Y}_{i\bullet k\bullet} -
\overline{Y}_{\bullet jk\bullet} +
\overline{Y}_{i\bullet\bullet\bullet} +
\overline{Y}_{\bullet j\bullet\bullet} +
\overline{Y}_{\bullet\bullet k\bullet} -
\overline{Y}_{\bullet\bullet\bullet\bullet}
\]

The predicted values simplify to the following values when the estimated
overall mean and corresponding main effects and interactions are added
together.

\[
\hat{Y}_{ijkl} = \hat{\mu}_{\bullet\bullet\bullet} +
\hat{\alpha}_i + \hat{\beta}_j + \hat{\gamma}_k
+\hat{(\alpha\beta)}_{ij}+\hat{(\alpha\gamma)}_{ik}+
\hat{(\beta\gamma)}_{jk} + \hat{(\alpha\beta\gamma)}_{ijk}=
\overline{Y}_{ijk\bullet}
\] The residuals are \(e_{ijkl}=Y_{ijkl}-\overline{Y}_{ijk\bullet}\).

The sums of squares, degrees of freedom and expected mean squares are
given below. To save space they will be given in terms of effects
estimates.

\[
\mbox{Total:} \quad SSTO=
\sum_{i=1}^a\sum_{j=1}^b\sum_{k=1}^c\sum_{l=1}^n
\left(Y_{ijkl}-\overline{Y}_{\bullet\bullet\bullet\bullet}\right)^2
\qquad df_{TO}=abcn-1
\] \[
\mbox{A:} \quad SSA=bcn\sum_{i=1}^a \hat{\alpha}_i^2 \qquad df_A=a-1 
\qquad E\{MSA\}=\sigma^2 + \frac{bcn\sum_{i=1}^a\alpha_i^2}{a-1}
\] \[
\mbox{B:} \quad SSB=acn\sum_{j=1}^b \hat{\beta}_j^2 \qquad df_B=b-1 
\qquad E\{MSB\}=\sigma^2 + \frac{acn\sum_{j=1}^b\beta_j^2}{b-1}
\] \[
\mbox{C:} \quad SSC=abn\sum_{k=1}^c \hat{\gamma}_k^2 \qquad df_C=c-1 
\qquad E\{MSC\}=\sigma^2 + \frac{abn\sum_{k=1}^c\gamma_k^2}{c-1}
\] \[
\mbox{AB:} \quad SSAB=cn\sum_{i=1}^a\sum_{j=1}^b \hat{(\alpha\beta)}_{ij}^2 \qquad df_{AB}=(a-1)(b-1) 
\] \[
E\{MSAB\}=\sigma^2 + 
\frac{cn\sum_{i=1}^a\sum_{j=1}^b(\alpha\beta)_{ij}^2}{(a-1)(b-1)}
\] \[
\mbox{AC:} \quad SSAC=bn\sum_{i=1}^a\sum_{k=1}^c \hat{(\alpha\gamma)}_{ik}^2 \qquad df_{AC}=(a-1)(c-1) 
\] \[
E\{MSAC\}=\sigma^2 + 
\frac{bn\sum_{i=1}^a\sum_{k=1}^c(\alpha\gamma)_{ik}^2}{(a-1)(c-1)}
\] \[
\mbox{BC:} \quad SSBC=an\sum_{j=1}^b\sum_{k=1}^c \hat{(\beta\gamma)}_{jk}^2 \qquad df_{BC}=(b-1)(c-1) 
\] \[
E\{MSBC\}=\sigma^2 + 
\frac{an\sum_{j=1}^b\sum_{k=1}^c(\beta\gamma)_{jk}^2}{(b-1)(c-1)}
\] \[
\mbox{ABC:} \quad SSABC=n\sum_{i=1}^a\sum_{j=1}^b\sum_{k=1}^c \hat{(\alpha\beta\gamma)}_{ijk}^2 \qquad df_{ABC}=(a-1)(b-1)(c-1) 
\] \[
E\{MSABC\}=\sigma^2 + 
\frac{n\sum_{i=1}^a\sum_{j=1}^b\sum_{k=1}^c
(\alpha\beta\gamma)_{ijk}^2}{(a-1)(b-1)(c-1)}
\]

\[
\mbox{Error:} \quad SSE =
\sum_{i=1}^a\sum_{j=1}^b\sum_{k=1}^c\sum_{l=1}^n
\left(Y_{ijkl}-\overline{Y}_{ijk\bullet}\right)^2
\quad df_{E}=abc(n-1) \quad
E\{MSE\}=\sigma^2
\]

To test for main effects and interactions, the following \(F\)-tests are
obtained (where TS is test statistic and RR is rejection region).

\[
H_0^{ABC}: (\alpha\beta\gamma)_{ijk}=0 \qquad i=1,\ldots,a; \quad 
j=1,\ldots,b; \quad k=1,\ldots,c
\] \[
TS:F_{ABC}^*=\frac{MSABC}{MSE}  \qquad RR:F_{ABC}^* \ge
F_{1-\alpha;(a-1)(b-1)(c-1),abc(n-1)}
\]

\[
H_0^{AB}: (\alpha\beta)_{ij}=0 \qquad i=1,\ldots,a; \quad 
j=1,\ldots,b;
\] \[
TS:F_{AB}^*=\frac{MSAB}{MSE}  \qquad RR:F_{AB}^* \ge
F_{1-\alpha;(a-1)(b-1),abc(n-1)}
\]

\[
H_0^{AC}: (\alpha\gamma)_{ik}=0 \qquad i=1,\ldots,a; \quad 
\quad k=1,\ldots,c
\] \[
TS:F_{AC}^*=\frac{MSAC}{MSE}  \qquad RR:F_{AC}^* \ge
F_{1-\alpha;(a-1)(c-1),abc(n-1)}
\]

\[
H_0^{BC}: (\beta\gamma)_{jk}=0 \qquad  
j=1,\ldots,b; \quad k=1,\ldots,c
\] \[
TS:F_{BC}^*=\frac{MSBC}{MSE}  \qquad RR:F_{BC}^* \ge
F_{1-\alpha;(b-1)(c-1),abc(n-1)}
\]

\[
H_0^{A}: \alpha_i=0 \qquad i=1,\ldots,a
\] \[
TS:F_{A}^*=\frac{MSA}{MSE}  \qquad RR:F_{A}^* \ge
F_{1-\alpha;a-1,abc(n-1)}
\]

\[
H_0^{B}: \beta_j=0 \qquad  
j=1,\ldots,b
\] \[
TS:F_{B}^*=\frac{MSB}{MSE}  \qquad RR:F_{B}^* \ge
F_{1-\alpha;b-1,abc(n-1)}
\]

\[
H_0^{C}: \gamma_k=0 \qquad  k=1,\ldots,c
\] \[
TS:F_{C}^*=\frac{MSC}{MSE}  \qquad RR:F_{C}^* \ge
F_{1-\alpha;c-1,abc(n-1)}
\]

\textbf{Example 8.2 - Finishing Treatments for Chef Jackets}

A study was conducted to observe the effects of three factors for
protecting chef jackets from scald injuries
(\citeproc{ref-devarajan_2017}{Deverajan, McQueen, and Wen 2017}). Three
responses were measured: absorbed energy (kj/m\^{}2) at 3 sensors
(upper, middle, and lower). This example will analyze the upper sensor
measurements. The three factors are as follow, there were \(n=5\)
replicates per treatment combination.

\begin{itemize}
\tightlist
\item
  Factor A \((a=4)\) - \textbf{Finish}: Regular \((i=1)\), Water
  \((i=2)\), Soil \((i=3)\), Teflon \((i=4)\)
\item
  Factor B \((b=2)\) - \textbf{Structure}: Plain \((j=1)\), Twill
  \((j=2)\)
\item
  Factor C \((c=3)\) - \textbf{Number of Layers}: 1 \((k=1)\), 2
  \((k=2)\), 3 \((k=3)\)
\end{itemize}

The R code and output are given below and the interaction plot is given
in Figure @ref(fig:ch0802).

\begin{verbatim}
##   finish structure layers trt absEnrgU absEnrgM absEnrgL
## 1      1         1      1   1  239.846  202.176  204.611
## 2      1         1      1   1  267.275  205.320  231.043
\end{verbatim}

\begin{verbatim}
##     finish structure layers trt absEnrgU absEnrgM absEnrgL
## 119      4         2      3  24  107.863   90.537   80.929
## 120      4         2      3  24  117.774  109.554  104.529
\end{verbatim}

\begin{figure}

{\centering \includegraphics[width=0.8\linewidth]{bookdownproj_STA4211Notes1_files/figure-latex/ch0802-1} 

}

\caption{Interaction plots for chef jacket study}\label{fig:ch0802}
\end{figure}

\begin{verbatim}
## 
## Call:
## aov(formula = absEnrgU ~ finish.f * structure.f * layers.f, data = cjb)
## 
## Residuals:
##     Min      1Q  Median      3Q     Max 
## -59.117  -7.176  -0.843   7.294  84.653 
## 
## Coefficients:
##                                  Estimate Std. Error t value Pr(>|t|)    
## (Intercept)                      224.4625     1.5753 142.486  < 2e-16 ***
## finish.f1                         25.7041     2.7285   9.420 2.62e-15 ***
## finish.f2                         19.0374     2.7285   6.977 3.87e-10 ***
## finish.f3                         28.8043     2.7285  10.557  < 2e-16 ***
## structure.f1                      -2.3958     1.5753  -1.521 0.131593    
## layers.f1                          7.8250     2.2278   3.512 0.000679 ***
## layers.f2                          2.6250     2.2278   1.178 0.241607    
## finish.f1:structure.f1            12.4625     2.7285   4.567 1.47e-05 ***
## finish.f2:structure.f1             2.3958     2.7285   0.878 0.382105    
## finish.f3:structure.f1           -22.5709     2.7285  -8.272 7.52e-13 ***
## finish.f1:layers.f1               -8.0917     3.8587  -2.097 0.038624 *  
## finish.f2:layers.f1               -7.6250     3.8587  -1.976 0.051021 .  
## finish.f3:layers.f1              -13.0416     3.8587  -3.380 0.001050 ** 
## finish.f1:layers.f2               -1.6415     3.8587  -0.425 0.671491    
## finish.f2:layers.f2                4.5249     3.8587   1.173 0.243841    
## finish.f3:layers.f2                4.0582     3.8587   1.052 0.295579    
## structure.f1:layers.f1            -0.7666     2.2278  -0.344 0.731527    
## structure.f1:layers.f2            -0.7667     2.2278  -0.344 0.731477    
## finish.f1:structure.f1:layers.f1  -9.7000     3.8587  -2.514 0.013609 *  
## finish.f2:structure.f1:layers.f1  -3.9333     3.8587  -1.019 0.310607    
## finish.f3:structure.f1:layers.f1  15.5833     3.8587   4.038 0.000108 ***
## finish.f1:structure.f1:layers.f2  -0.5500     3.8587  -0.143 0.886966    
## finish.f2:structure.f1:layers.f2  -2.0833     3.8587  -0.540 0.590524    
## finish.f3:structure.f1:layers.f2   0.5833     3.8587   0.151 0.880169    
## ---
## Signif. codes:  0 '***' 0.001 '**' 0.01 '*' 0.05 '.' 0.1 ' ' 1
## 
## Residual standard error: 17.26 on 96 degrees of freedom
## Multiple R-squared:  0.9055, Adjusted R-squared:  0.8829 
## F-statistic: 40.01 on 23 and 96 DF,  p-value: < 2.2e-16
\end{verbatim}

\begin{verbatim}
## Analysis of Variance Table
## 
## Response: absEnrgU
##                               Df Sum Sq Mean Sq  F value    Pr(>F)    
## finish.f                       3 217854   72618 243.8496 < 2.2e-16 ***
## structure.f                    1    689     689   2.3129 0.1315926    
## layers.f                       2   7093    3546  11.9090 2.399e-05 ***
## finish.f:structure.f           3  21899    7300  24.5126 7.328e-12 ***
## finish.f:layers.f              6  18696    3116  10.4634 6.608e-09 ***
## structure.f:layers.f           2    141      71   0.2368 0.7895699    
## finish.f:structure.f:layers.f  6   7680    1280   4.2983 0.0007001 ***
## Residuals                     96  28589     298                       
## ---
## Signif. codes:  0 '***' 0.001 '**' 0.01 '*' 0.05 '.' 0.1 ' ' 1
\end{verbatim}

\begin{verbatim}
## [1] "diff"  "lwr"   "upr"   "p.adj"
\end{verbatim}

\begin{verbatim}
## [1] 137
\end{verbatim}

With the exception of the Structure main effects and Structure/Layers
interaction effects, all terms are significant. Since the 3-way
interaction is significant, Structure and Structure/Layers should be
kept in the model. For any pair of treatment means, Tukey's HSD is
computed as follows (there are 4(2)(3)=24 means).

\[
MSE = 298 \quad df_E=24(5-1)=96 \qquad
HSD = q_{.95;24,96}\sqrt{\frac{298}{5}}=5.297(7.720)=40.893
\]

There are 137 (out of 24(23)/2=276) pairs of treatments that differ by
more than 40.893.

\chapter{Block Designs}\label{block-designs}

In many experimental settings, it is possible to group experimental
units into homogeneous groups or \textbf{blocks}, allowing each
treatment to be assigned to individual units within each block. It is
often possible to evaluate each treatment on the same individual unit or
subject. The goal is to remove block-to-block variation and obtain more
powerful tests for treatment effects. When individual units receive each
treatment (preferably in random order), these are often referred to as
\textbf{Repeated Measures}.

In other situations, it may be necessary to block on two or three
factors to obtain better precision for studying treatment effects. These
designs are \textbf{Latin Squares} for two blocking factors and
\textbf{Graeco-Latin Squares} for three blocking factors. In many
situations, these designs are used to obtain ``fair'' comparisons of
treatments, even if the factors are not concluded in the analysis. This
helps remove biases that can occur in the implementation of the
experiment, particularly when individual subjects receive each
treatment.

Some examples containing a single blocking factor are given below.

\begin{itemize}
\tightlist
\item
  \textbf{Agriculture} - An experiment to compare \(r\) varieties of
  fertilizer is conducted on \(b\) blocks of land. Within each block,
  \(r\) plots are obtained and one variety of fertilizer is applied to
  each plot.
\item
  \textbf{Industrial} - An experiment to compare \(r\) methods (or
  machines) to manufacture a product is conducted on \(b\) batches of
  raw material. Each batch is broken into \(r\) sub-batches and each
  treatment is applied to one sub-batch within the batch.
\item
  \textbf{Pharmaceutical} - An experiment is conducted to compare
  bioavailability among \(r\) formulations of a drug product, by having
  each product measured on \(b\) subjects. Subjects are given a
  particular sequence, and availability is measured for each
  formulation.
\end{itemize}

In many cases, blocks are simply individual replications of an
experiment on different days or at different locations. In other cases,
individuals may be ordered based on an external measurement
(e.g.~severity, ability) and blocks may be formed by that criteria and
the various treatments may be assigned to the group of individuals
within these blocks.

Block designs can have a single treatment factor or combinations of
multiple treatment factors. The larger the number of treatment levels,
the less homogeneous will be the units within a block. In these cases
\textbf{Incomplete Block Designs} can be utilized.

So far, when considering treatment factors, we have implicitly treated
them as \textbf{fixed factors}, the levels observed in the experiment
being the specific treatments of interest to researchers. In block
designs, the blocks are very often \textbf{random factors}, representing
a sample of all levels that could have been selected. These include:
plots of land, batches of raw material, subjects in bioequivalence
studies, and in virtually all repeated measures settings. The primary
goal in all of these experiments is to compare the treatment effects as
we have done in previous chapters. Fortunately, the \(F\)-test and
methods for making contrasts among treatments are the same whether
blocks are fixed or random.

\section{Randomized Block and Repeated Measures
Designs}\label{randomized-block-and-repeated-measures-designs}

In this section we introduce the model for the Randomized Complete Block
Design. This same model is used for the Repeated Measures Design, when
each subject receives each treatment. This is not the same as another
type of Repeated Measures Design, where each subject receives only one
treatment, but is measured longitudinally at multiple time points.

For this model, each of \(r\) treatments are observed, once within each
of \(b\) blocks or subjects, for a total of \(n_T=rb\) observations. For
this model, when treatments are applied to subjects in some time order,
we are assuming there are no order or carryover effects. More complex
designs can be used when this is the case. In this section, we will use
\textbf{blocks} and \textbf{subjects} interchangeably. In virtually all
cases, subjects are treated as random, except in settings where the only
subjects of interest are included in the experiment.

\textbf{Note: This notation is NOT consistent with Kutner, et al.}

\subsection{Model Structure and
Estimators}\label{model-structure-and-estimators}

\[
Y_{ij} = \mu_{\bullet\bullet} + \tau_i + \beta_j + \epsilon_{ij} \qquad
i=1,\ldots,r; \quad j=1,\ldots,b \qquad \epsilon_{ij} \sim NID\left(0,\sigma^2\right)
\]

We assume fixed treatment effects as before, and allow for either fixed
or random block effects.

\[
\mbox{Overall Mean:} \quad  \mu_{\bullet\bullet}  \qquad \qquad
\hat{\mu}_{\bullet\bullet} = \frac{\sum_{i=1}^r\sum_{j=1}^bY_{ij}}{rb} = 
\overline{Y}_{\bullet\bullet}
\] \[
\mbox{Treatment Effects:} \quad \tau_i = \mu_{i\bullet}-\mu_{\bullet\bullet}
\qquad \hat{\tau}_i = \overline{Y}_{i\bullet} - \overline{Y}_{\bullet\bullet}
\qquad  \sum_{i=1}^r\tau_i=\sum_{i=1}^r\hat{\tau}_i=0
\]

We consider the block effects separately, considering the fixed and
random effects cases.

\[
\mbox{Fixed Blocks:} \quad \beta_j=\mu_{\bullet j}-\mu_{\bullet\bullet}
\qquad \hat{\beta}_j = \overline{Y}_{\bullet j} - \overline{Y}_{\bullet\bullet}
\qquad  \sum_{j=1}^b\beta_j=\sum_{j=1}^b\hat{\beta}_j=0
\] For the fixed block case, we obtain the following model for
\(Y_{ij}\) and the treatment means \(\overline{Y}_{i\bullet}\).

\[
E\left\{Y_{ij}\right\} =\mu_{\bullet\bullet} + \tau_i + \beta_j \qquad
V\left\{Y_{ij}\right\} = \sigma^2\qquad 
\mbox{COV}\left\{Y_{ij},Y_{i'j'}\right\}=0 \quad \forall i\ne i',j\ne j'
\] \[
E\left\{\overline{Y}_{i\bullet}\right\}=
E\left\{\frac{1}{b}\sum_{j=1}^bY_{ij}\right\}=
\frac{1}{b}\sum_{j=1}^b
E\left\{Y_{ij}\right\}=
\frac{1}{b}\sum_{j=1}^b\left(\mu_{\bullet\bullet} + \tau_i + \beta_j\right)=
\mu_{\bullet\bullet} + \tau_i
\] \[
V\left\{\overline{Y}_{i\bullet}\right\}=
V\left\{\frac{1}{b}\sum_{j=1}^bY_{ij}\right\}=
\frac{1}{b^2}\sum_{j=1}^b
V\left\{Y_{ij}\right\}=
\frac{1}{b^2}\sum_{j=1}^b\sigma^2=\frac{\sigma^2}{b}
\]

\[
\mbox{Random Blocks:} \quad
\beta_j \sim NID\left(0,\sigma_{\beta}^2\right)
\qquad  \mbox{COV}\left\{\beta_j,\epsilon_{ij}\right\}=0
\qquad \hat{\beta}_j = \overline{Y}_{\bullet j} - \overline{Y}_{\bullet\bullet}
\qquad  \sum_{j=1}^b\hat{\beta}_j=0
\] For the random block case, we obtain the following model for
\(Y_{ij}\) and the treatment means \(\overline{Y}_{i\bullet}\).

The random block effects and the random errors are assumed independent.
By the structure of the estimators the estimated \(\beta_j\) sum to
zero, but it's the population of \(\beta_j\) that has mean 0. For the
random block case, we obtain the following model for \(Y_{ij}\).

\[
E\left\{Y_{ij}\right\} =\mu_{\bullet\bullet} + \tau_i  \qquad
V\left\{Y_{ij}\right\} = \sigma_{\beta}^2 + \sigma^2 
\qquad 
\mbox{COV}\left\{Y_{ij},Y_{i'j'}\right\}=
\left\{ \begin{array}{r@{\quad:\quad}l}
    \sigma_{\beta}^2+\sigma^2 & i=i',j=j' \\
    \sigma_{\beta}^2 & i\ne i',j=j' \\
    0 & \forall i,i', j\ne j'
    \end{array} \right.
\]

Two measurements (for different treatments) within the same block are
correlated. Measurements across different blocks are independent.

\[
E\left\{Y_{ij}\right\} =\mu_{\bullet\bullet} + \tau_i  \qquad
V\left\{Y_{ij}\right\} = 
\sigma_{\beta}^2+\sigma^2\qquad 
\mbox{COV}\left\{Y_{ij},Y_{i'j}\right\}=\sigma_{\beta}^2 \quad \forall i\ne i'
\]

\[
E\left\{\overline{Y}_{i\bullet}\right\}=
E\left\{\frac{1}{b}\sum_{j=1}^bY_{ij}\right\}=
\frac{1}{b}\sum_{j=1}^b
E\left\{Y_{ij}\right\}=
\frac{1}{b}\sum_{j=1}^b\left(\mu_{\bullet\bullet} + \tau_i\right)=
\mu_{\bullet\bullet} + \tau_i
\] \[
V\left\{\overline{Y}_{i\bullet}\right\}=
V\left\{\frac{1}{b}\sum_{j=1}^bY_{ij}\right\}=
\frac{1}{b^2}\sum_{j=1}^b
V\left\{Y_{ij}\right\} =
\frac{1}{b^2}\sum_{j=1}^b\left[\sigma_{\beta}^2+\sigma^2\right]=\frac{\sigma_{\beta}^2+\sigma^2}{b}
\] The Covariance of two different treatment means
\(\overline{Y}_{i\bullet}\) and \(\overline{Y}_{i'\bullet}\) is obtained
below. \[
\mbox{COV}\left\{\overline{Y}_{i\bullet},\overline{Y}_{i'\bullet}\right\}=
\mbox{COV}\left\{\frac{1}{b}\sum_{j=1}^bY_{ij},\frac{1}{b}\sum_{j=1}^bY_{i'j}\right\}=
\frac{1}{b^2}\mbox{COV}\left\{Y_{i1}+\cdots Y_{ib},Y_{i'1}+\cdots +Y_{i'b}\right\}=
\]

\[=\frac{1}{b^2}\left[b\sigma_{\beta}^2+b(b-1)(0)\right]=
\frac{\sigma_{\beta}^2}{b}
\] The fitted values and residuals are similar to other models described
previously.

\[
\hat{Y}_{ij} = \hat{\mu}_{\bullet\bullet} + \hat{\tau}_i + \hat{\beta_j}
= \overline{Y}_{i\bullet} + \overline{Y}_{\bullet j} - \overline{Y}_{\bullet\bullet}
\qquad  e_{ij} = Y_{ij}-\hat{Y}_{ij} =
Y_{ij}-\overline{Y}_{i\bullet} - \overline{Y}_{\bullet j} + \overline{Y}_{\bullet\bullet}
\]

Before setting up the Analysis of Variance, we state the mean and
variance structure for the treatment means for fixed and random blocks.
In each case \(E\left\{\overline{Y}_{i\bullet}\right\} = \mu+\tau_i\).

\[
\mbox{Fixed Blocks:} \quad 
V\left\{\overline{Y}_{i\bullet}\right\}=\frac{\sigma^2}{b} \qquad 
i\ne i': \quad 
\mbox{COV}\left\{\overline{Y}_{i\bullet},\overline{Y}_{i'\bullet}\right\}=0 \qquad
V\left\{\overline{Y}_{i\bullet}-\overline{Y}_{i'\bullet}\right\}=\frac{2\sigma^2}{b}
\]

\[
\mbox{Random Blocks:} \quad 
V\left\{\overline{Y}_{i\bullet}\right\}=\frac{\sigma_{\beta}^2+\sigma^2}{b} \qquad
i\ne i': \quad
\mbox{COV}\left\{\overline{Y}_{i\bullet},\overline{Y}_{i'\bullet}\right\}=
\frac{\sigma_{\beta}^2}{b}\]

\[
V\left\{\overline{Y}_{i\bullet}-\overline{Y}_{i'\bullet}\right\}=
2\left(\frac{\sigma_{\beta}^2+\sigma^2}{b}\right)-2\left(\frac{\sigma_{\beta}^2}{b}\right)=
\frac{2\sigma^2}{b}
\]

Thus, while the variances of the treatment means differ depending on
whether blocks are fixed or random, the \textbf{variances of the
difference in the treatment means are the same}.

\textbf{Example 9.1 - Sensory Analysis of Soy Sauce Recipes}

A study compared \(r=4\) soy sauce recipes, each judged by a panel of
\(b=8\) trained raters (\citeproc{ref-fidaleo_2012}{Fidaleo et al.
2012}). The \(r=4\) soy sauce recipes are given below, this being
treated as a fixed factor (only recipes of interest).

\begin{itemize}
\tightlist
\item
  \textbf{Recipe 1} - Original Soy Sauce
\item
  \textbf{Recipe 2} - Electrodialyzed desalted Soy Sauce
\item
  \textbf{Recipe 3} - Recipe 1 diluted at electric conductivity = 2.3
  S/m
\item
  \textbf{Recipe 4} - Recipe 2 re-salted to original level
\end{itemize}

In sensory studies such as this, there is a debate among researchers
whether raters (blocks) should be treated as a fixed or random factor.

\begin{itemize}
\tightlist
\item
  \textbf{Fixed Blocks} - These are the only raters of interest to the
  researchers (possibly the only raters employed by this company).
\item
  \textbf{Random Blocks} - The researchers are interested in
  generalizing these findings to a population of trained raters (these 8
  raters being a sample from that population).
\end{itemize}

Based on the sample values reported in this paper, we will construct
population models for the fixed and random blocks cases. These were
ratings on a 4-point scale (all numbers have been multiplied by 10 to
keep variances from being so small). The response these are based on was
appearance/color. Parameter values used are given here.

\[
\mu_{\bullet\bullet}=32 \qquad \tau_1=4 \quad \tau_2=3 \quad \tau_3=-13
\quad \tau_4=6  \qquad \sigma^2=11.56 \quad \sigma = 3.4
\] \[
\mbox{Fixed Blocks:} \quad \beta_1,\ldots, \beta_8 = \mbox{ -5.30 to 5.30 by 1.514} 
\qquad \mbox{Random Blocks:} \quad \sigma_{\beta}^2=13.76 \quad \sigma_{\beta}=3.71
\] The variance among block effects was chosen so that the following
equality holds. This makes the models comparable with respect to
variability in the block effects (see \(E\{MSBL\}\) in the next
sub-section). \[
\sigma_{\beta}^2 = \frac{\sum_{j=1}^8\beta_j^2}{8-1}=\frac{(-5.300)^2+\cdots+(5.300)^2}{7}=13.76
\] We will generate 100000 samples from this model (one for fixed block
effects, one for random block effects). The algorithms are described
here.

\textbf{Fixed Effects}

\begin{itemize}
\tightlist
\item
  Generate \(n_T=rb=4(8)\) pseudo
  \(\epsilon_{11}, \ldots , \epsilon_{48} \sim NID\left(0,3.4^2\right)\)
\item
  Assign
  \(Y_{ij}^F=\mu_{\bullet\bullet} + \tau_i + \beta_j^F + \epsilon_{ij} \qquad
   i=1,\ldots,4; \quad j=1,\ldots,8\)
\item
  Compute and save
  \(\overline{Y}_{1\bullet}^F,\ldots,\overline{Y}_{4\bullet}^F\)
\item
  Obtain (across samples) the mean and variance of
  \(\overline{Y}_{1\bullet}^F,\ldots,\overline{Y}_{4\bullet}^F\)
\item
  Obtain (across samples) the mean and variance of the differences among
  \(\overline{Y}_{1\bullet}^F,\ldots,\overline{Y}_{4\bullet}^F\)
\end{itemize}

\textbf{Random Effects}

\begin{itemize}
\tightlist
\item
  Generate \(n_T=rb=4(8)\) pseudo
  \(\epsilon_{11}, \ldots , \epsilon_{48} \sim NID\left(0,3.4^2\right)\)
\item
  Generate \(b=8\) pseudo
  \(\beta_1^R,\ldots , \beta_8^R \sim NID\left(0,3.71^2\right)\)
\item
  Assign
  \(Y_{ij}^R=\mu_{\bullet\bullet} + \tau_i + \beta_j^R + \epsilon_{ij} \qquad
   i=1,\ldots,4; \quad j=1,\ldots,8\)
\item
  Compute and save
  \(\overline{Y}_{1\bullet}^R,\ldots,\overline{Y}_{4\bullet}^R\)
\item
  Obtain (across samples) the mean and variance of
  \(\overline{Y}_{1\bullet}^R,\ldots,\overline{Y}_{4\bullet}^R\)
\item
  Obtain (across samples) the mean and variance of the differences among
  \(\overline{Y}_{1\bullet}^R,\ldots,\overline{Y}_{4\bullet}^R\)
\end{itemize}

For each model (fixed and random blocks), the theoretical means for the
recipe means and for differences among the means are given below.

\[
 E\{\overline{Y}_{i\bullet}\}=\mu_i=\mu_{\bullet\bullet}+ \tau_i  \qquad
 E\{\overline{Y}_{i\bullet}-\overline{Y}_{i'\bullet}\}=
 \mu_i-\mu_{i'} = \tau_i - \tau_{i'}  
\] For fixed block effects, the theoretical variances of the recipe
means and of the differences among the means are given below.

\[
 \sigma^2\{\overline{Y}_{i\bullet}\}=\frac{\sigma^2}{b}=\frac{11.56}{8}=1.445  \qquad
 \sigma^2\{\overline{Y}_{i\bullet}-\overline{Y}_{i'\bullet}\}=
 \frac{2\sigma^2}{b}=\frac{2(11.56)}{8}=2.90  
\] For random block effects, the theoretical variances of the recipe
means and of the differences among the means are given below.

\[
 \sigma^2\{\overline{Y}_{i\bullet}\}=\frac{\sigma_{\beta}^2+\sigma^2}{b}=\frac{13.76+11.56}{8}=3.165  \qquad
 \sigma^2\{\overline{Y}_{i\bullet}-\overline{Y}_{i'\bullet}\}=
 \frac{2\sigma^2}{b}=\frac{2(11.56)}{8}=2.90  
\]

\begin{verbatim}
## [1] -5.3000000 -3.7857143 -2.2714286 -0.7571429  0.7571429  2.2714286  3.7857143
## [8]  5.3000000
\end{verbatim}

\begin{verbatim}
## [1] 13.75837
\end{verbatim}

\begin{verbatim}
## [1] 3.709227
\end{verbatim}

\begin{verbatim}
##             Mean(Fixed) Mean(Random) Var(Fixed) Var(Random)
## Ybar1            36.006       36.006      1.438       3.167
## Ybar2            34.998       34.997      1.439       3.154
## Ybar3            19.001       19.001      1.441       3.148
## Ybar4            37.997       37.997      1.446       3.157
## Ybar2-Ybar1      -1.009       -1.009      2.867       2.867
## Ybar3-Ybar1     -17.005      -17.005      2.892       2.892
## Ybar3-Ybar2     -15.997      -15.997      2.888       2.888
## Ybar4-Ybar1       1.990        1.990      2.878       2.878
## Ybar4-Ybar2       2.999        2.999      2.874       2.874
## Ybar4-Ybar3      18.996       18.996      2.889       2.889
\end{verbatim}

All of the empirical results are very close to the theoretical results
(as they should be). The difference arises when considering
\(V\left\{\overline{Y}_{i\bullet}\right\}\) when block effects are
generated as fixed effects versus when they are generated as random
effects.

\[\nabla\]

\section{\texorpdfstring{Analysis of Variance and the
\(F\)-test}{Analysis of Variance and the F-test}}\label{analysis-of-variance-and-the-f-test}

We use the sums of squares as previously set up, with the Sum of Squares
Error \(SSE\), technically being the Sum of Squares for the Treatment by
Block Interaction as there is only one replicate per combination of
Treatment and Block. The sums of squares, degrees of freedom, and
expected mean squares are given here. In this model, we assume that
there is no Treatment/Block interaction. Tukey's One-Degree of Freedom
Test for Non-Additivity can be used to check this assumption.

\[
\mbox{Total:} \quad SSTO = \sum_{i=1}^r\sum_{j=1}^b \left(Y_{ij}-\overline{Y}_{\bullet\bullet}\right)^2 \qquad \qquad df_{TO}=rb-1 
\] \[
\mbox{Treatments:} \quad SSTR = 
b\sum_{i=1}^r \left(Y_{i\bullet}-\overline{Y}_{\bullet\bullet}\right)^2
\qquad \qquad df_{TR}=r-1 
\] \[
E\left\{MSTR\right\}=
\sigma^2+\frac{b\sum_{i=1}^r\tau_i^2}{r-1}
\]

\[
\mbox{Blocks:} \quad SSBL = 
r\sum_{j=1}^b \left(Y_{\bullet j}-\overline{Y}_{\bullet\bullet}\right)^2
\qquad \qquad df_{BL}=b-1 
\] \[
\mbox{Fixed Blocks:} \quad
E\left\{MSBL\right\}=
\sigma^2+\frac{r\sum_{j=1}^b\beta_j^2}{b-1}  \qquad
\mbox{Random Blocks:} \quad
E\left\{MSBL\right\}=
\sigma^2+ r\sigma_{\beta}^2
\]

\[
\mbox{Error (TRxBL):} \quad
SSE = \sum_{i=1}^a\sum_{j=1}^b\left(Y_{ij}-\overline{Y}_{i\bullet} - \overline{Y}_{\bullet j} + \overline{Y}_{\bullet\bullet}\right)^2 \]
\[df_E=(r-1)(b-1) \qquad
E\left\{MSE\right\}=
\sigma^2
\]

To test for treatment effects, we compute the ratio of \(MSTR\) to
\(MSE\) with the \(F\)-test. Software will automatically compute the
\(F\)-statistic for blocks as well, but the primary interest is
typically treatment effects. However, when blocks are random it may be
of interest to estimate \(\sigma_{\beta}^2\), the variance of the block
effects.

\[
H_0: \tau_1=\ldots = \tau_r=0 \qquad
TS: F_{TR}^*=\frac{MSTR}{MSE} \qquad
RR: F_{TR}^* \ge F_{1-\alpha;r-1,(r-1)(b-1)}
\]

Any contrasts or pairwise comparisons can be carried out as in the
Completely Randomized Design. For Tukey's HSD for all possible
comparisons, we compute HSD, which will be the same for all pairs
\(\left(n_i=b\right)\) and can obtain simultaneous \((1-\alpha)100\%\)
Confidence Intervals for \(\tau_i - \tau_{i`}\) as follows.

\[
HSD = q_{1-\alpha;r,(r-1)(b-1)}\sqrt{\frac{MSE}{b}}  
\] \[
(1-\alpha)100\% \mbox{ CI for } \tau_i-\tau_{i'}: \qquad
\left(\overline{Y}_{i\bullet}-\overline{Y}_{i'\bullet}\right) \pm HSD
\]

\textbf{Example 9.2 - Methods of Presenting Weather Information to
Pilots}

A study compared \(r=3\) cockpit weather displays (treatments) for
pilots in terms of information recall (\citeproc{ref-ohare_2009}{O'Hare
and Stenhouse 2009}) There were \(b=23\) pilots (blocks) who were
exposed to each display on a flight simulator. The three displays are
described below. The response \(Y\), is the percent of information
correctly recalled by the pilot.

\begin{itemize}
\tightlist
\item
  Ordinary text - English phrases \((i=1)\)
\item
  Redesigned graphical display \((i=2)\)
\item
  Old graphical display \((i=3)\)
\end{itemize}

Using R, we will set up the Analysis of Variance directly, then use the
\textbf{aov} function along with \textbf{TukeyHSD} function to confirm
results.

\begin{Shaded}
\begin{Highlighting}[]
\NormalTok{pwd }\OtherTok{\textless{}{-}} \FunctionTok{read.csv}\NormalTok{(}\StringTok{"data/pilot\_weatherdisplay.csv"}\NormalTok{)}
\FunctionTok{head}\NormalTok{(pwd,}\DecValTok{3}\NormalTok{)}
\end{Highlighting}
\end{Shaded}

\begin{verbatim}
##   trt.y blk.y corrRecall
## 1     1     1   45.68415
## 2     1     2   63.76311
## 3     1     3   92.90009
\end{verbatim}

\begin{Shaded}
\begin{Highlighting}[]
\FunctionTok{interaction.plot}\NormalTok{(pwd}\SpecialCharTok{$}\NormalTok{blk.y, pwd}\SpecialCharTok{$}\NormalTok{trt.y, pwd}\SpecialCharTok{$}\NormalTok{corrRecall,}
                 \AttributeTok{col=}\FunctionTok{c}\NormalTok{(}\StringTok{"red"}\NormalTok{, }\StringTok{"blue"}\NormalTok{, }\StringTok{"green4"}\NormalTok{), }\AttributeTok{lty=}\FunctionTok{c}\NormalTok{(}\DecValTok{1}\NormalTok{,}\DecValTok{1}\NormalTok{,}\DecValTok{1}\NormalTok{),}
                 \AttributeTok{xlab=}\StringTok{"subject"}\NormalTok{, }\AttributeTok{ylab=}\StringTok{"correct values"}\NormalTok{)}
\end{Highlighting}
\end{Shaded}

\pandocbounded{\includegraphics[keepaspectratio]{bookdownproj_STA4211Notes1_files/figure-latex/ch0902-1.pdf}}

\begin{Shaded}
\begin{Highlighting}[]
\DocumentationTok{\#\# Direct calculations}
\NormalTok{all.mean }\OtherTok{\textless{}{-}} \FunctionTok{mean}\NormalTok{(pwd}\SpecialCharTok{$}\NormalTok{corrRecall)}
\NormalTok{trt.mean }\OtherTok{\textless{}{-}} \FunctionTok{as.vector}\NormalTok{(}\FunctionTok{tapply}\NormalTok{(pwd}\SpecialCharTok{$}\NormalTok{corrRecall, pwd}\SpecialCharTok{$}\NormalTok{trt.y, mean))}
\NormalTok{blk.mean }\OtherTok{\textless{}{-}} \FunctionTok{as.vector}\NormalTok{(}\FunctionTok{tapply}\NormalTok{(pwd}\SpecialCharTok{$}\NormalTok{corrRecall, pwd}\SpecialCharTok{$}\NormalTok{blk.y, mean))}
\NormalTok{r }\OtherTok{\textless{}{-}} \FunctionTok{length}\NormalTok{(trt.mean)}
\NormalTok{b }\OtherTok{\textless{}{-}} \FunctionTok{length}\NormalTok{(blk.mean)}

\NormalTok{SSTO }\OtherTok{\textless{}{-}} \FunctionTok{sum}\NormalTok{((pwd}\SpecialCharTok{$}\NormalTok{corrRecall }\SpecialCharTok{{-}}\NormalTok{ all.mean)}\SpecialCharTok{\^{}}\DecValTok{2}\NormalTok{); df\_TO }\OtherTok{\textless{}{-}}\NormalTok{ r}\SpecialCharTok{*}\NormalTok{b}\DecValTok{{-}1}
\NormalTok{SSTR }\OtherTok{\textless{}{-}}\NormalTok{ b}\SpecialCharTok{*}\FunctionTok{sum}\NormalTok{((trt.mean}\SpecialCharTok{{-}}\NormalTok{all.mean)}\SpecialCharTok{\^{}}\DecValTok{2}\NormalTok{); df\_TR }\OtherTok{\textless{}{-}}\NormalTok{ r}\DecValTok{{-}1}
\NormalTok{MSTR }\OtherTok{\textless{}{-}}\NormalTok{ SSTR }\SpecialCharTok{/}\NormalTok{ df\_TR}
\NormalTok{SSBL }\OtherTok{\textless{}{-}}\NormalTok{ r}\SpecialCharTok{*}\FunctionTok{sum}\NormalTok{((blk.mean}\SpecialCharTok{{-}}\NormalTok{all.mean)}\SpecialCharTok{\^{}}\DecValTok{2}\NormalTok{); df\_BL }\OtherTok{\textless{}{-}}\NormalTok{ b}\DecValTok{{-}1}
\NormalTok{MSBL }\OtherTok{\textless{}{-}}\NormalTok{ SSBL }\SpecialCharTok{/}\NormalTok{ df\_BL}
\NormalTok{SSE }\OtherTok{\textless{}{-}}\NormalTok{ SSTO }\SpecialCharTok{{-}}\NormalTok{ SSTR }\SpecialCharTok{{-}}\NormalTok{ SSBL; df\_E }\OtherTok{\textless{}{-}}\NormalTok{ df\_TO }\SpecialCharTok{{-}}\NormalTok{ df\_TR }\SpecialCharTok{{-}}\NormalTok{ df\_BL}
\NormalTok{MSE }\OtherTok{\textless{}{-}}\NormalTok{ SSE }\SpecialCharTok{/}\NormalTok{ df\_E}
\NormalTok{Fstar }\OtherTok{=}\NormalTok{ MSTR }\SpecialCharTok{/}\NormalTok{ MSE}

\NormalTok{aov.out1 }\OtherTok{\textless{}{-}} \FunctionTok{cbind}\NormalTok{(df\_TR, SSTR, MSTR, Fstar, }\FunctionTok{qf}\NormalTok{(.}\DecValTok{95}\NormalTok{,df\_TR, df\_E),}
                  \DecValTok{1}\SpecialCharTok{{-}}\FunctionTok{pf}\NormalTok{(Fstar,df\_TR, df\_E))}
\NormalTok{aov.out2 }\OtherTok{\textless{}{-}} \FunctionTok{cbind}\NormalTok{(df\_BL, SSBL, MSBL, MSBL}\SpecialCharTok{/}\NormalTok{MSE, }\FunctionTok{df}\NormalTok{(.}\DecValTok{95}\NormalTok{,df\_BL,df\_E),}
                  \DecValTok{1}\SpecialCharTok{{-}}\FunctionTok{pf}\NormalTok{(MSBL}\SpecialCharTok{/}\NormalTok{MSE, df\_BL, df\_E))}
\NormalTok{aov.out3 }\OtherTok{\textless{}{-}} \FunctionTok{cbind}\NormalTok{(df\_E, SSE, MSE, }\ConstantTok{NA}\NormalTok{, }\ConstantTok{NA}\NormalTok{, }\ConstantTok{NA}\NormalTok{)}
\NormalTok{aov.out4 }\OtherTok{\textless{}{-}} \FunctionTok{cbind}\NormalTok{(df\_TO, SSTO, }\ConstantTok{NA}\NormalTok{, }\ConstantTok{NA}\NormalTok{, }\ConstantTok{NA}\NormalTok{, }\ConstantTok{NA}\NormalTok{)}
\NormalTok{aov.out }\OtherTok{\textless{}{-}} \FunctionTok{rbind}\NormalTok{(aov.out1, aov.out2, aov.out3, aov.out4)}
\FunctionTok{colnames}\NormalTok{(aov.out) }\OtherTok{\textless{}{-}} \FunctionTok{c}\NormalTok{(}\StringTok{"df"}\NormalTok{, }\StringTok{"SS"}\NormalTok{, }\StringTok{"MS"}\NormalTok{, }\StringTok{"F*"}\NormalTok{, }\StringTok{"F(.95)"}\NormalTok{, }\StringTok{"P(\textgreater{}F*)"}\NormalTok{)}
\FunctionTok{rownames}\NormalTok{(aov.out) }\OtherTok{\textless{}{-}} \FunctionTok{c}\NormalTok{(}\StringTok{"Treatments"}\NormalTok{, }\StringTok{"Blocks"}\NormalTok{, }\StringTok{"Error"}\NormalTok{, }\StringTok{"Total"}\NormalTok{)}
\FunctionTok{round}\NormalTok{(aov.out,}\DecValTok{4}\NormalTok{)}
\end{Highlighting}
\end{Shaded}

\begin{verbatim}
##            df       SS        MS     F* F(.95) P(>F*)
## Treatments  2  1582.86  791.4300 3.3300 3.2093  0.045
## Blocks     22 42814.57 1946.1167 8.1884 1.1164  0.000
## Error      44 10457.33  237.6667     NA     NA     NA
## Total      68 54854.76        NA     NA     NA     NA
\end{verbatim}

\begin{Shaded}
\begin{Highlighting}[]
\NormalTok{HSD }\OtherTok{\textless{}{-}} \FunctionTok{qtukey}\NormalTok{(.}\DecValTok{95}\NormalTok{, r, df\_E) }\SpecialCharTok{*} \FunctionTok{sqrt}\NormalTok{(MSE}\SpecialCharTok{/}\NormalTok{b)}
\NormalTok{Tukey.out }\OtherTok{\textless{}{-}} \FunctionTok{matrix}\NormalTok{(}\FunctionTok{rep}\NormalTok{(}\DecValTok{0}\NormalTok{,}\DecValTok{10}\SpecialCharTok{*}\NormalTok{r}\SpecialCharTok{*}\NormalTok{(r}\DecValTok{{-}1}\NormalTok{)}\SpecialCharTok{/}\DecValTok{2}\NormalTok{),}\AttributeTok{ncol=}\DecValTok{10}\NormalTok{)}

\NormalTok{row.count }\OtherTok{\textless{}{-}} \DecValTok{0}
\ControlFlowTok{for}\NormalTok{ (i1 }\ControlFlowTok{in} \DecValTok{2}\SpecialCharTok{:}\NormalTok{r) \{}
  \ControlFlowTok{for}\NormalTok{ (i2 }\ControlFlowTok{in} \DecValTok{1}\SpecialCharTok{:}\NormalTok{(i1}\DecValTok{{-}1}\NormalTok{)) \{}
\NormalTok{    row.count }\OtherTok{\textless{}{-}}\NormalTok{ row.count }\SpecialCharTok{+} \DecValTok{1}
\NormalTok{    Tukey.out[row.count,}\DecValTok{1}\NormalTok{] }\OtherTok{\textless{}{-}}\NormalTok{ i1}
\NormalTok{    Tukey.out[row.count,}\DecValTok{2}\NormalTok{] }\OtherTok{\textless{}{-}}\NormalTok{ i2}
\NormalTok{    Tukey.out[row.count,}\DecValTok{3}\NormalTok{] }\OtherTok{\textless{}{-}}\NormalTok{ trt.mean[i1]}
\NormalTok{    Tukey.out[row.count,}\DecValTok{4}\NormalTok{] }\OtherTok{\textless{}{-}}\NormalTok{ trt.mean[i2]}
\NormalTok{    Tukey.out[row.count,}\DecValTok{5}\NormalTok{] }\OtherTok{\textless{}{-}}\NormalTok{ trt.mean[i1] }\SpecialCharTok{{-}}\NormalTok{ trt.mean[i2]}
\NormalTok{    Tukey.out[row.count,}\DecValTok{6}\NormalTok{] }\OtherTok{\textless{}{-}} \FunctionTok{sqrt}\NormalTok{(}\DecValTok{2}\SpecialCharTok{*}\NormalTok{MSE}\SpecialCharTok{/}\NormalTok{b)}
\NormalTok{    Tukey.out[row.count,}\DecValTok{7}\NormalTok{] }\OtherTok{\textless{}{-}}\NormalTok{ HSD}
\NormalTok{    Tukey.out[row.count,}\DecValTok{8}\NormalTok{] }\OtherTok{\textless{}{-}}\NormalTok{ Tukey.out[row.count,}\DecValTok{5}\NormalTok{] }\SpecialCharTok{{-}}\NormalTok{ HSD}
\NormalTok{    Tukey.out[row.count,}\DecValTok{9}\NormalTok{] }\OtherTok{\textless{}{-}}\NormalTok{ Tukey.out[row.count,}\DecValTok{5}\NormalTok{] }\SpecialCharTok{+}\NormalTok{ HSD}
\NormalTok{    Tukey.out[row.count,}\DecValTok{10}\NormalTok{] }\OtherTok{\textless{}{-}} 
      \DecValTok{1}\SpecialCharTok{{-}}\FunctionTok{ptukey}\NormalTok{(}\FunctionTok{abs}\NormalTok{(Tukey.out[row.count,}\DecValTok{5}\NormalTok{])}\SpecialCharTok{/}\FunctionTok{sqrt}\NormalTok{(MSE}\SpecialCharTok{/}\NormalTok{b),r,df\_E)}
\NormalTok{  \}}
\NormalTok{\}}
\FunctionTok{colnames}\NormalTok{(Tukey.out) }\OtherTok{\textless{}{-}} \FunctionTok{c}\NormalTok{(}\StringTok{"i"}\NormalTok{, }\StringTok{"i\textquotesingle{}"}\NormalTok{, }\StringTok{"ybar\_i"}\NormalTok{, }\StringTok{"ybar\_i\textquotesingle{}"}\NormalTok{, }\StringTok{"diff"}\NormalTok{, }\StringTok{"Std Err"}\NormalTok{, }\StringTok{"HSD"}\NormalTok{,}
                         \StringTok{"LB"}\NormalTok{, }\StringTok{"UB"}\NormalTok{, }\StringTok{"p adj"}\NormalTok{)}
\FunctionTok{round}\NormalTok{(Tukey.out,}\DecValTok{4}\NormalTok{)}
\end{Highlighting}
\end{Shaded}

\begin{verbatim}
##      i i' ybar_i ybar_i' diff Std Err     HSD       LB      UB  p adj
## [1,] 2  1   63.7    52.3 11.4  4.5461 11.0264   0.3736 22.4264 0.0413
## [2,] 3  1   55.6    52.3  3.3  4.5461 11.0264  -7.7264 14.3264 0.7495
## [3,] 3  2   55.6    63.7 -8.1  4.5461 11.0264 -19.1264  2.9264 0.1875
\end{verbatim}

\begin{Shaded}
\begin{Highlighting}[]
\DocumentationTok{\#\# Using aov and TukeyHSD}

\NormalTok{pwd.mod1 }\OtherTok{\textless{}{-}} \FunctionTok{aov}\NormalTok{(corrRecall }\SpecialCharTok{\textasciitilde{}} \FunctionTok{factor}\NormalTok{(trt.y) }\SpecialCharTok{+} \FunctionTok{factor}\NormalTok{(blk.y), }\AttributeTok{data=}\NormalTok{pwd)}
\FunctionTok{anova}\NormalTok{(pwd.mod1)}
\end{Highlighting}
\end{Shaded}

\begin{verbatim}
## Analysis of Variance Table
## 
## Response: corrRecall
##               Df Sum Sq Mean Sq F value    Pr(>F)    
## factor(trt.y)  2   1583  791.43  3.3300   0.04501 *  
## factor(blk.y) 22  42815 1946.12  8.1884 2.253e-09 ***
## Residuals     44  10457  237.67                      
## ---
## Signif. codes:  0 '***' 0.001 '**' 0.01 '*' 0.05 '.' 0.1 ' ' 1
\end{verbatim}

\begin{Shaded}
\begin{Highlighting}[]
\NormalTok{HSD.trt.y }\OtherTok{\textless{}{-}} \FunctionTok{TukeyHSD}\NormalTok{(pwd.mod1, }\StringTok{"factor(trt.y)"}\NormalTok{)}
\NormalTok{HSD.trt.y}
\end{Highlighting}
\end{Shaded}

\begin{verbatim}
##   Tukey multiple comparisons of means
##     95% family-wise confidence level
## 
## Fit: aov(formula = corrRecall ~ factor(trt.y) + factor(blk.y), data = pwd)
## 
## $`factor(trt.y)`
##     diff         lwr       upr     p adj
## 2-1 11.4   0.3736119 22.426388 0.0413298
## 3-1  3.3  -7.7263883 14.326388 0.7495471
## 3-2 -8.1 -19.1263883  2.926388 0.1874886
\end{verbatim}

\begin{Shaded}
\begin{Highlighting}[]
\FunctionTok{plot}\NormalTok{(HSD.trt.y, }\AttributeTok{las=}\DecValTok{2}\NormalTok{, }\AttributeTok{col=}\StringTok{"green4"}\NormalTok{)}
\end{Highlighting}
\end{Shaded}

\pandocbounded{\includegraphics[keepaspectratio]{bookdownproj_STA4211Notes1_files/figure-latex/ch0902-2.pdf}}
Thus, there are differences among the three cockpit weather displays. In
particular, pilots had higher recall with the redesigned graphical
display \((i=2)\) than the ordinary text display \((i=1)\).

\[\nabla\]

\subsection{Checking Model
Assumptions}\label{checking-model-assumptions}

Some methods for checking model assumptions are given in this
sub-section, involving the error terms for the model.

\begin{itemize}
\tightlist
\item
  Stripchart of residuals versus blocks to visually observe whether the
  variance of errors is constant across blocks (each block received each
  treatment)
\item
  Plots of residuals versus fitted values to observe whether residual
  variance is related to mean
\item
  Plot of residuals versus time order when experiment is conducted
  sequentially or when blocks are days to observe whether serial
  correlation is present
\item
  Block x Treatment Interactions - Can be tested using Tukey's ODOFNA
  test. If significant, the remainder \(MSE^*\) can be used after
  removing the 1 degree of freedom \(SS\) for interaction
\end{itemize}

\textbf{Example 9.3 - Methods of Presenting Weather Information to
Pilots}

Here we give examples of the model checks, with the exception of
residuals versus time order, as that does not pertain to this study.

\begin{verbatim}
##   trt.y blk.y corrRecall
## 1     1     1   45.68415
## 2     1     2   63.76311
## 3     1     3   92.90009
\end{verbatim}

\pandocbounded{\includegraphics[keepaspectratio]{bookdownproj_STA4211Notes1_files/figure-latex/ch0903-1.pdf}}
\pandocbounded{\includegraphics[keepaspectratio]{bookdownproj_STA4211Notes1_files/figure-latex/ch0903-2.pdf}}

\begin{verbatim}
## 
## Tukey test on 5% alpha-level:
## 
## Test statistic: 0.05092 
## Critival value: 4.067 
## The additivity hypothesis cannot be rejected.
\end{verbatim}

While a few blocks (pilots) have little spread in his/her residuals
based on the stripchart, the spread is reasonably equal across the
\(b=23\) pilots. Plus, there are only \(r=3\) residuals per pilot.

The plot of the residuals versus predicted values does not show any
evidence of ``funneling'' out as the fitted values increase, so constant
variance seems reasonable.

Tukey's ODOFNA test gives a very small test statistic, well below its
critical value. This provides no evidence of an interaction between
treatment (display of weather information) and block (pilot).

\[\nabla\]

\subsection{Relative Efficiency and Within-Subject Variance-Covariance
Matrix}\label{relative-efficiency-and-within-subject-variance-covariance-matrix}

In this sub-section, we consider two measures that come up in Randomized
Block and Repeated Measures Designs.

\begin{itemize}
\tightlist
\item
  \textbf{Relative Efficiency} - A measure of the ratio of the error
  variance of the Completely Randomized Design (CRD) to the Randomized
  Block Design (RBD)
\item
  \textbf{Within Subject Variance-Covariance Matrix} - When
  blocks/subjects are random, the measurements within subjects are
  correlated. The assumptions are of equal variances for the \(r\)
  treatments and equal covariances among the \(r(r-1)/2\) pairs of
  treatments.
\end{itemize}

For \textbf{Relative Efficiency}, we label \(\sigma_{CRD}^2\) to be the
error variance for a model based on a Completely Randomized Design and
\(\sigma_{RBD}^2\) the error variance for the Randomized Block Design.
Recall that the goal of the RBD is to reduce the error variance, thus
making more precise estimates of contrasts (smaller standard errors).

We define the following ratios in terms of true (unknown) variances and
estimated variances from the RBD.

\[
E=\frac{\sigma_{CRD}^2}{\sigma_{RBD}^2}    \qquad \qquad
\hat{E}=\frac{s_{CRD}^2}{s_{RBD}^2} = 
\frac{(b-1)MSBL+b(r-1)MSE}{(br-1)MSE} 
\] Note that \((b-1)+b(r-1)=br-1\), so that \(\hat{E}\) is a weighted
average of \(MSBL/MSE\) and 1. The larger the ratio of \(MSBL/MSE\), the
larger will be \(\hat{E}\). Then, experiments with large variability
among blocks relative to error variability are more efficient.

The Relative Efficiency measures the multiple of the number of blocks
from the RBD that would need to be the treatment sample size in a CRD to
have comparable standard errors for contrasts.

Suppose that an experiment run as a RBD had \(b=5\) blocks had a
Relative Efficiency of \(\hat{E}=2\). Then an experiment run as a CRD
would need \(n=2(5)=10\) units per treatments to have comparable
standard errors for treatment contrasts.

Some authors describe a degrees of freedom adjustment for the estimated
Relative Efficiency. In many cases, it makes little difference, the
formula is given below, with \(r\) being the number of treatments and
\(b\) being the number of blocks.

\[
df_{CRD}=r(b-1) \quad df_{RBD}=(r-1)(b-1)  \qquad
\hat{E}^*=
\frac{\left(df_{RBD}+1\right)\left(df_{CRD}+3\right)}
     {\left(df_{RBD}+3\right)\left(df_{CRD}+1\right)}
     \hat{E}
\]

When blocks/subjects are random, then the responses \(Y\) within blocks
are correlated as described at the beginning of the chapter.

\[
V\left\{Y_{ij}\right\}=\sigma_{\beta}^2+\sigma^2 = \sigma_i^2
\qquad\qquad
i\ne i': \quad \mbox{COV}\left\{Y_{ij},Y_{i'j}\right\}=\sigma_{\beta}^2=\sigma_{ii'}
\]

For this model to be appropriate we need the following assumptions.

\[
\sigma_1^2=\cdots =\sigma_r^2 \qquad \sigma_{12}=\cdots =\sigma_{r-1,r}
\] In matrix form, the within-subject Variance-Covariance matrix can be
written as follows, with a similar set of assumptions.

\[
\Sigma = 
\begin{bmatrix}
\sigma_1^2 & \sigma_{12} & \cdots & \sigma_{1r} \\
\sigma_{21} & \sigma_2^2 & \cdots & \sigma_{2r} \\
\vdots & \vdots & \ddots & \vdots \\
\sigma_{r1} & \sigma_{r2} & \cdots & \sigma_r^2 \\
\end{bmatrix}
=
\begin{bmatrix}
\sigma^2 & \sigma_{ii'} & \cdots & \sigma_{ii'} \\
\sigma_{ii'} & \sigma^2 & \cdots & \sigma_{ii'} \\
\vdots & \vdots & \ddots & \vdots \\
\sigma_{ii'} & \sigma_{ii'} & \cdots & \sigma^2 \\
\end{bmatrix}
\] The sample version of the Variance-Covariance matrix is computed as
follows.

\[
S = 
\begin{bmatrix}
s_1^2 & s_{12} & \cdots & s_{1r} \\
s_{21} & s_2^2 & \cdots & s_{2r} \\
\vdots & \vdots & \ddots & \vdots \\
s_{r1} & s_{r2} & \cdots & s_r^2 \\
\end{bmatrix}
\] \[
s_i^2 = 
\frac{\sum_{j=1}^b\left(Y_{ij}-\overline{Y}_{i\bullet}\right)^2}{b-1}
\qquad
s_{ii'} = 
\frac{\sum_{j=1}^b\left(Y_{ij}-\overline{Y}_{i\bullet}\right)
                  \left(Y_{i'j}-\overline{Y}_{i'\bullet}\right)}{b-1}
\] There are tests of whether the Population Variance-Covariance meets
this assumption based on the Sample Variance-Covariance matrix, but are
beyond the scope of this course. A visual inspection of \(S\) should be
able to detect large discrepancies. Other structures (including an
unstructured form) can be implemented using mixed model software
packages. Also, transformations on \(Y\) can be implemented (such as the
Box-Cox transformation).

We will use the \textbf{cov} function by looping through the treatments
to obtain the Variance-Covariance matrix, keeping in mind that
\(\hat{\mbox{COV}}\{Y_{ij},Y_{ij}\}=\hat{V}\{Y_{ij}\}\).

\textbf{Example 9.4 - Methods of Presenting Weather Information to
Pilots}

Here we compute the Relative Efficiency and the Sample
Variance-Covariance matrix for the pilot weather information study based
on direct calculations.

\begin{Shaded}
\begin{Highlighting}[]
\NormalTok{pwd }\OtherTok{\textless{}{-}} \FunctionTok{read.csv}\NormalTok{(}\StringTok{"data/pilot\_weatherdisplay.csv"}\NormalTok{)}
\FunctionTok{head}\NormalTok{(pwd,}\DecValTok{3}\NormalTok{)}
\end{Highlighting}
\end{Shaded}

\begin{verbatim}
##   trt.y blk.y corrRecall
## 1     1     1   45.68415
## 2     1     2   63.76311
## 3     1     3   92.90009
\end{verbatim}

\begin{Shaded}
\begin{Highlighting}[]
\DocumentationTok{\#\# Direct calculations {-} Relative Efficiency}
\NormalTok{all.mean }\OtherTok{\textless{}{-}} \FunctionTok{mean}\NormalTok{(pwd}\SpecialCharTok{$}\NormalTok{corrRecall)}
\NormalTok{trt.mean }\OtherTok{\textless{}{-}} \FunctionTok{as.vector}\NormalTok{(}\FunctionTok{tapply}\NormalTok{(pwd}\SpecialCharTok{$}\NormalTok{corrRecall, pwd}\SpecialCharTok{$}\NormalTok{trt.y, mean))}
\NormalTok{blk.mean }\OtherTok{\textless{}{-}} \FunctionTok{as.vector}\NormalTok{(}\FunctionTok{tapply}\NormalTok{(pwd}\SpecialCharTok{$}\NormalTok{corrRecall, pwd}\SpecialCharTok{$}\NormalTok{blk.y, mean))}
\NormalTok{r }\OtherTok{\textless{}{-}} \FunctionTok{length}\NormalTok{(trt.mean)}
\NormalTok{b }\OtherTok{\textless{}{-}} \FunctionTok{length}\NormalTok{(blk.mean)}

\NormalTok{SSTO }\OtherTok{\textless{}{-}} \FunctionTok{sum}\NormalTok{((pwd}\SpecialCharTok{$}\NormalTok{corrRecall }\SpecialCharTok{{-}}\NormalTok{ all.mean)}\SpecialCharTok{\^{}}\DecValTok{2}\NormalTok{)}
\NormalTok{df\_TO }\OtherTok{\textless{}{-}}\NormalTok{ r}\SpecialCharTok{*}\NormalTok{b}\DecValTok{{-}1}
\NormalTok{SSTR }\OtherTok{\textless{}{-}}\NormalTok{ b}\SpecialCharTok{*}\FunctionTok{sum}\NormalTok{((trt.mean}\SpecialCharTok{{-}}\NormalTok{all.mean)}\SpecialCharTok{\^{}}\DecValTok{2}\NormalTok{)}
\NormalTok{df\_TR }\OtherTok{\textless{}{-}}\NormalTok{ r}\DecValTok{{-}1}
\NormalTok{MSTR }\OtherTok{\textless{}{-}}\NormalTok{ SSTR }\SpecialCharTok{/}\NormalTok{ df\_TR}
\NormalTok{SSBL }\OtherTok{\textless{}{-}}\NormalTok{ r}\SpecialCharTok{*}\FunctionTok{sum}\NormalTok{((blk.mean}\SpecialCharTok{{-}}\NormalTok{all.mean)}\SpecialCharTok{\^{}}\DecValTok{2}\NormalTok{)}
\NormalTok{df\_BL }\OtherTok{\textless{}{-}}\NormalTok{ b}\DecValTok{{-}1}
\NormalTok{MSBL }\OtherTok{\textless{}{-}}\NormalTok{ SSBL }\SpecialCharTok{/}\NormalTok{ df\_BL}
\NormalTok{SSE }\OtherTok{\textless{}{-}}\NormalTok{ SSTO }\SpecialCharTok{{-}}\NormalTok{ SSTR }\SpecialCharTok{{-}}\NormalTok{ SSBL}
\NormalTok{df\_E }\OtherTok{\textless{}{-}}\NormalTok{ df\_TO }\SpecialCharTok{{-}}\NormalTok{ df\_TR }\SpecialCharTok{{-}}\NormalTok{ df\_BL}
\NormalTok{MSE }\OtherTok{\textless{}{-}}\NormalTok{ SSE }\SpecialCharTok{/}\NormalTok{ df\_E}

\NormalTok{df\_CRD }\OtherTok{\textless{}{-}}\NormalTok{ b}\SpecialCharTok{*}\NormalTok{(r}\DecValTok{{-}1}\NormalTok{)}
\NormalTok{df\_RBD }\OtherTok{\textless{}{-}}\NormalTok{ df\_E}

\NormalTok{E.hat }\OtherTok{\textless{}{-}}\NormalTok{ ((b}\DecValTok{{-}1}\NormalTok{)}\SpecialCharTok{*}\NormalTok{MSBL }\SpecialCharTok{+}\NormalTok{ b}\SpecialCharTok{*}\NormalTok{(r}\DecValTok{{-}1}\NormalTok{)}\SpecialCharTok{*}\NormalTok{MSE) }\SpecialCharTok{/}\NormalTok{ ((b}\SpecialCharTok{*}\NormalTok{r}\DecValTok{{-}1}\NormalTok{)}\SpecialCharTok{*}\NormalTok{MSE)}
\NormalTok{E.hat.star }\OtherTok{\textless{}{-}}\NormalTok{  ( ((df\_RBD}\SpecialCharTok{+}\DecValTok{1}\NormalTok{)}\SpecialCharTok{*}\NormalTok{(df\_CRD}\SpecialCharTok{+}\DecValTok{3}\NormalTok{))}\SpecialCharTok{/}\NormalTok{((df\_RBD}\SpecialCharTok{+}\DecValTok{3}\NormalTok{)}\SpecialCharTok{*}\NormalTok{(df\_CRD}\SpecialCharTok{+}\DecValTok{1}\NormalTok{)) ) }\SpecialCharTok{*}\NormalTok{ E.hat}

\NormalTok{E.out }\OtherTok{\textless{}{-}} \FunctionTok{cbind}\NormalTok{(b, r, MSBL, MSE, E.hat, df\_CRD, df\_RBD, E.hat.star)}
\FunctionTok{colnames}\NormalTok{(E.out) }\OtherTok{\textless{}{-}} \FunctionTok{c}\NormalTok{(}\StringTok{"b"}\NormalTok{, }\StringTok{"r"}\NormalTok{, }\StringTok{"MSBL"}\NormalTok{, }\StringTok{"MSE"}\NormalTok{, }\StringTok{"E{-}hat"}\NormalTok{, }\StringTok{"df\_CRD"}\NormalTok{, }\StringTok{"df\_RBD"}\NormalTok{, }\StringTok{"E{-}hat*"}\NormalTok{)}
\FunctionTok{round}\NormalTok{(E.out,}\DecValTok{3}\NormalTok{)}
\end{Highlighting}
\end{Shaded}

\begin{verbatim}
##       b r     MSBL     MSE E-hat df_CRD df_RBD E-hat*
## [1,] 23 3 1946.117 237.667 3.326     46     44   3.32
\end{verbatim}

\begin{Shaded}
\begin{Highlighting}[]
\DocumentationTok{\#\# Setting up the Sample Variance{-}Covariance matrix}

\NormalTok{S.mat }\OtherTok{\textless{}{-}} \FunctionTok{matrix}\NormalTok{(}\FunctionTok{rep}\NormalTok{(}\DecValTok{0}\NormalTok{,r}\SpecialCharTok{\^{}}\DecValTok{2}\NormalTok{), }\AttributeTok{ncol=}\NormalTok{r)}

\ControlFlowTok{for}\NormalTok{ (i1 }\ControlFlowTok{in} \DecValTok{1}\SpecialCharTok{:}\NormalTok{r) \{}
  \ControlFlowTok{for}\NormalTok{ (i2 }\ControlFlowTok{in} \DecValTok{1}\SpecialCharTok{:}\NormalTok{r) \{}
\NormalTok{    S.mat[i1,i2] }\OtherTok{\textless{}{-}} \FunctionTok{cov}\NormalTok{(pwd}\SpecialCharTok{$}\NormalTok{corrRecall[pwd}\SpecialCharTok{$}\NormalTok{trt.y }\SpecialCharTok{==}\NormalTok{ i1],pwd}\SpecialCharTok{$}\NormalTok{corrRecall[pwd}\SpecialCharTok{$}\NormalTok{trt.y }\SpecialCharTok{==}\NormalTok{ i2])}
\NormalTok{  \}}
\NormalTok{\}}

\FunctionTok{round}\NormalTok{(S.mat, }\DecValTok{3}\NormalTok{)}
\end{Highlighting}
\end{Shaded}

\begin{verbatim}
##         [,1]    [,2]    [,3]
## [1,] 975.803 607.759 532.015
## [2,] 607.759 764.172 568.676
## [3,] 532.015 568.676 681.475
\end{verbatim}

The Relative Efficiency is \(\hat{E}=3.32\), which can be interpreted
as, if this had been run as a Completely Randomized Design, with each
pilot being only observed in one condition, there would need to be
\(b\hat{E}=23(3.32)=76.4\approx 77\) pilots per treatment
(\(n_T=3(77)=231\) total) to have comparable standard errors for
contrasts. This is an efficient design.

While the variances range from 681 to 976, and the covariances range
from 532 to 607, there is no strong evidence that the
Variance-Covariance matrix is far from the assumed model (standard
deviations range from 26 to 31).

\[\nabla\]

\section{Latin Square Designs}\label{latin-square-designs}

When there are two blocking factors, a \textbf{Latin Square Design} can
be used. There will be a row blocking factor, a column blocking factor,
and the treatment factor. In a ``true'' Latin Square, each factor will
have the same number of levels \((r)\). When \(r\) is small, this will
cause very low error degrees of freedom
\(\left(df_E=(r-1)(r-2)\right)\), so typically either the row and/or
column blocking factor may be replicated and have \(kr\) levels, where
\(k\) is an integer value.

The key idea in the Latin Square is that each treatment appears once in
each row, and once in each column. This way, we can directly remove row
and column effects to get more precise estimates of the treatment
effects and contrasts among them. An example of a Latin Square is given
below with \(r=6\), where rows and columns are blocking factors and
labels \(A,\ldots,F\) are the treatments.

\begin{longtable}[]{@{}lllllll@{}}
\caption{Latin Square Design (r=6)}\tabularnewline
\toprule\noalign{}
\endfirsthead
\endhead
\bottomrule\noalign{}
\endlastfoot
Row\Col & 1 & 2 & 3 & 4 & 5 & 6 \\
1 & A & B & C & D & E & F \\
2 & B & C & D & E & F & A \\
3 & C & D & E & F & A & B \\
4 & D & E & F & A & B & C \\
5 & E & F & A & B & C & D \\
6 & F & A & B & C & D & E \\
\end{longtable}

Keep in mind that while there are three factors, each with \(r\) levels,
there are only \(r^2\) observations. Once the row and columns have been
indexed, there is only one treatment that appears. This can make writing
the model confusing with an observation being \(Y_{ijk}\), because you
cannot cycle through all three subscripts.

However, there are still \(r\) means for each treatment, row, and
column, each based on \(r\) observations. Let \(T_{kl}\) represent the
measurement for the \(l^{th}\) replicate when treatment \(k\) is
assigned, \(R_{il}\) represent the \(l^{th}\) replicate in row \(i\),
and \(C_{jl}\) represent the \(l^{th}\) replicate in column \(j\).
Further let \(\overline{Y}\) represent the overall mean and \(SSTO\) be
the usual total sum of squares. We obtain the Analysis of Variance as
follows (note that when using software packages, this is a trivial
extension of the Randomized Block Design).

\[
\mbox{Total:} \quad SSTO = \sum_{\mbox{all data}}\left(Y-\overline{Y}\right)^2 \qquad
df_{TO}=r^2-1
\]

\[
\mbox{Treatments:} \quad \overline{T}_k=\frac{\sum_{l=1}^rT_{kl}}{r} \qquad
SSTR=r\sum_{k=1}^r\left(\overline{T}_k-\overline{Y}\right)^2 \qquad df_{TR}=r-1
\] \[
\mbox{Rows:} \quad \overline{R}_i=\frac{\sum_{l=1}^rR_{il}}{r} \qquad
SSR=r\sum_{i=1}^r\left(\overline{R}_i-\overline{Y}\right)^2 \qquad df_{R}=r-1
\] \[
\mbox{Columns:} \quad \overline{C}_j=\frac{\sum_{l=1}^rC_{jl}}{r} \qquad
SSC=r\sum_{j=1}^r\left(\overline{C}_j-\overline{Y}\right)^2 \qquad df_{C}=r-1
\] \[
\mbox{Error:} \quad SSE = SSTO-SSTR-SSR-SSC 
\] \[
df_E =\left(r^2-1\right)-3(r-1)=(r+1)(r-1)-3(r-1)=
(r-1)[(r+1)-3]=(r-1)(r-2)
\] The usual \(F\)-test for treatment effects and contrasts among
treatment levels can be conducted as before. For instance, for Tukey's
HSD, we would compute the following values for the Honest Significant
Difference and for simultaneous \((1-\alpha)100\%\) Confidence
Intervals. As in the RBD, these will all be based on common sample sizes
\((r)\), and \(\tau_k\) is the treatment effect for treatment \(k\).

\[
HSD = q(1-\alpha;r,(r-1)(r-2))\sqrt{\frac{MSE}{r}}
\] \[
(1-\alpha)100\% \mbox { Confidence Interval for } \tau_k-\tau_{k'}:
\left(\overline{T}_k-\overline{T}_{k'}\right) \pm HSD
\]

\textbf{Example 9.5 - Comparison of Medium (Psychic) Readings}

A study compared \(r=5\) Psychic Mediums (treatments) in a Latin Square
experiment (\citeproc{ref-okeeffe_2005}{O'Keeffe and Wiseman 2005}). We
will treat Medium as a fixed factor (only interest is comparing these
specific ones). There were 5 Sitters (row factor) and 5 Sessions (column
factor). Each Medium read each Sitter once, and each Medium read in each
Session once. This design would look like the previous example, with one
less row, one less column, and ``F'' removed. We will do computations
directly, then use the \textbf{aov} and \textbf{TukeyHSD} functions for
the analysis. The response was \(Y\), the number of statements by the
Medium on the Sitter.

\begin{Shaded}
\begin{Highlighting}[]
\NormalTok{mr }\OtherTok{\textless{}{-}} \FunctionTok{read.csv}\NormalTok{(}\StringTok{"data/medium\_reading.csv"}\NormalTok{)}
\FunctionTok{head}\NormalTok{(mr)}
\end{Highlighting}
\end{Shaded}

\begin{verbatim}
##   session medium sitter statements rating
## 1       1      1      1         55   3.33
## 2       2      2      1         92   3.72
## 3       3      3      1          6   1.52
## 4       4      4      1         24   3.67
## 5       5      5      1         80   5.24
## 6       2      1      2         62   2.88
\end{verbatim}

\begin{Shaded}
\begin{Highlighting}[]
\DocumentationTok{\#\# Direct Calculation}
\NormalTok{all.mean }\OtherTok{\textless{}{-}} \FunctionTok{mean}\NormalTok{(mr}\SpecialCharTok{$}\NormalTok{statements)}
\NormalTok{trt.mean }\OtherTok{\textless{}{-}} \FunctionTok{as.vector}\NormalTok{(}\FunctionTok{tapply}\NormalTok{(mr}\SpecialCharTok{$}\NormalTok{statements, mr}\SpecialCharTok{$}\NormalTok{medium, mean))}
\NormalTok{row.mean }\OtherTok{\textless{}{-}} \FunctionTok{as.vector}\NormalTok{(}\FunctionTok{tapply}\NormalTok{(mr}\SpecialCharTok{$}\NormalTok{statements, mr}\SpecialCharTok{$}\NormalTok{sitter, mean))}
\NormalTok{col.mean }\OtherTok{\textless{}{-}} \FunctionTok{as.vector}\NormalTok{(}\FunctionTok{tapply}\NormalTok{(mr}\SpecialCharTok{$}\NormalTok{statements, mr}\SpecialCharTok{$}\NormalTok{session, mean))}
\NormalTok{r }\OtherTok{\textless{}{-}} \FunctionTok{length}\NormalTok{(trt.mean)}

\NormalTok{SSTO }\OtherTok{\textless{}{-}} \FunctionTok{sum}\NormalTok{((mr}\SpecialCharTok{$}\NormalTok{statements}\SpecialCharTok{{-}}\NormalTok{all.mean)}\SpecialCharTok{\^{}}\DecValTok{2}\NormalTok{); df\_TO }\OtherTok{\textless{}{-}}\NormalTok{ r}\SpecialCharTok{\^{}}\DecValTok{2{-}1}
\NormalTok{SSTR }\OtherTok{\textless{}{-}}\NormalTok{ r}\SpecialCharTok{*}\FunctionTok{sum}\NormalTok{((trt.mean}\SpecialCharTok{{-}}\NormalTok{all.mean)}\SpecialCharTok{\^{}}\DecValTok{2}\NormalTok{);    df\_TR }\OtherTok{\textless{}{-}}\NormalTok{ r}\DecValTok{{-}1}\NormalTok{; MSTR }\OtherTok{\textless{}{-}}\NormalTok{ SSTR}\SpecialCharTok{/}\NormalTok{df\_TR}
\NormalTok{SSR }\OtherTok{\textless{}{-}}\NormalTok{ r}\SpecialCharTok{*}\FunctionTok{sum}\NormalTok{((row.mean}\SpecialCharTok{{-}}\NormalTok{all.mean)}\SpecialCharTok{\^{}}\DecValTok{2}\NormalTok{);    df\_R }\OtherTok{\textless{}{-}}\NormalTok{ r}\DecValTok{{-}1}\NormalTok{; MSR }\OtherTok{\textless{}{-}}\NormalTok{ SSR}\SpecialCharTok{/}\NormalTok{df\_R}
\NormalTok{SSC }\OtherTok{\textless{}{-}}\NormalTok{ r}\SpecialCharTok{*}\FunctionTok{sum}\NormalTok{((col.mean}\SpecialCharTok{{-}}\NormalTok{all.mean)}\SpecialCharTok{\^{}}\DecValTok{2}\NormalTok{);    df\_C }\OtherTok{\textless{}{-}}\NormalTok{ r}\DecValTok{{-}1}\NormalTok{; MSC }\OtherTok{\textless{}{-}}\NormalTok{ SSC}\SpecialCharTok{/}\NormalTok{df\_C}
\NormalTok{SSE }\OtherTok{\textless{}{-}}\NormalTok{ SSTO}\SpecialCharTok{{-}}\NormalTok{SSTR}\SpecialCharTok{{-}}\NormalTok{SSR}\SpecialCharTok{{-}}\NormalTok{SSC; df\_E }\OtherTok{\textless{}{-}}\NormalTok{ (r}\DecValTok{{-}1}\NormalTok{)}\SpecialCharTok{*}\NormalTok{(r}\DecValTok{{-}2}\NormalTok{); MSE }\OtherTok{\textless{}{-}}\NormalTok{ SSE}\SpecialCharTok{/}\NormalTok{df\_E}
\NormalTok{Fstar }\OtherTok{\textless{}{-}}\NormalTok{ MSTR}\SpecialCharTok{/}\NormalTok{MSE}
\NormalTok{pF }\OtherTok{\textless{}{-}} \DecValTok{1}\SpecialCharTok{{-}}\FunctionTok{pf}\NormalTok{(Fstar, df\_TR, df\_E)}

\NormalTok{aov.out1 }\OtherTok{\textless{}{-}} \FunctionTok{cbind}\NormalTok{(df\_TR, SSTR, MSTR, Fstar, }\FunctionTok{qf}\NormalTok{(.}\DecValTok{95}\NormalTok{,df\_TR,df\_E), pF)}
\NormalTok{aov.out2 }\OtherTok{\textless{}{-}} \FunctionTok{cbind}\NormalTok{(df\_R, SSR, MSR, MSR}\SpecialCharTok{/}\NormalTok{MSE, }\FunctionTok{qf}\NormalTok{(.}\DecValTok{95}\NormalTok{,df\_R,df\_E), }\DecValTok{1}\SpecialCharTok{{-}}\FunctionTok{pf}\NormalTok{(MSR}\SpecialCharTok{/}\NormalTok{MSE,df\_R,df\_E))}
\NormalTok{aov.out3 }\OtherTok{\textless{}{-}} \FunctionTok{cbind}\NormalTok{(df\_C, SSC, MSC, MSC}\SpecialCharTok{/}\NormalTok{MSE, }\FunctionTok{qf}\NormalTok{(.}\DecValTok{95}\NormalTok{,df\_C,df\_E), }\DecValTok{1}\SpecialCharTok{{-}}\FunctionTok{pf}\NormalTok{(MSC}\SpecialCharTok{/}\NormalTok{MSE,df\_C,df\_E))}
\NormalTok{aov.out4 }\OtherTok{\textless{}{-}} \FunctionTok{cbind}\NormalTok{(df\_E, SSE, MSE, }\ConstantTok{NA}\NormalTok{, }\ConstantTok{NA}\NormalTok{, }\ConstantTok{NA}\NormalTok{)}
\NormalTok{aov.out5 }\OtherTok{\textless{}{-}} \FunctionTok{cbind}\NormalTok{(df\_TO, SSTO, }\ConstantTok{NA}\NormalTok{, }\ConstantTok{NA}\NormalTok{, }\ConstantTok{NA}\NormalTok{, }\ConstantTok{NA}\NormalTok{)}
\NormalTok{aov.out }\OtherTok{\textless{}{-}} \FunctionTok{rbind}\NormalTok{(aov.out1,aov.out2,aov.out3,aov.out4,aov.out5)}
\FunctionTok{colnames}\NormalTok{(aov.out) }\OtherTok{\textless{}{-}} \FunctionTok{c}\NormalTok{(}\StringTok{"df"}\NormalTok{, }\StringTok{"SS"}\NormalTok{, }\StringTok{"MS"}\NormalTok{, }\StringTok{"F*"}\NormalTok{, }\StringTok{"F(.95)"}\NormalTok{, }\StringTok{"P(\textgreater{}F*)"}\NormalTok{)}
\FunctionTok{rownames}\NormalTok{(aov.out) }\OtherTok{\textless{}{-}} \FunctionTok{c}\NormalTok{(}\StringTok{"Medium"}\NormalTok{, }\StringTok{"Sitter"}\NormalTok{, }\StringTok{"Session"}\NormalTok{, }\StringTok{"Error"}\NormalTok{, }\StringTok{"Total"}\NormalTok{)}
\FunctionTok{round}\NormalTok{(aov.out,}\DecValTok{3}\NormalTok{)}
\end{Highlighting}
\end{Shaded}

\begin{verbatim}
##         df       SS      MS     F* F(.95) P(>F*)
## Medium   4 17280.56 4320.14 55.006  3.259  0.000
## Sitter   4   205.76   51.44  0.655  3.259  0.635
## Session  4   231.76   57.94  0.738  3.259  0.584
## Error   12   942.48   78.54     NA     NA     NA
## Total   24 18660.56      NA     NA     NA     NA
\end{verbatim}

\begin{Shaded}
\begin{Highlighting}[]
\NormalTok{HSD }\OtherTok{\textless{}{-}} \FunctionTok{qtukey}\NormalTok{(.}\DecValTok{95}\NormalTok{,r,df\_E) }\SpecialCharTok{*} \FunctionTok{sqrt}\NormalTok{(MSE}\SpecialCharTok{/}\NormalTok{r)}
\NormalTok{Tukey.out }\OtherTok{\textless{}{-}} \FunctionTok{matrix}\NormalTok{(}\FunctionTok{rep}\NormalTok{(}\DecValTok{0}\NormalTok{,}\DecValTok{8}\SpecialCharTok{*}\NormalTok{r}\SpecialCharTok{*}\NormalTok{(r}\DecValTok{{-}1}\NormalTok{)}\SpecialCharTok{/}\DecValTok{2}\NormalTok{),}\AttributeTok{ncol=}\DecValTok{8}\NormalTok{)}

\NormalTok{row.count }\OtherTok{\textless{}{-}} \DecValTok{0}
\ControlFlowTok{for}\NormalTok{ (i1 }\ControlFlowTok{in} \DecValTok{2}\SpecialCharTok{:}\NormalTok{r) \{}
  \ControlFlowTok{for}\NormalTok{ (i2 }\ControlFlowTok{in} \DecValTok{1}\SpecialCharTok{:}\NormalTok{(i1}\DecValTok{{-}1}\NormalTok{)) \{}
\NormalTok{    row.count }\OtherTok{\textless{}{-}}\NormalTok{ row.count }\SpecialCharTok{+} \DecValTok{1}
\NormalTok{    Tukey.out[row.count,}\DecValTok{1}\NormalTok{] }\OtherTok{\textless{}{-}}\NormalTok{ i1}
\NormalTok{    Tukey.out[row.count,}\DecValTok{2}\NormalTok{] }\OtherTok{\textless{}{-}}\NormalTok{ i2}
\NormalTok{    Tukey.out[row.count,}\DecValTok{3}\NormalTok{] }\OtherTok{\textless{}{-}}\NormalTok{ trt.mean[i1]}
\NormalTok{    Tukey.out[row.count,}\DecValTok{4}\NormalTok{] }\OtherTok{\textless{}{-}}\NormalTok{ trt.mean[i2]}
\NormalTok{    Tukey.out[row.count,}\DecValTok{5}\NormalTok{] }\OtherTok{\textless{}{-}}\NormalTok{ trt.mean[i1] }\SpecialCharTok{{-}}\NormalTok{ trt.mean[i2]}
\NormalTok{    Tukey.out[row.count,}\DecValTok{6}\NormalTok{] }\OtherTok{\textless{}{-}}\NormalTok{ Tukey.out[row.count,}\DecValTok{5}\NormalTok{] }\SpecialCharTok{{-}}\NormalTok{ HSD}
\NormalTok{    Tukey.out[row.count,}\DecValTok{7}\NormalTok{] }\OtherTok{\textless{}{-}}\NormalTok{ Tukey.out[row.count,}\DecValTok{5}\NormalTok{] }\SpecialCharTok{+}\NormalTok{ HSD}
\NormalTok{    Tukey.out[row.count,}\DecValTok{8}\NormalTok{] }\OtherTok{\textless{}{-}} 
      \DecValTok{1}\SpecialCharTok{{-}}\FunctionTok{ptukey}\NormalTok{(}\FunctionTok{abs}\NormalTok{(Tukey.out[row.count,}\DecValTok{5}\NormalTok{])}\SpecialCharTok{/}\FunctionTok{sqrt}\NormalTok{(MSE}\SpecialCharTok{/}\NormalTok{r),r,df\_E)}
\NormalTok{  \}}
\NormalTok{\}}

\FunctionTok{colnames}\NormalTok{(Tukey.out) }\OtherTok{\textless{}{-}} \FunctionTok{c}\NormalTok{(}\StringTok{"k"}\NormalTok{, }\StringTok{"k\textquotesingle{}"}\NormalTok{, }\StringTok{"ybar\_i"}\NormalTok{, }\StringTok{"ybar\_i\textquotesingle{}"}\NormalTok{, }\StringTok{"diff"}\NormalTok{,}
                         \StringTok{"LB"}\NormalTok{, }\StringTok{"UB"}\NormalTok{, }\StringTok{"p adj"}\NormalTok{)}
\FunctionTok{round}\NormalTok{(Tukey.out,}\DecValTok{4}\NormalTok{)}
\end{Highlighting}
\end{Shaded}

\begin{verbatim}
##       k k' ybar_i ybar_i'  diff       LB       UB  p adj
##  [1,] 2  1   75.4    59.8  15.6  -2.2656  33.4656 0.0983
##  [2,] 3  1    8.0    59.8 -51.8 -69.6656 -33.9344 0.0000
##  [3,] 3  2    8.0    75.4 -67.4 -85.2656 -49.5344 0.0000
##  [4,] 4  1   23.4    59.8 -36.4 -54.2656 -18.5344 0.0002
##  [5,] 4  2   23.4    75.4 -52.0 -69.8656 -34.1344 0.0000
##  [6,] 4  3   23.4     8.0  15.4  -2.4656  33.2656 0.1042
##  [7,] 5  1   67.2    59.8   7.4 -10.4656  25.2656 0.6848
##  [8,] 5  2   67.2    75.4  -8.2 -26.0656   9.6656 0.6026
##  [9,] 5  3   67.2     8.0  59.2  41.3344  77.0656 0.0000
## [10,] 5  4   67.2    23.4  43.8  25.9344  61.6656 0.0000
\end{verbatim}

\begin{Shaded}
\begin{Highlighting}[]
\DocumentationTok{\#\#\# aov and TukeyHSD functions}

\NormalTok{mr.mod1 }\OtherTok{\textless{}{-}} \FunctionTok{aov}\NormalTok{(statements }\SpecialCharTok{\textasciitilde{}} \FunctionTok{factor}\NormalTok{(medium) }\SpecialCharTok{+} \FunctionTok{factor}\NormalTok{(sitter) }\SpecialCharTok{+} \FunctionTok{factor}\NormalTok{(session), }\AttributeTok{data=}\NormalTok{mr)}
\FunctionTok{anova}\NormalTok{(mr.mod1)}
\end{Highlighting}
\end{Shaded}

\begin{verbatim}
## Analysis of Variance Table
## 
## Response: statements
##                 Df  Sum Sq Mean Sq F value   Pr(>F)    
## factor(medium)   4 17280.6  4320.1 55.0056 1.28e-07 ***
## factor(sitter)   4   205.8    51.4  0.6550   0.6346    
## factor(session)  4   231.8    57.9  0.7377   0.5840    
## Residuals       12   942.5    78.5                     
## ---
## Signif. codes:  0 '***' 0.001 '**' 0.01 '*' 0.05 '.' 0.1 ' ' 1
\end{verbatim}

\begin{Shaded}
\begin{Highlighting}[]
\NormalTok{HSD.med }\OtherTok{\textless{}{-}} \FunctionTok{TukeyHSD}\NormalTok{(mr.mod1, }\StringTok{"factor(medium)"}\NormalTok{)}
\NormalTok{HSD.med}
\end{Highlighting}
\end{Shaded}

\begin{verbatim}
##   Tukey multiple comparisons of means
##     95% family-wise confidence level
## 
## Fit: aov(formula = statements ~ factor(medium) + factor(sitter) + factor(session), data = mr)
## 
## $`factor(medium)`
##      diff        lwr        upr     p adj
## 2-1  15.6  -2.265551  33.465551 0.0983090
## 3-1 -51.8 -69.665551 -33.934449 0.0000069
## 4-1 -36.4 -54.265551 -18.534449 0.0002343
## 5-1   7.4 -10.465551  25.265551 0.6847850
## 3-2 -67.4 -85.265551 -49.534449 0.0000004
## 4-2 -52.0 -69.865551 -34.134449 0.0000066
## 5-2  -8.2 -26.065551   9.665551 0.6026044
## 4-3  15.4  -2.465551  33.265551 0.1042146
## 5-3  59.2  41.334449  77.065551 0.0000016
## 5-4  43.8  25.934449  61.665551 0.0000387
\end{verbatim}

\begin{Shaded}
\begin{Highlighting}[]
\FunctionTok{plot}\NormalTok{(HSD.med, }\AttributeTok{las=}\DecValTok{1}\NormalTok{, }\AttributeTok{cex.axis=}\FloatTok{0.75}\NormalTok{, }\AttributeTok{col=}\StringTok{"blue"}\NormalTok{)}
\end{Highlighting}
\end{Shaded}

\pandocbounded{\includegraphics[keepaspectratio]{bookdownproj_STA4211Notes1_files/figure-latex/ch0906-1.pdf}}

There are highly significant differences among the mediums in terms of
mean number of statements made. Of the 5(4)/2=10 pairs, six are
significantly different based on Tukey's HSD.

\[\nabla\]

\begin{verbatim}
## Warning: package 'lmerTest' was built under R version 4.5.3
\end{verbatim}

\begin{verbatim}
## Loading required package: lme4
\end{verbatim}

\begin{verbatim}
## Loading required package: Matrix
\end{verbatim}

\begin{verbatim}
## 
## Attaching package: 'Matrix'
\end{verbatim}

\begin{verbatim}
## The following objects are masked from 'package:tidyr':
## 
##     expand, pack, unpack
\end{verbatim}

\begin{verbatim}
## 
## Attaching package: 'lmerTest'
\end{verbatim}

\begin{verbatim}
## The following object is masked from 'package:lme4':
## 
##     lmer
\end{verbatim}

\begin{verbatim}
## The following object is masked from 'package:stats':
## 
##     step
\end{verbatim}

\chapter{Random and Mixed Effects
Models}\label{random-and-mixed-effects-models}

In some situations, experiments include random treatment factors. In
these cases, the treatment/group levels observed in the study represent
a sample from a larger population of such levels. Interest is in the
overall mean response, as well as in variation in responses across and
within treatments/groups.

In this chapter, we consider \textbf{1-Way Random Effects},
\textbf{2-Way Random Effects}, \textbf{2-Way Mixed Effects}, and
\textbf{3-Way Mixed Effects} Models.

\section{1-Way Random Effects Model}\label{way-random-effects-model}

This model has a single factor, with levels that represent a sample from
a population of such levels. The goal is to make inferences regarding
the overall mean and variances between groups and within groups. The
statistical model for the balanced case is given below.

\[
Y_{ij} = \mu_i + \epsilon_{ij} \qquad i=1,\ldots,r;\quad j=1,\ldots,n
\] \[
\mu_i \sim N\left(\mu_{\bullet},\sigma_{\mu}^2\right) \mbox{ independent }
\qquad
\epsilon_{ij} \sim N\left(0,\sigma^2\right) \mbox{ independent }
\qquad \{\mu_i\},\{\epsilon_{ij}\} \mbox{ independent }
\] Alternatively, we could formulate the model as follows. \[
Y_{ij} = \mu_{\bullet} + \tau_i + \epsilon_{ij} \qquad i=1,\ldots,r;\quad j=1,\ldots,n
\] \[
\tau_i \sim N\left(0,\sigma_{\tau}^2\right) \mbox{ independent }
\qquad
\epsilon_{ij} \sim N\left(0,\sigma^2\right) \mbox{ independent }
\qquad \{\tau_i\},\{\epsilon_{ij}\} \mbox{ independent }
\] For the original model, \(\mu_{\bullet}\) is the overall population
mean response, \(\sigma_{\mu}^2\) is the between group variance, and
\(\sigma^2\) is the within group variance. The experiment is conducted
as follows.

\begin{itemize}
\tightlist
\item
  Sample \(r\) groups from a population of groups.
\item
  Sample \(n\) units from within each group.
\item
  Observe the \(n_T=rn\) responses.
\end{itemize}

The mean and variance-covariance structure of the observations are given
below.

\[
E\left\{Y_{ij}\right\}=E\left\{\mu_i+\epsilon_{ij}\right\}=\mu_{\bullet}+0=\mu_{\bullet}
\]

\[
V\left\{Y_{ij}\right\}=V\left\{\mu_i+\epsilon_{ij}\right\}=
V\left\{\mu_i\right\} + V\left\{\epsilon_{ij}\right\}
+2\mbox{COV}\left\{\mu_i,\epsilon_{ij}\right\}=\]
\[=\sigma_{\mu}^2+\sigma^2+2(0)
=\sigma_{\mu}^2+\sigma^2=\sigma_Y^2
\] For two measurements within the same group \((j \ne j')\), we have
the following covariance, which differs from the fixed effects case
(where the observations are independent). \[
\mbox{COV}\left\{Y_{ij},Y_{ij'}\right\} = 
\mbox{COV}\left\{\mu_i+\epsilon_{ij},\mu_i+\epsilon_{ij'}\right\} =\]

\[
\mbox{COV}\left\{\mu_i,\mu_i\right\} +
\mbox{COV}\left\{\mu_i,\epsilon_{ij'}\right\} +
\mbox{COV}\left\{\epsilon_{ij},\mu_i\right\} +
\mbox{COV}\left\{\epsilon_{ij},\epsilon_{ij'}\right\} = \sigma_{\mu}^2+0+0+0=\sigma_{\mu}^2
\] \[
\mbox{COV}\left\{Y_{ij},Y_{i'j'}\right\} =0 \quad i \neq i' \quad \forall j,j'
\]

For the case, with \(r=n=2\), the mean and variance-covariance structure
for \(Y\) are given in the following matrix forms.

\[
Y=\begin{bmatrix}
Y_{11} \\ Y_{12} \\ Y_{21} \\ Y_{22} \\
\end{bmatrix} 
\qquad 
E\{Y\} =\begin{bmatrix}
\mu_{\bullet} \\
\mu_{\bullet} \\
\mu_{\bullet} \\
\mu_{\bullet} \\
\end{bmatrix}
\qquad
\sigma^2\{Y\} =\begin{bmatrix}
\sigma_Y^2 & \sigma_{\mu}^2 & 0 & 0 \\
\sigma_{\mu}^2 & \sigma_Y^2 & 0 & 0 \\
0 & 0 & \sigma_Y^2 & \sigma_{\mu}^2 \\
0 & 0 & \sigma_{\mu}^2 & \sigma_Y^2 \\
\end{bmatrix}
\]

The \textbf{Intraclass Correlation Coefficient} measures the proportion
of the total variation in responses that is due to variation in group
means.

\[
\rho_I = \frac{\mbox{COV}\left\{Y_{ij},Y_{ij'}\right\}}{
 \sqrt{V\left\{Y_{ij}\right\}V\left\{Y_{ij'}\right\}}}=
 \frac{\sigma_{\mu}^2}{\sigma_Y^2}
\]

\subsection{Analysis of Variance}\label{analysis-of-variance}

In this section, we set up the Analysis of Variance and derive the
expected mean squares for \(MSTR\) and \(MSE\). Note that the quantities
are the same for the random effects model as for the fixed effects
model, but the expectations differ for the two models.

\[
SSTR = n\sum_{i=1}^r\left(\overline{Y}_{i\bullet}-\overline{Y}_{\bullet\bullet}\right)^2 =
n\sum_{i=1}^r\overline{Y}_{i\bullet}^2-nr\overline{Y}_{\bullet\bullet}^2
\]

\[
SSE=\sum_{i=1}^r\sum_{j=1}^n\left(Y_{ij}-\overline{Y}_{i\bullet}\right)^2 =
\sum_{i=1}^r\sum_{j=1}^n Y_{ij}^2 - n\sum_{i=1}^r\overline{Y}_{i\bullet}^2
\]

\[
E\left\{Y_{ij}\right\} = E\left\{\overline{Y}_{i\bullet}\right\} =
E\left\{\overline{Y}_{\bullet\bullet}\right\} = \mu_{\bullet}
\]

\[
V\left\{Y_{ij}\right\}=\sigma_{\mu}^2+\sigma^2=\sigma_Y^2
\]

\[
V\left\{\overline{Y}_{i\bullet}\right\}
=\frac{1}{n^2}\left[\sum_{j=1}^nV\left\{Y_{ij}\right\} + 
2\sum_{j=1}^{n-1}\sum_{j'=j
+1}^n\mbox{COV}\left\{Y_{ij},Y_{ij'}\right\}\right] =
\] \[
=\frac{1}{n^2}\left[n\left(\sigma_{\mu}^2+\sigma^2\right) +
2\frac{n(n-1)}{2}\sigma_{\mu}^2\right]=
\frac{1}{n}\left[\sigma_{\mu}^2(1+n-1) + \sigma^2\right]=
\frac{n\sigma_{\mu}^2+\sigma^2}{n}
\]

Note that the sample means for the various groups are independent, which
is helpful in deriving the variance of the overall mean.

\[
V\left\{\overline{Y}_{\bullet\bullet}\right\} = 
V\left\{\frac{1}{r}\sum_{i=1}^r\overline{Y}_{i\bullet}\right\}
=\frac{1}{r^2}r\frac{n\sigma_{\mu}^2+\sigma^2}{n}=
\frac{n\sigma_{\mu}^2+\sigma^2}{rn}
\]

Now, we obtain the expected sums of squares and mean squares.

\[
E\left\{Y_{ij}^2\right\} = \sigma_{\mu}^2+\sigma^2 + \mu_{\bullet}^2 
\quad \Rightarrow \quad
\sum_{i=1}^r\sum_{j=1}^n E\left\{Y_{ij}^2\right\} =
rn\left(\sigma_{\mu}^2+\sigma^2 + \mu_{\bullet}^2\right)
\]

\[
E\left\{\overline{Y}_{i\bullet}^2\right\} =\frac{n\sigma_{\mu}^2+\sigma^2}{n}+\mu_{\bullet}^2
\quad \Rightarrow \quad
n\sum_{i=1}^rE\left\{\overline{Y}_{i\bullet}^2\right\}
=nr\sigma_{\mu}^2 + r\sigma^2 + nr\mu_{\bullet}^2
\] \[
E\left\{\overline{Y}_{\bullet\bullet}^2\right\} =
\frac{n\sigma_{\mu}^2+\sigma^2}{rn} + \mu_{\bullet}^2
\quad \Rightarrow \quad
rnE\left\{\overline{Y}_{\bullet\bullet}^2\right\} = n\sigma_{\mu}^2+\sigma^2 +
rn\mu_{\bullet}^2
\]

\[
E\{SSTR\} = \left(nr\sigma_{\mu}^2 + r\sigma^2 + nr\mu_{\bullet}^2\right) -
\left(n\sigma_{\mu}^2+\sigma^2 +
rn\mu_{\bullet}^2\right) = n(r-1)\sigma_{\mu}^2 + (r-1)\sigma^2
\]

\[
E\{SSE\} = \left(rn\left(\sigma_{\mu}^2+\sigma^2 + \mu_{\bullet}^2\right)\right) -
\left(nr\sigma_{\mu}^2 + r\sigma^2 + nr\mu_{\bullet}^2\right) = r(n-1)\sigma^2
\] The expected mean squares are the expected sums of squares divided by
the degrees of freedom with \(df_{TR}=r-1\) and \(df_E=r(n-1)\).

\[
E\{MSTR\}=E\left\{\frac{SSTR}{df_{TR}}\right\} = \sigma^2 + n\sigma_{\mu}^2 
\qquad \qquad
E\{MSE\}=E\left\{\frac{SSE}{df_E}\right\} = \sigma^2
\]

The \(F\)-test is used to test whether the variance of the population of
group means is 0.

\[
H_0:\sigma_{\mu}^2=0 \quad H_A:\sigma_{\mu}^2>0 \qquad TS: F^*=\frac{MSTR}{MSE}
\qquad RR:F^*\geq F_{1-\alpha;r-1,r(n-1)}
\]

\textbf{Example 10.1 - Alpha Acids in Varieties of Beer}

A study was conducted to measure alpha acids in a sample of \(r=10\)
varieties of beer (\citeproc{ref-meilgaard_1960}{Meilgaard 1960}). There
were \(n=10\) replicates per variety. The following R code obtains the
ANOVA based on direct computations and a plot of the measurements is
given in Figure @ref(fig:ch1001).

\begin{Shaded}
\begin{Highlighting}[]
\NormalTok{hops }\OtherTok{\textless{}{-}} \FunctionTok{read.csv}\NormalTok{(}\StringTok{"data/hop\_alphaacid.csv"}\NormalTok{)}
\FunctionTok{head}\NormalTok{(hops,}\DecValTok{3}\NormalTok{)}
\end{Highlighting}
\end{Shaded}

\begin{verbatim}
##   variety replicate alphaAcid
## 1       1         1       708
## 2       1         2       708
## 3       1         3       728
\end{verbatim}

\begin{Shaded}
\begin{Highlighting}[]
\FunctionTok{plot}\NormalTok{(alphaAcid }\SpecialCharTok{\textasciitilde{}} \FunctionTok{factor}\NormalTok{(variety), }\AttributeTok{data=}\NormalTok{hops, }\AttributeTok{col=}\StringTok{"orange"}\NormalTok{)}
\end{Highlighting}
\end{Shaded}

\begin{figure}

{\centering \includegraphics[width=0.8\linewidth]{bookdownproj_STA4211Notes1_files/figure-latex/ch1001-1} 

}

\caption{Interaction plots for chef jacket study}\label{fig:ch1001}
\end{figure}

\begin{Shaded}
\begin{Highlighting}[]
\NormalTok{ybar }\OtherTok{\textless{}{-}} \FunctionTok{mean}\NormalTok{(hops}\SpecialCharTok{$}\NormalTok{alphaAcid)}
\NormalTok{ybar.trt }\OtherTok{\textless{}{-}} \FunctionTok{as.vector}\NormalTok{(}\FunctionTok{tapply}\NormalTok{(hops}\SpecialCharTok{$}\NormalTok{alphaAcid,hops}\SpecialCharTok{$}\NormalTok{variety,mean))}
\NormalTok{var.trt }\OtherTok{\textless{}{-}} \FunctionTok{as.vector}\NormalTok{(}\FunctionTok{tapply}\NormalTok{(hops}\SpecialCharTok{$}\NormalTok{alphaAcid,hops}\SpecialCharTok{$}\NormalTok{variety,var))}
\NormalTok{n.trt }\OtherTok{\textless{}{-}} \FunctionTok{as.vector}\NormalTok{(}\FunctionTok{tapply}\NormalTok{(hops}\SpecialCharTok{$}\NormalTok{alphaAcid,hops}\SpecialCharTok{$}\NormalTok{variety,length))}
\NormalTok{r }\OtherTok{\textless{}{-}} \FunctionTok{length}\NormalTok{(ybar.trt)}

\NormalTok{sstr }\OtherTok{\textless{}{-}} \FunctionTok{sum}\NormalTok{(n.trt}\SpecialCharTok{*}\NormalTok{(ybar.trt}\SpecialCharTok{{-}}\NormalTok{ybar)}\SpecialCharTok{\^{}}\DecValTok{2}\NormalTok{)}
\NormalTok{sse }\OtherTok{\textless{}{-}} \FunctionTok{sum}\NormalTok{((n.trt}\DecValTok{{-}1}\NormalTok{)}\SpecialCharTok{*}\NormalTok{var.trt)}
\NormalTok{df.trt }\OtherTok{\textless{}{-}}\NormalTok{ r}\DecValTok{{-}1}
\NormalTok{df.err }\OtherTok{\textless{}{-}} \FunctionTok{sum}\NormalTok{(n.trt)}\SpecialCharTok{{-}}\NormalTok{r}
\NormalTok{mse }\OtherTok{\textless{}{-}}\NormalTok{ sse}\SpecialCharTok{/}\NormalTok{df.err}
\NormalTok{mstr }\OtherTok{\textless{}{-}}\NormalTok{ sstr}\SpecialCharTok{/}\NormalTok{df.trt}
\NormalTok{f\_star }\OtherTok{\textless{}{-}}\NormalTok{ mstr}\SpecialCharTok{/}\NormalTok{mse}
\NormalTok{p\_f\_star }\OtherTok{\textless{}{-}} \DecValTok{1}\SpecialCharTok{{-}}\FunctionTok{pf}\NormalTok{(f\_star,df.trt,df.err)}

\NormalTok{aov.out1 }\OtherTok{\textless{}{-}} \FunctionTok{cbind}\NormalTok{(df.trt,sstr,mstr,f\_star,}\FunctionTok{qf}\NormalTok{(.}\DecValTok{95}\NormalTok{,df.trt,df.err),p\_f\_star)}
\NormalTok{aov.out2 }\OtherTok{\textless{}{-}} \FunctionTok{cbind}\NormalTok{(df.err,sse,mse,}\ConstantTok{NA}\NormalTok{,}\ConstantTok{NA}\NormalTok{,}\ConstantTok{NA}\NormalTok{)}
\NormalTok{aov.out3 }\OtherTok{\textless{}{-}} \FunctionTok{cbind}\NormalTok{(df.trt}\SpecialCharTok{+}\NormalTok{df.err,sstr}\SpecialCharTok{+}\NormalTok{sse,}\ConstantTok{NA}\NormalTok{,}\ConstantTok{NA}\NormalTok{,}\ConstantTok{NA}\NormalTok{,}\ConstantTok{NA}\NormalTok{)}
\NormalTok{aov.out }\OtherTok{\textless{}{-}} \FunctionTok{rbind}\NormalTok{(aov.out1, aov.out2,aov.out3)}
\FunctionTok{colnames}\NormalTok{(aov.out) }\OtherTok{\textless{}{-}} \FunctionTok{c}\NormalTok{(}\StringTok{"df"}\NormalTok{,}\StringTok{"SS"}\NormalTok{,}\StringTok{"MS"}\NormalTok{,}\StringTok{"F*"}\NormalTok{,}\StringTok{"F(.95)"}\NormalTok{,}\StringTok{"P(\textgreater{}F*)"}\NormalTok{)}
\FunctionTok{rownames}\NormalTok{(aov.out) }\OtherTok{\textless{}{-}} \FunctionTok{c}\NormalTok{(}\StringTok{"Variety"}\NormalTok{, }\StringTok{"Error"}\NormalTok{, }\StringTok{"Total"}\NormalTok{)}
\FunctionTok{round}\NormalTok{(aov.out,}\DecValTok{4}\NormalTok{)}
\end{Highlighting}
\end{Shaded}

\begin{verbatim}
##         df        SS          MS      F* F(.95) P(>F*)
## Variety  9 1463915.9 162657.3211 235.728 1.9856      0
## Error   90   62101.9    690.0211      NA     NA     NA
## Total   99 1526017.8          NA      NA     NA     NA
\end{verbatim}

There is strong evidence that the population means (of all varieties)
are not all equal, that is \(\sigma_{\mu}^2 > 0\).

\[\nabla \]

\subsection{\texorpdfstring{Estimating the Population Mean
\(\mu_{\bullet}\)}{Estimating the Population Mean \textbackslash mu\_\{\textbackslash bullet\}}}\label{estimating-the-population-mean-mu_bullet}

In this section, we consider estimating the population mean response.
The expectations of the individual measurements, group means, and
overall mean are all \(\mu_{\bullet}\). However, the overall mean has
the smallest variance. Note also that
\(E\{MSTR\}=\sigma^2+n\sigma_{\mu}^2\).

\[
\hat{\mu}_{\bullet}=\overline{Y}_{\bullet\bullet}  \qquad
V\left\{\overline{Y}_{\bullet\bullet}\right\}=
\frac{n\sigma_{\mu}^2+\sigma^2}{rn} \qquad
\hat{V}\left\{\overline{Y}_{\bullet\bullet}\right\} =
\frac{MSTR}{rn}
\qquad
\hat{SE}\left\{\overline{Y}_{\bullet\bullet}\right\} =
\sqrt{\frac{MSTR}{rn}}
\]

\[
t=\frac{\overline{Y}_{\bullet\bullet}-\mu_{\bullet}}{\hat{SE}\left\{\overline{Y}_{\bullet\bullet}\right\}} \sim t_{r-1}  \quad \Rightarrow \quad
(1-\alpha)100\% \mbox{ CI for } \mu_{\bullet} \equiv
\overline{Y}_{\bullet\bullet} \pm t_{1-\alpha/2;r-1}\sqrt{\frac{MSTR}{rn}}
\]

\textbf{Example 10.2 - Alpha Acids in Varieties of Beer}

The following R code obtains the point estimate of the overall mean and
95\% Confidence Interval for \(\mu_{\bullet}\) based on direct
computations.

\begin{Shaded}
\begin{Highlighting}[]
\NormalTok{hops }\OtherTok{\textless{}{-}} \FunctionTok{read.csv}\NormalTok{(}\StringTok{"data/hop\_alphaacid.csv"}\NormalTok{)}
\NormalTok{ybar }\OtherTok{\textless{}{-}} \FunctionTok{mean}\NormalTok{(hops}\SpecialCharTok{$}\NormalTok{alphaAcid)}
\NormalTok{ybar.trt }\OtherTok{\textless{}{-}} \FunctionTok{as.vector}\NormalTok{(}\FunctionTok{tapply}\NormalTok{(hops}\SpecialCharTok{$}\NormalTok{alphaAcid,hops}\SpecialCharTok{$}\NormalTok{variety,mean))}
\NormalTok{var.trt }\OtherTok{\textless{}{-}} \FunctionTok{as.vector}\NormalTok{(}\FunctionTok{tapply}\NormalTok{(hops}\SpecialCharTok{$}\NormalTok{alphaAcid,hops}\SpecialCharTok{$}\NormalTok{variety,var))}
\NormalTok{n.trt }\OtherTok{\textless{}{-}} \FunctionTok{as.vector}\NormalTok{(}\FunctionTok{tapply}\NormalTok{(hops}\SpecialCharTok{$}\NormalTok{alphaAcid,hops}\SpecialCharTok{$}\NormalTok{variety,length))}
\NormalTok{r }\OtherTok{\textless{}{-}} \FunctionTok{length}\NormalTok{(ybar.trt)}

\NormalTok{sstr }\OtherTok{\textless{}{-}} \FunctionTok{sum}\NormalTok{(n.trt}\SpecialCharTok{*}\NormalTok{(ybar.trt}\SpecialCharTok{{-}}\NormalTok{ybar)}\SpecialCharTok{\^{}}\DecValTok{2}\NormalTok{)}
\NormalTok{sse }\OtherTok{\textless{}{-}} \FunctionTok{sum}\NormalTok{((n.trt}\DecValTok{{-}1}\NormalTok{)}\SpecialCharTok{*}\NormalTok{var.trt)}
\NormalTok{df.trt }\OtherTok{\textless{}{-}}\NormalTok{ r}\DecValTok{{-}1}
\NormalTok{df.err }\OtherTok{\textless{}{-}} \FunctionTok{sum}\NormalTok{(n.trt)}\SpecialCharTok{{-}}\NormalTok{r}
\NormalTok{mse }\OtherTok{\textless{}{-}}\NormalTok{ sse}\SpecialCharTok{/}\NormalTok{df.err}
\NormalTok{mstr }\OtherTok{\textless{}{-}}\NormalTok{ sstr}\SpecialCharTok{/}\NormalTok{df.trt}

\NormalTok{s.ybar }\OtherTok{\textless{}{-}} \FunctionTok{sqrt}\NormalTok{(mstr}\SpecialCharTok{/}\FunctionTok{sum}\NormalTok{(n.trt))}
\NormalTok{LB }\OtherTok{\textless{}{-}}\NormalTok{ ybar }\SpecialCharTok{{-}} \FunctionTok{qt}\NormalTok{(.}\DecValTok{975}\NormalTok{,df.trt)}\SpecialCharTok{*}\NormalTok{s.ybar}
\NormalTok{UB }\OtherTok{\textless{}{-}}\NormalTok{ ybar }\SpecialCharTok{+} \FunctionTok{qt}\NormalTok{(.}\DecValTok{975}\NormalTok{,df.trt)}\SpecialCharTok{*}\NormalTok{s.ybar}
\NormalTok{mu.ci.out }\OtherTok{\textless{}{-}} \FunctionTok{cbind}\NormalTok{(ybar, s.ybar, df.trt, }\FunctionTok{qt}\NormalTok{(.}\DecValTok{975}\NormalTok{,df.trt), LB, UB)}
\FunctionTok{colnames}\NormalTok{(mu.ci.out) }\OtherTok{\textless{}{-}} \FunctionTok{c}\NormalTok{(}\StringTok{"Estimate"}\NormalTok{,}\StringTok{"Std Err"}\NormalTok{,}\StringTok{"df"}\NormalTok{,}\StringTok{"t(.975)"}\NormalTok{,}\StringTok{"LB"}\NormalTok{, }\StringTok{"UB"}\NormalTok{)}
\FunctionTok{rownames}\NormalTok{(mu.ci.out) }\OtherTok{\textless{}{-}} \FunctionTok{c}\NormalTok{(}\StringTok{"Mean"}\NormalTok{)}
\FunctionTok{round}\NormalTok{(mu.ci.out,}\DecValTok{4}\NormalTok{)}
\end{Highlighting}
\end{Shaded}

\begin{verbatim}
##      Estimate Std Err df t(.975)       LB       UB
## Mean   579.61 40.3308  9  2.2622 488.3754 670.8446
\end{verbatim}

The point estimate is about 580, with a 95\% Confidence Interval of 488
to 671. Even though there were \(rn=10(10)=100\) observations, the
between treatment mean square is very large, yielding a wide confidence
interval.

\[\nabla \]

\subsection{Estimating the Intraclass Correlation
Coefficient}\label{estimating-the-intraclass-correlation-coefficient}

The Intraclass Correlation Coefficient is
\(\rho_I=\sigma_{\mu}^2/\sigma_Y^2\), the ratio of the between group
variance to the total variance. To obtain a Confidence Interval for
\(\rho_I\), we make use of the following distributional results and the
fact that \(SSE\) and \(SSTR\) are independent.

\[
E\{MSE\}=\sigma^2 \qquad \frac{SSE}{\sigma^2}=\frac{r(n-1)MSE}{\sigma^2} \sim \chi_{r(n-1)}^2
\]

\[
E\{MSTR\}=\sigma^2 + n\sigma_{\mu}^2 \qquad
\frac{SSTR}{\sigma^2 + n\sigma_{\mu}^2}=
\frac{(r-1)MSTR}{\sigma^2 + n\sigma_{\mu}^2} \sim \chi_{r-1}^2
\]

\[
\frac{\left[\frac{\frac{SSTR}{\sigma^2 + n\sigma_{\mu}^2}}{r-1}\right]}  {
\left[\frac{\frac{SSE}{\sigma^2}}{r(n-1)}\right]} =
\frac{\left[\frac{MSTR}{\sigma^2 + n\sigma_{\mu}^2}\right]} {
\left[\frac{MSE}{\sigma^2}\right] } = F^*\left(\frac{\sigma^2}{\sigma^2 + n\sigma_{\mu}^2}
\right) \sim F_{r-1,r(n-1)}
\] \[
\quad \Rightarrow \quad
P\left(F_{\alpha/2;r-1,r(n-1)} \leq F^*\left(\frac{\sigma^2}{\sigma^2 + n\sigma_{\mu}^2}
\right) \leq F_{1-\alpha/2;r-1,r(n-1)}\right) = 1-\alpha
\]

The goal is to isolate \(\rho_I=\sigma_{\mu}^2/\sigma_Y^2\) in the
middle of the probability statement. Define \(L\) and \(U\) as follow.

\[
L= \frac{1}{n}\left[\frac{F^*}{F_{1-\alpha/2;r-1,r(n-1)}}-1\right]  \qquad \qquad
U= \frac{1}{n}\left[\frac{F^*}{F_{\alpha/2;r-1,r(n-1)}}-1\right]
\] Then the \((1-\alpha)100\%\) Confidence Interval for \(\rho_I\) is
obtained as follows.

\[
(1-\alpha)100\% \mbox{ Confidence Interval for }
\rho_I=\frac{\sigma_{\mu}^2}{\sigma_{\mu}^2+\sigma^2}: \quad
\left[L^*=\frac{L}{1+L},U^*=\frac{U}{1+U}\right]
\]

\textbf{Example 10.3 - Alpha Acids in Varieties of Beer}

The following R code obtains the 95\% Confidence Interval for \(\rho_I\)
based on direct computations.

\begin{Shaded}
\begin{Highlighting}[]
\NormalTok{hops }\OtherTok{\textless{}{-}} \FunctionTok{read.csv}\NormalTok{(}\StringTok{"data/hop\_alphaacid.csv"}\NormalTok{)}
\FunctionTok{head}\NormalTok{(hops, }\DecValTok{3}\NormalTok{)}
\end{Highlighting}
\end{Shaded}

\begin{verbatim}
##   variety replicate alphaAcid
## 1       1         1       708
## 2       1         2       708
## 3       1         3       728
\end{verbatim}

\begin{Shaded}
\begin{Highlighting}[]
\NormalTok{ybar }\OtherTok{\textless{}{-}} \FunctionTok{mean}\NormalTok{(hops}\SpecialCharTok{$}\NormalTok{alphaAcid)}
\NormalTok{ybar.trt }\OtherTok{\textless{}{-}} \FunctionTok{as.vector}\NormalTok{(}\FunctionTok{tapply}\NormalTok{(hops}\SpecialCharTok{$}\NormalTok{alphaAcid,hops}\SpecialCharTok{$}\NormalTok{variety,mean))}
\NormalTok{var.trt }\OtherTok{\textless{}{-}} \FunctionTok{as.vector}\NormalTok{(}\FunctionTok{tapply}\NormalTok{(hops}\SpecialCharTok{$}\NormalTok{alphaAcid,hops}\SpecialCharTok{$}\NormalTok{variety,var))}
\NormalTok{n.trt }\OtherTok{\textless{}{-}} \FunctionTok{as.vector}\NormalTok{(}\FunctionTok{tapply}\NormalTok{(hops}\SpecialCharTok{$}\NormalTok{alphaAcid,hops}\SpecialCharTok{$}\NormalTok{variety,length))}
\NormalTok{r }\OtherTok{\textless{}{-}} \FunctionTok{length}\NormalTok{(ybar.trt)}
\NormalTok{n }\OtherTok{\textless{}{-}}\NormalTok{ n.trt[}\DecValTok{1}\NormalTok{]}

\NormalTok{sstr }\OtherTok{\textless{}{-}} \FunctionTok{sum}\NormalTok{(n.trt}\SpecialCharTok{*}\NormalTok{(ybar.trt}\SpecialCharTok{{-}}\NormalTok{ybar)}\SpecialCharTok{\^{}}\DecValTok{2}\NormalTok{)}
\NormalTok{sse }\OtherTok{\textless{}{-}} \FunctionTok{sum}\NormalTok{((n.trt}\DecValTok{{-}1}\NormalTok{)}\SpecialCharTok{*}\NormalTok{var.trt)}
\NormalTok{df.trt }\OtherTok{\textless{}{-}}\NormalTok{ r}\DecValTok{{-}1}
\NormalTok{df.err }\OtherTok{\textless{}{-}} \FunctionTok{sum}\NormalTok{(n.trt)}\SpecialCharTok{{-}}\NormalTok{r}
\NormalTok{mse }\OtherTok{\textless{}{-}}\NormalTok{ sse}\SpecialCharTok{/}\NormalTok{df.err}
\NormalTok{mstr }\OtherTok{\textless{}{-}}\NormalTok{ sstr}\SpecialCharTok{/}\NormalTok{df.trt}
\NormalTok{f\_star }\OtherTok{\textless{}{-}}\NormalTok{ mstr}\SpecialCharTok{/}\NormalTok{mse}
\NormalTok{f}\FloatTok{.025} \OtherTok{\textless{}{-}} \FunctionTok{qf}\NormalTok{(.}\DecValTok{025}\NormalTok{,df.trt,df.err)}
\NormalTok{f}\FloatTok{.975} \OtherTok{\textless{}{-}} \FunctionTok{qf}\NormalTok{(.}\DecValTok{975}\NormalTok{,df.trt,df.err)}

\NormalTok{L }\OtherTok{\textless{}{-}}\NormalTok{ (}\DecValTok{1}\SpecialCharTok{/}\NormalTok{n)}\SpecialCharTok{*}\NormalTok{((f\_star}\SpecialCharTok{/}\NormalTok{f}\FloatTok{.975}\NormalTok{) }\SpecialCharTok{{-}} \DecValTok{1}\NormalTok{)}
\NormalTok{U }\OtherTok{\textless{}{-}}\NormalTok{ (}\DecValTok{1}\SpecialCharTok{/}\NormalTok{n)}\SpecialCharTok{*}\NormalTok{((f\_star}\SpecialCharTok{/}\NormalTok{f}\FloatTok{.025}\NormalTok{) }\SpecialCharTok{{-}} \DecValTok{1}\NormalTok{)}

\NormalTok{L\_star }\OtherTok{\textless{}{-}}\NormalTok{ L}\SpecialCharTok{/}\NormalTok{(}\DecValTok{1}\SpecialCharTok{+}\NormalTok{L)}
\NormalTok{U\_star }\OtherTok{\textless{}{-}}\NormalTok{ U}\SpecialCharTok{/}\NormalTok{(}\DecValTok{1}\SpecialCharTok{+}\NormalTok{U)}

\NormalTok{rho.ci.out }\OtherTok{\textless{}{-}} \FunctionTok{cbind}\NormalTok{(f\_star, f}\FloatTok{.025}\NormalTok{, f}\FloatTok{.975}\NormalTok{, L, U, L\_star, U\_star)}
\FunctionTok{colnames}\NormalTok{(rho.ci.out) }\OtherTok{\textless{}{-}} \FunctionTok{c}\NormalTok{(}\StringTok{"F*"}\NormalTok{,}\StringTok{"F(.025)"}\NormalTok{,}\StringTok{"F(.975)"}\NormalTok{,}\StringTok{"L"}\NormalTok{, }\StringTok{"U"}\NormalTok{, }\StringTok{"LB"}\NormalTok{, }\StringTok{"UB"}\NormalTok{)}
\FunctionTok{rownames}\NormalTok{(rho.ci.out) }\OtherTok{\textless{}{-}} \FunctionTok{c}\NormalTok{(}\StringTok{"rho\_I"}\NormalTok{)}
\FunctionTok{round}\NormalTok{(rho.ci.out,}\DecValTok{4}\NormalTok{)}
\end{Highlighting}
\end{Shaded}

\begin{verbatim}
##            F* F(.025) F(.975)       L       U     LB     UB
## rho_I 235.728  0.2932  2.2588 10.3361 80.3097 0.9118 0.9877
\end{verbatim}

\[\nabla \]

\subsection{Estimating the Within Group
Variance}\label{estimating-the-within-group-variance}

The withing group variance, \(\sigma^2\) can be estimated by the mean
square error, \(MSE\), as it is an unbiased estimator. Further, a
\((1-\alpha)100\%\) Confidence Interval can be obtained as
\(SSE/\sigma^2\) is distributed as Chi-square with \(r(n-1)\) degrees of
freedom.

\[
E\{MSE\}=\sigma^2 \qquad  \frac{SSE}{\sigma^2} =
\frac{r(n-1)MSE}{\sigma^2}\sim \chi_{r(n-1)}^2
\]

\[
\quad \Rightarrow \quad P\left(\chi_{\alpha/2;r(n-1)}^2 \leq \frac{SSE}{\sigma^2}
\leq \chi_{1-\alpha/2;r(n-1)}^2\right)=1-\alpha
\]

\[
\quad \Rightarrow \quad P\left(\frac{\chi_{\alpha/2;r(n-1)}^2}{SSE} \leq \frac{1}{\sigma^2}
\leq \frac{\chi_{1-\alpha/2;r(n-1)}^2}{SSE}\right)=1-\alpha
\]

\[
\quad \Rightarrow \quad P\left(\frac{SSE}{\chi_{1-\alpha/2;r(n-1)}^2} \leq \sigma^2
\leq \frac{SSE}{\chi_{\alpha/2;r(n-1)}^2}\right)=
P\left(\frac{r(n-1)MSE}{\chi_{1-\alpha/2;r(n-1)}^2} \leq \sigma^2
\leq \frac{r(n-1)MSE}{\chi_{\alpha/2;r(n-1)}^2}\right)
=1-\alpha
\]

The Confidence Interval is simply from \(SSE\) divided by the upper
critical value to \(SSE\) divided by the lower critical value.

\textbf{Example 10.4 - Alpha Acids in Varieties of Beer}

The following R code obtains the 95\% Confidence Interval for
\(\sigma^2\) based on direct computations.

\begin{Shaded}
\begin{Highlighting}[]
\NormalTok{hops }\OtherTok{\textless{}{-}} \FunctionTok{read.csv}\NormalTok{(}\StringTok{"data/hop\_alphaacid.csv"}\NormalTok{)}
\NormalTok{ybar }\OtherTok{\textless{}{-}} \FunctionTok{mean}\NormalTok{(hops}\SpecialCharTok{$}\NormalTok{alphaAcid)}
\NormalTok{ybar.trt }\OtherTok{\textless{}{-}} \FunctionTok{as.vector}\NormalTok{(}\FunctionTok{tapply}\NormalTok{(hops}\SpecialCharTok{$}\NormalTok{alphaAcid,hops}\SpecialCharTok{$}\NormalTok{variety,mean))}
\NormalTok{var.trt }\OtherTok{\textless{}{-}} \FunctionTok{as.vector}\NormalTok{(}\FunctionTok{tapply}\NormalTok{(hops}\SpecialCharTok{$}\NormalTok{alphaAcid,hops}\SpecialCharTok{$}\NormalTok{variety,var))}
\NormalTok{n.trt }\OtherTok{\textless{}{-}} \FunctionTok{as.vector}\NormalTok{(}\FunctionTok{tapply}\NormalTok{(hops}\SpecialCharTok{$}\NormalTok{alphaAcid,hops}\SpecialCharTok{$}\NormalTok{variety,length))}
\NormalTok{r }\OtherTok{\textless{}{-}} \FunctionTok{length}\NormalTok{(ybar.trt)}
\NormalTok{n }\OtherTok{\textless{}{-}}\NormalTok{ n.trt[}\DecValTok{1}\NormalTok{]}

\NormalTok{sse }\OtherTok{\textless{}{-}} \FunctionTok{sum}\NormalTok{((n.trt}\DecValTok{{-}1}\NormalTok{)}\SpecialCharTok{*}\NormalTok{var.trt)}
\NormalTok{df.err }\OtherTok{\textless{}{-}} \FunctionTok{sum}\NormalTok{(n.trt)}\SpecialCharTok{{-}}\NormalTok{r}

\NormalTok{X2}\FloatTok{.025} \OtherTok{\textless{}{-}} \FunctionTok{qchisq}\NormalTok{(.}\DecValTok{025}\NormalTok{,df.err)}
\NormalTok{X2}\FloatTok{.975} \OtherTok{\textless{}{-}} \FunctionTok{qchisq}\NormalTok{(.}\DecValTok{975}\NormalTok{,df.err)}

\NormalTok{sigma2.ci.out }\OtherTok{\textless{}{-}} \FunctionTok{cbind}\NormalTok{(sse}\SpecialCharTok{/}\NormalTok{df.err, X2}\FloatTok{.025}\NormalTok{, X2}\FloatTok{.975}\NormalTok{,  sse}\SpecialCharTok{/}\NormalTok{X2}\FloatTok{.975}\NormalTok{, sse}\SpecialCharTok{/}\NormalTok{X2}\FloatTok{.025}\NormalTok{)}
\FunctionTok{colnames}\NormalTok{(sigma2.ci.out) }\OtherTok{\textless{}{-}} \FunctionTok{c}\NormalTok{(}\StringTok{"MSE"}\NormalTok{,}\StringTok{"X2(.025)"}\NormalTok{,}\StringTok{"X2(.975)"}\NormalTok{,}\StringTok{"LB"}\NormalTok{, }\StringTok{"UB"}\NormalTok{)}
\FunctionTok{rownames}\NormalTok{(sigma2.ci.out) }\OtherTok{\textless{}{-}} \FunctionTok{c}\NormalTok{(}\StringTok{"sigma\^{}2"}\NormalTok{)}
\FunctionTok{round}\NormalTok{(sigma2.ci.out,}\DecValTok{4}\NormalTok{)}
\end{Highlighting}
\end{Shaded}

\begin{verbatim}
##              MSE X2(.025) X2(.975)       LB      UB
## sigma^2 690.0211  65.6466 118.1359 525.6819 946.003
\end{verbatim}

The point estimate is 690, while the 95\% Confidence Interval ranges
from 525 to 946. In terms of the within group standard deviation, the
95\% Confidence Interval ranges from 23 to 31.

\[\nabla \]

\subsection{Estimating the Between Group
Variance}\label{estimating-the-between-group-variance}

Recall the expected mean squares derived previously.

\[
E\{MSTR\}=\sigma^2 + n\sigma_{\mu}^2 \qquad \qquad E\{MSE\}=\sigma^2
\]

\[
\quad \Rightarrow \quad E\left\{\frac{MSTR-MSE}{n}\right\} = 
\left(\frac{1}{n}\right)E\{MSTR\} +\left(-\frac{1}{n}\right)E\{MSE\}
=\sigma_{\mu}^2
\] This represents a linear combination of mean squares, each with a
known degrees of freedom. \textbf{Satterthwaite's Approximation} is
based on obtaining an approximation to a chi-square distribution based
on linear combinations of mean squares. The approximation is for the
degrees of freedom for the (approximate) chi-square distribution. Note
that this estimator can be negative, and is usually treated as if it is
0 in practice.

Let \(MS_1,\ldots,MS_h\) be \(h\) mean squares, with degrees of freedom
\(df_1,\ldots,df_h\), respectively. Further, let \(c_1,\ldots,c_h\), be
fixed constants. Consider the following quantities.

\[
L=\sum_{i=1}^hc_iE\{MS_i\}   \qquad\qquad
\hat{L} =\sum_{i=1}^hc_i MS_i \qquad\qquad
\frac{\left(df\right)\hat{L}}{L}  \stackrel{\bullet}{\sim} \chi_{df}^2
\] This method matches the first two moments of the approximate
chi-square distribution to solve for the approximate degrees of freedom.
For the chi-square distribution with \(\nu\) degrees of freedom, the
mean is \(\nu\) and the variance is \(2\nu\). Also, the mean squares are
independent.

\[
E\left\{\hat{L}\right\}=L \quad \Rightarrow \quad E\left\{\frac{\left(df\right)\hat{L}}{L}\right\}=df
\]

\[
\frac{\left(df_i\right)MS_i}{E\left\{MS_i\right\}} \sim \chi_{df_i}^2
\quad \Rightarrow \quad 
E\left\{\frac{\left(df_i\right)MS_i}{E\left\{MS_i\right\}}\right\}=df_i
\qquad
V\left\{\frac{\left(df_i\right)MS_i}{E\left\{MS_i\right\}}\right\}=2df_i
\] \[
\quad \Rightarrow \quad
V\left\{MS_i\right\}=\frac{2\left(df_i\right)\left(E\left\{MS_i\right\}\right)^2}{
\left(df_i\right)^2}
\quad \Rightarrow \quad
V\left\{\hat{L}\right\}=2\sum_{i=1}^hc_i^2\frac{\left(E\left\{MS_i\right\}\right)^2}{df_i}
\] \[
V\left\{\frac{\left(df\right)\hat{L}}{L}\right\}=2df
\quad \Rightarrow \quad
V\left\{\hat{L}\right\} = \frac{2(df)L^2}{(df)^2}=\frac{2L^2}{df}
\]

Now, equating the two versions of \(V\left\{\hat{L}\right\}\) and
solving for \(df\) gives the following result. The unknown expected mean
squares in \(L\) are replaced with the observed mean squares.

\[
\frac{2L^2}{df}=
2\sum_{i=1}^hc_i^2\frac{\left(E\left\{MS_i\right\}\right)^2}{df_i}
\quad \Rightarrow \quad
df = \frac{\hat{L}^2}{\sum_{i=1}^hc_i^2\frac{\left(MS_i\right)^2}{df_i}}=
\frac{\left(\sum_{i=1}^h c_iMS_i\right)^2}{\sum_{i=1}^hc_i^2\frac{\left(MS_i\right)^2}{df_i}}
\] The approximate \((1-\alpha)100\%\) Confidence Interval is computed
as follows.

\[
(1-\alpha)100\% \mbox{ Confidence Interval for } \sigma_{\mu}^2: \quad
\left[\frac{(df)\hat{L}}{\chi_{1-\alpha/2;df}},\frac{(df)\hat{L}}{\chi_{\alpha/2;df}}\right]
\]

\textbf{Example 10.5 - Alpha Acids in Varieties of Beer}

The following R code obtains the approximate 95\% Confidence Interval
for \(\sigma_{\mu}^2\) based on direct computations. Note that
\(c_1=1/n\) and \(c_2=-1/n\).

\begin{Shaded}
\begin{Highlighting}[]
\NormalTok{hops }\OtherTok{\textless{}{-}} \FunctionTok{read.csv}\NormalTok{(}\StringTok{"data/hop\_alphaacid.csv"}\NormalTok{)}
\NormalTok{ybar }\OtherTok{\textless{}{-}} \FunctionTok{mean}\NormalTok{(hops}\SpecialCharTok{$}\NormalTok{alphaAcid)}
\NormalTok{ybar.trt }\OtherTok{\textless{}{-}} \FunctionTok{as.vector}\NormalTok{(}\FunctionTok{tapply}\NormalTok{(hops}\SpecialCharTok{$}\NormalTok{alphaAcid,hops}\SpecialCharTok{$}\NormalTok{variety,mean))}
\NormalTok{var.trt }\OtherTok{\textless{}{-}} \FunctionTok{as.vector}\NormalTok{(}\FunctionTok{tapply}\NormalTok{(hops}\SpecialCharTok{$}\NormalTok{alphaAcid,hops}\SpecialCharTok{$}\NormalTok{variety,var))}
\NormalTok{n.trt }\OtherTok{\textless{}{-}} \FunctionTok{as.vector}\NormalTok{(}\FunctionTok{tapply}\NormalTok{(hops}\SpecialCharTok{$}\NormalTok{alphaAcid,hops}\SpecialCharTok{$}\NormalTok{variety,length))}
\NormalTok{r }\OtherTok{\textless{}{-}} \FunctionTok{length}\NormalTok{(ybar.trt)}
\NormalTok{n }\OtherTok{\textless{}{-}}\NormalTok{ n.trt[}\DecValTok{1}\NormalTok{]}

\NormalTok{sstr }\OtherTok{\textless{}{-}} \FunctionTok{sum}\NormalTok{(n.trt}\SpecialCharTok{*}\NormalTok{(ybar.trt}\SpecialCharTok{{-}}\NormalTok{ybar)}\SpecialCharTok{\^{}}\DecValTok{2}\NormalTok{)}
\NormalTok{sse }\OtherTok{\textless{}{-}} \FunctionTok{sum}\NormalTok{((n.trt}\DecValTok{{-}1}\NormalTok{)}\SpecialCharTok{*}\NormalTok{var.trt)}
\NormalTok{df.trt }\OtherTok{\textless{}{-}}\NormalTok{ r}\DecValTok{{-}1}
\NormalTok{df.err }\OtherTok{\textless{}{-}} \FunctionTok{sum}\NormalTok{(n.trt)}\SpecialCharTok{{-}}\NormalTok{r}
\NormalTok{mse }\OtherTok{\textless{}{-}}\NormalTok{ sse}\SpecialCharTok{/}\NormalTok{df.err}
\NormalTok{mstr }\OtherTok{\textless{}{-}}\NormalTok{ sstr}\SpecialCharTok{/}\NormalTok{df.trt}

\NormalTok{L.hat }\OtherTok{\textless{}{-}}\NormalTok{ (mstr}\SpecialCharTok{{-}}\NormalTok{mse) }\SpecialCharTok{/}\NormalTok{ n}
\NormalTok{df.L.hat }\OtherTok{\textless{}{-}}\NormalTok{ (L.hat}\SpecialCharTok{\^{}}\DecValTok{2}\NormalTok{) }\SpecialCharTok{/}\NormalTok{ ((}\DecValTok{1}\SpecialCharTok{/}\NormalTok{n)}\SpecialCharTok{\^{}}\DecValTok{2}\SpecialCharTok{*}\NormalTok{mstr}\SpecialCharTok{\^{}}\DecValTok{2}\SpecialCharTok{/}\NormalTok{df.trt }\SpecialCharTok{+}\NormalTok{ (}\SpecialCharTok{{-}}\DecValTok{1}\SpecialCharTok{/}\NormalTok{n)}\SpecialCharTok{\^{}}\DecValTok{2}\SpecialCharTok{*}\NormalTok{mse}\SpecialCharTok{\^{}}\DecValTok{2}\SpecialCharTok{/}\NormalTok{df.err)}

\NormalTok{X2}\FloatTok{.025} \OtherTok{\textless{}{-}} \FunctionTok{qchisq}\NormalTok{(.}\DecValTok{025}\NormalTok{,df.L.hat)}
\NormalTok{X2}\FloatTok{.975} \OtherTok{\textless{}{-}} \FunctionTok{qchisq}\NormalTok{(.}\DecValTok{975}\NormalTok{,df.L.hat)}
\NormalTok{sigma2.mu.ci.out }\OtherTok{\textless{}{-}} \FunctionTok{cbind}\NormalTok{(L.hat, df.L.hat, X2}\FloatTok{.025}\NormalTok{, X2}\FloatTok{.975}\NormalTok{, }
\NormalTok{                           df.L.hat}\SpecialCharTok{*}\NormalTok{L.hat}\SpecialCharTok{/}\NormalTok{X2}\FloatTok{.975}\NormalTok{, df.L.hat}\SpecialCharTok{*}\NormalTok{L.hat}\SpecialCharTok{/}\NormalTok{X2}\FloatTok{.025}\NormalTok{)}
\FunctionTok{colnames}\NormalTok{(sigma2.mu.ci.out) }\OtherTok{\textless{}{-}} \FunctionTok{c}\NormalTok{(}\StringTok{"s2\_mu"}\NormalTok{,}\StringTok{"df"}\NormalTok{, }\StringTok{"X2(.025)"}\NormalTok{,}\StringTok{"X2(.975)"}\NormalTok{,}\StringTok{"LB"}\NormalTok{, }\StringTok{"UB"}\NormalTok{)}
\FunctionTok{rownames}\NormalTok{(sigma2.mu.ci.out) }\OtherTok{\textless{}{-}} \FunctionTok{c}\NormalTok{(}\StringTok{"sigma2\_mu"}\NormalTok{)}
\FunctionTok{round}\NormalTok{(sigma2.mu.ci.out,}\DecValTok{4}\NormalTok{)}
\end{Highlighting}
\end{Shaded}

\begin{verbatim}
##              s2_mu     df X2(.025) X2(.975)       LB       UB
## sigma2_mu 16196.73 8.9238   2.6597  18.9104 7643.216 54342.17
\end{verbatim}

The point estimate is 16197, and the approximate 95\% Confidence
Interval goes from 7643 to 54342. In terms of the between group standard
deviation, the range is approximately from 87 to 233.

We will use the \textbf{lmerTest} package and the \textbf{lmer} function
to fit the model. With this function, we must include a ``1'' for the
fixed mean, and we include ``(1\textbar variety.f)'' to include random
effects for varieties.

\begin{Shaded}
\begin{Highlighting}[]
\NormalTok{hops }\OtherTok{\textless{}{-}} \FunctionTok{read.csv}\NormalTok{(}\StringTok{"data/hop\_alphaacid.csv"}\NormalTok{)}
\NormalTok{hops}\SpecialCharTok{$}\NormalTok{variety.f }\OtherTok{\textless{}{-}} \FunctionTok{factor}\NormalTok{(hops}\SpecialCharTok{$}\NormalTok{variety)}
\FunctionTok{options}\NormalTok{(}\AttributeTok{contrasts=}\FunctionTok{c}\NormalTok{(}\StringTok{"contr.sum"}\NormalTok{, }\StringTok{"contr.poly"}\NormalTok{))}

\FunctionTok{library}\NormalTok{(lmerTest)}
\NormalTok{hops.mod1 }\OtherTok{\textless{}{-}} \FunctionTok{lmer}\NormalTok{(alphaAcid }\SpecialCharTok{\textasciitilde{}} \DecValTok{1} \SpecialCharTok{+}\NormalTok{ (}\DecValTok{1}\SpecialCharTok{|}\NormalTok{variety.f), }\AttributeTok{data=}\NormalTok{hops)}
\FunctionTok{summary}\NormalTok{(hops.mod1)}
\end{Highlighting}
\end{Shaded}

\begin{verbatim}
## Linear mixed model fit by REML. t-tests use Satterthwaite's method [
## lmerModLmerTest]
## Formula: alphaAcid ~ 1 + (1 | variety.f)
##    Data: hops
## 
## REML criterion at convergence: 981.9
## 
## Scaled residuals: 
##      Min       1Q   Median       3Q      Max 
## -2.63193 -0.46012  0.00626  0.56199  3.06408 
## 
## Random effects:
##  Groups    Name        Variance Std.Dev.
##  variety.f (Intercept) 16197    127.27  
##  Residual                690     26.27  
## Number of obs: 100, groups:  variety.f, 10
## 
## Fixed effects:
##             Estimate Std. Error     df t value Pr(>|t|)    
## (Intercept)   579.61      40.33   9.00   14.37 1.64e-07 ***
## ---
## Signif. codes:  0 '***' 0.001 '**' 0.01 '*' 0.05 '.' 0.1 ' ' 1
\end{verbatim}

\[ \nabla \]

\section{Two-Way Random Effects
Model}\label{two-way-random-effects-model}

Some studies include two (or more) random factors. In these cases we
want to generalize across populations of units for the two factors and
their interactions. In \textbf{Measurement Systems Analysis}, variation
in parts and operators who measure the parts are studied. The
experiments are often referred to as \textbf{Gage Repeatability \&
Reproducibility (R\&R)} studies. Each of \(a\) parts are measured by
\(b\) operators \(n\) times in random order. The variation in part
sizes, operator measurements, and interactions between parts and
operators are of interest. Repeatability refers to the variation when
the same part is measured by the same operator multiple times.

The statistical model for the balanced case is given below.

\[
Y_{ijk} = \mu_{\bullet\bullet} + \alpha_i + \beta_j + (\alpha\beta)_{ij} + \epsilon_{ijk} \qquad i=1,\ldots,a; \quad j=1,\ldots,b; \quad k=1,\ldots,n
\]

It is assumed that \(\left\{\alpha_i\right\}\),
\(\left\{\beta_j\right\}\), \(\left\{(\alpha\beta)_{ij}\right\}\), and
\(\left\{\epsilon_{ijk}\right\}\) are independent and normally
distributed with means 0 and variances \(\sigma_{\alpha}^2\),
\(\sigma_{\beta}^2\), \(\sigma_{\alpha\beta}^2\), and \(\sigma^2\),
respectively. The mean, variance and covariance structure are given
below.

\[
E\left\{Y_{ijk}\right\} = E\left\{\mu_{\bullet\bullet} + \alpha_i + \beta_j + (\alpha\beta)_{ij} + \epsilon_{ijk}\right\}
= \mu_{\bullet\bullet} + 0 + 0 + 0 + 0 = \mu_{\bullet\bullet}
\] \[
V\left\{Y_{ijk}\right\} = V\left\{\mu_{\bullet\bullet} + \alpha_i + \beta_j + (\alpha\beta)_{ij} + \epsilon_{ijk}\right\} =
\sigma_{\alpha}^2 + \sigma_{\beta}^2 + \sigma_{\alpha\beta}^2 +  \sigma^2
\] \[
\mbox{COV}\left\{Y_{ijk},Y_{i'j'k'}\right\}=
\left\{ \begin{array}{r@{\quad:\quad}l}
    \sigma_{\alpha}^2 + \sigma_{\beta}^2 + \sigma_{\alpha\beta}^2 +  \sigma^2 & i=i',j=j',k=k' \\
    \sigma_{\alpha}^2 + \sigma_{\beta}^2 + \sigma_{\alpha\beta}^2 & i=i', j=j',k\ne k' \\
    \sigma_{\alpha}^2 & i=i',j\ne j', \forall k,k'\\
     \sigma_{\beta}^2&  i\ne i', j= j',\forall k,k' \\
     0 & i\ne i', j\ne j', \forall k,k' \\
    \end{array} \right.
\] The expected mean squares for each source of variation can be derived
from this covariance structure. The results are given here.

\[
\mbox{Factor A:} \quad E\{MSA\} = \sigma^2 + n\sigma_{\alpha\beta}^2+ bn\sigma_{\alpha}^2 
\qquad 
\mbox{Factor B:} \quad E\{MSB\} = \sigma^2 + n\sigma_{\alpha\beta}^2+ an\sigma_{\beta}^2 
\] \[
\mbox{AB Interaction:} \quad E\{MSAB\}=\sigma^2 + n\sigma_{\alpha\beta}^2
\qquad
\mbox{Error:} \quad E\{MSE\}=\sigma^2
\]

The tests for interaction and main effects are obtained from the
expected mean squares. Note that the ``error'' terms for factors A and B
are the AB interaction, unlike the fixed effects model, which uses
\(MSE\) for the error term.

\[
H_0^{AB}:\sigma_{\alpha\beta}^2 = 0 \qquad H_A^{AB}:\sigma_{\alpha\beta}^2 > 0 \qquad
TS:F_{AB}=\frac{MSAB}{MSE} \qquad RR:F_{AB}\geq F_{1-\alpha;(a-1)(b-1),ab(n-1)}
\] \[
H_0^{A}:\sigma_{\alpha}^2 = 0 \qquad H_A^{A}:\sigma_{\alpha}^2 > 0 \qquad
TS:F_{A}=\frac{MSA}{MSAB} \qquad RR:F_{A}\geq F_{1-\alpha;a-1,(a-1)(b-1)}
\] \[
H_0^{B}:\sigma_{\beta}^2 = 0 \qquad H_A^{B}:\sigma_{\beta}^2 > 0 \qquad
TS:F_{B}=\frac{MSB}{MSAB} \qquad RR:F_{B}\geq F_{1-\alpha;b-1,(a-1)(b-1)}
\]

Point estimates for the variance components are obtained from the
expected mean squares. Some of these can be negative, and in practice
are typically replaced by 0.

\[
s^2=MSE \qquad s_{\alpha\beta}^2=\frac{MSAB-MSE}{n}  \qquad
s_{\alpha}^2=\frac{MSA-MSAB}{bn}  \qquad s_{\beta}^2=\frac{MSB-MSAB}{an}
\]

As these are all linear functions of mean squares, Satterthwaite's
approximation can be used to obtain approximate degrees' of freedom and
Confidence Intervals for \(\sigma_{\alpha\beta}^2\),
\(\sigma_{\alpha}^2\), and \(\sigma_{\beta}^2\).

\[
df_{AB}=\frac{\left(\frac{1}{n}MSAB + \frac{-1}{n}MSE\right)^2} {
\frac{\left(\frac{1}{n}MSAB\right)^2}{(a-1)(b-1)} + 
\frac{\left(\frac{-1}{n}MSE\right)^2}{ab(n-1)}
}
\]

\[
df_{A}=\frac{\left(\frac{1}{bn}MSA + \frac{-1}{bn}MSAB\right)^2} {
\frac{\left(\frac{1}{bn}MSA\right)^2}{a-1} + 
\frac{\left(\frac{-1}{bn}MSAB\right)^2}{(a-1)(b-1)}
}
\qquad
df_{B}=\frac{\left(\frac{1}{an}MSB + \frac{-1}{an}MSAB\right)^2} {
\frac{\left(\frac{1}{an}MSB\right)^2}{b-1} + 
\frac{\left(\frac{-1}{an}MSAB\right)^2}{(a-1)(b-1)}
}
\]

\textbf{Example 10.6 - Measurements of Skulls by Observers}

An anthropological study was conducted to measure variation in frontal
chord measurements of skulls by various observers
(\citeproc{ref-hanihara_1999}{Hanihara et al. 1999}). There were
\(a=10\) skulls, \(b=6\) observers, and \(n=2\) replicates per
skull/observer. The R code for the analysis and output are included
here.

\begin{Shaded}
\begin{Highlighting}[]
\NormalTok{sm }\OtherTok{\textless{}{-}} \FunctionTok{read.csv}\NormalTok{(}\StringTok{"data/skull\_measure.csv"}\NormalTok{)}
\FunctionTok{head}\NormalTok{(sm, }\DecValTok{3}\NormalTok{)}
\end{Highlighting}
\end{Shaded}

\begin{verbatim}
##   observer skull time chord
## 1        1     1    1 100.9
## 2        1     1    2 100.8
## 3        1     2    1  91.1
\end{verbatim}

\begin{Shaded}
\begin{Highlighting}[]
\FunctionTok{tail}\NormalTok{(sm, }\DecValTok{3}\NormalTok{)}
\end{Highlighting}
\end{Shaded}

\begin{verbatim}
##     observer skull time chord
## 118        6     9    2  91.8
## 119        6    10    1  93.7
## 120        6    10    2  93.8
\end{verbatim}

\begin{Shaded}
\begin{Highlighting}[]
\NormalTok{sm}\SpecialCharTok{$}\NormalTok{skull.f }\OtherTok{\textless{}{-}} \FunctionTok{factor}\NormalTok{(sm}\SpecialCharTok{$}\NormalTok{skull)}
\NormalTok{sm}\SpecialCharTok{$}\NormalTok{observer.f }\OtherTok{\textless{}{-}} \FunctionTok{factor}\NormalTok{(sm}\SpecialCharTok{$}\NormalTok{observer)}

\NormalTok{mod1 }\OtherTok{\textless{}{-}} \FunctionTok{aov}\NormalTok{(chord}\SpecialCharTok{\textasciitilde{}}\NormalTok{ skull.f}\SpecialCharTok{*}\NormalTok{observer.f, }\AttributeTok{data=}\NormalTok{sm)}
\FunctionTok{anova}\NormalTok{(mod1)}
\end{Highlighting}
\end{Shaded}

\begin{verbatim}
## Analysis of Variance Table
## 
## Response: chord
##                    Df  Sum Sq Mean Sq   F value    Pr(>F)    
## skull.f             9 1612.01 179.112 3074.8786 < 2.2e-16 ***
## observer.f          5    7.32   1.463   25.1207 1.449e-13 ***
## skull.f:observer.f 45    5.26   0.117    2.0065  0.005956 ** 
## Residuals          60    3.49   0.058                        
## ---
## Signif. codes:  0 '***' 0.001 '**' 0.01 '*' 0.05 '.' 0.1 ' ' 1
\end{verbatim}

\begin{Shaded}
\begin{Highlighting}[]
\NormalTok{a }\OtherTok{\textless{}{-}} \FunctionTok{max}\NormalTok{(sm}\SpecialCharTok{$}\NormalTok{skull); b }\OtherTok{\textless{}{-}} \FunctionTok{max}\NormalTok{(sm}\SpecialCharTok{$}\NormalTok{observer); n }\OtherTok{\textless{}{-}} \FunctionTok{max}\NormalTok{(sm}\SpecialCharTok{$}\NormalTok{time)}
\NormalTok{msa }\OtherTok{\textless{}{-}} \FunctionTok{anova}\NormalTok{(mod1)[}\DecValTok{1}\NormalTok{,}\DecValTok{3}\NormalTok{];  df\_a }\OtherTok{\textless{}{-}} \FunctionTok{anova}\NormalTok{(mod1)[}\DecValTok{1}\NormalTok{,}\DecValTok{1}\NormalTok{]}
\NormalTok{msb }\OtherTok{\textless{}{-}} \FunctionTok{anova}\NormalTok{(mod1)[}\DecValTok{2}\NormalTok{,}\DecValTok{3}\NormalTok{];  df\_b }\OtherTok{\textless{}{-}} \FunctionTok{anova}\NormalTok{(mod1)[}\DecValTok{2}\NormalTok{,}\DecValTok{1}\NormalTok{]}
\NormalTok{msab }\OtherTok{\textless{}{-}} \FunctionTok{anova}\NormalTok{(mod1)[}\DecValTok{3}\NormalTok{,}\DecValTok{3}\NormalTok{];  df\_ab }\OtherTok{\textless{}{-}} \FunctionTok{anova}\NormalTok{(mod1)[}\DecValTok{3}\NormalTok{,}\DecValTok{1}\NormalTok{]}
\NormalTok{mse }\OtherTok{\textless{}{-}} \FunctionTok{anova}\NormalTok{(mod1)[}\DecValTok{4}\NormalTok{,}\DecValTok{3}\NormalTok{];  df\_e }\OtherTok{\textless{}{-}} \FunctionTok{anova}\NormalTok{(mod1)[}\DecValTok{4}\NormalTok{,}\DecValTok{1}\NormalTok{]}

\NormalTok{s2\_ab }\OtherTok{\textless{}{-}}\NormalTok{ (msab}\SpecialCharTok{{-}}\NormalTok{mse)}\SpecialCharTok{/}\NormalTok{n}
\NormalTok{s2\_a }\OtherTok{\textless{}{-}}\NormalTok{ (msa}\SpecialCharTok{{-}}\NormalTok{msab)}\SpecialCharTok{/}\NormalTok{(b}\SpecialCharTok{*}\NormalTok{n)}
\NormalTok{s2\_b }\OtherTok{\textless{}{-}}\NormalTok{ (msb}\SpecialCharTok{{-}}\NormalTok{msab)}\SpecialCharTok{/}\NormalTok{(a}\SpecialCharTok{*}\NormalTok{n)}
\NormalTok{df\_s2\_ab }\OtherTok{\textless{}{-}}\NormalTok{ s2\_ab}\SpecialCharTok{\^{}}\DecValTok{2} \SpecialCharTok{/}\NormalTok{ ((msab}\SpecialCharTok{/}\NormalTok{n)}\SpecialCharTok{\^{}}\DecValTok{2}\SpecialCharTok{/}\NormalTok{df\_ab }\SpecialCharTok{+}\NormalTok{ (}\SpecialCharTok{{-}}\NormalTok{mse}\SpecialCharTok{/}\NormalTok{n)}\SpecialCharTok{\^{}}\DecValTok{2}\SpecialCharTok{/}\NormalTok{df\_e)}
\NormalTok{df\_s2\_a }\OtherTok{\textless{}{-}}\NormalTok{ s2\_a}\SpecialCharTok{\^{}}\DecValTok{2} \SpecialCharTok{/}\NormalTok{ ((msa}\SpecialCharTok{/}\NormalTok{(b}\SpecialCharTok{*}\NormalTok{n))}\SpecialCharTok{\^{}}\DecValTok{2}\SpecialCharTok{/}\NormalTok{df\_a }\SpecialCharTok{+}\NormalTok{ (}\SpecialCharTok{{-}}\NormalTok{msab}\SpecialCharTok{/}\NormalTok{(b}\SpecialCharTok{*}\NormalTok{n))}\SpecialCharTok{\^{}}\DecValTok{2}\SpecialCharTok{/}\NormalTok{df\_ab)}
\NormalTok{df\_s2\_b }\OtherTok{\textless{}{-}}\NormalTok{ s2\_b}\SpecialCharTok{\^{}}\DecValTok{2} \SpecialCharTok{/}\NormalTok{ ((msb}\SpecialCharTok{/}\NormalTok{(a}\SpecialCharTok{*}\NormalTok{n))}\SpecialCharTok{\^{}}\DecValTok{2}\SpecialCharTok{/}\NormalTok{df\_b }\SpecialCharTok{+}\NormalTok{ (}\SpecialCharTok{{-}}\NormalTok{msab}\SpecialCharTok{/}\NormalTok{(a}\SpecialCharTok{*}\NormalTok{n))}\SpecialCharTok{\^{}}\DecValTok{2}\SpecialCharTok{/}\NormalTok{df\_ab)}

\NormalTok{s2.out1 }\OtherTok{\textless{}{-}} \FunctionTok{cbind}\NormalTok{(msa, msa}\SpecialCharTok{/}\NormalTok{msab, }\FunctionTok{qf}\NormalTok{(.}\DecValTok{95}\NormalTok{,df\_a,df\_ab), }\DecValTok{1}\SpecialCharTok{{-}}\FunctionTok{pf}\NormalTok{(msa}\SpecialCharTok{/}\NormalTok{msab,df\_a,df\_ab),}
\NormalTok{                 s2\_a, df\_s2\_a, df\_s2\_a}\SpecialCharTok{*}\NormalTok{s2\_a}\SpecialCharTok{/}\FunctionTok{qchisq}\NormalTok{(.}\DecValTok{975}\NormalTok{,df\_s2\_a),df\_s2\_a}\SpecialCharTok{*}\NormalTok{s2\_a}\SpecialCharTok{/}\FunctionTok{qchisq}\NormalTok{(.}\DecValTok{025}\NormalTok{,df\_s2\_a))}
\NormalTok{s2.out2 }\OtherTok{\textless{}{-}} \FunctionTok{cbind}\NormalTok{(msb, msb}\SpecialCharTok{/}\NormalTok{msab, }\FunctionTok{qf}\NormalTok{(.}\DecValTok{95}\NormalTok{,df\_b,df\_ab), }\DecValTok{1}\SpecialCharTok{{-}}\FunctionTok{pf}\NormalTok{(msb}\SpecialCharTok{/}\NormalTok{msab,df\_b,df\_ab),}
\NormalTok{                 s2\_b, df\_s2\_b, df\_s2\_b}\SpecialCharTok{*}\NormalTok{s2\_b}\SpecialCharTok{/}\FunctionTok{qchisq}\NormalTok{(.}\DecValTok{975}\NormalTok{,df\_s2\_b),df\_s2\_b}\SpecialCharTok{*}\NormalTok{s2\_b}\SpecialCharTok{/}\FunctionTok{qchisq}\NormalTok{(.}\DecValTok{025}\NormalTok{,df\_s2\_b))}
\NormalTok{s2.out3 }\OtherTok{\textless{}{-}} \FunctionTok{cbind}\NormalTok{(msab, msab}\SpecialCharTok{/}\NormalTok{mse, }\FunctionTok{qf}\NormalTok{(.}\DecValTok{95}\NormalTok{,df\_ab,df\_e), }\DecValTok{1}\SpecialCharTok{{-}}\FunctionTok{pf}\NormalTok{(msab}\SpecialCharTok{/}\NormalTok{mse,df\_ab,df\_e),}
\NormalTok{                 s2\_ab, df\_s2\_ab, df\_s2\_ab}\SpecialCharTok{*}\NormalTok{s2\_ab}\SpecialCharTok{/}\FunctionTok{qchisq}\NormalTok{(.}\DecValTok{975}\NormalTok{,df\_s2\_ab),df\_s2\_ab}\SpecialCharTok{*}\NormalTok{s2\_ab}\SpecialCharTok{/}\FunctionTok{qchisq}\NormalTok{(.}\DecValTok{025}\NormalTok{,df\_s2\_ab))}
\NormalTok{s2.out4 }\OtherTok{\textless{}{-}} \FunctionTok{cbind}\NormalTok{(mse, }\ConstantTok{NA}\NormalTok{, }\ConstantTok{NA}\NormalTok{, }\ConstantTok{NA}\NormalTok{, mse, df\_e, df\_e}\SpecialCharTok{*}\NormalTok{mse}\SpecialCharTok{/}\FunctionTok{qchisq}\NormalTok{(.}\DecValTok{975}\NormalTok{,df\_e),}
\NormalTok{                 df\_e}\SpecialCharTok{*}\NormalTok{mse}\SpecialCharTok{/}\FunctionTok{qchisq}\NormalTok{(.}\DecValTok{025}\NormalTok{,df\_e))}
\NormalTok{s2.out }\OtherTok{\textless{}{-}} \FunctionTok{rbind}\NormalTok{(s2.out1,s2.out2,s2.out3,s2.out4)}
\FunctionTok{colnames}\NormalTok{(s2.out) }\OtherTok{\textless{}{-}} \FunctionTok{c}\NormalTok{(}\StringTok{"MS"}\NormalTok{, }\StringTok{"F*"}\NormalTok{, }\StringTok{"F(.95)"}\NormalTok{, }\StringTok{"P(\textgreater{}F*)"}\NormalTok{, }\StringTok{"s\^{}2"}\NormalTok{, }\StringTok{"Satt df"}\NormalTok{, }\StringTok{"LB"}\NormalTok{, }\StringTok{"UB"}\NormalTok{)}
\FunctionTok{rownames}\NormalTok{(s2.out) }\OtherTok{\textless{}{-}} \FunctionTok{c}\NormalTok{(}\StringTok{"sigma\_a\^{}2"}\NormalTok{, }\StringTok{"sigma\_b\^{}2"}\NormalTok{, }\StringTok{"sigma\_ab\^{}2"}\NormalTok{, }\StringTok{"sigma\^{}2"}\NormalTok{)}
\FunctionTok{round}\NormalTok{(s2.out,}\DecValTok{4}\NormalTok{)}
\end{Highlighting}
\end{Shaded}

\begin{verbatim}
##                  MS        F* F(.95) P(>F*)     s^2 Satt df     LB      UB
## sigma_a^2  179.1117 1532.4942 2.0958  0.000 14.9162  8.9883 7.0543 49.7643
## sigma_b^2    1.4633   12.5200 2.4221  0.000  0.0673  4.2302 0.0247  0.5102
## sigma_ab^2   0.1169    2.0065 1.5749  0.006  0.0293  9.5444 0.0141  0.0934
## sigma^2      0.0582        NA     NA     NA  0.0582 60.0000 0.0420  0.0863
\end{verbatim}

\begin{Shaded}
\begin{Highlighting}[]
\FunctionTok{library}\NormalTok{(lmerTest)}
\FunctionTok{options}\NormalTok{(}\AttributeTok{contrasts=}\FunctionTok{c}\NormalTok{(}\StringTok{"contr.sum"}\NormalTok{, }\StringTok{"contr.poly"}\NormalTok{))}
\NormalTok{mod2 }\OtherTok{\textless{}{-}} \FunctionTok{lmer}\NormalTok{(chord }\SpecialCharTok{\textasciitilde{}} \DecValTok{1} \SpecialCharTok{+}\NormalTok{  (}\DecValTok{1}\SpecialCharTok{|}\NormalTok{skull.f) }\SpecialCharTok{+}\NormalTok{ (}\DecValTok{1}\SpecialCharTok{|}\NormalTok{observer.f) }\SpecialCharTok{+}\NormalTok{ (}\DecValTok{1}\SpecialCharTok{|}\NormalTok{skull.f}\SpecialCharTok{:}\NormalTok{observer.f), }\AttributeTok{data=}\NormalTok{sm)}
\FunctionTok{summary}\NormalTok{(mod2)}
\end{Highlighting}
\end{Shaded}

\begin{verbatim}
## Linear mixed model fit by REML. t-tests use Satterthwaite's method [
## lmerModLmerTest]
## Formula: 
## chord ~ 1 + (1 | skull.f) + (1 | observer.f) + (1 | skull.f:observer.f)
##    Data: sm
## 
## REML criterion at convergence: 123.9
## 
## Scaled residuals: 
##      Min       1Q   Median       3Q      Max 
## -2.90775 -0.49619  0.01282  0.41430  2.21863 
## 
## Random effects:
##  Groups             Name        Variance Std.Dev.
##  skull.f:observer.f (Intercept)  0.02931 0.1712  
##  skull.f            (Intercept) 14.91283 3.8617  
##  observer.f         (Intercept)  0.06732 0.2595  
##  Residual                        0.05825 0.2414  
## Number of obs: 120, groups:  skull.f:observer.f, 60; skull.f, 10; observer.f, 6
## 
## Fixed effects:
##             Estimate Std. Error     df t value Pr(>|t|)    
## (Intercept)   96.464      1.226  9.138   78.67 2.97e-14 ***
## ---
## Signif. codes:  0 '***' 0.001 '**' 0.01 '*' 0.05 '.' 0.1 ' ' 1
\end{verbatim}

The variance of the skulls is by far the largest (not surprising) while
the variances among observers, interactions and error are much smaller.

In the measurement systems analysis, the variances are broken into
various components.

\begin{itemize}
\tightlist
\item
  \textbf{Repeatability} - Variation within same part/operator
  (skull/observer): \(\sigma^2\)
\item
  \textbf{Reproducibility} - Variation across operators and
  part/operator interactions:
  \(\sigma_{\beta}^2+\sigma_{\alpha\beta}^2\)
\item
  \textbf{Parts} - Variation across parts: \(\sigma_{\alpha}^2\)
\item
  \textbf{Gage} - Repeatability + Reproducibility
\item
  \textbf{Total} - Parts + Gage
\item
  \textbf{\% R\&R} - \(100*\sigma_{\mbox{Gage}}/\sigma_{\mbox{Total}}\)
  in terms of standard deviations, not variances.
\end{itemize}

For this example (based on the lmer results), we get the following
estimates.

\[
s_{\mbox{Repeatability}}^2= s^2=0.05825 \qquad 
s_{\mbox{Reproducibility}}^2= s_{\beta}^2+s_{\alpha\beta}^2=0.06732+0.02931=0.09663
\] \[
s_{\mbox{Parts}}^2=14.92183  \qquad
s_{\mbox{Gage}}^2=0.05825+0.09963=0.015488
\] \[
s_{\mbox{Total}}^2=14.92183+0.015488=15.07671 \qquad
 \% \mbox{ R&R:} \quad
100\sqrt{\frac{0.015488}{15.07671}}=3.21\%
\]

Only about 3.2\% of the total variation (standard deviation) is due to
measurement error.

\[ \nabla \]

\section{Two-Way Mixed Effects Model}\label{two-way-mixed-effects-model}

Models can include both fixed and random effects. In this section, we
describe the \textbf{unrestricted} two-way mixed model, with factor A
fixed and factor B random. The \textbf{restricted} model has the
two-factor interactions sum to 0 over the fixed factor. Statistical
software packages typically fit the unrestricted model that allows for
unbalanced data. The model is given below.

\[
Y_{ijk}= \mu_{\bullet\bullet} + \alpha_i + \beta_j + (\alpha\beta)_{ij} + \epsilon_{ijk}
\qquad i=1,\ldots,a;\quad j=1,\ldots,b; \quad k=1,\ldots,n 
\]

\[
\sum_{i=1}^a \alpha_i = 0 \qquad \beta_j\sim NID\left(0,\sigma_{\beta}^2\right) \qquad
(\alpha\beta)_{ij} \sim NID\left(0,\sigma_{\alpha\beta}^2\right) \qquad
\epsilon_{ijk} \sim NID\left(0,\sigma^2\right)
\]

The random effects \(\left\{\beta_j\right\}\),
\(\left\{(\alpha\beta)_{ij}\right\}\), and error terms
\(\left\{\epsilon_{ijk}\right\}\) are assumed pairwise independent.

The mean, covariance structure, expected mean squares and estimators for
the unrestricted mixed model are given below.

\[
E\left\{Y_{ijk}\right\}=\mu_{\bullet\bullet}+ \alpha_i \qquad
\mbox{COV}\left\{Y_{ijk},Y_{i'j'k'}\right\} =
\left\{ \begin{array}{r@{\quad:\quad}l}
   \sigma_{\beta}^2 + \sigma_{\alpha\beta}^2 +  \sigma^2 & i=i',j=j',k=k' \\
   \sigma_{\beta}^2 + \sigma_{\alpha\beta}^2 & i=i', j=j',k\ne k' \\
     \sigma_{\beta}^2&  i\ne i', j= j',\forall k,k' \\
     0 & \forall i,i', j\ne j', \forall k,k' \\
    \end{array} \right.
\]

\[
E\{MSA\} = \sigma^2 +n\sigma_{\alpha\beta}^2+\frac{bn\sum_{i=1}^a\alpha_i^2}{a-1} \qquad
E\{MSB\} = \sigma^2 +n\sigma_{\alpha\beta}^2+ an\sigma_{\beta}^2
\] \[
E\{MSAB\}= \sigma^2 +n\sigma_{\alpha\beta}^2  \qquad E\{MSE\}=\sigma^2
\]

\[
s^2=MSE \qquad s_{\alpha\beta}^2=\frac{MSAB-MSE}{n} 
\qquad s_{\beta}^2=\frac{MSB-MSAB}{an}
\] The variances and covariances of the means are developed below, which
can be used to make inferences regarding the fixed effects in the model.

\[
V\left\{\overline{Y}_{ij\bullet}\right\} = 
V\left\{\frac{1}{n}\sum_{k=1}^nY_{ijk}\right\}=
\frac{1}{n^2}\left[\sum_{k=1}^n V\left\{Y_{ijk}\right\} +
2\sum_{k=1}^{n-1}\sum_{k'=k+1}^n\mbox{COV}\left\{Y_{ijk},Y_{ijk'}\right\}\right]
=
\] \[
= \frac{1}{n^2}\left[n\left(\sigma_{\beta}^2 + \sigma_{\alpha\beta}^2 +  \sigma^2\right)
+n(n-1)\left(\sigma_{\beta}^2 + \sigma_{\alpha\beta}^2\right)\right] = 
\frac{n\sigma_{\beta}^2 + n\sigma_{\alpha\beta}^2 + \sigma^2}{n}
\] \[
\mbox{COV}\left\{\overline{Y}_{ij\bullet},\overline{Y}_{i'j'\bullet}\right\}=0 \quad j\neq j' \quad \forall i,i'
\] \[
\mbox{COV}\left\{\overline{Y}_{ij\bullet},\overline{Y}_{i'j\bullet}\right\}=
\frac{1}{n^2} \mbox{COV}\left\{\sum_{k=1}^nY_{ijk},\sum_{k=1}^nY_{i'jk}\right\}=
\frac{1}{n^2}n^2\sigma_{\beta}^2
\quad i\neq i' \quad \forall j
\]

\[
V\left\{\overline{Y}_{i\bullet\bullet}\right\} = 
V\left\{\frac{1}{b} \sum_{j=1}^b\overline{Y}_{ij\bullet}\right\} =
\frac{1}{b^2}\left[\sum_{j=1}^b V\left\{\overline{Y}_{ij\bullet}\right\} +
2\sum_{j=1}^{b-1}\sum_{j'=j+1}^b
\mbox{COV}\left\{\overline{Y}_{ij\bullet},\overline{Y}_{ij'\bullet}\right\}\right] =
\]

\[
= \frac{1}{b^2}\left[b\frac{n\sigma_{\beta}^2 + n\sigma_{\alpha\beta}^2 + \sigma^2}{n}
+b(b-1)(0)\right] = \frac{n\sigma_{\beta}^2 + n\sigma_{\alpha\beta}^2 + \sigma^2}{bn}
\] \[
\mbox{COV}\left\{\overline{Y}_{i\bullet\bullet},\overline{Y}_{i'\bullet\bullet}\right\}=
\frac{1}{b^2} \mbox{COV}\left\{\sum_{j=1}^b\overline{Y}_{ij\bullet},\sum_{j=1}^b
\overline{Y}_{i'j\bullet}\right\}=
\] \[
\frac{1}{b^2}\left[ b\mbox{COV}\left\{ \overline{Y}_{ij\bullet},
\overline{Y}_{i'j\bullet}\right\} +
b(b-1)\mbox{COV}\left\{ \overline{Y}_{ij\bullet},
\overline{Y}_{i'j'\bullet}\right\}\right] = 
\]

\[
=\frac{1}{b^2}\left[b\sigma_\beta^2 + b(b-1)(0)\right]=
\frac{\sigma_{\beta}^2}{b}=\frac{n\sigma_{\beta}^2}{nb}
\quad i\neq i' 
\]

\[
V\left\{\overline{Y}_{\bullet j\bullet}\right\} = 
V\left\{\frac{1}{a} \sum_{i=1}^a\overline{Y}_{ij\bullet}\right\} =
\frac{1}{a^2}\left[\sum_{i=1}^a V\left\{\overline{Y}_{ij\bullet}\right\} +
2\sum_{i=1}^{a-1}\sum_{i'=i+1}^a
\mbox{COV}\left\{\overline{Y}_{ij\bullet},\overline{Y}_{i'j\bullet}\right\}\right] =
\] \[
=\frac{1}{a^2}\left[a\frac{n\sigma_{\beta}^2 + n\sigma_{\alpha\beta}^2 + \sigma^2}{n} +
 a(a-1)\frac{n\sigma_{\beta}^2}{n}\right] =
\frac{an\sigma_{\beta}^2 + n\sigma_{\alpha\beta}^2 + \sigma^2}{an}
\] \[
\mbox{COV}\left\{\overline{Y}_{\bullet j\bullet},\overline{Y}_{\bullet j'\bullet}\right\}=
\frac{1}{a^2} \mbox{COV}\left\{\sum_{i=1}^a\overline{Y}_{ij\bullet},\sum_{i=1}^a
\overline{Y}_{ij'\bullet}\right\}=
\frac{1}{a^2}\left[ a^2(0)\right] = 0 \qquad j\neq j'
\] \[
V\left\{\overline{Y}_{\bullet\bullet\bullet}\right\}= 
V\left\{\frac{1}{b}\sum_{j=1}^b\overline{Y}_{\bullet j\bullet}\right\}
\frac{1}{b^2}\left[bV\left\{\overline{Y}_{\bullet j\bullet}\right\} +
b(b-1)\mbox{COV}\left\{\overline{Y}_{\bullet j\bullet},
\overline{Y}_{\bullet j'\bullet}\right\}\right]=
\] \[
=\frac{1}{b^2}\left[b\frac{an\sigma_{\beta}^2 + n\sigma_{\alpha\beta}^2 + \sigma^2}{an} +
 b(b-1)(0)\right] =
\frac{an\sigma_{\beta}^2 + n\sigma_{\alpha\beta}^2 + \sigma^2}{abn} 
\] The numerator of the variance of the mean is the expected value of
the mean square for factor B. This leads to the following results and a
Confidence Interval for the overall population mean.

\[
V\left\{\overline{Y}_{\bullet\bullet\bullet}\right\}=
\frac{an\sigma_{\beta}^2 + n\sigma_{\alpha\beta}^2 + \sigma^2}{abn} \qquad \qquad
\hat{V}\left\{\overline{Y}_{\bullet\bullet\bullet}\right\}=\frac{MSB}{abn}
\] \[
(1-\alpha)100\% \mbox{ Confidence Interval for } \mu: \qquad
\overline{Y}_{\bullet\bullet\bullet} \pm t_{1-\alpha/2;b-1}
\sqrt{\frac{MSB}{n}}
\]

For any contrast, \(L=\sum_{i=1}^ac_i\alpha_i\) where
\(\sum_{i=1}^ac_i=0\), among the means of factor A, we obtain the
following variance for its estimator.

\[
L=\sum_{i=1}c_i\left(\mu_{\bullet\bullet}+\alpha_i\right) \qquad \qquad 
\hat{L}=\sum_{i=1}c_i\overline{Y}_{i\bullet\bullet}
\] \[
V\left\{\hat{L} \right\}=
\sum_{i=1}^ac_i^2V\left\{\overline{Y}_{i\bullet\bullet}\right\} + 
2\sum_{i=1}^{a-1}\sum_{i'=i+1}^ac_ic_{i'}
\mbox{COV}\left\{\overline{Y}_{i\bullet\bullet},\overline{Y}_{i'\bullet\bullet}\right\}=
\] \[
\sum_{i=1}^ac_i^2\left(\frac{n\sigma_{\beta}^2 + n\sigma_{\alpha\beta}^2 + \sigma^2}{bn}\right) + 
2\sum_{i=1}^{a-1}\sum_{i'=i+1}^ac_ic_{i'}\left(\frac{n\sigma_{\beta}^2}{bn}\right) = 
\] \[
\sum_{i=1}^ac_i^2\left(\frac{n\sigma_{\alpha\beta}^2 + \sigma^2}{bn}\right)  +
\sum_{i=1}^a\sum_{i'=1}^ac_ic_{i'}\left(\frac{n\sigma_{\beta}^2}{bn}\right) = 
\sum_{i=1}^ac_i^2\left(\frac{n\sigma_{\alpha\beta}^2 + \sigma^2}{bn}\right) +
\left(\frac{n\sigma_{\beta}^2}{bn}\right)\sum_{i=1}^ac_i\sum_{i'=1}^ac_{i'}
\] Since, by definition, the contrast coefficients sum to zero, the last
term is 0 and we obtain the following results.

\[
V\left\{\hat{L} \right\}=\sum_{i=1}^ac_i^2\left(\frac{n\sigma_{\alpha\beta}^2 + \sigma^2}{bn}\right)   \qquad  E\{MSAB\} = n\sigma_{\alpha\beta}^2 + \sigma^2  \qquad
\hat{V}\left\{\hat{L} \right\}=\sum_{i=1}^ac_i^2\left(\frac{MSAB}{bn}\right)
\] These results can be used to test whether a contrast is equal to 0
and/or obtain a Confidence Interval for the contrast.

\[
H_0: L=0 \quad H_A: L \neq 0 \qquad  
TS: t^*=\frac{\hat{L}}{\sqrt{\sum_{i=1}^ac_i^2\left(\frac{MSAB}{bn}\right)}}
\qquad RR:|t^*| \geq t_{1-\alpha/2;(a-1)(b-1)}
\]

\[
(1-\alpha)100\% \mbox{ Confidence Interval for } L: \qquad
\hat{L} \pm t_{1-\alpha/2;(a-1)(b-1)} 
\sqrt{\sum_{i=1}^ac_i^2\left(\frac{MSAB}{bn}\right)}
\]

In particular, when making pairwise comparisons among treatment means,
we have \(\sum_{i=1}^ac_i^2=2\) and we have the following results.

\[
V\left\{\overline{Y}_{i\bullet\bullet} - \overline{Y}_{i'\bullet\bullet}\right\}=
\frac{2\left(n\sigma_{\alpha\beta}^2 + \sigma^2\right)}{bn}  \qquad \qquad
\hat{V}\left\{\overline{Y}_{i\bullet\bullet} - \overline{Y}_{i'\bullet\bullet}\right\}=
\frac{2MSAB}{bn}
\]

Tukey's method can be used to compare all levels of factor A.

\[
HSD = q_{1-\alpha;a,(a-1)(b-1)}\sqrt{\frac{MSAB}{bn}}
\] \[
(1-\alpha)100\% \mbox{ Confidence Interval for } \alpha_i-\alpha_{i'}: \quad
\left(\overline{Y}_{i\bullet\bullet}-\overline{Y}_{i'\bullet\bullet}\right) \pm HSD
\]

\textbf{Example 10.7 - Breath Alcohol Concentration Measurement
Comparison}

A study compared \(a=6\) (fixed) instruments among \(b=3\) (random)
subjects in measuring breath alcohol levels
(\citeproc{ref-gullberg_2008}{Gullberg 2008}). Each subject was measured
by each machine \(n=10\) times. Scores have been multiplied by 100 to
avoid very small variances.

R code for direct calculations and the \textbf{lmer} function are used
for the analysis. The confidence level for the \textbf{difflsmeans} has
been adjusted by the Bonferroni method to construct simultaneous 95\%
Confidence Intervals among the \(6(6-1)/2=15\) pairs of instruments.

\[
\hat{SE}\left\{\overline{Y}_{i\bullet\bullet}-\overline{Y}_{i'\bullet\bullet}\right\}=
\sqrt{\frac{2MSAB}{bn}}=\sqrt{\frac{2(0.3083)}{3(10)}}=0.1434
\qquad t_{1-.05/(2(15));10} = 3.827
\] \[
95\% \mbox{ CI for } \alpha_i-\alpha_{i'}: \quad
\left(\overline{Y}_{i\bullet\bullet}-\overline{Y}_{i'\bullet\bullet}\right)
\pm 3.827(0.1434) 
\] \[
\quad \equiv \quad
95\% \mbox{ CI for } \alpha_i-\alpha_{i'}: \quad
\left(\overline{Y}_{i\bullet\bullet}-\overline{Y}_{i'\bullet\bullet}\right)
\pm 0.549
\] Note that the \(P\)-values are not adjusted for the multiple tests.

\begin{Shaded}
\begin{Highlighting}[]
\NormalTok{ba1 }\OtherTok{\textless{}{-}} \FunctionTok{read.csv}\NormalTok{(}\StringTok{"data/breathalyzer.csv"}\NormalTok{)}
\FunctionTok{head}\NormalTok{(ba1,}\DecValTok{3}\NormalTok{); }\FunctionTok{tail}\NormalTok{(ba1,}\DecValTok{3}\NormalTok{)}
\end{Highlighting}
\end{Shaded}

\begin{verbatim}
##   instrument subject repNum breathAlc
## 1          1       1      1    0.0938
## 2          1       1      2    0.0946
## 3          1       1      3    0.0944
\end{verbatim}

\begin{verbatim}
##     instrument subject repNum breathAlc
## 178          6       2      8    0.0597
## 179          6       2      9    0.0599
## 180          6       2     10    0.0590
\end{verbatim}

\begin{Shaded}
\begin{Highlighting}[]
\DocumentationTok{\#\# Create a new data frame that orders first by instrument then subject within instrument}
\NormalTok{ba }\OtherTok{\textless{}{-}}\NormalTok{ ba1[}\FunctionTok{order}\NormalTok{(ba1}\SpecialCharTok{$}\NormalTok{instrument, ba1}\SpecialCharTok{$}\NormalTok{subject),]}
\FunctionTok{head}\NormalTok{(ba,}\DecValTok{3}\NormalTok{)}
\end{Highlighting}
\end{Shaded}

\begin{verbatim}
##   instrument subject repNum breathAlc
## 1          1       1      1    0.0938
## 2          1       1      2    0.0946
## 3          1       1      3    0.0944
\end{verbatim}

\begin{Shaded}
\begin{Highlighting}[]
\FunctionTok{tail}\NormalTok{(ba,}\DecValTok{3}\NormalTok{)}
\end{Highlighting}
\end{Shaded}

\begin{verbatim}
##     instrument subject repNum breathAlc
## 168          6       3      8    0.0703
## 169          6       3      9    0.0697
## 170          6       3     10    0.0745
\end{verbatim}

\begin{Shaded}
\begin{Highlighting}[]
\NormalTok{ba}\SpecialCharTok{$}\NormalTok{instrument.f }\OtherTok{\textless{}{-}} \FunctionTok{factor}\NormalTok{(ba}\SpecialCharTok{$}\NormalTok{instrument)}
\NormalTok{ba}\SpecialCharTok{$}\NormalTok{subject.f }\OtherTok{\textless{}{-}} \FunctionTok{factor}\NormalTok{(ba}\SpecialCharTok{$}\NormalTok{subject)}

\NormalTok{Y }\OtherTok{\textless{}{-}} \DecValTok{100}\SpecialCharTok{*}\NormalTok{ba}\SpecialCharTok{$}\NormalTok{breathAlc}
\NormalTok{instMean }\OtherTok{\textless{}{-}} \FunctionTok{as.vector}\NormalTok{(}\FunctionTok{tapply}\NormalTok{(Y, ba}\SpecialCharTok{$}\NormalTok{instrument, mean))}
\NormalTok{subjMean }\OtherTok{\textless{}{-}} \FunctionTok{as.vector}\NormalTok{(}\FunctionTok{tapply}\NormalTok{(Y, ba}\SpecialCharTok{$}\NormalTok{subject, mean))}
\NormalTok{inSuMean }\OtherTok{\textless{}{-}} \FunctionTok{as.vector}\NormalTok{(}\FunctionTok{tapply}\NormalTok{(Y, }\FunctionTok{list}\NormalTok{(ba}\SpecialCharTok{$}\NormalTok{subject,ba}\SpecialCharTok{$}\NormalTok{instrument), mean))}
\NormalTok{a }\OtherTok{\textless{}{-}} \FunctionTok{length}\NormalTok{(instMean)}
\NormalTok{b }\OtherTok{\textless{}{-}} \FunctionTok{length}\NormalTok{(subjMean)}
\NormalTok{n }\OtherTok{\textless{}{-}} \FunctionTok{length}\NormalTok{(ba}\SpecialCharTok{$}\NormalTok{breathAlc)}\SpecialCharTok{/}\NormalTok{(a}\SpecialCharTok{*}\NormalTok{b)}

\NormalTok{Ybar\_i.. }\OtherTok{\textless{}{-}} \FunctionTok{rep}\NormalTok{(instMean, }\AttributeTok{each=}\NormalTok{b}\SpecialCharTok{*}\NormalTok{n)}
\NormalTok{Ybar\_.j. }\OtherTok{\textless{}{-}} \FunctionTok{rep}\NormalTok{(subjMean, }\AttributeTok{each=}\NormalTok{n, }\AttributeTok{times=}\NormalTok{a)}
\NormalTok{Ybar\_ij. }\OtherTok{\textless{}{-}} \FunctionTok{rep}\NormalTok{(inSuMean, }\AttributeTok{each=}\NormalTok{n)}
\NormalTok{Ybar\_... }\OtherTok{\textless{}{-}} \FunctionTok{rep}\NormalTok{(}\FunctionTok{mean}\NormalTok{(Y), a}\SpecialCharTok{*}\NormalTok{b}\SpecialCharTok{*}\NormalTok{n)}
\DocumentationTok{\#\# cbind(Y,Ybar\_i..,Ybar\_.j.,Ybar\_ij., Ybar\_...)}
\NormalTok{SSA }\OtherTok{\textless{}{-}} \FunctionTok{sum}\NormalTok{((Ybar\_i..}\SpecialCharTok{{-}}\NormalTok{Ybar\_...)}\SpecialCharTok{\^{}}\DecValTok{2}\NormalTok{); df\_A }\OtherTok{\textless{}{-}}\NormalTok{ a}\DecValTok{{-}1}\NormalTok{; MSA }\OtherTok{\textless{}{-}}\NormalTok{ SSA}\SpecialCharTok{/}\NormalTok{df\_A}
\NormalTok{SSB }\OtherTok{\textless{}{-}} \FunctionTok{sum}\NormalTok{((Ybar\_.j.}\SpecialCharTok{{-}}\NormalTok{Ybar\_...)}\SpecialCharTok{\^{}}\DecValTok{2}\NormalTok{); df\_B }\OtherTok{\textless{}{-}}\NormalTok{ b}\DecValTok{{-}1}\NormalTok{; MSB }\OtherTok{\textless{}{-}}\NormalTok{ SSB}\SpecialCharTok{/}\NormalTok{df\_B}
\NormalTok{SSAB }\OtherTok{\textless{}{-}} \FunctionTok{sum}\NormalTok{((Ybar\_ij.}\SpecialCharTok{{-}}\NormalTok{Ybar\_i..}\SpecialCharTok{{-}}\NormalTok{Ybar\_.j.}\SpecialCharTok{+}\NormalTok{Ybar\_...)}\SpecialCharTok{\^{}}\DecValTok{2}\NormalTok{); df\_AB }\OtherTok{\textless{}{-}}\NormalTok{ (a}\DecValTok{{-}1}\NormalTok{)}\SpecialCharTok{*}\NormalTok{(b}\DecValTok{{-}1}\NormalTok{)}
\NormalTok{MSAB }\OtherTok{\textless{}{-}}\NormalTok{ SSAB}\SpecialCharTok{/}\NormalTok{df\_AB}
\NormalTok{SSE }\OtherTok{\textless{}{-}} \FunctionTok{sum}\NormalTok{((Y}\SpecialCharTok{{-}}\NormalTok{Ybar\_ij.)}\SpecialCharTok{\^{}}\DecValTok{2}\NormalTok{);  df\_E }\OtherTok{\textless{}{-}}\NormalTok{ a}\SpecialCharTok{*}\NormalTok{b}\SpecialCharTok{*}\NormalTok{(n}\DecValTok{{-}1}\NormalTok{); MSE }\OtherTok{\textless{}{-}}\NormalTok{ SSE}\SpecialCharTok{/}\NormalTok{df\_E}
\NormalTok{SSTO }\OtherTok{\textless{}{-}} \FunctionTok{sum}\NormalTok{((Y}\SpecialCharTok{{-}}\NormalTok{Ybar\_...)}\SpecialCharTok{\^{}}\DecValTok{2}\NormalTok{); df\_TO }\OtherTok{\textless{}{-}}\NormalTok{ a}\SpecialCharTok{*}\NormalTok{b}\SpecialCharTok{*}\NormalTok{n}\DecValTok{{-}1}

\NormalTok{aov.out1 }\OtherTok{\textless{}{-}} \FunctionTok{cbind}\NormalTok{(df\_A, SSA, MSA, MSA}\SpecialCharTok{/}\NormalTok{MSAB, }\FunctionTok{qf}\NormalTok{(.}\DecValTok{95}\NormalTok{,df\_A,df\_AB),}
                  \DecValTok{1}\SpecialCharTok{{-}}\FunctionTok{pf}\NormalTok{(MSA}\SpecialCharTok{/}\NormalTok{MSAB,df\_A,df\_AB))}
\NormalTok{aov.out2 }\OtherTok{\textless{}{-}} \FunctionTok{cbind}\NormalTok{(df\_B, SSB, MSB, MSB}\SpecialCharTok{/}\NormalTok{MSAB, }\FunctionTok{qf}\NormalTok{(.}\DecValTok{95}\NormalTok{,df\_B,df\_AB),}
                  \DecValTok{1}\SpecialCharTok{{-}}\FunctionTok{pf}\NormalTok{(MSB}\SpecialCharTok{/}\NormalTok{MSAB,df\_B,df\_AB))}
\NormalTok{aov.out3 }\OtherTok{\textless{}{-}} \FunctionTok{cbind}\NormalTok{(df\_AB, SSAB, MSAB, MSAB}\SpecialCharTok{/}\NormalTok{MSE, }\FunctionTok{qf}\NormalTok{(.}\DecValTok{95}\NormalTok{,df\_AB,df\_E),}
                  \DecValTok{1}\SpecialCharTok{{-}}\FunctionTok{pf}\NormalTok{(MSAB}\SpecialCharTok{/}\NormalTok{MSE,df\_AB,df\_E))}
\NormalTok{aov.out4 }\OtherTok{\textless{}{-}} \FunctionTok{cbind}\NormalTok{(df\_E, SSE, MSE,}\ConstantTok{NA}\NormalTok{,}\ConstantTok{NA}\NormalTok{,}\ConstantTok{NA}\NormalTok{)}
\NormalTok{aov.out5 }\OtherTok{\textless{}{-}} \FunctionTok{cbind}\NormalTok{(df\_TO,SSTO,}\ConstantTok{NA}\NormalTok{,}\ConstantTok{NA}\NormalTok{,}\ConstantTok{NA}\NormalTok{,}\ConstantTok{NA}\NormalTok{)}
\NormalTok{aov.out }\OtherTok{\textless{}{-}} \FunctionTok{rbind}\NormalTok{(aov.out1,aov.out2,aov.out3,aov.out4, aov.out5)}
\FunctionTok{colnames}\NormalTok{(aov.out) }\OtherTok{\textless{}{-}} \FunctionTok{c}\NormalTok{(}\StringTok{"df"}\NormalTok{,}\StringTok{"SS"}\NormalTok{,}\StringTok{"MS"}\NormalTok{,}\StringTok{"F*"}\NormalTok{,}\StringTok{"F(.95)"}\NormalTok{, }\StringTok{"P(\textgreater{}F*)"}\NormalTok{)}
\FunctionTok{rownames}\NormalTok{(aov.out) }\OtherTok{\textless{}{-}} \FunctionTok{c}\NormalTok{(}\StringTok{"Instrument"}\NormalTok{, }\StringTok{"Subject"}\NormalTok{, }\StringTok{"InstxSubj"}\NormalTok{, }\StringTok{"Error"}\NormalTok{, }\StringTok{"Total"}\NormalTok{)}
\FunctionTok{round}\NormalTok{(aov.out,}\DecValTok{4}\NormalTok{)}
\end{Highlighting}
\end{Shaded}

\begin{verbatim}
##             df       SS       MS       F* F(.95) P(>F*)
## Instrument   5  19.0059   3.8012  12.3290 3.3258  5e-04
## Subject      2 219.0440 109.5220 355.2307 4.1028  0e+00
## InstxSubj   10   3.0831   0.3083   5.5595 1.8896  0e+00
## Error      162   8.9840   0.0555       NA     NA     NA
## Total      179 250.1171       NA       NA     NA     NA
\end{verbatim}

\begin{Shaded}
\begin{Highlighting}[]
\NormalTok{s2 }\OtherTok{\textless{}{-}}\NormalTok{ MSE}
\NormalTok{s2\_AB }\OtherTok{\textless{}{-}}\NormalTok{ (MSAB}\SpecialCharTok{{-}}\NormalTok{MSE)}\SpecialCharTok{/}\NormalTok{n}
\NormalTok{s2\_B }\OtherTok{\textless{}{-}}\NormalTok{ (MSB}\SpecialCharTok{{-}}\NormalTok{MSAB)}\SpecialCharTok{/}\NormalTok{(a}\SpecialCharTok{*}\NormalTok{n)}
\NormalTok{df\_s2\_AB }\OtherTok{\textless{}{-}}\NormalTok{ s2\_AB}\SpecialCharTok{\^{}}\DecValTok{2} \SpecialCharTok{/}\NormalTok{ ((MSAB}\SpecialCharTok{/}\NormalTok{n)}\SpecialCharTok{\^{}}\DecValTok{2}\SpecialCharTok{/}\NormalTok{df\_AB }\SpecialCharTok{+}\NormalTok{ (}\SpecialCharTok{{-}}\NormalTok{MSE}\SpecialCharTok{/}\NormalTok{n)}\SpecialCharTok{\^{}}\DecValTok{2}\SpecialCharTok{/}\NormalTok{df\_E)}
\NormalTok{df\_s2\_B }\OtherTok{\textless{}{-}}\NormalTok{ s2\_B}\SpecialCharTok{\^{}}\DecValTok{2} \SpecialCharTok{/}\NormalTok{ ((MSB}\SpecialCharTok{/}\NormalTok{(a}\SpecialCharTok{*}\NormalTok{n))}\SpecialCharTok{\^{}}\DecValTok{2}\SpecialCharTok{/}\NormalTok{df\_B }\SpecialCharTok{+}\NormalTok{ (}\SpecialCharTok{{-}}\NormalTok{MSAB}\SpecialCharTok{/}\NormalTok{(a}\SpecialCharTok{*}\NormalTok{n))}\SpecialCharTok{\^{}}\DecValTok{2}\SpecialCharTok{/}\NormalTok{df\_AB)}

\NormalTok{s2.out1 }\OtherTok{\textless{}{-}} \FunctionTok{cbind}\NormalTok{(s2\_B, df\_s2\_B, df\_s2\_B}\SpecialCharTok{*}\NormalTok{s2\_B}\SpecialCharTok{/}\FunctionTok{qchisq}\NormalTok{(.}\DecValTok{975}\NormalTok{,df\_s2\_B), }
\NormalTok{                 df\_s2\_B}\SpecialCharTok{*}\NormalTok{s2\_B}\SpecialCharTok{/}\FunctionTok{qchisq}\NormalTok{(.}\DecValTok{025}\NormalTok{,df\_s2\_B))}
\NormalTok{s2.out2 }\OtherTok{\textless{}{-}} \FunctionTok{cbind}\NormalTok{(s2\_AB, df\_s2\_AB, df\_s2\_AB}\SpecialCharTok{*}\NormalTok{s2\_AB}\SpecialCharTok{/}\FunctionTok{qchisq}\NormalTok{(.}\DecValTok{975}\NormalTok{,df\_s2\_AB), }
\NormalTok{                 df\_s2\_AB}\SpecialCharTok{*}\NormalTok{s2\_AB}\SpecialCharTok{/}\FunctionTok{qchisq}\NormalTok{(.}\DecValTok{025}\NormalTok{,df\_s2\_AB))}
\NormalTok{s2.out3 }\OtherTok{\textless{}{-}} \FunctionTok{cbind}\NormalTok{(s2, df\_E, df\_E}\SpecialCharTok{*}\NormalTok{s2}\SpecialCharTok{/}\FunctionTok{qchisq}\NormalTok{(.}\DecValTok{975}\NormalTok{,df\_E), }
\NormalTok{                 df\_E}\SpecialCharTok{*}\NormalTok{s2}\SpecialCharTok{/}\FunctionTok{qchisq}\NormalTok{(.}\DecValTok{025}\NormalTok{,df\_E))}
\NormalTok{s2.out }\OtherTok{\textless{}{-}} \FunctionTok{rbind}\NormalTok{(s2.out1,s2.out2,s2.out3)}
\FunctionTok{colnames}\NormalTok{(s2.out) }\OtherTok{\textless{}{-}} \FunctionTok{c}\NormalTok{(}\StringTok{"s\^{}2"}\NormalTok{, }\StringTok{"Satt df"}\NormalTok{, }\StringTok{"LB"}\NormalTok{, }\StringTok{"UB"}\NormalTok{)}
\FunctionTok{rownames}\NormalTok{(s2.out) }\OtherTok{\textless{}{-}} \FunctionTok{c}\NormalTok{(}\StringTok{"Subject"}\NormalTok{, }\StringTok{"InstxSubj"}\NormalTok{, }\StringTok{"Error"}\NormalTok{)}
\FunctionTok{round}\NormalTok{(s2.out,}\DecValTok{4}\NormalTok{)}
\end{Highlighting}
\end{Shaded}

\begin{verbatim}
##              s^2  Satt df     LB      UB
## Subject   1.8202   1.9888 0.4922 73.1909
## InstxSubj 0.0253   6.7127 0.0109  0.1091
## Error     0.0555 162.0000 0.0451  0.0698
\end{verbatim}

\begin{Shaded}
\begin{Highlighting}[]
\NormalTok{HSD }\OtherTok{\textless{}{-}} \FunctionTok{qtukey}\NormalTok{(.}\DecValTok{95}\NormalTok{,a,df\_AB) }\SpecialCharTok{*} \FunctionTok{sqrt}\NormalTok{(MSAB}\SpecialCharTok{/}\NormalTok{(b}\SpecialCharTok{*}\NormalTok{n))}
\NormalTok{HSD.out }\OtherTok{\textless{}{-}} \FunctionTok{matrix}\NormalTok{(}\FunctionTok{rep}\NormalTok{(}\DecValTok{0}\NormalTok{,}\DecValTok{8}\SpecialCharTok{*}\NormalTok{a}\SpecialCharTok{*}\NormalTok{(a}\DecValTok{{-}1}\NormalTok{)}\SpecialCharTok{/}\DecValTok{2}\NormalTok{), }\AttributeTok{ncol=}\DecValTok{8}\NormalTok{)}
\NormalTok{row.count }\OtherTok{\textless{}{-}} \DecValTok{0}
\ControlFlowTok{for}\NormalTok{ (i1 }\ControlFlowTok{in} \DecValTok{2}\SpecialCharTok{:}\NormalTok{a) \{}
  \ControlFlowTok{for}\NormalTok{ (i2 }\ControlFlowTok{in} \DecValTok{1}\SpecialCharTok{:}\NormalTok{(i1}\DecValTok{{-}1}\NormalTok{)) \{}
\NormalTok{    row.count }\OtherTok{\textless{}{-}}\NormalTok{ row.count }\SpecialCharTok{+} \DecValTok{1}
\NormalTok{    HSD.out[row.count,}\DecValTok{1}\NormalTok{] }\OtherTok{\textless{}{-}}\NormalTok{ i1}
\NormalTok{    HSD.out[row.count,}\DecValTok{2}\NormalTok{] }\OtherTok{\textless{}{-}}\NormalTok{ i2}
\NormalTok{    HSD.out[row.count,}\DecValTok{3}\NormalTok{] }\OtherTok{\textless{}{-}}\NormalTok{ instMean[i1]}
\NormalTok{    HSD.out[row.count,}\DecValTok{4}\NormalTok{] }\OtherTok{\textless{}{-}}\NormalTok{ instMean[i2]}
\NormalTok{    HSD.out[row.count,}\DecValTok{5}\NormalTok{] }\OtherTok{\textless{}{-}}\NormalTok{ instMean[i1] }\SpecialCharTok{{-}}\NormalTok{ instMean[i2]}
\NormalTok{    HSD.out[row.count,}\DecValTok{6}\NormalTok{] }\OtherTok{\textless{}{-}}\NormalTok{ (instMean[i1] }\SpecialCharTok{{-}}\NormalTok{ instMean[i2]) }\SpecialCharTok{{-}}\NormalTok{ HSD}
\NormalTok{    HSD.out[row.count,}\DecValTok{7}\NormalTok{] }\OtherTok{\textless{}{-}}\NormalTok{ (instMean[i1] }\SpecialCharTok{{-}}\NormalTok{ instMean[i2]) }\SpecialCharTok{+}\NormalTok{ HSD}
\NormalTok{    HSD.out[row.count,}\DecValTok{8}\NormalTok{] }\OtherTok{\textless{}{-}} 
      \DecValTok{1}\SpecialCharTok{{-}}\FunctionTok{ptukey}\NormalTok{(}\FunctionTok{abs}\NormalTok{(instMean[i1] }\SpecialCharTok{{-}}\NormalTok{ instMean[i2])}\SpecialCharTok{/}\FunctionTok{sqrt}\NormalTok{(MSAB}\SpecialCharTok{/}\NormalTok{(b}\SpecialCharTok{*}\NormalTok{n)),a,df\_AB)}
\NormalTok{  \}}
\NormalTok{\}}

\FunctionTok{colnames}\NormalTok{(HSD.out) }\OtherTok{\textless{}{-}} \FunctionTok{c}\NormalTok{(}\StringTok{"i"}\NormalTok{, }\StringTok{"i\textquotesingle{}"}\NormalTok{, }\StringTok{"Ybar\_i"}\NormalTok{, }\StringTok{"Ybar\_i\textquotesingle{}"}\NormalTok{, }\StringTok{"Diff"}\NormalTok{, }\StringTok{"LB"}\NormalTok{, }\StringTok{"UB"}\NormalTok{, }\StringTok{"p adj"}\NormalTok{)}
\FunctionTok{round}\NormalTok{(HSD.out,}\DecValTok{4}\NormalTok{)}
\end{Highlighting}
\end{Shaded}

\begin{verbatim}
##       i i' Ybar_i Ybar_i'    Diff      LB      UB  p adj
##  [1,] 2  1 7.3547  7.8340 -0.4793 -0.9773  0.0186 0.0609
##  [2,] 3  1 7.6247  7.8340 -0.2093 -0.7073  0.2886 0.6941
##  [3,] 3  2 7.6247  7.3547  0.2700 -0.2280  0.7680 0.4621
##  [4,] 4  1 6.8353  7.8340 -0.9987 -1.4966 -0.5007 0.0004
##  [5,] 4  2 6.8353  7.3547 -0.5193 -1.0173 -0.0214 0.0399
##  [6,] 4  3 6.8353  7.6247 -0.7893 -1.2873 -0.2914 0.0026
##  [7,] 5  1 7.5943  7.8340 -0.2397 -0.7376  0.2583 0.5760
##  [8,] 5  2 7.5943  7.3547  0.2397 -0.2583  0.7376 0.5760
##  [9,] 5  3 7.5943  7.6247 -0.0303 -0.5283  0.4676 0.9999
## [10,] 5  4 7.5943  6.8353  0.7590  0.2610  1.2570 0.0034
## [11,] 6  1 7.2090  7.8340 -0.6250 -1.1230 -0.1270 0.0131
## [12,] 6  2 7.2090  7.3547 -0.1457 -0.6436  0.3523 0.9023
## [13,] 6  3 7.2090  7.6247 -0.4157 -0.9136  0.0823 0.1188
## [14,] 6  4 7.2090  6.8353  0.3737 -0.1243  0.8716 0.1818
## [15,] 6  5 7.2090  7.5943 -0.3853 -0.8833  0.1126 0.1618
\end{verbatim}

\begin{Shaded}
\begin{Highlighting}[]
\DocumentationTok{\#\# The error terms are incorrect for instrument and subject using the aov command}
\NormalTok{mod1 }\OtherTok{\textless{}{-}} \FunctionTok{aov}\NormalTok{(}\DecValTok{100}\SpecialCharTok{*}\NormalTok{breathAlc }\SpecialCharTok{\textasciitilde{}} \FunctionTok{factor}\NormalTok{(instrument)}\SpecialCharTok{*}\FunctionTok{factor}\NormalTok{(subject), }\AttributeTok{data=}\NormalTok{ba)}
\FunctionTok{anova}\NormalTok{(mod1)}
\end{Highlighting}
\end{Shaded}

\begin{verbatim}
## Analysis of Variance Table
## 
## Response: 100 * breathAlc
##                                     Df  Sum Sq Mean Sq   F value    Pr(>F)    
## factor(instrument)                   5  19.006   3.801   68.5430 < 2.2e-16 ***
## factor(subject)                      2 219.044 109.522 1974.9028 < 2.2e-16 ***
## factor(instrument):factor(subject)  10   3.083   0.308    5.5595 4.177e-07 ***
## Residuals                          162   8.984   0.055                        
## ---
## Signif. codes:  0 '***' 0.001 '**' 0.01 '*' 0.05 '.' 0.1 ' ' 1
\end{verbatim}

\begin{Shaded}
\begin{Highlighting}[]
\DocumentationTok{\#\# Fit unrestricted mixed model }
\FunctionTok{options}\NormalTok{(}\AttributeTok{contrasts=}\FunctionTok{c}\NormalTok{(}\StringTok{"contr.sum"}\NormalTok{, }\StringTok{"contr.poly"}\NormalTok{))}
\NormalTok{mod2 }\OtherTok{\textless{}{-}} \FunctionTok{lmer}\NormalTok{(}\DecValTok{100}\SpecialCharTok{*}\NormalTok{breathAlc }\SpecialCharTok{\textasciitilde{}}\NormalTok{ instrument.f }\SpecialCharTok{+}\NormalTok{ (}\DecValTok{1}\SpecialCharTok{|}\NormalTok{subject.f) }\SpecialCharTok{+}\NormalTok{ (}\DecValTok{1}\SpecialCharTok{|}\NormalTok{instrument.f}\SpecialCharTok{:}\NormalTok{subject.f), }\AttributeTok{data=}\NormalTok{ba)}
\FunctionTok{summary}\NormalTok{(mod2)}
\end{Highlighting}
\end{Shaded}

\begin{verbatim}
## Linear mixed model fit by REML. t-tests use Satterthwaite's method [
## lmerModLmerTest]
## Formula: 
## 100 * breathAlc ~ instrument.f + (1 | subject.f) + (1 | instrument.f:subject.f)
##    Data: ba
## 
## REML criterion at convergence: 46.9
## 
## Scaled residuals: 
##     Min      1Q  Median      3Q     Max 
## -2.7063 -0.5975  0.0655  0.5337  3.5359 
## 
## Random effects:
##  Groups                 Name        Variance Std.Dev.
##  instrument.f:subject.f (Intercept) 0.02529  0.1590  
##  subject.f              (Intercept) 1.82025  1.3492  
##  Residual                           0.05546  0.2355  
## Number of obs: 180, groups:  instrument.f:subject.f, 18; subject.f, 3
## 
## Fixed effects:
##               Estimate Std. Error       df t value Pr(>|t|)    
## (Intercept)    7.40867    0.78004  1.99996   9.498 0.010905 *  
## instrument.f1  0.42533    0.09254  9.99985   4.596 0.000986 ***
## instrument.f2 -0.05400    0.09254  9.99985  -0.584 0.572476    
## instrument.f3  0.21600    0.09254  9.99985   2.334 0.041762 *  
## instrument.f4 -0.57333    0.09254  9.99985  -6.195 0.000102 ***
## instrument.f5  0.18567    0.09254  9.99985   2.006 0.072627 .  
## ---
## Signif. codes:  0 '***' 0.001 '**' 0.01 '*' 0.05 '.' 0.1 ' ' 1
## 
## Correlation of Fixed Effects:
##             (Intr) inst.1 inst.2 inst.3 inst.4
## instrmnt.f1  0.000                            
## instrmnt.f2  0.000 -0.200                     
## instrmnt.f3  0.000 -0.200 -0.200              
## instrmnt.f4  0.000 -0.200 -0.200 -0.200       
## instrmnt.f5  0.000 -0.200 -0.200 -0.200 -0.200
\end{verbatim}

\begin{Shaded}
\begin{Highlighting}[]
\FunctionTok{anova}\NormalTok{(mod2)}
\end{Highlighting}
\end{Shaded}

\begin{verbatim}
## Type III Analysis of Variance Table with Satterthwaite's method
##              Sum Sq Mean Sq NumDF  DenDF F value    Pr(>F)    
## instrument.f 3.4186 0.68372     5 9.9999  12.329 0.0005164 ***
## ---
## Signif. codes:  0 '***' 0.001 '**' 0.01 '*' 0.05 '.' 0.1 ' ' 1
\end{verbatim}

\begin{Shaded}
\begin{Highlighting}[]
\FunctionTok{difflsmeans}\NormalTok{(mod2, }\AttributeTok{level=}\DecValTok{1}\FloatTok{{-}.05}\SpecialCharTok{/}\NormalTok{(a}\SpecialCharTok{*}\NormalTok{(a}\DecValTok{{-}1}\NormalTok{)}\SpecialCharTok{/}\DecValTok{2}\NormalTok{)) }\DocumentationTok{\#\# Bonferroni adjustment for CI\textquotesingle{}s (not P{-}values)}
\end{Highlighting}
\end{Shaded}

\begin{verbatim}
## Least Squares Means table:
## 
##                                Estimate Std. Error df t value     lower
## instrument.f1 - instrument.f2  0.479333   0.143368 10  3.3434 -0.069387
## instrument.f1 - instrument.f3  0.209333   0.143368 10  1.4601 -0.339387
## instrument.f1 - instrument.f4  0.998667   0.143368 10  6.9658  0.449947
## instrument.f1 - instrument.f5  0.239667   0.143368 10  1.6717 -0.309053
## instrument.f1 - instrument.f6  0.625000   0.143368 10  4.3594  0.076280
## instrument.f2 - instrument.f3 -0.270000   0.143368 10 -1.8833 -0.818720
## instrument.f2 - instrument.f4  0.519333   0.143368 10  3.6224 -0.029387
## instrument.f2 - instrument.f5 -0.239667   0.143368 10 -1.6717 -0.788387
## instrument.f2 - instrument.f6  0.145667   0.143368 10  1.0160 -0.403053
## instrument.f3 - instrument.f4  0.789333   0.143368 10  5.5056  0.240613
## instrument.f3 - instrument.f5  0.030333   0.143368 10  0.2116 -0.518387
## instrument.f3 - instrument.f6  0.415667   0.143368 10  2.8993 -0.133053
## instrument.f4 - instrument.f5 -0.759000   0.143368 10 -5.2941 -1.307720
## instrument.f4 - instrument.f6 -0.373667   0.143368 10 -2.6063 -0.922387
## instrument.f5 - instrument.f6  0.385333   0.143368 10  2.6877 -0.163387
##                                   upper  Pr(>|t|)    
## instrument.f1 - instrument.f2  1.028053 0.0074481 ** 
## instrument.f1 - instrument.f3  0.758053 0.1749512    
## instrument.f1 - instrument.f4  1.547387 3.873e-05 ***
## instrument.f1 - instrument.f5  0.788387 0.1255364    
## instrument.f1 - instrument.f6  1.173720 0.0014225 ** 
## instrument.f2 - instrument.f3  0.278720 0.0890443 .  
## instrument.f2 - instrument.f4  1.068053 0.0046713 ** 
## instrument.f2 - instrument.f5  0.309053 0.1255364    
## instrument.f2 - instrument.f6  0.694387 0.3335656    
## instrument.f3 - instrument.f4  1.338053 0.0002597 ***
## instrument.f3 - instrument.f5  0.579053 0.8366871    
## instrument.f3 - instrument.f6  0.964387 0.0158528 *  
## instrument.f4 - instrument.f5 -0.210280 0.0003505 ***
## instrument.f4 - instrument.f6  0.175053 0.0262048 *  
## instrument.f5 - instrument.f6  0.934053 0.0227889 *  
## ---
## Signif. codes:  0 '***' 0.001 '**' 0.01 '*' 0.05 '.' 0.1 ' ' 1
## 
##   Confidence level: 100%
##   Degrees of freedom method: Satterthwaite
\end{verbatim}

\[ \nabla \]

\section{Three-Way Mixed Effects
Models}\label{three-way-mixed-effects-models}

In this section, the Three-Way mixed model with two fixed factors and a
single random factor (typically subjects) is described. In many cases,
there is a single replicate per combination of factors \((n=1)\) and
this can be thought of as an extension of a Randomized Block/Repeated
Measures Design. Here is the case when \(n=1\), factors A and B are
fixed, and factor C is random (Subject).

\[
SSA = bc\sum_{i=1}^a\left(\overline{Y}_{i\bullet\bullet}-
\overline{Y}_{\bullet\bullet\bullet}\right)^2
\qquad df_A=a-1  
\] \[
E\{MSA\}=\sigma^2 + \sigma_{\alpha\beta\gamma}^2+
b\sigma_{\alpha\gamma}^2+\frac{bc\sum_{i=1}^a\alpha_i^2}{a-1}
=E\{MSAC\}+\frac{bc\sum_{i=1}^a\alpha_i^2}{a-1}
\] \[
SSB = ac\sum_{j=1}^b\left(\overline{Y}_{\bullet j\bullet}-
\overline{Y}_{\bullet\bullet\bullet}\right)^2
\qquad df_B=b-1  
\] \[
E\{MSB\}=\sigma^2 + \sigma_{\alpha\beta\gamma}^2+
a\sigma_{\beta\gamma}^2+\frac{ac\sum_{j=1}^b\beta_j^2}{b-1}
= E\{MSBC\}+\frac{ac\sum_{j=1}^b\beta_j^2}{b-1}
\] \[
SSAB=c\sum_{i=1}^a \sum_{j=1}^b \left(
\overline{Y}_{ij\bullet}-\overline{Y}_{i\bullet\bullet}-
\overline{Y}_{\bullet j \bullet} +
\overline{Y}_{\bullet\bullet\bullet} \right)^2
\qquad df_{AB}=(a-1)(b-1)
\] \[
E\{MSAB\}=\sigma^2 + \sigma_{\alpha\beta\gamma}^2 + 
\frac{c\sum_{i=1}^a\sum_{j=1}^b(\alpha\beta)_{ij}^2}{(a-1)(b-1)}
=E\{MSABC\}+
\frac{c\sum_{i=1}^a\sum_{j=1}^b(\alpha\beta)_{ij}^2}{(a-1)(b-1)}
\]

The error terms for the fixed main effects and their 2-way interaction
are the interactions between each term and the random factor C.

Tukey's HSD or the Bonferroni method can be used to compare the levels
of factors A and B, or the \(ab\) treatment means if the interaction is
important.

\textbf{Example 10.8 - Navigational Techniques for Web Maps}

A study was conducted to measure effects of two factors on subjects'
skill on reading web maps (\citeproc{ref-wu_2011}{Wu, Lin, and You
2011}). The factors included are given below.

\begin{itemize}
\tightlist
\item
  \textbf{Navigational Technique} - Combined Panning Buttons (CPB,
  (\(i=1\))), Distributed Panning Buttons (DPB, (\(i=2\))), Enhanced
  Navigaor with Continuous Control (ENCC, (\(i=3\))), and Grab \& Drag
  (GD, (\(i=4\)))
\item
  \textbf{Input Method} - Touch Screen (DT, (\(j=1\))), Mouse (M,
  (\(j=2\)))
\item
  \textbf{Subject} - 36 subjects who each performed once on the 8
  combinations of Input Method and Navigational Technique
\end{itemize}

One response was task completion time (seconds). We will analyze the
data directly and with the \textbf{lmer} function below.

\begin{Shaded}
\begin{Highlighting}[]
\NormalTok{ntm }\OtherTok{\textless{}{-}} \FunctionTok{read.csv}\NormalTok{(}\StringTok{"data/navtechmap1.csv"}\NormalTok{)}
\FunctionTok{head}\NormalTok{(ntm, }\DecValTok{3}\NormalTok{); }\FunctionTok{tail}\NormalTok{(ntm, }\DecValTok{3}\NormalTok{)}
\end{Highlighting}
\end{Shaded}

\begin{verbatim}
##   navTech inputMeth subject taskTime
## 1       1         1       1   163.30
## 2       1         1       2   214.95
## 3       1         1       3   179.73
\end{verbatim}

\begin{verbatim}
##     navTech inputMeth subject taskTime
## 286       4         2      34   105.90
## 287       4         2      35    73.22
## 288       4         2      36    78.10
\end{verbatim}

\begin{Shaded}
\begin{Highlighting}[]
\NormalTok{ntm}\SpecialCharTok{$}\NormalTok{inputMet.f }\OtherTok{\textless{}{-}} \FunctionTok{factor}\NormalTok{(ntm}\SpecialCharTok{$}\NormalTok{inputMeth, }\AttributeTok{levels=}\DecValTok{1}\SpecialCharTok{:}\DecValTok{2}\NormalTok{, }\AttributeTok{labels=}\FunctionTok{c}\NormalTok{(}\StringTok{"DT"}\NormalTok{,}\StringTok{"M"}\NormalTok{))}
\NormalTok{ntm}\SpecialCharTok{$}\NormalTok{navTech.f }\OtherTok{\textless{}{-}} \FunctionTok{factor}\NormalTok{(ntm}\SpecialCharTok{$}\NormalTok{navTech, }\AttributeTok{levels=}\DecValTok{1}\SpecialCharTok{:}\DecValTok{4}\NormalTok{, }
                        \AttributeTok{labels=}\FunctionTok{c}\NormalTok{(}\StringTok{"CPB"}\NormalTok{,}\StringTok{"DPB"}\NormalTok{,}\StringTok{"ENCC"}\NormalTok{,}\StringTok{"GD"}\NormalTok{))}
\NormalTok{ntm}\SpecialCharTok{$}\NormalTok{subject.f }\OtherTok{\textless{}{-}} \FunctionTok{factor}\NormalTok{(ntm}\SpecialCharTok{$}\NormalTok{subject)}
\NormalTok{allmean }\OtherTok{\textless{}{-}} \FunctionTok{mean}\NormalTok{(ntm}\SpecialCharTok{$}\NormalTok{taskTime)}
\NormalTok{A\_mean }\OtherTok{\textless{}{-}} \FunctionTok{as.vector}\NormalTok{(}\FunctionTok{tapply}\NormalTok{(ntm}\SpecialCharTok{$}\NormalTok{taskTime, ntm}\SpecialCharTok{$}\NormalTok{navTech.f, mean))}
\NormalTok{B\_mean }\OtherTok{\textless{}{-}} \FunctionTok{as.vector}\NormalTok{(}\FunctionTok{tapply}\NormalTok{(ntm}\SpecialCharTok{$}\NormalTok{taskTime, ntm}\SpecialCharTok{$}\NormalTok{inputMet.f, mean))}
\NormalTok{C\_mean }\OtherTok{\textless{}{-}} \FunctionTok{as.vector}\NormalTok{(}\FunctionTok{tapply}\NormalTok{(ntm}\SpecialCharTok{$}\NormalTok{taskTime, ntm}\SpecialCharTok{$}\NormalTok{subject.f, mean))}
\NormalTok{AB\_mean }\OtherTok{\textless{}{-}} \FunctionTok{as.vector}\NormalTok{(}\FunctionTok{tapply}\NormalTok{(ntm}\SpecialCharTok{$}\NormalTok{taskTime, }\FunctionTok{list}\NormalTok{(ntm}\SpecialCharTok{$}\NormalTok{inputMet.f,ntm}\SpecialCharTok{$}\NormalTok{navTech.f), mean))}
\NormalTok{AC\_mean }\OtherTok{\textless{}{-}} \FunctionTok{as.vector}\NormalTok{(}\FunctionTok{tapply}\NormalTok{(ntm}\SpecialCharTok{$}\NormalTok{taskTime, }\FunctionTok{list}\NormalTok{(ntm}\SpecialCharTok{$}\NormalTok{subject.f,ntm}\SpecialCharTok{$}\NormalTok{navTech.f), mean))}
\NormalTok{BC\_mean }\OtherTok{\textless{}{-}} \FunctionTok{as.vector}\NormalTok{(}\FunctionTok{tapply}\NormalTok{(ntm}\SpecialCharTok{$}\NormalTok{taskTime, }\FunctionTok{list}\NormalTok{(ntm}\SpecialCharTok{$}\NormalTok{subject.f,ntm}\SpecialCharTok{$}\NormalTok{inputMet.f), mean))}
\NormalTok{ABC\_mean }\OtherTok{\textless{}{-}}\NormalTok{ ntm}\SpecialCharTok{$}\NormalTok{taskTime}
\NormalTok{a }\OtherTok{\textless{}{-}} \FunctionTok{length}\NormalTok{(A\_mean); b }\OtherTok{\textless{}{-}} \FunctionTok{length}\NormalTok{(B\_mean); c }\OtherTok{\textless{}{-}} \FunctionTok{length}\NormalTok{(C\_mean)}

\NormalTok{Ybar\_... }\OtherTok{\textless{}{-}} \FunctionTok{rep}\NormalTok{(allmean,a}\SpecialCharTok{*}\NormalTok{b}\SpecialCharTok{*}\NormalTok{c)}
\NormalTok{Ybar\_i.. }\OtherTok{\textless{}{-}} \FunctionTok{rep}\NormalTok{(A\_mean, }\AttributeTok{each=}\NormalTok{b}\SpecialCharTok{*}\NormalTok{c)}
\NormalTok{Ybar\_.j. }\OtherTok{\textless{}{-}} \FunctionTok{rep}\NormalTok{(B\_mean, }\AttributeTok{each=}\NormalTok{c, }\AttributeTok{times=}\NormalTok{a)}
\NormalTok{Ybar\_..k }\OtherTok{\textless{}{-}} \FunctionTok{rep}\NormalTok{(C\_mean, }\AttributeTok{times=}\NormalTok{a}\SpecialCharTok{*}\NormalTok{b)}
\NormalTok{Ybar\_ij. }\OtherTok{\textless{}{-}} \FunctionTok{rep}\NormalTok{(AB\_mean, }\AttributeTok{each=}\NormalTok{c)}
\NormalTok{Ybar\_i.k }\OtherTok{\textless{}{-}} \FunctionTok{c}\NormalTok{(}\FunctionTok{rep}\NormalTok{(AC\_mean[}\DecValTok{1}\SpecialCharTok{:}\NormalTok{c], }\AttributeTok{times=}\NormalTok{b),}
              \FunctionTok{rep}\NormalTok{(AC\_mean[(c}\SpecialCharTok{+}\DecValTok{1}\NormalTok{)}\SpecialCharTok{:}\NormalTok{(}\DecValTok{2}\SpecialCharTok{*}\NormalTok{c)], }\AttributeTok{times=}\NormalTok{b),}
              \FunctionTok{rep}\NormalTok{(AC\_mean[(}\DecValTok{2}\SpecialCharTok{*}\NormalTok{c}\SpecialCharTok{+}\DecValTok{1}\NormalTok{)}\SpecialCharTok{:}\NormalTok{(}\DecValTok{3}\SpecialCharTok{*}\NormalTok{c)], }\AttributeTok{times=}\NormalTok{b),}
              \FunctionTok{rep}\NormalTok{(AC\_mean[(}\DecValTok{3}\SpecialCharTok{*}\NormalTok{c}\SpecialCharTok{+}\DecValTok{1}\NormalTok{)}\SpecialCharTok{:}\NormalTok{(}\DecValTok{4}\SpecialCharTok{*}\NormalTok{c)], }\AttributeTok{times=}\NormalTok{b))}
\NormalTok{Ybar\_.jk }\OtherTok{\textless{}{-}} \FunctionTok{rep}\NormalTok{(BC\_mean, }\AttributeTok{times=}\NormalTok{a)}
\NormalTok{Ybar\_ijk }\OtherTok{\textless{}{-}}\NormalTok{ ABC\_mean}

\NormalTok{SSA }\OtherTok{\textless{}{-}} \FunctionTok{sum}\NormalTok{((Ybar\_i..}\SpecialCharTok{{-}}\NormalTok{Ybar\_...)}\SpecialCharTok{\^{}}\DecValTok{2}\NormalTok{); df\_A }\OtherTok{\textless{}{-}}\NormalTok{ a}\DecValTok{{-}1}\NormalTok{; MSA }\OtherTok{\textless{}{-}}\NormalTok{ SSA}\SpecialCharTok{/}\NormalTok{df\_A}
\NormalTok{SSB }\OtherTok{\textless{}{-}} \FunctionTok{sum}\NormalTok{((Ybar\_.j.}\SpecialCharTok{{-}}\NormalTok{Ybar\_...)}\SpecialCharTok{\^{}}\DecValTok{2}\NormalTok{); df\_B }\OtherTok{\textless{}{-}}\NormalTok{ b}\DecValTok{{-}1}\NormalTok{; MSB }\OtherTok{\textless{}{-}}\NormalTok{ SSB}\SpecialCharTok{/}\NormalTok{df\_B}
\NormalTok{SSC }\OtherTok{\textless{}{-}} \FunctionTok{sum}\NormalTok{((Ybar\_..k}\SpecialCharTok{{-}}\NormalTok{Ybar\_...)}\SpecialCharTok{\^{}}\DecValTok{2}\NormalTok{); df\_C }\OtherTok{\textless{}{-}}\NormalTok{ c}\DecValTok{{-}1}\NormalTok{; MSC }\OtherTok{\textless{}{-}}\NormalTok{ SSC}\SpecialCharTok{/}\NormalTok{df\_C}
\NormalTok{SSAB }\OtherTok{\textless{}{-}} \FunctionTok{sum}\NormalTok{((Ybar\_ij.}\SpecialCharTok{{-}}\NormalTok{Ybar\_i..}\SpecialCharTok{{-}}\NormalTok{Ybar\_.j.}\SpecialCharTok{+}\NormalTok{Ybar\_...)}\SpecialCharTok{\^{}}\DecValTok{2}\NormalTok{)}
\NormalTok{df\_AB }\OtherTok{\textless{}{-}}\NormalTok{ (a}\DecValTok{{-}1}\NormalTok{)}\SpecialCharTok{*}\NormalTok{(b}\DecValTok{{-}1}\NormalTok{); MSAB }\OtherTok{\textless{}{-}}\NormalTok{ SSAB}\SpecialCharTok{/}\NormalTok{df\_AB}
\NormalTok{SSAC }\OtherTok{\textless{}{-}} \FunctionTok{sum}\NormalTok{((Ybar\_i.k}\SpecialCharTok{{-}}\NormalTok{Ybar\_i..}\SpecialCharTok{{-}}\NormalTok{Ybar\_..k}\SpecialCharTok{+}\NormalTok{Ybar\_...)}\SpecialCharTok{\^{}}\DecValTok{2}\NormalTok{)}
\NormalTok{df\_AC }\OtherTok{\textless{}{-}}\NormalTok{ (a}\DecValTok{{-}1}\NormalTok{)}\SpecialCharTok{*}\NormalTok{(c}\DecValTok{{-}1}\NormalTok{); MSAC }\OtherTok{\textless{}{-}}\NormalTok{ SSAC}\SpecialCharTok{/}\NormalTok{df\_AC}
\NormalTok{SSBC }\OtherTok{\textless{}{-}} \FunctionTok{sum}\NormalTok{((Ybar\_.jk}\SpecialCharTok{{-}}\NormalTok{Ybar\_.j.}\SpecialCharTok{{-}}\NormalTok{Ybar\_..k}\SpecialCharTok{+}\NormalTok{Ybar\_...)}\SpecialCharTok{\^{}}\DecValTok{2}\NormalTok{)}
\NormalTok{df\_BC }\OtherTok{\textless{}{-}}\NormalTok{ (b}\DecValTok{{-}1}\NormalTok{)}\SpecialCharTok{*}\NormalTok{(c}\DecValTok{{-}1}\NormalTok{); MSBC }\OtherTok{\textless{}{-}}\NormalTok{ SSBC}\SpecialCharTok{/}\NormalTok{df\_BC}
\NormalTok{SSABC }\OtherTok{\textless{}{-}} \FunctionTok{sum}\NormalTok{((Ybar\_ijk}\SpecialCharTok{{-}}\NormalTok{Ybar\_ij.}\SpecialCharTok{{-}}\NormalTok{Ybar\_i.k}\SpecialCharTok{{-}}\NormalTok{Ybar\_.jk}\SpecialCharTok{+}\NormalTok{Ybar\_i..}\SpecialCharTok{+}\NormalTok{Ybar\_.j.}\SpecialCharTok{+}\NormalTok{Ybar\_..k}\SpecialCharTok{{-}}\NormalTok{Ybar\_...)}\SpecialCharTok{\^{}}\DecValTok{2}\NormalTok{)}
\NormalTok{df\_ABC }\OtherTok{\textless{}{-}}\NormalTok{ (a}\DecValTok{{-}1}\NormalTok{)}\SpecialCharTok{*}\NormalTok{(b}\DecValTok{{-}1}\NormalTok{)}\SpecialCharTok{*}\NormalTok{(c}\DecValTok{{-}1}\NormalTok{); MSABC }\OtherTok{\textless{}{-}}\NormalTok{ SSABC}\SpecialCharTok{/}\NormalTok{df\_ABC}

\NormalTok{aov.out1 }\OtherTok{\textless{}{-}} \FunctionTok{cbind}\NormalTok{(df\_A, SSA, MSA, df\_AC, MSAC,MSA}\SpecialCharTok{/}\NormalTok{MSAC, }\FunctionTok{qf}\NormalTok{(.}\DecValTok{95}\NormalTok{,df\_A,df\_AC),}
                  \DecValTok{1}\SpecialCharTok{{-}}\FunctionTok{pf}\NormalTok{(MSA}\SpecialCharTok{/}\NormalTok{MSAC,df\_A,df\_AC))}
\NormalTok{aov.out2 }\OtherTok{\textless{}{-}} \FunctionTok{cbind}\NormalTok{(df\_B, SSB, MSB, df\_BC, MSBC, MSB}\SpecialCharTok{/}\NormalTok{MSBC, }\FunctionTok{qf}\NormalTok{(.}\DecValTok{95}\NormalTok{,df\_B,df\_BC),}
                  \DecValTok{1}\SpecialCharTok{{-}}\FunctionTok{pf}\NormalTok{(MSB}\SpecialCharTok{/}\NormalTok{MSBC,df\_B,df\_BC))}
\NormalTok{aov.out3 }\OtherTok{\textless{}{-}} \FunctionTok{cbind}\NormalTok{(df\_AB, SSAB, MSAB, df\_ABC, MSABC, MSAB}\SpecialCharTok{/}\NormalTok{MSABC, }\FunctionTok{qf}\NormalTok{(.}\DecValTok{95}\NormalTok{,df\_AB,df\_ABC),}
                  \DecValTok{1}\SpecialCharTok{{-}}\FunctionTok{pf}\NormalTok{(MSAB}\SpecialCharTok{/}\NormalTok{MSABC,df\_AB,df\_ABC))}
\NormalTok{aov.out }\OtherTok{\textless{}{-}} \FunctionTok{rbind}\NormalTok{(aov.out1,aov.out2,aov.out3)}
\FunctionTok{colnames}\NormalTok{(aov.out) }\OtherTok{\textless{}{-}} \FunctionTok{c}\NormalTok{(}\StringTok{"df"}\NormalTok{, }\StringTok{"SS"}\NormalTok{, }\StringTok{"MS"}\NormalTok{, }\StringTok{"Err df"}\NormalTok{, }\StringTok{"Err MS"}\NormalTok{, }\StringTok{"F*"}\NormalTok{,}
                       \StringTok{"F(.95)"}\NormalTok{, }\StringTok{"P(\textgreater{}F*)"}\NormalTok{)}
\FunctionTok{rownames}\NormalTok{(aov.out) }\OtherTok{\textless{}{-}} \FunctionTok{c}\NormalTok{(}\StringTok{"Navigation"}\NormalTok{,}\StringTok{"Input"}\NormalTok{,}\StringTok{"Inp/Nav"}\NormalTok{)}
\FunctionTok{round}\NormalTok{(aov.out,}\DecValTok{4}\NormalTok{)}
\end{Highlighting}
\end{Shaded}

\begin{verbatim}
##            df       SS        MS Err df    Err MS      F* F(.95) P(>F*)
## Navigation  3 66995.51 22331.838    105 1417.1179 15.7586 2.6911  0e+00
## Input       1 30635.76 30635.757     35 1960.1337 15.6294 4.1213  4e-04
## Inp/Nav     3 18710.29  6236.762    105  926.4971  6.7316 2.6911  3e-04
\end{verbatim}

\begin{Shaded}
\begin{Highlighting}[]
\NormalTok{mod1 }\OtherTok{\textless{}{-}} \FunctionTok{aov}\NormalTok{(taskTime }\SpecialCharTok{\textasciitilde{}}\NormalTok{ navTech.f }\SpecialCharTok{+}\NormalTok{ inputMet.f }\SpecialCharTok{+}\NormalTok{ subject.f }\SpecialCharTok{+}
\NormalTok{              navTech.f}\SpecialCharTok{:}\NormalTok{inputMet.f }\SpecialCharTok{+}\NormalTok{ navTech.f}\SpecialCharTok{:}\NormalTok{subject.f }\SpecialCharTok{+}
\NormalTok{              inputMet.f}\SpecialCharTok{:}\NormalTok{subject.f, }\AttributeTok{data=}\NormalTok{ntm)}
\FunctionTok{anova}\NormalTok{(mod1)}
\end{Highlighting}
\end{Shaded}

\begin{verbatim}
## Analysis of Variance Table
## 
## Response: taskTime
##                       Df Sum Sq Mean Sq F value    Pr(>F)    
## navTech.f              3  66996 22331.8 24.1035 6.033e-12 ***
## inputMet.f             1  30636 30635.8 33.0662 8.840e-08 ***
## subject.f             35  84008  2400.2  2.5906 0.0001025 ***
## navTech.f:inputMet.f   3  18710  6236.8  6.7316 0.0003376 ***
## navTech.f:subject.f  105 148797  1417.1  1.5295 0.0152294 *  
## inputMet.f:subject.f  35  68605  1960.1  2.1156 0.0018646 ** 
## Residuals            105  97282   926.5                      
## ---
## Signif. codes:  0 '***' 0.001 '**' 0.01 '*' 0.05 '.' 0.1 ' ' 1
\end{verbatim}

In this case the main effects and the interaction are all very
significant, so we will use Tukey's HSD and the Bonferroni minimum
significant difference method to compare the 8 treatments (8(8-1)/2=28
comparisons). These will use \(MSABC\) as the ``error term'' with
\(df=(4-1)(2-1)(36-1)=105\).

\[
\mbox{Tukey:} \quad HSD = q_{.95,8,105}\sqrt{\frac{MSABC}{c}}
=4.374\sqrt{\frac{926.4971}{36}}=22.19 
\]

\[
\mbox{Bonferroni:} \quad MSD=t_{1-.05/(2(28));105}\sqrt{\frac{2MSABC}{c}}
=3.206\sqrt{\frac{2(926.4971)}{36}}=23.00
\]

\begin{Shaded}
\begin{Highlighting}[]
\NormalTok{ntm }\OtherTok{\textless{}{-}} \FunctionTok{read.csv}\NormalTok{(}\StringTok{"data/navtechmap1.csv"}\NormalTok{)}
\NormalTok{ntm}\SpecialCharTok{$}\NormalTok{inputMet.f }\OtherTok{\textless{}{-}} \FunctionTok{factor}\NormalTok{(ntm}\SpecialCharTok{$}\NormalTok{inputMeth, }\AttributeTok{levels=}\DecValTok{1}\SpecialCharTok{:}\DecValTok{2}\NormalTok{, }\AttributeTok{labels=}\FunctionTok{c}\NormalTok{(}\StringTok{"DT"}\NormalTok{,}\StringTok{"M"}\NormalTok{))}
\NormalTok{ntm}\SpecialCharTok{$}\NormalTok{navTech.f }\OtherTok{\textless{}{-}} \FunctionTok{factor}\NormalTok{(ntm}\SpecialCharTok{$}\NormalTok{navTech, }\AttributeTok{levels=}\DecValTok{1}\SpecialCharTok{:}\DecValTok{4}\NormalTok{, }
                        \AttributeTok{labels=}\FunctionTok{c}\NormalTok{(}\StringTok{"CPB"}\NormalTok{,}\StringTok{"DPB"}\NormalTok{,}\StringTok{"ENCC"}\NormalTok{,}\StringTok{"GD"}\NormalTok{))}
\NormalTok{ntm}\SpecialCharTok{$}\NormalTok{subject.f }\OtherTok{\textless{}{-}} \FunctionTok{factor}\NormalTok{(ntm}\SpecialCharTok{$}\NormalTok{subject)}

\FunctionTok{options}\NormalTok{(}\AttributeTok{contrasts=}\FunctionTok{c}\NormalTok{(}\StringTok{"contr.sum"}\NormalTok{,}\StringTok{"contr.poly"}\NormalTok{))}
\NormalTok{mod2 }\OtherTok{\textless{}{-}} \FunctionTok{lmer}\NormalTok{(taskTime }\SpecialCharTok{\textasciitilde{}}\NormalTok{ navTech.f }\SpecialCharTok{*}\NormalTok{ inputMet.f }\SpecialCharTok{+}\NormalTok{ (}\DecValTok{1}\SpecialCharTok{|}\NormalTok{subject.f) }\SpecialCharTok{+}
\NormalTok{               (}\DecValTok{1}\SpecialCharTok{|}\NormalTok{navTech.f}\SpecialCharTok{:}\NormalTok{subject.f) }\SpecialCharTok{+}\NormalTok{ (}\DecValTok{1}\SpecialCharTok{|}\NormalTok{inputMet.f}\SpecialCharTok{:}\NormalTok{subject.f),    }\AttributeTok{data=}\NormalTok{ntm)}
\end{Highlighting}
\end{Shaded}

\begin{verbatim}
## boundary (singular) fit: see help('isSingular')
\end{verbatim}

\begin{Shaded}
\begin{Highlighting}[]
\FunctionTok{summary}\NormalTok{(mod2)}
\end{Highlighting}
\end{Shaded}

\begin{verbatim}
## Linear mixed model fit by REML. t-tests use Satterthwaite's method [
## lmerModLmerTest]
## Formula: 
## taskTime ~ navTech.f * inputMet.f + (1 | subject.f) + (1 | navTech.f:subject.f) +  
##     (1 | inputMet.f:subject.f)
##    Data: ntm
## 
## REML criterion at convergence: 2851.3
## 
## Scaled residuals: 
##      Min       1Q   Median       3Q      Max 
## -2.14082 -0.60108  0.01166  0.59746  2.75578 
## 
## Random effects:
##  Groups               Name        Variance Std.Dev.
##  navTech.f:subject.f  (Intercept) 243.0    15.59   
##  inputMet.f:subject.f (Intercept) 253.6    15.92   
##  subject.f            (Intercept)   0.0     0.00   
##  Residual                         927.9    30.46   
## Number of obs: 288, groups:  
## navTech.f:subject.f, 144; inputMet.f:subject.f, 72; subject.f, 36
## 
## Fixed effects:
##                        Estimate Std. Error      df t value Pr(>|t|)    
## (Intercept)             147.964      2.904  78.527  50.957  < 2e-16 ***
## navTech.f1               15.511      3.838 112.449   4.042 9.75e-05 ***
## navTech.f2               14.616      3.838 112.449   3.809 0.000228 ***
## navTech.f3              -18.413      3.838 112.449  -4.798 4.96e-06 ***
## inputMet.f1              10.314      2.597  54.542   3.972 0.000211 ***
## navTech.f1:inputMet.f1    9.371      3.109 107.587   3.014 0.003214 ** 
## navTech.f2:inputMet.f1    5.476      3.109 107.587   1.762 0.080985 .  
## navTech.f3:inputMet.f1  -11.414      3.109 107.587  -3.671 0.000378 ***
## ---
## Signif. codes:  0 '***' 0.001 '**' 0.01 '*' 0.05 '.' 0.1 ' ' 1
## 
## Correlation of Fixed Effects:
##             (Intr) nvTc.1 nvTc.2 nvTc.3 inpM.1 nT.1:M nT.2:M
## navTech.f1   0.000                                          
## navTech.f2   0.000 -0.333                                   
## navTech.f3   0.000 -0.333 -0.333                            
## inputMet.f1  0.000  0.000  0.000  0.000                     
## nvTch.1:M.1  0.000  0.000  0.000  0.000  0.000              
## nvTch.2:M.1  0.000  0.000  0.000  0.000  0.000 -0.333       
## nvTch.3:M.1  0.000  0.000  0.000  0.000  0.000 -0.333 -0.333
## optimizer (nloptwrap) convergence code: 0 (OK)
## boundary (singular) fit: see help('isSingular')
\end{verbatim}

\begin{Shaded}
\begin{Highlighting}[]
\FunctionTok{anova}\NormalTok{(mod2)}
\end{Highlighting}
\end{Shaded}

\begin{verbatim}
## Type III Analysis of Variance Table with Satterthwaite's method
##                      Sum Sq Mean Sq NumDF   DenDF F value    Pr(>F)    
## navTech.f             43964 14654.6     3 112.449 15.7940 1.228e-08 ***
## inputMet.f            14636 14635.9     1  54.542 15.7738 0.0002109 ***
## navTech.f:inputMet.f  18710  6236.8     3 107.587  6.7216 0.0003360 ***
## ---
## Signif. codes:  0 '***' 0.001 '**' 0.01 '*' 0.05 '.' 0.1 ' ' 1
\end{verbatim}

\[ \nabla \]

\begin{verbatim}
## 
## Attaching package: 'nlme'
\end{verbatim}

\begin{verbatim}
## The following object is masked from 'package:lme4':
## 
##     lmList
\end{verbatim}

\begin{verbatim}
## The following object is masked from 'package:dplyr':
## 
##     collapse
\end{verbatim}

\chapter{Nested Designs}\label{nested-designs}

So far, factorial designs have contained \textbf{crossed} factors. The
levels of factor B were the same within the levels of factor A, and vice
versa. In this chapter the levels of factor B will be different under
the various levels of factor A, and are said to be \textbf{nested}.
These models are often referred to as hierarchical or multilevel.

In a nested model, with factor A at the ``top,'' the levels of factor B
will differ within the levels of factor A. Consider a study to compare
NCAA conferences, with the conference as factor A, and schools are
factor B, with different schools being in the different conferences
(although they are changing constantly).

\section{Estimators and the Analysis of
Variance}\label{estimators-and-the-analysis-of-variance}

The factors can be fixed (all levels included), random (levels are
sampled), or mixed (one factor fixed, the other random). We will write
out the general model, and whether effects are fixed or random should be
clear from the setting of the experiment. In the two-factor balanced
design, factor A will have \(a\) levels, within each level of A, factor
B will have \(b\) levels, and there will be \(n\) replicates within each
of the \(ab\) treatments. The model is given here.

\[
\mbox{Factor A:}\quad \mu_{i\bullet} = \mu_{\bullet\bullet}+\alpha_i \qquad
i=1,\ldots,a
\] \[
\mbox{Factor B within A:} \quad
\mu_{ij}=\mu_{\bullet\bullet}+\alpha_i+\beta_{j(i)} \qquad
i=1,\ldots,a; 
\quad j=1,\ldots,b
\] \[
\quad \Rightarrow \quad \beta_{j(i)} =
\mu_{ij}-\left(\mu_{\bullet\bullet}+\alpha_i\right)=\mu_{ij}-\mu_{i\bullet}
\]

When A and B are both fixed factors, all levels of interest are included
and \(\left\{\alpha_i\right\}\) and \(\left\{\beta_{j(i)}\right\}\) are
all fixed (unknown) parameters and we can use the following restrictions
and model.

\[
\sum_{i=1}^a \alpha_i=0 \qquad \sum_{j=1}^b\beta_{j(i)}=0 \quad i=1,\ldots,a
\] \[
Y_{ijk}=\mu_{\bullet\bullet}+\alpha_i+\beta_{j(i)} +\epsilon_{ijk}
\qquad i=1,\ldots,a;\quad j=1,\ldots,b; \quad k=1,\ldots,n
\] \[
\epsilon_{ijk} \sim NID\left(0,\sigma^2\right) \qquad
E\left\{Y_{ijk}\right\}=\mu_{\bullet\bullet}+\alpha_i+\beta_{j(i)} \qquad
V\left\{Y_{ijk}\right\}=\sigma^2
\]

When factor A is fixed and factor B is random, the model changes as
follows.

\[
\sum_{i=1}^a \alpha_i=0 \qquad \qquad 
\beta_{j(i)} \sim NID\left(0,\sigma_{\beta(\alpha)}^2\right)
\qquad \left\{\beta_{j(i)}\right\} \bot
\left\{\epsilon_{ijk}\right\}
\] \[
E\left\{Y_{ijk}\right\} = \mu_{\bullet\bullet}+\alpha_i  \qquad
V\left\{Y_{ijk}\right\} = \sigma_{\beta(\alpha)}^2 + \sigma^2 
\] \[
\mbox{COV}\left\{Y_{ijk},Y_{i'j'k'}\right\} = 
\left\{ \begin{array}{r@{\quad:\quad}l}
    \sigma_{\beta(\alpha)}^2 +  \sigma^2 & i=i',j=j',k=k' \\
    \sigma_{\beta(\alpha)}^2  & i=i', j=j',k\ne k' \\
    0 & \mbox{otherwise}\\
    \end{array} \right.
\]

When factors A and B are both random, the model changes as follows.

\[
\alpha_i\sim NID\left(0,\sigma_{\alpha}^2\right) \qquad
\beta_{j(i)} \sim NID\left(0,\sigma_{\beta(\alpha)}^2\right)
\qquad \left\{\alpha_i\right\} \bot\left\{\beta_{j(i)}\right\} \bot
\left\{\epsilon_{ijk}\right\}
\] \[
E\left\{Y_{ijk}\right\} = \mu_{\bullet\bullet}  \qquad
V\left\{Y_{ijk}\right\} = \sigma_{\alpha}^2+\sigma_{\beta(\alpha)}^2 + \sigma^2 
\] \[
\mbox{COV}\left\{Y_{ijk},Y_{i'j'k'}\right\} = 
\left\{ \begin{array}{r@{\quad:\quad}l}
    \sigma_{\alpha}^2 + \sigma_{\beta(\alpha)}^2 +  \sigma^2 & i=i',j=j',k=k' \\
    \sigma_{\alpha}^2 +\sigma_{\beta(\alpha)}^2  & i=i', j=j',k\ne k' \\
    \sigma_{\alpha}^2  & i=i',j\ne j', \forall k,k'\\
    0 & i\neq i', \forall j,j', \forall k,k'\\
    \end{array} \right.
\]

Estimators, fitted values, residuals, sums of squares and expected mean
squares are given below.

\[
\hat{\mu}_{\bullet\bullet}=\overline{Y}_{\bullet\bullet\bullet}
\qquad \hat{\alpha}_i=
\overline{Y}_{i\bullet\bullet}-\overline{Y}_{\bullet\bullet\bullet}
\qquad \hat{\beta}_{j(i)}
=\overline{Y}_{ij\bullet}-\overline{Y}_{i\bullet\bullet}
\] \[
\hat{Y}_{ijk}=\hat{\mu}_{\bullet\bullet}+\hat{\alpha}_i+\hat{\beta}_{j(i)}
=\overline{Y}_{ij\bullet} \qquad
e_{ijk}=Y_{ijk}-\hat{Y}_{ijk}=Y_{ijk}-\overline{Y}_{ij\bullet}
\]

\[
\mbox{Error:} \quad SSE = \sum_{i=1}^a\sum_{j=1}^b\sum_{k=1}^n
\left(Y_{ijk}-\overline{Y}_{ij\bullet}\right)^2
\qquad  df_E=ab(n-1)
\] \[
\mbox{Factor A:} \quad SSA =bn\sum_{i=1}^a
\left(Y_{i\bullet\bullet}-Y_{\bullet\bullet\bullet}\right)^2
\qquad df_A =a-1
\] \[
\mbox{Factor B within A:} \quad SSB(A) =n\sum_{i=1}^a\sum_{j=1}^b
\left(Y_{ij\bullet}-Y_{i\bullet\bullet}\right)^2
\qquad df_{B(A)} =a(b-1)
\]

The expected mean squares depend on whether the factors are fixed or
random. In all cases, \(E\{MSE\}=\sigma^2\). When both A and B are
fixed, we obtain the following results.

\[
E\{MSA\}=\sigma^2 + \frac{bn\sum_{i=1}^a\alpha_i^2}{a-1}
\qquad
E\{MSB(A)\} = \sigma^2 + 
\frac{n\sum_{i=1}^a\sum_{j=1}^b \beta_{j(i)}^2}{a(b-1)}
\]

The denominator for the \(F\)-tests for factors A and B in the fixed
case is \(MSE\).

When factor A is fixed and factor B is random, the following expected
mean squares are obtained.

\[
E\{MSA\}=\sigma^2 + n\sigma_{\beta(\alpha)}^2 + \frac{bn\sum_{i=1}^a\alpha_i^2}{a-1}
\qquad
E\{MSB(A)\} = \sigma^2 + 
n\sigma_{\beta(\alpha)}^2
\]

In this mixed case, the error term for factor B is \(MSE\), but the
error term for factor A is \(MSB(A)\) due to the covariance structure.

When factors A and B are both random, we obtain the following expected
mean squares.

\[
E\{MSA\}=\sigma^2 + n\sigma_{\beta(\alpha)}^2 + bn\sigma_{\alpha}^2
\qquad
E\{MSB(A)\} = \sigma^2 + 
n\sigma_{\beta(\alpha)}^2
\]

The error terms are the same in the random effects case are the same as
in the mixed effects case. Variance components can be estimated as in
the previous chapter.

\section{Tests, Contrasts, and Pairwise
Comparisons}\label{tests-contrasts-and-pairwise-comparisons}

When both factors are fixed, \(MSE\) is an unbiased estimator of
\(V\left\{Y_{ijk}\right\}\) and we can estimate/test contrasts among
levels of factor A and B within A.

First, we can test for effects among the levels of factor A.

\[
H_0^A: \alpha_1=\cdots = \alpha_a=0 \qquad
TS: F_A^*=\frac{MSA}{MSE}   \qquad
RR: F_A^* \geq F_{1-\alpha;a-1,ab(n-1)}
\] The estimators of the population means for levels of factor A and
their properties are given below.

\[
E\left\{\overline{Y}_{i\bullet\bullet}\right\}=\mu_i=
\mu_{\bullet\bullet}+\alpha_i   \qquad
V\left\{\overline{Y}_{i\bullet\bullet}\right\}=\frac{\sigma^2}{bn}
\qquad
\hat{V}\left\{\overline{Y}_{i\bullet\bullet}\right\}=\frac{MSE}{bn}
\] Any contrast among the levels of factor A can be written as follows,
with \(\sum_{i=1}^ac_i=0\).

\[
L_A = \sum_{i=1}^ac_i\mu_i \qquad
\hat{L}_A=\sum_{i=1}^ac_i \overline{Y}_{i\bullet\bullet}
\qquad
\hat{V}\left\{\hat{L}_A\right\}=\frac{MSE}{bn}\sum_{i=1}^ac_i^2
\]

\[
(1-\alpha)100\% \mbox{ Confidence Interval for } L_A:
\quad 
\hat{L}_A \pm t_{1-\alpha/2; ab(n-1)}\sqrt{\frac{MSE}{bn}\sum_{i=1}^ac_i^2}
\]

\[
\mbox{Tukey's:} \quad HSD=q_{1-\alpha;a,ab(n-1)}\sqrt{\frac{MSE}{bn}}
\qquad
\mbox{Simultaneous CI's:} \quad
\left(\overline{Y}_{i\bullet\bullet}-\overline{Y}_{i'\bullet\bullet}\right)
\pm HSD
\]

Next, consider inference for factor B when both factors are fixed.

\[
H_0^B(A): \beta_{1(1)}=\cdots = \beta_{b(a)}=0 \qquad
TS: F_{B(A)}^*=\frac{MSB(A)}{MSE}   
\] \[
RR: F_{B(A)}^* \geq F_{1-\alpha;a(b-1),ab(n-1)}
\]

\[
E\left\{\overline{Y}_{ij\bullet}\right\}=
\mu_{\bullet\bullet}+\alpha_i + \beta_{j(i)}  \qquad
V\left\{\overline{Y}_{ij\bullet}\right\}=\frac{\sigma^2}{n}
\qquad
\hat{V}\left\{\overline{Y}_{ij\bullet}\right\}=\frac{MSE}{n}
\]

Tukey's method can be applied to compare the \(b\) means within each
level of factor A as well.

\[
 HSD_B=q_{1-\alpha;b,ab(n-1)}\sqrt{\frac{MSE}{n}}
\qquad
\mbox{Simultaneous CI's:} \quad
\left(\overline{Y}_{ij\bullet}-\overline{Y}_{ij'\bullet}\right)
\pm HSD_B
\]

\textbf{Example 11.1 - Determination of Alcohol Content in Liquor}

A study was conducted to measure alcohol content in 3 types of spirits
(\citeproc{ref-oliveira_2017}{Oliveira et al. 2017}). The researchers
bought 3 brands each of Vodka, Whisky, and Cachaca. Factor A with
\(a=3\) is the spirit type (Vodka, Whisky, Cachaca). The brands are
nested within the spirit types with \(b=3\). Each of the brands
\((ab=3(3)=9)\) was measured \(n=24\) times. The response is \(Y\), the
difference between the measured alcohol content and the content provided
on the bottle's label. We will treat both factors as fixed in this
analysis. Note that the \textbf{TukeyHSD} function does not work well
with the nested (brand) factor.

\begin{Shaded}
\begin{Highlighting}[]
\NormalTok{sac }\OtherTok{\textless{}{-}} \FunctionTok{read.csv}\NormalTok{(}\StringTok{"data/spirits\_alccont.csv"}\NormalTok{)}
\FunctionTok{head}\NormalTok{(sac); }\FunctionTok{tail}\NormalTok{(sac)}
\end{Highlighting}
\end{Shaded}

\begin{verbatim}
##   spiritType brandSprt labelAC alcCntnt
## 1          1         1    37.5    36.81
## 2          1         1    37.5    36.43
## 3          1         1    37.5    37.02
## 4          1         1    37.5    37.38
## 5          1         1    37.5    36.70
## 6          1         1    37.5    36.97
\end{verbatim}

\begin{verbatim}
##     spiritType brandSprt labelAC alcCntnt
## 211          3         9      39    38.89
## 212          3         9      39    38.78
## 213          3         9      39    38.41
## 214          3         9      39    39.27
## 215          3         9      39    39.03
## 216          3         9      39    38.94
\end{verbatim}

\begin{Shaded}
\begin{Highlighting}[]
\NormalTok{sac}\SpecialCharTok{$}\NormalTok{spiritType.f }\OtherTok{\textless{}{-}} \FunctionTok{factor}\NormalTok{(sac}\SpecialCharTok{$}\NormalTok{spiritType)}
\NormalTok{sac}\SpecialCharTok{$}\NormalTok{brandSprt.f }\OtherTok{\textless{}{-}} \FunctionTok{factor}\NormalTok{(sac}\SpecialCharTok{$}\NormalTok{brandSprt)}
\NormalTok{sac}\SpecialCharTok{$}\NormalTok{Y }\OtherTok{\textless{}{-}}\NormalTok{ sac}\SpecialCharTok{$}\NormalTok{alcCntnt }\SpecialCharTok{{-}}\NormalTok{ sac}\SpecialCharTok{$}\NormalTok{labelAC}

\FunctionTok{plot}\NormalTok{(sac}\SpecialCharTok{$}\NormalTok{Y }\SpecialCharTok{\textasciitilde{}}\NormalTok{ sac}\SpecialCharTok{$}\NormalTok{spiritType, }\AttributeTok{pch=}\DecValTok{16}\NormalTok{, }\AttributeTok{col=}\NormalTok{sac}\SpecialCharTok{$}\NormalTok{brandSprt, }\AttributeTok{cex=}\FloatTok{0.6}\NormalTok{)}
\end{Highlighting}
\end{Shaded}

\pandocbounded{\includegraphics[keepaspectratio]{bookdownproj_STA4211Notes1_files/figure-latex/ch1101-1.pdf}}

\begin{Shaded}
\begin{Highlighting}[]
\NormalTok{A\_mean }\OtherTok{\textless{}{-}} \FunctionTok{as.vector}\NormalTok{(}\FunctionTok{tapply}\NormalTok{(sac}\SpecialCharTok{$}\NormalTok{Y, sac}\SpecialCharTok{$}\NormalTok{spiritType.f, mean))}
\NormalTok{BA\_mean }\OtherTok{\textless{}{-}} \FunctionTok{as.vector}\NormalTok{(}\FunctionTok{tapply}\NormalTok{(sac}\SpecialCharTok{$}\NormalTok{Y, sac}\SpecialCharTok{$}\NormalTok{brandSprt.f, mean))}
\NormalTok{all\_mean }\OtherTok{\textless{}{-}} \FunctionTok{mean}\NormalTok{(sac}\SpecialCharTok{$}\NormalTok{Y)}
\NormalTok{a }\OtherTok{\textless{}{-}} \FunctionTok{length}\NormalTok{(A\_mean)}
\NormalTok{b }\OtherTok{\textless{}{-}} \FunctionTok{length}\NormalTok{(BA\_mean) }\SpecialCharTok{/}\NormalTok{ a}
\NormalTok{n }\OtherTok{\textless{}{-}} \FunctionTok{length}\NormalTok{(sac}\SpecialCharTok{$}\NormalTok{Y) }\SpecialCharTok{/}\NormalTok{ (a}\SpecialCharTok{*}\NormalTok{b)}

\NormalTok{Ybar\_i.. }\OtherTok{\textless{}{-}} \FunctionTok{rep}\NormalTok{(A\_mean, }\AttributeTok{each=}\NormalTok{b}\SpecialCharTok{*}\NormalTok{n)}
\NormalTok{Ybar\_ij. }\OtherTok{\textless{}{-}} \FunctionTok{rep}\NormalTok{(BA\_mean, }\AttributeTok{each=}\NormalTok{n)}
\NormalTok{Ybar\_... }\OtherTok{\textless{}{-}} \FunctionTok{rep}\NormalTok{(all\_mean, a}\SpecialCharTok{*}\NormalTok{b}\SpecialCharTok{*}\NormalTok{n)}

\NormalTok{SSA }\OtherTok{\textless{}{-}} \FunctionTok{sum}\NormalTok{((Ybar\_i..}\SpecialCharTok{{-}}\NormalTok{Ybar\_...)}\SpecialCharTok{\^{}}\DecValTok{2}\NormalTok{); df\_A }\OtherTok{\textless{}{-}}\NormalTok{ a}\DecValTok{{-}1}\NormalTok{; MSA }\OtherTok{\textless{}{-}}\NormalTok{ SSA}\SpecialCharTok{/}\NormalTok{df\_A}
\NormalTok{SSBA }\OtherTok{\textless{}{-}} \FunctionTok{sum}\NormalTok{((Ybar\_ij.}\SpecialCharTok{{-}}\NormalTok{Ybar\_i..)}\SpecialCharTok{\^{}}\DecValTok{2}\NormalTok{); df\_BA }\OtherTok{\textless{}{-}}\NormalTok{ a}\SpecialCharTok{*}\NormalTok{(b}\DecValTok{{-}1}\NormalTok{); MSBA }\OtherTok{\textless{}{-}}\NormalTok{ SSBA}\SpecialCharTok{/}\NormalTok{df\_BA}
\NormalTok{SSE }\OtherTok{\textless{}{-}} \FunctionTok{sum}\NormalTok{((sac}\SpecialCharTok{$}\NormalTok{Y}\SpecialCharTok{{-}}\NormalTok{Ybar\_ij.)}\SpecialCharTok{\^{}}\DecValTok{2}\NormalTok{); df\_E }\OtherTok{\textless{}{-}}\NormalTok{ a}\SpecialCharTok{*}\NormalTok{b}\SpecialCharTok{*}\NormalTok{(n}\DecValTok{{-}1}\NormalTok{); MSE }\OtherTok{\textless{}{-}}\NormalTok{ SSE}\SpecialCharTok{/}\NormalTok{df\_E}
\NormalTok{SSTO }\OtherTok{\textless{}{-}} \FunctionTok{sum}\NormalTok{((sac}\SpecialCharTok{$}\NormalTok{Y}\SpecialCharTok{{-}}\NormalTok{Ybar\_...)}\SpecialCharTok{\^{}}\DecValTok{2}\NormalTok{); df\_TO }\OtherTok{\textless{}{-}}\NormalTok{ a}\SpecialCharTok{*}\NormalTok{b}\SpecialCharTok{*}\NormalTok{n}\DecValTok{{-}1}

\NormalTok{aov.out1 }\OtherTok{\textless{}{-}} \FunctionTok{cbind}\NormalTok{(df\_A,SSA,MSA,MSA}\SpecialCharTok{/}\NormalTok{MSE,}\FunctionTok{qf}\NormalTok{(.}\DecValTok{95}\NormalTok{,df\_A,df\_E),}\DecValTok{1}\SpecialCharTok{{-}}\FunctionTok{pf}\NormalTok{(MSA}\SpecialCharTok{/}\NormalTok{MSE,df\_A,df\_E),SSA}\SpecialCharTok{/}\NormalTok{SSTO, SSA}\SpecialCharTok{/}\NormalTok{(SSA}\SpecialCharTok{+}\NormalTok{SSE))}
\NormalTok{aov.out2 }\OtherTok{\textless{}{-}} \FunctionTok{cbind}\NormalTok{(df\_BA,SSBA,MSBA,MSBA}\SpecialCharTok{/}\NormalTok{MSE,}\FunctionTok{qf}\NormalTok{(.}\DecValTok{95}\NormalTok{,df\_BA,df\_E),}
                  \DecValTok{1}\SpecialCharTok{{-}}\FunctionTok{pf}\NormalTok{(MSBA}\SpecialCharTok{/}\NormalTok{MSE,df\_BA,df\_E),SSBA}\SpecialCharTok{/}\NormalTok{SSTO,SSBA}\SpecialCharTok{/}\NormalTok{(SSBA}\SpecialCharTok{+}\NormalTok{SSE))}
\NormalTok{aov.out3 }\OtherTok{\textless{}{-}} \FunctionTok{cbind}\NormalTok{(df\_E, SSE, MSE,}\ConstantTok{NA}\NormalTok{,}\ConstantTok{NA}\NormalTok{,}\ConstantTok{NA}\NormalTok{,}\ConstantTok{NA}\NormalTok{,}\ConstantTok{NA}\NormalTok{)}
\NormalTok{aov.out4 }\OtherTok{\textless{}{-}} \FunctionTok{cbind}\NormalTok{(df\_TO,SSTO,}\ConstantTok{NA}\NormalTok{,}\ConstantTok{NA}\NormalTok{,}\ConstantTok{NA}\NormalTok{,}\ConstantTok{NA}\NormalTok{,}\ConstantTok{NA}\NormalTok{,}\ConstantTok{NA}\NormalTok{)}
\NormalTok{aov.out }\OtherTok{\textless{}{-}} \FunctionTok{rbind}\NormalTok{(aov.out1,aov.out2,aov.out3,aov.out4)}
\FunctionTok{colnames}\NormalTok{(aov.out) }\OtherTok{\textless{}{-}} \FunctionTok{c}\NormalTok{(}\StringTok{"df"}\NormalTok{,}\StringTok{"SS"}\NormalTok{,}\StringTok{"MS"}\NormalTok{,}\StringTok{"F*"}\NormalTok{,}\StringTok{"F(.95)"}\NormalTok{,}\StringTok{"P(\textgreater{}F*)"}\NormalTok{,}\StringTok{"eta\^{}2"}\NormalTok{, }\StringTok{"Part eta\^{}2"}\NormalTok{)}
\FunctionTok{rownames}\NormalTok{(aov.out) }\OtherTok{\textless{}{-}} \FunctionTok{c}\NormalTok{(}\StringTok{"Type"}\NormalTok{,}\StringTok{"Brand(Type)"}\NormalTok{,}\StringTok{"Error"}\NormalTok{,}\StringTok{"Total"}\NormalTok{)}
\FunctionTok{round}\NormalTok{(aov.out,}\DecValTok{4}\NormalTok{)}
\end{Highlighting}
\end{Shaded}

\begin{verbatim}
##              df      SS     MS      F* F(.95) P(>F*)  eta^2 Part eta^2
## Type          2  3.1967 1.5983 13.2669 3.0395      0 0.0992     0.1136
## Brand(Type)   6  4.0872 0.6812  5.6543 2.1426      0 0.1268     0.1408
## Error       207 24.9385 0.1205      NA     NA     NA     NA         NA
## Total       215 32.2224     NA      NA     NA     NA     NA         NA
\end{verbatim}

\begin{Shaded}
\begin{Highlighting}[]
\NormalTok{HSD\_A }\OtherTok{\textless{}{-}} \FunctionTok{qtukey}\NormalTok{(.}\DecValTok{95}\NormalTok{,a,df\_E)}\SpecialCharTok{*}\FunctionTok{sqrt}\NormalTok{(MSE}\SpecialCharTok{/}\NormalTok{(b}\SpecialCharTok{*}\NormalTok{n))}
\NormalTok{HSD\_B }\OtherTok{\textless{}{-}} \FunctionTok{qtukey}\NormalTok{(.}\DecValTok{95}\NormalTok{,b,df\_E)}\SpecialCharTok{*}\FunctionTok{sqrt}\NormalTok{(MSE}\SpecialCharTok{/}\NormalTok{n)}

\NormalTok{HSD\_A.out }\OtherTok{\textless{}{-}} \FunctionTok{matrix}\NormalTok{(}\FunctionTok{rep}\NormalTok{(}\DecValTok{0}\NormalTok{,}\DecValTok{8}\SpecialCharTok{*}\NormalTok{a}\SpecialCharTok{*}\NormalTok{(a}\DecValTok{{-}1}\NormalTok{)}\SpecialCharTok{/}\DecValTok{2}\NormalTok{), }\AttributeTok{ncol=}\DecValTok{8}\NormalTok{)}
\NormalTok{HSD\_B.out }\OtherTok{\textless{}{-}} \FunctionTok{matrix}\NormalTok{(}\FunctionTok{rep}\NormalTok{(}\DecValTok{0}\NormalTok{,}\DecValTok{9}\SpecialCharTok{*}\NormalTok{a}\SpecialCharTok{*}\NormalTok{b}\SpecialCharTok{*}\NormalTok{(b}\DecValTok{{-}1}\NormalTok{)}\SpecialCharTok{/}\DecValTok{2}\NormalTok{), }\AttributeTok{ncol=}\DecValTok{9}\NormalTok{)}
\NormalTok{row\_A.count }\OtherTok{\textless{}{-}} \DecValTok{0}
\NormalTok{row\_B.count }\OtherTok{\textless{}{-}} \DecValTok{0}

\ControlFlowTok{for}\NormalTok{ (i1 }\ControlFlowTok{in} \DecValTok{2}\SpecialCharTok{:}\NormalTok{a) \{}
  \ControlFlowTok{for}\NormalTok{ (i2 }\ControlFlowTok{in} \DecValTok{1}\SpecialCharTok{:}\NormalTok{(i1}\DecValTok{{-}1}\NormalTok{)) \{}
\NormalTok{    row\_A.count }\OtherTok{\textless{}{-}}\NormalTok{ row\_A.count }\SpecialCharTok{+} \DecValTok{1}
\NormalTok{    HSD\_A.out[row\_A.count,}\DecValTok{1}\NormalTok{] }\OtherTok{\textless{}{-}}\NormalTok{ i1}
\NormalTok{    HSD\_A.out[row\_A.count,}\DecValTok{2}\NormalTok{] }\OtherTok{\textless{}{-}}\NormalTok{ i2}
\NormalTok{    HSD\_A.out[row\_A.count,}\DecValTok{3}\NormalTok{] }\OtherTok{\textless{}{-}}\NormalTok{ A\_mean[i1]}
\NormalTok{    HSD\_A.out[row\_A.count,}\DecValTok{4}\NormalTok{] }\OtherTok{\textless{}{-}}\NormalTok{ A\_mean[i2]}
\NormalTok{    HSD\_A.out[row\_A.count,}\DecValTok{5}\NormalTok{] }\OtherTok{\textless{}{-}}\NormalTok{ A\_mean[i1]}\SpecialCharTok{{-}}\NormalTok{A\_mean[i2]}
\NormalTok{    HSD\_A.out[row\_A.count,}\DecValTok{6}\NormalTok{] }\OtherTok{\textless{}{-}}\NormalTok{ (A\_mean[i1]}\SpecialCharTok{{-}}\NormalTok{A\_mean[i2])}\SpecialCharTok{{-}}\NormalTok{HSD\_A}
\NormalTok{    HSD\_A.out[row\_A.count,}\DecValTok{7}\NormalTok{] }\OtherTok{\textless{}{-}}\NormalTok{ (A\_mean[i1]}\SpecialCharTok{{-}}\NormalTok{A\_mean[i2])}\SpecialCharTok{+}\NormalTok{HSD\_A}
\NormalTok{    HSD\_A.out[row\_A.count,}\DecValTok{8}\NormalTok{] }\OtherTok{\textless{}{-}} 
      \DecValTok{1}\SpecialCharTok{{-}}\FunctionTok{ptukey}\NormalTok{(}\FunctionTok{abs}\NormalTok{((A\_mean[i1]}\SpecialCharTok{{-}}\NormalTok{A\_mean[i2])}\SpecialCharTok{/}\FunctionTok{sqrt}\NormalTok{(MSE}\SpecialCharTok{/}\NormalTok{(b}\SpecialCharTok{*}\NormalTok{n))),a,df\_E)}
\NormalTok{  \}}
\NormalTok{\}}
\FunctionTok{colnames}\NormalTok{(HSD\_A.out) }\OtherTok{\textless{}{-}} \FunctionTok{c}\NormalTok{(}\StringTok{"i"}\NormalTok{,}\StringTok{"i\textquotesingle{}"}\NormalTok{,}\StringTok{"Ybar\_i"}\NormalTok{,}\StringTok{"Ybar\_i\textquotesingle{}"}\NormalTok{,}\StringTok{"Diff"}\NormalTok{,}\StringTok{"LB"}\NormalTok{,}\StringTok{"UB"}\NormalTok{,}\StringTok{"p adj"}\NormalTok{)}
\FunctionTok{round}\NormalTok{(HSD\_A.out,}\DecValTok{4}\NormalTok{)}
\end{Highlighting}
\end{Shaded}

\begin{verbatim}
##      i i'  Ybar_i Ybar_i'    Diff      LB     UB p adj
## [1,] 2  1 -0.3062 -0.3037 -0.0025 -0.1391 0.1341 0.999
## [2,] 3  1 -0.0469 -0.3037  0.2568  0.1202 0.3934 0.000
## [3,] 3  2 -0.0469 -0.3062  0.2593  0.1227 0.3959 0.000
\end{verbatim}

\begin{Shaded}
\begin{Highlighting}[]
\ControlFlowTok{for}\NormalTok{ (i1 }\ControlFlowTok{in} \DecValTok{1}\SpecialCharTok{:}\NormalTok{a) \{}
  \ControlFlowTok{for}\NormalTok{ (i2 }\ControlFlowTok{in}\NormalTok{ ((i1}\DecValTok{{-}1}\NormalTok{)}\SpecialCharTok{*}\NormalTok{a}\SpecialCharTok{+}\DecValTok{2}\NormalTok{)}\SpecialCharTok{:}\NormalTok{((i1}\DecValTok{{-}1}\NormalTok{)}\SpecialCharTok{*}\NormalTok{a}\SpecialCharTok{+}\NormalTok{b)) \{}
    \ControlFlowTok{for}\NormalTok{ (i3 }\ControlFlowTok{in}\NormalTok{ ((i1}\DecValTok{{-}1}\NormalTok{)}\SpecialCharTok{*}\NormalTok{a}\SpecialCharTok{+}\DecValTok{1}\NormalTok{)}\SpecialCharTok{:}\NormalTok{(i2}\DecValTok{{-}1}\NormalTok{)) \{}
\NormalTok{    row\_B.count }\OtherTok{\textless{}{-}}\NormalTok{ row\_B.count }\SpecialCharTok{+} \DecValTok{1}
\NormalTok{    HSD\_B.out[row\_B.count,}\DecValTok{1}\NormalTok{] }\OtherTok{\textless{}{-}}\NormalTok{ i1}
\NormalTok{    HSD\_B.out[row\_B.count,}\DecValTok{2}\NormalTok{] }\OtherTok{\textless{}{-}}\NormalTok{ i2}
\NormalTok{    HSD\_B.out[row\_B.count,}\DecValTok{3}\NormalTok{] }\OtherTok{\textless{}{-}}\NormalTok{ i3}
\NormalTok{    HSD\_B.out[row\_B.count,}\DecValTok{4}\NormalTok{] }\OtherTok{\textless{}{-}}\NormalTok{ BA\_mean[i2]}
\NormalTok{    HSD\_B.out[row\_B.count,}\DecValTok{5}\NormalTok{] }\OtherTok{\textless{}{-}}\NormalTok{ BA\_mean[i3]}
\NormalTok{    HSD\_B.out[row\_B.count,}\DecValTok{6}\NormalTok{] }\OtherTok{\textless{}{-}}\NormalTok{ BA\_mean[i2]}\SpecialCharTok{{-}}\NormalTok{BA\_mean[i3]}
\NormalTok{    HSD\_B.out[row\_B.count,}\DecValTok{7}\NormalTok{] }\OtherTok{\textless{}{-}}\NormalTok{ (BA\_mean[i2]}\SpecialCharTok{{-}}\NormalTok{BA\_mean[i3])}\SpecialCharTok{{-}}\NormalTok{HSD\_B}
\NormalTok{    HSD\_B.out[row\_B.count,}\DecValTok{8}\NormalTok{] }\OtherTok{\textless{}{-}}\NormalTok{ (BA\_mean[i2]}\SpecialCharTok{{-}}\NormalTok{BA\_mean[i3])}\SpecialCharTok{+}\NormalTok{HSD\_B}
\NormalTok{    HSD\_B.out[row\_B.count,}\DecValTok{9}\NormalTok{] }\OtherTok{\textless{}{-}} 
      \DecValTok{1}\SpecialCharTok{{-}}\FunctionTok{ptukey}\NormalTok{(}\FunctionTok{abs}\NormalTok{((BA\_mean[i2]}\SpecialCharTok{{-}}\NormalTok{BA\_mean[i3])}\SpecialCharTok{/}\FunctionTok{sqrt}\NormalTok{(MSE}\SpecialCharTok{/}\NormalTok{n)),b,df\_E)}
\NormalTok{    \}}
\NormalTok{  \}}
\NormalTok{\}}
\FunctionTok{colnames}\NormalTok{(HSD\_B.out) }\OtherTok{\textless{}{-}} \FunctionTok{c}\NormalTok{(}\StringTok{"i"}\NormalTok{,}\StringTok{"j"}\NormalTok{,}\StringTok{"j\textquotesingle{}"}\NormalTok{,}\StringTok{"Ybar\_j(i)"}\NormalTok{,}\StringTok{"Ybar\_j\textquotesingle{}(i)"}\NormalTok{,}\StringTok{"Diff"}\NormalTok{,}\StringTok{"LB"}\NormalTok{,}\StringTok{"UB"}\NormalTok{,}\StringTok{"p adj"}\NormalTok{)}
\FunctionTok{round}\NormalTok{(HSD\_B.out,}\DecValTok{4}\NormalTok{)}
\end{Highlighting}
\end{Shaded}

\begin{verbatim}
##       i j j' Ybar_j(i) Ybar_j'(i)    Diff      LB     UB  p adj
##  [1,] 1 2  1   -0.2400    -0.4308  0.1908 -0.0457 0.4274 0.1399
##  [2,] 1 3  1   -0.2404    -0.4308  0.1904 -0.0461 0.4270 0.1411
##  [3,] 1 3  2   -0.2404    -0.2400 -0.0004 -0.2370 0.2361 1.0000
##  [4,] 2 5  4   -0.1388    -0.5892  0.4504  0.2139 0.6870 0.0000
##  [5,] 2 6  4   -0.1908    -0.5892  0.3983  0.1618 0.6349 0.0003
##  [6,] 2 6  5   -0.1908    -0.1388 -0.0521 -0.2886 0.1845 0.8618
##  [7,] 3 8  7    0.0792    -0.0900  0.1692 -0.0674 0.4057 0.2121
##  [8,] 3 9  7   -0.1300    -0.0900 -0.0400 -0.2765 0.1965 0.9159
##  [9,] 3 9  8   -0.1300     0.0792 -0.2092 -0.4457 0.0274 0.0949
\end{verbatim}

\begin{Shaded}
\begin{Highlighting}[]
\NormalTok{mod1 }\OtherTok{\textless{}{-}} \FunctionTok{aov}\NormalTok{(Y }\SpecialCharTok{\textasciitilde{}}\NormalTok{ spiritType.f }\SpecialCharTok{+}\NormalTok{ spiritType.f}\SpecialCharTok{/}\NormalTok{brandSprt.f, }\AttributeTok{data=}\NormalTok{sac)}
\FunctionTok{anova}\NormalTok{(mod1)}
\end{Highlighting}
\end{Shaded}

\begin{verbatim}
## Analysis of Variance Table
## 
## Response: Y
##                           Df  Sum Sq Mean Sq F value    Pr(>F)    
## spiritType.f               2  3.1967 1.59834 13.2669 3.791e-06 ***
## spiritType.f:brandSprt.f   6  4.0872 0.68120  5.6543 1.850e-05 ***
## Residuals                207 24.9385 0.12048                      
## ---
## Signif. codes:  0 '***' 0.001 '**' 0.01 '*' 0.05 '.' 0.1 ' ' 1
\end{verbatim}

\begin{Shaded}
\begin{Highlighting}[]
\FunctionTok{TukeyHSD}\NormalTok{(mod1, }\StringTok{"spiritType.f"}\NormalTok{)}
\end{Highlighting}
\end{Shaded}

\begin{verbatim}
##   Tukey multiple comparisons of means
##     95% family-wise confidence level
## 
## Fit: aov(formula = Y ~ spiritType.f + spiritType.f/brandSprt.f, data = sac)
## 
## $spiritType.f
##           diff        lwr       upr     p adj
## 2-1 -0.0025000 -0.1390648 0.1340648 0.9989709
## 3-1  0.2568056  0.1202407 0.3933704 0.0000435
## 3-2  0.2593056  0.1227407 0.3958704 0.0000362
\end{verbatim}

\begin{Shaded}
\begin{Highlighting}[]
\FunctionTok{options}\NormalTok{(}\AttributeTok{es.use\_symbols=}\ConstantTok{TRUE}\NormalTok{)}

\FunctionTok{eta\_squared}\NormalTok{(mod1)}
\end{Highlighting}
\end{Shaded}

\begin{verbatim}
## # Effect Size for ANOVA (Type I)
## 
## Parameter                | η² (partial) |       95% CI
## ------------------------------------------------------
## spiritType.f             |         0.11 | [0.05, 1.00]
## spiritType.f:brandSprt.f |         0.14 | [0.06, 1.00]
## 
## - One-sided CIs: upper bound fixed at [1.00].
\end{verbatim}

\begin{Shaded}
\begin{Highlighting}[]
\FunctionTok{eta\_squared}\NormalTok{(mod1, }\AttributeTok{partial=}\ConstantTok{FALSE}\NormalTok{)}
\end{Highlighting}
\end{Shaded}

\begin{verbatim}
## # Effect Size for ANOVA (Type I)
## 
## Parameter                |   η² |       95% CI
## ----------------------------------------------
## spiritType.f             | 0.10 | [0.04, 1.00]
## spiritType.f:brandSprt.f | 0.13 | [0.05, 1.00]
## 
## - One-sided CIs: upper bound fixed at [1.00].
\end{verbatim}

There is evidence for differences among the \(a=3\) types of spirits.
The differences between measured and reported proofs are significantly
higher for Cachaca \((i=3)\) than for Whisky \((i=2)\) and for Vodka
\((i=1)\).

There are also differences among brands within spirits types. In
particular, brands 5 and 6 (within Whisky) are significantly higher than
brand 4. \[\nabla \] When factor A is fixed and B is random, we would
like to compare the levels of factor A, and estimate the variance for
factor B. The error term for factor A is now \(MSB(A)\) as opposed to
\(MSE\).

\[
H_0^A: \alpha_1=\cdots = \alpha_a=0 \qquad
TS: F_A^*=\frac{MSA}{MSB(A)}   \qquad
RR: F_A^* \geq F_{1-\alpha;a-1,a(b-1)}
\] The estimators of the population means for levels of factor A and
their properties are given below.

\[
E\left\{\overline{Y}_{i\bullet\bullet}\right\}=\mu_i=
\mu_{\bullet\bullet}+\alpha_i   \qquad
V\left\{\overline{Y}_{i\bullet\bullet}\right\}=
\frac{\sigma^2+\sigma_{\beta(\alpha)}^2}{bn}
\qquad
\hat{V}\left\{\overline{Y}_{i\bullet\bullet}\right\}=\frac{MSB(A)}{bn}
\]

For Tukey's method, relative to the fixed effects case, we replace
\(MSE\) with \(MSB(A)\) and adjust the degrees of freedom.

\[
HSD = q_{1-\alpha;a,a(b-1)}\sqrt{\frac{MSB(A)}{bn}}
\] \[
(1-\alpha)100\% \mbox{ CI for } \quad \alpha_i-\alpha_{i'}:
\qquad 
\left(\overline{Y}_{i\bullet\bullet}-\overline{Y}_{i'\bullet\bullet}
\right) \pm HSD
\] A point estimate for the variance among the levels of factor A and
its corresponding approximate Confidence Interval based on
Satterthwaite's approximation are given below. First, we test whether
the variance component for factor B is 0.

\[
H_0^{B(A)}:\sigma_{\beta(\alpha)}^2=0 \qquad
TS: F_{B(A)}^*=\frac{MSB(A)}{MSE}   \qquad
RR: F_{B(A)}^* \geq F_{1-\alpha;a(b-1),ab(n-1)}
\] \[
E\{MSB(A)\}=\sigma^2+n\sigma_{\beta(\alpha)}^2 \qquad E\{MSE\}=\sigma^2
\quad \Rightarrow \quad s_{\beta(\alpha)}^2=\frac{MSB(A)-MSE}{n}
\] \[
df_{\beta(\alpha)} = \frac{\left(s_{\beta(\alpha)}^2\right)^2}{
\left[\frac{\left(\frac{1}{n}MSB(A)\right)^2}{a(b-1)} +
\frac{\left(-\frac{1}{n}MSE\right)^2}{ab(n-1)} \right]
}
\] \[
\mbox{Approximate } (1-\alpha)100\% \mbox{ CI for }
\sigma_{\beta(\alpha)}^2: \quad
\left[
\frac{df_{\beta(\alpha)}s_{\beta(\alpha)}^2}{
\chi_{1-\alpha/2; df_{\beta(\alpha)}}^2},
\frac{df_{\beta(\alpha)}s_{\beta(\alpha)}^2}{
\chi_{\alpha/2; df_{\beta(\alpha)}}^2}
\right]
\]

\textbf{Example 11.2 - Momentum of Animal Traps}

A study compared \(a=8\) models of animal traps in terms of momentum
(\citeproc{ref-cook_1989}{Cook and Proulx 1989}). Within each model,
\(b=3\) traps were built. Each trap was measured \(n=10\) times. We will
treat the models as a fixed factor, and the traps within models as a
random factor (any number could have been assembled for each model). The
response measured was the momentum at HDISP (when both jaws displayed
half way). The response has been multiplied by 100 to avoid small
variance estimates.

\begin{Shaded}
\begin{Highlighting}[]
\NormalTok{at }\OtherTok{\textless{}{-}} \FunctionTok{read.csv}\NormalTok{(}\StringTok{"data/animal\_trap.csv"}\NormalTok{)}
\FunctionTok{head}\NormalTok{(at); }\FunctionTok{tail}\NormalTok{(at)}
\end{Highlighting}
\end{Shaded}

\begin{verbatim}
##   trapGrp model trapModel momentum
## 1       1     1        11  0.51423
## 2       1     1        11  0.53202
## 3       1     1        11  0.51550
## 4       1     1        11  0.46825
## 5       1     1        11  0.48496
## 6       1     1        11  0.51743
\end{verbatim}

\begin{verbatim}
##     trapGrp model trapModel momentum
## 235      24     8        83  0.92003
## 236      24     8        83  0.91039
## 237      24     8        83  0.95544
## 238      24     8        83  0.84232
## 239      24     8        83  0.92145
## 240      24     8        83  0.87547
\end{verbatim}

\begin{Shaded}
\begin{Highlighting}[]
\NormalTok{at}\SpecialCharTok{$}\NormalTok{model.f }\OtherTok{\textless{}{-}} \FunctionTok{factor}\NormalTok{(at}\SpecialCharTok{$}\NormalTok{model)}
\NormalTok{at}\SpecialCharTok{$}\NormalTok{trapModel.f }\OtherTok{\textless{}{-}} \FunctionTok{factor}\NormalTok{(at}\SpecialCharTok{$}\NormalTok{trapModel)}
\NormalTok{at}\SpecialCharTok{$}\NormalTok{Y }\OtherTok{\textless{}{-}} \DecValTok{100}\SpecialCharTok{*}\NormalTok{at}\SpecialCharTok{$}\NormalTok{momentum}

\FunctionTok{plot}\NormalTok{(at}\SpecialCharTok{$}\NormalTok{Y }\SpecialCharTok{\textasciitilde{}}\NormalTok{ at}\SpecialCharTok{$}\NormalTok{model, }\AttributeTok{pch=}\DecValTok{16}\NormalTok{, }\AttributeTok{col=}\NormalTok{at}\SpecialCharTok{$}\NormalTok{trapModel, }\AttributeTok{cex=}\FloatTok{0.6}\NormalTok{)}
\end{Highlighting}
\end{Shaded}

\pandocbounded{\includegraphics[keepaspectratio]{bookdownproj_STA4211Notes1_files/figure-latex/ch1102-1.pdf}}

\begin{Shaded}
\begin{Highlighting}[]
\NormalTok{A\_mean }\OtherTok{\textless{}{-}} \FunctionTok{as.vector}\NormalTok{(}\FunctionTok{tapply}\NormalTok{(at}\SpecialCharTok{$}\NormalTok{Y, at}\SpecialCharTok{$}\NormalTok{model.f, mean))}
\NormalTok{BA\_mean }\OtherTok{\textless{}{-}} \FunctionTok{as.vector}\NormalTok{(}\FunctionTok{tapply}\NormalTok{(at}\SpecialCharTok{$}\NormalTok{Y, at}\SpecialCharTok{$}\NormalTok{trapModel.f, mean))}
\NormalTok{all\_mean }\OtherTok{\textless{}{-}} \FunctionTok{mean}\NormalTok{(at}\SpecialCharTok{$}\NormalTok{Y)}
\NormalTok{a }\OtherTok{\textless{}{-}} \FunctionTok{length}\NormalTok{(A\_mean)}
\NormalTok{b }\OtherTok{\textless{}{-}} \FunctionTok{length}\NormalTok{(BA\_mean) }\SpecialCharTok{/}\NormalTok{ a}
\NormalTok{n }\OtherTok{\textless{}{-}} \FunctionTok{length}\NormalTok{(at}\SpecialCharTok{$}\NormalTok{Y) }\SpecialCharTok{/}\NormalTok{ (a}\SpecialCharTok{*}\NormalTok{b)}

\NormalTok{Ybar\_i.. }\OtherTok{\textless{}{-}} \FunctionTok{rep}\NormalTok{(A\_mean, }\AttributeTok{each=}\NormalTok{b}\SpecialCharTok{*}\NormalTok{n)}
\NormalTok{Ybar\_ij. }\OtherTok{\textless{}{-}} \FunctionTok{rep}\NormalTok{(BA\_mean, }\AttributeTok{each=}\NormalTok{n)}
\NormalTok{Ybar\_... }\OtherTok{\textless{}{-}} \FunctionTok{rep}\NormalTok{(all\_mean, a}\SpecialCharTok{*}\NormalTok{b}\SpecialCharTok{*}\NormalTok{n)}

\NormalTok{SSA }\OtherTok{\textless{}{-}} \FunctionTok{sum}\NormalTok{((Ybar\_i..}\SpecialCharTok{{-}}\NormalTok{Ybar\_...)}\SpecialCharTok{\^{}}\DecValTok{2}\NormalTok{); df\_A }\OtherTok{\textless{}{-}}\NormalTok{ a}\DecValTok{{-}1}\NormalTok{; MSA }\OtherTok{\textless{}{-}}\NormalTok{ SSA}\SpecialCharTok{/}\NormalTok{df\_A}
\NormalTok{SSBA }\OtherTok{\textless{}{-}} \FunctionTok{sum}\NormalTok{((Ybar\_ij.}\SpecialCharTok{{-}}\NormalTok{Ybar\_i..)}\SpecialCharTok{\^{}}\DecValTok{2}\NormalTok{); df\_BA }\OtherTok{\textless{}{-}}\NormalTok{ a}\SpecialCharTok{*}\NormalTok{(b}\DecValTok{{-}1}\NormalTok{); MSBA }\OtherTok{\textless{}{-}}\NormalTok{ SSBA}\SpecialCharTok{/}\NormalTok{df\_BA}
\NormalTok{SSE }\OtherTok{\textless{}{-}} \FunctionTok{sum}\NormalTok{((at}\SpecialCharTok{$}\NormalTok{Y}\SpecialCharTok{{-}}\NormalTok{Ybar\_ij.)}\SpecialCharTok{\^{}}\DecValTok{2}\NormalTok{); df\_E }\OtherTok{\textless{}{-}}\NormalTok{ a}\SpecialCharTok{*}\NormalTok{b}\SpecialCharTok{*}\NormalTok{(n}\DecValTok{{-}1}\NormalTok{); MSE }\OtherTok{\textless{}{-}}\NormalTok{ SSE}\SpecialCharTok{/}\NormalTok{df\_E}
\NormalTok{SSTO }\OtherTok{\textless{}{-}} \FunctionTok{sum}\NormalTok{((at}\SpecialCharTok{$}\NormalTok{Y}\SpecialCharTok{{-}}\NormalTok{Ybar\_...)}\SpecialCharTok{\^{}}\DecValTok{2}\NormalTok{); df\_TO }\OtherTok{\textless{}{-}}\NormalTok{ a}\SpecialCharTok{*}\NormalTok{b}\SpecialCharTok{*}\NormalTok{n}\DecValTok{{-}1}

\NormalTok{aov.out1 }\OtherTok{\textless{}{-}} \FunctionTok{cbind}\NormalTok{(df\_A,SSA,MSA,MSA}\SpecialCharTok{/}\NormalTok{MSBA,}\FunctionTok{qf}\NormalTok{(.}\DecValTok{95}\NormalTok{,df\_A,df\_BA),}\DecValTok{1}\SpecialCharTok{{-}}\FunctionTok{pf}\NormalTok{(MSA}\SpecialCharTok{/}\NormalTok{MSBA,df\_A,df\_BA),SSA}\SpecialCharTok{/}\NormalTok{SSTO, SSA}\SpecialCharTok{/}\NormalTok{(SSA}\SpecialCharTok{+}\NormalTok{SSE))}
\NormalTok{aov.out2 }\OtherTok{\textless{}{-}} \FunctionTok{cbind}\NormalTok{(df\_BA,SSBA,MSBA,MSBA}\SpecialCharTok{/}\NormalTok{MSE,}\FunctionTok{qf}\NormalTok{(.}\DecValTok{95}\NormalTok{,df\_BA,df\_E),}
                  \DecValTok{1}\SpecialCharTok{{-}}\FunctionTok{pf}\NormalTok{(MSBA}\SpecialCharTok{/}\NormalTok{MSE,df\_BA,df\_E),SSBA}\SpecialCharTok{/}\NormalTok{SSTO,SSBA}\SpecialCharTok{/}\NormalTok{(SSBA}\SpecialCharTok{+}\NormalTok{SSE))}
\NormalTok{aov.out3 }\OtherTok{\textless{}{-}} \FunctionTok{cbind}\NormalTok{(df\_E, SSE, MSE,}\ConstantTok{NA}\NormalTok{,}\ConstantTok{NA}\NormalTok{,}\ConstantTok{NA}\NormalTok{,}\ConstantTok{NA}\NormalTok{,}\ConstantTok{NA}\NormalTok{)}
\NormalTok{aov.out4 }\OtherTok{\textless{}{-}} \FunctionTok{cbind}\NormalTok{(df\_TO,SSTO,}\ConstantTok{NA}\NormalTok{,}\ConstantTok{NA}\NormalTok{,}\ConstantTok{NA}\NormalTok{,}\ConstantTok{NA}\NormalTok{,}\ConstantTok{NA}\NormalTok{,}\ConstantTok{NA}\NormalTok{)}
\NormalTok{aov.out }\OtherTok{\textless{}{-}} \FunctionTok{rbind}\NormalTok{(aov.out1,aov.out2,aov.out3,aov.out4)}
\FunctionTok{colnames}\NormalTok{(aov.out) }\OtherTok{\textless{}{-}} \FunctionTok{c}\NormalTok{(}\StringTok{"df"}\NormalTok{,}\StringTok{"SS"}\NormalTok{,}\StringTok{"MS"}\NormalTok{,}\StringTok{"F*"}\NormalTok{,}\StringTok{"F(.95)"}\NormalTok{,}\StringTok{"P(\textgreater{}F*)"}\NormalTok{,}\StringTok{"eta\^{}2"}\NormalTok{, }\StringTok{"Part eta\^{}2"}\NormalTok{)}
\FunctionTok{rownames}\NormalTok{(aov.out) }\OtherTok{\textless{}{-}} \FunctionTok{c}\NormalTok{(}\StringTok{"Model"}\NormalTok{,}\StringTok{"Trap(Model)"}\NormalTok{,}\StringTok{"Error"}\NormalTok{,}\StringTok{"Total"}\NormalTok{)}
\FunctionTok{round}\NormalTok{(aov.out,}\DecValTok{4}\NormalTok{)}
\end{Highlighting}
\end{Shaded}

\begin{verbatim}
##              df        SS        MS      F* F(.95) P(>F*)  eta^2 Part eta^2
## Model         7 35845.161 5120.7372 48.6523 2.6572      0 0.9013     0.9411
## Trap(Model)  16  1684.028  105.2518 10.1408 1.6905      0 0.0423     0.4290
## Error       216  2241.867   10.3790      NA     NA     NA     NA         NA
## Total       239 39771.056        NA      NA     NA     NA     NA         NA
\end{verbatim}

\begin{Shaded}
\begin{Highlighting}[]
\NormalTok{HSD\_A }\OtherTok{\textless{}{-}} \FunctionTok{qtukey}\NormalTok{(.}\DecValTok{95}\NormalTok{,a,df\_BA)}\SpecialCharTok{*}\FunctionTok{sqrt}\NormalTok{(MSBA}\SpecialCharTok{/}\NormalTok{(b}\SpecialCharTok{*}\NormalTok{n))}

\NormalTok{HSD\_A.out }\OtherTok{\textless{}{-}} \FunctionTok{matrix}\NormalTok{(}\FunctionTok{rep}\NormalTok{(}\DecValTok{0}\NormalTok{,}\DecValTok{8}\SpecialCharTok{*}\NormalTok{a}\SpecialCharTok{*}\NormalTok{(a}\DecValTok{{-}1}\NormalTok{)}\SpecialCharTok{/}\DecValTok{2}\NormalTok{), }\AttributeTok{ncol=}\DecValTok{8}\NormalTok{)}
\NormalTok{row\_A.count }\OtherTok{\textless{}{-}} \DecValTok{0}

\ControlFlowTok{for}\NormalTok{ (i1 }\ControlFlowTok{in} \DecValTok{2}\SpecialCharTok{:}\NormalTok{a) \{}
  \ControlFlowTok{for}\NormalTok{ (i2 }\ControlFlowTok{in} \DecValTok{1}\SpecialCharTok{:}\NormalTok{(i1}\DecValTok{{-}1}\NormalTok{)) \{}
\NormalTok{    row\_A.count }\OtherTok{\textless{}{-}}\NormalTok{ row\_A.count }\SpecialCharTok{+} \DecValTok{1}
\NormalTok{    HSD\_A.out[row\_A.count,}\DecValTok{1}\NormalTok{] }\OtherTok{\textless{}{-}}\NormalTok{ i1}
\NormalTok{    HSD\_A.out[row\_A.count,}\DecValTok{2}\NormalTok{] }\OtherTok{\textless{}{-}}\NormalTok{ i2}
\NormalTok{    HSD\_A.out[row\_A.count,}\DecValTok{3}\NormalTok{] }\OtherTok{\textless{}{-}}\NormalTok{ A\_mean[i1]}
\NormalTok{    HSD\_A.out[row\_A.count,}\DecValTok{4}\NormalTok{] }\OtherTok{\textless{}{-}}\NormalTok{ A\_mean[i2]}
\NormalTok{    HSD\_A.out[row\_A.count,}\DecValTok{5}\NormalTok{] }\OtherTok{\textless{}{-}}\NormalTok{ A\_mean[i1]}\SpecialCharTok{{-}}\NormalTok{A\_mean[i2]}
\NormalTok{    HSD\_A.out[row\_A.count,}\DecValTok{6}\NormalTok{] }\OtherTok{\textless{}{-}}\NormalTok{ (A\_mean[i1]}\SpecialCharTok{{-}}\NormalTok{A\_mean[i2])}\SpecialCharTok{{-}}\NormalTok{HSD\_A}
\NormalTok{    HSD\_A.out[row\_A.count,}\DecValTok{7}\NormalTok{] }\OtherTok{\textless{}{-}}\NormalTok{ (A\_mean[i1]}\SpecialCharTok{{-}}\NormalTok{A\_mean[i2])}\SpecialCharTok{+}\NormalTok{HSD\_A}
\NormalTok{    HSD\_A.out[row\_A.count,}\DecValTok{8}\NormalTok{] }\OtherTok{\textless{}{-}} 
      \DecValTok{1}\SpecialCharTok{{-}}\FunctionTok{ptukey}\NormalTok{(}\FunctionTok{abs}\NormalTok{((A\_mean[i1]}\SpecialCharTok{{-}}\NormalTok{A\_mean[i2])}\SpecialCharTok{/}\FunctionTok{sqrt}\NormalTok{(MSBA}\SpecialCharTok{/}\NormalTok{(b}\SpecialCharTok{*}\NormalTok{n))),a,df\_BA)}
\NormalTok{  \}}
\NormalTok{\}}
\FunctionTok{colnames}\NormalTok{(HSD\_A.out) }\OtherTok{\textless{}{-}} \FunctionTok{c}\NormalTok{(}\StringTok{"i"}\NormalTok{,}\StringTok{"i\textquotesingle{}"}\NormalTok{,}\StringTok{"Ybar\_i"}\NormalTok{,}\StringTok{"Ybar\_i\textquotesingle{}"}\NormalTok{,}\StringTok{"Diff"}\NormalTok{,}\StringTok{"LB"}\NormalTok{,}\StringTok{"UB"}\NormalTok{,}\StringTok{"p adj"}\NormalTok{)}
\FunctionTok{round}\NormalTok{(HSD\_A.out,}\DecValTok{4}\NormalTok{)}
\end{Highlighting}
\end{Shaded}

\begin{verbatim}
##       i i'  Ybar_i Ybar_i'    Diff      LB      UB  p adj
##  [1,] 2  1 55.8800 53.8400  2.0400 -7.1310 11.2110 0.9925
##  [2,] 3  1 59.1267 53.8400  5.2866 -3.8843 14.4576 0.5132
##  [3,] 3  2 59.1267 55.8800  3.2466 -5.9243 12.4176 0.9124
##  [4,] 4  1 68.8299 53.8400 14.9899  5.8189 24.1609 0.0007
##  [5,] 4  2 68.8299 55.8800 12.9499  3.7789 22.1209 0.0032
##  [6,] 4  3 68.8299 59.1267  9.7033  0.5323 18.8742 0.0342
##  [7,] 5  1 77.1966 53.8400 23.3566 14.1856 32.5276 0.0000
##  [8,] 5  2 77.1966 55.8800 21.3166 12.1456 30.4876 0.0000
##  [9,] 5  3 77.1966 59.1267 18.0700  8.8990 27.2409 0.0001
## [10,] 5  4 77.1966 68.8299  8.3667 -0.8043 17.5377 0.0874
## [11,] 6  1 81.1433 53.8400 27.3033 18.1323 36.4742 0.0000
## [12,] 6  2 81.1433 55.8800 25.2633 16.0923 34.4342 0.0000
## [13,] 6  3 81.1433 59.1267 22.0166 12.8457 31.1876 0.0000
## [14,] 6  4 81.1433 68.8299 12.3134  3.1424 21.4843 0.0051
## [15,] 6  5 81.1433 77.1966  3.9467 -5.2243 13.1176 0.8020
## [16,] 7  1 82.5467 53.8400 28.7066 19.5357 37.8776 0.0000
## [17,] 7  2 82.5467 55.8800 26.6666 17.4957 35.8376 0.0000
## [18,] 7  3 82.5467 59.1267 23.4200 14.2490 32.5910 0.0000
## [19,] 7  4 82.5467 68.8299 13.7167  4.5458 22.8877 0.0018
## [20,] 7  5 82.5467 77.1966  5.3500 -3.8209 14.5210 0.4993
## [21,] 7  6 82.5467 81.1433  1.4034 -7.7676 10.5743 0.9993
## [22,] 8  1 86.8833 53.8400 33.0433 23.8723 42.2143 0.0000
## [23,] 8  2 86.8833 55.8800 31.0033 21.8323 40.1743 0.0000
## [24,] 8  3 86.8833 59.1267 27.7567 18.5857 36.9276 0.0000
## [25,] 8  4 86.8833 68.8299 18.0534  8.8824 27.2244 0.0001
## [26,] 8  5 86.8833 77.1966  9.6867  0.5157 18.8577 0.0346
## [27,] 8  6 86.8833 81.1433  5.7400 -3.4309 14.9110 0.4177
## [28,] 8  7 86.8833 82.5467  4.3367 -4.8343 13.5076 0.7234
\end{verbatim}

\begin{Shaded}
\begin{Highlighting}[]
\NormalTok{s2\_ba }\OtherTok{\textless{}{-}}\NormalTok{ (MSBA}\SpecialCharTok{{-}}\NormalTok{MSE)}\SpecialCharTok{/}\NormalTok{n}
\NormalTok{df\_ba }\OtherTok{\textless{}{-}}\NormalTok{ s2\_ba}\SpecialCharTok{\^{}}\DecValTok{2} \SpecialCharTok{/}\NormalTok{ ((MSBA}\SpecialCharTok{/}\NormalTok{n)}\SpecialCharTok{\^{}}\DecValTok{2}\SpecialCharTok{/}\NormalTok{df\_BA }\SpecialCharTok{+}\NormalTok{ (}\SpecialCharTok{{-}}\NormalTok{MSE}\SpecialCharTok{/}\NormalTok{n)}\SpecialCharTok{\^{}}\DecValTok{2}\SpecialCharTok{/}\NormalTok{df\_E)}
\NormalTok{s2\_ba\_L }\OtherTok{\textless{}{-}}\NormalTok{ df\_ba}\SpecialCharTok{*}\NormalTok{s2\_ba}\SpecialCharTok{/}\FunctionTok{qchisq}\NormalTok{(}\DecValTok{1}\FloatTok{{-}.05}\SpecialCharTok{/}\DecValTok{2}\NormalTok{,df\_ba)}
\NormalTok{s2\_ba\_U }\OtherTok{\textless{}{-}}\NormalTok{ df\_ba}\SpecialCharTok{*}\NormalTok{s2\_ba}\SpecialCharTok{/}\FunctionTok{qchisq}\NormalTok{(.}\DecValTok{05}\SpecialCharTok{/}\DecValTok{2}\NormalTok{,df\_ba)}
\NormalTok{s2\_L }\OtherTok{\textless{}{-}}\NormalTok{ SSE}\SpecialCharTok{/}\FunctionTok{qchisq}\NormalTok{(}\DecValTok{1}\FloatTok{{-}.05}\SpecialCharTok{/}\DecValTok{2}\NormalTok{,df\_E)}
\NormalTok{s2\_U }\OtherTok{\textless{}{-}}\NormalTok{ SSE}\SpecialCharTok{/}\FunctionTok{qchisq}\NormalTok{(.}\DecValTok{05}\SpecialCharTok{/}\DecValTok{2}\NormalTok{,df\_E)}

\NormalTok{s2.out1 }\OtherTok{\textless{}{-}} \FunctionTok{cbind}\NormalTok{(df\_ba, s2\_ba, s2\_ba\_L, s2\_ba\_U)}
\NormalTok{s2.out2 }\OtherTok{\textless{}{-}} \FunctionTok{cbind}\NormalTok{(df\_E, MSE, s2\_L, s2\_U)}
\NormalTok{s2.out }\OtherTok{\textless{}{-}} \FunctionTok{rbind}\NormalTok{(s2.out1, s2.out2)}
\FunctionTok{colnames}\NormalTok{(s2.out) }\OtherTok{\textless{}{-}} \FunctionTok{c}\NormalTok{(}\StringTok{"df"}\NormalTok{, }\StringTok{"s\^{}2"}\NormalTok{, }\StringTok{"LB"}\NormalTok{, }\StringTok{"UB"}\NormalTok{)}
\FunctionTok{rownames}\NormalTok{(s2.out) }\OtherTok{\textless{}{-}} \FunctionTok{c}\NormalTok{(}\StringTok{"Trap(Model)"}\NormalTok{, }\StringTok{"Error"}\NormalTok{)}
\FunctionTok{round}\NormalTok{(s2.out,}\DecValTok{4}\NormalTok{)}
\end{Highlighting}
\end{Shaded}

\begin{verbatim}
##                   df     s^2     LB      UB
## Trap(Model)  12.9907  9.4873 4.9852 24.6343
## Error       216.0000 10.3790 8.6693 12.6524
\end{verbatim}

\begin{Shaded}
\begin{Highlighting}[]
\NormalTok{mod1 }\OtherTok{\textless{}{-}} \FunctionTok{aov}\NormalTok{(Y }\SpecialCharTok{\textasciitilde{}}\NormalTok{ model.f }\SpecialCharTok{+}\NormalTok{ model.f}\SpecialCharTok{/}\NormalTok{trapModel.f, }\AttributeTok{data=}\NormalTok{at)}
\FunctionTok{anova}\NormalTok{(mod1)}
\end{Highlighting}
\end{Shaded}

\begin{verbatim}
## Analysis of Variance Table
## 
## Response: Y
##                      Df Sum Sq Mean Sq F value    Pr(>F)    
## model.f               7  35845  5120.7 493.374 < 2.2e-16 ***
## model.f:trapModel.f  16   1684   105.3  10.141 < 2.2e-16 ***
## Residuals           216   2242    10.4                      
## ---
## Signif. codes:  0 '***' 0.001 '**' 0.01 '*' 0.05 '.' 0.1 ' ' 1
\end{verbatim}

\begin{Shaded}
\begin{Highlighting}[]
\FunctionTok{options}\NormalTok{(}\AttributeTok{contrasts=}\FunctionTok{c}\NormalTok{(}\StringTok{"contr.sum"}\NormalTok{, }\StringTok{"contr.poly"}\NormalTok{))}
\NormalTok{mod2 }\OtherTok{\textless{}{-}} \FunctionTok{lmer}\NormalTok{(Y }\SpecialCharTok{\textasciitilde{}}\NormalTok{ model.f }\SpecialCharTok{+}\NormalTok{ (}\DecValTok{1}\SpecialCharTok{|}\NormalTok{model.f}\SpecialCharTok{:}\NormalTok{trapModel.f), }\AttributeTok{data=}\NormalTok{at)}
\FunctionTok{summary}\NormalTok{(mod2)}
\end{Highlighting}
\end{Shaded}

\begin{verbatim}
## Linear mixed model fit by REML. t-tests use Satterthwaite's method [
## lmerModLmerTest]
## Formula: Y ~ model.f + (1 | model.f:trapModel.f)
##    Data: at
## 
## REML criterion at convergence: 1269.7
## 
## Scaled residuals: 
##      Min       1Q   Median       3Q      Max 
## -1.92671 -0.74501 -0.08267  0.72379  2.48273 
## 
## Random effects:
##  Groups              Name        Variance Std.Dev.
##  model.f:trapModel.f (Intercept)  9.487   3.080   
##  Residual                        10.379   3.222   
## Number of obs: 240, groups:  model.f:trapModel.f, 24
## 
## Fixed effects:
##             Estimate Std. Error       df t value Pr(>|t|)    
## (Intercept)  70.6808     0.6622  16.0000 106.731  < 2e-16 ***
## model.f1    -16.8408     1.7521  16.0000  -9.612 4.75e-08 ***
## model.f2    -14.8008     1.7521  16.0000  -8.447 2.72e-07 ***
## model.f3    -11.5542     1.7521  16.0000  -6.594 6.17e-06 ***
## model.f4     -1.8509     1.7521  16.0000  -1.056  0.30648    
## model.f5      6.5158     1.7521  16.0000   3.719  0.00187 ** 
## model.f6     10.4625     1.7521  16.0000   5.971 1.96e-05 ***
## model.f7     11.8658     1.7521  16.0000   6.772 4.48e-06 ***
## ---
## Signif. codes:  0 '***' 0.001 '**' 0.01 '*' 0.05 '.' 0.1 ' ' 1
## 
## Correlation of Fixed Effects:
##          (Intr) mdl.f1 mdl.f2 mdl.f3 mdl.f4 mdl.f5 mdl.f6
## model.f1  0.000                                          
## model.f2  0.000 -0.143                                   
## model.f3  0.000 -0.143 -0.143                            
## model.f4  0.000 -0.143 -0.143 -0.143                     
## model.f5  0.000 -0.143 -0.143 -0.143 -0.143              
## model.f6  0.000 -0.143 -0.143 -0.143 -0.143 -0.143       
## model.f7  0.000 -0.143 -0.143 -0.143 -0.143 -0.143 -0.143
\end{verbatim}

\begin{Shaded}
\begin{Highlighting}[]
\FunctionTok{anova}\NormalTok{(mod2)}
\end{Highlighting}
\end{Shaded}

\begin{verbatim}
## Type III Analysis of Variance Table with Satterthwaite's method
##         Sum Sq Mean Sq NumDF DenDF F value    Pr(>F)    
## model.f 3534.7  504.96     7    16  48.652 1.319e-09 ***
## ---
## Signif. codes:  0 '***' 0.001 '**' 0.01 '*' 0.05 '.' 0.1 ' ' 1
\end{verbatim}

There are many differences among the 28 pairs of trap models. Also,
within models there is relatively high variation among the individual
traps. \[ \nabla \]

When factors A and B are both random, we are generally interested in
estimating their variance components and testing whether they are 0.

\[
H_0^A:\sigma_{\alpha}^2=0  \qquad
TS:F_A^*=\frac{MSA}{MSB(A)} \qquad
RR:F_A^* \geq F_{1-\alpha;a-1,a(b-1)}
\]

\[
E\{MSA\}=\sigma^2 + n\sigma_{\beta(\alpha)}^2 + bn\sigma_{\alpha}^2
\qquad 
E\{MSB(A)\}=\sigma^2 + n\sigma_{\beta(\alpha)}^2
\] \[
s_{\alpha}^2 = \frac{MSA-MSB(A)}{bn}
\qquad
df_{\alpha}=\frac{\left(s_{\alpha}^2\right)^2}{
\left[\frac{\left(\frac{1}{bn}MSA\right)^2}{a-1} +
\frac{\left(-\frac{1}{bn}MSB(A)\right)^2}{a(b-1)} \right]
}
\]

\[
\mbox{Approximate } (1-\alpha)100\% \mbox{ CI for }
\sigma_{\alpha}^2: \quad
\left[
\frac{df_{\alpha}s_{\alpha}^2}{
\chi_{1-\alpha/2; df_{\alpha}}^2},
\frac{df_{\alpha}s_{\alpha}^2}{
\chi_{\alpha/2; df_{\alpha}}^2}
\right]
\] The test and Confidence Interval for \(\sigma_{\beta(\alpha)}^2\) are
the same for the random effects model as for the mixed effects model.

\textbf{Example 11.3 - Measurements on Semiconductor Wafers}

An instructional paper described making measurements on semiconducter
wafers (\citeproc{ref-jensen_2002}{Jensen 2002}). There were \(a=20\)
lots (batches) sampled, within each lot, there were \(b=2\) wafers
sampled, and each wafer was measured at \(n=9\) locations (replicates).
The response variable was not given for proprietary reasons. In this
example, we treat lot and wafer(lot) as random factors. There are two
wafer variables, \textbf{wafer1} takes on the values 1 and 2 within each
lot, \textbf{wafer2} gives each wafer a unique value (1-40) and is less
risky of creating an error in computing.

\begin{verbatim}
##   lotNum wafer wafer1 replicate       Y
## 1      1     1      1         1 181.247
## 2      1     1      1         2 181.280
## 3      1     1      1         3 185.021
## 4      1     1      1         4 180.144
## 5      1     1      1         5 192.570
## 6      1     1      1         6 178.741
\end{verbatim}

\begin{verbatim}
##     lotNum wafer wafer1 replicate       Y
## 355     20     2     40         4 161.783
## 356     20     2     40         5 160.013
## 357     20     2     40         6 160.981
## 358     20     2     40         7 162.384
## 359     20     2     40         8 163.720
## 360     20     2     40         9 161.716
\end{verbatim}

\pandocbounded{\includegraphics[keepaspectratio]{bookdownproj_STA4211Notes1_files/figure-latex/ch1103-1.pdf}}

\begin{verbatim}
##             df        SS       MS     F* F(.95) P(>F*)  eta^2 Part eta^2
## Lot         19 14050.871 739.5195 7.5475 2.1370      0 0.6358     0.6976
## Wafer(Lot)  20  1959.633  97.9817 5.1488 1.6035      0 0.0887     0.2435
## Error      320  6089.584  19.0299     NA     NA     NA     NA         NA
## Total      359 22100.088       NA     NA     NA     NA     NA         NA
\end{verbatim}

\begin{verbatim}
##                  df     s^2      LB      UB
## Lot         14.0642 35.6410 19.1268 88.4214
## Wafer(Lot)  12.9551  8.7724  4.6061 22.8150
## Error      320.0000 19.0299 16.3941 22.3605
\end{verbatim}

\begin{verbatim}
## Analysis of Variance Table
## 
## Response: Y
##                    Df  Sum Sq Mean Sq F value    Pr(>F)    
## lotNum.f           19 14050.9  739.52 38.8608 < 2.2e-16 ***
## lotNum.f:wafer1.f  20  1959.6   97.98  5.1488 3.708e-11 ***
## Residuals         320  6089.6   19.03                      
## ---
## Signif. codes:  0 '***' 0.001 '**' 0.01 '*' 0.05 '.' 0.1 ' ' 1
\end{verbatim}

\begin{verbatim}
## Linear mixed model fit by REML. t-tests use Satterthwaite's method [
## lmerModLmerTest]
## Formula: Y ~ 1 + (1 | lotNum.f) + (1 | lotNum.f:wafer1.f)
##    Data: scq
## 
## REML criterion at convergence: 2184.6
## 
## Scaled residuals: 
##     Min      1Q  Median      3Q     Max 
## -2.2897 -0.5826 -0.1729  0.4287  3.7529 
## 
## Random effects:
##  Groups            Name        Variance Std.Dev.
##  lotNum.f:wafer1.f (Intercept)  8.772   2.962   
##  lotNum.f          (Intercept) 35.641   5.970   
##  Residual                      19.030   4.362   
## Number of obs: 360, groups:  lotNum.f:wafer1.f, 40; lotNum.f, 20
## 
## Fixed effects:
##             Estimate Std. Error      df t value Pr(>|t|)    
## (Intercept)  174.342      1.433  19.000   121.6   <2e-16 ***
## ---
## Signif. codes:  0 '***' 0.001 '**' 0.01 '*' 0.05 '.' 0.1 ' ' 1
\end{verbatim}

The largest source of variation is among lots
\(\left(s_{\alpha}^2=35.64\right)\), followed by error (within wafer)
\(\left(s^2=19.03\right)\), followed by between wafers within lots
\(\left(s_{\beta(\alpha)}^2=8.77\right)\).

\[ \nabla \]

\section{Repeated Measures Design}\label{repeated-measures-design}

In this section we cover a second type of repeated measures design where
subjects are randomized to treatments as in the Completely Randomized
Design, then are measured at multiple time points. Recall in the chapter
on block designs, there was a repeated measures design where each
subject (block) received each treatment. In this case, each subject
receives only one of the treatments.

\textbf{Note that the notation used here is not the same as Kutner, et
al}.

\subsection{Model Structure}\label{model-structure}

The elements are described below for the case with a single treatment
factor (it can be extended to multiple factors) and a balanced design.

\begin{itemize}
\tightlist
\item
  \textbf{Treatments} - The primary factor of interest, with \(a\)
  (fixed) levels.
\item
  \textbf{Subjects} - Individuals included in the experiment \(b\)
  (random) units per treatment.
\item
  \textbf{Time Points} - When the measurements are obtained on each
  subject, with \(c\) (fixed) time points per individual.
\end{itemize}

We will use \(i\) to represent treatment, \(j\) represent the subject
within treatment and \(k\) to represent time point. The model is given
below.

\[
Y_{ijk}=\mu_{\bullet\bullet\bullet} + \alpha_i + \beta_{j(i)} +
\gamma_k + (\alpha\gamma)_{ik} + \epsilon_{ijk}  \quad
i=1,\ldots,a; \quad j=1,\ldots,b; \quad k=1,\ldots,c
\]

\[
\sum_{i=1}^a\alpha_i=\sum_{k=1}^c\gamma_k=
\sum_{i=1}^a(\alpha\gamma)_{ik}=\sum_{k=1}^c(\alpha\gamma)_{ik}=0
\] \[
\beta_{j(i)} \sim NID\left(0,\sigma_{\beta(\alpha)}^2\right) \qquad
\epsilon_{ijk} \sim NID\left(0,\sigma^2\right) \qquad
\left\{\beta_{j(i)}\right\} \bot
\left\{\epsilon_{ijk}\right\}
\]

\[
E\left\{Y_{ijk}\right\}=
\mu_{\bullet\bullet\bullet} + \alpha_i + \gamma_k + (\alpha\gamma)_{ik}
\qquad
V\left\{Y_{ijk}\right\}=\sigma_{\beta(\alpha)}^2+\sigma^2
\] \[
\mbox{COV}\left\{Y_{ijk},Y_{i'j'k'}\right\} = 
\left\{ \begin{array}{r@{\quad:\quad}l}
   \sigma_{\beta(\alpha)}^2 +  \sigma^2 & i=i',j=j',k=k' \\
    \sigma_{\beta(\alpha)}^2  & i=i', j=j',k\ne k' \\
    0 & \mbox{otherwise}\\
    \end{array} \right.
\] The variance-covariance within subjects has
\(\sigma_{\beta(\alpha)}^2 +  \sigma^2\) on the main diagonal and
\(\sigma_{\beta(\alpha)}^2\) on the off diagonals. In practice,
particularly when there are many time points, this structure is too
simple and a more complex structure can be fit.

\[
\mbox{COV}\left\{
\begin{bmatrix}
Y_{ij1} \\ Y_{ij2} \\ \vdots \\Y_{ijc} \\
\end{bmatrix}
\right\}
= 
\begin{bmatrix}
\sigma_{\beta(\alpha)}^2 +  \sigma^2 &
\sigma_{\beta(\alpha)}^2  &  \cdots &
\sigma_{\beta(\alpha)}^2 \\
\sigma_{\beta(\alpha)}^2 & 
\sigma_{\beta(\alpha)}^2 +  \sigma^2 &
\cdots & \sigma_{\beta(\alpha)}^2 \\
\vdots & \vdots & \ddots & \vdots \\
\sigma_{\beta(\alpha)}^2 & \sigma_{\beta(\alpha)}^2 & 
\cdots & \sigma_{\beta(\alpha)}^2 +  \sigma^2\\
\end{bmatrix}
\]

\subsection{\texorpdfstring{Analysis of Variance and
\(F\)-tests}{Analysis of Variance and F-tests}}\label{analysis-of-variance-and-f-tests}

For the various factors, we have he following sums of squares, degrees
of freedom, and expected mean squares.

\[
\mbox{Treatments:} \quad
SSA = bc\sum_{i=1}^a\left(\overline{Y}_{i\bullet\bullet}-
\overline{Y}_{\bullet\bullet\bullet}\right)^2 \qquad  df_A=a-1
\] \[
E\{MSA\}=\sigma^2 + c\sigma_{\beta(\alpha)}^2 + \frac{bc\sum_{i=1}^a\alpha_i^2}{a-1}
\]

\[
\mbox{Subjects(Trts):} \quad
SSB(A) = c\sum_{i=1}^a\sum_{j=1}^b \left(
\overline{Y}_{ij\bullet}-\overline{Y}_{i\bullet\bullet}\right)^2
\qquad
df_{B(A)}=a(b-1)
\] \[
E\{MSB(A)\}=\sigma^2 + c\sigma_{\beta(\alpha)}^2 
\] \[
\mbox{Time:} \quad
SSC=ab\sum_{k=1}^c
\left(\overline{Y}_{\bullet\bullet k}-
\overline{Y}_{\bullet\bullet\bullet}\right)^2 \qquad  df_C=c-1
\] \[
E\{MSC\} = \sigma^2 + \frac{ab\sum_{k=1}^c\gamma_k^2}{c-1}
\] \[
\mbox{TreatmentxTime:} \quad SSAC=
b\sum_{i=1}^a\sum_{k=1}^c\left(
\overline{Y}_{i\bullet k}-
\overline{Y}_{i\bullet\bullet}-
\overline{Y}_{\bullet\bullet k}+
\overline{Y}_{\bullet\bullet\bullet}\right)^2
\] \[
df_{AC}=(a-1)(c-1)  \qquad
E\{MSAC\} =\sigma^2 + 
\frac{b\sum_{i=1}^a\sum_{k=1}^c(\alpha\gamma)_{ik}^2}{(a-1)(c-1)}
\]

The error sum of squares represents the Subject(Treatment)xTime
interaction and has expectation \(E\{MSE\}=\sigma^2\).

\[
SSE=\sum_{i=1}^a\sum_{j=1}^b\sum_{k=1}^c \left(
Y_{ijk}-\overline{Y}_{ij\bullet}-\overline{Y}_{i\bullet k}+
\overline{Y}_{i\bullet\bullet}
\right)^2 \qquad
df_E = a(b-1)(c-1)
\]

To test for treatment effects, compare \(MSA\) to \(MSB(A)\). To test
for Time effects and TimexTreatment interaction, compare their mean
squares to \(MSE\).

\[
H_0^A: \alpha_1=\cdots=\alpha_a 
\qquad
TS: F_A^*=\frac{MSA}{MSB(A)} \qquad
RR: F_A^*\geq F_{1-\alpha;a-1,a(b-1)}
\]

\[
H_0^C:\gamma_1=\cdots =\gamma_k=0
\qquad
TS:F_C^*=\frac{MSC}{MSE} \qquad
RR: F_C^* \geq F_{1-\alpha;c-1,a(b-1)(c-1)}
\] \[
H_0^{AC}:(\alpha\gamma)_{11}=\cdots =(\alpha\gamma)_{ac}=0
\qquad
TS:F_{AC}^*=\frac{MSAC}{MSE} \qquad
RR: F_{AC}^* \geq F_{1-\alpha;(a-1)(c-1),a(b-1)(c-1)}
\] Depending on whether the TreatmentxTime interaction is important,
treatments can be compared within time points (interaction important) or
marginally (interaction unimportant).

\subsection{Treatment Comparisons}\label{treatment-comparisons}

First, consider the case where the interaction is not important, when
the ``time effect'' is consistent across treatments. Then we can use
Tukey's method for comparing all pairs of treatments, keeping in mind
that for treatments, \(MSB(A)\) is the error term.

\[
HSD = q_{1-\alpha;a,a(b-1)}\sqrt{\frac{MSB(A)}{bc}}
\] \[
(1-\alpha)100\% \mbox{ CI for } \alpha_i-\alpha_{i'}:
\quad
\left(\overline{Y}_{i\bullet\bullet}-
\overline{Y}_{i\bullet\bullet}\right) \pm HSD
\] When the TreatmentxTime interaction is important, we may wish to
compare the treatments within the time periods. This involves taking the
standard error for the difference between means within the time periods
as a linear function of \(MSB(A)\) and \(MSE\) and using Satterthwaite's
approximation for the degrees of freedom.

\[
V\left\{\overline{Y}_{i\bullet k} - 
\overline{Y}_{i'\bullet k}\right\} =
\frac{2\left(\sigma^2+\sigma_{\beta(\alpha)}^2\right)}{b} \qquad
E\{MSB(A)+(c-1)MSE\}=c\left(\sigma^2+\sigma_{\beta(\alpha)}^2\right)
\]

So if we divide \(MSB(A)+(c-1)MSE\) by \(bc\) and multiply by 2, this
has expectation equal to the variance of the mean difference.

\[
\hat{V}\left\{\overline{Y}_{i\bullet k} - 
\overline{Y}_{i'\bullet k}\right\} = 
\frac{2\left(MSB(A)+(c-1)MSE\right)}{bc}
\] \[
df^* = \frac{\left(MSB(A)+(c-1)MSE\right)^2}{
\left[ \frac{\left(MSB(A)\right)^2}{a(b-1)} + 
\frac{\left((c-1)MSE\right)^2}{a(b-1)(c-1)} \right]
}
\]

Then we can use Tukey's method to compare all pairs of treatments within
specific time points.

\[
HSD_k= q_{1-\alpha;a,df^*}\sqrt{\frac{MSB(A)+(c-1)MSE}{bc}}
\] \[
(1-\alpha)100\% \mbox{ Confidence Intervals:} \quad
\left(\overline{Y}_{i\bullet k}-
\overline{Y}_{i'\bullet k}\right) \pm HSD_k
\] \textbf{Example 11.4 - Treating Cats with Anxiety}

A study was conducted to test a drug zylkene versus placebo in cats with
anxiety (\citeproc{ref-beata_2007}{Beata, Cordel, and Marlois 2007}).
There were \(a=2\) treatments (zylkene \(i=1\) and placebo \(i=2\)),
with \(b=17\) cats randomized to each treatment (34 total) and each cat
was observed at each of \(c=5\) time points. The response was a clinical
globel impression (higher scores mean less anxiety and are better). R
code and output are given below.

\begin{Shaded}
\begin{Highlighting}[]
\NormalTok{caz }\OtherTok{\textless{}{-}} \FunctionTok{read.csv}\NormalTok{(}\StringTok{"data/cat\_zyl\_uni.csv"}\NormalTok{)}
\FunctionTok{head}\NormalTok{(caz); }\FunctionTok{tail}\NormalTok{(caz)}
\end{Highlighting}
\end{Shaded}

\begin{verbatim}
##   id weight age gender trt timepnt y
## 1  2    4.0  67      1   0       1 9
## 2  2    4.0  67      1   0       2 9
## 3  2    4.0  67      1   0       3 9
## 4  2    4.0  67      1   0       4 9
## 5  2    4.0  67      1   0       5 9
## 6  4    3.5  78      2   0       1 9
\end{verbatim}

\begin{verbatim}
##     id weight age gender trt timepnt  y
## 165 31    3.1  90      1   1       5 23
## 166 32    4.0  56      1   1       1  8
## 167 32    4.0  56      1   1       2 11
## 168 32    4.0  56      1   1       3 13
## 169 32    4.0  56      1   1       4 14
## 170 32    4.0  56      1   1       5 17
\end{verbatim}

\begin{Shaded}
\begin{Highlighting}[]
\NormalTok{caz}\SpecialCharTok{$}\NormalTok{trt.f }\OtherTok{\textless{}{-}} \FunctionTok{factor}\NormalTok{(caz}\SpecialCharTok{$}\NormalTok{trt, }\AttributeTok{levels=}\DecValTok{0}\SpecialCharTok{:}\DecValTok{1}\NormalTok{, }\AttributeTok{labels=}\FunctionTok{c}\NormalTok{(}\StringTok{"placebo"}\NormalTok{, }\StringTok{"zylkene"}\NormalTok{))}
\NormalTok{caz}\SpecialCharTok{$}\NormalTok{timepnt.f }\OtherTok{\textless{}{-}} \FunctionTok{factor}\NormalTok{(caz}\SpecialCharTok{$}\NormalTok{timepnt)}
\NormalTok{caz}\SpecialCharTok{$}\NormalTok{id1.f }\OtherTok{\textless{}{-}} \FunctionTok{factor}\NormalTok{(}\FunctionTok{rep}\NormalTok{(}\DecValTok{1}\SpecialCharTok{:}\DecValTok{34}\NormalTok{,}\AttributeTok{each=}\DecValTok{5}\NormalTok{))}


\NormalTok{A\_mean }\OtherTok{\textless{}{-}} \FunctionTok{as.vector}\NormalTok{(}\FunctionTok{tapply}\NormalTok{(caz}\SpecialCharTok{$}\NormalTok{y, caz}\SpecialCharTok{$}\NormalTok{trt, mean))}
\NormalTok{BA\_mean }\OtherTok{\textless{}{-}} \FunctionTok{as.vector}\NormalTok{(}\FunctionTok{tapply}\NormalTok{(caz}\SpecialCharTok{$}\NormalTok{y, caz}\SpecialCharTok{$}\NormalTok{id1.f, mean))}
\NormalTok{C\_mean }\OtherTok{\textless{}{-}} \FunctionTok{as.vector}\NormalTok{(}\FunctionTok{tapply}\NormalTok{(caz}\SpecialCharTok{$}\NormalTok{y, caz}\SpecialCharTok{$}\NormalTok{timepnt.f, mean))}
\NormalTok{AC\_mean }\OtherTok{\textless{}{-}} \FunctionTok{as.vector}\NormalTok{(}\FunctionTok{tapply}\NormalTok{(caz}\SpecialCharTok{$}\NormalTok{y,}\FunctionTok{list}\NormalTok{(caz}\SpecialCharTok{$}\NormalTok{timepnt.f,caz}\SpecialCharTok{$}\NormalTok{trt.f),mean))}
\NormalTok{all\_mean }\OtherTok{\textless{}{-}} \FunctionTok{mean}\NormalTok{(caz}\SpecialCharTok{$}\NormalTok{y)}
\NormalTok{a }\OtherTok{\textless{}{-}} \FunctionTok{length}\NormalTok{(A\_mean)}
\NormalTok{b }\OtherTok{\textless{}{-}} \FunctionTok{length}\NormalTok{(BA\_mean) }\SpecialCharTok{/}\NormalTok{ a}
\NormalTok{c }\OtherTok{\textless{}{-}} \FunctionTok{length}\NormalTok{(C\_mean)}

\NormalTok{Ybar\_i.. }\OtherTok{\textless{}{-}} \FunctionTok{rep}\NormalTok{(A\_mean, }\AttributeTok{each=}\NormalTok{b}\SpecialCharTok{*}\NormalTok{c)}
\NormalTok{Ybar\_ij. }\OtherTok{\textless{}{-}} \FunctionTok{rep}\NormalTok{(BA\_mean, }\AttributeTok{each=}\NormalTok{c)}
\NormalTok{Ybar\_..k }\OtherTok{\textless{}{-}} \FunctionTok{rep}\NormalTok{(C\_mean, }\AttributeTok{times=}\NormalTok{a}\SpecialCharTok{*}\NormalTok{b)}
\NormalTok{Ybar\_i.k }\OtherTok{\textless{}{-}} \FunctionTok{c}\NormalTok{(}\FunctionTok{rep}\NormalTok{(AC\_mean[}\DecValTok{1}\SpecialCharTok{:}\NormalTok{c],}\AttributeTok{times=}\NormalTok{b),}\FunctionTok{rep}\NormalTok{(AC\_mean[(c}\SpecialCharTok{+}\DecValTok{1}\NormalTok{)}\SpecialCharTok{:}\NormalTok{(}\DecValTok{2}\SpecialCharTok{*}\NormalTok{c)], }\AttributeTok{times=}\NormalTok{b))}
\NormalTok{Ybar\_... }\OtherTok{\textless{}{-}} \FunctionTok{rep}\NormalTok{(all\_mean, a}\SpecialCharTok{*}\NormalTok{b}\SpecialCharTok{*}\NormalTok{c)}

\NormalTok{SSA }\OtherTok{\textless{}{-}} \FunctionTok{sum}\NormalTok{((Ybar\_i..}\SpecialCharTok{{-}}\NormalTok{Ybar\_...)}\SpecialCharTok{\^{}}\DecValTok{2}\NormalTok{); df\_A }\OtherTok{\textless{}{-}}\NormalTok{ a}\DecValTok{{-}1}\NormalTok{; MSA }\OtherTok{\textless{}{-}}\NormalTok{ SSA}\SpecialCharTok{/}\NormalTok{df\_A}
\NormalTok{SSBA }\OtherTok{\textless{}{-}} \FunctionTok{sum}\NormalTok{((Ybar\_ij.}\SpecialCharTok{{-}}\NormalTok{Ybar\_i..)}\SpecialCharTok{\^{}}\DecValTok{2}\NormalTok{); df\_BA }\OtherTok{\textless{}{-}}\NormalTok{ a}\SpecialCharTok{*}\NormalTok{(b}\DecValTok{{-}1}\NormalTok{); MSBA }\OtherTok{\textless{}{-}}\NormalTok{ SSBA}\SpecialCharTok{/}\NormalTok{df\_BA}
\NormalTok{SSC }\OtherTok{\textless{}{-}} \FunctionTok{sum}\NormalTok{((Ybar\_..k}\SpecialCharTok{{-}}\NormalTok{Ybar\_...)}\SpecialCharTok{\^{}}\DecValTok{2}\NormalTok{); df\_C }\OtherTok{\textless{}{-}}\NormalTok{ c}\DecValTok{{-}1}\NormalTok{; MSC }\OtherTok{\textless{}{-}}\NormalTok{ SSC}\SpecialCharTok{/}\NormalTok{df\_C}
\NormalTok{SSAC }\OtherTok{\textless{}{-}} \FunctionTok{sum}\NormalTok{((Ybar\_i.k}\SpecialCharTok{{-}}\NormalTok{Ybar\_i..}\SpecialCharTok{{-}}\NormalTok{Ybar\_..k}\SpecialCharTok{+}\NormalTok{Ybar\_...)}\SpecialCharTok{\^{}}\DecValTok{2}\NormalTok{); df\_AC }\OtherTok{\textless{}{-}}\NormalTok{ (a}\DecValTok{{-}1}\NormalTok{)}\SpecialCharTok{*}\NormalTok{(c}\DecValTok{{-}1}\NormalTok{)}
\NormalTok{   MSAC }\OtherTok{\textless{}{-}}\NormalTok{ SSAC}\SpecialCharTok{/}\NormalTok{df\_AC}
\NormalTok{SSE }\OtherTok{\textless{}{-}} \FunctionTok{sum}\NormalTok{((caz}\SpecialCharTok{$}\NormalTok{y}\SpecialCharTok{{-}}\NormalTok{Ybar\_ij.}\SpecialCharTok{{-}}\NormalTok{Ybar\_i.k}\SpecialCharTok{+}\NormalTok{Ybar\_i..)}\SpecialCharTok{\^{}}\DecValTok{2}\NormalTok{); }
\NormalTok{   df\_E }\OtherTok{\textless{}{-}}\NormalTok{ a}\SpecialCharTok{*}\NormalTok{(b}\DecValTok{{-}1}\NormalTok{)}\SpecialCharTok{*}\NormalTok{(c}\DecValTok{{-}1}\NormalTok{); MSE }\OtherTok{\textless{}{-}}\NormalTok{ SSE}\SpecialCharTok{/}\NormalTok{df\_E}
\NormalTok{SSTO }\OtherTok{\textless{}{-}} \FunctionTok{sum}\NormalTok{((caz}\SpecialCharTok{$}\NormalTok{y}\SpecialCharTok{{-}}\NormalTok{Ybar\_...)}\SpecialCharTok{\^{}}\DecValTok{2}\NormalTok{); df\_TO }\OtherTok{\textless{}{-}}\NormalTok{ a}\SpecialCharTok{*}\NormalTok{b}\SpecialCharTok{*}\NormalTok{c}\DecValTok{{-}1}  

\NormalTok{aov.out1 }\OtherTok{\textless{}{-}} \FunctionTok{cbind}\NormalTok{(df\_A,SSA,MSA,MSA}\SpecialCharTok{/}\NormalTok{MSBA,}\FunctionTok{qf}\NormalTok{(.}\DecValTok{95}\NormalTok{,df\_A,df\_BA),}
                  \DecValTok{1}\SpecialCharTok{{-}}\FunctionTok{pf}\NormalTok{(MSA}\SpecialCharTok{/}\NormalTok{MSBA,df\_A,df\_BA),SSA}\SpecialCharTok{/}\NormalTok{SSTO)}
\NormalTok{aov.out2 }\OtherTok{\textless{}{-}} \FunctionTok{cbind}\NormalTok{(df\_BA,SSBA,MSBA,MSBA}\SpecialCharTok{/}\NormalTok{MSE,}\FunctionTok{qf}\NormalTok{(.}\DecValTok{95}\NormalTok{,df\_BA,df\_E),}
                  \DecValTok{1}\SpecialCharTok{{-}}\FunctionTok{pf}\NormalTok{(MSBA}\SpecialCharTok{/}\NormalTok{MSE,df\_BA,df\_E),SSBA}\SpecialCharTok{/}\NormalTok{SSTO)}
\NormalTok{aov.out3 }\OtherTok{\textless{}{-}} \FunctionTok{cbind}\NormalTok{(df\_C,SSC,MSC,MSC}\SpecialCharTok{/}\NormalTok{MSE,}\FunctionTok{qf}\NormalTok{(.}\DecValTok{95}\NormalTok{,df\_C,df\_E),}
                  \DecValTok{1}\SpecialCharTok{{-}}\FunctionTok{pf}\NormalTok{(MSC}\SpecialCharTok{/}\NormalTok{MSE,df\_C,df\_E),SSC}\SpecialCharTok{/}\NormalTok{SSTO)}
\NormalTok{aov.out4 }\OtherTok{\textless{}{-}} \FunctionTok{cbind}\NormalTok{(df\_AC,SSAC,MSAC,MSAC}\SpecialCharTok{/}\NormalTok{MSE,}\FunctionTok{qf}\NormalTok{(.}\DecValTok{95}\NormalTok{,df\_AC,df\_E),}
                  \DecValTok{1}\SpecialCharTok{{-}}\FunctionTok{pf}\NormalTok{(MSAC}\SpecialCharTok{/}\NormalTok{MSE,df\_AC,df\_E),SSAC}\SpecialCharTok{/}\NormalTok{SSTO)}
\NormalTok{aov.out5 }\OtherTok{\textless{}{-}} \FunctionTok{cbind}\NormalTok{(df\_E, SSE, MSE,}\ConstantTok{NA}\NormalTok{,}\ConstantTok{NA}\NormalTok{,}\ConstantTok{NA}\NormalTok{,}\ConstantTok{NA}\NormalTok{)}
\NormalTok{aov.out6 }\OtherTok{\textless{}{-}} \FunctionTok{cbind}\NormalTok{(df\_TO,SSTO,}\ConstantTok{NA}\NormalTok{,}\ConstantTok{NA}\NormalTok{,}\ConstantTok{NA}\NormalTok{,}\ConstantTok{NA}\NormalTok{,}\ConstantTok{NA}\NormalTok{)}
\NormalTok{aov.out }\OtherTok{\textless{}{-}} \FunctionTok{rbind}\NormalTok{(aov.out1,aov.out2,aov.out3,aov.out4,aov.out5,aov.out6)}
\FunctionTok{colnames}\NormalTok{(aov.out) }\OtherTok{\textless{}{-}} \FunctionTok{c}\NormalTok{(}\StringTok{"df"}\NormalTok{,}\StringTok{"SS"}\NormalTok{,}\StringTok{"MS"}\NormalTok{,}\StringTok{"F*"}\NormalTok{,}\StringTok{"F(.95)"}\NormalTok{,}\StringTok{"P(\textgreater{}F*)"}\NormalTok{,}\StringTok{"eta\^{}2"}\NormalTok{)}
\FunctionTok{rownames}\NormalTok{(aov.out) }\OtherTok{\textless{}{-}} \FunctionTok{c}\NormalTok{(}\StringTok{"Drug"}\NormalTok{,}\StringTok{"Cat(Drug)"}\NormalTok{,}\StringTok{"Time"}\NormalTok{,}
                       \StringTok{"DrugxTime"}\NormalTok{, }\StringTok{"Error"}\NormalTok{,}\StringTok{"Total"}\NormalTok{)}
\FunctionTok{round}\NormalTok{(aov.out,}\DecValTok{4}\NormalTok{)}
\end{Highlighting}
\end{Shaded}

\begin{verbatim}
##            df        SS       MS      F* F(.95) P(>F*)  eta^2
## Drug        1  382.5000 382.5000  5.7412 4.1491 0.0226 0.1167
## Cat(Drug)  32 2131.9765  66.6243 21.9512 1.5341 0.0000 0.6505
## Time        4  324.1176  81.0294 26.6974 2.4425 0.0000 0.0989
## DrugxTime   4   50.5882  12.6471  4.1669 2.4425 0.0033 0.0154
## Error     128  388.4941   3.0351      NA     NA     NA     NA
## Total     169 3277.6765       NA      NA     NA     NA     NA
\end{verbatim}

\begin{Shaded}
\begin{Highlighting}[]
\NormalTok{mod1 }\OtherTok{\textless{}{-}} \FunctionTok{aov}\NormalTok{(y }\SpecialCharTok{\textasciitilde{}}\NormalTok{ trt.f }\SpecialCharTok{+}\NormalTok{ trt.f}\SpecialCharTok{/}\NormalTok{id1.f }\SpecialCharTok{+}\NormalTok{ timepnt.f }\SpecialCharTok{+}\NormalTok{ trt.f}\SpecialCharTok{:}\NormalTok{timepnt.f,}
            \AttributeTok{data=}\NormalTok{caz)}
\FunctionTok{anova}\NormalTok{(mod1)}
\end{Highlighting}
\end{Shaded}

\begin{verbatim}
## Analysis of Variance Table
## 
## Response: y
##                  Df  Sum Sq Mean Sq  F value    Pr(>F)    
## trt.f             1  382.50  382.50 126.0251 < 2.2e-16 ***
## timepnt.f         4  324.12   81.03  26.6974 4.137e-16 ***
## trt.f:id1.f      32 2131.98   66.62  21.9512 < 2.2e-16 ***
## trt.f:timepnt.f   4   50.59   12.65   4.1669  0.003316 ** 
## Residuals       128  388.49    3.04                       
## ---
## Signif. codes:  0 '***' 0.001 '**' 0.01 '*' 0.05 '.' 0.1 ' ' 1
\end{verbatim}

\begin{Shaded}
\begin{Highlighting}[]
\DocumentationTok{\#\# This model uses correct error term for treatment (B(A))}
\NormalTok{mod2 }\OtherTok{\textless{}{-}} \FunctionTok{aov}\NormalTok{(y }\SpecialCharTok{\textasciitilde{}}\NormalTok{ trt.f}\SpecialCharTok{*}\NormalTok{timepnt.f }\SpecialCharTok{+} \FunctionTok{Error}\NormalTok{(id1.f), }\AttributeTok{data=}\NormalTok{caz)}
\FunctionTok{summary}\NormalTok{(mod2)}
\end{Highlighting}
\end{Shaded}

\begin{verbatim}
## 
## Error: id1.f
##           Df Sum Sq Mean Sq F value Pr(>F)  
## trt.f      1  382.5   382.5   5.741 0.0226 *
## Residuals 32 2132.0    66.6                 
## ---
## Signif. codes:  0 '***' 0.001 '**' 0.01 '*' 0.05 '.' 0.1 ' ' 1
## 
## Error: Within
##                  Df Sum Sq Mean Sq F value   Pr(>F)    
## timepnt.f         4  324.1   81.03  26.697 4.14e-16 ***
## trt.f:timepnt.f   4   50.6   12.65   4.167  0.00332 ** 
## Residuals       128  388.5    3.04                     
## ---
## Signif. codes:  0 '***' 0.001 '**' 0.01 '*' 0.05 '.' 0.1 ' ' 1
\end{verbatim}

\begin{Shaded}
\begin{Highlighting}[]
\FunctionTok{options}\NormalTok{(}\AttributeTok{es.use\_symbols=}\ConstantTok{TRUE}\NormalTok{)}

\FunctionTok{eta\_squared}\NormalTok{(mod1)}
\end{Highlighting}
\end{Shaded}

\begin{verbatim}
## # Effect Size for ANOVA (Type I)
## 
## Parameter       | η² (partial) |       95% CI
## ---------------------------------------------
## trt.f           |         0.50 | [0.40, 1.00]
## timepnt.f       |         0.45 | [0.34, 1.00]
## trt.f:id1.f     |         0.85 | [0.80, 1.00]
## trt.f:timepnt.f |         0.12 | [0.03, 1.00]
## 
## - One-sided CIs: upper bound fixed at [1.00].
\end{verbatim}

\begin{Shaded}
\begin{Highlighting}[]
\FunctionTok{eta\_squared}\NormalTok{(mod1, }\AttributeTok{partial=}\ConstantTok{FALSE}\NormalTok{)}
\end{Highlighting}
\end{Shaded}

\begin{verbatim}
## # Effect Size for ANOVA (Type I)
## 
## Parameter       |   η² |       95% CI
## -------------------------------------
## trt.f           | 0.12 | [0.04, 1.00]
## timepnt.f       | 0.10 | [0.02, 1.00]
## trt.f:id1.f     | 0.65 | [0.53, 1.00]
## trt.f:timepnt.f | 0.02 | [0.00, 1.00]
## 
## - One-sided CIs: upper bound fixed at [1.00].
\end{verbatim}

Based on the ANOVA model, fit based on the formulas described in this
section, we observe a significant interaction between Treatment and Time
\(\left(F_{AC}^*=4.167, P=.0033\right)\), as well as significant main
effects for Treatment and Time. Here, we will compare the two treatments
within the five time points using 95\% Confidence Intervals for
\(\mu_Z-\mu_P\) within each time point (1 comparison per time). We will
use \(i=1\) to represent Zylkene and \(i=2\) to represent Placebo.

\[
\hat{V}\left\{\overline{Y}_{1\bullet k}-\overline{Y}_{2\bullet k}\right\} =
\frac{2\left(MSB(A)+(c-1)MSE\right)}{bc} =
\frac{2\left(66.62+(5-1)(3.04)\right)}{17(5)}=1.854
\] \[
df^* = \frac{\left(66.62+4(3.04)\right)^2}{
       \left[\frac{(66.62)^2}{32}+\frac{(4(3.04))^2}{128}\right]}=
       \frac{6206.29}{139.85}=44.4
\] \[
t_{.975;44.4}=2.015 \quad \Rightarrow 
\quad
\left(\overline{Y}_{1\bullet k}-\overline{Y}_{2\bullet k}\right)\pm
2.015\sqrt{1.854}
 \equiv  
\left(\overline{Y}_{1\bullet k}-\overline{Y}_{2\bullet k}\right)\pm
2.74
\] For the 5 time points, we compute the following Confidence Intervals.

\[
k=1: \quad (10.94-9.06) \pm 2.74 \quad \equiv \quad 1.88\pm 2.74
\equiv (-0.86,4.62)
\] \[
k=2: \quad (12.12-10.18) \pm 2.74  \equiv  1.94\pm 2.74
\equiv (-0.80,4.68)
\] \[
k=3: \quad (13.59-10.88) \pm 2.74  \equiv  2.71\pm 2.74
\equiv (-0.03,5.45)
\] \[
k=4: \quad (15.18-11.41) \pm 3.77  \equiv  3.77\pm 2.74
\equiv (1.03,6.51)
\] \[
k=5: \quad (16.12-11.41) \pm 4.71  \equiv  4.71\pm 2.74
\equiv (1.97,7.45)
\] At time point 3, they are very close to being significantly different
as the interval just contains 0. At time points 4 and 5, Zylkene gives
significantly higher scores than placebo.

Next, we consider two other covariance structures among measurements
within cats: \textbf{unstructured} and \textbf{AR(1)} which is
autoregressive of order 1. We will use the \textbf{lme} function in the
\textbf{nlme} package to fit them.

The unstructured case has no restrictions on the variances and
covariances.

\[
\mbox{COV}\left\{
\begin{bmatrix}
Y_{ij1} \\ Y_{ij2} \\ \vdots \\Y_{ijc} \\
\end{bmatrix}
\right\}
= 
\begin{bmatrix}
\sigma_{1}^2   &
\sigma_{12}  &  \cdots &
\sigma_{1c} \\
\sigma_{12} & 
\sigma_2^2 &
\cdots & \sigma_{2c} \\
\vdots & \vdots & \ddots & \vdots \\
\sigma_{1c} & \sigma_{2c} & 
\cdots & \sigma_c^2\\
\end{bmatrix}
\]

For the AR(1) case, the variances are all equal, and the covariances
(and correlations) decrease multiplicatively as the observations are
further apart in time. This can be extended to allow for heterogeneity
of variances, which is not covered here. In the following, we assume
\(|\phi |<1\).

\[
\mbox{COV}\left\{
\begin{bmatrix}
Y_{ij1} \\ Y_{ij2} \\ \vdots \\Y_{ijc} \\
\end{bmatrix}
\right\}
= 
\begin{bmatrix}
\sigma^2 &
\phi\sigma^2  &  \cdots &
\phi^{c-1}\sigma^2 \\
\phi\sigma^2 & 
\sigma^2 &
\cdots & \phi^{c-2}\sigma^2 \\
\vdots & \vdots & \ddots & \vdots \\
\phi^{c-1}\sigma^2 & \phi^{c-2}\sigma^2 & 
\cdots & \sigma^2\\
\end{bmatrix}
\]

\begin{Shaded}
\begin{Highlighting}[]
\NormalTok{caz }\OtherTok{\textless{}{-}} \FunctionTok{read.csv}\NormalTok{(}\StringTok{"data/cat\_zyl\_uni.csv"}\NormalTok{)}
\FunctionTok{head}\NormalTok{(caz); }\FunctionTok{tail}\NormalTok{(caz)}
\end{Highlighting}
\end{Shaded}

\begin{verbatim}
##   id weight age gender trt timepnt y
## 1  2    4.0  67      1   0       1 9
## 2  2    4.0  67      1   0       2 9
## 3  2    4.0  67      1   0       3 9
## 4  2    4.0  67      1   0       4 9
## 5  2    4.0  67      1   0       5 9
## 6  4    3.5  78      2   0       1 9
\end{verbatim}

\begin{verbatim}
##     id weight age gender trt timepnt  y
## 165 31    3.1  90      1   1       5 23
## 166 32    4.0  56      1   1       1  8
## 167 32    4.0  56      1   1       2 11
## 168 32    4.0  56      1   1       3 13
## 169 32    4.0  56      1   1       4 14
## 170 32    4.0  56      1   1       5 17
\end{verbatim}

\begin{Shaded}
\begin{Highlighting}[]
\NormalTok{caz}\SpecialCharTok{$}\NormalTok{trt.f }\OtherTok{\textless{}{-}} \FunctionTok{factor}\NormalTok{(caz}\SpecialCharTok{$}\NormalTok{trt, }\AttributeTok{levels=}\DecValTok{0}\SpecialCharTok{:}\DecValTok{1}\NormalTok{, }\AttributeTok{labels=}\FunctionTok{c}\NormalTok{(}\StringTok{"placebo"}\NormalTok{, }\StringTok{"zylkene"}\NormalTok{))}
\NormalTok{caz}\SpecialCharTok{$}\NormalTok{timepnt.f }\OtherTok{\textless{}{-}} \FunctionTok{factor}\NormalTok{(caz}\SpecialCharTok{$}\NormalTok{timepnt)}
\NormalTok{caz}\SpecialCharTok{$}\NormalTok{id1.f }\OtherTok{\textless{}{-}} \FunctionTok{factor}\NormalTok{(}\FunctionTok{rep}\NormalTok{(}\DecValTok{1}\SpecialCharTok{:}\DecValTok{34}\NormalTok{,}\AttributeTok{each=}\DecValTok{5}\NormalTok{))}

\DocumentationTok{\#\#\# Uses linear mixed effect model for various covariance structures}
\DocumentationTok{\#\#\#   within cats}

\DocumentationTok{\#\# Compound Symmetry that prints marginal V{-}Cov matrix = ZGZ\textquotesingle{}+R w/in cats}
\DocumentationTok{\#\#   Same as anova model above}
\NormalTok{fit.cs1 }\OtherTok{\textless{}{-}} \FunctionTok{lme}\NormalTok{(y }\SpecialCharTok{\textasciitilde{}}\NormalTok{ trt.f}\SpecialCharTok{*}\NormalTok{timepnt.f, }\AttributeTok{data=}\NormalTok{caz,}
  \AttributeTok{random=} \FunctionTok{list}\NormalTok{(}\AttributeTok{id =} \FunctionTok{pdCompSymm}\NormalTok{(}\SpecialCharTok{\textasciitilde{}}\NormalTok{ trt.f }\SpecialCharTok{{-}} \DecValTok{1}\NormalTok{)))}
\FunctionTok{summary}\NormalTok{(fit.cs1)}
\end{Highlighting}
\end{Shaded}

\begin{verbatim}
## Linear mixed-effects model fit by REML
##   Data: caz 
##        AIC      BIC    logLik
##   798.2436 838.2209 -386.1218
## 
## Random effects:
##  Formula: ~trt.f - 1 | id
##  Structure: Compound Symmetry
##              StdDev   Corr
## trt.fplacebo 3.566192     
## trt.fzylkene 3.566192 0   
## Residual     1.742159     
## 
## Fixed effects:  y ~ trt.f * timepnt.f 
##                       Value Std.Error  DF   t-value p-value
## (Intercept)       12.088235 0.6260226 128 19.309583  0.0000
## trt.f1            -1.500000 0.6260226  32 -2.396080  0.0226
## timepnt.f1        -2.088235 0.2672349 128 -7.814231  0.0000
## timepnt.f2        -0.941176 0.2672349 128 -3.521907  0.0006
## timepnt.f3         0.147059 0.2672349 128  0.550298  0.5831
## timepnt.f4         1.205882 0.2672349 128  4.512443  0.0000
## trt.f1:timepnt.f1  0.558824 0.2672349 128  2.091132  0.0385
## trt.f1:timepnt.f2  0.529412 0.2672349 128  1.981073  0.0497
## trt.f1:timepnt.f3  0.147059 0.2672349 128  0.550298  0.5831
## trt.f1:timepnt.f4 -0.382353 0.2672349 128 -1.430775  0.1549
##  Correlation: 
##                   (Intr) trt.f1 tmpn.1 tmpn.2 tmpn.3 tmpn.4 t.1:.1 t.1:.2
## trt.f1             0.00                                                  
## timepnt.f1         0.00   0.00                                           
## timepnt.f2         0.00   0.00  -0.25                                    
## timepnt.f3         0.00   0.00  -0.25  -0.25                             
## timepnt.f4         0.00   0.00  -0.25  -0.25  -0.25                      
## trt.f1:timepnt.f1  0.00   0.00   0.00   0.00   0.00   0.00               
## trt.f1:timepnt.f2  0.00   0.00   0.00   0.00   0.00   0.00  -0.25        
## trt.f1:timepnt.f3  0.00   0.00   0.00   0.00   0.00   0.00  -0.25  -0.25 
## trt.f1:timepnt.f4  0.00   0.00   0.00   0.00   0.00   0.00  -0.25  -0.25 
##                   t.1:.3
## trt.f1                  
## timepnt.f1              
## timepnt.f2              
## timepnt.f3              
## timepnt.f4              
## trt.f1:timepnt.f1       
## trt.f1:timepnt.f2       
## trt.f1:timepnt.f3       
## trt.f1:timepnt.f4 -0.25 
## 
## Standardized Within-Group Residuals:
##         Min          Q1         Med          Q3         Max 
## -3.21774790 -0.51423747 -0.04419265  0.71213362  2.31966939 
## 
## Number of Observations: 170
## Number of Groups: 34
\end{verbatim}

\begin{Shaded}
\begin{Highlighting}[]
\FunctionTok{anova}\NormalTok{(fit.cs1)}
\end{Highlighting}
\end{Shaded}

\begin{verbatim}
##                 numDF denDF  F-value p-value
## (Intercept)         1   128 372.8600  <.0001
## trt.f               1    32   5.7412  0.0226
## timepnt.f           4   128  26.6973  <.0001
## trt.f:timepnt.f     4   128   4.1669  0.0033
\end{verbatim}

\begin{Shaded}
\begin{Highlighting}[]
\FunctionTok{AIC}\NormalTok{(fit.cs1)}
\end{Highlighting}
\end{Shaded}

\begin{verbatim}
## [1] 798.2436
\end{verbatim}

\begin{Shaded}
\begin{Highlighting}[]
\FunctionTok{getVarCov}\NormalTok{(fit.cs1, }\AttributeTok{individual=}\StringTok{"1"}\NormalTok{,}\AttributeTok{type=}\StringTok{"marginal"}\NormalTok{)}
\end{Highlighting}
\end{Shaded}

\begin{verbatim}
## id 1 
## Marginal variance covariance matrix
##        1      2      3      4      5
## 1 15.753 12.718 12.718 12.718 12.718
## 2 12.718 15.753 12.718 12.718 12.718
## 3 12.718 12.718 15.753 12.718 12.718
## 4 12.718 12.718 12.718 15.753 12.718
## 5 12.718 12.718 12.718 12.718 15.753
##   Standard Deviations: 3.969 3.969 3.969 3.969 3.969
\end{verbatim}

\begin{Shaded}
\begin{Highlighting}[]
\DocumentationTok{\#\# Unstructured Covariance Structure }
\NormalTok{fit.un }\OtherTok{\textless{}{-}} \FunctionTok{lme}\NormalTok{(y }\SpecialCharTok{\textasciitilde{}}\NormalTok{ trt.f}\SpecialCharTok{*}\NormalTok{timepnt.f, }\AttributeTok{data=}\NormalTok{caz,}
  \AttributeTok{random=} \SpecialCharTok{\textasciitilde{}}\DecValTok{1} \SpecialCharTok{|}\NormalTok{ id1.f,}
  \AttributeTok{corr =} \FunctionTok{corSymm}\NormalTok{(}\AttributeTok{form=} \SpecialCharTok{\textasciitilde{}} \DecValTok{1} \SpecialCharTok{|}\NormalTok{ id1.f),}
  \AttributeTok{weights =} \FunctionTok{varIdent}\NormalTok{(}\AttributeTok{form =} \SpecialCharTok{\textasciitilde{}} \DecValTok{1} \SpecialCharTok{|}\NormalTok{ timepnt.f), }
  \AttributeTok{control=}\FunctionTok{lmeControl}\NormalTok{(}\AttributeTok{maxIter=}\DecValTok{100}\NormalTok{, }\AttributeTok{opt=}\StringTok{"optim"}\NormalTok{))}
\FunctionTok{summary}\NormalTok{(fit.un)}
\end{Highlighting}
\end{Shaded}

\begin{verbatim}
## Linear mixed-effects model fit by REML
##   Data: caz 
##        AIC      BIC    logLik
##   749.8307 829.7852 -348.9153
## 
## Random effects:
##  Formula: ~1 | id1.f
##         (Intercept) Residual
## StdDev:    3.073923 1.391268
## 
## Correlation Structure: General
##  Formula: ~1 | id1.f 
##  Parameter estimate(s):
##  Correlation: 
##   1      2      3      4     
## 2  0.031                     
## 3 -0.011  0.621              
## 4  0.219  0.697  0.886       
## 5 -0.041  0.550  0.805  0.929
## Variance function:
##  Structure: Different standard deviations per stratum
##  Formula: ~1 | timepnt.f 
##  Parameter estimates:
##        1        2        3        4        5 
## 1.000000 1.139761 1.712557 2.374562 2.335903 
## Fixed effects:  y ~ trt.f * timepnt.f 
##                       Value Std.Error  DF   t-value p-value
## (Intercept)       12.088235 0.6264483 128 19.296463  0.0000
## trt.f1            -1.500000 0.6264483  32 -2.394452  0.0227
## timepnt.f1        -2.088235 0.3723377 128 -5.608445  0.0000
## timepnt.f2        -0.941176 0.2328305 128 -4.042325  0.0001
## timepnt.f3         0.147059 0.1790910 128  0.821140  0.4131
## timepnt.f4         1.205882 0.2300381 128  5.242098  0.0000
## trt.f1:timepnt.f1  0.558824 0.3723377 128  1.500851  0.1359
## trt.f1:timepnt.f2  0.529412 0.2328305 128  2.273808  0.0246
## trt.f1:timepnt.f3  0.147059 0.1790910 128  0.821140  0.4131
## trt.f1:timepnt.f4 -0.382353 0.2300381 128 -1.662129  0.0989
##  Correlation: 
##                   (Intr) trt.f1 tmpn.1 tmpn.2 tmpn.3 tmpn.4 t.1:.1 t.1:.2
## trt.f1             0.000                                                 
## timepnt.f1        -0.421  0.000                                          
## timepnt.f2        -0.325  0.000  0.380                                   
## timepnt.f3         0.091  0.000 -0.415 -0.198                            
## timepnt.f4         0.533  0.000 -0.742 -0.683  0.090                     
## trt.f1:timepnt.f1  0.000 -0.421  0.000  0.000  0.000  0.000              
## trt.f1:timepnt.f2  0.000 -0.325  0.000  0.000  0.000  0.000  0.380       
## trt.f1:timepnt.f3  0.000  0.091  0.000  0.000  0.000  0.000 -0.415 -0.198
## trt.f1:timepnt.f4  0.000  0.533  0.000  0.000  0.000  0.000 -0.742 -0.683
##                   t.1:.3
## trt.f1                  
## timepnt.f1              
## timepnt.f2              
## timepnt.f3              
## timepnt.f4              
## trt.f1:timepnt.f1       
## trt.f1:timepnt.f2       
## trt.f1:timepnt.f3       
## trt.f1:timepnt.f4  0.090
## 
## Standardized Within-Group Residuals:
##        Min         Q1        Med         Q3        Max 
## -2.3359672 -0.6444285 -0.1169039  0.6566682  2.7544055 
## 
## Number of Observations: 170
## Number of Groups: 34
\end{verbatim}

\begin{Shaded}
\begin{Highlighting}[]
\FunctionTok{anova}\NormalTok{(fit.un)}
\end{Highlighting}
\end{Shaded}

\begin{verbatim}
##                 numDF denDF  F-value p-value
## (Intercept)         1   128 385.9196  <.0001
## trt.f               1    32   3.9783  0.0547
## timepnt.f           4   128   9.7978  <.0001
## trt.f:timepnt.f     4   128   2.7241  0.0323
\end{verbatim}

\begin{Shaded}
\begin{Highlighting}[]
\FunctionTok{AIC}\NormalTok{(fit.un)}
\end{Highlighting}
\end{Shaded}

\begin{verbatim}
## [1] 749.8307
\end{verbatim}

\begin{Shaded}
\begin{Highlighting}[]
\FunctionTok{getVarCov}\NormalTok{(fit.un, }\AttributeTok{individual=}\StringTok{"1"}\NormalTok{,}\AttributeTok{type=}\StringTok{"conditional"}\NormalTok{)}
\end{Highlighting}
\end{Shaded}

\begin{verbatim}
## id1.f 1 
## Conditional variance covariance matrix
##           1        2         3       4        5
## 1  1.935600 0.067913 -0.037752  1.0086 -0.18467
## 2  0.067913 2.514500  2.348000  3.6492  2.83350
## 3 -0.037752 2.348000  5.676900  6.9764  6.23720
## 4  1.008600 3.649200  6.976400 10.9140  9.97360
## 5 -0.184670 2.833500  6.237200  9.9736 10.56200
##   Standard Deviations: 1.3913 1.5857 2.3826 3.3037 3.2499
\end{verbatim}

\begin{Shaded}
\begin{Highlighting}[]
\DocumentationTok{\#\# AR(1) Covariance Structure}
\NormalTok{fit.ar1 }\OtherTok{\textless{}{-}} \FunctionTok{lme}\NormalTok{(y }\SpecialCharTok{\textasciitilde{}}\NormalTok{ trt.f}\SpecialCharTok{*}\NormalTok{timepnt.f, }\AttributeTok{data=}\NormalTok{caz,}
\AttributeTok{random=} \SpecialCharTok{\textasciitilde{}}\DecValTok{1} \SpecialCharTok{|}\NormalTok{ id1.f,}
\AttributeTok{corr =} \FunctionTok{corAR1}\NormalTok{(, }\AttributeTok{form=} \SpecialCharTok{\textasciitilde{}} \DecValTok{1} \SpecialCharTok{|}\NormalTok{ id1.f))}
\FunctionTok{summary}\NormalTok{(fit.ar1)}
\end{Highlighting}
\end{Shaded}

\begin{verbatim}
## Linear mixed-effects model fit by REML
##   Data: caz 
##        AIC      BIC    logLik
##   744.0077 783.9849 -359.0038
## 
## Random effects:
##  Formula: ~1 | id1.f
##         (Intercept) Residual
## StdDev: 0.001180249 3.963392
## 
## Correlation Structure: AR(1)
##  Formula: ~1 | id1.f 
##  Parameter estimate(s):
##       Phi 
## 0.9047053 
## Fixed effects:  y ~ trt.f * timepnt.f 
##                       Value Std.Error  DF   t-value p-value
## (Intercept)       12.088235 0.6296157 128 19.199388  0.0000
## trt.f1            -1.500000 0.6296157  32 -2.382406  0.0233
## timepnt.f1        -2.088235 0.3074310 128 -6.792533  0.0000
## timepnt.f2        -0.941176 0.2261102 128 -4.162469  0.0001
## timepnt.f3         0.147059 0.1916021 128  0.767522  0.4442
## timepnt.f4         1.205882 0.2261102 128  5.333163  0.0000
## trt.f1:timepnt.f1  0.558824 0.3074310 128  1.817720  0.0714
## trt.f1:timepnt.f2  0.529412 0.2261102 128  2.341389  0.0208
## trt.f1:timepnt.f3  0.147059 0.1916021 128  0.767522  0.4442
## trt.f1:timepnt.f4 -0.382353 0.2261102 128 -1.691003  0.0933
##  Correlation: 
##                   (Intr) trt.f1 tmpn.1 tmpn.2 tmpn.3 tmpn.4 t.1:.1 t.1:.2
## trt.f1             0.000                                                 
## timepnt.f1        -0.075  0.000                                          
## timepnt.f2         0.051  0.000  0.414                                   
## timepnt.f3         0.120  0.000 -0.310 -0.003                            
## timepnt.f4         0.051  0.000 -0.677 -0.640 -0.003                     
## trt.f1:timepnt.f1  0.000 -0.075  0.000  0.000  0.000  0.000              
## trt.f1:timepnt.f2  0.000  0.051  0.000  0.000  0.000  0.000  0.414       
## trt.f1:timepnt.f3  0.000  0.120  0.000  0.000  0.000  0.000 -0.310 -0.003
## trt.f1:timepnt.f4  0.000  0.051  0.000  0.000  0.000  0.000 -0.677 -0.640
##                   t.1:.3
## trt.f1                  
## timepnt.f1              
## timepnt.f2              
## timepnt.f3              
## timepnt.f4              
## trt.f1:timepnt.f1       
## trt.f1:timepnt.f2       
## trt.f1:timepnt.f3       
## trt.f1:timepnt.f4 -0.003
## 
## Standardized Within-Group Residuals:
##         Min          Q1         Med          Q3         Max 
## -2.37467405 -0.60851025 -0.04452508  0.72724395  2.41919926 
## 
## Number of Observations: 170
## Number of Groups: 34
\end{verbatim}

\begin{Shaded}
\begin{Highlighting}[]
\FunctionTok{anova}\NormalTok{(fit.ar1)}
\end{Highlighting}
\end{Shaded}

\begin{verbatim}
##                 numDF denDF  F-value p-value
## (Intercept)         1   128 369.4090  <.0001
## trt.f               1    32   6.7869  0.0138
## timepnt.f           4   128  12.4477  <.0001
## trt.f:timepnt.f     4   128   2.0486  0.0914
\end{verbatim}

\begin{Shaded}
\begin{Highlighting}[]
\FunctionTok{AIC}\NormalTok{(fit.ar1)}
\end{Highlighting}
\end{Shaded}

\begin{verbatim}
## [1] 744.0077
\end{verbatim}

\begin{Shaded}
\begin{Highlighting}[]
\FunctionTok{getVarCov}\NormalTok{(fit.ar1, }\AttributeTok{individual=}\StringTok{"1"}\NormalTok{,}\AttributeTok{type=}\StringTok{"conditional"}\NormalTok{)}
\end{Highlighting}
\end{Shaded}

\begin{verbatim}
## id1.f 1 
## Conditional variance covariance matrix
##        1      2      3      4      5
## 1 15.708 14.212 12.857 11.632 10.524
## 2 14.212 15.708 14.212 12.857 11.632
## 3 12.857 14.212 15.708 14.212 12.857
## 4 11.632 12.857 14.212 15.708 14.212
## 5 10.524 11.632 12.857 14.212 15.708
##   Standard Deviations: 3.9634 3.9634 3.9634 3.9634 3.9634
\end{verbatim}

The AR(1) model allows the covariance of the observations within cats to
decrease as they are farther apart in time. Based on having the smallest
\(AIC\) value, it gives the best fit of the three models without having
excess variance parameters.

\[
\hat{\mbox{COV}}\left\{Y_{ijk},Y_{ijk'}\right\} = 15.708(0.905)^{|k-k'|}
\]

\chapter{Analysis of Covariance
(ANCOVA)}\label{analysis-of-covariance-ancova}

In various cases, researchers are interested in comparing treatments
after adjusting for a numeric predictor that is associated with an
experimental unit. That is often, but not necessarily, a baseline
pre-treatment score on the unit. The goal is to determine whether
treatment effects are present after controlling for the covariate.

When the Analysis of Covariance was first developed, the computations
used were very complicated. Now, it's just as simple as fitting a
regression with one or more predictors and dummy variables for the
treatment conditions. Models can have one or more factors. We will
consider the case with one treatment factor and one covariate.

\section{Additive Model}\label{additive-model}

In the additive model, the slope relating the covariate to the response
is assumed to be the same for each treatment. As with the 1-Way ANOVA
model, we will use \(r\) as the number of treatments with \(n_i\)
replicates for treatment \(i\). We will use centered \(X\) values for
the covariate for ease of interpretation when describing adjusted means.
We will also include \(r-1\) indicators for treatments \(1,\ldots,r-1\)
assuming again that treatment effects sum to 0.

\[
Y_{ij}= \mu_{\bullet} + \tau_1 I_{ij1} + \cdots +
\tau_{r-1}I_{ij,r-1} + 
\gamma \left(X_{ij}-\overline{X}_{\bullet\bullet}\right) +\epsilon_{ijk}
\]

\[
\epsilon_{ijk} \sim NID\left(0,\sigma^2\right) \qquad
\sum_{i=1}^r\tau_i=0 \quad \Rightarrow \quad \tau_r=-\sum_{i=1}^{r-1}\tau_i
\] \[
I_{ij1} = 
\left\{ \begin{array}{r@{\quad:\quad}l}
    1 & i=1 \\
    0 & i=2,\ldots,r-1 \\
    -1 & i=r\\
    \end{array} \right.
  \qquad \cdots
   \qquad
I_{ij,r-1} = 
\left\{ \begin{array}{r@{\quad:\quad}l}
    1 & i=r-1 \\
    0 & i=1,\ldots,r-2 \\
    -1 & i=r\\
     \end{array} \right.
\]

To test for treatment effects, controlling for the covariate, we test
\(H_0:\tau_1=\cdots = \tau_r=0\). This is easily conducted as a general
linear test, where the full model contains the \(r-1\) treatment
indicator variables and the centered covariate. The reduced model
contains only the centered covariate. The test is as follows.

\[
TS: F^*=\frac{\left[\frac{SSE(R)-SSE(C)}{df_R-df_C}\right]}{
\left[\frac{SSE(C)}{df_C}\right]
}
\qquad RR:F^*
\geq F_{1-\alpha;df_R-df_F,df_F}
\] Note that for the current case, \(df_R=n_T-2\) and
\(df_F=n_T-2-(r-1)=n_T-4\).

\textbf{Example 12.1 - Comparison of Skin Softeners}

A study compared two treatments and a control \((r=3)\) for skin
softening (\citeproc{ref-maor_1997}{Ma'Or, Yehuda, and Voss 1997}). The
active treatments were as follows, and the covariate was the
pre-treatment softness score for the subjects. There were 20 subjects
per treatment, so \(n_T=3(20)=60\). The response is skin roughness
(lower scores are better).

\begin{itemize}
\tightlist
\item
  Treatment 1 - Gel Formulation
\item
  Treatment 2 - Gel Formulation + 1\% Dead Sea Concentrate
\item
  Treatment 3 - Placebo
\end{itemize}

The R code and output are given below.

\begin{Shaded}
\begin{Highlighting}[]
\NormalTok{dsm }\OtherTok{\textless{}{-}} \FunctionTok{read.csv}\NormalTok{(}\StringTok{"data/deadseaminerals.csv"}\NormalTok{)}
\FunctionTok{head}\NormalTok{(dsm,}\DecValTok{3}\NormalTok{); }\FunctionTok{tail}\NormalTok{(dsm,}\DecValTok{3}\NormalTok{)}
\end{Highlighting}
\end{Shaded}

\begin{verbatim}
##   trt_Grp  pre_x post_y gel gelDS
## 1       1 121.10  76.30   1     0
## 2       1 375.40 265.10   1     0
## 3       1 138.05  82.92   1     0
\end{verbatim}

\begin{verbatim}
##    trt_Grp  pre_x post_y gel gelDS
## 58       3 138.39 132.16   0     0
## 59       3 138.59 107.22   0     0
## 60       3 134.89 125.20   0     0
\end{verbatim}

\begin{Shaded}
\begin{Highlighting}[]
\NormalTok{dsm}\SpecialCharTok{$}\NormalTok{I1 }\OtherTok{\textless{}{-}} \FunctionTok{ifelse}\NormalTok{(dsm}\SpecialCharTok{$}\NormalTok{trt\_Grp}\SpecialCharTok{==}\DecValTok{1}\NormalTok{,}\DecValTok{1}\NormalTok{,}\FunctionTok{ifelse}\NormalTok{(dsm}\SpecialCharTok{$}\NormalTok{trt\_Grp}\SpecialCharTok{==}\DecValTok{3}\NormalTok{,}\SpecialCharTok{{-}}\DecValTok{1}\NormalTok{,}\DecValTok{0}\NormalTok{))}
\NormalTok{dsm}\SpecialCharTok{$}\NormalTok{I2 }\OtherTok{\textless{}{-}} \FunctionTok{ifelse}\NormalTok{(dsm}\SpecialCharTok{$}\NormalTok{trt\_Grp}\SpecialCharTok{==}\DecValTok{2}\NormalTok{,}\DecValTok{1}\NormalTok{,}\FunctionTok{ifelse}\NormalTok{(dsm}\SpecialCharTok{$}\NormalTok{trt\_Grp}\SpecialCharTok{==}\DecValTok{3}\NormalTok{,}\SpecialCharTok{{-}}\DecValTok{1}\NormalTok{,}\DecValTok{0}\NormalTok{))}
\NormalTok{dsm}\SpecialCharTok{$}\NormalTok{Xc }\OtherTok{\textless{}{-}}\NormalTok{ dsm}\SpecialCharTok{$}\NormalTok{pre\_x }\SpecialCharTok{{-}} \FunctionTok{mean}\NormalTok{(dsm}\SpecialCharTok{$}\NormalTok{pre\_x)}

\NormalTok{dsm.mod1 }\OtherTok{\textless{}{-}} \FunctionTok{lm}\NormalTok{(post\_y }\SpecialCharTok{\textasciitilde{}}\NormalTok{ I1 }\SpecialCharTok{+}\NormalTok{ I2 }\SpecialCharTok{+}\NormalTok{ Xc, }\AttributeTok{data=}\NormalTok{dsm)}
\FunctionTok{summary}\NormalTok{(dsm.mod1)}
\end{Highlighting}
\end{Shaded}

\begin{verbatim}
## 
## Call:
## lm(formula = post_y ~ I1 + I2 + Xc, data = dsm)
## 
## Residuals:
##    Min     1Q Median     3Q    Max 
## -31.13 -10.35  -3.93  12.29  35.34 
## 
## Coefficients:
##              Estimate Std. Error t value Pr(>|t|)    
## (Intercept) 137.53667    2.08924  65.831  < 2e-16 ***
## I1           -2.66382    2.95473  -0.902    0.371    
## I2          -26.91735    2.95595  -9.106 1.22e-12 ***
## Xc            0.70732    0.03257  21.716  < 2e-16 ***
## ---
## Signif. codes:  0 '***' 0.001 '**' 0.01 '*' 0.05 '.' 0.1 ' ' 1
## 
## Residual standard error: 16.18 on 56 degrees of freedom
## Multiple R-squared:  0.9121, Adjusted R-squared:  0.9073 
## F-statistic: 193.6 on 3 and 56 DF,  p-value: < 2.2e-16
\end{verbatim}

\begin{Shaded}
\begin{Highlighting}[]
\FunctionTok{anova}\NormalTok{(dsm.mod1)}
\end{Highlighting}
\end{Shaded}

\begin{verbatim}
## Analysis of Variance Table
## 
## Response: post_y
##           Df Sum Sq Mean Sq F value    Pr(>F)    
## I1         1   9853    9853  37.622 9.205e-08 ***
## I2         1  18744   18744  71.572 1.361e-11 ***
## Xc         1 123501  123501 471.566 < 2.2e-16 ***
## Residuals 56  14666     262                      
## ---
## Signif. codes:  0 '***' 0.001 '**' 0.01 '*' 0.05 '.' 0.1 ' ' 1
\end{verbatim}

\begin{Shaded}
\begin{Highlighting}[]
\NormalTok{dsm.mod2 }\OtherTok{\textless{}{-}} \FunctionTok{lm}\NormalTok{(post\_y }\SpecialCharTok{\textasciitilde{}}\NormalTok{ Xc, }\AttributeTok{data=}\NormalTok{dsm)}
\FunctionTok{summary}\NormalTok{(dsm.mod2)}
\end{Highlighting}
\end{Shaded}

\begin{verbatim}
## 
## Call:
## lm(formula = post_y ~ Xc, data = dsm)
## 
## Residuals:
##     Min      1Q  Median      3Q     Max 
## -50.220 -19.668  -4.885  21.923  63.420 
## 
## Coefficients:
##              Estimate Std. Error t value Pr(>|t|)    
## (Intercept) 137.53667    3.66612   37.52   <2e-16 ***
## Xc            0.69686    0.05713   12.20   <2e-16 ***
## ---
## Signif. codes:  0 '***' 0.001 '**' 0.01 '*' 0.05 '.' 0.1 ' ' 1
## 
## Residual standard error: 28.4 on 58 degrees of freedom
## Multiple R-squared:  0.7195, Adjusted R-squared:  0.7147 
## F-statistic: 148.8 on 1 and 58 DF,  p-value: < 2.2e-16
\end{verbatim}

\begin{Shaded}
\begin{Highlighting}[]
\FunctionTok{anova}\NormalTok{(dsm.mod2)}
\end{Highlighting}
\end{Shaded}

\begin{verbatim}
## Analysis of Variance Table
## 
## Response: post_y
##           Df Sum Sq Mean Sq F value    Pr(>F)    
## Xc         1 119991  119991  148.79 < 2.2e-16 ***
## Residuals 58  46773     806                      
## ---
## Signif. codes:  0 '***' 0.001 '**' 0.01 '*' 0.05 '.' 0.1 ' ' 1
\end{verbatim}

\begin{Shaded}
\begin{Highlighting}[]
\FunctionTok{anova}\NormalTok{(dsm.mod2, dsm.mod1)}
\end{Highlighting}
\end{Shaded}

\begin{verbatim}
## Analysis of Variance Table
## 
## Model 1: post_y ~ Xc
## Model 2: post_y ~ I1 + I2 + Xc
##   Res.Df   RSS Df Sum of Sq      F    Pr(>F)    
## 1     58 46773                                  
## 2     56 14666  2     32107 61.297 7.888e-15 ***
## ---
## Signif. codes:  0 '***' 0.001 '**' 0.01 '*' 0.05 '.' 0.1 ' ' 1
\end{verbatim}

\begin{Shaded}
\begin{Highlighting}[]
\FunctionTok{summary}\NormalTok{(dsm}\SpecialCharTok{$}\NormalTok{Xc)}
\end{Highlighting}
\end{Shaded}

\begin{verbatim}
##    Min. 1st Qu.  Median    Mean 3rd Qu.    Max. 
## -150.41  -48.35  -19.85    0.00   42.88  188.51
\end{verbatim}

\begin{Shaded}
\begin{Highlighting}[]
\NormalTok{Xc.grid }\OtherTok{\textless{}{-}} \FunctionTok{seq}\NormalTok{(}\SpecialCharTok{{-}}\DecValTok{200}\NormalTok{,}\DecValTok{200}\NormalTok{,}\FloatTok{0.1}\NormalTok{)}
\FunctionTok{plot}\NormalTok{(dsm}\SpecialCharTok{$}\NormalTok{post\_y }\SpecialCharTok{\textasciitilde{}}\NormalTok{ dsm}\SpecialCharTok{$}\NormalTok{Xc, }\AttributeTok{pch=}\DecValTok{16}\NormalTok{, }\AttributeTok{col=}\NormalTok{dsm}\SpecialCharTok{$}\NormalTok{trt\_Grp, }\AttributeTok{cex=}\FloatTok{0.6}\NormalTok{,}
     \AttributeTok{xlab=}\StringTok{"Pre{-}Treatment Score"}\NormalTok{, }\AttributeTok{ylab=}\StringTok{"Post{-}Treatment Score"}\NormalTok{,}
     \AttributeTok{xlim=}\FunctionTok{c}\NormalTok{(}\SpecialCharTok{{-}}\DecValTok{200}\NormalTok{,}\DecValTok{200}\NormalTok{), }\AttributeTok{ylim=}\FunctionTok{c}\NormalTok{(}\DecValTok{0}\NormalTok{,}\DecValTok{320}\NormalTok{), }\AttributeTok{main=}\StringTok{"Skin Roughness"}\NormalTok{)}
\FunctionTok{lines}\NormalTok{(Xc.grid, }
\NormalTok{      dsm.mod1}\SpecialCharTok{$}\NormalTok{coef[}\DecValTok{1}\NormalTok{]}\SpecialCharTok{+}\NormalTok{dsm.mod1}\SpecialCharTok{$}\NormalTok{coef[}\DecValTok{2}\NormalTok{]}\SpecialCharTok{+}\NormalTok{Xc.grid}\SpecialCharTok{*}\NormalTok{dsm.mod1}\SpecialCharTok{$}\NormalTok{coef[}\DecValTok{4}\NormalTok{],}
      \AttributeTok{col=}\DecValTok{1}\NormalTok{)}
\FunctionTok{lines}\NormalTok{(Xc.grid, }
\NormalTok{      dsm.mod1}\SpecialCharTok{$}\NormalTok{coef[}\DecValTok{1}\NormalTok{]}\SpecialCharTok{+}\NormalTok{dsm.mod1}\SpecialCharTok{$}\NormalTok{coef[}\DecValTok{3}\NormalTok{]}\SpecialCharTok{+}\NormalTok{Xc.grid}\SpecialCharTok{*}\NormalTok{dsm.mod1}\SpecialCharTok{$}\NormalTok{coef[}\DecValTok{4}\NormalTok{],}
      \AttributeTok{col=}\DecValTok{2}\NormalTok{)}
\FunctionTok{lines}\NormalTok{(Xc.grid, }
\NormalTok{      dsm.mod1}\SpecialCharTok{$}\NormalTok{coef[}\DecValTok{1}\NormalTok{]}\SpecialCharTok{{-}}\NormalTok{dsm.mod1}\SpecialCharTok{$}\NormalTok{coef[}\DecValTok{2}\NormalTok{]}\SpecialCharTok{{-}}\NormalTok{dsm.mod1}\SpecialCharTok{$}\NormalTok{coef[}\DecValTok{3}\NormalTok{]}\SpecialCharTok{+}\NormalTok{Xc.grid}\SpecialCharTok{*}\NormalTok{dsm.mod1}\SpecialCharTok{$}\NormalTok{coef[}\DecValTok{4}\NormalTok{],}
      \AttributeTok{col=}\DecValTok{3}\NormalTok{)}
\FunctionTok{legend}\NormalTok{(}\StringTok{"topleft"}\NormalTok{, }\FunctionTok{c}\NormalTok{(}\StringTok{"Gel"}\NormalTok{, }\StringTok{"Gel+DeadSea"}\NormalTok{,}\StringTok{"Placebo"}\NormalTok{), }\AttributeTok{col=}\DecValTok{1}\SpecialCharTok{:}\DecValTok{3}\NormalTok{, }\AttributeTok{lty=}\DecValTok{1}\NormalTok{)}
\end{Highlighting}
\end{Shaded}

\begin{figure}

{\centering \includegraphics[width=0.8\linewidth]{bookdownproj_STA4211Notes1_files/figure-latex/ch1201-1} 

}

\caption{Analysis of Covariance for skin softener study - Additive Model}\label{fig:ch1201}
\end{figure}

The \(F^*\)-statistic is 61.297 with 2 and 56 degrees of freedom and
\(P<.0001\). There is strong evidence of differences among the
treatments, controlling for baseline score.

\[\nabla \]

\subsection{Adjusted means}\label{adjusted-means}

Where the lines cross at the centered X level of 0 are the
\textbf{adjusted means}. That is, they are the predicted scores for the
treatments at the ``average'' value of the uncentered covariate. Based
on this additive model, we obtain the following adjusted means.

\[
i=1,\ldots , r-1:\quad   \hat{\mu}_i = \hat{\mu}_{\bullet} +\hat{\tau}_i \qquad 
\hat{\mu}_r=\hat{\mu}_{\bullet} -\sum_{i=1}^{r-1}\hat{\tau}_i
\]

The goal is to typically compare the adjusted means, as would be done
for the 1-Way ANOVA. The standard errors depend on the
variance-covariance matrix of the regression coefficients, which is
obtained with the \textbf{vcov} function from the regression model fit.
The variance-covariance matrix for this example is given below.

\[
\begin{bmatrix}
\hat{V}\left\{\hat{\mu}_{\bullet}\right\} &
\hat{\mbox{COV}}\left\{\hat{\mu}_{\bullet},\hat{\tau}_1\right\} &
\hat{\mbox{COV}}\left\{\hat{\mu}_{\bullet},\hat{\tau}_2\right\} &
\hat{\mbox{COV}}\left\{\hat{\mu}_{\bullet},\hat{\gamma}\right\} \\
\hat{\mbox{COV}}\left\{\hat{\mu}_{\bullet},\hat{\tau}_1\right\} &
\hat{V}\left\{\hat{\tau}_1\right\} &
\hat{\mbox{COV}}\left\{\hat{\tau}_1,\hat{\tau}_2\right\} &
\hat{\mbox{COV}}\left\{\hat{\tau}_1,\hat{\gamma}\right\} \\
\hat{\mbox{COV}}\left\{\hat{\mu}_{\bullet},\hat{\tau}_2\right\} &
\hat{\mbox{COV}}\left\{\hat{\tau}_1,\hat{\tau}_2\right\} &
\hat{V}\left\{\hat{\tau}_2\right\} &
\hat{\mbox{COV}}\left\{\hat{\tau}_2,\hat{\gamma}\right\} \\
\hat{\mbox{COV}}\left\{\hat{\mu}_{\bullet},\hat{\gamma}\right\} &
\hat{\mbox{COV}}\left\{\hat{\tau}_1,\hat{\gamma}\right\} &
\hat{\mbox{COV}}\left\{\hat{\tau}_2,\hat{\gamma}\right\} &
\hat{V}\left\{\hat{\gamma}\right\}
\end{bmatrix}
\]

\[
i,i' < r  \quad
\hat{SE}\left\{\hat{\tau}_i-\hat{\tau}_{i'}\right\} =
\sqrt{\hat{V}\left\{\hat{\tau}_i\right\}+\hat{V}\left\{\hat{\tau}_{i'}\right\}
-2\hat{\mbox{COV}}\left\{\hat{\tau}_i,\hat{\tau}_{i'}\right\}}
\]

When comparing the first \(r-1\) treatments with the \(r^{th}\), we use
the following result.

\[
\hat{\tau}_i-\hat{\tau}_r=\hat{\tau}_i-\left(-\sum_{i=1}^{r-1}\hat{\tau}_i\right)=
2\hat{\tau}_i + \sum_{\stackrel{i'=1}{i'\ne i}}^{r-1}\hat{\tau}_{i'}
\]

Then the (messy) standard errors for these differences are as follow.

\[
\hat{SE}\left\{\hat{\tau}_i-\hat{\tau}_r\right\} =
\sqrt{4\hat{V}\left\{\hat{\tau}_i\right\} +
\sum_{\stackrel{i'=1}{i'\ne i}}^{r-1}\hat{V}\left\{\hat{\tau}_{i'}\right\}
+4\sum_{\stackrel{i'=1}{i'\ne i}}^{r-1}\hat{\mbox{COV}}\left\{\hat{\tau}_i,\hat{\tau}_{i'}\right\}
+2\sum_{\stackrel{i'=1}{i'\ne i}}^{r-2}\sum_{\stackrel{i''=i'+1}{i''\ne i}}^{r-1}
\hat{\mbox{COV}}\left\{\hat{\tau}_{i'},\hat{\tau}_{i''}\right\}
}
\] Here, we consider a few cases to observe these variance calculations.
When \(r=2\), we have the following simple case.

\[
\hat{\tau}_2 = -\hat{\tau}_1 \quad \Rightarrow \quad 
\hat{\tau}_1-\hat{\tau}_2=
2\hat{\tau}_1 
\quad \Rightarrow \quad 
\hat{V}\left\{\hat{\tau}_1-\hat{\tau}_2\right\}
=4\hat{V}\left\{\hat{\tau}_1\right\}
\] When \(r=3\), as in the skin softener study, we have the following
results.

\[
\hat{\tau}_3 = -\hat{\tau}_1-\hat{\tau}_2 \quad \Rightarrow \quad
\hat{\tau}_1-\hat{\tau}_3 = \hat{\tau}_1-\left[\hat{\tau}_1-\hat{\tau}_2\right]=
2\hat{\tau}_1+\hat{\tau}_2 
\] \[\quad \Rightarrow \quad
\hat{V}\left\{\hat{\tau}_1-\hat{\tau}_3\right\}=
4\hat{V}\left\{\hat{\tau}_1\right\}+\hat{V}\left\{\hat{\tau}_2\right\}+
4\hat{\mbox{COV}}\left\{\hat{\tau}_1,\hat{\tau}_2\right\}
\] Finally, when \(r=4\), have the following results. \[
\hat{\tau}_1-\hat{\tau}_4 = 
\hat{\tau}_1-\left[\hat{\tau}_1-\hat{\tau}_2-\hat{\tau}_3\right]=
2\hat{\tau}_1+\hat{\tau}_2+\hat{\tau}_3 \quad \Rightarrow \quad  
\] \[
\hat{V}\left\{\hat{\tau}_1-\hat{\tau}_4\right\}=
4\hat{V}\left\{\hat{\tau}_1\right\}+\hat{V}\left\{\hat{\tau}_2\right\}+
\hat{V}\left\{\hat{\tau}_3\right\}+
4\hat{\mbox{COV}}\left\{\hat{\tau}_1,\hat{\tau}_2\right\}+
4\hat{\mbox{COV}}\left\{\hat{\tau}_1,\hat{\tau}_3\right\}+
2\hat{\mbox{COV}}\left\{\hat{\tau}_2,\hat{\tau}_3\right\}
\]

\textbf{Example 12.2 - Comparison of Skin Softeners}

We obtain the adjusted means and their differences along with multiple
comparisons with the following R code.

\begin{Shaded}
\begin{Highlighting}[]
\NormalTok{dsm }\OtherTok{\textless{}{-}} \FunctionTok{read.csv}\NormalTok{(}\StringTok{"data/deadseaminerals.csv"}\NormalTok{)}
\FunctionTok{head}\NormalTok{(dsm,}\DecValTok{2}\NormalTok{); }\FunctionTok{tail}\NormalTok{(dsm,}\DecValTok{2}\NormalTok{)}
\end{Highlighting}
\end{Shaded}

\begin{verbatim}
##   trt_Grp pre_x post_y gel gelDS
## 1       1 121.1   76.3   1     0
## 2       1 375.4  265.1   1     0
\end{verbatim}

\begin{verbatim}
##    trt_Grp  pre_x post_y gel gelDS
## 59       3 138.59 107.22   0     0
## 60       3 134.89 125.20   0     0
\end{verbatim}

\begin{Shaded}
\begin{Highlighting}[]
\NormalTok{dsm}\SpecialCharTok{$}\NormalTok{I1 }\OtherTok{\textless{}{-}} \FunctionTok{ifelse}\NormalTok{(dsm}\SpecialCharTok{$}\NormalTok{trt\_Grp}\SpecialCharTok{==}\DecValTok{1}\NormalTok{,}\DecValTok{1}\NormalTok{,}\FunctionTok{ifelse}\NormalTok{(dsm}\SpecialCharTok{$}\NormalTok{trt\_Grp}\SpecialCharTok{==}\DecValTok{3}\NormalTok{,}\SpecialCharTok{{-}}\DecValTok{1}\NormalTok{,}\DecValTok{0}\NormalTok{))}
\NormalTok{dsm}\SpecialCharTok{$}\NormalTok{I2 }\OtherTok{\textless{}{-}} \FunctionTok{ifelse}\NormalTok{(dsm}\SpecialCharTok{$}\NormalTok{trt\_Grp}\SpecialCharTok{==}\DecValTok{2}\NormalTok{,}\DecValTok{1}\NormalTok{,}\FunctionTok{ifelse}\NormalTok{(dsm}\SpecialCharTok{$}\NormalTok{trt\_Grp}\SpecialCharTok{==}\DecValTok{3}\NormalTok{,}\SpecialCharTok{{-}}\DecValTok{1}\NormalTok{,}\DecValTok{0}\NormalTok{))}
\NormalTok{dsm}\SpecialCharTok{$}\NormalTok{Xc }\OtherTok{\textless{}{-}}\NormalTok{ dsm}\SpecialCharTok{$}\NormalTok{pre\_x }\SpecialCharTok{{-}} \FunctionTok{mean}\NormalTok{(dsm}\SpecialCharTok{$}\NormalTok{pre\_x)}

\NormalTok{dsm.mod1 }\OtherTok{\textless{}{-}} \FunctionTok{lm}\NormalTok{(post\_y }\SpecialCharTok{\textasciitilde{}}\NormalTok{ I1 }\SpecialCharTok{+}\NormalTok{ I2 }\SpecialCharTok{+}\NormalTok{ Xc, }\AttributeTok{data=}\NormalTok{dsm)}
\FunctionTok{summary}\NormalTok{(dsm.mod1)}
\end{Highlighting}
\end{Shaded}

\begin{verbatim}
## 
## Call:
## lm(formula = post_y ~ I1 + I2 + Xc, data = dsm)
## 
## Residuals:
##    Min     1Q Median     3Q    Max 
## -31.13 -10.35  -3.93  12.29  35.34 
## 
## Coefficients:
##              Estimate Std. Error t value Pr(>|t|)    
## (Intercept) 137.53667    2.08924  65.831  < 2e-16 ***
## I1           -2.66382    2.95473  -0.902    0.371    
## I2          -26.91735    2.95595  -9.106 1.22e-12 ***
## Xc            0.70732    0.03257  21.716  < 2e-16 ***
## ---
## Signif. codes:  0 '***' 0.001 '**' 0.01 '*' 0.05 '.' 0.1 ' ' 1
## 
## Residual standard error: 16.18 on 56 degrees of freedom
## Multiple R-squared:  0.9121, Adjusted R-squared:  0.9073 
## F-statistic: 193.6 on 3 and 56 DF,  p-value: < 2.2e-16
\end{verbatim}

\begin{Shaded}
\begin{Highlighting}[]
\FunctionTok{anova}\NormalTok{(dsm.mod1)}
\end{Highlighting}
\end{Shaded}

\begin{verbatim}
## Analysis of Variance Table
## 
## Response: post_y
##           Df Sum Sq Mean Sq F value    Pr(>F)    
## I1         1   9853    9853  37.622 9.205e-08 ***
## I2         1  18744   18744  71.572 1.361e-11 ***
## Xc         1 123501  123501 471.566 < 2.2e-16 ***
## Residuals 56  14666     262                      
## ---
## Signif. codes:  0 '***' 0.001 '**' 0.01 '*' 0.05 '.' 0.1 ' ' 1
\end{verbatim}

\begin{Shaded}
\begin{Highlighting}[]
\NormalTok{vcov1 }\OtherTok{\textless{}{-}} \FunctionTok{vcov}\NormalTok{(dsm.mod1)}
\NormalTok{r }\OtherTok{\textless{}{-}} \FunctionTok{length}\NormalTok{(}\FunctionTok{unique}\NormalTok{(dsm}\SpecialCharTok{$}\NormalTok{trt\_Grp))}

\NormalTok{adjmean }\OtherTok{\textless{}{-}} \FunctionTok{rep}\NormalTok{(}\DecValTok{0}\NormalTok{,r)}
\ControlFlowTok{for}\NormalTok{ (i1 }\ControlFlowTok{in} \DecValTok{1}\SpecialCharTok{:}\NormalTok{(r}\DecValTok{{-}1}\NormalTok{)) adjmean[i1] }\OtherTok{\textless{}{-}}\NormalTok{ dsm.mod1}\SpecialCharTok{$}\NormalTok{coef[}\DecValTok{1}\NormalTok{] }\SpecialCharTok{+}\NormalTok{ dsm.mod1}\SpecialCharTok{$}\NormalTok{coef[i1}\SpecialCharTok{+}\DecValTok{1}\NormalTok{]}
\NormalTok{adjmean[r] }\OtherTok{\textless{}{-}}\NormalTok{ dsm.mod1}\SpecialCharTok{$}\NormalTok{coef[}\DecValTok{1}\NormalTok{]}\SpecialCharTok{{-}}\FunctionTok{sum}\NormalTok{(dsm.mod1}\SpecialCharTok{$}\NormalTok{coef[}\DecValTok{2}\SpecialCharTok{:}\NormalTok{r])}
\NormalTok{adjmean}
\end{Highlighting}
\end{Shaded}

\begin{verbatim}
## [1] 134.8728 110.6193 167.1178
\end{verbatim}

\begin{Shaded}
\begin{Highlighting}[]
\NormalTok{se.diff }\OtherTok{\textless{}{-}} \FunctionTok{rep}\NormalTok{(}\DecValTok{0}\NormalTok{,r}\SpecialCharTok{*}\NormalTok{(r}\DecValTok{{-}1}\NormalTok{)}\SpecialCharTok{/}\DecValTok{2}\NormalTok{)}
\NormalTok{se.diff[}\DecValTok{1}\NormalTok{] }\OtherTok{\textless{}{-}} \FunctionTok{sqrt}\NormalTok{(vcov1[}\DecValTok{2}\NormalTok{,}\DecValTok{2}\NormalTok{]}\SpecialCharTok{+}\NormalTok{vcov1[}\DecValTok{3}\NormalTok{,}\DecValTok{3}\NormalTok{]}\SpecialCharTok{{-}}\DecValTok{2}\SpecialCharTok{*}\NormalTok{vcov1[}\DecValTok{2}\NormalTok{,}\DecValTok{3}\NormalTok{])}
\NormalTok{se.diff[}\DecValTok{2}\NormalTok{] }\OtherTok{\textless{}{-}} \FunctionTok{sqrt}\NormalTok{(}\DecValTok{4}\SpecialCharTok{*}\NormalTok{vcov1[}\DecValTok{2}\NormalTok{,}\DecValTok{2}\NormalTok{]}\SpecialCharTok{+}\NormalTok{vcov1[}\DecValTok{3}\NormalTok{,}\DecValTok{3}\NormalTok{]}\SpecialCharTok{+}\DecValTok{4}\SpecialCharTok{*}\NormalTok{vcov1[}\DecValTok{2}\NormalTok{,}\DecValTok{3}\NormalTok{])}
\NormalTok{se.diff[}\DecValTok{3}\NormalTok{] }\OtherTok{\textless{}{-}} \FunctionTok{sqrt}\NormalTok{(}\DecValTok{4}\SpecialCharTok{*}\NormalTok{vcov1[}\DecValTok{3}\NormalTok{,}\DecValTok{3}\NormalTok{]}\SpecialCharTok{+}\NormalTok{vcov1[}\DecValTok{2}\NormalTok{,}\DecValTok{2}\NormalTok{]}\SpecialCharTok{+}\DecValTok{4}\SpecialCharTok{*}\NormalTok{vcov1[}\DecValTok{2}\NormalTok{,}\DecValTok{3}\NormalTok{])}

\NormalTok{HSD.q }\OtherTok{\textless{}{-}} \FunctionTok{qtukey}\NormalTok{(.}\DecValTok{95}\NormalTok{,r,}\FunctionTok{df.residual}\NormalTok{(dsm.mod1)) }
\NormalTok{HSD.out }\OtherTok{\textless{}{-}} \FunctionTok{matrix}\NormalTok{(}\FunctionTok{rep}\NormalTok{(}\DecValTok{0}\NormalTok{,}\DecValTok{10}\SpecialCharTok{*}\NormalTok{r}\SpecialCharTok{*}\NormalTok{(r}\DecValTok{{-}1}\NormalTok{)}\SpecialCharTok{/}\DecValTok{2}\NormalTok{), }\AttributeTok{ncol=}\DecValTok{10}\NormalTok{)}
\NormalTok{row.count }\OtherTok{\textless{}{-}} \DecValTok{0}
\ControlFlowTok{for}\NormalTok{ (i1 }\ControlFlowTok{in} \DecValTok{2}\SpecialCharTok{:}\NormalTok{r) \{}
  \ControlFlowTok{for}\NormalTok{ (i2 }\ControlFlowTok{in} \DecValTok{1}\SpecialCharTok{:}\NormalTok{(i1}\DecValTok{{-}1}\NormalTok{)) \{}
\NormalTok{    row.count }\OtherTok{\textless{}{-}}\NormalTok{ row.count }\SpecialCharTok{+} \DecValTok{1}
\NormalTok{    HSD.out[row.count,}\DecValTok{1}\NormalTok{] }\OtherTok{\textless{}{-}}\NormalTok{ i1}
\NormalTok{    HSD.out[row.count,}\DecValTok{2}\NormalTok{] }\OtherTok{\textless{}{-}}\NormalTok{ i2}
\NormalTok{    HSD.out[row.count,}\DecValTok{3}\NormalTok{] }\OtherTok{\textless{}{-}}\NormalTok{ adjmean[i1]}
\NormalTok{    HSD.out[row.count,}\DecValTok{4}\NormalTok{] }\OtherTok{\textless{}{-}}\NormalTok{ adjmean[i2]}
\NormalTok{    HSD.out[row.count,}\DecValTok{5}\NormalTok{] }\OtherTok{\textless{}{-}}\NormalTok{ adjmean[i1] }\SpecialCharTok{{-}}\NormalTok{ adjmean[i2]}
\NormalTok{    HSD.out[row.count,}\DecValTok{6}\NormalTok{] }\OtherTok{\textless{}{-}}\NormalTok{ se.diff[row.count]}
\NormalTok{    HSD.out[row.count,}\DecValTok{7}\NormalTok{] }\OtherTok{\textless{}{-}}\NormalTok{ HSD.q}\SpecialCharTok{*}\NormalTok{se.diff[row.count]}
\NormalTok{    HSD.out[row.count,}\DecValTok{8}\NormalTok{] }\OtherTok{\textless{}{-}} 
\NormalTok{      (adjmean[i1] }\SpecialCharTok{{-}}\NormalTok{ adjmean[i2]) }\SpecialCharTok{{-}}\NormalTok{ HSD.q}\SpecialCharTok{*}\NormalTok{se.diff[row.count]}
\NormalTok{    HSD.out[row.count,}\DecValTok{9}\NormalTok{] }\OtherTok{\textless{}{-}} 
\NormalTok{      (adjmean[i1] }\SpecialCharTok{{-}}\NormalTok{ adjmean[i2]) }\SpecialCharTok{+}\NormalTok{ HSD.q}\SpecialCharTok{*}\NormalTok{se.diff[row.count]}
\NormalTok{    HSD.out[row.count,}\DecValTok{10}\NormalTok{] }\OtherTok{\textless{}{-}} 
      \DecValTok{1}\SpecialCharTok{{-}}\FunctionTok{ptukey}\NormalTok{(}\FunctionTok{sqrt}\NormalTok{(}\DecValTok{2}\NormalTok{)}\SpecialCharTok{*}\FunctionTok{abs}\NormalTok{(adjmean[i1] }\SpecialCharTok{{-}}\NormalTok{ adjmean[i2])}\SpecialCharTok{/}\NormalTok{se.diff[row.count],}
\NormalTok{               r,}\FunctionTok{df.residual}\NormalTok{(dsm.mod1))}
\NormalTok{  \}}
\NormalTok{\}}

\FunctionTok{colnames}\NormalTok{(HSD.out) }\OtherTok{\textless{}{-}} \FunctionTok{c}\NormalTok{(}\StringTok{"i"}\NormalTok{, }\StringTok{"i\textquotesingle{}"}\NormalTok{, }\StringTok{"AdjMn\_i"}\NormalTok{, }\StringTok{"AdjMn\_i\textquotesingle{}"}\NormalTok{, }\StringTok{"Diff"}\NormalTok{, }\StringTok{"Std Err"}\NormalTok{, }
                      \StringTok{"HSD"}\NormalTok{, }\StringTok{"LB"}\NormalTok{, }\StringTok{"UB"}\NormalTok{, }\StringTok{"p adj"}\NormalTok{)}
\FunctionTok{round}\NormalTok{(HSD.out,}\DecValTok{4}\NormalTok{)}
\end{Highlighting}
\end{Shaded}

\begin{verbatim}
##      i i'  AdjMn_i AdjMn_i'     Diff Std Err     HSD       LB      UB p adj
## [1,] 2  1 110.6193 134.8728 -24.2535  5.1188 17.4286 -41.6821 -6.8250     0
## [2,] 3  1 167.1178 134.8728  32.2450  5.1177 17.4248  14.8202 49.6698     0
## [3,] 3  2 167.1178 110.6193  56.4985  5.1198 17.4320  39.0665 73.9306     0
\end{verbatim}

All three pairs of treatments are significantly different. Treatment 2
(Gel+DeadSea) has significantly lower roughness scores than the other
treatments and Gel has significantly lower scores than Placebo.

Here we use the factor variable for treatment and the \textbf{emmeans}
package to obtain and compare the adjusted means.

\begin{Shaded}
\begin{Highlighting}[]
\NormalTok{dsm }\OtherTok{\textless{}{-}} \FunctionTok{read.csv}\NormalTok{(}\StringTok{"data/deadseaminerals.csv"}\NormalTok{)}
\FunctionTok{head}\NormalTok{(dsm,}\DecValTok{2}\NormalTok{)}
\end{Highlighting}
\end{Shaded}

\begin{verbatim}
##   trt_Grp pre_x post_y gel gelDS
## 1       1 121.1   76.3   1     0
## 2       1 375.4  265.1   1     0
\end{verbatim}

\begin{Shaded}
\begin{Highlighting}[]
\NormalTok{dsm}\SpecialCharTok{$}\NormalTok{trt\_Grp.f }\OtherTok{\textless{}{-}} \FunctionTok{factor}\NormalTok{(dsm}\SpecialCharTok{$}\NormalTok{trt\_Grp, }\AttributeTok{levels=}\DecValTok{1}\SpecialCharTok{:}\DecValTok{3}\NormalTok{,}\AttributeTok{labels=}\FunctionTok{c}\NormalTok{(}\StringTok{"gel"}\NormalTok{,}\StringTok{"gelDS"}\NormalTok{, }\StringTok{"control"}\NormalTok{))}
\FunctionTok{options}\NormalTok{(}\AttributeTok{contrasts=}\FunctionTok{c}\NormalTok{(}\StringTok{"contr.sum"}\NormalTok{, }\StringTok{"contr.poly"}\NormalTok{))}
\NormalTok{dsm.mod.add }\OtherTok{\textless{}{-}} \FunctionTok{aov}\NormalTok{(post\_y }\SpecialCharTok{\textasciitilde{}}\NormalTok{ pre\_x }\SpecialCharTok{+}\NormalTok{ trt\_Grp.f, }\AttributeTok{data=}\NormalTok{dsm)}
\FunctionTok{summary.lm}\NormalTok{(dsm.mod.add)}
\end{Highlighting}
\end{Shaded}

\begin{verbatim}
## 
## Call:
## aov(formula = post_y ~ pre_x + trt_Grp.f, data = dsm)
## 
## Residuals:
##    Min     1Q Median     3Q    Max 
## -31.13 -10.35  -3.93  12.29  35.34 
## 
## Coefficients:
##              Estimate Std. Error t value Pr(>|t|)    
## (Intercept)   5.34449    6.43598   0.830    0.410    
## pre_x         0.70732    0.03257  21.716  < 2e-16 ***
## trt_Grp.f1   -2.66382    2.95473  -0.902    0.371    
## trt_Grp.f2  -26.91735    2.95595  -9.106 1.22e-12 ***
## ---
## Signif. codes:  0 '***' 0.001 '**' 0.01 '*' 0.05 '.' 0.1 ' ' 1
## 
## Residual standard error: 16.18 on 56 degrees of freedom
## Multiple R-squared:  0.9121, Adjusted R-squared:  0.9073 
## F-statistic: 193.6 on 3 and 56 DF,  p-value: < 2.2e-16
\end{verbatim}

\begin{Shaded}
\begin{Highlighting}[]
\FunctionTok{anova}\NormalTok{(dsm.mod.add)}
\end{Highlighting}
\end{Shaded}

\begin{verbatim}
## Analysis of Variance Table
## 
## Response: post_y
##           Df Sum Sq Mean Sq F value    Pr(>F)    
## pre_x      1 119991  119991 458.166 < 2.2e-16 ***
## trt_Grp.f  2  32107   16053  61.297 7.888e-15 ***
## Residuals 56  14666     262                      
## ---
## Signif. codes:  0 '***' 0.001 '**' 0.01 '*' 0.05 '.' 0.1 ' ' 1
\end{verbatim}

\begin{Shaded}
\begin{Highlighting}[]
\NormalTok{dsm.emm1 }\OtherTok{\textless{}{-}} \FunctionTok{emmeans}\NormalTok{(dsm.mod.add, }\SpecialCharTok{\textasciitilde{}}\NormalTok{ trt\_Grp.f)}
\NormalTok{dsm.emm1}
\end{Highlighting}
\end{Shaded}

\begin{verbatim}
##  trt_Grp.f emmean   SE df lower.CL upper.CL
##  gel          135 3.62 56      128      142
##  gelDS        111 3.62 56      103      118
##  control      167 3.62 56      160      174
## 
## Confidence level used: 0.95
\end{verbatim}

\begin{Shaded}
\begin{Highlighting}[]
\FunctionTok{contrast}\NormalTok{(dsm.emm1, }\AttributeTok{method=}\StringTok{"pairwise"}\NormalTok{, }\AttributeTok{adjust=}\StringTok{"tukey"}\NormalTok{)}
\end{Highlighting}
\end{Shaded}

\begin{verbatim}
##  contrast        estimate   SE df t.ratio p.value
##  gel - gelDS         24.3 5.12 56   4.738 <0.0001
##  gel - control      -32.2 5.12 56  -6.301 <0.0001
##  gelDS - control    -56.5 5.12 56 -11.035 <0.0001
## 
## P value adjustment: tukey method for comparing a family of 3 estimates
\end{verbatim}

\[\nabla  \]

\section{Interaction Model}\label{interaction-model}

When the slope relating the response to the covariate differs among
treatments, the model contains interactions. This can be tested by
including cross-product terms for the treatment indicator variables and
the centered covariate. Then the full model containing these extra terms
can be compared with the additive (reduced) model. The interaction model
is given below.

\[
Y_{ij}= \mu_{\bullet} + \sum_{k=1}^{r-1}\tau_k I_{ijk}  + 
\gamma \left(X_{ij}-\overline{X}_{\bullet\bullet}\right) +
\sum_{k=1}^{r-1}\beta_kI_{ijk}\left(X_{ij}-\overline{X}_{\bullet\bullet}\right) +
\epsilon_{ijk}
\]

Comparing adjusted means is more difficult when there is an interaction,
as the difference among treatment means differs at different \(X\)
levels (unlike the additive model). A method was developed by Johnson
and Neyman to find regions of \(X\) values where two treatments are
significantly different and where they are not. We will not pursue that
method here, but will run the general linear test for interactions.

\textbf{Example 12.3 - Comparison of Skin Softeners}

We obtain the adjusted means and their differences along with multiple
comparisons with the following R code.

\begin{verbatim}
##   trt_Grp pre_x post_y gel gelDS
## 1       1 121.1   76.3   1     0
## 2       1 375.4  265.1   1     0
\end{verbatim}

\begin{verbatim}
##    trt_Grp  pre_x post_y gel gelDS
## 59       3 138.59 107.22   0     0
## 60       3 134.89 125.20   0     0
\end{verbatim}

\begin{verbatim}
## 
## Call:
## lm(formula = post_y ~ I1 + I2 + Xc, data = dsm)
## 
## Residuals:
##    Min     1Q Median     3Q    Max 
## -31.13 -10.35  -3.93  12.29  35.34 
## 
## Coefficients:
##              Estimate Std. Error t value Pr(>|t|)    
## (Intercept) 137.53667    2.08924  65.831  < 2e-16 ***
## I1           -2.66382    2.95473  -0.902    0.371    
## I2          -26.91735    2.95595  -9.106 1.22e-12 ***
## Xc            0.70732    0.03257  21.716  < 2e-16 ***
## ---
## Signif. codes:  0 '***' 0.001 '**' 0.01 '*' 0.05 '.' 0.1 ' ' 1
## 
## Residual standard error: 16.18 on 56 degrees of freedom
## Multiple R-squared:  0.9121, Adjusted R-squared:  0.9073 
## F-statistic: 193.6 on 3 and 56 DF,  p-value: < 2.2e-16
\end{verbatim}

\begin{verbatim}
## Analysis of Variance Table
## 
## Response: post_y
##           Df Sum Sq Mean Sq F value    Pr(>F)    
## I1         1   9853    9853  37.622 9.205e-08 ***
## I2         1  18744   18744  71.572 1.361e-11 ***
## Xc         1 123501  123501 471.566 < 2.2e-16 ***
## Residuals 56  14666     262                      
## ---
## Signif. codes:  0 '***' 0.001 '**' 0.01 '*' 0.05 '.' 0.1 ' ' 1
\end{verbatim}

\begin{verbatim}
## 
## Call:
## lm(formula = post_y ~ I1 + I2 + Xc + I(I1 * Xc) + I(I2 * Xc), 
##     data = dsm)
## 
## Residuals:
##     Min      1Q  Median      3Q     Max 
## -28.385  -7.771   0.046   7.394  35.485 
## 
## Coefficients:
##              Estimate Std. Error t value Pr(>|t|)    
## (Intercept) 137.83352    1.64131  83.978  < 2e-16 ***
## I1           -2.99032    2.32058  -1.289   0.2030    
## I2          -26.78674    2.32170 -11.538 3.41e-16 ***
## Xc            0.72542    0.02632  27.564  < 2e-16 ***
## I(I1 * Xc)   -0.05744    0.03431  -1.674   0.0999 .  
## I(I2 * Xc)   -0.17549    0.03784  -4.638 2.28e-05 ***
## ---
## Signif. codes:  0 '***' 0.001 '**' 0.01 '*' 0.05 '.' 0.1 ' ' 1
## 
## Residual standard error: 12.71 on 54 degrees of freedom
## Multiple R-squared:  0.9477, Adjusted R-squared:  0.9429 
## F-statistic: 195.8 on 5 and 54 DF,  p-value: < 2.2e-16
\end{verbatim}

\begin{verbatim}
## Analysis of Variance Table
## 
## Response: post_y
##            Df Sum Sq Mean Sq F value    Pr(>F)    
## I1          1   9853    9853  61.026 1.985e-10 ***
## I2          1  18744   18744 116.096 4.560e-15 ***
## Xc          1 123501  123501 764.925 < 2.2e-16 ***
## I(I1 * Xc)  1   2475    2475  15.330 0.0002553 ***
## I(I2 * Xc)  1   3472    3472  21.507 2.279e-05 ***
## Residuals  54   8719     161                      
## ---
## Signif. codes:  0 '***' 0.001 '**' 0.01 '*' 0.05 '.' 0.1 ' ' 1
\end{verbatim}

\begin{verbatim}
## Analysis of Variance Table
## 
## Model 1: post_y ~ I1 + I2 + Xc
## Model 2: post_y ~ I1 + I2 + Xc + I(I1 * Xc) + I(I2 * Xc)
##   Res.Df     RSS Df Sum of Sq      F    Pr(>F)    
## 1     56 14666.1                                  
## 2     54  8718.5  2    5947.5 18.419 7.971e-07 ***
## ---
## Signif. codes:  0 '***' 0.001 '**' 0.01 '*' 0.05 '.' 0.1 ' ' 1
\end{verbatim}

\begin{figure}

{\centering \includegraphics[width=0.8\linewidth]{bookdownproj_STA4211Notes1_files/figure-latex/ch1204-1} 

}

\caption{Analysis of Covariance for skin softener study - Interaction Model}\label{fig:ch1204}
\end{figure}

The interaction is significant \((F^*=18.42, P<.0001)\). Based on the
plot, there appears to be larger treatment differences among patients
with higher baseline skin roughness.

\[\nabla \]

\begin{verbatim}
## Warning: package 'rsm' was built under R version 4.5.3
\end{verbatim}

\begin{verbatim}
## Warning: package 'mixexp' was built under R version 4.5.3
\end{verbatim}

\begin{verbatim}
## Loading required package: lattice
\end{verbatim}

\begin{verbatim}
## Loading required package: grid
\end{verbatim}

\begin{verbatim}
## Loading required package: daewr
\end{verbatim}

\begin{verbatim}
## Warning: package 'daewr' was built under R version 4.5.3
\end{verbatim}

\chapter{Response Surface and Mixture
Designs}\label{response-surface-and-mixture-designs}

Response surface and mixture designs are used in a wide range of
engineering fields. The goal is to choose levels of a group of numeric
predictor variables that optimize the response variable(s). Response
surfaces are based on \(k\) predictors at several levels and models are
typically based on a second-order model with linear, interaction, and
quadratic terms among the predictors. Mixture designs have several
inputs, but these are restricted to sum to 1 (or 100\%).

\section{Response Surface Designs}\label{response-surface-designs}

The model for a second order response surface with \(k\) factors is
given below and can be implemented easily using the \textbf{rsm} package
in R.

\[
Y = \beta_0 + \sum_{i=1}^k \beta_{i}X_i + \sum_{i=1}^{k-1}\sum_{i'=i+1}^k\beta_{ii'}X_iX_i'
+ \sum_{i=1}^k\beta_{ii}X_i^2 + \epsilon
\] The model has \(1+k+k(k-1)/2+k\) parameters. Once the regression
model is fit, individual parameters and groups (linear terms, 2-factor
interactions, quadratic terms) can be tested for significance. The
fitted model can be written in scalar or matrix form.

\[
\hat{Y}=\hat{\beta}_0 + \sum_{i=1}^k \hat{\beta}_{i}X_i + \sum_{i=1}^{k-1}\sum_{i'=i+1}^k\hat{\beta}_{ii'}X_iX_i'
+ \sum_{i=1}^k\hat{\beta}_{ii}X_i^2 =
\hat{\beta}_0+ B_1'x + x'B_2x
\] where \(x\), \(B_1\), \(B_2\) are defined as follow. \[
x=\begin{bmatrix}
X_1 \\ \vdots \\ X_k \\
\end{bmatrix}
\qquad
B_1 =
\begin{bmatrix}
\hat{\beta}_1 \\ \vdots \\ \hat{\beta}_k\\
\end{bmatrix}
\qquad
B_2 = 
\begin{bmatrix}
\hat{\beta}_{11} & \hat{\beta}_{12}/2 & \cdots & \hat{\beta}_{1k}/2 \\
\hat{\beta}_{12}/2 & \hat{\beta}_{22} & \cdots & \hat{\beta}_{2k}/2 \\
\vdots & \vdots & \ddots & \vdots \\
\hat{\beta}_{1k}/2 & \hat{\beta}_{2k}/2 & \cdots & \hat{\beta}_{kk}\\
\end{bmatrix}
\] Taking the derivative of \(\hat{Y}\) with respect to \(x\), setting
it equal to zero, gives the optimal set of inputs for the \(k\) factors
(which may fall outside of possible values). This is printed directly in
the summary of the \textbf{rsm} fit.

\[
\frac{\partial \hat{Y}}{\partial x}= B_1 + 2B_2x \quad \stackrel{\rm set}{=}0 \quad
\Rightarrow \quad x^* = -\frac{1}{2}B_2^{-1}B_1
\]

Two commonly used designs for response surfaces are the \textbf{Central
Composite Design} and the \textbf{Box-Behnken Design}.

In the Central Composite Design, there are \(k\) factors, each with
three equally spaced levels coded \(-1,0,+1\). A \(2^k\) factorial is
set up with each factor at \(+/-1\). Then \(2k\) observations are made
with \(k-1\) variables at their 0 levels and the remaining factor at
\(\pm \alpha\) where \(\alpha\) is commonly \(\sqrt{2}\), \(\sqrt{3}\),
or 2, these are labelled ``axial points.'' Finally, there are \(c\)
observations with each of the \(k\) factors at their center (0) levels,
which permits a goodness of fit test.

\textbf{Example 13.1 - Solar Drying of Avocado Pulp}

A study was conducted as a response surface design with \(k=4\) factors
to measure solar drying of avocado pulp
(\citeproc{ref-kowarit_2024}{Kowarit, Sathpornprasath, and Jansri
2024}). The 4 factors are described below, and the response was moisture
content of the avocado pulp (lower values are better). This is a central
composite design with \(\alpha=2\).

\begin{itemize}
\tightlist
\item
  \textbf{Hot Air Temperature} (\(X_1\), Celsius): -1=40, 0=55, +1=70
  (axial points at 25, 85)
\item
  \textbf{Drying Time} (\(X_2\), hours): -1=13, 0=16.5, +1=20 (axial
  points at 9.5, 23.5)
\item
  \textbf{Raw Material Thickness} (\(X_3\), cm): -1=0.5, 0=0.75, +1=1.00
  (axial points at 0.25, 1.25)
\item
  \textbf{Wind Speed} (\(X_4\), meters/second): -1=0.12, 0=0.19, +1=0.26
  (axial points at 0.05, 0.33)
\end{itemize}

There were \(2^4=16\) runs with the 4 factors at +/-1, \(4(2)=8\) runs
at the axial points for each factor (with the other 3 factors at their 0
level), and 6 runs at the center points, for a total of 30 experimental
runs. The R program and output are given below, with response surface
and contour plots in Figure @ref(fig:ch1301) and Figure
@ref(fig:ch1301a), respectively.

\begin{Shaded}
\begin{Highlighting}[]
\NormalTok{am }\OtherTok{\textless{}{-}} \FunctionTok{read.csv}\NormalTok{(}\StringTok{"data/avocado.csv"}\NormalTok{)}
\FunctionTok{head}\NormalTok{(am); }\FunctionTok{tail}\NormalTok{(am)}
\end{Highlighting}
\end{Shaded}

\begin{verbatim}
##   run hotAirTemp dryTime rawMatThck windSpeed moisture
## 1   1         70    20.0       1.00      0.26    20.33
## 2   2         70    13.0       1.00      0.12    20.15
## 3   3         85    16.5       0.75      0.19    18.26
## 4   4         55    16.5       0.75      0.19    18.90
## 5   5         55    16.5       0.75      0.19    17.99
## 6   6         70    20.0       0.50      0.12    18.54
\end{verbatim}

\begin{verbatim}
##    run hotAirTemp dryTime rawMatThck windSpeed moisture
## 25  25         70    13.0       1.00      0.26    20.29
## 26  26         55     9.5       0.75      0.19    19.03
## 27  27         40    20.0       1.00      0.26    21.61
## 28  28         55    16.5       0.75      0.19    18.90
## 29  29         55    16.5       0.75      0.19    18.90
## 30  30         40    13.0       0.50      0.12    17.73
\end{verbatim}

\begin{Shaded}
\begin{Highlighting}[]
\NormalTok{mod1 }\OtherTok{\textless{}{-}} \FunctionTok{rsm}\NormalTok{(moisture }\SpecialCharTok{\textasciitilde{}} \FunctionTok{SO}\NormalTok{(hotAirTemp, dryTime,rawMatThck,windSpeed), }\AttributeTok{data=}\NormalTok{am)}
\FunctionTok{summary}\NormalTok{(mod1)}
\end{Highlighting}
\end{Shaded}

\begin{verbatim}
## 
## Call:
## rsm(formula = moisture ~ SO(hotAirTemp, dryTime, rawMatThck, 
##     windSpeed), data = am)
## 
##                          Estimate  Std. Error t value Pr(>|t|)   
## (Intercept)            2.1996e+01  1.1195e+01  1.9649 0.068230 . 
## hotAirTemp             2.6835e-01  1.5944e-01  1.6831 0.113053   
## dryTime               -4.8311e-01  7.3777e-01 -0.6548 0.522495   
## rawMatThck            -8.5945e+00  9.1563e+00 -0.9386 0.362780   
## windSpeed             -6.4616e+01  3.2146e+01 -2.0101 0.062761 . 
## hotAirTemp:dryTime    -6.8333e-03  5.4205e-03 -1.2607 0.226691   
## hotAirTemp:rawMatThck -5.9667e-02  7.5887e-02 -0.7863 0.443959   
## hotAirTemp:windSpeed   2.1786e-01  2.7102e-01  0.8038 0.434049   
## dryTime:rawMatThck    -1.4143e-01  3.2523e-01 -0.4349 0.669856   
## dryTime:windSpeed      1.5561e+00  1.1615e+00  1.3397 0.200282   
## rawMatThck:windSpeed  -2.4071e+01  1.6261e+01 -1.4803 0.159495   
## hotAirTemp^2          -1.2204e-03  9.6599e-04 -1.2633 0.225752   
## dryTime^2              2.7585e-02  1.7743e-02  1.5547 0.140853   
## rawMatThck^2           1.2027e+01  3.4776e+00  3.4584 0.003511 **
## windSpeed^2            1.3274e+02  4.4357e+01  2.9925 0.009110 **
## ---
## Signif. codes:  0 '***' 0.001 '**' 0.01 '*' 0.05 '.' 0.1 ' ' 1
## 
## Multiple R-squared:  0.7643, Adjusted R-squared:  0.5444 
## F-statistic: 3.475 on 14 and 15 DF,  p-value: 0.01123
## 
## Analysis of Variance Table
## 
## Response: moisture
##                                                 Df  Sum Sq Mean Sq F value
## FO(hotAirTemp, dryTime, rawMatThck, windSpeed)   4 23.1131  5.7783  4.4595
## TWI(hotAirTemp, dryTime, rawMatThck, windSpeed)  6  9.1073  1.5179  1.1715
## PQ(hotAirTemp, dryTime, rawMatThck, windSpeed)   4 30.8123  7.7031  5.9450
## Residuals                                       15 19.4359  1.2957        
## Lack of fit                                     10 18.7458  1.8746 13.5823
## Pure error                                       5  0.6901  0.1380        
##                                                   Pr(>F)
## FO(hotAirTemp, dryTime, rawMatThck, windSpeed)  0.014231
## TWI(hotAirTemp, dryTime, rawMatThck, windSpeed) 0.371796
## PQ(hotAirTemp, dryTime, rawMatThck, windSpeed)  0.004511
## Residuals                                               
## Lack of fit                                     0.005031
## Pure error                                              
## 
## Stationary point of response surface:
## hotAirTemp    dryTime rawMatThck  windSpeed 
## 67.1141331 14.0922992  0.7835529  0.1767657 
## 
## Eigenanalysis:
## eigen() decomposition
## $values
## [1] 133.931086216  10.838362960   0.023680706  -0.002003314
## 
## $vectors
##                     [,1]          [,2]          [,3]         [,4]
## hotAirTemp  0.0008310985  0.0017517386 -0.1599582827  0.987121871
## dryTime     0.0058342613 -0.0005625942  0.9871073880  0.159952022
## rawMatThck -0.0982548808 -0.9951596092 -0.0002787034  0.001803565
## windSpeed   0.9951438334 -0.0982546031 -0.0056810736 -0.001584081
\end{verbatim}

\begin{Shaded}
\begin{Highlighting}[]
\FunctionTok{par}\NormalTok{(}\AttributeTok{mfrow=}\FunctionTok{c}\NormalTok{(}\DecValTok{2}\NormalTok{,}\DecValTok{3}\NormalTok{))}
\FunctionTok{persp}\NormalTok{(mod1, }\SpecialCharTok{\textasciitilde{}}\NormalTok{ hotAirTemp}\SpecialCharTok{+}\NormalTok{dryTime}\SpecialCharTok{+}\NormalTok{rawMatThck}\SpecialCharTok{+}\NormalTok{windSpeed, }
   \AttributeTok{at=}\FunctionTok{summary}\NormalTok{(mod1)}\SpecialCharTok{$}\NormalTok{canonical}\SpecialCharTok{$}\NormalTok{xs)}
\end{Highlighting}
\end{Shaded}

\begin{figure}

{\centering \includegraphics[width=0.8\linewidth]{bookdownproj_STA4211Notes1_files/figure-latex/ch1301-1} 

}

\caption{Response Surface for avocado pulp study - 3D Plots}\label{fig:ch1301}
\end{figure}

\begin{Shaded}
\begin{Highlighting}[]
\FunctionTok{par}\NormalTok{(}\AttributeTok{mfrow=}\FunctionTok{c}\NormalTok{(}\DecValTok{1}\NormalTok{,}\DecValTok{1}\NormalTok{))}
\end{Highlighting}
\end{Shaded}

\begin{Shaded}
\begin{Highlighting}[]
\FunctionTok{par}\NormalTok{(}\AttributeTok{mfrow=}\FunctionTok{c}\NormalTok{(}\DecValTok{2}\NormalTok{,}\DecValTok{3}\NormalTok{))}
\FunctionTok{contour}\NormalTok{(mod1, }\SpecialCharTok{\textasciitilde{}}\NormalTok{ hotAirTemp}\SpecialCharTok{+}\NormalTok{dryTime}\SpecialCharTok{+}\NormalTok{rawMatThck}\SpecialCharTok{+}\NormalTok{windSpeed, }
   \AttributeTok{at=}\FunctionTok{summary}\NormalTok{(mod1)}\SpecialCharTok{$}\NormalTok{canonical}\SpecialCharTok{$}\NormalTok{xs)}
\end{Highlighting}
\end{Shaded}

\begin{figure}

{\centering \includegraphics[width=0.8\linewidth]{bookdownproj_STA4211Notes1_files/figure-latex/ch1301a-1} 

}

\caption{Response Surface for avocado pulp study - Contour Plots}\label{fig:ch1301a}
\end{figure}

\begin{Shaded}
\begin{Highlighting}[]
\FunctionTok{par}\NormalTok{(}\AttributeTok{mfrow=}\FunctionTok{c}\NormalTok{(}\DecValTok{1}\NormalTok{,}\DecValTok{1}\NormalTok{))}
\end{Highlighting}
\end{Shaded}

The linear (first order) terms and polynomial quadratic terms are
significant as groups with \(P\)-values of .0142 and .0045,
respectively. The two-factor interactions are not, with \(P\)=.3718. The
stationary (optimal) point is \(x^*\)=(67.11,14.09,0.78,0.18).

\[\nabla \]

The Box-Behnken design places runs at every combination of \(\pm 1\) for
each pair of factors, with all other factors at their 0 (central) level.
There are multiple runs for each factor at their central level. The
analysis is the same as for the Central Composite Design.

\textbf{Example 13.2 -- Gels of Diclofenac and Curcumin for Transdermal
Drug Delivery}

A study (Chaudhary, et al (2011)) had 3 factors, each at 3 levels
(\citeproc{ref-chaudhary_2011}{Chaudhary et al. 2011}). The factors were
as follow.

\begin{itemize}
\tightlist
\item
  \textbf{Polymer Concentration} - (0.5, 1.0, 1.5 \% w/w)
\item
  \textbf{Ethanol Concentration} (10, 15, 20 \% w/w)
\item
  \textbf{PG Concentration} (5, 10, 15 \%w/w).
\end{itemize}

There were three response variables: \(Y1\) = Flux of DDEA (mg/(cm2 h)),
\(Y2\) = Flux of CRM (mg/(cm2 h)), and \(Y3\) = Viscosity of Gel (cp).
The experiment had \(n\) = 17 runs. The design and data are given below.

\begin{longtable}[]{@{}rrrrrrr@{}}
\caption{Box-Behnken Design for Drug Delivery Study}\tabularnewline
\toprule\noalign{}
runnum & polyconc & ethnconc & pgconc & fluxddea & fluxcrm & viscgel \\
\midrule\noalign{}
\endfirsthead
\toprule\noalign{}
runnum & polyconc & ethnconc & pgconc & fluxddea & fluxcrm & viscgel \\
\midrule\noalign{}
\endhead
\bottomrule\noalign{}
\endlastfoot
1 & -1 & -1 & 0 & 0.67 & 1.48 & 185 \\
2 & 0 & 0 & 0 & 0.24 & 1.90 & 1924 \\
3 & 0 & 1 & -1 & 0.25 & 3.31 & 2018 \\
4 & 0 & -1 & 1 & 0.22 & 2.88 & 2310 \\
5 & 1 & 0 & -1 & 0.11 & 3.30 & 3227 \\
6 & -1 & 1 & 0 & 0.67 & 1.72 & 145 \\
7 & -1 & 0 & -1 & 0.69 & 1.37 & 143 \\
8 & 0 & 0 & 0 & 0.23 & 1.87 & 1923 \\
9 & -1 & 0 & 1 & 0.67 & 1.52 & 176 \\
10 & 0 & 0 & 0 & 0.24 & 1.91 & 1800 \\
11 & 0 & 0 & 0 & 0.24 & 1.92 & 1921 \\
12 & 1 & 0 & 1 & 0.17 & 2.98 & 3071 \\
13 & 0 & 1 & 1 & 0.23 & 2.01 & 1783 \\
14 & 0 & 0 & 0 & 0.24 & 1.87 & 1922 \\
15 & 1 & -1 & 0 & 0.11 & 3.07 & 3320 \\
16 & 1 & 1 & 0 & 0.16 & 3.45 & 2801 \\
17 & 0 & -1 & -1 & 0.20 & 1.71 & 2245 \\
\end{longtable}

\[\nabla\]

\section{Mixture Models}\label{mixture-models}

Mixture models are similar to response surface models, with the
restriction that the sum of the \(X\) levels is equal to 1, that is the
\(X\) variables are components of a mixture. There are four widely used
models that can be implemented in the \textbf{mixexp} package in R. Note
that these models do not include an intercept, as the X's sum to 1, and
we will use as \(k\) as the number of mixture components.

\[
\mbox{Linear:} \quad E\{Y\}=\sum_{i=1}^k \beta_iX_i 
\] \[
\mbox{Quadratic:} \quad  E\{Y\}=\sum_{i=1}^k \beta_iX_i + 
\sum_{i=1}^{k-1}\sum_{i'=i+1}^k \beta_{ii'}X_iX_{i'}
\] \[
\mbox{Full Cubic:} \quad
E\{Y\}=\sum_{i=1}^k \beta_iX_i + 
\sum_{i=1}^{k-1}\sum_{i'=i+1}^k \beta_{ii'}X_iX_{i'} 
+ \sum_{i=1}^{k-1}\sum_{i'=i+1}^k\delta_{ii'}X_iX_{i'}\left(X_i-X_{i'}\right) + \]

\[ + \sum_{i=1}^{k-2}\sum_{i'=i+1}^{k-1}\sum_{i''=i'+1}^k\beta_{ii'i''}X_iX_{i'}X_{i''}
\] \[
\mbox{Special Cubic:} \quad
E\{Y\}=\sum_{i=1}^k \beta_iX_i + 
\sum_{i=1}^{k-1}\sum_{i'=i+1}^k \beta_{ii'}X_iX_{i'} +
\sum_{i=1}^{k-2}\sum_{i'=i+1}^{k-1}\sum_{i''=i'+1}^k\beta_{ii'i''}X_iX_{i'}X_{i''}
\] The goal is to choose the mixture that optimizes the output.

\textbf{Example 13.3 - Breaking Strength of Lipstick Blends}

A mixture experiment considered \(k=3\) inputs for lipstick
(\citeproc{ref-hui_2017}{Hui, Tamburic, and Grant-Poss 2017}). One
response was breaking strength (grams), with the following inputs, this
was Stage 1C in the paper, and the Full Cubic model was fit.

\begin{itemize}
\tightlist
\item
  \textbf{Sweet Almond Oil} \((X_1)\)
\item
  \textbf{Hydrogenated Polyisobutene} \((X_2)\)
\item
  \textbf{Octydodecanol} \((X_3)\)
\end{itemize}

The R code and output are given below.

\begin{Shaded}
\begin{Highlighting}[]
\NormalTok{lsa }\OtherTok{\textless{}{-}} \FunctionTok{read.csv}\NormalTok{(}\StringTok{"data/lipstick\_sweetalmond.csv"}\NormalTok{)}
\FunctionTok{attach}\NormalTok{(lsa)}
\NormalTok{lsa}
\end{Highlighting}
\end{Shaded}

\begin{verbatim}
##    ExpNum      X1      X2      X3 Break Soft
## 1       1 1.00000 0.00000 0.00000 325.8 62.0
## 2       2 1.00000 0.00000 0.00000 239.2 53.0
## 3       3 0.00000 1.00000 0.00000 332.9 61.3
## 4       4 0.00000 0.00000 1.00000 242.1 58.8
## 5       5 0.66667 0.33333 0.00000 247.3 53.3
## 6       6 0.33333 0.66667 0.00000 258.9 53.3
## 7       7 0.66667 0.00000 0.33333 342.1 61.5
## 8       8 0.66667 0.00000 0.33333 343.0 57.8
## 9       9 0.33333 0.33333 0.33333 292.5 53.8
## 10     10 0.00000 0.66667 0.33333 299.9 52.8
## 11     11 0.33333 0.00000 0.66667 399.2 61.3
## 12     12 0.00000 0.33333 0.66667 407.6 58.0
## 13     13 0.66667 0.16667 0.16667 258.9 53.3
## 14     14 0.16667 0.66667 0.16667 326.0 55.0
## 15     15 0.16667 0.16667 0.66667 325.1 53.8
## 16     16 0.16667 0.16667 0.66667 436.8 59.0
\end{verbatim}

\begin{Shaded}
\begin{Highlighting}[]
\NormalTok{mixvars }\OtherTok{\textless{}{-}} \FunctionTok{c}\NormalTok{(}\StringTok{"X1"}\NormalTok{,}\StringTok{"X2"}\NormalTok{,}\StringTok{"X3"}\NormalTok{) }\DocumentationTok{\#\#\# Creates group of mixvars}
\NormalTok{ls.mod1 }\OtherTok{\textless{}{-}} \FunctionTok{MixModel}\NormalTok{(lsa,}\StringTok{"Break"}\NormalTok{,mixvars,}\DecValTok{3}\NormalTok{)}
\end{Highlighting}
\end{Shaded}

\begin{verbatim}
##      
##               coefficients    Std.err    t.value         Prob
## X1                283.4746   33.32269  8.5069535 0.0001444655
## X2                331.9844   47.06751  7.0533671 0.0004063754
## X3                244.3524   47.08335  5.1897831 0.0020349071
## cubic(X1, X2)    -201.1627  393.63036 -0.5110446 0.6275740863
## cubic(X1, X3)    -437.2026  349.65399 -1.2503865 0.2577122053
## cubic(X2, X3)    -667.6002  388.25675 -1.7194813 0.1363220821
## X2:X1            -222.8240  196.31796 -1.1350157 0.2996670546
## X3:X1             464.4747  180.58235  2.5720936 0.0422159307
## X2:X3             314.0409  210.14454  1.4944045 0.1856908476
## X2:X3:X1        -1152.8906 1320.00853 -0.8733963 0.4160293996
##      
## Residual standard error:  47.18515  on  6 degrees of freedom 
## Corrected Multiple R-squared:  0.7582724
\end{verbatim}

\begin{Shaded}
\begin{Highlighting}[]
\DocumentationTok{\#\#\# Fits model, must specify: dataframe,"Y",mixvars,}
\DocumentationTok{\#\#\# model (1=linear, 2=quadratic, 3=full cubic, 4=special cubic)}
\FunctionTok{anova}\NormalTok{(ls.mod1)}
\end{Highlighting}
\end{Shaded}

\begin{verbatim}
## Analysis of Variance Table
## 
## Response: Break
##               Df Sum Sq Mean Sq  F value    Pr(>F)    
## X1             1 822487  822487 369.4245 1.283e-06 ***
## X2             1 418251  418251 187.8594 9.375e-06 ***
## X3             1 380430  380430 170.8724 1.236e-05 ***
## cubic(X1, X2)  1    802     802   0.3604   0.57023    
## cubic(X1, X3)  1   1624    1624   0.7293   0.42589    
## cubic(X2, X3)  1   5300    5300   2.3807   0.17378    
## X1:X2          1   7691    7691   3.4546   0.11243    
## X1:X3          1  11369   11369   5.1064   0.06457 .  
## X2:X3          1   3437    3437   1.5437   0.26042    
## X1:X2:X3       1   1698    1698   0.7629   0.41602    
## Residuals      6  13358    2226                       
## ---
## Signif. codes:  0 '***' 0.001 '**' 0.01 '*' 0.05 '.' 0.1 ' ' 1
\end{verbatim}

\begin{Shaded}
\begin{Highlighting}[]
\FunctionTok{summary}\NormalTok{(ls.mod1)}
\end{Highlighting}
\end{Shaded}

\begin{verbatim}
## 
## Call:
## lm(formula = mixmodnI, data = frame)
## 
## Residuals:
##     Min      1Q  Median      3Q     Max 
## -60.722 -15.433   0.876  12.755  50.978 
## 
## Coefficients:
##               Estimate Std. Error t value Pr(>|t|)    
## X1              283.47      33.32   8.507 0.000144 ***
## X2              331.98      47.07   7.053 0.000406 ***
## X3              244.35      47.08   5.190 0.002035 ** 
## cubic(X1, X2)  -201.16     393.63  -0.511 0.627570    
## cubic(X1, X3)  -437.20     349.65  -1.250 0.257712    
## cubic(X2, X3)  -667.60     388.25  -1.719 0.136321    
## X1:X2          -222.83     196.32  -1.135 0.299661    
## X1:X3           464.47     180.58   2.572 0.042215 *  
## X2:X3           314.04     210.14   1.494 0.185689    
## X1:X2:X3      -1152.90    1320.00  -0.873 0.416021    
## ---
## Signif. codes:  0 '***' 0.001 '**' 0.01 '*' 0.05 '.' 0.1 ' ' 1
## 
## Residual standard error: 47.18 on 6 degrees of freedom
## Multiple R-squared:  0.992,  Adjusted R-squared:  0.9786 
## F-statistic: 74.25 on 10 and 6 DF,  p-value: 1.76e-05
\end{verbatim}

\begin{Shaded}
\begin{Highlighting}[]
\FunctionTok{MixturePlot}\NormalTok{(X3,X2,X1,Break,}\AttributeTok{x3lab=}\StringTok{"Octydodecanol"}\NormalTok{, }\AttributeTok{x2lab=}\StringTok{"Hydrogenated Polyisobutene"}\NormalTok{,}
\AttributeTok{x1lab=}\StringTok{"Sweet Almond Oil"}\NormalTok{, }\AttributeTok{corner.labs=}\FunctionTok{c}\NormalTok{(}\StringTok{"Oct"}\NormalTok{,}\StringTok{"HP"}\NormalTok{,}\StringTok{"SAO"}\NormalTok{),}
\AttributeTok{constrts=}\ConstantTok{FALSE}\NormalTok{,}\AttributeTok{contrs=}\ConstantTok{TRUE}\NormalTok{,}\AttributeTok{cols=}\ConstantTok{TRUE}\NormalTok{, }\AttributeTok{mod=}\DecValTok{2}\NormalTok{,}\AttributeTok{n.breaks=}\DecValTok{9}\NormalTok{)}
\end{Highlighting}
\end{Shaded}

\begin{figure}

{\centering \includegraphics[width=0.8\linewidth]{bookdownproj_STA4211Notes1_files/figure-latex/ch1303-1} 

}

\caption{Mixture Plot for lipstick study}\label{fig:ch1303}
\end{figure}

\begin{Shaded}
\begin{Highlighting}[]
\FunctionTok{cbind}\NormalTok{(X1,X2,X3,Break, }\FunctionTok{predict}\NormalTok{(ls.mod1))}
\end{Highlighting}
\end{Shaded}

\begin{verbatim}
##         X1      X2      X3 Break         
## 1  1.00000 0.00000 0.00000 325.8 283.4746
## 2  1.00000 0.00000 0.00000 239.2 283.4746
## 3  0.00000 1.00000 0.00000 332.9 331.9841
## 4  0.00000 0.00000 1.00000 242.1 244.3521
## 5  0.66667 0.33333 0.00000 247.3 235.2266
## 6  0.33333 0.66667 0.00000 258.9 281.1992
## 7  0.66667 0.00000 0.33333 342.1 341.2641
## 8  0.66667 0.00000 0.33333 343.0 341.2641
## 9  0.33333 0.33333 0.33333 292.5 305.6438
## 10 0.00000 0.66667 0.33333 299.9 323.1077
## 11 0.33333 0.00000 0.66667 399.2 392.9941
## 12 0.00000 0.33333 0.66667 407.6 392.8011
## 13 0.66667 0.16667 0.16667 258.9 263.7996
## 14 0.16667 0.66667 0.16667 326.0 285.0699
## 15 0.16667 0.16667 0.66667 325.1 385.8225
## 16 0.16667 0.16667 0.66667 436.8 385.8225
\end{verbatim}

The highest breaking strengths \((>360)\) appear at blends of
approximately 50/50\% between Sweet Almond Oil and Octydodecanol, based
on the plot of the simplex.

\[\nabla \]

\phantomsection\label{refs}
\begin{CSLReferences}{1}{0}
\bibitem[\citeproctext]{ref-aksu_2007}
Aksu, M. I., H. Imik, and M. Karoglu. 2007. {``Influence of Dietary
Sorghum (Sorghum Vulgare) and Corn Supplemented with Methionine on
Cut-up Pieces Weights of Broiler Carcass and Quality Properties of
Breast and Drumsticks Meat.''} \emph{Food Science and Technology
International} 13 (5): 361--67.

\bibitem[\citeproctext]{ref-beata_2007}
Beata, C., J. Cordel, and N. Marlois. 2007. {``Effect of
Alpha-Casozepine (Zylkene) on Cats with Anxiety.''} \emph{Journal of
Veterinary Behavior: Clinical Applications and Research} 2 (2): 40--46.

\bibitem[\citeproctext]{ref-bishop_1996}
Bishop, D., C. E. Frazier, L. Lanza-Kaduce, and L. Winner. 1996. {``The
Transfer of Juveniles to Criminal Court: Does It Make a Difference?''}
\emph{Crime \& Delinquency} 42: 171--91.

\bibitem[\citeproctext]{ref-bronet_2015}
Bronet, F., M-C. Nogales, E. Martinez, M. Ariza, C. Rubio, J-A.
Garcia-Velasco, and M. Meseguer. 2015. {``Is There a Relationship
Between Time-Lapse Parameters and Embryo Sex?''} \emph{Fertility and
Sterility} 103 (2): 396--401.

\bibitem[\citeproctext]{ref-carr_2012}
Carr, D. J., C. Lankester, A. Peare, N. Fabri, and N. Gridley. 2012.
{``Does Quilting Improve the Fragment Protective Performance of Body
Armour?''} \emph{Textile Research Journal} 82 (9): 883--88.

\bibitem[\citeproctext]{ref-carraro_2019}
Carraro, A., C. A. Elliott, and E. Gobbi. 2019. {``Perceived Treadmill
Function Is Correlated with Enjoyment of Use in Trained Runners: A
User-Centered Approach.''} \emph{Applied Ergonomics} 74: 37--40.

\bibitem[\citeproctext]{ref-chang_2013}
Chang, P.-C., S.-Y. Chou, and K.-K. Shieh. 2013. {``Reading Performance
and Visual Fatigue When Using Electronic Displays in Long-Duration
Reading Tasks Under Various Lighting Conditions.''} \emph{Displays} 34:
208--14.

\bibitem[\citeproctext]{ref-chaudhary_2011}
Chaudhary, H., K. Kohli, S. Amin, P. Rathee, and V. Kumar. 2011.
{``Optimization and Formulation Design of Gels of Diclfenac and Curcumin
for Transdermal Drug Delivery by Box-Behnken Statistical Design.''}
\emph{Journal of Pharmaceutical Sciences} 100 (2): 580--93.

\bibitem[\citeproctext]{ref-cook_1989}
Cook, S. R., and G. Proulx. 1989. {``Mechanical Evaluation and
Performance Improvement of the Rotating Jaw Conibear 120 Trap.''}
\emph{Journal of Testing and Evaluation} 17 (3): 190--95.

\bibitem[\citeproctext]{ref-devarajan_2017}
Deverajan, S., R. H. McQueen, and S. Wen. 2017. {``Can Common Finishing
Treatments Used in Chef Jacket Fabrics Improve Protection Against Scald
Injury?''} \emph{Fashion \& Textiles} 4 (19): 1:15.
\url{https://doi.org/10.1186/s40691-017-0103-3}.

\bibitem[\citeproctext]{ref-fidaleo_2012}
Fidaleo, M., M. Moresi, A. Cammaroto, M. Ladrange, and R. Nardi. 2012.
{``Soy Sauce Desalting by Electrodialysis.''} \emph{Journal of Food
Engineering} 110: 175--81.

\bibitem[\citeproctext]{ref-garcia_1976}
Garcia, M. L., and J. Diaz. 1976. {``Combability Measurements on Human
Hair.''} \emph{Journal of the Society of Cosmetic Chemists} 27: 379--98.

\bibitem[\citeproctext]{ref-genova_2012}
Genova, G., P. Iacopini, M. Baldi, A. Ranieri, P. Storchi, and L.
Sebastiani. 2012. {``Temperature and Storage Effects on Antioxidant
Activity of Juice from Red and White Grapes.''} \emph{International
Journal of Food Science \& Technology} 47: 13--23.

\bibitem[\citeproctext]{ref-grand_1997}
Grand, A. N., and L. N. Bell. 1997. {``Caffeine Content of Fountain and
Private Label Store Brand Carbonated Beverages.''} \emph{Journal of the
American Dietetic Association} 97 (2): 179--82.

\bibitem[\citeproctext]{ref-gullberg_2008}
Gullberg, R. G. 2008. {``Employing Components-of-Variance to Evaluate
Forensic Breath Test Instruments.''} \emph{Science and Justice} 48:
2--7.

\bibitem[\citeproctext]{ref-hanihara_1999}
Hanihara, T., Y. Dodo, O. Kondo, T. Nara, N. Doi, and N. Sensui. 1999.
{``Intra- and Inter-Observer Errors in Facial Flatness Measurements.''}
\emph{Anthropological Science} 107 (1): 25--39.

\bibitem[\citeproctext]{ref-hinneh_2020}
Hinneh, M., E. E. Abotsi, D. Van de Walle, D. A. Tzompa-Sosa, A. De
Winne, J. Simonis, J. Messens K. anf Van Durme, E. O. Afoakwa, L. De
Cooman, and K. Dewettink. 2020. {``Pod Storage with Roasting: A Tool to
Diversifying the Flavor Profiles of Dark Chocolates Produced from
???bulk??? Cocoa Beans? (Part II: Quality and Sensory Profiling of
Chocolates).''} \emph{Food Research International} 132 (109116).

\bibitem[\citeproctext]{ref-hsu_1991}
Hsu, S.-H., and S.-P. Wu. 1991. {``An Investigation for Determining the
Optimal Length of Chopsticks.''} \emph{Applied Ergonomics} 22 (6):
395--400.

\bibitem[\citeproctext]{ref-hui_2017}
Hui, W. N., S. Tamburic, and P. Grant-Poss. 2017. {``Lip-Smacking
Results.''} \emph{Cosmetics \& Toiletries} 132 (3): 48--68.

\bibitem[\citeproctext]{ref-hyllegard_2009}
Hyllegard, K., J. Ogle, and R.-N. Yan. 2009. {``The Impact of
Advertising Message Strategy - Fair Labour v Sex Appeal - Upon Gen y
Consumers' Intent to Patronize an Apparel Retailer.''} \emph{Journal of
Fashion Marketing and Management} 13 (1): 109--27.

\bibitem[\citeproctext]{ref-ismail_2014}
Ismail, S., S. Siddiqui, F. Shafiq, M. Ishaq, and S. Khan. 2014.
{``Blood Transfusion in Patients Having Caesarean Section: A Prospective
Multicentre Observational Study of Practice in Three Pakistan
Hospitals.''} \emph{International Journal of Obstetric Anesthesia} 23:
253--59.

\bibitem[\citeproctext]{ref-jensen_2002}
Jensen, C. R. 2002. {``Variance Components Calculations: Common Methods
and Misapplications in the Semiconductor Industry.''} \emph{Quality
Engineering} 14 (4): 647--57.

\bibitem[\citeproctext]{ref-jung_2018}
Jung, J., and Y. J. Ahn. 2018. {``Effects of Interface on Procedural
Skill Transfer in Virtual Training: Lifeboat Launching Operation
Study.''} \emph{Computer Animation and Virtual Worlds} 29: e1812.
\url{https://doi.org/10.1002/cav.1812}.

\bibitem[\citeproctext]{ref-keshtkaran_2014}
Keshtkaran, Z., F. Sharif, and M. Rambod. 2014. {``Students' Readiness
for and Perception of Inter-Professional Learning: A Cross-Sectional
Study.''} \emph{Nurse Education Today} 34: 991--98.

\bibitem[\citeproctext]{ref-kowarit_2024}
Kowarit, S., K. Sathpornprasath, and S. N. Jansri. 2024. {``Application
of Hot Air-Derived RSM Conditions and Shading for Solar Drying of
Avocado Pulp and Its Properties.''} \emph{Solar Energy} 278: 1--12.
\url{https://doi.org/10.1016/j.solener.2024.112768}.

\bibitem[\citeproctext]{ref-landy_1974}
Landy, D., and H. Sigall. 1974. {``Beauty Is Talent: Task Evaluation as
a Function of the Performer's Physical Attraction.''} \emph{Journal of
Personality and Social Psychology} 29 (3): 299--304.

\bibitem[\citeproctext]{ref-maor_1997}
Ma'Or, Z., S. Yehuda, and W. Voss. 1997. {``Skin Smoothing Effects of
Dead Sea Minerals: Comparative Profilometric Evaluation of Skin
Surface.''} \emph{International Journal of Cosmetic Science} 19:
105--10.

\bibitem[\citeproctext]{ref-marchiori_2012}
Marchiori, D., O. Cornielle, and O. Klein. 2012. {``Container Size
Influences Snack Food Intake Independently of Portion Size.''}
\emph{Appetite} 58: 814--17.

\bibitem[\citeproctext]{ref-meilgaard_1960}
Meilgaard, M. 1960. {``Hop Analysis, Cohumulone Factor and the
Bitterness of Beer: Review and Critical Evaluation.''} \emph{Journal of
the Institute of Brewing} 66 (1): 35--50.

\bibitem[\citeproctext]{ref-ohare_2009}
O'Hare, D., and N. Stenhouse. 2009. {``Under the Weather: An Evaluation
of Different Modes of Presenting Meteorological Information for
Pilots.''} \emph{Applied Ergonomics} 40: 688--93.

\bibitem[\citeproctext]{ref-okeeffe_2005}
O'Keeffe, C., and R. Wiseman. 2005. {``Testing Alleged Mediumship:
Methods and Results.''} \emph{British Journal of Psychology} 96:
165--79.

\bibitem[\citeproctext]{ref-oliveira_2017}
Oliveira, A. S., F. M. Dalla Nora, R. O. Mello, P. A. Mello, B. Tischer,
A. B. Costa, and J. S. Barin. 2017. {``One-Shot, Reagent-Free
Determination of the Alcohol Content of Distilled Beverages by Thermal
Infrared Enthalpimetry.''} \emph{Talanta} 171: 335--40.
\url{https://doi.org/10.1016/j.talanta.2017.05.011}.

\bibitem[\citeproctext]{ref-ouedrhiria_2016}
Ouedrhiria, W., M. Balouiri, S. Bouhdid, M., S. Moja, F. O. Chahdie, M.
Taleb, and H. Greche. 2016. {``Mixture Designs of Origanum Compactum,
Origanum Majorana and Thymus Serpyllum Essential Oils: Optimization of
Their Antibacterial Effect.''} \emph{Industrial Crops and Products} 89:
1--9.

\bibitem[\citeproctext]{ref-ranthinam_2011}
Ranthinam, A., J. R. Rao, and B. U. Nair. 2011. {``Adsorption of Phenol
onto Activated Carbon from Seaweed: Determination of the Optimal
Experimental Parameters Using Factorial Design.''} \emph{Journal of the
Taiwan Institute of Chemical Engineers} 42: 952--56.

\bibitem[\citeproctext]{ref-schultz_1952}
Schultz, A. M., and H. H. Biswell. 1952. {``Competition Between Grasses
Reseeded on Burned Brushlands in California.''} \emph{Journal of Range
Management} 5 (5): 338--45.

\bibitem[\citeproctext]{ref-simonson_2007}
Simonson, D. K. 2007. {``Creative Life Cycles in Literature: Poets Vs
Novelists or Conceptualists Versus Experimentalists?''} \emph{Psychology
of Aesthetics, Creativity, and the Arts} 1 (3): 133--39.

\bibitem[\citeproctext]{ref-smith_1969}
Smith, K. V. 1969. {``Stock Price and Economic Indexes for Generating
Efficient Portfolios.''} \emph{The Journal of Business} 42 (3): 332--36.

\bibitem[\citeproctext]{ref-smith_2010}
Smith, P. K., T. Niiler, and P. W. McCullough. 2010. {``Evaluating
Makiwara Punching Board Performance.''} \emph{Journal of Asian Martial
Arts} 19 (2): 34--45.

\bibitem[\citeproctext]{ref-snedecor_1967}
Snedecor, George W., and William G. Cochran. 1967. \emph{Statistical
Methods}. 6th ed. Ames, Iowa: The Iowa State University Press.

\bibitem[\citeproctext]{ref-song_2007}
Song, X-j., M. Zhang, and A. S. Mujumdar. 2007. {``Optimization of
Vacuum Microwave Predrying and Vacuum Frying Conditions to Produce Fried
Potato Chips.''} \emph{Drying Technology} 25: 2027--34.

\bibitem[\citeproctext]{ref-suaria_2014}
Suaria, G., and S. Aliani. 2014. {``Floating Debris in the Mediterranean
Sea.''} \emph{Marine Pollution Bulletin} 86: 494--504.

\bibitem[\citeproctext]{ref-tan_2009}
Tan, K. S., S. V. Wong, R. S. Radin Umar, A. M. S. Hamouda, and N. K.
Gupta. 2009. {``Impact Behavior Modeling of Motorcycle Front Wheel-Tire
Assembly.''} \emph{International Journal of Automotive Technology} 10
(3): 329--39.

\bibitem[\citeproctext]{ref-to_2007}
To, W. M., T. W. Yu, T. M. Lai, and S. P. Li. 2007. {``Characterization
of Commercial Clothes Dryers Based on Energy-Efficiency Analysis.''}
\emph{International Journal of Clothing Science and Technology} 19 (5):
270--90.

\bibitem[\citeproctext]{ref-wohlgemuth_2024}
Wohlgemuth, K. J., A. Jesko, V. Frost, M. J. Conner, and J. A. Mota.
2024. {``Characteristics of Slow and Fast Performers on a Firefighter
Air Consumption Test.''} \emph{Applied Ergonomics} 118.
\url{https://doi.org/10.1016/j.apergo.2024.104262}.

\bibitem[\citeproctext]{ref-wu_2011}
Wu, F-G., H. Lin, and M. You. 2011. {``The Enhanced Navigator for the
Touch Screen: A Comparative Study on Navigational Techniques of Web
Maps.''} \emph{Displays} 32: 284--95.

\bibitem[\citeproctext]{ref-zeng_2015}
Zeng, R., J. Lin, S. Wu, L. Chen, S. Chen, H. Gao, Y. Zheng, and H. Ma.
2015. {``A Randomized Controlled Trial: Preoperative Home-Based Combined
Tai Chi and Strength Training (TCST) to Improve Balance and Aerobic
Capacity in Patients with Total Hip Arthroplasty (THA).''}
\emph{Archives of Gerontology and Geriatrics} 60: 265--71.

\bibitem[\citeproctext]{ref-zouid_2013}
Zouid, I., R. Siret, F. Jourjon, E. Mehinagic, and L. Rolle. 2013.
{``Impact of Grapes Heterogeneity According to Sugar Level on Both
Physical and Mechanical Berries Properties and Their Anthocyanins
Extractability at Harvest.''} \emph{Journal of Texture Studies} 44:
95--103.

\end{CSLReferences}

\backmatter
\end{document}
